\documentclass[10pt]{amsart}
\usepackage{amsmath,amssymb,amsthm,mathtools,mathrsfs}
\usepackage[a4paper,hmargin=0.76in,vmargin=0.60in]{geometry}
\usepackage{booktabs}
\usepackage{array}
\usepackage{longtable}
\usepackage{tabularx}
\usepackage[hypertexnames=false]{hyperref}
\hypersetup{hidelinks}
\usepackage[final,expansion=false]{microtype}
\usepackage{url}
\emergencystretch=2.5em

\newtheorem{theorem}{Theorem}[section]
\newtheorem{proposition}[theorem]{Proposition}
\newtheorem{lemma}[theorem]{Lemma}
\newtheorem{corollary}[theorem]{Corollary}
\theoremstyle{definition}
\newtheorem{definition}[theorem]{Definition}
\newtheorem{problem}[theorem]{Problem}
\theoremstyle{remark}
\newtheorem{remark}[theorem]{Remark}

\newcommand{\C}{\mathbb C}
\newcommand{\R}{\mathbb R}
\newcommand{\N}{\mathbb N}
\newcommand{\Z}{\mathbb Z}
\newcommand{\Q}{\mathbb Q}
\newcommand{\Rel}{\operatorname{Re}}
\newcommand{\Iml}{\operatorname{Im}}
\newcommand{\supp}{\operatorname{supp}}
\newcommand{\TV}{\operatorname{TV}}
\newcommand{\Kp}{\mathcal K^{\mathrm{prim}}_{2T}}
\newcommand{\Kw}{\mathcal K_{2T}}
\newcommand{\PT}{\mathcal P_T}
\newcommand{\IT}{I_T}
\newcommand{\lamP}{\lambda^{\mathrm{prim}}_T}
\newcommand{\lamW}{\lambda^{\mathrm W}_T}
\newcommand{\lamA}{\lambda_{\mathrm A}(T)}
\newcommand{\sT}{\sigma_T}
\newcommand{\gh}[1]{\widehat g_{#1}}
\newcommand{\nuh}{\widehat\nu}
\newcommand{\wh}{\widehat w}
\newcommand{\Fh}{\widehat F}
\newcommand{\cF}{c_F}
\newcommand{\Lev}{\mathfrak E}
\newcommand{\Wcl}{\mathcal W_T}
\newcommand{\NT}{\mathcal N_T}
\newcommand{\Aconst}{\mathsf A_T}
\newcommand{\Ainf}{\mathcal A_\infty}
\newcommand{\PiT}{\Pi_T}
\newcommand{\mT}{\mathfrak m_T}
\newcommand{\laminf}{\lambda^{\infty}_{\min}}
\newcommand{\laminfP}{\lambda^{\infty,\mathrm P}_{\min}}
\newcolumntype{Y}{>{\raggedright\arraybackslash}X}

\providecommand{\anCP}{C_P}
\providecommand{\anPa}{\Phi_a}
\providecommand{\anPb}{\Phi_b}
\providecommand{\anPl}{\Phi}
\providecommand{\ancl}{\mathrm{cl}}
\providecommand{\aneu}{\mathrm{eu}_d}
\providecommand{\anrd}{\mathrm{rd}}
\providecommand{\auF}{\mathbb F}
\providecommand{\auGL}{\mathrm{GL}}
\providecommand{\auPb}{\Phi_b}
\providecommand{\aubar}[1]{\overline{#1}}
\providecommand{\aubud}{\mathrm{bud}}
\providecommand{\pgG}{G_E}
\providecommand{\pgNch}{N_{\mathrm{ch}}}
\providecommand{\pgS}{S_{\mathrm{bad}}}
\providecommand{\pgarec}{\operatorname{rec}}
\providecommand{\pgcodd}{c_{\mathrm{odd}}}
\providecommand{\pgcyc}{\varepsilon_\ell}
\providecommand{\pgloc}{\mathrm{loc}}
\providecommand{\pgtr}{\operatorname{tr}}
\providecommand{\pgv}{v}
\providecommand{\pllB}{B_{\mathrm{cris}}}
\providecommand{\pllD}{\mathrm{dB}}
\providecommand{\pllperp}{\perp}
\providecommand{\prics}{\mathrm{cs}}
\providecommand{\prirl}{\operatorname{rl}}
\newenvironment{quotedthm}[2]{\par\medskip\noindent\begingroup
  \textbf{#1}#2.\ \itshape\ignorespaces}{\endgroup\par\medskip}
\newenvironment{quoteddef}[2]{\par\medskip\noindent\begingroup
  \textbf{#1}#2.\ \ignorespaces}{\endgroup\par\medskip}
\usepackage{xurl}
\begin{document}
\addtolength{\abovedisplayskip}{-3.4pt}
\addtolength{\belowdisplayskip}{-3.4pt}
\addtolength{\belowdisplayshortskip}{-2pt}
\title[Explicit cup exponents]{Explicit cup exponents for local Galois representations beyond the Fontaine--Laffaille range}
\author{David Alejandro Trejo Pizzo}
\date{}
\begin{abstract}
Let $\Gamma$ be the absolute Galois group of a finite extension of $\mathbb Q_\ell$. For two
lattices in crystalline representations of $\Gamma$ that are in duality only modulo a power of
$\ell$, the cup product of two crystalline classes need not vanish. We bound the power of $\ell$
that kills all such cup products, for every prime $\ell$ and with no restriction of
Fontaine--Laffaille type, when the Hodge--Tate weights lie in a low window around $0$ and a far
window, possibly much larger than $\ell$. The bound is explicit and is sensitive to the number of
low weights rather than to the size of the largest weight; for large $\ell$ it is zero. We also
treat potentially crystalline lattices of tame type. The proofs use lattices in
$W(\mathbb C_\ell^\flat)\otimes L$ defined by orders of vanishing along the Frobenius translates of
an Eisenstein element, and a lemma that converts chains of moves of points of these lattices into
exponents for cup products.
\end{abstract}
\maketitle
\setcounter{tocdepth}{1}\tableofcontents

\section{Introduction}\label{C4:intro}

Let $\ell$ be a prime, $F$ a finite extension of $\mathbb Q_\ell$ and
$\Gamma=\operatorname{Gal}(\overline{\mathbb Q}_\ell/F)$. Let $\mathcal O$ be the ring of integers
of a finite extension of $\mathbb Q_\ell$, large enough for what follows. A \emph{lattice} is a
finite free $\mathcal O$-module $T$ with a continuous $\mathcal O$-linear action of $\Gamma$; its
Tate dual is $T^*=\operatorname{Hom}_{\mathcal O}(T,\mathcal O(1))$, and its Hodge--Tate weights,
when $T[1/\ell]$ is potentially crystalline, are indexed by the embeddings
$\gamma\colon\mathcal O\to\mathbb C_\ell$; we call them the \emph{jumps} of $T$ at $\gamma$. A
continuous cocycle $\Gamma\to T$ is \emph{crystalline} if its class lies in the preimage
$H^1_f(\Gamma,T)$ of the subspace $H^1_f(\Gamma,T[1/\ell])$ of Bloch and Kato~\cite{BlochKato};
the crystalline cocycles form an $\mathcal O$-module $Z_f(T)$.

By a theorem of Bloch and Kato, $H^1_f$ of a de Rham representation and $H^1_f$ of its Tate dual
are orthogonal under local Tate duality. The integral question studied here arises when the
duality itself exists only modulo a power of $\ell$. Let $Y$ and $L$ be lattices. A
\emph{pairing modulo $\ell^m$} of $Y$ and $L$ is an $\mathcal O$-bilinear map
$\psi\colon Y\times L\to\mathcal O(1)$ whose failure to be equivariant is divisible by $\ell^m$,
$g\psi-\psi=\ell^m c_\psi(g)$, with $c_\psi$ a crystalline cocycle of the lattice of bilinear
maps; it induces a $\Gamma$-equivariant pairing $b_\psi\colon Y\times L\to\mathbb Z/\ell^m(1)$.
For crystalline cocycles $a$ of $Y$ and $c$ of $L$ the cup product of their reductions through
$b_\psi$ is a class in $H^2(\Gamma,\mathbb Z/\ell^m(1))=\mathbb Z/\ell^m$. If $\psi$ were
equivariant this class would vanish. In general it does not, and the question is how far it is
from vanishing. For a class $\mathcal C$ of lattices we say that $\mathcal C$ \emph{has cup
exponent $\kappa$ on $L$} if, for all $Y\in\mathcal C$, all $m$, all pairings modulo $\ell^m$ of $Y$
and $L$ and all $a\in Z_f(Y)$, $c\in Z_f(L)$, the cup product of $a$ and $\ell^\kappa c$
vanishes. More precisely, the cocycles $c$ of $L$ whose cup product with every such $a$ vanishes
form a module $L^{\perp\mathcal C}$, and the \emph{dual defect} is the normalised length
\[
\delta^\vee_{\mathcal C}(L)=\operatorname{len}\bigl(Z_f(L)/(Z_f(L)\cap L^{\perp\mathcal C})\bigr),
\qquad\operatorname{len}(\mathcal O/\ell)=1 ;
\]
a cup exponent $\kappa$ gives $\delta^\vee_{\mathcal C}(L)\le2\kappa\,[F:\mathbb Q_\ell]\operatorname{rk}L$
({\cite[Lemma~5.6]{P55}}). These are the notions of {\cite[Definition~5.1]{P55}}.
By local Tate duality, $\delta^\vee_{\mathcal C}(T^*)$ bounds the following more familiar
quantity: the length of the module of cocycles $\Gamma\to T/\ell^M$ that are images of
crystalline cocycles of members of $\mathcal C$ under equivariant maps $Y\to T/\ell^M$, modulo
those that are reductions of crystalline cocycles of $T$ ({\cite[Proposition~5.4]{P55}}). Thus the dual defect
measures the failure of $H^1_f$ to be functorial for maps that exist only modulo a power of
$\ell$.

When $F$ is unramified and all jumps of all lattices involved lie in an interval of length at
most $\ell-2$, the theory of Fontaine and Laffaille~\cite{FontaineLaffaille}, with
\cite[\S4]{BlochKato}, leads one to expect that the defect is zero. We are concerned with the
opposite situation, in which the weights are much larger than $\ell$, the field may be ramified,
and $\ell$ may be~$2$. Bounds for torsion crystalline representations in large weight are
usually expressed through the largest Hodge--Tate weight and the ramification index. For the
application that motivates this paper such bounds are too weak by a large factor, and what is
needed is a bound sensitive to the distribution of the weights
(Remark~\ref{pki:rem-published}).

The hypothesis on the weights is the following. For integers $\eta,W,a\ge0$ and a sign
$\varepsilon$ put
\[
\operatorname{Win}(\eta,\varepsilon,a,W)=[-\eta,\eta]\,\cup\,\varepsilon\,[a,a+W]\subset\mathbb Z ,
\]
the union of a \emph{low window} and a \emph{far window}. Each embedding $\gamma$ is declared
\emph{thick} or \emph{thin} and carries a sign $\varepsilon(\gamma)$; the far windows are placed
at $a_\gamma+\delta$, where $a_\gamma=b$ if $\gamma$ is thick and $a_\gamma=G$ if it is thin.
We are given sets $U(\gamma),U'(\gamma)\subset\mathbb Z$ with
\[
U(\gamma)\subseteq\operatorname{Win}\bigl(\eta,-\varepsilon(\gamma),a_\gamma+\delta,W\bigr),\qquad
U'(\gamma)\subseteq\operatorname{Win}\bigl(\eta,\varepsilon(\gamma),a_\gamma+\delta,W\bigr),
\]
the class $\mathcal C(U)$ of all crystalline lattices whose jumps at each $\gamma$ lie in
$U(\gamma)$, and a crystalline lattice $L$ with jumps in $U'$. The integers
$Q,m,h,G,b,\delta,\eta,W$ are subject to the size conditions (S1)--(S4) of
Definition~\ref{pki:def-sizes}, which say in substance that $Q^2$, $WQ$ and $m$ are small
against $h$, that $h^3$ is small against $G$ and $Gh^3$ against $b$, that $\log b$ is small
against $m$, and that $\delta<\eta$ and $\eta Q$ is small against $h^3$. A condition
relates the thick embeddings and the signs to the action of Frobenius on the embeddings of the
residue field (\emph{the letters read thickness}, Definition~\ref{pki:def-reads}). It requires
in particular that the ramification index of $F$ be at most $h$; that $F$ be unramified and
$\ell\ge7$ when $\ell>Q$; and, when $\ell\le Q$, that the thick embeddings be sparse along the
Frobenius orbits and that the signs be constant unless $F$ is unramified and $\ell\ge7$. Finally $L$ \emph{has few low jumps} for $(h,\mu,\eta)$ if every negative
element $u$ of the sets $U(\gamma)$, $U'(\gamma)$ satisfies $2u\le-h$ and at most $\mu$ jumps of
$L$, with multiplicity, lie in the low window at each~$\gamma$. Put
\[
\lambda=v_\ell\Bigl((2b)!\cdot\Bigl(\prod_{b-30h^3<n\le b+30h^3}n\Bigr)^{h^3}\Bigr),\qquad
K=\Bigl\lfloor\frac{2\lambda}{Q}\Bigr\rfloor .
\]

\medskip\noindent\textbf{Theorem} (Theorem~\ref{pki:crys}). \emph{Under these hypotheses,
suppose that $L$ has few low jumps for $(h,\mu,\eta)$, or that every $U'(\gamma)$ is contained in
$[-\eta,0]$, in which case we put $\mu=0$. Then
\[
\delta^\vee_{\mathcal C(U)}(L)\le[F:\mathbb Q_\ell]\cdot10^5\bigl(\operatorname{rk}L\cdot K\,m+\mu\,\lambda\bigr).
\]
More precisely, if $\ell\le Q$ then $\mathcal C(U)$ has cup exponent
$10^4K(\lfloor\log_\ell h^3\rfloor+1)$ on $L$; if $\ell>Q$ and $\lambda\ge h^3$ it has cup exponent
$10^4Km$; and if $\ell>Q$ and $\lambda<h^3$, then $\delta^\vee_{\mathcal C(U)}(L)=0$ unless
$\lambda\ge\lfloor h/5\rfloor$ and $L$ has few low jumps.}\medskip

The number $\lambda$ is about $2b/(\ell-1)$ for small $\ell$ and vanishes for $\ell>2b$, so that
there is no defect at all for large $\ell$, and none for $\ell>Q$ unless $\ell$ is small against
$b$ or divides an integer close to $b$. For the remaining primes the bound has two terms. The
first is $\operatorname{rk}L$ times $m\lambda/Q$; the second is $\lambda$ times the number $\mu$ of
low jumps. In the application $Q=kt$ and $\mu\le\operatorname{rk}L\cdot m/k$ with $m$ much
smaller than $k$, so that both terms are at most of the order of
$\operatorname{rk}L\cdot\lambda\cdot m/k$; the factor $m/k$ is what the application requires, and
a bound that charges every embedding an amount growing with the largest weight does not give it
(Remark~\ref{pki:rem-published}). The hypotheses allow every prime, including $\ell=2$, and
ramified $F$ when $\ell\le Q$; no Fontaine--Laffaille hypothesis is made.

Theorem~\ref{pki:crys} is deduced in Section~\ref{pki:sec} from three statements on crystalline
lattices, each valid under its own hypotheses, which may be of independent use.
Proposition~\ref{pki:generic} treats \emph{generic primes}: $F$ unramified, $\ell\ge5$, and no
multiple of $\ell$ near the far windows; then the defect vanishes when there are no jumps in a gap
below $0$, and is bounded in terms of the number of low jumps otherwise. Proposition~\ref{pki:any} gives, for unramified $F$ and $\ell\ge7$, a
cup exponent of the order of $b/\ell^2+G/\ell$ for an arbitrary set of thick embeddings, and
Proposition~\ref{pki:sparse} one of the order of $e\,b/\ell^{J+1}$, with $e$ the ramification
index, when a Frobenius orbit contains no three thick letters with consecutive distances at most
$J$; it allows ramified $F$ when the signs are constant. A fourth statement,
Proposition~\ref{pki:tame}, concerns $F=\mathbb Q_\ell$ and potentially crystalline lattices that
become crystalline over a tame extension, with a condition on the action of inertia on
$D_{\mathrm{pcris}}$: for $\ell$ larger than the weights the class has a cup exponent that
depends only on the rank, on the degree of the polynomials that kill inertia, and on the widths
of the windows.

The method is the following. We fix a \emph{Kummer chart} of $\Gamma$
(Definition~\ref{pkp:def-chart}): a compatible system of $\ell$-power roots of a uniformiser,
which gives a coordinate $z$ in $A_{\mathrm{inf}}=W(\mathcal O_{\mathbb C_\ell}^\flat)$, the ring
$\mathfrak S=\Lambda[[z]]$ and an Eisenstein element $E\in\mathfrak S$ generating the kernel of
Fontaine's map $\theta$. For a lattice $L$ we work inside $W(\mathbb C_\ell^\flat)\otimes L$,
decomposed according to the embeddings of the residue field (the \emph{letters}), and consider
the submodules of elements with prescribed orders of vanishing along the primes
$\varphi^a(E)$, $a\ge0$, measured against the Hodge filtration through maps to
$B_{\mathrm{dR}}$ (Definition~\ref{pkp:def-orders}). These modules play the role of the
Breuil--Kisin module of $L$ with the elements $\varphi^a(E)$ inverted, but they are defined
directly from $L$ and the period rings. A
crystalline cocycle of $L$ is encoded by the extension of $L$ by a trivial lattice, and a class is
represented by \emph{points}: elements of the module of the extension that lie over a multiple
$\ell^ec$ of a constant. The \emph{pairing lemma} (Theorem~\ref{pkp:pairing}) is the link with cup
products. If $x$ and $x'$ are points over $\ell^{e_1}c$ and $\ell^{e_2}c'$ on the two sides of a
pairing modulo $\ell^m$, and if one can pass from $(\varphi x,\varphi x')$ to $(x,x')$ by a chain
of four moves each of which changes the pairing only by elements of a small ideal after
multiplication by $\ell^\kappa$, then $\ell^{e_1+e_2+\kappa}$ kills the cup product of the classes
of $c$ and $c'$. A cup exponent is therefore obtained from points with good orders over small
multiples of the constants. Such points are produced by a walk through the letters along the
Frobenius orbit (\S\ref{pkm:sec}), by elements that kill the cokernels which obstruct the steps,
and by \emph{deepening} (Theorem~\ref{pke:I2}): the depth of a point at the prime $\varphi^a(E)$
can be raised past an integer $n$ at the price of the factor $\ell^{v_\ell(n)+d+3a}$ in the
multiple of the constant, where $d=0$ if $F$ is unramified and $d=v_\ell(e)+1$ otherwise. For
$a=0$ and unramified $F$ a step is therefore free when $\ell$ does not divide $n$, and this is why
the far windows cost nothing when $\ell$ has no multiple near them. For such primes the constants
that admit points of the required depth without any multiple form a module whose colength is
bounded in terms of the number of low jumps (Theorem~\ref{pkh:I4}). This is the source of the
second term of the theorem, and of its dependence on $\mu$.

The parts of the argument that require most care are the following. To compute with the modules
above we prove a structure theorem, the existence of \emph{frames} (Theorem~\ref{pkf:main}); it
uses properties of Fontaine's rings \cite{pkfColmez98,pkfFontaineII,pkfFontaineIII}, the theorems
of Tate~\cite{pklTate67} and Ax~\cite{pkfAx70}, and tilting in the sense of Fontaine--Wintenberger
and Scholze~\cite{pkfWintenberger83,pkfScholze12}, in the form established in
\S\ref{pkr:sec}. When the signs $\varepsilon(\gamma)$ change along a Frobenius orbit the
far windows of consecutive letters lie on opposite sides of $0$, and the walk must cross a
\emph{long gap}; this is done in \S\ref{pkg:sec}. For ramified $F$ the different enters
the costs (\S\S\ref{pkl:sec} and~\ref{pke:sec}). The tame case requires a chart adapted to a tamely
ramified extension and an analysis of inertia on $D_{\mathrm{pcris}}$ (Section~\ref{pkw:sec}).

This paper belongs to a series whose main paper~\cite{P46} deduces the Riemann hypothesis for
Dirichlet $L$-functions from three theorems on integral algebraic modular forms on a chain of
theta lifts between unitary groups over a CM field $E$, the Short Forms Theorem, the Spanning
Theorem and the Covolume Theorem. One input of the third is the existence, for every prime
$\ell$, of a twist for which the congruences between theta lifts and other forms are bounded;
this is proved in~\cite{P55} by Ribet's method, which produces torsion cocycles above $\ell$
whose distance from the crystalline ones is bounded by a defect of the kind studied here. There
$F$ is a completion $E_v$ of $E$ (in the body of the paper we write $E_v$ for $F$ and $D_v$ for
$\Gamma$), the lattices are those of the Galois representations of the lifts,
the thick embeddings correspond to a sparse set of infinite places of $E$ at which the weight
$b$ is large, and $m$ is the number of lifts. The statement required is
{\cite[Theorem~5.8]{P55}}. We deduce it in \S\ref{pki:ss-deduce} from Theorem~\ref{pki:crys} and
Proposition~\ref{pki:tame}, after showing that the data of~\cite{P46} satisfy the hypotheses
(Propositions~\ref{pki:chart} and~\ref{pki:data}); only these two propositions involve the
global situation, and the reader interested in the local results may ignore them. What we take
from~\cite{P55} is the following: the definitions of crystalline cocycles, of the dual defect and
of cup exponents ({\cite[Definition~5.1]{P55}}), with the elementary {\cite[Lemma~5.6]{P55}}, and, for
Proposition~\ref{pki:data} only, the Hodge--Tate weights and the inertial types of the Galois
representations of the lifts ({\cite[Definition~2.6]{P55}} and the classes of lattices defined before {\cite[Proposition~5.2]{P55}}). None of this depends on {\cite[Theorem~5.8]{P55}}, which is stated in~\cite{P55} and proved here, and through which, with Theorem~\ref{pkj:chart}, the
results of this paper enter~\cite{P55}. Where we mention the local loss itself ({\cite[Propositions~2.28, 4.5 and~5.2]{P55}}), it is to say what the theorem is used for, or to name the lattices to which it is applied. So the two papers do not form a circle. The
statements of~\cite{P46} and~\cite{P55} to which we refer are reproduced in
Appendix~\ref{C4:quoted}. Apart from them, and from the published results quoted with precise
references where they are used, the paper is self-contained.

Section~\ref{pki:sec} states the core propositions and deduces Theorem~\ref{pki:crys} and
{\cite[Theorem~5.8]{P55}}. Section~\ref{C4:g2} introduces Kummer charts, letters, orders and
points and proves the pairing lemma (\S\ref{pkp:sec}), and contains the walk and the cokernel
killers (\S\ref{pkm:sec}). Sections~\ref{C4:g3} and~\ref{C4:g4} develop the modules: their
definition, the existence of frames and the period rings in Section~\ref{C4:g3}; standard orders
and Galois stability, and the behaviour of orders under products and Frobenius, in
Section~\ref{C4:g4}. Section~\ref{C4:g5} proves the three quantitative inputs: deepening with
costs, the long gap, and the colength of the deep constants. Section~\ref{pkj:sec} constructs the
chart for the data of~\cite{P46}, Section~\ref{C4:g7} proves the four core propositions, and
Section~\ref{pkw:sec} treats tame charts.


\providecommand{\auO}{\mathcal O}
\providecommand{\auQlb}{\overline{\Q}_\ell}
\providecommand{\pglen}{\operatorname{len}}
\providecommand{\pllZf}{Z_f}
\providecommand{\pllrk}{\operatorname{rk}}
\providecommand{\plldef}{\delta^{\vee}}
\providecommand{\pkiWin}{\operatorname{Win}}
\providecommand{\pkiC}{\mathcal C}
\providecommand{\pkilog}{\operatorname{lg}}

\section{The explicit cup exponents}\label{pki:sec}

This section concerns {\cite[Theorem~5.8]{P55}}. In \S\ref{pki:ss-need} we determine what is needed of this theorem and explain why we do not replace the argument by a published theorem. In
\S\ref{pki:ss-arch} we state the four propositions on which everything rests
(Propositions~\ref{pki:generic}, \ref{pki:any}, \ref{pki:sparse}, \ref{pki:tame}), together with two
inputs on the data (Propositions~\ref{pki:chart}, \ref{pki:data}), and we deduce from them a bound
uniform in $\ell$ (Theorem~\ref{pki:crys}) and {\cite[Theorem~5.8]{P55}} itself
(\S\ref{pki:ss-deduce}).

\emph{Dependencies.} Proved here in full: Theorem~\ref{pki:crys} and the deduction of
{\cite[Theorem~5.8]{P55}} from Propositions~\ref{pki:generic}--\ref{pki:data}. These six
propositions are proved in \S\S\ref{pkc:sec}, \ref{pkd:sec}, \ref{pkj:sec} (and Lemma~\ref{pki:data-sizes}), from the theory of Kummer
charts and the pairing lemma written in \S\S\ref{pkp:sec}--\ref{pkw:sec}; the chain is recorded at the
end of \S\ref{pki:ss-arch}.

\emph{Conventions.} As in {\cite[\S5]{P55}}: $\ell$ a prime, $v\mid\ell$ a place of $E$,
$\Gamma=D_v=\operatorname{Gal}(\auQlb/E_v)$, $\auO$ a large coefficient ring, lengths normalised so
that $\pglen\auO/\ell=1$, $\gamma$ runs over the embeddings $\auO\to\C_\ell$. For a lattice $L$, its
crystalline cocycles $\pllZf(L)$, the cup-orthogonal $L^{\perp\pkiC}$, the dual defect
$\plldef_\pkiC(L)$ and ``cup exponent'' are as in {\cite[Definition~5.1]{P55}}. For finite sets
$U(\gamma)\subset\Z$, $\pkiC(U)$ is the class of crystalline lattices whose jumps at each $\gamma$
lie in $U(\gamma)$. We write $\pkilog_\ell H=\lfloor\log H/\log\ell\rfloor$ for $H\ge1$.

\subsection{What is needed}\label{pki:ss-need}

{\cite[Theorem~5.8]{P55}} is used only through {\cite[(9)]{P55}}, summed over $v\mid\ell$ and
over the summands ({\cite[Lemma~5.10]{P55}}), and then over $\ell$ in {\cite[Lemma~2.5]{P55}}: what must hold
is $\sum_{\ell\in T}\pgloc_\ell\log\ell\le c\,\aubud$ for a small absolute fraction $c$, where
$\aubud=f\,x\,b(\log b)^3/(mk)$ is the budget and $\pgloc_\ell$ is the loss per unit of rank of the
main summand. {\cite[Lemma~2.5]{P55}} obtains $\sum_\ell\pgloc_\ell\log\ell\le 10^8f(m/k)\,b\log b$,
leaving a margin of a factor $x(\log b)^2/(10^8m^2)\ge(\log b)^2/(4\cdot10^8)$.

\begin{remark}[Why a published bound does not suffice]\label{pki:rem-published}
The decisive feature of the required bound is the factor $1/k$, with $k=r+1-m\ge x^5/2$ by
{\cite[(36)]{P46}}. Any bound that charges, at each $\ell$ and each embedding, an amount growing
with the \emph{largest} Hodge--Tate weight $\asymp b$ fails by a power of $x$. Indeed, the bounds
in the literature for torsion crystalline (or semistable) representations in large weight
(Fontaine--Laffaille outside its range, the results of T.~Liu and of Caruso--Liu on lifting maps
between torsion crystalline representations, Gee--Liu--Savitt type results) are expressed through
the maximal weight $H$ and the ramification index $e$; even under the most favourable reading,
an exponent of size $\asymp\pkilog_\ell(eH)+1$ for every $\ell\le eH$, the resulting contribution
is, by the argument of {\cite[Lemma~5.6]{P55}} and of {\cite[Lemma~5.10]{P55}},
\[
\sum_{\ell\le 4b}4f\,r'\,(\pkilog_\ell(4b)+1)\log\ell\ \ge\ 4f\,r'\sum_{\ell\le4b}\log 4b\cdot\tfrac12
\ \gg\ f\,r'\,b,
\]
whereas the budget allows $r'\cdot c\,\aubud\le r'f\,b\,x^{-4.4}$. The argument below gains the
factor $1/k$ from two facts specific to the data: at all but a sparse set of embeddings (the
\emph{thin} ones, not $b$-places) the far weights are of size $G\ll b$, and the main summand has at
most $\lfloor r'm/k\rfloor$ jumps in the low window. No published theorem we know bounds the
torsion of crystalline extension classes embedding by embedding in this way. We therefore write a
direct argument. (For $\ell\ge4b+2$, $\ell\nmid q_1dM$, see also {\cite[Remark~5.7]{P55}} for another route;
the argument below covers all $\ell$ uniformly and does not use it.)
\end{remark}

\subsection{Architecture}\label{pki:ss-arch}

\subsubsection*{Windows, letters, sizes}

For integers $\eta,W\ge0$, $a\ge0$ and a sign $\varepsilon\in\{\pm1\}$, the \emph{windows} are
\[
\pkiWin(\eta,\varepsilon,a,W)=[-\eta,\eta]\cup\varepsilon\,[a,a+W]\subset\Z
\]
(the \emph{low window} and the \emph{far window}). For $n\ge0$ and $a\in\Z$ we say that $\ell$ has
\emph{no multiple near $a$ at distance $n$}, written $\mathrm{NM}(\ell,n,a)$, if $\ell$ divides no
nonzero integer $y$ with $|y|\le n$, $|y-a|\le n$ or $|y+a|\le n$.

Let $e$, $f$ be the ramification index and residue degree of $E_v/\Q_\ell$, and $d_v=0$ if $e=1$,
$d_v=v_\ell(e)+1$ otherwise. Fix $\zeta\in\auO$ of exact order $\ell^f-1$ (possible as $\auO$ is
large). The \emph{letters} are the $(\ell^f-1)$-th roots of unity $z\in\C_\ell$; the \emph{letter}
of $\gamma$ is $\gamma(\zeta)$; the Frobenius acts on letters by $z\mapsto z^\ell$. For jump sets
$U$, the \emph{letter jump set} at $z$ is $U_z=\{0\}\cup\bigcup_{\gamma(\zeta)=z}U(\gamma)$.

\begin{definition}[Sizes]\label{pki:def-sizes}
Integers $Q,m,h,G,b,\delta,\eta,W\ge0$ satisfy the \emph{size conditions} if
\begin{enumerate}
\item[(S1)] $Q\ge1$ and $99\,(Q^2+(W+1)Q+Q+W+m)\le h$;
\item[(S2)] $99\,h^3\le G$ and $99\,(Gh^3+G)\le b$;
\item[(S3)] $2\log(4b)<m$;
\item[(S4)] $\delta<\eta$ and $99\,((\eta+W)Q+\eta+h)\le h^3$.
\end{enumerate}
The \emph{loss} at $\ell$ is $\lambda=\lambda_\ell(h,b)=v_\ell\bigl((2b)!\cdot(\prod_{b-30h^3<n\le
b+30h^3}n)^{h^3}\bigr)$; with $h=r^8$ it is the $\lambda_\ell$ of {\cite[\S2.1]{P55}}.
\end{definition}

\begin{lemma}\label{pki:elem}
Let $\lambda=\lambda_\ell(h,b)$. (i) $2b<\ell(\lambda+1)$. (ii) If $\ell$ divides some
$n\in(b-30h^3,b+30h^3]$, then $\lambda\ge h^3$. (iii) If (S3) holds and $1\le H\le4b$ then
$\pkilog_\ell H<m$. (iv) Under (S1)--(S4), $\pkilog_\ell h^3<m$, $d_v\le\pkilog_\ell e+1$, and if
$e\le h$ then $\pkilog_\ell h^3+1+d_v\le h$.
\end{lemma}

\begin{proof}
(i) $\lambda\ge v_\ell((2b)!)\ge\lfloor2b/\ell\rfloor$, and $2b<\ell(\lfloor2b/\ell\rfloor+1)$.
(ii) $n^{h^3}$ divides the product. (iii) $\log\ell\ge\log2>1/2$, so $\pkilog_\ell H\le2\log
H\le2\log 4b<m$. (iv) $h^3\le b\le4b$ by (S2) and $e\le h$, so (iii) applies to $h^3$ and to $e$;
$d_v\le v_\ell(e)+1\le\pkilog_\ell e+1$. Hence $\pkilog_\ell h^3+1+d_v<2m+2\le h$ by (S1).
\end{proof}

\begin{definition}[Letters that read thickness]\label{pki:def-reads}
Let each embedding $\gamma$ be \emph{thick} or \emph{thin} and carry a sign $\varepsilon(\gamma)$, and
let $T$ (a set of letters) and $s$ (a sign on letters) be given. We say that \emph{the letters read
thickness} (for $Q$, $h$) if: $e\le h$; for every $\gamma$ that is thick, and for every $\gamma$ if
$\ell>Q$, one has [$\gamma$ thick $\Leftrightarrow\gamma(\zeta)\in T$] and
$\varepsilon(\gamma)=s(\gamma(\zeta))$; either $e=1$ and $\ell\ge7$, or $\ell\le Q$ and $s$ is
constant; and, if $\ell\le Q$, there is no letter $z$ with $z$, $z^{\ell^{a}}$, $z^{\ell^{a+a'}}$ all
in $T$ for some $a,a'\in[1,\pkilog_\ell h^3]$.
\end{definition}

A lattice $L$ \emph{has few low jumps} (for $h,\mu,\eta$ and jump sets $U,U'$) if every negative
element $u$ of every $U(\gamma)\cup U'(\gamma)$ satisfies $2u\le-h$, and at every $\gamma$ at most
$\mu$ jumps of $L$, with multiplicity, lie in $[-\eta,\eta]$.

\subsubsection*{The four core propositions}

In Propositions~\ref{pki:generic}--\ref{pki:sparse}, $U$ and $U'$ are jump sets, $L$ is a crystalline
lattice with jumps in $U'$, and the class is $\pkiC(U)$; at each letter $z$ an integer $b_z\ge0$
and a sign $\varepsilon_z$ are given.

\begin{proposition}[Generic primes]\label{pki:generic}
Let $e=1$, $\ell\ge5$, and $a,\eta,W\ge0$ with $\eta\ge1$, such that for every letter $z$:
$b_z+W+\eta\le(\ell-3)a$; $6(a+\eta+W)<b_z$; $\mathrm{NM}(\ell,6(a+\eta+W),b_z)$; $U_z\subseteq
\pkiWin(\eta,-\varepsilon_z,b_z,W)$ and $U'_z\subseteq\pkiWin(\eta,\varepsilon_z,b_z,W)$.
\begin{enumerate}
\item[(i)] If no element of any $U(\gamma)\cup U'(\gamma)$ lies in $[-a,0)$, or if every
$U'(\gamma)\subseteq[-\eta,0]$, then $\plldef_{\pkiC(U)}(L)=0$.
\item[(ii)] If $L$ has at most $\mu$ jumps in $[-\eta,\eta]$ at each $\gamma$, then
$\plldef_{\pkiC(U)}(L)\le11f\bigl(\mu(a+1)+\pllrk_\auO L\cdot(W+1)\bigr)$.
\end{enumerate}
\end{proposition}

\begin{proposition}[Any pattern of $b$-letters]\label{pki:any}
Let $e=1$, $\ell\ge7$, and $\eta\ge1$, $V,M,G,b,K\ge0$ with $V+4\eta\le M\le K$, $0<G\le b$,
$b<K\ell^2$, $2G<K\ell$. Let $B$ be a set of letters, $b_z=b$ for $z\in B$ and $b_z=G$ otherwise,
$U_z\subseteq\pkiWin(\eta,-\varepsilon_z,b_z,V)$ and $U'_z\subseteq\pkiWin(\eta,\varepsilon_z,b_z,V)$.
Then $\pkiC(U)$ has cup exponent $10^4\bigl(K+\pkilog_\ell(2(b+M))\bigr)$ on $L$.
\end{proposition}

\begin{proposition}[Sparse $b$-letters]\label{pki:sparse}
Let $\eta,J\ge1$, $V,G,b,M,\vartheta\ge0$ with $G+V\le\eta$. Let $B$, $b_z$ and the window
conditions be as in Proposition~\ref{pki:any} (with width $V$). Assume: there is no letter $z$ with
$z,z^{\ell^s},z^{\ell^{s+t}}\in B$ for some $s,t\in[1,J]$; either [$e=1$, $\ell\ge7$ and
$10\ell^J\eta\le b$] or [$\varepsilon_z$ is constant and $8e\le\ell^{J+1}$];
$2e(b+5\eta)\le\ell^{J+1}\vartheta$; and $20(J+1+d_v)\eta+\pkilog_\ell(2(b+5\eta))\le M$. Then
$\pkiC(U)$ has cup exponent $10^3\bigl(\vartheta+(J+1)M\bigr)$ on $L$.
\end{proposition}

\begin{proposition}[Tame types]\label{pki:tame}
Let $E_v=\Q_\ell$, $r,N\ge1$, $e_0=\operatorname{lcm}_{1\le i\le r}(\ell^i-1)$. For each inertial
$s\in\Gamma$ let $K(s)$ be a set of polynomials over $\auO$ containing $X^{e_0}-1$ and the image of a
monic polynomial over $\Z_\ell$ of degree $\le N$; let $\pkiC(U,K)$ be the class of potentially
crystalline lattices with jumps in $U$ on whose $D_{\mathrm{pcris}}$ each $P\in K(s)$ kills $s$, for
every inertial $s$. Let $L\in\pkiC(U',K^{-1})$, where $K^{-1}(s)=K(s^{-1})$. Let $\eta,W,b\ge0$ and
a sign $\varepsilon$ with $U(\gamma)\subseteq\pkiWin(\eta,-\varepsilon,b,W)$ and
$U'(\gamma)\subseteq\pkiWin(\eta,\varepsilon,b,W)$ for all $\gamma$, and
$\ell>b+6(W+\eta)+r^2(N+1)^2+2$. Then $\pkiC(U,K)$ has cup exponent
$3r^2(N+1)^2+150(\eta+W+1)$ on $L$.
\end{proposition}

\subsubsection*{Two inputs on the data}

\begin{proposition}[A chart for the data]\label{pki:chart}
Let $Q=kt$, $h=r^8$. Call $\gamma$ thick if the complex embedding $\sigma_\gamma$ of $E$ attached to
$\gamma$ through $\iota$ and $v$ (Convention~\ref{pkj:conv}) lies over a $b$-place, and let
$\varepsilon(\gamma)$ be the sign given by the CM type at $\sigma_\gamma$. For every $\ell$ and $v\mid\ell$,
every $\auO\supseteq\auO_{E_v}$ and every $\zeta\in\auO$ of exact order $\ell^f-1$ there are $T$ and $s$
for which the letters read thickness (Definition~\ref{pki:def-reads}).
\end{proposition}

\begin{proposition}[The summands fit the windows]\label{pki:data}
Let $h=r^8$, $Q=kt$, $G$ the weight at the places that are not $b$-places, $\eta=h^2+1$, $W=4t$, and
for $0\le j\le m$ let $\delta_j=j\Delta+h-2t$. Then (S1)--(S4) hold for $(Q,m,h,G,b,\delta_j,\eta,W)$.
Let $L_c$ and $\mathcal C_v$ be as in {\cite[Proposition~5.2]{P55}} and let $U(\gamma)$ be the jump set
of $\mathcal C_v$, $U'(\gamma)$ that of $L_c^*$. There are $j\le m$ and, at each $\gamma$, a sign
$\varepsilon(\gamma)$, with $U(\gamma)\subseteq\pkiWin(\eta,-\varepsilon(\gamma),a_\gamma+\delta_j,W)$
and $U'(\gamma)\subseteq\pkiWin(\eta,\varepsilon(\gamma),a_\gamma+\delta_j,W)$, where $a_\gamma=b$ if
$\gamma$ is thick and $a_\gamma=G$ otherwise. Moreover $L_{\mathrm{main}}^*$ has few low jumps for
$(h,\lfloor r'm/k\rfloor,\eta)$, and for a character $L_e$ every $U'(\gamma)\subseteq[-\eta,0]$.
At $\ell=p_i$, $\mathcal C_v\subseteq\pkiC(U,K)$ and $L_c^*\in\pkiC(U',K^{-1})$ for a $K$ as in
Proposition~\ref{pki:tame} with $N=\max(r+t,2r)\,r$, $\varepsilon(\gamma)$ and $a_\gamma$ do not depend on $\gamma$ (write $a$ for the common value), and
$\ell>a+\delta_j+6(W+\eta)+r^2(N+1)^2+2$ for $a\in\{b,G\}$.
\end{proposition}

\begin{lemma}[Sizes of the data]\label{pki:data-sizes}
With $h=r^8$, $Q=kt$, $G=a/x_1$ (the weight $\kappa(\tau)$ at the places $\tau\notin\auPb$),
$\eta=h^2+1$, $W=4t$ and $\delta_j=j\Delta+h-2t$, $0\le j\le m$, conditions (S1)--(S4) hold ($G$ is an integer since $x_1\mid b$, (D4) of {\cite[Definition~13.3]{P46}}). This is
the first assertion of Proposition~\ref{pki:data}.
\end{lemma}

\begin{proof}
By {\cite[(36)]{P46}}: $x^5\le r\le3x^5$, $t=5r$, $k\le r$, $\sqrt x\le m\le2\sqrt x$, $\Delta=2h$,
$(\log b)^{100}\le x$, $x^{10^6}\le x_1$, $x_1^{10^6}\le b$, $a=xb$. (S1): $Q\le5r^2$, $W=20r$, so
$Q^2+(W+1)Q+Q+W+m\le25r^4+105r^3+25r^2+m\le200r^4$, and $99\cdot200r^4\le r^8=h$ since
$r\ge x^5\ge 10^2$. (S2): $h^3=r^{24}\le3^{24}x^{120}$, while $G=xb/x_1\ge b/x_1\ge
x_1^{10^6-1}\ge x^{10^{11}}$, so $99h^3\le G$; and $99(Gh^3+G)\le 198\,h^3\,xb/x_1\le b$ because
$198\cdot3^{24}x^{121}\le x_1$. (S3): $2\log(4b)\le3\log b\le3x^{1/100}<\sqrt x\le m$. (S4):
$\delta_j\le2mh+h\le(4\sqrt x+1)h<h^2+1=\eta$; and $(\eta+W)Q+\eta+h\le(r^{16}+20r+1)\cdot5r^2+
r^{16}+r^8+1\le7r^{18}$, with $99\cdot7r^{18}\le r^{24}=h^3$.
\end{proof}

\subsubsection*{Where the propositions are proved}

Proposition~\ref{pki:generic} is Theorem~\ref{pkc:generic}; Proposition~\ref{pki:tame} is
Theorem~\ref{pkc:tame}; Proposition~\ref{pki:sparse} is Theorem~\ref{pkd:sparse};
Proposition~\ref{pki:any} is Theorem~\ref{pkd:any}; Proposition~\ref{pki:chart} is
Theorem~\ref{pkj:chart}; Proposition~\ref{pki:data} is Lemma~\ref{pki:data-sizes} together with
Theorem~\ref{pkd:data}. The four cases rest on the pairing lemma and the walk
(\S\S\ref{pkp:sec}, \ref{pkm:sec}) and on the inputs (I1) (Proposition~\ref{pkt:I1}), (I2)
(Theorem~\ref{pke:I2}), (I3) (Proposition~\ref{pkt:I3}), (I4) (Theorem~\ref{pkh:I4}), (I5)
(Theorem~\ref{pkg:I5}), (I6) (Lemma~\ref{pkw:dual}) and (T1)--(T4) (Proposition~\ref{pkw:T1},
Theorems~\ref{pkw:T2}, \ref{pkw:T3}, \ref{pkw:T4}); these in turn rest on the modules $\mathcal N$ of
\S\ref{pkn:sec}, whose properties (Lemma~\ref{pkm:inputs}, Lemmas~\ref{pkp:orders-add},
\ref{pkp:tau}) are proved in \S\S\ref{pkn:sec}--\ref{pkt:sec}, the existence of frames
(Theorem~\ref{pkf:main}) and the ring statements of \S\ref{pkr:sec}. Published results enter only
through \S\S\ref{pkf:sec}, \ref{pkr:sec}, \ref{pkl:sec}, \ref{pkw:sec} (Fontaine's period rings,
Tate's theorem, Fontaine--Wintenberger/Scholze tilting, and the facts on $D_{\mathrm{pcris}}$ listed
there).

\subsection{The uniform crystalline bound and the deduction of \texorpdfstring{{\cite[Theorem~5.8]{P55}}}{[Galois, Theorem~5.8]}}\label{pki:ss-deduce}

\begin{lemma}[Letter windows]\label{pki:letterwin}
Let the letters read thickness, let $\eta\le G$ and $\delta+W\le G$, and let $U(\gamma)\subseteq
\pkiWin(\eta,-\varepsilon(\gamma),a_\gamma+\delta,W)$ with $a_\gamma=b$ for thick and $G$ for thin
$\gamma$. Put $\eta_\ell=2G$ if $\ell\le Q$ and $\eta_\ell=\eta$ if $\ell>Q$, and $b_z=b+\delta$ if
$z\in T$, $b_z=G+\delta$ otherwise. Then $U_z\subseteq\pkiWin(\eta_\ell,-s(z),b_z,W)$ for every
letter $z$. The same holds with all signs reversed.
\end{lemma}

\begin{proof}
$0$ lies in the low window. Let $u\in U(\gamma)$ with $\gamma(\zeta)=z$. If $\gamma$ is thick, or
$\ell>Q$, then thickness and sign of $\gamma$ are those read from $z$, so $u$ lies in
$[-\eta,\eta]\subseteq[-\eta_\ell,\eta_\ell]$ or in $-s(z)[b_z,b_z+W]$. Otherwise $\gamma$ is thin and
$\ell\le Q$; then $|u|\le\max(\eta,G+\delta+W)\le2G=\eta_\ell$.
\end{proof}

\begin{theorem}[Uniform crystalline bound]\label{pki:crys}
Assume (S1)--(S4) for $(Q,m,h,G,b,\delta,\eta,W)$, that the letters read thickness, that
$U(\gamma)\subseteq\pkiWin(\eta,-\varepsilon(\gamma),a_\gamma+\delta,W)$ and
$U'(\gamma)\subseteq\pkiWin(\eta,\varepsilon(\gamma),a_\gamma+\delta,W)$ for all $\gamma$ (with
$a_\gamma$ as in Lemma~\ref{pki:letterwin}), that $L$ is crystalline with jumps in $U'$, and that
either $L$ has few low jumps for $(h,\mu,\eta)$, or every $U'(\gamma)\subseteq[-\eta,0]$. Let
$\lambda=\lambda_\ell(h,b)$ and $K=\lfloor2\lambda/Q\rfloor$. Then:
\begin{enumerate}
\item[(b)] if $\ell\le Q$, $\pkiC(U)$ has cup exponent $10^4K(\pkilog_\ell h^3+1)$ on $L$;
\item[(c)] if $\ell>Q$ and $\lambda\ge h^3$, it has cup exponent $10^4Km$;
\item[(d)] if $\ell>Q$ and $\lambda<h^3$, then $\plldef_{\pkiC(U)}(L)=0$ unless $\lambda\ge\lfloor
h/5\rfloor$ and $L$ has few low jumps; and always
$\plldef_{\pkiC(U)}(L)\le[E_v:\Q_\ell]\cdot10^5(\pllrk L\cdot Km+\mu\lambda)$.
\end{enumerate}
In all cases $\plldef_{\pkiC(U)}(L)\le[E_v:\Q_\ell]\cdot10^5(\pllrk L\cdot Km+\mu\lambda)$.
\end{theorem}

\begin{proof}
By (S2), (S4): $\eta\le h^3\le G$, $\delta+W<\eta+W\le G$, so Lemma~\ref{pki:letterwin} applies
to $U$ and $U'$. By Lemma~\ref{pki:elem}(i), $2b<\ell(\lambda+1)$; by definition of $K$,
$2\lambda<(K+1)Q$.

\emph{(b) $\ell\le Q$.} First, $K$ is large: from $\lambda+1>2b/\ell\ge2b/Q$ and
$2\lambda<(K+1)Q$ we get $K+3>4b/Q^2$, and $Q^2\le h/99$ by (S1), so
$h^2K>396\,bh-3h^2\ge 2(b+\delta)$ and $K\ge 396\,b/h-3\ge41Gh$ by (S2), (S4) (note $\delta<\eta\le
h^3\le b$). Apply Proposition~\ref{pki:sparse} with low radius $2G$, $b\to b+\delta$, $G\to G+\delta$,
$V=W$, $B=T$, $\varepsilon_z=s(z)$ (and the reversed sign for $U$), $J=\pkilog_\ell h^3$,
$\vartheta=2K$, $M=K$. The windows are those of Lemma~\ref{pki:letterwin}, and $G+\delta+W\le2G$.
$J\ge1$ as $\ell\le Q\le h\le h^3$. The sparseness is part of Definition~\ref{pki:def-reads}. If
$e=1$ and $\ell\ge7$: $10\ell^J\cdot2G\le20h^3G\le b+\delta$. If $s$ is constant:
$8e\le8h\le h^3<\ell^{J+1}$. Depth: $e\cdot2(b+\delta+10G)\le h\cdot4(b+\delta)<h^3\cdot2K\le
\ell^{J+1}\cdot2K$, using $10G\le b$, $e\le h$, $2(b+\delta)<h^2K$ and $h^3<\ell^{J+1}$. Last
condition: by Lemma~\ref{pki:elem}(iii)--(iv), $20(J+1+d_v)\cdot2G+\pkilog_\ell(2(b+\delta+10G))\le
40hG+m\le41hG\le K$. The exponent is $10^3(2K+(J+1)K)\le10^4K(J+1)$.

\emph{From now on $\ell>Q$}; by Definition~\ref{pki:def-reads}, $e=1$ and $\ell\ge7$, and the low
radius is $\eta$.

\emph{(c) $\lambda\ge h^3$.} Then $K+1>2h^3/Q$, so $K\ge h^3/Q\ge99h^2\ge2$ (as $Q\le h/99$).
Apply Proposition~\ref{pki:any} with $b\to b+\delta$, $G\to G+\delta$, $V=W$, $M=W+4\eta$, and $K$.
Conditions: $W+4\eta\le4(\eta+W)\le4h^3/(99Q)\le K$ by (S4). $0<G+\delta\le b+\delta$. Next,
$b<\ell(\lambda+1)/2<\ell((K+1)Q/2+1)/2\le\ell^2(K+1)/4+\ell/2$ since $Q<\ell$; and
$\ell^2K/4\ge\ell/2$, $\ell^2K/4\ge 12K\ge12h^3/Q>\delta$ (as $\delta<\eta\le h^3/(99Q)$); so
$b+\delta<\ell^2K$. Finally $K\ell\ge(2\lambda/Q-1)\ell\ge\lambda\ell/Q\ge\ell(\lambda+1)/(2Q)>b/Q$
(using $\lambda\ge h^3\ge Q$), and $b/Q\ge99Gh^3/Q>2(G+\delta)$. The exponent is
$10^4(K+\pkilog_\ell(2(b+\delta+W+4\eta)))\le10^4(K+m-1)\le10^4Km$ by Lemma~\ref{pki:elem}(iii)
(the argument is $\le4b$) and $K\ge1$.

\emph{(d) $\lambda<h^3$.} By Lemma~\ref{pki:elem}(ii), $\ell$ divides no integer in
$(b-30h^3,b+30h^3]$. Put $a=2(\lambda+1)\le2h^3$ and $n=6(a+\eta+W)\le12h^3+6h^3/99<13h^3$.
\emph{Gap.} $\ell>2b/(\lambda+1)\ge2b/h^3\ge198G$. A nonzero multiple $y$ of $\ell$ has $|y|\ge\ell>
G+\delta+n$ (as $n<13h^3\le G/7$ and $\delta<G/99$), so it is not within $n$ of $0$ or
$\pm(G+\delta)$; if it were within $n$ of $\pm(b+\delta)$, then $|y|\in(b-30h^3,b+30h^3]$ since
$\delta+n<30h^3$, a contradiction. So $\mathrm{NM}(\ell,n,b_z)$ for every letter. \emph{Heights.}
$(\ell-3)a=2\ell(\lambda+1)-6(\lambda+1)>4b-6h^3\ge b_z+W+\eta$, and $6(a+\eta+W)<13h^3<G\le b_z$.
So Proposition~\ref{pki:generic} applies with this $a$, radius $\eta$, $b_z$ and
$\varepsilon_z=s(z)$.

If $\lambda\ge\lfloor h/5\rfloor$ and $L$ has few low jumps, (ii) gives $\plldef\le
11f(\mu(2\lambda+3)+\pllrk L(W+1))$. Here $\lambda\ge1$, so $2\lambda+3\le6\lambda$; and
$(K+1)Q>2\lambda\ge h/3$ with $(W+1)Q\le h/99$ gives $K\ge W+1$; and $m\ge1$. Hence
$\plldef\le f\cdot10^5(\pllrk L\cdot Km+\mu\lambda)$, and $f=[E_v:\Q_\ell]$ as $e=1$.

Otherwise either every $U'(\gamma)\subseteq[-\eta,0]$, or $L$ has few low jumps and
$\lambda<\lfloor h/5\rfloor$; in the latter case every negative $u$ in $U(\gamma)\cup U'(\gamma)$
has $u\le-h/2<-2(\lambda+1)=-a$. In both cases (i) gives $\plldef=0$.

\emph{The final bound.} In cases (b), (c), {\cite[Lemma~5.6]{P55}} gives $\plldef\le2k'[E_v:\Q_\ell]
\pllrk L$ for the exponent $k'$; in case (b) $2\cdot10^4K(\pkilog_\ell h^3+1)\le10^5Km$ since
$\pkilog_\ell h^3+1\le m$ (Lemma~\ref{pki:elem}(iv)), and in case (c) $2\cdot10^4Km\le10^5Km$.
\end{proof}

\begin{proof}[Proof of {\cite[Theorem~5.8]{P55}}, given Propositions~\ref{pki:generic}--\ref{pki:data}]
Fix $\ell$, $v$, $c$ and put $L=L_c^*$, $\rho=\pllrk L_c$, $x=x_c$. By Proposition~\ref{pki:data}
(with its $j$), the size conditions and windows of Theorem~\ref{pki:crys} hold with
$\delta=\delta_j$; by Proposition~\ref{pki:chart} the letters read thickness; and $L$ has few low
jumps with $\mu=\lfloor\rho m/k\rfloor$ ($c=\mathrm{main}$, $x=1$) or all $U'(\gamma)\subseteq
[-\eta,0]$ ($c$ a character). If $\ell\notin\{p_0,p_1\}$, $\mathcal C_v=\pkiC(U)$ (the summands are crystalline above $\ell\nmid M=p_0p_1$ by (M2) of {\cite[Definition~2.6]{P55}} and (G2), the twist entering only through the cocycle conditions at $Q\cup\aubar Q$) and Theorem~\ref{pki:crys} gives (b), (c), (d) with the exponents printed in
{\cite[Theorem~5.8]{P55}}; in (d), the defect of a character vanishes, so the bound holds with
$\mu=\lfloor\rho m/(kx)\rfloor$ for both kinds of summand. If $\ell=p_i$, $E_v=\Q_\ell$ ($p_i$ splits
completely in $E$) and Proposition~\ref{pki:tame} with $N=\max(r+t,2r)r$, $W=4t$, $\eta=h^2+1$ and
$b\to a+\delta_j$, where $a=a_\gamma$ is the common value (Proposition~\ref{pki:data}), gives the exponent $3r^2(N+1)^2+150(h^2+4t+2)$ of (a), using the last clause
of Proposition~\ref{pki:data} (a cup exponent of $\pkiC(U,K)$ on $L$ is one of its subclass
$\mathcal C_v$). The passage to {\cite[(9)]{P55}} is {\cite[Remark~5.9]{P55}}.
\end{proof}

\begin{remark}[Comparison with {\cite[Proposition~4.5]{P55}}]\label{pki:align}
{\cite[Proposition~4.5(b)]{P55}} bounds $\sum_L\sum_v\delta_v(L)/e_\auO$, the quantity of
{\cite[Proposition~5.2]{P55}} with lengths normalised so that $\pglen\auO/\ell=1$; its part (a) is
{\cite[Proposition~5.4]{P55}} applied to the module of Ribet cocycles, whose elements come from
$\mathcal C_v$ by {\cite[Lemma~5.5]{P55}} (applied with $\Gamma=D_v$ inside $\pgG$ to the lattices
$u_x$, whose restrictions to $D_v$ lie in $\mathcal C_v$ by {\cite[Proposition~4.5(b)]{P55}}).
\end{remark}



\section{Kummer charts, the pairing lemma and the walk}\label{C4:g2}
\providecommand{\auO}{\mathcal O}
\providecommand{\auQlb}{\overline{\Q}_\ell}
\providecommand{\panot}[1]{\par\smallskip{\raggedright\noindent\textbf{[Not proved here:} #1\textbf{]}\par}\smallskip}
\providecommand{\pkpA}{A_{\mathrm{inf}}}
\providecommand{\pkpW}{W(C^\flat)}
\providecommand{\pkpm}{W(\mathfrak m^\flat)}
\providecommand{\pkpM}{M_m}
\providecommand{\pkpsm}{\mathrm{Sm}_m}
\providecommand{\pkpAmb}{\mathbb A}

\subsection{Kummer charts and the pairing lemma}\label{pkp:sec}

This section supplies the engine on which Propositions~\ref{pki:generic}--\ref{pki:tame} rest:
the Kummer charts of a local Galois group, their splitting into letters, the moves of a pair of
lattice points and the \emph{pairing lemma}, which turns a short chain of safe moves into a bound
for the exponent of a cup product. \S\ref{pkp:ss-arch} states every definition and result cited by
the four cases of \S\ref{pki:ss-arch}. The proofs follow in
\S\ref{pkp:ss-AS}--\S\ref{pkp:ss-walk}.

Lemmas~\ref{pkp:orders-add} and \ref{pkp:tau} are proved in \S\ref{pkt:sec} (Lemma~\ref{pkt:add},
Proposition~\ref{pkt:tau-i}, Lemmas~\ref{pkt:tau-iv}, \ref{pkt:tau-v}), Lemma~\ref{pkp:tau}(iii) is
Lemma~\ref{pkm:tau-iii}, and Proposition~\ref{pki:chart} is Theorem~\ref{pkj:chart}.

\emph{Conventions.} As in \S\ref{pki:sec}: $\ell$ a prime, $\Gamma=D_v$ for a place $v\mid\ell$ of
$E$, $j\colon\Gamma\to G_{\Q_\ell}=\operatorname{Gal}(\auQlb/\Q_\ell)$ the inclusion (an open
embedding), $\auO$ a large coefficient ring (finite free over $\Z_\ell$). $C=\C_\ell$, $O_C$ its
ring of integers, $O^\flat=\varprojlim_{x\mapsto x^\ell}O_C/\ell$ the tilt, $C^\flat$ its fraction
field (algebraically closed, of characteristic $\ell$), $\mathfrak m^\flat$ the maximal ideal of
$O^\flat$. Witt vectors are $\ell$-typical. $\pkpA=W(O^\flat)\subset\pkpW$; $\varphi$ is the Witt
vector Frobenius; $\theta\colon\pkpA\to O_C$ is Fontaine's map; $\varepsilon=(1,\zeta_\ell,
\zeta_{\ell^2},\dots)\in O^\flat$, $\tau=[\varepsilon^{1/\ell}]-1\in\pkpA$, $t=\log[\varepsilon]\in
B_{\mathrm{dR}}^+$. We write $\pkpm\subset\pkpA$ for the ideal of Witt vectors all of whose
coordinates lie in $\mathfrak m^\flat$ (the kernel of $\pkpA\to W(\overline{\mathbb F}_\ell)$). The
group $\Z\times\Gamma$ acts on $\pkpW$: the generator $\varphi=(1,1)$ by Frobenius, $g\in\Gamma$
through $j$; the two actions commute.

\subsubsection{Architecture}\label{pkp:ss-arch}

\paragraph*{Kummer charts and letters}

\begin{definition}[Kummer chart]\label{pkp:def-chart}
A \emph{Kummer chart} of $\Gamma$ consists of the data below; $K'=\auQlb^{\,j(\Gamma')}$ denotes the
fixed field of its subgroup $\Gamma'$, $e$ and $f$ its ramification index and residue degree over
$\Q_\ell$, and $\Lambda=W(\mathbb F_{\ell^f})\subset\pkpA$. The data are: an
element $\varpi^\flat\in\mathfrak m^\flat$ whose untilt $\pi=\theta([\varpi^\flat])$ is a uniformiser of
$K'$, the coordinate $z=[\varpi^\flat]\in\pkpA$, the ring $\mathfrak S=\Lambda[[z]]\subset\pkpA$ (the
closure of $\Lambda[z]$; $\Lambda[[X]]\to\mathfrak S$, $X\mapsto z$, is an isomorphism), an element
$E=z^eA+\ell\in\mathfrak S$ with $A\in\mathfrak S^\times$ generating $\ker\theta$ (so
$E(X)$ is an Eisenstein polynomial times a unit), and an open subgroup $\Gamma'\subseteq\Gamma$ of index prime to $\ell$
such that: $j(\Gamma')$ fixes $\Lambda$ and $\theta(z)$; every element of $G_{\Q_\ell}$ fixing $\Lambda$
pointwise and $\theta(z)$ lies in $j(\Gamma')$; and for every $l\ge0$ some $g\in\Gamma'$ has
$gz=[\varepsilon]^lz$. We write $p_a=\varphi^a(E)$ ($a\ge0$); $\theta\circ\varphi^{-a}$ has kernel
$p_a\pkpA$. The chart is \emph{unramified} if $e=1$ and $E=z+\ell$. We put $d_v=0$ if $e=1$ and
$d_v=v_\ell(e)+1$ otherwise.
\end{definition}

\begin{lemma}[Existence]\label{pkp:chart-exists}
Kummer charts of $\Gamma$ exist, with $\Gamma'=\Gamma$; if $j(\Gamma)$ contains the inertia group,
the chart can be taken unramified.
\end{lemma}

\begin{definition}[Letters]\label{pkp:def-letters}
The \emph{letters} are the elements $i\in\mathbb F_{\ell^f}^\times$; $[i]\in\Lambda$ is the
Teichm\"uller lift. Fix $\zeta\in\auO$ with $\zeta^{\ell^f-1}=1$ (Lemma~\ref{pkw:letters} shows that the $e_i$ below are then orthogonal idempotents with sum $1$; in the cases one takes $\zeta$ of exact order $\ell^f-1$, or $\zeta=1$ for a tame chart). For a finite free $\auO$-module $L$
put $\pkpAmb(L)=\pkpW\otimes_{\Z_\ell}L$; its \emph{letter part} at $i$ is
$\pkpAmb(L)_i=\{x:(1\otimes\zeta)x=[i]x\}$, and $e_i=\prod_{i'\ne i}\frac{(1\otimes\zeta)-[i']}{[i]-[i']}$
is the projector onto it (the differences $[i]-[i']$ are units of $\Lambda$, being Teichm\"uller lifts
of nonzero residues); $\sum_ie_i=1$ and $\pkpAmb(L)=\bigoplus_i\pkpAmb(L)_i$. Since
$\varphi[i]=[i^\ell]$ and $\Gamma'$ fixes $\Lambda$, $\varphi$ maps the part at $i$ to the part at
$i^\ell$ and $\Gamma'$ preserves each part. For $a\ge0$, $i^{(-a)}$ (``$a$ steps before $i$'') is the
letter with $(i^{(-a)})^{\ell^a}=i$. Under an embedding $\gamma\colon\auO\to C$ the letter of
$\gamma$ is the residue of $\gamma(\zeta)$; this matches the letters of \S\ref{pki:ss-arch}.
\end{definition}

\begin{definition}[Lattices with jump sets; orders]\label{pkp:def-orders}
A \emph{lattice with jump sets} is a finite free $\auO$-module $L$ with a continuous $\auO$-linear
action of $\Gamma$, crystalline on $\Gamma'$, together with finite sets $U(i)\subset\Z$ (one for each
letter) containing the Hodge--Tate jumps of $L$ at all embeddings of letter $i$. Let
$\pkpA\mathcal N=\pkpA\mathcal N(L)$ and $\mathcal N=\mathcal N(L)$ be the modules of
Definition~\ref{pkn:def-N} (they play the role of the Breuil--Kisin module of $L$ with the $p_a$
inverted). A \emph{form of type
$u$ at $i$} is an $\pkpA$-linear map $\mu\colon\pkpAmb(L)\to B_{\mathrm{dR}}$, $\Gamma'$-equivariant on
$1\otimes L$, with $\mu(e_i(1\otimes L))\subseteq t^{-u}B^+_{\mathrm{dR}}$. For a function
$o\colon\Z_{\ge0}\times\Z\to\Z$ (``orders''), $x$ \emph{has orders $o$ at $i$} if $x\in\pkpA\mathcal
N\cap\pkpAmb(L)_i$ and $\mu(\varphi^{-a}x)\in t^{o(a,u)}B^+_{\mathrm{dR}}$ for all $a\ge0$, all $u\in
U(i^{(-a)})$ and all forms $\mu$ of type $u$ at $i^{(-a)}$. These $x$ form an $\pkpA$-module
$L_i(o)$. ($a$ indexes the prime $p_a$, $u$ the jump.) The \emph{standard orders} are
$\mathrm{std}_l(a,u)=0$ for $a<l$ and $-u$ for $a\ge l$; we also use
\[
\rho_t(a,u)=t\cdot[a=0]-u,\qquad \mathrm{sh}_{l,n}(o)(a,u)=\begin{cases}n&a<l\\ o(a-l,u)&a\ge
l,\end{cases}
\]
(``raised at $0$'' and ``shifted by $l$''). An element $x\in\mathcal N$ is \emph{deep for $o'$} if
$gx-x\in L_i(o')$ for all $g\in\Gamma'$.
\end{definition}

\paragraph*{Sides, pairings, points}

\begin{definition}[Side]\label{pkp:def-side}
A \emph{side} $S$ is a lattice $Y$ with jump sets $U(i)\ni0$, a finite free $\auO$-module $C_S$ with
trivial action, and an $\auO$-linear map $\alpha\colon C_S\to Z^1_f(\Gamma,Y)$ (crystalline
cocycles, {\cite[Definition~5.1]{P55}}). Its \emph{extension} is $X_S=Y\oplus C_S$ with
$g(y,c)=(gy+\alpha(c)(g),c)$, a lattice with the same jump sets (its jumps are those of $Y$ and $0$).
Let $\mathrm{pr}\colon\pkpAmb(X_S)\to\pkpAmb(C_S)$ and $\mathrm{fst}\colon\pkpAmb(X_S)\to\pkpAmb(Y)$
be the two projections. For $c\in C_S$ and a letter $i$, the \emph{constant of $c$ at $i$} is
$e_i(1\otimes c)$; a vector $x$ \emph{lies over $\ell^ec$ at $i$} if $x\in\pkpAmb(X_S)_i$ and
$\mathrm{pr}(x)=\ell^ee_i(1\otimes c)$.
\end{definition}

\begin{definition}[Pairing modulo $\ell^m$]\label{pkp:def-pairing}
For sides $S,S'$ with lattices $Y,Y'$, a \emph{pairing modulo $\ell^m$} is a $\Z_\ell$-bilinear,
$\Gamma$-equivariant map $P\colon Y\times Y'\to\Z/\ell^m(1)$ which is the reduction (through a
$\Z_\ell$-linear $\auO\to\Z_\ell$) of an $\auO$-bilinear $\psi\colon Y\times Y'\to\auO(1)$ whose
failure of equivariance $g\psi-\psi=\ell^m c_P(g)$ is a crystalline cocycle $c_P$ of
$\operatorname{Hom}(Y\otimes Y',\auO(1))$. Put $\pkpM=\pkpW\otimes_{\Z_\ell}\Z/\ell^m(1)$, with
$\Z\times\Gamma$ acting diagonally ($\varphi$ trivially on $\Z/\ell^m(1)$), and let $P$ also denote
the $\pkpW$-bilinear extension $\pkpAmb(Y)\times\pkpAmb(Y')\to\pkpM$; it commutes with $\Z\times\Gamma$.
The \emph{small part} is $\pkpsm=\pkpm\cdot(1\otimes\epsilon_m)\subset\pkpM$, where $\epsilon_m$ is a generator of $\Z/\ell^m(1)$; it does not depend on the choice of $\epsilon_m$, since the generators differ by units of $\Z/\ell^m$. The \emph{Tate-dual side} of
$P$ has lattice $\operatorname{Hom}(Y\otimes Y',\auO(1))$, jump sets $\{0\}\cup\{-u-u'-1:u\in U(i),
u'\in U'(i)\}$, $C=\auO$ and $\alpha(1)=c_P$.
\end{definition}

\begin{definition}[Safe pairs, moves, paths]\label{pkp:def-safe}
Let $\kappa\ge0$. A pair $(w,y)\in\pkpAmb(X_S)\times\pkpAmb(X_{S'})$ is \emph{$\kappa$-safe} if
$\ell^\kappa P(\mathrm{fst}\,w,\mathrm{fst}\,y)\in\pkpsm$. The \emph{move of the first side by $y$
with $x'$ at rest} is safe if $(y,gx'-x')$ is safe for all $g\in\Gamma'$; the \emph{move of the second
side by $y'$ with $x$ at rest} is safe if $(gx-x,gy')$ is safe for all $g\in\Gamma'$. A \emph{safe
path} from $(y,y')$ to $(x,x')$ through $(w,w')$ consists of four safe moves:
\begin{center}
first side by $y-w$ with $y'$ at rest;\quad second side by $y'-w'$ with $w$ at rest;\\
first side by $w-x$ with $w'$ at rest;\quad second side by $w'-x'$ with $x$ at rest.
\end{center}
\end{definition}

\begin{definition}[Scaled cup coboundary]\label{pkp:def-scc}
For $c\in C_S$, $c'\in C_{S'}$ and $N\ge0$, we say $\mathrm{SCC}_N(c,c')$ if the continuous 2-cocycle
$(g,h)\mapsto\zeta_{\ell^m}^{\,P(\alpha(c)(g),\,g\,\ell^N\alpha'(c')(h))}$ of $\Gamma$ with values in
the discrete module $\auQlb^{\times}$ is the coboundary of a continuous 1-cochain. (Here
$\Z/\ell^m(1)\cong\mu_{\ell^m}$, $x\mapsto\zeta_{\ell^m}^x$.) In other words, $\ell^N$ kills the cup
product of the two classes in $H^2(\Gamma,\auQlb^\times)$ read through $P$.
\end{definition}

\begin{theorem}[Pairing lemma]\label{pkp:pairing}
Let $c\in C_S$, $c'\in C_{S'}$, $e_1,e_2,\kappa\ge0$, and let $x,w\in\pkpAmb(X_S)$ satisfy
$\mathrm{pr}(x)=\mathrm{pr}(w)=\ell^{e_1}(1\otimes c)$ and $x',w'\in\pkpAmb(X_{S'})$ satisfy
$\mathrm{pr}(x')=\mathrm{pr}(w')=\ell^{e_2}(1\otimes c')$. Assume:
\begin{enumerate}
\item[(a)] for every $h\in\Gamma'$ the move of the second side by $hx'-x'$ with $x$ at rest is
$\kappa$-safe;
\item[(b)] there is a $\kappa$-safe path from $(\varphi x,\varphi x')$ through $(w,w')$ to $(x,x')$.
\end{enumerate}
Then $\mathrm{SCC}_N(c,c')$ for every $N\ge e_1+e_2+\kappa$.
\end{theorem}

\begin{corollary}[Letter by letter]\label{pkp:letters}
Suppose that for every letter $i$ vectors $y_i,w_i,x_i$ lie over $\ell^ec$ at $i$ and
$y'_i,w'_i,x'_i$ lie over $\ell^{e'}c'$ at $i$, that $\varphi\sum_ix_i=\sum_iy_i$ and
$\varphi\sum_ix'_i=\sum_iy'_i$, that for all $i$ and $h\in\Gamma'$ the move of the second side by
$hx'_i-x'_i$ with $x_i$ at rest is $\kappa$-safe, and that for every $i$ there is a $\kappa$-safe path
from $(y_i,y'_i)$ through $(w_i,w'_i)$ to $(x_i,x'_i)$. Then $\mathrm{SCC}_N(c,c')$ for $N\ge
e+e'+\kappa$.
\end{corollary}

\paragraph*{Points, cokernel killers, the rule and the walk}

\begin{definition}[Points; killing the cokernel]\label{pkp:def-points}
Let $S$ be a side, $c\in C_S$, $i$ a letter, $e\ge0$, $o,o'$ orders. The \emph{points over $\ell^ec$ at
$i$ with orders $o$ and depth $o'$} are the $x\in\mathcal N(X_S)$ that lie over $\ell^ec$ at $i$, have
orders $o$ at $i$ and are deep for $o'$. Put, for $t,\eta,l\ge0$,
\[
\mathrm{dep}_{t,\eta,l}(a,u)=-u+t[a=0]+\eta[a<l],\qquad
\mathrm{stop}_{\eta,l}(a,u)=\max\bigl(\mathrm{std}_l(a,u),\,-u-\eta[a<l]\bigr),
\]
\[
\mathrm{rule}_{l,t}(a,s)=\begin{cases}1&a<l,\\ t-s&a=l,\\ -s&a>l.\end{cases}
\]
\emph{Standard points} (for $t,\eta,l$) are points with orders $\rho_{-\eta}$ and depth
$\mathrm{dep}_{t,\eta,l}$; \emph{stops} (for $\eta,l$) are points with orders and depth
$\mathrm{stop}_{\eta,l}$. ``$x\in L_i(o)$ lies over $0$'' means $\mathrm{pr}(x)=0$; write
$L^0_i(o)$. We say that \emph{$\ell^\kappa E^t$ kills the cokernel at $i$} for the side $S$ if for
every $c\in C_S$ there is $n\in(X_S)_i(\mathrm{std}_0)$ (that is, $n$ in the letter-$i$ part of
$\pkpA\mathcal N(X_S)$ with orders $\mathrm{std}_0$; $n$ need not lie in $\mathcal N(X_S)$, but it can
be moved into it, Lemma~\ref{pkt:lift}) and $\mathrm{pr}(n)=\ell^\kappa E^t
e_i(1\otimes c)$.
\end{definition}

\begin{proposition}[The rule]\label{pkp:rule}
Let $P$ pair $S$, $S'$, let $l,t,\kappa\ge0$ and suppose $\ell^\kappa E^t$ kills the cokernel of
the Tate-dual side at the letter $i^{(-l)}$. Let $y\in L^0_i(o)$ (for $X_S$) and let $x'$ be a point
of $S'$ over $\ell^{e'}c'$ at $i$ with depth $d'$ (any orders). If
$\mathrm{rule}_{l,t}(a,u+u')\le o(a,u)+d'(a,u')$ for all $a\ge0$, $u\in U(i^{(-a)})$, $u'\in
U'(i^{(-a)})$, then the move of the first side by $y$ with $x'$ at rest is $\kappa$-safe. Symmetrically,
for $y'\in L^0_i(o')$ (for $X_{S'}$) and a point $x$ of $S$ over $\ell^ec$ at $i$ with depth $d$, if
$\mathrm{rule}_{l,t}(a,u+u')\le d(a,u)+o'(a,u')$ for all such $a,u,u'$, the move of the second side
by $y'$ with $x$ at rest is $\kappa$-safe.
\end{proposition}

\begin{proposition}[The walk at one cut]\label{pkp:cut}
Let $0<l$, $0<l_0$, let $\ell^\kappa E^t$ kill the cokernel of the Tate-dual side at $i$ and
$\ell^\kappa E^{t_0}$ at $i^{(-l)}$. Let $x,x'$ be standard points (for $t,\eta,l$) of $S,S'$ over
$\ell^ec,\ell^{e'}c'$ at $i$, and $x_0,x'_0$ standard points (for $t_0,\eta,l_0$) over the same at
$i^{(-l)}$. Suppose there is a stop $q$ (for $\eta,l$) of $S$ over $\ell^ec$ at $i$ such that the move
of the second side by $\varphi^lx'_0-x'$ with $q$ at rest is $\kappa$-safe, or a stop $q'$ of $S'$
over $\ell^{e'}c'$ at $i$ such that the move of the first side by $\varphi^lx_0-x$ with $q'$ at rest
is $\kappa$-safe. Then there are $w,w'$ lying over $\ell^ec,\ell^{e'}c'$ at $i$ and a $\kappa$-safe
path from $(\varphi^lx_0,\varphi^lx'_0)$ through $(w,w')$ to $(x,x')$. Moreover, for every $d\ge0$,
every $g\in\Gamma'$ and standard points $x,x'$ (for $t,\eta,l$) at $i^{(-d)}$, with $\ell^\kappa E^t$
killing the cokernel at $i^{(-d)}$, the move of the second side by
$g\varphi^dx'-\varphi^dx'$ with $\varphi^dx$ at rest is $\kappa$-safe.
\end{proposition}

The two inputs from the theory of Breuil--Kisin modules used in the proof of the rule are:

\begin{lemma}[Orders add]\label{pkp:orders-add}
Let $L,L',L''$ be lattices with jump sets $U,U',U''$ and $b\colon L\times L'\to L''$ a
$\Gamma'$-equivariant $\auO$-bilinear map, with $\pkpW$-bilinear extension $\tilde b$. If $w\in L_i(o)$,
$y\in L'_i(o')$ and $o''(a,v)\le o(a,u)+o'(a,u')$ whenever $u+u'\le v$ (for all $a\ge0$,
$u\in U(i^{(-a)})$, $u'\in U'(i^{(-a)})$, $v\in U''(i^{(-a)})$), then $\tilde b(w,y)\in L''_i(o'')$.
In particular, for $L''=L\otimes_\auO L'$ with jump sets $U(i)+U'(i)$ and $b=\otimes$,
$w\otimes y\in(L\otimes L')_i(o'')$.
\end{lemma}

\begin{lemma}[Integrality in $\auO(1)$]\label{pkp:tau}
(i) For the lattice $\auO(1)$ (jump set $\{-1\}$), an element $\xi$ with orders $\mathrm{std}_0$ at $i$
satisfies $\tau\xi\in\pkpA\otimes_{\Z_\ell}\auO(1)$. (ii) Moreover $z\in\pkpm$, and for $\xi\in
\pkpA\otimes\auO(1)$ the image of $z\xi$ in $\pkpM$ lies in $\pkpsm$. (iii) If $\ell^\kappa E^t$ kills the
cokernel at $i$ for a side, then for every $c$ there is $\nu\in(X_S)_i(\rho_{-t})$ with
$\mathrm{pr}(\nu)=\ell^\kappa e_i(1\otimes c)$ (Lemma~\ref{pkm:tau-iii}; by Proposition~\ref{pkt:I1}
$\nu$ can be taken to be a point, with depth $\rho_{-t}$). For every lattice with jump sets: (iv) if $x\in\mathcal N$ lies at $i$,
is deep for $o'$ and $\mathrm{pr}(x)$ is $\Gamma'$-invariant (for an extension), then for $g\in\Gamma'$
the vector $(z\tau)^{-1}(gx-x)$ (note $z\tau\in\pkpW^\times$) has orders $o'$ at $i$ and lies over
$0$; (v) if $x\in L_{i^{(-l)}}(o)$ then $\varphi^lx\in L_i(\mathrm{sh}_{l,0}(o))$; if moreover $x$ is a point at
$i^{(-l)}$ with orders $o$ and depth $o'$, then $\varphi^lx$ is a point at $i$ (over the same constant)
with orders $\mathrm{sh}_{l,0}(o)$ and depth $\mathrm{sh}_{l,1}(o')$.
\end{lemma}

The last input concerns the data and is stated in \S\ref{pki:ss-arch} as
Proposition~\ref{pki:chart}.

\begin{proof}[Proof of Lemma~\ref{pkp:chart-exists}]
Let $\pi$ be a uniformiser of $K$ (take $\pi=-\ell$ if $K/\Q_\ell$ is unramified, which is the
case when $j(\Gamma)$ contains inertia), choose $\pi_n\in\auQlb$ with $\pi_0=\pi$, $\pi_{n+1}^\ell=\pi_n$,
and let $\varpi^\flat=(\pi_n\bmod\ell)_n\in\mathfrak m^\flat$; its untilt is $\pi$. Put
$\Gamma'=\Gamma$, so $j(\Gamma')=G_K$. $G_K$ fixes $\pi=\theta(z)$, and fixes $\Lambda$ because it acts on
Witt coordinates through $O^\flat$ and fixes the Teichm\"uller lifts of $\mathbb F_{\ell^f}$ (the
residue field of $K$). Conversely an element fixing $\Lambda$ and $\pi$ fixes $K_0=\Lambda[1/\ell]$ and
$\pi$, hence $K=K_0(\pi)$ ($K/K_0$ is totally ramified, generated by $\pi$). Write $\pi^e=-\ell v$ with
$v\in O_K^\times$; $\theta$ maps $\mathfrak S$ onto $O_K$ ($\Lambda\to W(\mathbb F_{\ell^f})$, $z\mapsto\pi$),
so choose $A\in\mathfrak S$ with $\theta(A)=v^{-1}$; $A$ is a unit of the local ring $\mathfrak S$, and
$\theta(z^eA+\ell)=0$. That $E=z^eA+\ell$ generates $\ker\theta$ is Corollary~\ref{pkn:E-gen}, whose proof uses of the present one only the equality $\theta(E)=0$ just obtained.
If $e=1$ and $\pi=-\ell$, $v=1$ and we may take $A=1$, $E=z+\ell$.

The last property. Fix $l\ge0$ and $\zeta_{\ell^n}$ with $\varepsilon=(\zeta_{\ell^n}\bmod\ell)_n$.
Let $X_n=\{g\in G_K:g\pi_n=\zeta_{\ell^n}^l\pi_n\}$. Each $X_n$ is closed (a stabiliser condition on
one algebraic element). $X_{n+1}\subseteq X_n$: raise $g\pi_{n+1}=\zeta_{\ell^{n+1}}^l\pi_{n+1}$ to the
$\ell$-th power. $X_n\ne\emptyset$: $X^{\ell^n}-\pi$ is Eisenstein over $O_K$, hence irreducible over
$K$, so $G_K$ is transitive on its roots, among which are $\pi_n$ and $\zeta_{\ell^n}^l\pi_n$. As
$G_K$ is compact, $\bigcap_nX_n$ contains some $g$. Then $g$ sends every coordinate of $\varpi^\flat$ to
the corresponding coordinate of $\varepsilon^l\varpi^\flat$, so $g\varpi^\flat=\varepsilon^l\varpi^\flat$
and $gz=[\varepsilon]^lz$ (Teichm\"uller lifts are multiplicative and equivariant).
\end{proof}

Lemma~\ref{pkp:tau}(ii) is Lemma~\ref{pkp:AS}(vii).

\subsubsection{Artin--Schreier facts in $\pkpM$}\label{pkp:ss-AS}

\begin{lemma}\label{pkp:AS}
Let $m\ge0$ and $\pkpM$, $\pkpsm$ be as in Definition~\ref{pkp:def-pairing}. Fix a generator $\epsilon_m$ of $\Z/\ell^m(1)$ and write $\iota(x)=x\otimes\epsilon_m$
for $x\in\pkpW$; $\iota$ is onto and $\iota(x)=\iota(y)$ iff $x\equiv y\pmod{\ell^m\pkpW}$.
\begin{enumerate}
\item[(i)] $\ell^m\pkpM=0$; $\pkpsm$ is stable under $\Gamma$.
\item[(ii)] $\varphi-1$ is onto $\pkpM$.
\item[(iii)] If $(\varphi-1)d\in\pkpsm$, then $d-1\otimes k\in\pkpsm$ for some $k\in\Z/\ell^m(1)$.
\item[(iv)] If $1\otimes k\in\pkpsm$ with $k\in\Z/\ell^m(1)$, then $k=0$.
\item[(v)] Let $\rho$ be a $\Z_\ell[\Gamma]$-module with $gy-y\in\ell^m\rho$ for $g$ near $1$ (depending
on $y$), and $L\colon\pkpW\otimes\rho\to\pkpM$ a $\pkpW$-linear map. For each $v$, $L(gv-v)\in\pkpsm$
for all $g$ in a neighbourhood of $1$.
\item[(vi)] (Used in (v).) For $x,d\in\pkpW$, the first $m$ Witt coordinates of $(gx-x)d$ lie in
$\mathfrak m^\flat$ for $g$ near $1$.
\item[(vii)] $z\in\pkpm$, and the image in $\pkpM$ of $z\cdot(\pkpA\otimes_{\Z_\ell}\auO(1))$ under any
$\Z_\ell$-linear reduction $\auO(1)\to\Z/\ell^m(1)$ lies in $\pkpsm$.
\end{enumerate}
\end{lemma}

\begin{proof}
Since $\Z/\ell^m(1)$ is free of rank one over $\Z/\ell^m$, $\pkpM\cong\pkpW/\ell^m$ as a
$\pkpW$-module via $\iota$, which gives the first assertions. $\varphi$ acts by $\iota(x)\mapsto
\iota(\varphi x)$ and $g\in\Gamma$ by $\iota(x)\mapsto\chi(g)\iota(gx)$. Also $\iota(x)\in\pkpsm$ iff
$x\equiv a\pmod{\ell^m}$ for some $a\in\pkpm$.

(i) The first claim is clear. $g$ acts on Witt coordinates through the automorphism $g$ of $O^\flat$,
which preserves $\mathfrak m^\flat$; so $g\pkpm=\pkpm$ and $g(a\otimes\epsilon_m)=\chi(g)ga\otimes\epsilon_m$.

(ii) Since $C^\flat$ is algebraically closed, every $X^\ell-X-c$ ($c\in C^\flat$) has a root, and
by induction on the Witt length (the $n$-th coordinate of $\varphi y-y$ is $y_n^\ell-y_n$ plus a
polynomial in $y_0,\dots,y_{n-1}$) $\varphi-1$ is onto $\pkpW$; apply $\iota$.

(iii) Write $d=\iota(x)$; by hypothesis $\varphi x-x=a+\ell^mw$ with $a\in\pkpm$, $w\in\pkpW$. As
$O^\flat$ is integrally closed in the algebraically closed $C^\flat$, the roots of $X^\ell-X-c$ for
$c\in O^\flat$ lie in $O^\flat$, so as in (ii) $a=\varphi b-b$ with $b\in\pkpA$, and $w=\varphi
y-y$ with $y\in\pkpW$. The image $\bar b$ of $b$ in $W(\overline{\mathbb F}_\ell)$ satisfies
$\varphi\bar b=\bar b$ (as $\bar a=0$), so $\bar b=c\in\Z_\ell$ and $b-c\in\pkpm$. Then
$x-\ell^my-(b-c)$ is fixed by $\varphi$, hence lies in $W(\mathbb F_\ell)=\Z_\ell$ (the
$\varphi$-fixed points of $W(C^\flat)$ are $W((C^\flat)^{x^\ell=x})$); call it $c'$. Then
$d-1\otimes c'=\iota(\ell^my+(b-c))\in\pkpsm$.

(iv) Let $k$ be the class of $c\in\Z_\ell$ and suppose $c\equiv a\pmod{\ell^m\pkpW}$ with
$a\in\pkpm$. Congruence modulo $\ell^m$ in $W(C^\flat)$ means equality of the first $m$ Witt
coordinates. Those of $c$ lie in $\mathbb F_\ell$, those of $a$ in $\mathfrak m^\flat$; so they
vanish, i.e.\ $\ell^m\mid c$ in $W(\mathbb F_\ell)=\Z_\ell$, $k=0$.

(vi) The Witt coordinates of $gx-x$ are universal polynomials in those of $gx$ and $x$ vanishing on
the diagonal, and $g\mapsto gx_n$ is continuous for the valuation topology of $C^\flat$; so the
first $m$ coordinates of $gx-x$ tend to $0$, and so do those of $(gx-x)d$ (again universal
polynomials, vanishing when the first factor is $0$). Hence they lie in $\mathfrak m^\flat$ for $g$
near $1$.

(v) Both sides are additive in $v$, so let $v=x\otimes y$. Write $L(1\otimes y)=\iota(d)$; then
$L(x'\otimes y)=\iota(x'd)$. For $g$ near $1$, $gy=y+\ell^mw$, and $L(g x\otimes \ell^mw)=0$ by (i),
so $L(gv-v)=\iota((gx-x)d)\in\pkpsm$ by (vi).

(vii) $z=[\varpi^\flat]$ has coordinates $(\varpi^\flat,0,0,\dots)$ in $\mathfrak m^\flat$. An element of
$\pkpA\otimes\auO(1)$ maps to some $\iota(b)$ with $b\in\pkpA$, and $z b\in\pkpm$ since $\pkpm$ is an
ideal.
\end{proof}

\subsubsection{Proof of the pairing lemma}\label{pkp:ss-pairing}

Throughout, $\langle v,v'\rangle=\ell^\kappa P(\mathrm{fst}\,v,\mathrm{fst}\,v')\in\pkpM$ for
$v\in\pkpAmb(X_S)$, $v'\in\pkpAmb(X_{S'})$; it is $\pkpW$-bilinear. Put $A_g=\ell^{e_1}\otimes
\alpha(c)(g)$ and $A'_g=\ell^{e_2}\otimes\alpha'(c')(g)$ in $\pkpAmb(Y)$, $\pkpAmb(Y')$. By the
definition of the extension, for $v$ with $\mathrm{pr}(v)=\ell^{e_1}(1\otimes c)$ and $g\in\Gamma$,
\begin{equation}\label{pkp:eq-fst}
\mathrm{fst}(gv)=g\,\mathrm{fst}(v)+A_g,\qquad \mathrm{fst}(\varphi v)=\varphi\,\mathrm{fst}(v),
\end{equation}
and $\mathrm{fst}(gv)=g\,\mathrm{fst}(v)$ if $\mathrm{pr}(v)=0$; similarly on the second side.
$\mathrm{pr}$ commutes with $\Z\times\Gamma$, and $\varphi(1\otimes c)=1\otimes c$, $g(1\otimes
c)=1\otimes c$. Hence $\mathrm{pr}(\varphi x')=\mathrm{pr}(x')$, so $\varphi x'-w'$, $w'-x'$, $gx'-x'$
lie over $0$. The map $g\mapsto A_g$ is a cocycle: $A_{gh}=gA_h+A_g$.

\begin{proof}[Proof of Theorem~\ref{pkp:pairing}]
\emph{Step 0 (the cochain).} By Lemma~\ref{pkp:AS}(ii) choose $y\in\pkpM$ with
$(\varphi-1)y=\langle w,\varphi x'-w'\rangle+\langle x,w'-x'\rangle$. For $g\in\Gamma$ put
\[
f(g)=\ell^\kappa P(A_g,\,g\,\mathrm{fst}\,x')+(g-1)y-\langle x,gx'-x'\rangle .
\]

\emph{Step 1: $(\varphi-1)f(g)\in\pkpsm$ for $g\in\Gamma'$.} Since $\varphi$ commutes with $g$,
fixes $A_g$ and commutes with $P$ and $\mathrm{fst}$,
\[
(\varphi-1)f(g)=\ell^\kappa P(A_g,g\,\mathrm{fst}(\varphi x'-x'))+(g-1)\bigl(\langle
w,a'\rangle+\langle x,b'\rangle\bigr)-\langle\varphi x,g\varphi x'-\varphi x'\rangle+\langle
x,gx'-x'\rangle,
\]
with $a'=\varphi x'-w'$, $b'=w'-x'$, which lie over $0$. By \eqref{pkp:eq-fst}, for $v\in\{w,x\}$ and
$v'$ over $0$, $g\langle v,v'\rangle=\langle gv,gv'\rangle-\ell^\kappa P(A_g,g\,\mathrm{fst}\,v')$. As
$a'+b'=\varphi x'-x'$, the $A_g$-terms cancel and
\[
(\varphi-1)f(g)=\langle gw,ga'\rangle-\langle w,a'\rangle+\langle gx,gb'\rangle-\langle
x,b'\rangle-\langle\varphi x,g\varphi x'-\varphi x'\rangle+\langle x,gx'-x'\rangle .
\]
Expanding $\langle gv,gv'\rangle-\langle v,v'\rangle=\langle gv-v,gv'\rangle+\langle v,gv'-v'\rangle$
and using $b'+x'=w'$, $a'=\varphi x'-w'$, this equals
\[
\langle gw-w,\,g(\varphi x'-w')\rangle-\langle\varphi x-w,\,g\varphi x'-\varphi x'\rangle-\langle
w-x,\,gw'-w'\rangle+\langle gx-x,\,g(w'-x')\rangle .
\]
The four terms are the values $\langle\cdot,\cdot\rangle$ on the pairs tested by the four moves of the path
(b) (second move, first move, third move, fourth move, in this order), each in $\pkpsm$ by
Definition~\ref{pkp:def-safe}.

\emph{Step 2.} By Lemma~\ref{pkp:AS}(iii), for each $g\in\Gamma'$ there is $k(g)\in\Z/\ell^m(1)$ with
$\delta(g)=f(g)-1\otimes k(g)\in\pkpsm$.

\emph{Step 3: the cup product on $\Gamma'$.} For $h\in\Gamma$ put $F_h=\mathrm{fst}(hx'-x')=(h-1)
\mathrm{fst}\,x'+A'_h$, a cocycle. Write $f=T_1+(g-1)y+T_2$ with $T_1(g)=\ell^\kappa P(A_g,g\,
\mathrm{fst}\,x')$ and $T_2(g)=-\ell^\kappa P(\mathrm{fst}\,x,F_g)$. Using the cocycle relations,
\[
gT_1(h)-T_1(gh)+T_1(g)=-\ell^\kappa P(A_g,g(h-1)\mathrm{fst}\,x'),\qquad
gT_2(h)-T_2(gh)+T_2(g)=-\ell^\kappa P((g-1)\mathrm{fst}\,x,gF_h),
\]
and the $y$-part has zero coboundary. On the other hand $\mathrm{fst}(gx-x)=(g-1)\mathrm{fst}\,x+A_g$, so
\[
\langle gx-x,g(hx'-x')\rangle=\ell^\kappa P((g-1)\mathrm{fst}\,x,gF_h)+\ell^\kappa
P(A_g,g(h-1)\mathrm{fst}\,x')+\ell^\kappa P(A_g,gA'_h).
\]
Adding, and noting $\ell^\kappa P(A_g,gA'_h)=1\otimes\ell^{e_1+e_2+\kappa}P(\alpha(c)(g),g\alpha'(c')(h))$,
we get for $g,h\in\Gamma$
\[
1\otimes\ell^{e_1+e_2+\kappa}P(\alpha(c)(g),g\alpha'(c')(h))=\langle
gx-x,g(hx'-x')\rangle+gf(h)-f(gh)+f(g).
\]
For $g,h\in\Gamma'$, substitute $f=1\otimes k+\delta$: the difference
$1\otimes[\ell^{e_1+e_2+\kappa}P(\alpha(c)(g),g\alpha'(c')(h))-(\chi(g)k(h)-k(gh)+k(g))]$ equals
$\langle gx-x,g(hx'-x')\rangle+(g\delta(h)-\delta(gh)+\delta(g))$, which lies in $\pkpsm$ by (a)
(the move by $hx'-x'$ with $x$ at rest, tested at $g$) and Lemma~\ref{pkp:AS}(i). By
Lemma~\ref{pkp:AS}(iv), with $r=\ell^{N-e_1-e_2-\kappa}$,
\begin{equation}\label{pkp:eq-cob}
P(\alpha(c)(g),g\,\ell^N\alpha'(c')(h))=\chi(g)\,rk(h)-rk(gh)+rk(g)\qquad(g,h\in\Gamma').
\end{equation}

\emph{Step 4: $k$ vanishes near $1$.} Cocycles and the action on lattices are continuous, so for $h$
in a neighbourhood $V$ of $1$ (inside the open subgroup $\Gamma'$): $\alpha(c)(h)\in\ell^m Y$,
$\alpha'(c')(h)\in\ell^mY'$, $hy'-y'\in\ell^mY'$ for $y'$ in a fixed finite generating set, and
$\chi(h)\equiv1\pmod{\ell^m}$. By Lemma~\ref{pkp:AS}(i), terms containing $A_h$ or $A'_h$ vanish, so
$f(h)=(h-1)y-\ell^\kappa P(\mathrm{fst}\,x,(h-1)\mathrm{fst}\,x')$. By Lemma~\ref{pkp:AS}(v)
(applied to $\rho=\Z/\ell^m(1)$, $L=\mathrm{id}$, and to $\rho=Y'$, $L=\ell^\kappa
P(\mathrm{fst}\,x,\cdot)$), after shrinking $V$, $f(h)\in\pkpsm$. Then $1\otimes k(h)\in\pkpsm$ and
$k(h)=0$ by Lemma~\ref{pkp:AS}(iv).

\emph{Step 5: from $\Gamma'$ to $\Gamma$.} Let $b(g,h)=P(\alpha(c)(g),g\ell^N\alpha'(c')(h))$, a
2-cocycle of $\Gamma$ with values in $\Z/\ell^m(1)$ (the cup product of two 1-cocycles through an
equivariant pairing). It is locally constant: near any $(g,h)$, the changes of $\alpha(c)$,
$\alpha'(c')$ and of the action are multiples of $\ell^m$, by continuity and the cocycle relations.
Read it in $\auQlb^\times$ through $\Z/\ell^m(1)\cong\mu_{\ell^m}$ and let $\beta$ be its class in the
continuous cohomology $H^2(\Gamma,\auQlb^\times)$ (discrete coefficients). The cochain $rk$ on
$\Gamma'$ is locally constant: for $h\in V$, $rk(h)=0$ and $b(g,h)=0$ (as $\alpha'(c')(h)\in\ell^mY'$),
so \eqref{pkp:eq-cob} gives $rk(gh)=rk(g)$. Hence $\beta|_{\Gamma'}=0$. Since $\Gamma'$ is open of
finite index, $\mathrm{cor}\circ\mathrm{res}=[\Gamma:\Gamma']$ on $H^2(\Gamma,-)$, so
$[\Gamma:\Gamma']\beta=0$; also $\ell^m\beta=0$ since $\ell^m b=0$. As $[\Gamma:\Gamma']$ is prime to
$\ell$, $\beta=0$, which is $\mathrm{SCC}_N(c,c')$.
\end{proof}

\begin{proof}[Proof of Corollary~\ref{pkp:letters}]
\emph{Letters do not interact.} If $v\in\pkpAmb(Y)_i$ and $v'\in\pkpAmb(Y')_{i'}$ with $i\ne i'$,
then $P(v,v')=0$: since $\psi$ is $\auO$-bilinear, $P(\zeta v,v')=P(v,\zeta v')$, i.e.\
$[i]P(v,v')=[i']P(v,v')$, and $[i]-[i']$ is a unit. The projection $\mathrm{fst}$ preserves letter
parts, and $\Gamma'$ preserves them. Hence, if $(v_i,v'_i)$ are $\kappa$-safe pairs with $v_i,v'_i$
at the letter $i$, then $(\sum_iv_i,\sum_iv'_i)$ is $\kappa$-safe (the cross terms vanish and
$\pkpsm$ is a group). Applying this to each move of the paths and to the moves in the cocycle
hypothesis, with $x=\sum x_i$, $w=\sum w_i$, $x'=\sum x'_i$, $w'=\sum w'_i$, we obtain (a) and
(b) of Theorem~\ref{pkp:pairing} (the path starts at $(\sum y_i,\sum y'_i)=(\varphi x,\varphi
x')$). Finally $\mathrm{pr}(x)=\ell^e\sum_ie_i(1\otimes c)=\ell^e(1\otimes c)$, and likewise for
$w,x',w'$. Theorem~\ref{pkp:pairing} applies with $e_1=e$, $e_2=e'$.
\end{proof}

\subsubsection{The rule and the walk at one cut}\label{pkp:ss-walk}

\begin{lemma}[Extension form]\label{pkp:extform}
Let $T=\operatorname{Hom}(Y\otimes Y',\auO(1))$ and $X_T=T\oplus\auO$ the extension of the Tate-dual
side, $g(\phi,a)=(g\phi+a\,c_P(g),a)$. The $\auO$-bilinear form $\epsilon((\phi,a),v)=\ell^m\phi(v)+a
\psi(v)$ on $X_T\times(Y\otimes Y')$ with values in $\auO(1)$ is $\Gamma$-equivariant, and its
reduction modulo $\ell^m$ (through $\auO\to\Z_\ell$) is $((\phi,a),v)\mapsto P(av)$, with $a$ acting on the $\auO$-module $Y\otimes Y'$; on
$\pkpAmb$ the same holds with $a\in\pkpW\otimes\auO$.
\end{lemma}

\begin{proof}
With $(g\phi)(gv)=g(\phi(v))$ and $(g\psi)(gv)=g(\psi(v))$, $g\psi=\psi+\ell^mc_P(g)$:
$\epsilon(g(\phi,a),gv)=\ell^mg(\phi(v))+a\ell^mc_P(g)(gv)+a\psi(gv)=\ell^mg(\phi(v))+a\,g(\psi(v))$.
The reduction is clear.
\end{proof}

\begin{proof}[Proof of Proposition~\ref{pkp:rule}]
We prove the first assertion; the second is symmetric (see the end). Fix $g\in\Gamma'$; we must show
$\ell^\kappa P(\mathrm{fst}\,y,\mathrm{fst}(gx'-x'))\in\pkpsm$. By Lemma~\ref{pkp:tau}(iv),
$v'=(z\tau)^{-1}(gx'-x')$ has orders $d'$ at $i$ and lies over $0$; so $gx'-x'=z\tau v'$ and it
suffices to show $z\tau\,\ell^\kappa P(\mathrm{fst}\,y,\mathrm{fst}\,v')\in\pkpsm$ for $y\in L^0_i(o)$,
$v'\in L'^0_i(d')$. By Lemma~\ref{pkp:tau}(iii) (Lemma~\ref{pkm:tau-iii}) there is $\nu$ in the Tate-dual extension, lying
over $\ell^\kappa\cdot1$ at $i^{(-l)}$, with orders $\rho_{-t}$; by Lemma~\ref{pkp:tau}(v) (first
assertion, Lemma~\ref{pkt:tau-v}), $\varphi^l\nu$ lies
over $\ell^\kappa\cdot1$ at $i$ with orders $o_\nu=\mathrm{sh}_{l,0}(\rho_{-t})$. One checks from the
definitions that
\begin{equation}\label{pkp:eq-ruleid}
o_\nu(a,-u-u'-1)+\mathrm{rule}_{l,t}(a,u+u')=1
\end{equation}
(with $s=u+u'$: for $a<l$ the terms are $0$ and $1$; for $a=l$ they are $\rho_{-t}(0,-s-1)=s+1-t$ and
$t-s$; for $a>l$ they are $\rho_{-t}(a-l,-s-1)=s+1$ and $-s$). Since $y$, $v'$ lie over $0$ they lie in $\pkpAmb(Y)$,
$\pkpAmb(Y')$, and by Lemma~\ref{pkt:sub} they have orders $o$, $d'$ at $i$ there. By Lemma~\ref{pkp:orders-add} (that is, Lemma~\ref{pkt:add}(ii)), $y\otimes v'$ has orders $o''(a,v)=1-o_\nu(a,-v-1)$ at
$i$ in $Y\otimes Y'$: indeed if $u+u'\le v$ then, $o_\nu(a,\cdot)$ being non-increasing,
$1-o_\nu(a,-v-1)\le1-o_\nu(a,-u-u'-1)=\mathrm{rule}_{l,t}(a,u+u')\le o(a,u)+d'(a,u')$. Pair
$\varphi^l\nu$ with $y\otimes v'$ by $\epsilon$ (Lemma~\ref{pkp:extform}); by
Lemma~\ref{pkp:orders-add} again (for $b=\epsilon$, $L''=\auO(1)$, $U''=\{-1\}$), the value $\xi\in\pkpAmb(\auO(1))$ has orders at least
$o_\nu(a,u)+o''(a,u'')\ge1$ on the jump $-1$ whenever $u+u''\le-1$, that is, orders
$\mathrm{std}_0$ (which is $1$ on the jump $-1$). By Lemma~\ref{pkp:tau}, $\tau\xi$ is integral and the
image of $z\tau\xi$ in $\pkpM$ lies in $\pkpsm$. By Lemma~\ref{pkp:extform}, that image is
$z\tau\,P(\mathrm{pr}(\varphi^l\nu)(y\otimes v'))$ with $\mathrm{pr}(\varphi^l\nu)=\ell^\kappa e_i(1\otimes
1)$; since $y$ lies at the letter $i$, $e_iy=y$, so this is $z\tau\,\ell^\kappa P(e_i(y\otimes v'))=z\tau\,\ell^\kappa P(y,v')$. So $z\tau\ell^\kappa P(y,v')\in\pkpsm$.

For the second assertion: for $g\in\Gamma'$ we need $(gx-x,gy')$ safe. Since $y'$ lies over $0$ and
$L'_i(o')$ is $\Gamma'$-stable (Lemma~\ref{pkm:inputs}(G)), $gy'\in L'^0_i(o')$; and $gx-x=z\tau v$ with $v\in
L^0_i(d)$ by Lemma~\ref{pkp:tau}(iv). The argument above, with the roles of the sides exchanged,
gives the claim.
\end{proof}

\begin{lemma}[Arithmetic of the walk]\label{pkp:walk-arith}
Let $l,l_0>0$, $t,t_0,\eta\ge0$, $a\ge0$, $u,v\in\Z$. Then
\begin{enumerate}
\item[(i)] $\mathrm{rule}_{l,t_0}(a,u+v)\le\min\bigl(\mathrm{stop}_{\eta,l},\mathrm{sh}_{l,0}(\rho_{-\eta})
\bigr)(a,u)+\mathrm{sh}_{l,1}(\mathrm{dep}_{t_0,\eta,l_0})(a,v)$;
\item[(ii)] $\mathrm{rule}_{0,t}(a,u+v)\le\min\bigl(\rho_{-\eta},\mathrm{stop}_{\eta,l}\bigr)(a,u)+
\mathrm{dep}_{t,\eta,l}(a,v)$;
\item[(iii)] for $d\ge0$: $\mathrm{rule}_{d,t}(a,u+v)\le\mathrm{sh}_{d,1}(\mathrm{dep}_{t,\eta,l})(a,u)+
\mathrm{sh}_{d,1}(\mathrm{dep}_{t,\eta,l})(a,v)$.
\end{enumerate}
\end{lemma}

\begin{proof}
(i) If $a<l$: $\mathrm{stop}=\max(0,-u-\eta)\ge0$ and $\mathrm{sh}_{l,0}(\rho_{-\eta})=0$, so the minimum
is $0$, and the last term is $1$; both sides are $1$. If $a=l$: $\mathrm{stop}=-u$,
$\rho_{-\eta}(0,u)=-\eta-u$, the minimum is $-u-\eta$, and $\mathrm{dep}_{t_0,\eta,l_0}(0,v)=
-v+t_0+\eta$; both sides are $t_0-u-v$. If $a>l$: the minimum is $-u$ and
$\mathrm{dep}_{t_0,\eta,l_0}(a-l,v)\ge-v$; the left side is $-u-v$.
(ii) If $a=0$: $\rho_{-\eta}(0,u)=-\eta-u\le\max(0,-u-\eta)=\mathrm{stop}(0,u)$ (as $l>0$), and
$\mathrm{dep}(0,v)=-v+t+\eta$; both sides are $t-u-v$. If $0<a<l$: the minimum is
$\ge-u-\eta$ and $\mathrm{dep}=-v+\eta$; the left side is $-u-v$. If $a\ge l$: the minimum is $-u$,
$\mathrm{dep}=-v$.
(iii) If $a<d$: $1\le1+1$. If $a=d$: $t-u-v\le(-u+t+\eta[0<l])+(-v+t+\eta[0<l])$. If $a>d$:
$-u-v\le(-u+\eta[\cdot])+(-v+\eta[\cdot])$.
\end{proof}

\begin{proof}[Proof of Proposition~\ref{pkp:cut}]
By Lemma~\ref{pkp:tau}(v), $\varphi^lx_0$ and $\varphi^lx'_0$ are points at $i$ over the same
constants with orders $\mathrm{sh}_{l,0}(\rho_{-\eta})$ and depth $\mathrm{sh}_{l,1}(\mathrm{dep}_{t_0,
\eta,l_0})$. The difference of two points over the same constant at $i$ lies over $0$ and has the
minimum of their orders (the modules $L_i(o)$ grow as $o$ decreases).

\emph{Case of a stop $q$ of $S$.} Take $w=q$, $w'=x'$. First move: $\varphi^lx_0-q\in
L^0_i(\min(\mathrm{stop},\mathrm{sh}_{l,0}\rho_{-\eta}))$ with $\varphi^lx'_0$ at rest; by
Proposition~\ref{pkp:rule} (killing at $i^{(-l)}$, rule for $l,t_0$) and Lemma~\ref{pkp:walk-arith}(i) it
is safe. Second move: by $\varphi^lx'_0-x'$ with $q$ at rest: the hypothesis. Third move: $q-x\in
L^0_i(\min(\rho_{-\eta},\mathrm{stop}))$ with $x'$ (depth $\mathrm{dep}_{t,\eta,l}$) at rest: rule for
$0,t$ (killing at $i$) and Lemma~\ref{pkp:walk-arith}(ii). Fourth move: by $x'-x'=0$, trivially safe.

\emph{Case of a stop $q'$ of $S'$.} Take $w=\varphi^lx_0$, $w'=q'$. The first move is by $0$; the
second, by $\varphi^lx'_0-q'$ with $\varphi^lx_0$ at rest, is safe by the second half of
Proposition~\ref{pkp:rule} and Lemma~\ref{pkp:walk-arith}(i); the third is the hypothesis; the fourth,
by $q'-x'$ with $x$ at rest, by the rule for $0,t$ and Lemma~\ref{pkp:walk-arith}(ii). In both cases
$w,w'$ lie over the right constants at $i$.

\emph{The cocycle moves.} By Lemma~\ref{pkp:tau}(v), $\varphi^dx'$ is a point at $i$ with depth
$\mathrm{sh}_{d,1}(\mathrm{dep}_{t,\eta,l})$, so $g\varphi^dx'-\varphi^dx'$ has those orders and lies over
$0$ ($\mathrm{pr}$ of a point is $\Gamma'$-invariant). With $\varphi^dx$ (same depth) at rest,
Proposition~\ref{pkp:rule} (rule for $d,t$, killing at $i^{(-d)}$) and Lemma~\ref{pkp:walk-arith}(iii)
give the claim.
\end{proof}

\begin{remark}[How the cases use this]\label{pkp:rem-use}
In each case one chooses cuts among the letters (a nonempty set of letters with the distance $l(i)>0$
from each letter back to the previous cut), standard points at the cuts, carries them forward by
$\varphi$ to the letters between cuts, and applies Proposition~\ref{pkp:cut} at each cut (the
non-cut letters use the trivial path) and Corollary~\ref{pkp:letters}; the exponent bound is
$N=e+e'+\kappa$, where $\ell^\kappa E^t$ kills the cokernels. The existence of standard points and
stops at the required depth, and the cokernel killers, are the case-specific inputs.
\end{remark}



\providecommand{\auO}{\mathcal O}
\providecommand{\auQlb}{\overline{\Q}_\ell}
\providecommand{\pkpA}{A_{\mathrm{inf}}}
\providecommand{\pkpW}{W(C^\flat)}
\providecommand{\pkpm}{W(\mathfrak m^\flat)}
\providecommand{\pkpM}{M_m}
\providecommand{\pkpsm}{\mathrm{Sm}_m}
\providecommand{\pkpAmb}{\mathbb A}
\providecommand{\pkmJ}[3]{J^{#1}_{#2}(#3)}
\providecommand{\pkmI}[2]{I_{#1}(#2)}
\providecommand{\pkmv}{v_\ell}

\subsection{The global walk and the cokernel killers}\label{pkm:sec}

This section continues \S\ref{pkp:sec}, whose conventions, definitions and labels we use. In
\S\ref{pkm:ss-state} we state everything that the four cases (Propositions~\ref{pki:generic},
\ref{pki:any}, \ref{pki:tame}, \ref{pki:sparse}) cite beyond \S\ref{pkp:ss-arch}: the global walk,
the stops at plain cuts, the improvement of orders, and the cokernel killers. Proofs follow in
\S\ref{pkm:ss-walk}--\S\ref{pkm:ss-kill}. The inputs from the theory of the modules $\mathcal N$
that these proofs use, beyond Lemmas~\ref{pkp:orders-add} and~\ref{pkp:tau}, are collected in
Lemma~\ref{pkm:inputs}.

\emph{Dependencies.} Proved in full, given Lemma~\ref{pkp:tau}(iv),(v), items (M), (E), (G) of Lemma~\ref{pkm:inputs} and the statements of
\S\ref{pkp:ss-arch} proved there: the global walk (Theorem~\ref{pkm:walk},
Corollary~\ref{pkm:walk-all}), the stops at plain cuts (Lemmas~\ref{pkm:down}, \ref{pkm:plain}), the
improvement of orders (Lemma~\ref{pkm:improved}), and Lemma~\ref{pkp:tau}(iii). Proved in full
\emph{given Lemma~\ref{pkm:inputs}}: the cokernel killers (Propositions~\ref{pkm:kill-small},
\ref{pkm:kill-unram}, \ref{pkm:kill-J} and Lemma~\ref{pkm:zpow}). The items of
Lemma~\ref{pkm:inputs}, Lemma~\ref{pkp:orders-add} and Lemma~\ref{pkp:tau}(i),(iv),(v) are proved in
\S\S\ref{pkn:sec}--\ref{pkt:sec} (see the paragraph at the end of \S\ref{pkm:ss-kill}).

\subsubsection{Statements used by the cases}\label{pkm:ss-state}

Throughout, $S,S'$ are sides with lattices $Y,Y'$ and parameter lattices $C_S,C_{S'}$, $P$ is a
pairing modulo $\ell^m$ of $S$ and $S'$, $T_P$ is its Tate-dual side (Definition~\ref{pkp:def-pairing}),
$c\in C_S$, $c'\in C_{S'}$. The Frobenius $\sigma(i)=i^\ell$ permutes the (finitely many) letters, and
$i^{(-a)}=\sigma^{-a}(i)$. For a side $S$, ``the jumps of $S$ at $i$'' are the elements of its jump
set $U(i)$ (which contains $0$).

\paragraph*{The global walk}

\begin{theorem}[The walk]\label{pkm:walk}
Let $\mathrm{cut}$ be a set of letters and $l,t$ functions from letters to $\Z_{\ge0}$ such that for
every letter $i$, $l(i)$ is the least $a>0$ with $i^{(-a)}\in\mathrm{cut}$. Let $\kappa,e,e',\eta,N\ge0$
and, for every $i\in\mathrm{cut}$, let $x(i)\in\pkpAmb(X_S)$, $x'(i)\in\pkpAmb(X_{S'})$. Assume, for every
$i\in\mathrm{cut}$:
\begin{enumerate}
\item[(K)] $\ell^\kappa E^{t(i)}$ kills the cokernel of $T_P$ at $i$;
\item[(P)] $x(i)$ is a standard point (for $t(i),\eta,l(i)$) of $S$ over $\ell^ec$ at $i$, and $x'(i)$ a
standard point (for $t(i),\eta,l(i)$) of $S'$ over $\ell^{e'}c'$ at $i$;
\item[(St)] either there is a stop $q$ (for $\eta,l(i)$) of $S$ over $\ell^ec$ at $i$ such that the
move of the second side by $\varphi^{l(i)}x'(i^{(-l(i))})-x'(i)$ with $q$ at rest is $\kappa$-safe, or
there is a stop $q'$ (for $\eta,l(i)$) of $S'$ over $\ell^{e'}c'$ at $i$ such that the move of the
first side by $\varphi^{l(i)}x(i^{(-l(i))})-x(i)$ with $q'$ at rest is $\kappa$-safe.
\end{enumerate}
If $N\ge e+e'+\kappa$, then $\mathrm{SCC}_N(c,c')$.
\end{theorem}

(The hypothesis on $l$ forces $\mathrm{cut}\ne\emptyset$.)

\begin{corollary}[Cuts at all letters]\label{pkm:walk-all}
Let $\kappa,e,e',t,\eta,N\ge0$, and for every letter $i$ let $\ell^\kappa E^t$ kill the cokernel of
$T_P$ at $i$, let $x(i)$, $x'(i)$ be standard points (for $t,\eta,1$) of $S,S'$ over $\ell^ec$,
$\ell^{e'}c'$ at $i$, and assume (St) with $l(i)=1$. If $N\ge e+e'+\kappa$, then $\mathrm{SCC}_N(c,c')$.
\end{corollary}

\paragraph*{Improved orders and stops at plain cuts}

For orders $o_{\mathrm{lo}},o_{\mathrm{hi}}$ put
\[
\mathrm{imp}(o_{\mathrm{lo}},o_{\mathrm{hi}})(a,u)=\begin{cases}o_{\mathrm{hi}}(a,u)&\text{if }
o_{\mathrm{lo}}(a,u)>0,\\ \min(o_{\mathrm{hi}}(a,u),0)&\text{otherwise,}\end{cases}
\qquad
\mathrm{pd}_{t,\eta,l,\eta_0}(a,u)=\mathrm{dep}_{t,\eta,l}(a,u)+\eta_0[a=l].
\]

\begin{lemma}[Improved orders]\label{pkm:improved}
Let $L$ be a lattice with jump sets and $i$ a letter. If $x\in L_i(o_{\mathrm{lo}})$ and $gx-x\in
L_i(o_{\mathrm{hi}})$ for every $g\in\Gamma'$, then $x\in L_i(\mathrm{imp}(o_{\mathrm{lo}},o_{\mathrm{hi}}))$.
\end{lemma}

\begin{lemma}[A standard point of a side that is down is a stop]\label{pkm:down}
Let $0<l$, and suppose every jump of $S$ at $i^{(-a)}$, $a<l$, is $\le\eta$. Then every standard point
(for $t,\eta,l$) of $S$ over $\ell^ec$ at $i$ is a stop (for $\eta,l$) of $S$ over $\ell^ec$ at $i$.
\end{lemma}

\begin{lemma}[Stops at plain cuts]\label{pkm:plain}
Let $0<l$, $t,\eta,\eta_0\ge0$, and let $\ell^\kappa E^t$ kill the cokernel of $T_P$ at $i$. Suppose that
every jump of $S$ at $i^{(-a)}$ ($a<l$) is $\le\eta$ and every jump of $S'$ at $i^{(-a)}$ ($a<l$) is
$\ge-\eta$. Let $x$ be a point of $S$ over $\ell^ec$ at $i$ with orders $\rho_{-\eta}$ and depth
$\mathrm{pd}_{t,\eta,l,\eta_0}$; let $x'$ be a point of $S'$ over $\ell^{e'}c'$ at $i$ with orders
$\rho_{-\eta}$ and $x'_0$ a point of $S'$ over $\ell^{e'}c'$ at $i^{(-l)}$ with orders $\rho_{-\eta_0}$
(any depths). Then $x$ is a stop (for $\eta,l$) of $S$ over $\ell^ec$ at $i$, and the move of the
second side by $\varphi^lx'_0-x'$ with $x$ at rest is $\kappa$-safe. The same holds with the roles of
$S$ and $S'$ (and of the two moves) exchanged.
\end{lemma}

In particular, at a plain cut the alternative (St) of Theorem~\ref{pkm:walk} holds as soon as the
standard point of the side that is down has depth $\mathrm{pd}_{t,\eta,l,\eta_0}$ and the points of
the other side have orders $\rho_{-\eta}$ (at $i$) and $\rho_{-\eta_0}$ (at $i^{(-l)}$); this is
Lemma~\ref{pkm:plain}, since a point with depth $\mathrm{pd}_{t,\eta,l,\eta_0}\ge\mathrm{dep}_{t,\eta,l}$
and orders $\rho_{-\eta}$ is standard.

\paragraph*{The lift ideals and the cokernel killers}

\begin{definition}[Lift ideals]\label{pkm:def-J}
Let $S$ be a side, $i$ a letter, $l\ge0$. The ideal $\pkmI li\subseteq\pkpA$ consists of the
$s\in\pkpA$ such that for every $c\in C_S$ there is $x\in(X_S)_i(\mathrm{std}_l)$ (the module of
Definition~\ref{pkp:def-orders} for the extension $X_S$) with $\mathrm{pr}(x)=s\,e_i(1\otimes c)$.
For $r\in\Z_{\ge0}$ put $\pkmJ rli=\pkmI li+r\pkpA$. Thus $\ell^\kappa E^t$ kills the cokernel at $i$
(Definition~\ref{pkp:def-points}; see Remark~\ref{pkm:rem-kill}) if and only if
$\ell^\kappa E^t\in\pkmI0i$.
\end{definition}

$\pkmI li$ is an ideal, because $(X_S)_i(\mathrm{std}_l)$ is an $\pkpA$-module and $\mathrm{pr}$ is
$\pkpW$-linear; the $s$ with the lifting property for one $c$ form an ideal, and $\pkmI li$ is their
intersection.

\begin{remark}[Killing the cokernel: lifts in $\pkpA\mathcal N$ versus $\mathcal N$]\label{pkm:rem-kill}
In Definition~\ref{pkp:def-points} the lift $n$ is required to lie in $(X_S)_i(\mathrm{std}_0)$, i.e.\
in the letter-$i$ part of $\pkpA\mathcal N$ with orders $\mathrm{std}_0$, not in $\mathcal N(X_S)$
itself. This is the condition proved below, and it is all that the proof of
Proposition~\ref{pkp:rule} uses (through Lemma~\ref{pkp:tau}(iii), proved in
Lemma~\ref{pkm:tau-iii}). For a multiplier $f\in\mathfrak S$ (in particular $f=\ell^\kappa E^t$) the
stronger reading follows: by the retraction (R) a lift in $(X_S)_i(\mathrm{std}_l)$ can be replaced
by one in $\mathcal N(X_S)$ (Lemma~\ref{pkt:lift}). This is how the input (I1) produces points in
$\mathcal N(X_S)$ itself (Proposition~\ref{pkt:I1}).
\end{remark}

Let $d_v$ be as in Definition~\ref{pkp:def-chart} and $e$ the ramification index of the chart. All
jump bounds below refer to the jump sets of the side $S$ (in the cases, $S$ is the Tate-dual side
$T_P$, whose jump sets are $\{0\}\cup\{-u-u'-1\}$).

\begin{proposition}[Small jumps]\label{pkm:kill-small}
If every jump of $S$ at every letter is $\le\omega$ and $e\,\omega\le\ell-2$, then $1=\ell^0E^0$ kills
the cokernel of $S$ at every letter.
\end{proposition}

\begin{proposition}[Unramified charts]\label{pkm:kill-unram}
Let the chart be unramified ($e=1$, $E=z+\ell$), $\ell\ge5$, and $s,t\ge0$. If every jump of $S$ at
every letter is $\le(\ell-3)(s+t)$, then $\ell^sE^t$ kills the cokernel of $S$ at every letter.
\end{proposition}

\begin{lemma}[A power of $z$ in the lift ideals]\label{pkm:zpow}
If every jump of $S$ at every letter is $\le\omega$ and $\kappa\ge0$, there is $a\ge0$ with
$(\ell-1)a\le e(\omega+\kappa)$ and $z^a\in\pkmJ{\ell^{\kappa+1}}0i$ for every letter $i$.
\end{lemma}

\begin{proposition}[Killing by a power of $\ell$]\label{pkm:kill-J}
Let $i$ be a letter, $\omega,\omega',J,\kappa\ge0$. Assume every jump of $S$ at every letter is
$\le\omega$, every jump of $S$ at $\sigma^s(i)$ for $1\le s\le J$ is $\le\omega'$, $8e\le\ell^{J+1}$,
and
\[
2Jd_v\omega'+\frac{(2J+4)\,\omega'}{\ell-1}+\frac{8e\,\omega}{\ell^{J+1}}\le\kappa .
\]
Then $\ell^\kappa$ kills the cokernel of $S$ at $i$.
\end{proposition}

\subsubsection{Proof of the walk}\label{pkm:ss-walk}

\begin{lemma}[Distance to the previous cut]\label{pkm:dist}
In the situation of Theorem~\ref{pkm:walk} put $d(i)=0$ if $i\in\mathrm{cut}$ and $d(i)=l(i)$
otherwise. Then $i^{(-d(i))}\in\mathrm{cut}$ and $l(i)=d(i^{(-1)})+1$ for every letter $i$.
\end{lemma}

\begin{proof}
The first claim is the definition of $l$. Let $i_1=i^{(-1)}$; note $i_1^{(-a)}=i^{(-a-1)}$. If
$i_1\in\mathrm{cut}$, then $a=1$ is admissible, so $l(i)=1=d(i_1)+1$. Otherwise $l(i)\ge2$ and
$l(i)-1>0$ is admissible for $i_1$, so $l(i_1)\le l(i)-1$; conversely $l(i_1)+1$ is admissible for
$i$, so $l(i)\le l(i_1)+1$. Hence $l(i)=l(i_1)+1=d(i_1)+1$.
\end{proof}

\begin{proof}[Proof of Theorem~\ref{pkm:walk}]
For every letter $i$ put $\tilde x_i=\varphi^{d(i)}x(i^{(-d(i))})$, $y_i=\varphi^{l(i)}x(i^{(-l(i))})$,
and likewise $\tilde x'_i$, $y'_i$; these make sense by Lemma~\ref{pkm:dist}. By
Lemma~\ref{pkp:tau}(v), $\varphi^a$ of a point at $i^{(-a)}$ over $\ell^ec$ is a point at $i$ over
$\ell^ec$; hence $\tilde x_i,y_i$ lie over $\ell^ec$ at $i$, and $\tilde x'_i,y'_i$ over $\ell^{e'}c'$
at $i$.

{\emergencystretch=6em \emph{Frobenius.} Since $i\mapsto i^{(-1)}$ is a bijection of the letters and, by
Lemma~\ref{pkm:dist}, $\varphi\tilde x_{i^{(-1)}}=\varphi^{d(i^{(-1)})+1}x(i^{(-1-d(i^{(-1)}))})=y_i$,
\[
\varphi\sum_i\tilde x_i=\sum_i\varphi\tilde x_{i^{(-1)}}=\sum_iy_i,
\]
and likewise $\varphi\sum_i\tilde x'_i=\sum_iy'_i$.\par}

\emph{Paths.} Let $i\in\mathrm{cut}$, so $\tilde x_i=x(i)$, $\tilde x'_i=x'(i)$. The letter
$i^{(-l(i))}$ is a cut with $l(i^{(-l(i))})>0$; by (K) at $i$ and at $i^{(-l(i))}$, (P) at both and
(St) at $i$, Proposition~\ref{pkp:cut} (with $l=l(i)$, $l_0=l(i^{(-l(i))})$, $t=t(i)$,
$t_0=t(i^{(-l(i))})$, $x_0=x(i^{(-l(i))})$, $x'_0=x'(i^{(-l(i))})$) gives $w_i,w'_i$ over $\ell^ec$,
$\ell^{e'}c'$ at $i$ and a $\kappa$-safe path from $(y_i,y'_i)$ through $(w_i,w'_i)$ to
$(\tilde x_i,\tilde x'_i)$. Let $i\notin\mathrm{cut}$; then $\tilde x_i=y_i$, $\tilde x'_i=y'_i$. Take
$w_i=y_i$, $w'_i=y'_i$: each of the four moves is a move by $0$, and $(0,\cdot)$, $(\cdot,0)$ are
safe pairs since $P$ is bilinear and $0\in\pkpsm$.

\emph{Cocycle moves.} Let $i$ be a letter and $j=i^{(-d(i))}\in\mathrm{cut}$. By (K) and (P) at $j$,
the last assertion of Proposition~\ref{pkp:cut} (with the letter $j$ in place of $i^{(-d)}$, $d=d(i)$)
says that for every $h\in\Gamma'$ the move of the second side by $h\tilde x'_i-\tilde x'_i$ with
$\tilde x_i$ at rest is $\kappa$-safe.

All hypotheses of Corollary~\ref{pkp:letters} hold for the vectors $y_i,w_i,\tilde x_i$ and
$y'_i,w'_i,\tilde x'_i$, with $e,e',\kappa$. Hence $\mathrm{SCC}_N(c,c')$ for $N\ge e+e'+\kappa$.
\end{proof}

\begin{proof}[Proof of Corollary~\ref{pkm:walk-all}]
Take $\mathrm{cut}$ = all letters, $l\equiv1$ (the least $a>0$ with $i^{(-a)}$ a cut is $1$), $t$
constant, and apply Theorem~\ref{pkm:walk}.
\end{proof}

\subsubsection{Improved orders and stops}\label{pkm:ss-stops}

Recall (Definition~\ref{pkp:def-chart}) that $j(\Gamma')=G_{K'}$ for a finite extension $K'$ of
$\Q_\ell$ (the fixed field of $j(\Gamma')$); $\auQlb\subset B^+_{\mathrm{dR}}$ and $\mathrm{Fil}^n
B_{\mathrm{dR}}=t^nB^+_{\mathrm{dR}}$, with $\mathrm{gr}^nB_{\mathrm{dR}}\cong C(n)$ equivariantly.

\begin{lemma}[Invariants modulo the filtration]\label{pkm:BdR}
Let $y\in B_{\mathrm{dR}}$ and $n\in\Z$ with $gy-y\in\mathrm{Fil}^nB_{\mathrm{dR}}$ for all $g\in G_{K'}$.
Then there is $c\in K'$ with $y-c\in\mathrm{Fil}^nB_{\mathrm{dR}}$. If moreover $y\in\mathrm{Fil}^m$
with $m\ge1$, then $c$ can be taken $0$; and an element of $K'$ lying in $\mathrm{Fil}^1$ is $0$.
\end{lemma}

\begin{proof}
The last claim: $\theta$ is injective on $\auQlb\subset B^+_{\mathrm{dR}}$, and $\mathrm{Fil}^1=\ker\theta$.
By Tate's theorem, $C(j)^{G_{K'}}=0$ for $j\ne0$ and $C^{G_{K'}}=K'$ (\cite[\S3, Theorems~1 and~2, p.~176]{pklTate67}; also quoted in \cite[1.5.7]{pkfFontaineII}). Let $y\in\mathrm{Fil}^m$ (some $m$ exists). We show by induction
on $n-m$ that $y-c\in\mathrm{Fil}^n$ for some $c\in K'$, with $c=0$ if $m\ge1$. If $m\ge n$ take
$c=0$. Otherwise the image $\bar y$ of $y$ in $\mathrm{gr}^m=C(m)$ is $G_{K'}$-invariant (as
$gy-y\in\mathrm{Fil}^n\subseteq\mathrm{Fil}^{m+1}$). If $m\ne0$, then $\bar y=0$, so $y\in\mathrm{Fil}^{m+1}$
and we conclude by induction. If $m=0$, then $\bar y=\theta(y)\in K'$; put $c_0=\theta(y)\in K'\subset
B^+_{\mathrm{dR}}$; then $y-c_0\in\mathrm{Fil}^1$ is still invariant modulo $\mathrm{Fil}^n$ ($g$ fixes
$c_0$), and by induction (with $m\ge1$) $y-c_0\in\mathrm{Fil}^n$.
\end{proof}

\begin{lemma}[Values are equivariant]\label{pkm:equiv}
Let $L$ be a lattice with jump sets, $x\in\pkpA\mathcal N(L)$, $a\ge0$, $\mu$ an $\pkpA$-linear map
$\pkpAmb(L)\to B_{\mathrm{dR}}$ that is $\Gamma'$-equivariant on $1\otimes L$, and $g\in\Gamma'$. Then
$\mu(\varphi^{-a}gx)=g\,\mu(\varphi^{-a}x)$.
\end{lemma}

\begin{proof}
There is $r\in\pkpA$, a product of powers of $\tau$ and of the $p_b$, with $rx\in\pkpA\otimes L$
(Lemma~\ref{pkm:inputs}(M)). Write $\varphi^{-a}(rx)=\sum_kb_k\otimes y_k$ with $b_k\in\pkpA$. Then
$\varphi^{-a}(g(rx))=g\varphi^{-a}(rx)=\sum_k g(b_k)\otimes gy_k$, and by linearity and equivariance
$\mu(\varphi^{-a}g(rx))=\sum_kg(b_k)\,g\mu(1\otimes y_k)=g\mu(\varphi^{-a}(rx))$. Since
$\mu(\varphi^{-a}(rv))=\varphi^{-a}(r)\mu(\varphi^{-a}v)$ for all $v$, and $g$ commutes with
$\varphi^{-a}$, this reads $g(\varphi^{-a}r)\cdot\mu(\varphi^{-a}gx)=g(\varphi^{-a}r)\cdot
g\mu(\varphi^{-a}x)$. The image of $\varphi^{-a}(r)$ in the field $B_{\mathrm{dR}}$ is nonzero
($\pkpA\to B^+_{\mathrm{dR}}$ is injective and $\pkpA$ is a domain), hence so is that of $g\varphi^{-a}(r)$;
divide.
\end{proof}

\begin{proof}[Proof of Lemma~\ref{pkm:improved}]
Let $a\ge0$, $u\in U(i^{(-a)})$, $\mu$ a form of type $u$ at $i^{(-a)}$, $y=\mu(\varphi^{-a}x)$, and
$n=o_{\mathrm{hi}}(a,u)$. By Lemma~\ref{pkm:equiv}, $gy-y=\mu(\varphi^{-a}(gx-x))\in\mathrm{Fil}^n$ for all
$g\in\Gamma'$, so by Lemma~\ref{pkm:BdR} there is $c\in K'$ with $y-c\in\mathrm{Fil}^n$. Also
$y\in\mathrm{Fil}^{o_{\mathrm{lo}}(a,u)}$, so $c\in\mathrm{Fil}^{\min(o_{\mathrm{lo}},o_{\mathrm{hi}})(a,u)}$.
If $o_{\mathrm{lo}}(a,u)>0$ and $n>0$, then $c=0$ by Lemma~\ref{pkm:BdR}; hence $c\in\mathrm{Fil}^{\mathrm{imp}(a,u)}$
in all cases ($\mathrm{imp}\le0$ when $o_{\mathrm{lo}}\le0$, $c\in K'\subset\mathrm{Fil}^0$; and
$\mathrm{imp}=n\le0$ or $c=0$ when $o_{\mathrm{lo}}>0$). As $\mathrm{imp}\le n$, $y=(y-c)+c\in
\mathrm{Fil}^{\mathrm{imp}(a,u)}$. Membership in $\pkpA\mathcal N$ and in the letter part is unchanged.
\end{proof}

\begin{lemma}[Comparison]\label{pkm:le}
(i) $L_i(o)\subseteq L_i(o')$ if $o'(a,u)\le o(a,u)$ for all $a$ and $u\in U(i^{(-a)})$; hence a point
with orders $o_1$ and depth $d_1$ has orders $o_2$ and depth $d_2$ when $o_2\le o_1$ and $d_2\le
d_1$ on the jumps. (ii) If $q\in\mathcal N(X_S)$ lies over $\ell^ec$ at $i$ and has orders $o\ge
\mathrm{stop}_{\eta,l}$ on the jumps, then $q$ is a stop (for $\eta,l$).
\end{lemma}

\begin{proof}
(i) is immediate from Definition~\ref{pkp:def-orders} ($\mathrm{Fil}^n$ decreases in $n$). (ii) By (i),
$q\in L_i(\mathrm{stop}_{\eta,l})$; since $L_i(\mathrm{stop}_{\eta,l})$ is $\Gamma'$-stable
(Lemma~\ref{pkm:inputs}(G)), $gq-q\in L_i(\mathrm{stop}_{\eta,l})$ for $g\in\Gamma'$, i.e.\ $q$ is deep
for $\mathrm{stop}_{\eta,l}$.
\end{proof}

\begin{proof}[Proof of Lemma~\ref{pkm:down}]
Let $x$ be standard: orders $\rho_{-\eta}$, depth $\mathrm{dep}=\mathrm{dep}_{t,\eta,l}$. By
Lemma~\ref{pkm:improved}, $x$ has orders $\mathrm{imp}(\rho_{-\eta},\mathrm{dep})$. By
Lemma~\ref{pkm:le}(ii) it suffices to show $\mathrm{imp}(\rho_{-\eta},\mathrm{dep})(a,u)\ge
\mathrm{stop}_{\eta,l}(a,u)$ for $u\in U(i^{(-a)})$. \emph{$a=0$} ($<l$): $\rho_{-\eta}=-u-\eta$,
$\mathrm{dep}=-u+t+\eta$, $\mathrm{stop}=\max(0,-u-\eta)$. If $-u-\eta>0$, $\mathrm{imp}=-u+t+\eta\ge-u-\eta=
\mathrm{stop}$. Otherwise $\mathrm{stop}=0$ and $\mathrm{imp}=\min(-u+t+\eta,0)=0$ as $u\le\eta$.
\emph{$0<a<l$}: $\rho_{-\eta}=-u$, $\mathrm{dep}=-u+\eta\ge0$ (as $u\le\eta$), $\mathrm{stop}=
\max(0,-u-\eta)$. If $u<0$, $\mathrm{imp}=-u+\eta\ge\max(0,-u-\eta)$; if $u\ge0$, $\mathrm{imp}=
\min(-u+\eta,0)=0=\mathrm{stop}$. \emph{$a\ge l$}: $\rho_{-\eta}=\mathrm{dep}=\mathrm{stop}=-u$, and
$\mathrm{imp}=-u$ in both cases ($\min(-u,0)=-u$ when $-u\le0$).
\end{proof}

\begin{proof}[Proof of Lemma~\ref{pkm:plain}]
As $\mathrm{pd}_{t,\eta,l,\eta_0}\ge\mathrm{dep}_{t,\eta,l}$, $x$ is a standard point (Lemma~\ref{pkm:le}(i)),
hence a stop by Lemma~\ref{pkm:down}. By Lemma~\ref{pkp:tau}(v), $\varphi^lx'_0$ is a point at $i$ over
$\ell^{e'}c'$ with orders $\mathrm{sh}_{l,0}(\rho_{-\eta_0})$; so $y'=\varphi^lx'_0-x'$ lies over $0$
and has orders $o'=\min(\rho_{-\eta},\mathrm{sh}_{l,0}(\rho_{-\eta_0}))$. By the second half of
Proposition~\ref{pkp:rule} with $l=0$ (killing at $i^{(0)}=i$), it suffices that
$\mathrm{rule}_{0,t}(a,u+v)\le o'(a,u)+\mathrm{pd}_{t,\eta,l,\eta_0}(a,v)$ for $u$ a jump of $S'$ at
$i^{(-a)}$ and $v$ a jump of $S$ there. \emph{$a=0$}: left $t-u-v$; right
$\min(-u-\eta,0)+(-v+t+\eta)\ge t-u-v$ since $-u-\eta+\eta=-u$ and $0+\eta\ge-u$ ($u\ge-\eta$).
\emph{$0<a<l$}: left $-u-v$; right $\min(-u,0)+(-v+\eta)\ge-u-v$ as $\eta\ge-u$. \emph{$a=l$}: left
$-u-v$; $o'(l,u)=\min(-u,-u-\eta_0)=-u-\eta_0$ and $\mathrm{pd}(l,v)=-v+\eta_0$. \emph{$a>l$}: $o'=-u$,
$\mathrm{pd}=-v$. The statement with the sides exchanged is proved in the same way with the first
half of Proposition~\ref{pkp:rule}.
\end{proof}

\begin{lemma}[Lemma~\ref{pkp:tau}(iii)]\label{pkm:tau-iii}
If $\ell^\kappa E^t$ kills the cokernel of a side $S$ at $i$, then for every $c\in C_S$ there is
$\nu\in(X_S)_i(\rho_{-t})$ with $\mathrm{pr}(\nu)=\ell^\kappa e_i(1\otimes c)$.
\end{lemma}

\begin{proof}
Take $x\in(X_S)_i(\mathrm{std}_0)$ with $\mathrm{pr}(x)=\ell^\kappa E^te_i(1\otimes c)$ and put
$\nu=E^{-t}x$ ($E$ is a unit of $\pkpW$). By Lemma~\ref{pkm:inputs}(E), $\nu$ has orders
$\mathrm{std}_0-t[a=0]=\rho_{-t}$; and $\mathrm{pr}$ is $\pkpW$-linear.
\end{proof}

This is the form in which Lemma~\ref{pkp:tau}(iii) is used in the proof of
Proposition~\ref{pkp:rule}: there $\nu$ need only have orders $\rho_{-t}$ and lie over $\ell^\kappa\cdot1$
at $i^{(-l)}$, and Lemma~\ref{pkp:tau}(v) is applied to its orders only.

\subsubsection{The cokernel killers}\label{pkm:ss-kill}

\paragraph*{Inputs}

\begin{lemma}[Inputs from the theory of the modules $\mathcal N$]\label{pkm:inputs}
Let $L$ be a lattice with jump sets, $i$ a letter, $o$ orders, and $g\in\Gamma'$. We write
$\mathcal N_i(o)=\mathcal N(L)\cap L_i(o)$, an $\mathfrak S$-module.
\begin{enumerate}
\item[(M)] Every $x\in\pkpA\mathcal N(L)$ has $rx\in\pkpA\otimes L$ for some product $r$ of powers of
$\tau$ and of the $p_b$; $\pkpA\mathcal N(L)$ and $\mathcal N(L)$ are stable under multiplication by
$p_b^{\pm1}$ (in $\pkpW^\times$), and $\pkpA\mathcal N(L)$ is the $\pkpA$-span of $\mathcal N(L)$.
\item[(E)] For $b\ge0$, $n\in\Z$ and $x\in L_i(o)$: $p_b^nx\in L_i(o+n[a=b])$.
\item[(G)] $L_i(o)$ and $\pkpA\mathcal N(L)$ are $\Gamma'$-stable; if $x\in\mathcal N(L)$
and $gx-x\in L_i(o)$, then $(z\tau)^{-1}(gx-x)\in L_i(o)$ (this is Lemma~\ref{pkp:tau}(iv)). The second assertion follows from the first and from
\begin{enumerate}
\item[(G$_0$)] for $x\in\mathcal N(L)$ and $g\in\Gamma'$: $(z\tau)^{-1}(gx-x)\in\pkpA\mathcal N(L)$
\end{enumerate}
(Proposition~\ref{pkl:G0}).
\item[($\Phi$)] $\varphi(\mathcal N(L))\subseteq\mathcal N(L)$.
\item[(F)] For $l\ge0$, $\varphi^l$ maps $L_i(\mathrm{std}_0)$ bijectively onto
$L_{\sigma^l(i)}(\mathrm{std}_l)$.
\item[(H)] (Frame.) There is $H\ge0$ with $\tau^H(1\otimes y)\in\pkpA\mathcal N(L)$ and
$e_i\tau^H(1\otimes y)\in L_i(\mathrm{std}_0)$ for all $y\in L$. If $\Gamma$ acts trivially on $L$, then
$\mathcal N(L)\cap(\pkpA\otimes L)=\mathfrak S\otimes L$, $L_i(\mathrm{std}_0)\subseteq\pkpA\otimes L$, and
$e_i(1\otimes y)\in\mathcal N_i(\mathrm{std}_l)$ for all $l\ge0$.
\item[(R)] (Retraction.) There are an $\mathfrak S$-linear map $\lambda_0\colon\pkpA\to\mathfrak S$ that is the
identity on $\mathfrak S$, and, for every lattice, an $\mathfrak S$-linear map $\mathrm{ret}\colon
\pkpA\mathcal N\to\mathcal N$ with $\mathrm{ret}(bx)=\lambda_0(b)x$ for $b\in\pkpA$, $x\in\mathcal N$,
mapping $L_i(\mathrm{std}_l)$ into $\mathcal N_i(\mathrm{std}_l)$ and commuting with $\mathrm{pr}$ (for an
extension $X_S$).
\item[(D)] (Level one.) If $x\in\mathcal N_i(\mathrm{std}_1)$, then $gx-x\in L_i(o_1)$, where
$o_1(0,u)=1$ and $o_1(a,u)=-u$ for $a\ge1$.
\item[(C)] (Chart identities.) (a) $z\tau E$ divides $gs-s$ in $\pkpA$ for all $s\in\mathfrak S$.
(b) Some $t_1\in\Gamma'$ has $t_1z=[\varepsilon]z$. (c) For every $n\ge1$ there are
$\alpha,\beta\in\pkpA$ with $\tau E(\ell^{\pkmv(n)}+\beta z)=\alpha([\varepsilon]^n-1)$. (d) For every
$n\ge1$ there are $g\in\Gamma'$, $\omega\in\pkpA^\times$ and $\upsilon,q\in\pkpA$ with $gE=\omega E$,
$E\mid\upsilon q-\ell^{\pkmv(n)+d_v}$ and $z\tau q=\omega^n-1$.
\end{enumerate}
\end{lemma}

Also used, and elementary: $z,\tau,E$ are units of $\pkpW$ (their first Witt coordinates are nonzero
elements of the field $C^\flat$); $\varphi(z)=z^\ell$ and $\varphi(\mathfrak S)\subseteq\mathfrak S$; $z^e\in
E\mathfrak S+\ell\mathfrak S$ and $z\mid E-\ell$ in $\mathfrak S$ (from $E=z^eA+\ell$ with $A\in\mathfrak
S^\times$); $\varphi(e_i(1\otimes c))=e_{\sigma(i)}(1\otimes c)$ and $g\,e_i(1\otimes c)=e_i(1\otimes c)$
for $g\in\Gamma'$ (Definition~\ref{pkp:def-letters}); and $1-\ell b\in\pkpA^\times$ for $b\in\pkpA$.

\paragraph*{Operations on the lift ideals}

Fix a side $S$; $\pkmI li$, $\pkmJ rli$ refer to $S$, and ``jumps'' to its jump sets. We write
$\varepsilon_i(c)=e_i(1\otimes c)$ for $c\in C_S$ (the constant of $c$ at $i$).

\begin{lemma}\label{pkm:J-basic}
(i) (Nakayama.) If $\ell^c\in\pkmJ{\ell^{c+1}}li$, then $\pkmJ{\ell^{c+1}}li=\pkmI li$; in particular
$\ell^c\in\pkmI li$.
(ii) (Up.) If $s\in\pkmJ r0i$ and every jump at $i$ is $\le n$, then $sE^n\in\pkmJ r1i$.
(iii) (Frobenius.) $\varphi^l(s)\in\pkmJ rl{\sigma^l(i)}$ if and only if $s\in\pkmJ r0i$.
(iv) (Integral witness.) If $f\in\mathfrak S\cap\pkmJ rli$, there is $t\in\mathfrak S$ such that for every
$c$ some $x\in\mathcal N_i(\mathrm{std}_l)$ (for $X_S$) has $\mathrm{pr}(x)=(f-rt)\varepsilon_i(c)$.
\end{lemma}

\begin{proof}
(i) Write $\ell^c=s+\ell^{c+1}b$ with $s\in\pkmI li$. Then $\ell^c(1-\ell b)=s$, so
$\ell^c\in\pkmI li$ as $1-\ell b$ is a unit; hence $\ell^{c+1}\pkpA\subseteq\pkmI li$.
(ii) Write $s=s_0+rb$, $s_0\in\pkmI0i$, and let $x\in(X_S)_i(\mathrm{std}_0)$ with
$\mathrm{pr}(x)=s_0\varepsilon_i(c)$. By Lemma~\ref{pkm:inputs}(E), $E^nx$ has orders $\mathrm{std}_0+n[a=0]$,
which at $a=0$ is $n-u\ge0=\mathrm{std}_1(0,u)$ on the jumps $u\le n$ at $i$, and equals
$\mathrm{std}_1$ for $a\ge1$; so $E^nx\in(X_S)_i(\mathrm{std}_1)$ (Lemma~\ref{pkm:le}(i)) with
$\mathrm{pr}(E^nx)=s_0E^n\varepsilon_i(c)$. Thus $sE^n=s_0E^n+rbE^n\in\pkmJ r1i$.
(iii) $\varphi^l$ is a ring automorphism of $\pkpA$, commutes with $\mathrm{pr}$, and maps
$\varepsilon_i(c)$ to $\varepsilon_{\sigma^l(i)}(c)$. By Lemma~\ref{pkm:inputs}(F), $x\mapsto\varphi^lx$ is a
bijection $(X_S)_i(\mathrm{std}_0)\to(X_S)_{\sigma^l(i)}(\mathrm{std}_l)$, and $\mathrm{pr}(x)=
s_0\varepsilon_i(c)$ iff $\mathrm{pr}(\varphi^lx)=\varphi^l(s_0)\varepsilon_{\sigma^l(i)}(c)$. Hence
$\varphi^l(\pkmI0i)=\pkmI l{\sigma^l(i)}$, and the claim follows as $\varphi^l(r\pkpA)=r\pkpA$.
(iv) Write $f=s+rb$ with $s\in\pkmI li$, $b\in\pkpA$, and put $t=\lambda_0(b)$. For $c\in C_S$ let
$x\in(X_S)_i(\mathrm{std}_l)$ with $\mathrm{pr}(x)=s\varepsilon_i(c)$, and $x_1=\mathrm{ret}(x)\in
\mathcal N_i(\mathrm{std}_l)$ (Lemma~\ref{pkm:inputs}(R)). By (H), $\varepsilon_i(c)\in\mathcal N(C_S)$, so
$\mathrm{pr}(x_1)=\mathrm{ret}(s\varepsilon_i(c))=\lambda_0(s)\varepsilon_i(c)$, and
$\lambda_0(s)=\lambda_0(f-rb)=f-rt$.
\end{proof}

\begin{lemma}[Operators]\label{pkm:op}
Let $f\in\mathfrak S\cap\pkmJ rli$, $g\in\Gamma'$, $\alpha,\beta,s\in\pkpA$ and $d\in\pkpA\cap\pkpW^\times$.
Assume: $d\mid gt-t$ for all $t\in\mathfrak S$; for every $x\in\mathcal N_i(\mathrm{std}_l)$ (for $X_S$),
$d^{-1}(gx-x)\in(X_S)_i(\mathrm{std}_l)$; and $d(s+\beta f)=\alpha(gf-f)$. Then $s\in\pkmJ rli$.
\end{lemma}

\begin{proof}
Take $t$ as in Lemma~\ref{pkm:J-basic}(iv) and write $gt-t=dv$. For $c\in C_S$ let $x\in\mathcal
N_i(\mathrm{std}_l)$ with $\mathrm{pr}(x)=(f-rt)\varepsilon_i(c)$, and $y=d^{-1}(gx-x)\in
(X_S)_i(\mathrm{std}_l)$. As $g$ fixes $\varepsilon_i(c)$ and commutes with $\mathrm{pr}$,
$\mathrm{pr}(gx-x)=(gf-f-r(gt-t))\varepsilon_i(c)$. Hence
\[
d\,\mathrm{pr}(\alpha y)=\alpha(gf-f)\varepsilon_i(c)-r\alpha dv\,\varepsilon_i(c)=d\,(s+\beta f-r\alpha
v)\,\varepsilon_i(c),
\]
and, $d$ being a unit of $\pkpW$, $\mathrm{pr}(\alpha y)=(s+\beta f-r\alpha v)\varepsilon_i(c)$ with $\alpha
y\in(X_S)_i(\mathrm{std}_l)$. So $s+\beta f-r\alpha v\in\pkmI li$, whence $s+\beta f\in\pkmJ rli$, and
$s\in\pkmJ rli$ since $\beta f\in\pkmJ rli$.
\end{proof}

\begin{corollary}[Derivations]\label{pkm:der}
Let $f\in\mathfrak S$, $g\in\Gamma'$, $\alpha,\beta,s\in\pkpA$. (i) If $f\in\pkmJ rli$ and
$z\tau(s+\beta f)=\alpha(gf-f)$, then $s\in\pkmJ rli$. (ii) If $f\in\pkmJ r1i$ and $z\tau E(s+\beta
f)=\alpha(gf-f)$, then $s\in\pkmJ r1i$.
\end{corollary}

\begin{proof}
(i) Apply Lemma~\ref{pkm:op} with $d=z\tau$: $z\tau\mid gt-t$ by Lemma~\ref{pkm:inputs}(C)(a); for
$x\in\mathcal N_i(\mathrm{std}_l)$, $gx-x\in(X_S)_i(\mathrm{std}_l)$ by (G), and then $(z\tau)^{-1}(gx-x)\in
(X_S)_i(\mathrm{std}_l)$ by (G). (ii) Apply Lemma~\ref{pkm:op} with $d=z\tau E$: the divisibility is
(C)(a); for $x\in\mathcal N_i(\mathrm{std}_1)$, $gx-x\in(X_S)_i(o_1)$ by (D), $(z\tau)^{-1}(gx-x)\in
(X_S)_i(o_1)$ by (G), and by (E) $E^{-1}(z\tau)^{-1}(gx-x)$ has orders $o_1-[a=0]$, which is
$0=\mathrm{std}_1$ at $a=0$ and $-u$ for $a\ge1$.
\end{proof}

\paragraph*{Monomials in $\ell$, $z$, $E$}

\begin{lemma}[Coming down a level]\label{pkm:down-level}
If $\ell^cz^n\in\pkmJ r1{\sigma(i)}$, $\ell a\le n$ and $c+\pkmv(n!)\le\kappa+\pkmv((\ell a)!)$, then
$\ell^\kappa z^a\in\pkmJ r0i$.
\end{lemma}

\begin{proof}
Induction on $n\ge\ell a$. If $n=\ell a$: $\ell^cz^{\ell a}=\varphi(\ell^cz^a)$, so $\ell^cz^a\in\pkmJ r0i$
by Lemma~\ref{pkm:J-basic}(iii), and $c\le\kappa$. Step $n\to n+1$ (with $n\ge\ell a$): let
$f=\ell^cz^{n+1}\in\mathfrak S\cap\pkmJ r1{\sigma(i)}$, $t_1$ as in (C)(b) and $\alpha,\beta$ as in
(C)(c) for $n+1$. Then $t_1f-f=\ell^cz^{n+1}([\varepsilon]^{n+1}-1)$ and
\[
z\tau E\bigl(\ell^{c+\pkmv(n+1)}z^n+\beta f\bigr)=\ell^cz^{n+1}\,\tau E(\ell^{\pkmv(n+1)}+\beta z)=
\alpha(t_1f-f).
\]
By Corollary~\ref{pkm:der}(ii), $\ell^{c+\pkmv(n+1)}z^n\in\pkmJ r1{\sigma(i)}$, and
$c+\pkmv(n+1)+\pkmv(n!)=c+\pkmv((n+1)!)\le\kappa+\pkmv((\ell a)!)$; apply the induction hypothesis.
\end{proof}

\begin{lemma}[Trading $E$ for $z$]\label{pkm:trade}
If $\ell^cz^aE^n\in\pkmJ r1i$, $c+d_vn+\pkmv(n!)\le b$ and $a+n\le\kappa$, then $\ell^bz^\kappa\in\pkmJ r1i$.
\end{lemma}

\begin{proof}
Induction on $n$; $n=0$ is clear ($\ell^bz^\kappa$ is a multiple of $\ell^cz^a$). Step $n+1$: let
$f=\ell^cz^aE^{n+1}$, and $g,\omega,\upsilon,q$ as in (C)(d) for $n+1$; write $\upsilon
q-\ell^{\gamma}=Eb_1$ with $\gamma=\pkmv(n+1)+d_v$. By (C)(a), $gz=z(1+\tau Ev)$ for some $v$, and
$(1+\tau Ev)^a-1=\tau Evb_2$ for some $b_2$. Hence $gf=\ell^cz^a(1+\tau Evb_2)\omega^{n+1}E^{n+1}$, and,
using $\omega^{n+1}-1=z\tau q$,
\[
z\upsilon(gf-f)=\ell^cz^{a+1}E^{n+1}\tau\bigl(z\upsilon q+E\upsilon vb_2\omega^{n+1}\bigr)
=z\tau E\bigl(\ell^{c+\gamma}z^{a+1}E^n+\beta f\bigr),\quad \beta=zb_1+\upsilon vb_2\omega^{n+1},
\]
where we substituted $\upsilon q=\ell^\gamma+Eb_1$. By Corollary~\ref{pkm:der}(ii) (with $\alpha=z\upsilon$),
$\ell^{c+\gamma}z^{a+1}E^n\in\pkmJ r1i$. Since $(c+\gamma)+d_vn+\pkmv(n!)=c+d_v(n+1)+\pkmv((n+1)!)\le b$
and $(a+1)+n\le\kappa$, the induction hypothesis applies.
\end{proof}

\begin{lemma}[Lowering the power of $z$]\label{pkm:lower}
(i) If $\ell^cz^{n}\in\pkmJ r0i$ with $n\ge1$, then $\ell^{c+\pkmv(n)+1}z^{n-1}\in\pkmJ r0i$.
(ii) If $\ell^cz^n\in\pkmJ r0i$, $m\le n$ and $c+n+\pkmv(n!)\le\kappa+m+\pkmv(m!)$, then $\ell^\kappa
z^m\in\pkmJ r0i$.
\end{lemma}

\begin{proof}
(i) Let $\alpha,\beta$ be as in (C)(c) for $n$, $v=\pkmv(n)$, and write $E-\ell=ze'$ with $e'\in\mathfrak S$.
Then $\tau\bigl(\ell^{v+1}+z(\ell^ve'+\beta E)\bigr)=\tau E(\ell^v+\beta z)=\alpha([\varepsilon]^n-1)$.
With $f=\ell^cz^n$ and $t_1$ as in (C)(b), multiplying by $\ell^cz^n$ gives $z\tau(\ell^{c+v+1}z^{n-1}+
(\ell^ve'+\beta E)f)=\alpha(t_1f-f)$; apply Corollary~\ref{pkm:der}(i).
(ii) Induction on $n-m$ using (i): a step from $n+1$ to $n$ costs $\pkmv(n+1)+1$, and
$\pkmv((n+1)!)=\pkmv(n+1)+\pkmv(n!)$.
\end{proof}

\begin{lemma}[One letter back]\label{pkm:div}
If every jump at $\sigma(j)$ is $\le\omega$ and $\ell^cz^n\in\pkmJ r0{\sigma(j)}$, then
$\ell^{c+d_v\omega+\pkmv(\omega!)}z^{\lfloor(n+\omega)/\ell\rfloor}\in\pkmJ r0j$.
\end{lemma}

\begin{proof}
By Lemma~\ref{pkm:J-basic}(ii), $\ell^cz^nE^\omega\in\pkmJ r1{\sigma(j)}$. By Lemma~\ref{pkm:trade} (with
$a=n$, $n=\omega$, $\kappa=n+\omega$, $b=c+d_v\omega+\pkmv(\omega!)$), $\ell^bz^{n+\omega}\in\pkmJ
r1{\sigma(j)}$. Put $a'=\lfloor(n+\omega)/\ell\rfloor$; no multiple of $\ell$ lies in $(\ell a',n+\omega]$,
so $\pkmv((n+\omega)!)=\pkmv((\ell a')!)$, and Lemma~\ref{pkm:down-level} (with $c=\kappa=b$) gives
$\ell^bz^{a'}\in\pkmJ r0j$.
\end{proof}

\begin{lemma}[Back through $J$ letters]\label{pkm:back}
Let $z^a\in\pkmJ r0j$ for every letter $j$, and let every jump at $\sigma^s(i)$, $1\le s\le J$, be
$\le\omega$. Then $\ell^cz^n\in\pkmJ r0i$ for some $c\le J(d_v\omega+\pkmv(\omega!))$ and $n\le
a/\ell^J+\omega/(\ell-1)$.
\end{lemma}

\begin{proof}
Induction on $J$; $J=0$: $c=0$, $n=a$. Step: by the induction hypothesis at the letter $\sigma(i)$
(whose next $J$ letters are $\sigma^s(i)$, $2\le s\le J+1$) there are $c,n$ for $\sigma(i)$; by
Lemma~\ref{pkm:div} with $j=i$, $\ell^{c+d_v\omega+\pkmv(\omega!)}z^{\lfloor(n+\omega)/\ell\rfloor}\in\pkmJ
r0i$, and $\lfloor(n+\omega)/\ell\rfloor\le(a/\ell^J+\omega/(\ell-1)+\omega)/\ell=a/\ell^{J+1}+
\omega/(\ell-1)$.
\end{proof}

\paragraph*{Proofs of the killers}

\begin{proof}[Proof of Lemma~\ref{pkm:zpow}]
\emph{Some power exists.} Let $H$ be as in Lemma~\ref{pkm:inputs}(H) for $X_S$. For $c\in C_S$,
$x=e_i\tau^H(1\otimes(0,c))\in(X_S)_i(\mathrm{std}_0)$ and $\mathrm{pr}(x)=\tau^H\varepsilon_i(c)$; so
$\tau^H\in\pkmI0i$ for all $i$. Also $\ell^{\kappa+1}\in\pkmJ{\ell^{\kappa+1}}0i$. In $\pkpA/\ell\pkpA=
O^\flat$ the images of $\tau$ and $z$ are $\varepsilon^{1/\ell}-1$ and $\varpi^\flat$, both of positive
valuation; so $z^M\in\tau\pkpA+\ell\pkpA$ for some $M$, and $z^{M(H+\kappa+1)}\in(\tau,\ell)^{H+\kappa+1}
\subseteq\bigcap_i\pkmJ{\ell^{\kappa+1}}0i$.

\emph{The bound.} Let $a$ be the least integer with $z^a\in\pkmJ{\ell^{\kappa+1}}0i$ for every $i$. Write
$z^e=u_1E+v_1\ell$ with $u_1,v_1\in\mathfrak S$, and $K=\omega+\kappa$. For each $i$,
$z^aE^\omega\in\pkmJ{\ell^{\kappa+1}}1{\sigma(i)}$ by Lemma~\ref{pkm:J-basic}(ii), and
$z^{a+eK}=z^a(u_1E+v_1\ell)^K$ is a combination of the $z^aE^j\ell^{K-j}$; for $j\ge\omega$ these are
multiples of $z^aE^\omega$, and for $j<\omega$ multiples of $\ell^{\kappa+1}$. So $z^{a+eK}\in\pkmJ
{\ell^{\kappa+1}}1{\sigma(i)}$, and by Lemma~\ref{pkm:down-level} (with $c=\kappa=0$, $a'=\lfloor(a+eK)/\ell
\rfloor$, as in Lemma~\ref{pkm:div}) $z^{a'}\in\pkmJ{\ell^{\kappa+1}}0i$ for all $i$. By minimality $a\le a'$,
so $\ell a\le a+eK$, i.e.\ $(\ell-1)a\le e(\omega+\kappa)$.
\end{proof}

\begin{proof}[Proof of Proposition~\ref{pkm:kill-small}]
Lemma~\ref{pkm:zpow} with $\kappa=0$ gives $a$ with $(\ell-1)a\le e\omega\le\ell-2$, so $a=0$ and
$1\in\pkmJ\ell0i$; by Lemma~\ref{pkm:J-basic}(i) with $c=0$, $1\in\pkmI0i$.
\end{proof}

\begin{proof}[Proof of Proposition~\ref{pkm:kill-unram}]
Only $\ell\ge3$ is used: it gives $\omega=(\ell-3)K\ge0$ and $\ell-2>0$. Let $K=s+t$. Lemma~\ref{pkm:zpow} (with $\omega=(\ell-3)K$, $\kappa=K$, $e=1$) gives $b$ with $(\ell-1)b\le
(\ell-2)K$ and $z^b\in\pkmJ{\ell^{K+1}}0i$ for all $i$. By Legendre's formula $(\ell-1)\pkmv(b!)\le b$, so
$(\ell-2)(b+\pkmv(b!))\le(\ell-2)b+(\ell-1)\pkmv(b!)\le(\ell-1)b\le(\ell-2)K$, i.e.\ $b+\pkmv(b!)\le K$.
For every $m\ge0$, $\ell^{\max(K-m,0)}z^m\in\pkmJ{\ell^{K+1}}0i$: from $z^{\max(b,m)}$ (a multiple of $z^b$)
apply Lemma~\ref{pkm:lower}(ii); the cost condition $\max(b,m)+\pkmv(\max(b,m)!)\le\max(K-m,0)+m+
\pkmv(m!)$ is trivial if $m\ge b$, and follows from $b+\pkmv(b!)\le K$ if $m<b$. With $m=0$,
$\ell^K\in\pkmJ{\ell^{K+1}}0i$, so $\pkmJ{\ell^{K+1}}0i=\pkmI0i$ by Lemma~\ref{pkm:J-basic}(i). Finally
$\ell^sE^t=\ell^s(z+\ell)^t=\sum_{m\le t}\binom tm\ell^{K-m}z^m\in\pkmI0i$.
\end{proof}

\begin{lemma}[Arithmetic]\label{pkm:arith}
Let $J,d,e,\omega,\omega',\kappa,a,c,n\ge0$ with $c\le J(d\omega'+\pkmv(\omega'!))$, $n\le
a/\ell^J+\omega'/(\ell-1)$, $(\ell-1)a\le e(\omega+\kappa)$, $8e\le\ell^{J+1}$ and $2Jd\omega'+
(2J+4)\omega'/(\ell-1)+8e\omega/\ell^{J+1}\le\kappa$. Then $c+n+\pkmv(n!)\le\kappa$.
\end{lemma}

\begin{proof}
By Legendre, $\pkmv(\omega'!)\le\omega'/(\ell-1)$ and $\pkmv(n!)\le n$; so $c+n+\pkmv(n!)\le
J(d\omega'+\omega'/(\ell-1))+2n$. As $\ell\le2(\ell-1)$, $\ell a\le2e(\omega+\kappa)$, so
$a/\ell^J=\ell a/\ell^{J+1}\le2e\omega/\ell^{J+1}+(8e/\ell^{J+1})\kappa/4\le2e\omega/\ell^{J+1}+\kappa/4$.
Hence $c+n+\pkmv(n!)\le\tfrac12\bigl(2Jd\omega'+(2J+4)\omega'/(\ell-1)+8e\omega/\ell^{J+1}\bigr)+
\kappa/2\le\kappa$.
\end{proof}

\begin{proof}[Proof of Proposition~\ref{pkm:kill-J}]
Lemma~\ref{pkm:zpow} gives $a$ with $(\ell-1)a\le e(\omega+\kappa)$ and $z^a\in\pkmJ{\ell^{\kappa+1}}0j$ for
all $j$. Lemma~\ref{pkm:back} (with $\omega'$) gives $\ell^cz^n\in\pkmJ{\ell^{\kappa+1}}0i$ with $c\le
J(d_v\omega'+\pkmv(\omega'!))$ and $n\le a/\ell^J+\omega'/(\ell-1)$. By Lemma~\ref{pkm:arith},
$c+n+\pkmv(n!)\le\kappa$, so Lemma~\ref{pkm:lower}(ii) with $m=0$ gives $\ell^\kappa\in\pkmJ{\ell^{\kappa+1}}0i$,
and Lemma~\ref{pkm:J-basic}(i) gives $\ell^\kappa\in\pkmI0i$.
\end{proof}

\emph{Where the inputs are proved.} Everything else in \S\ref{pkm:sec} is proved from
Lemma~\ref{pkm:inputs}, Lemma~\ref{pkp:orders-add} and Lemma~\ref{pkp:tau}(i), (iv), (v). The items of
Lemma~\ref{pkm:inputs}: (M), first two assertions, Lemma~\ref{pkn:M}, third assertion,
Proposition~\ref{pkn:frame}(iv) with Theorem~\ref{pkf:main}; (E) Lemma~\ref{pkn:E}; (G), first assertion,
Lemma~\ref{pkl:G1}, second assertion and (G$_0$), Proposition~\ref{pkl:G0}; ($\Phi$)
Corollary~\ref{pkl:Phi}; (F) Corollary~\ref{pkl:F}; (H) Corollary~\ref{pkl:H1}, Proposition~\ref{pkl:H2}
and Corollary~\ref{pkl:H3}; (R) Proposition~\ref{pkl:R}; (D) Proposition~\ref{pkl:D}; (C)(a)
Proposition~\ref{pkn:Ca}, (C)(b) \S\ref{pkn:ss-ME}, (C)(c) Proposition~\ref{pkn:Cc}, (C)(d)
Proposition~\ref{pkl:Cd}. Lemma~\ref{pkp:orders-add} is Lemma~\ref{pkt:add}, and
Lemma~\ref{pkp:tau}(i), (iv), (v) are Proposition~\ref{pkt:tau-i}, Lemmas~\ref{pkt:tau-iv},
\ref{pkt:tau-v}.



\section{The modules $\mathcal N$, frames and the period rings}\label{C4:g3}
\providecommand{\auO}{\mathcal O}
\providecommand{\auQlb}{\overline{\Q}_\ell}
\providecommand{\panot}[1]{\par\smallskip{\raggedright\noindent\textbf{[Not proved here:} #1\textbf{]}\par}\smallskip}
\providecommand{\pkpA}{A_{\mathrm{inf}}}
\providecommand{\pkpW}{W(C^\flat)}
\providecommand{\pkpm}{W(\mathfrak m^\flat)}
\providecommand{\pkpAmb}{\mathbb A}
\providecommand{\pknMinf}{\mathfrak M_{\mathrm{inf}}}
\providecommand{\pknK}{\mathcal K_0}
\providecommand{\pknBp}{B^+_{\mathrm{dR}}}
\providecommand{\pknB}{B_{\mathrm{dR}}}

\subsection{The modules $\mathcal N$: definition and first inputs}\label{pkn:sec}

This section continues \S\ref{pkp:sec} and \S\ref{pkm:sec}, whose conventions and labels we use. It
fixes the definition of the modules $\mathcal N$ used in Definition~\ref{pkp:def-orders}
(\S\ref{pkn:ss-def}), proves that a distinguished element
generates $\ker\theta$ (which completes the proof of Lemma~\ref{pkp:chart-exists}), and proves the
items (M) (first two assertions), (E), (C)(a), (C)(b) and (C)(c) of Lemma~\ref{pkm:inputs}.

\emph{Dependencies.} Proved in full: Lemma~\ref{pkn:dist} (which completes the proof of
Lemma~\ref{pkp:chart-exists}), Lemmas~\ref{pkn:mu}, \ref{pkn:BdR}, \ref{pkn:M}, \ref{pkn:E} and
Propositions~\ref{pkn:Ca}, \ref{pkn:Cc} (items (M) first two assertions, (E), (C)(a), (C)(b), (C)(c) of
Lemma~\ref{pkm:inputs}), modulo published facts on $\pknBp$ (that $t$ is a uniformiser,
Fontaine \cite[1.5.2, 1.5.4]{pkfFontaineII}; that $\pkpA\to\pknBp$ is injective, \cite[4.1.1]{pkfFontaineII}) and Tate's theorem (through Lemma~\ref{pkm:BdR}). Lemma~\ref{pkn:frame-id} is proved for any given frame. Proved modulo the existence of a frame (Definition~\ref{pkn:def-frame}), which is
Theorem~\ref{pkf:main}: Proposition~\ref{pkn:frame} (item (M), third assertion, and the first assertion
of (H)). The other items are proved in \S\S\ref{pkf:sec}--\ref{pkl:sec}; see the paragraph at the end
of the section.

\subsubsection{The modules $\mathcal N$}\label{pkn:ss-def}

Fix a Kummer chart (Definition~\ref{pkp:def-chart}). Recall that $\pknBp$ is the completion of
$\pkpA[1/\ell]$ for the $\ker(\theta[1/\ell])$-adic topology, $\pknB$ its fraction field, and
$\mathrm{Fil}^n=t^n\pknBp$ ($n\in\Z$); $\pkpA\to\pknBp$ is the natural map. The group $\Gamma$ acts on
$\pknB$ through $j$, compatibly with this map.

\begin{definition}[The modules $\mathcal N$]\label{pkn:def-N}
Let $L$ be a lattice with jump sets (Definition~\ref{pkp:def-orders}, first sentence). A
\emph{form} on $L$ is an $\pkpA$-linear map $\mu\colon\pkpAmb(L)\to\pknB$ with $\mu(1\otimes gy)=
g\,\mu(1\otimes y)$ for $g\in\Gamma'$, $y\in L$; a form \emph{of type $u$ at $i$} moreover satisfies
$\mu(e_i(1\otimes L))\subseteq\mathrm{Fil}^{-u}$. Put
\begin{align*}
\pknMinf(L)&=\{x\in\pkpAmb(L):\ \tau^nx\in\pkpA\otimes L\text{ for some }n\ge0,\ \text{and }
\mu(\varphi^{n}x)\in\pknBp\text{ for all }n\ge1\text{ and all forms }\mu\},\\
\pkpA\mathcal N(L)&=\{x\in\pkpAmb(L):\ rx\in\pknMinf(L)\text{ for some finite product }r\text{ of
elements }p_a,\ a\ge0\},\\
\mathcal N(L)&=\{x\in\pkpA\mathcal N(L):\ \mu(x)\in\pknK\text{ for all forms }\mu\},
\end{align*}
where $\pkpA\otimes L$ is the $\pkpA$-span of the $1\otimes y$, and $\pknK\subseteq\pknB$ is the
subfield generated over $K'=\auQlb^{\,j(\Gamma')}$ by the images of the products $bs$, $b\in\mathfrak
S_B$, $s\in\mathfrak S$, with $\mathfrak S_B\subseteq B^+$ the bicommutant of the chart (Definition~\ref{pkf:def-SB}; a subring
containing $1$, mapped to $\pknBp$ by $B^+\to\pknBp$, see \S\ref{pkn:ss-frame}; we use only that $\pknK$ is a subfield of $\pknB$ containing the image of $\mathfrak
S$). The orders of Definition~\ref{pkp:def-orders} are defined with this $\pkpA\mathcal N(L)$: for
orders $o$, $L_i(o)$ is the set of $x\in\pkpA\mathcal N(L)\cap\pkpAmb(L)_i$ with $\mu(\varphi^{-a}x)\in
\mathrm{Fil}^{o(a,u)}$ for all $a\ge0$, $u\in U(i^{(-a)})$ and forms $\mu$ of type $u$ at $i^{(-a)}$.
\end{definition}

The notation $\pkpA\mathcal N$ is justified by Lemma~\ref{pkm:inputs}(M) (third assertion, Proposition~\ref{pkn:frame}(iv) with
Theorem~\ref{pkf:main}). No comparison with Breuil--Kisin modules is used anywhere in
\S\S\ref{pkp:sec}--\ref{pkn:sec}.

With this definition, two statements of \S\ref{pkp:sec} are to be read as follows.

\emph{Killing the cokernel} (Definition~\ref{pkp:def-points}): $\ell^\kappa E^t$ kills the cokernel
at $i$ for the side $S$ if for every $c\in C_S$ there is $n\in(X_S)_i(\mathrm{std}_0)$ with
$\mathrm{pr}(n)=\ell^\kappa E^te_i(1\otimes c)$; $n$ lies in $\pkpA\mathcal N(X_S)$, not necessarily in
$\mathcal N(X_S)$.

\emph{Reading of Lemma~\ref{pkp:tau}.} (i) For the lattice $\auO(1)$ (jump set $\{-1\}$), every $\xi\in
\auO(1)_i(\mathrm{std}_0)$ has $\tau\xi\in\pkpA\otimes\auO(1)$. (ii) $z\in\pkpm$, and the image of
$z(\pkpA\otimes\auO(1))$ in $\pkpM$ lies in $\pkpsm$ (Lemma~\ref{pkp:AS}(vii)). (iii) If
$\ell^\kappa E^t$ kills the cokernel at $i$ for a side, then for every $c$ there is $\nu\in
(X_S)_i(\rho_{-t})$ with $\mathrm{pr}(\nu)=\ell^\kappa e_i(1\otimes c)$ (Lemma~\ref{pkm:tau-iii}).
(iv) If $x\in\mathcal N$ lies at $i$, is deep for $o'$ and $\mathrm{pr}(x)$ is $\Gamma'$-invariant, then
for $g\in\Gamma'$, $(z\tau)^{-1}(gx-x)\in L_i(o')$ and lies over $0$. (v) If $x\in L_{i^{(-l)}}(o)$
then $\varphi^lx\in L_i(\mathrm{sh}_{l,0}(o))$; if $x$ is a point at $i^{(-l)}$ with orders $o$ and
depth $o'$, then $\varphi^lx$ is a point at $i$ over the same constant with orders $\mathrm{sh}_{l,0}(o)$
and depth $\mathrm{sh}_{l,1}(o')$. Assertions (i), (iv), (v) are proved in \S\ref{pkt:sec}.

\subsubsection{Distinguished elements}\label{pkn:ss-dist}

For $x\in\pkpA=W(O^\flat)$ write $x=\sum_{n\ge0}\ell^n[x_n^{1/\ell^n}]$ with $x_n\in O^\flat$ the Witt
coordinates ($O^\flat$ is perfect); so $x\equiv[x_0]\pmod{\ell\pkpA}$, $\pkpA/\ell=O^\flat$ via $x\mapsto
x_0$, and $\theta(x)=\sum_n\ell^n(x_n^{1/\ell^n})^\sharp$, where $b\mapsto b^\sharp$ is the
multiplicative map $O^\flat\to O_C$, $b^\sharp=\lim b_n^{\ell^n}$ for lifts $b_n$ of the coordinates.
Let $v^\flat(b)=v_\ell(b^\sharp)$, a valuation on $O^\flat$ with $O^\flat$ its valuation ring. An element
$x$ is a unit of $\pkpA$ iff $x_0\in O^{\flat\times}$ (if $x_0$ is a unit, $x=[x_0](1-\ell y)$ with
$y\in\pkpA$, and $1-\ell y$ is invertible by $\ell$-adic completeness); equivalently iff its image in
$W(\overline{\mathbb F}_\ell)$ under $O^\flat\to O^\flat/\mathfrak m^\flat=\overline{\mathbb F}_\ell$ is a unit.

\begin{definition}\label{pkn:def-dist}
$\xi\in\pkpA$ is \emph{distinguished} if $\xi_0\in\mathfrak m^\flat$ and its image in
$W(\overline{\mathbb F}_\ell)$ is $\ell$ times a unit; equivalently $\xi=[\xi_0]+\ell\xi'$ with
$\xi_0\in\mathfrak m^\flat$, $\xi'\in\pkpA^\times$.
\end{definition}

(Equivalence: the image of $[\xi_0]$ in $W(\overline{\mathbb F}_\ell)$ is $0$, so the image of $\xi$
is $\ell$ times that of $\xi'$, and $W(\overline{\mathbb F}_\ell)$ is torsion free.)

\begin{lemma}\label{pkn:dist}
Let $\xi$ be distinguished with $\theta(\xi)=0$. Then:
\begin{enumerate}
\item[(i)] $v^\flat(\xi_0)=1$ ($=v_\ell(\ell)$);
\item[(ii)] $\ker\theta=\xi\pkpA$;
\item[(iii)] for $c\in\Z\smallsetminus\{0\}$, $\varphi^c(\xi)$ is distinguished and
$\theta(\varphi^c\xi)\ne0$;
\item[(iv)] if $\xi''$ is another distinguished element of $\ker\theta$, then $\xi''=\xi u$ with
$u\in\pkpA^\times$.
\end{enumerate}
\end{lemma}

\begin{proof}
(i) $0=\theta(\xi)=\xi_0^\sharp+\ell\theta(\xi')$, and $\theta(\xi')\in O_C^\times$ because
$\theta(\xi')\equiv(\xi'_0)^\sharp\pmod\ell$ with $\xi'_0$ a unit. Hence $v_\ell(\xi_0^\sharp)=1$.

(ii) Clearly $\xi\pkpA\subseteq\ker\theta$. \emph{Step 1: if $x\in\pkpA$ and $\theta(x)\in\ell O_C$,
then $x\in\xi\pkpA+\ell\pkpA$.} Indeed $\theta(x)\equiv x_0^\sharp\pmod\ell$, so $v^\flat(x_0)\ge1=
v^\flat(\xi_0)$ and $x_0=\xi_0c$ with $c\in O^\flat$. Then $x-\xi[c]\equiv[x_0]-[\xi_0][c]=0
\pmod{\ell\pkpA}$. \emph{Step 2.} Let $x\in\ker\theta$. Put $x^{(0)}=x$. Given $x^{(k)}\in\ker\theta$,
Step 1 gives $y_k,x^{(k+1)}\in\pkpA$ with $x^{(k)}=\xi y_k+\ell x^{(k+1)}$; then
$\ell\theta(x^{(k+1)})=0$, and as $O_C$ is torsion free, $x^{(k+1)}\in\ker\theta$. Thus
$x=\xi(y_0+\ell y_1+\dots+\ell^ky_k)+\ell^{k+1}x^{(k+1)}$ for all $k$. As $\pkpA$ is $\ell$-adically
complete and separated, $y=\sum_k\ell^ky_k$ converges and $x=\xi y$.

(iii) $\varphi$ is the automorphism of $\pkpA$ induced by the $\ell$-th power map on coordinates, so
$\varphi^c(\xi)_0=\xi_0^{\ell^c}\in\mathfrak m^\flat$, and $\varphi^c$ induces the Frobenius power on
$W(\overline{\mathbb F}_\ell)$, which fixes $\ell$ and maps units to units: $\varphi^c\xi$ is
distinguished. If $\theta(\varphi^c\xi)=0$, (i) applied to $\varphi^c\xi$ would give
$v^\flat(\xi_0^{\ell^c})=1$, while $v^\flat(\xi_0^{\ell^c})=\ell^c\ne1$.

(iv) By (ii), $\xi''=\xi u$ with $u\in\pkpA$. In $W(\overline{\mathbb F}_\ell)$: $\ell w''=\ell w\bar u$
with $w,w''$ units, so $\bar u$ is a unit, hence $u\in\pkpA^\times$.
\end{proof}

\begin{corollary}[Completion of Lemma~\ref{pkp:chart-exists}]\label{pkn:E-gen}
In Definition~\ref{pkp:def-chart} and the proof of Lemma~\ref{pkp:chart-exists}, $E=z^eA+\ell$ is
distinguished and generates $\ker\theta$; for every $a\in\Z$, $\varphi^a(E)$ generates
$\ker(\theta\circ\varphi^{-a})$, and $\theta(\varphi^c E)\ne0$ for $c\ne0$.
\end{corollary}

\begin{proof}
$E_0=(\varpi^\flat)^e\bar A_0\in\mathfrak m^\flat$ (the zeroth coordinate is multiplicative and additive
modulo $\ell$, and $z=[\varpi^\flat]$), and the image of $E$ in $W(\overline{\mathbb F}_\ell)$ is $\ell$
because $z\mapsto[\bar\varpi^\flat]=0$. As $\theta(E)=0$ (proof of Lemma~\ref{pkp:chart-exists}),
Lemma~\ref{pkn:dist}(ii),(iii) apply; $\varphi^{-a}$ is an automorphism carrying $\varphi^a(E)\pkpA$ to
$E\pkpA$.
\end{proof}

This removes the appeal to ``the standard criterion'' in the proof of Lemma~\ref{pkp:chart-exists},
which is therefore proved in full.

\begin{lemma}[$\tau E$ and $\mu$]\label{pkn:mu}
Let $\mu_\varepsilon=[\varepsilon]-1=\varphi(\tau)$ and $\xi_\varepsilon=\sum_{k=0}^{\ell-1}[\varepsilon^{k/\ell}]$,
so that $\mu_\varepsilon=\tau\xi_\varepsilon$. Then $\xi_\varepsilon$ is distinguished, lies in
$\ker\theta$, and $\xi_\varepsilon=Eu$ with $u\in\pkpA^\times$; hence $\mu_\varepsilon=\tau Eu$.
Moreover, if $x\in\pkpA$ satisfies $\theta(\varphi^jx)=0$ for all $j\ge0$, then $\mu_\varepsilon\mid x$.
\end{lemma}

\begin{proof}
$\theta(\xi_\varepsilon)=\sum_{k<\ell}\zeta_\ell^k=0$. Since $x\mapsto x_0$ is the reduction $\pkpA\to\pkpA/\ell=O^\flat$, a ring map, its zeroth coordinate is
$\sum_k\varepsilon^{k/\ell}=(\varepsilon-1)/(\varepsilon^{1/\ell}-1)$, which lies in $\mathfrak m^\flat$ because
$v^\flat(\varepsilon-1)=\ell/(\ell-1)>v^\flat(\varepsilon^{1/\ell}-1)=1/(\ell-1)$ (as
$(\varepsilon^{1/\ell^n}-1)^\sharp=\lim_m(\zeta_{\ell^{m+n}}-1)^{\ell^m}$ has valuation
$\ell^{-n}\cdot\ell/(\ell-1)$). Its image in $W(\overline{\mathbb F}_\ell)$ is $\sum_k1=\ell$. By
Lemma~\ref{pkn:dist}(iv) and Corollary~\ref{pkn:E-gen}, $\xi_\varepsilon=Eu$ with $u$ a unit.

For the last claim, put $\xi_j=\varphi^{-j}(\xi_\varepsilon)$, a distinguished generator of
$\ker(\theta\circ\varphi^j)$ (Lemma~\ref{pkn:dist}(iii) and the argument of Corollary~\ref{pkn:E-gen}).
\emph{Claim: $\prod_{j=0}^n\xi_j$ divides $x$.} By induction on $n$: if $x=\xi_0\cdots\xi_{n-1}y$, apply
$\theta\circ\varphi^n$, whose kernel is $\xi_n\pkpA$: $\theta\varphi^n(\xi_j)=\theta(\varphi^{n-j}\xi_\varepsilon)
\ne0$ for $j<n$ (Lemma~\ref{pkn:dist}(iii)) and $O_C$ is a domain, so $\theta\varphi^n(y)=0$ and
$\xi_n\mid y$. Since $\varphi^{-j}(\mu_\varepsilon)=\varphi^{-j-1}(\mu_\varepsilon)\xi_j$, the product is
$\mu_\varepsilon/\varphi^{-n-1}(\mu_\varepsilon)$, so $\varphi^{-n-1}(\mu_\varepsilon)x\in\mu_\varepsilon\pkpA$.
Reduce modulo $\ell$: in $O^\flat$, $(\varepsilon-1)$ divides $(\varepsilon^{1/\ell^{n+1}}-1)x_0$, whence
$v^\flat(x_0)\ge\frac{\ell}{\ell-1}(1-\ell^{-n-1})$ for all $n$, so $v^\flat(x_0)\ge v^\flat(\varepsilon-1)$
and $x_0=(\varepsilon-1)c$. Then $x=\mu_\varepsilon[c]+\ell x'$ for some $x'$, and $x'$ again satisfies the
hypothesis: $\theta\varphi^j(\mu_\varepsilon)=\theta([\varepsilon^{\ell^j}])-1=1-1=0$, so $\ell\,\theta\varphi^j(x')=0$, and $O_C$ is torsion free. Iterating and using $\ell$-adic
completeness as in Lemma~\ref{pkn:dist}(ii), $x\in\mu_\varepsilon\pkpA$.
\end{proof}

\begin{lemma}[The primes $p_a$ in $\pknB$]\label{pkn:BdR}
The image of $E$ in $\pknBp$ generates $\mathrm{Fil}^1$, so $E^n\pknBp=\mathrm{Fil}^n$ ($n\ge0$) and
$E^n\mathrm{Fil}^m=\mathrm{Fil}^{m+n}$ ($n,m\in\Z$, $E$ being invertible in $\pknB$). For $c\in\Z\smallsetminus
\{0\}$ the image of $\varphi^c(E)$ is a unit of $\pknBp$. Consequently, for $a,b\ge0$, $n\in\Z$ and
$y\in\pknB$: $\varphi^{-a}(p_b)^n y\in\mathrm{Fil}^{m+n[a=b]}$ iff $y\in\mathrm{Fil}^m$.
\end{lemma}

\begin{proof}
By Corollary~\ref{pkn:E-gen}, $\ker(\theta[1/\ell])=E\pkpA[1/\ell]$, so $\pknBp$ is the $E$-adic
completion of $\pkpA[1/\ell]$ and its maximal ideal $\ker(\theta\colon\pknBp\to C)$ is $E\pknBp$. By
Fontaine's theorem, $\pknBp$ is a complete discrete valuation ring with uniformiser $t$ \cite[1.5.2, 1.5.4]{pkfFontaineII}, so $E$ and
$t$ are both uniformisers and $E^n\pknBp=t^n\pknBp=\mathrm{Fil}^n$. If $c\ne0$, $\theta(\varphi^cE)\ne0$
(Corollary~\ref{pkn:E-gen}), so $\varphi^cE\notin\ker\theta$ and its image is a unit of the local ring
$\pknBp$. Finally $\varphi^{-a}(p_b)=\varphi^{b-a}(E)$: it is $E$ if $a=b$ and a unit otherwise.
\end{proof}

\subsubsection{Items (M), (E) and (C)(b)}\label{pkn:ss-ME}

\begin{lemma}[Lemma~\ref{pkm:inputs}(M), first two assertions]\label{pkn:M}
Let $L$ be a lattice with jump sets. (i) $\pknMinf(L)$ and $\pkpA\mathcal N(L)$ are $\pkpA$-submodules of
$\pkpAmb(L)$, and $\mathcal N(L)$ is an $\mathfrak S$-submodule. (ii) Every $x\in\pkpA\mathcal N(L)$ has
$rx\in\pkpA\otimes L$ for some product $r$ of a power of $\tau$ and of elements $p_b$. (iii)
$\pkpA\mathcal N(L)$ and $\mathcal N(L)$ are stable under multiplication by $p_b$ and by $p_b^{-1}\in
\pkpW^\times$ ($b\ge0$).
\end{lemma}

\begin{proof}
(i) For $b\in\pkpA$ and $x,x'\in\pknMinf$ with $\tau^nx,\tau^{n'}x'\in\pkpA\otimes L$:
$\tau^{n+n'}(bx+x')\in\pkpA\otimes L$; and $\mu(\varphi^k(bx))=\varphi^k(b)\mu(\varphi^kx)\in\pknBp$, as
$\varphi^k$ is semilinear on $\pkpAmb(L)=\pkpW\otimes L$ and $\mu$ is $\pkpA$-linear. For
$\pkpA\mathcal N$ use products $rr'$ of the two multipliers. For $\mathcal N$ and $s\in\mathfrak S$:
$\mu(sx)=s\mu(x)\in\pknK$ since $\pknK$ is a field containing the image of $\mathfrak S$.
(ii) If $rx\in\pknMinf$ then $\tau^nrx\in\pkpA\otimes L$ for some $n$. (iii) $p_b\in\mathfrak
S\subseteq\pkpA$, so stability under $p_b$ is (i). The $p_b$ are units of $\pkpW$ (their zeroth
coordinate $\varphi^b(E)_0=E_0^{\ell^b}$ is a nonzero element of the field $C^\flat$). If $rx\in\pknMinf$,
then $(rp_b)(p_b^{-1}x)=rx$, so $p_b^{-1}x\in\pkpA\mathcal N$. If moreover $x\in\mathcal N$, then for every
form $\mu$, $p_b\,\mu(p_b^{-1}x)=\mu(x)\in\pknK$, and the image of $p_b$ in $\pknB$ is a nonzero element
of $\pknK$ (Lemma~\ref{pkn:BdR}), so $\mu(p_b^{-1}x)\in\pknK$.
\end{proof}

\begin{lemma}[Lemma~\ref{pkm:inputs}(E)]\label{pkn:E}
For $b\ge0$, $n\in\Z$ and $x\in L_i(o)$, $p_b^nx\in L_i(o+n[a=b])$, where $(o+n[a=b])(a,u)=
o(a,u)+n[a=b]$.
\end{lemma}

\begin{proof}
$p_b^nx\in\pkpA\mathcal N(L)$ by Lemma~\ref{pkn:M}(iii), and it lies in $\pkpAmb(L)_i$ because $p_b$ acts
on the first tensor factor and commutes with $1\otimes\zeta$ and with the scalars $[i']$. Let $a\ge0$,
$u\in U(i^{(-a)})$ and $\mu$ a form of type $u$ at $i^{(-a)}$. Since $\varphi^{-a}$ is semilinear and
$\mu$ is $\pkpA$-linear, $\varphi^{-a}(p_b)^{|n|}$ (an element of $\pkpA$) satisfies, for $n\ge0$,
$\mu(\varphi^{-a}(p_b^nx))=\varphi^{-a}(p_b)^n\mu(\varphi^{-a}x)$, and for $n<0$,
$\varphi^{-a}(p_b)^{-n}\mu(\varphi^{-a}(p_b^nx))=\mu(\varphi^{-a}x)$; in both cases
$\mu(\varphi^{-a}(p_b^nx))=\varphi^{-a}(p_b)^n\mu(\varphi^{-a}x)$ in the field $\pknB$ (the image of
$\varphi^{-a}(p_b)$ is nonzero by Lemma~\ref{pkn:BdR}). By Lemma~\ref{pkn:BdR} this lies in
$\mathrm{Fil}^{o(a,u)+n[a=b]}$.
\end{proof}

\emph{Item (C)(b) of Lemma~\ref{pkm:inputs}} is the case $l=1$ of the last condition in
Definition~\ref{pkp:def-chart}.

\subsubsection{Item (C)(a)}\label{pkn:ss-Ca}

\begin{proposition}[Lemma~\ref{pkm:inputs}(C)(a)]\label{pkn:Ca}
For $g\in\Gamma'$ and $s\in\mathfrak S$, $z\tau E$ divides $gs-s$ in $\pkpA$.
\end{proposition}

\begin{proof}
\emph{Step 1: $z\mu_\varepsilon\mid gz-z$.} $g$ acts on $\pkpA=W(O^\flat)$ through its action on
$O^\flat$, so $gz=[g\varpi^\flat]$. As $\theta$ is $\Gamma$-equivariant and $j(\Gamma')$ fixes
$\theta(z)$, $(g\varpi^\flat)^\sharp=g(\varpi^{\flat\sharp})=\varpi^{\flat\sharp}=\pi\ne0$. So $\eta=
g\varpi^\flat/\varpi^\flat\in C^\flat$ has $\eta^\sharp=1$ ($\sharp$ is multiplicative), hence
$v^\flat(\eta)=0$, $\eta\in O^{\flat\times}$, and $gz-z=z([\eta]-1)$. For $j\ge0$,
$\theta(\varphi^j([\eta]-1))=(\eta^{\ell^j})^\sharp-1=0$, so $\mu_\varepsilon\mid[\eta]-1$ by
Lemma~\ref{pkn:mu}.

\emph{Step 2.} By Lemma~\ref{pkn:mu}, $\mu_\varepsilon=\tau Eu$ with $u$ a unit, so $z\tau E\mid gz-z$.
Write $s=\sum_{k\ge0}\lambda_kz^k$ with $\lambda_k\in\Lambda$ (Definition~\ref{pkp:def-chart}). $g$ fixes
$\Lambda$ and is continuous, so $gs-s=\sum_{k\ge1}\lambda_k((gz)^k-z^k)=(gz-z)\sum_{k\ge1}\lambda_k
\sum_{j=0}^{k-1}(gz)^jz^{k-1-j}$, where the last series converges in $\pkpA$: its $k$-th term is
$\lambda_k$ times a sum of $k$ Teichm\"uller lifts $[(g\varpi^\flat)^j(\varpi^\flat)^{k-1-j}]=[\eta^j
(\varpi^\flat)^{k-1}]$, i.e.\ $\lambda_k[\varpi^\flat]^{k-1}w_k$ with $w_k=\sum_{j<k}[\eta]^j\in\pkpA$, and
$\sum_k\lambda_kz^{k-1}w_k$ converges in the same topology in which $\sum_k\lambda_kz^k$ does (the
$(\ell,z)$-adic one, for which $\pkpA$ is complete and separated) because $\lambda_kw_k\in\pkpA$. Hence
$gz-z$, and so $z\tau E$, divides $gs-s$.
\end{proof}

\subsubsection{Frames and their consequences: (M), first half of (H)}\label{pkn:ss-frame}

\emph{Period rings.} Besides $\pknBp\subset\pknB$ we use two rings $B^+$ and $B_{\max}$ with
$\Gamma$-equivariant ring maps $\pkpA\to B^+\to B_{\max}\to\pknB$ and $B^+\to\pknBp$, the two maps
$B^+\to\pknB$ agreeing, the composite $\pkpA\to B^+\to\pknBp$ being the natural map (so that an element of $\pkpA$, such as $E$, has the same image in $\pknBp$ whether or not it is first viewed in $B^+$; for the rings of \S\ref{pkf:ss-rings} this is part of (P1)), and Frobenius endomorphisms $\varphi$ of $B^+$ and $B_{\max}$ such that
$\pkpA\to B^+\to B_{\max}$ are $\varphi$-equivariant, commuting with $\Gamma$ (any choice with these properties works in this subsection; the existence of frames,
\S\ref{pkf:sec}, uses $B^+=A_{\max}[1/\ell]$ and $B_{\max}=B^+[1/\mu]$ of \S\ref{pkf:ss-rings}). The
bicommutant $\mathfrak S_B\subseteq B^+$ of the chart is defined in Definition~\ref{pkf:def-SB}; by
Proposition~\ref{pkf:closure} it is the $\mathfrak S$-subalgebra obtained from the closure of
$\mathfrak S[E/\ell]$ (for the topology with basis $\ell^nA_{\max}$) by inverting $\ell$; all
that is used below is that $\mathfrak S_B$ is a subring of $B^+$ whose image in $\pknB$ lies in
$\pknK$ (true by the definition of $\pknK$). The map $\pkpA\to\pknBp$ is injective (Fontaine \cite[4.1.1]{pkfFontaineII}), so a nonzero element of $\pkpA$ has nonzero, hence invertible, image in the field
$\pknB$.

Fix a lattice with jump sets $L$, a $\Z_\ell$-basis $e_1,\dots,e_d$ of $L$, and write $\rho(g)\in
GL_d(\Z_\ell)$ for the matrix of $g\in\Gamma$ ($ge_j=\sum_k\rho(g)_{kj}e_k$). For a form $\mu$ let
$\mu(\mathbf e)=(\mu(1\otimes e_1),\dots,\mu(1\otimes e_d))$, a row vector over $\pknB$. Matrices over
$B^+$, $B_{\max}$, $\pkpA$ are mapped to $\pknB$ entrywise when multiplied with such rows.

\begin{definition}[Frame]\label{pkn:def-frame}
A \emph{frame} of $L$ (for an integer $a\ge0$) consists of $M\in GL_d(B_{\max})$, $U,Y\in
M_d(B^+)$, $Z\in M_d(\mathfrak S_B)$, $Q,Q'\in M_d(\pkpA)$ and vectors $n_1,\dots,n_d\in\pkpAmb(L)$ with:
\begin{enumerate}
\item[(Fr1)] $\rho(g)\,g(M)=M$ for $g\in\Gamma'$;
\item[(Fr2)] the image of $U$ in $B_{\max}$ is $(E\tau)^aM$;
\item[(Fr3)] $QQ'=\tau^{2a}\cdot1$, $UZ=E^{2a}Q$ (in $B^+$), $YZ=E^{2a}\cdot1$;
\item[(Fr4)] $\tau^an_i=\sum_jQ_{ji}(1\otimes e_j)$ for every $i$.
\end{enumerate}
\end{definition}

($\tau$ is a unit of $\pkpW$, so given $Q$ the $n_i$ exist and are unique.)

Fix a frame. For $m\ge0$ and a form $\mu$ put $c^{(m)}(\mu)=\mu(\mathbf e)\,\varphi^m(M)$ (a row over
$\pknB$).

\begin{lemma}\label{pkn:frame-id}
(i) For every form $\mu$ and $m\ge0$, the entries of $c^{(m)}(\mu)$ lie in $K'$; in particular in
$\pknBp\cap\pknK$.
(ii) For every form $\mu$, $m\ge0$ and $t\in\pkpA^d$:
$\ \varphi^m(E)^a\,\mu\bigl(\sum_it_i\varphi^m(n_i)\bigr)=c^{(m)}(\mu)\,\varphi^m(Z)\,t.$
(iii) For every $\nu$ and $m$ there is a form $\sigma$ with $c^{(m)}(\sigma)$ the $\nu$-th unit row.
(iv) (Coefficient test.) Let $W\subseteq\pknB$ be a $B^+$-submodule, $m\ge0$, $t\in\pkpA^d$, and suppose
$\mu(\sum_it_i\varphi^m(n_i))\in W$ for all forms $\mu$. Then the image of $\varphi^m(E)^{2a}t_i$ lies in
$W$ for every $i$.
\end{lemma}

\begin{proof}
(i) Applying $\varphi^m$ to (Fr1) ($\varphi$ commutes with $\Gamma$ and fixes $\rho(g)$) gives
$\rho(g)g(M_m)=M_m$ for $M_m=\varphi^m(M)$. For $g\in\Gamma'$, $g(\mu(1\otimes e_j))=\mu(1\otimes ge_j)=
\sum_k\mu(1\otimes e_k)\rho(g)_{kj}$, i.e.\ $g(\mu(\mathbf e))=\mu(\mathbf e)\rho(g)$; hence
$g(c^{(m)}(\mu))=\mu(\mathbf e)\rho(g)g(M_m)=c^{(m)}(\mu)$. An element $y\in\pknB$ fixed by $G_{K'}$ lies in
$K'$: by Lemma~\ref{pkm:BdR} with $n=1$ there is $c\in K'$ with $y-c\in\mathrm{Fil}^1$; $y-c$ is fixed, so
by the second assertion of Lemma~\ref{pkm:BdR} (with $m=1$, $c=0$) $y-c\in\mathrm{Fil}^n$ for every $n$,
and $\bigcap_n\mathrm{Fil}^n=0$. Finally $K'\subset\pknBp$ and $K'\subseteq\pknK$.

(ii) Apply $\varphi^m$ to (Fr4); $\varphi$ fixes $1\otimes e_j$: $\varphi^m(\tau)^a\sum_it_i\varphi^m(n_i)=
\sum_j(\varphi^m(Q)t)_j(1\otimes e_j)$. Apply $\mu$ ($\pkpA$-linear) and multiply by
$\varphi^m(E)^{2a}$; by (Fr3), $\varphi^m(E)^{2a}\varphi^m(Q)=\varphi^m(U)\varphi^m(Z)$, and by (Fr2) the
image of $\varphi^m(U)$ in $\pknB$ is that of $\varphi^m(E\tau)^a\varphi^m(M)$. So
$\varphi^m(E)^{2a}\varphi^m(\tau)^a\mu(\sum_it_i\varphi^m n_i)=\varphi^m(E\tau)^a\,c^{(m)}(\mu)\varphi^m(Z)t$.
Cancel $\varphi^m(E\tau)^a$, nonzero in $\pknB$.

(iii) Let $\lambda$ be a $\pknB$-linear functional on $V=\pknB\otimes_{\pkpA}\pkpW$ with $\lambda(1\otimes1)
=1$; it exists as soon as $1\otimes1\ne0$. For this, let $T=\pkpA\smallsetminus\{0\}$; $T$ maps to
nonzero elements of the domains $\pknB$ and $\pkpW$, so $T^{-1}\pknB$ and $T^{-1}\pkpW$ are nonzero vector
spaces over $F=\operatorname{Frac}\pkpA$, and $T^{-1}V=T^{-1}\pknB\otimes_FT^{-1}\pkpW$ is a tensor product of nonzero vector spaces over a
field, in which the image of $1\otimes1$ is a tensor of two nonzero vectors, hence nonzero; so
$1\otimes1\ne0$ in $V$. Put $\lambda_0(w)=\lambda(1\otimes w)$, an
$\pkpA$-linear map $\pkpW\to\pknB$ with $\lambda_0(1)=1$, and $\sigma(\sum_jw_j\otimes e_j)=\sum_j
\lambda_0(w_j)(M_m^{-1})_{\nu j}$ ($\pkpAmb(L)=\bigoplus_j\pkpW\otimes e_j$). Then $\sigma(\mathbf e)$ is
the $\nu$-th row of $M_m^{-1}$, so $c^{(m)}(\sigma)$ is the $\nu$-th unit row. $\sigma$ is a form: from
$\rho(g)g(M_m)=M_m$, $g(M_m^{-1})=M_m^{-1}\rho(g)$, i.e.\ $g(\sigma(\mathbf e))=\sigma(\mathbf e)\rho(g)$,
which says $\sigma(1\otimes gy)=g\sigma(1\otimes y)$ for $y$ in the basis, hence for all $y\in L$ by
$\Z_\ell$-linearity ($\rho(g)$ has entries in $\Z_\ell$).

(iv) Let $\sigma_\nu$ be as in (iii) and $x=\sum_it_i\varphi^m(n_i)$. By (ii),
$\varphi^m(E)^a\sigma_\nu(x)=(\varphi^m(Z)t)_\nu$. Hence
$\sum_\nu\varphi^m(Y)_{i\nu}\allowbreak \varphi^m(E)^a\sigma_\nu(x)=(\varphi^m(YZ)t)_i=\varphi^m(E)^{2a}t_i$ by (Fr3).
The left side lies in $W$, as $\varphi^m(Y)_{i\nu}\varphi^m(E)^a\in B^+$ and $\sigma_\nu(x)\in W$.
\end{proof}

\begin{proposition}[Frame theorem, part]\label{pkn:frame}
Suppose $L$ has a frame (for some $a$). Then:
\begin{enumerate}
\item[(i)] $n_i\in\pknMinf(L)\cap\mathcal N(L)$, and the $n_i$ are linearly independent over $\pkpA$;
\item[(ii)] $\pknMinf(L)=\sum_i\pkpA n_i$;
\item[(iii)] $\tau^a\pknMinf(L)\subseteq\pkpA\otimes L$ and $\tau^a(\pkpA\otimes L)\subseteq\pknMinf(L)$;
\item[(iv)] $\pkpA\mathcal N(L)$ is the $\pkpA$-span of $\mathcal N(L)$ (Lemma~\ref{pkm:inputs}(M), third
assertion);
\item[(v)] $\tau^a(1\otimes y)\in\pknMinf(L)\subseteq\pkpA\mathcal N(L)$ for $y\in L$
(Lemma~\ref{pkm:inputs}(H), first assertion, with $H=a$).
\end{enumerate}
\end{proposition}

\begin{proof}
\emph{$n_i\in\pknMinf$.} $\tau^an_i\in\pkpA\otimes L$ by (Fr4). For $m\ge1$ and a form $\mu$,
Lemma~\ref{pkn:frame-id}(ii) with $t$ the $i$-th unit vector gives $p_m^a\mu(\varphi^mn_i)=
c^{(m)}(\mu)\varphi^m(Z)_{\cdot i}\in\pknBp$ (entries of $c^{(m)}$ in $\pknBp$, entries of $\varphi^m(Z)$ in
$B^+$); $p_m$ is a unit of $\pknBp$ (Lemma~\ref{pkn:BdR}), so $\mu(\varphi^mn_i)\in\pknBp$.
\emph{$n_i\in\mathcal N$.} With $m=0$: $E^a\mu(n_i)=c^{(0)}(\mu)Z_{\cdot i}\in\pknK$ (Lemma~\ref{pkn:frame-id}(i)
and $Z\in M_d(\mathfrak S_B)$), and $E$ is a nonzero element of $\pknK$. Also $n_i\in\pknMinf\subseteq
\pkpA\mathcal N$ (take $r=1$).
\emph{Independence.} If $\sum_it_in_i=0$, Lemma~\ref{pkn:frame-id}(iv) with $W=0$, $m=0$ gives
$E^{2a}t_i\mapsto0$ in $\pknB$, so $t_i=0$ by injectivity.

\emph{$\tau$-test.} Let $t\in\pkpA^d$ and $x=\sum_it_in_i$, and suppose $\mu(\varphi^{m}x)\in\mathrm{Fil}^1$ for
all $m\ge1$ and all forms $\mu$. Then $\tau\mid t_i$ for all $i$. Indeed $\varphi^mx=\sum_i\varphi^m(t_i)
\varphi^m(n_i)$, and Lemma~\ref{pkn:frame-id}(iv) with $W=\mathrm{Fil}^1$ (a $B^+$-submodule, being a
$\pknBp$-submodule) gives $p_m^{2a}\varphi^m(t_i)\in\mathrm{Fil}^1$; for $r\in\pkpA$, $r\in\mathrm{Fil}^1$ iff
$\theta(r)=0$ ($\mathrm{Fil}^1=E\pknBp=\ker(\theta\colon\pknBp\to C)$, Lemma~\ref{pkn:BdR}), and
$\theta(p_m)\ne0$ (Corollary~\ref{pkn:E-gen}), so $\theta(\varphi^mt_i)=0$ for all $m\ge1$. Then $y=\varphi(t_i)$
has $\theta(\varphi^jy)=0$ for all $j\ge0$, so $\mu_\varepsilon=\varphi(\tau)$ divides $\varphi(t_i)$
(Lemma~\ref{pkn:mu}), and $\tau\mid t_i$.

(ii) $\supseteq$: (i) and Lemma~\ref{pkn:M}(i). $\subseteq$: first, for $x'=\sum_j\alpha_j(1\otimes e_j)\in
\pkpA\otimes L$, put $t=Q'\alpha$; by (Fr4) and (Fr3), $\tau^a\sum_it_in_i=\sum_j(QQ'\alpha)_j(1\otimes e_j)
=\tau^{2a}x'$, so $\tau^ax'=\sum_it_in_i$ ($\pkpAmb(L)$ has no $\tau$-torsion, $\tau\in\pkpW^\times$).
Now let $x\in\pknMinf$, with $\tau^Nx\in\pkpA\otimes L$; then $\tau^{a+N}x\in\sum_i\pkpA n_i$. We show:
if $k\ge0$ and $\tau^{k+1}x=\sum_it_in_i$, then $\tau^kx\in\sum_i\pkpA n_i$; by induction $x\in\sum\pkpA n_i$.
For $m\ge1$ and a form $\mu$, $\mu(\varphi^m(\tau^{k+1}x))=\varphi^m(\tau)\mu(\varphi^m(\tau^kx))$; the second
factor is in $\pknBp$ ($\tau^kx\in\pknMinf$), and $\theta(\varphi^m\tau)=\zeta_\ell^{\ell^m}-1=0$, so the first
is in $\mathrm{Fil}^1$. By the $\tau$-test, $t_i=\tau t'_i$, and $\tau(\tau^kx-\sum_it'_in_i)=0$ gives
$\tau^kx=\sum_it'_in_i$.

(iii) The first inclusion follows from (ii) and (Fr4); the second was shown in the proof of (ii)
(together with (i) and Lemma~\ref{pkn:M}(i)).

(iv) The span of $\mathcal N$ is contained in the $\pkpA$-module $\pkpA\mathcal N$. Conversely let
$x\in\pkpA\mathcal N$ and $r=p_{b_1}\cdots p_{b_k}$ with $rx\in\pknMinf$. By (ii), $rx=\sum_it_in_i$, so
$x=\sum_it_i\,(r^{-1}n_i)$ with $r^{-1}n_i\in\mathcal N$ by (i) and Lemma~\ref{pkn:M}(iii).

(v) is the second inclusion of (iii) with $x'=1\otimes y$.
\end{proof}

\subsubsection{Item (C)(c)}\label{pkn:ss-Cc}

\begin{proposition}[Lemma~\ref{pkm:inputs}(C)(c)]\label{pkn:Cc}
For every $n\ge1$ there are $\alpha,\beta\in\pkpA$ with $\tau E(\ell^{v_\ell(n)}+\beta z)=
\alpha([\varepsilon]^n-1)$.
\end{proposition}

\begin{proof}
Write $n=m\ell^v$ with $v=v_\ell(n)$, $\ell\nmid m$, and $x=[\varepsilon]^{\ell^v}$.
\emph{Step 1: $x-1=\tau\prod_{a=0}^{v}\varphi^a(E)\cdot u$ with $u\in\pkpA^\times$.} With
$\xi_\varepsilon$ as in Lemma~\ref{pkn:mu}, $\varphi(\tau)=\mu_\varepsilon=\tau\xi_\varepsilon$, so
$\varphi^{a+1}(\tau)=\varphi^a(\tau)\varphi^a(\xi_\varepsilon)$ and, by induction,
$\varphi^{v+1}(\tau)=\tau\prod_{a=0}^v\varphi^a(\xi_\varepsilon)$. Now $\varphi^{v+1}(\tau)=[\varepsilon]^{\ell^v}-1
=x-1$ ($\varphi[b]=[b^\ell]$), and $\varphi^a(\xi_\varepsilon)=\varphi^a(E)\varphi^a(u_0)$ with $u_0$ the unit
of Lemma~\ref{pkn:mu}; $\varphi^a$ maps units to units.
\emph{Step 2.} $[\varepsilon]^n-1=x^m-1=(x-1)w$ with $w=\sum_{k<m}x^k$, and $w$ is a unit: its image in
$W(\overline{\mathbb F}_\ell)$ is $m$ (as $\varepsilon\mapsto1$ in $O^\flat/\mathfrak m^\flat$), a unit since
$\ell\nmid m$. Hence $\tau\prod_{a=0}^v\varphi^a(E)=\alpha([\varepsilon]^n-1)$ with $\alpha=(uw)^{-1}$.
\emph{Step 3.} For $a\ge1$, $\varphi^a(E)=\varphi^a(z)^e\varphi^a(A)+\ell=z^{e\ell^a}\varphi^a(A)+\ell$ with
$\varphi^a(A)\in\mathfrak S$ ($\varphi(z)=z^\ell$, $\varphi(\Lambda)=\Lambda$), so $\varphi^a(E)\equiv\ell
\pmod{z\mathfrak S}$ and $\prod_{a=1}^v\varphi^a(E)=\ell^v+\beta z$ with $\beta\in\mathfrak S$. Since
$\varphi^0(E)=E$, Step 2 reads $\tau E(\ell^v+\beta z)=\alpha([\varepsilon]^n-1)$.
\end{proof}

\emph{Where the remaining items are proved.} The existence of a frame for every lattice with jump
sets, with $a$ the jump bound, is Theorem~\ref{pkf:main}, proved in \S\ref{pkf:sec} modulo published
facts of $\ell$-adic Hodge theory and the ring statements (A1)--(A8), (P4) proved in \S\ref{pkr:sec};
hence items (M) (third assertion) and (H) (first assertion) of Lemma~\ref{pkm:inputs} hold. Items (G)
and (H) (in the form stated in \S\ref{pkl:sec}), (F), (R), (D), (C)(d), (G$_0$) and ($\Phi$) are proved in
\S\ref{pkl:sec}; Lemma~\ref{pkp:orders-add} and Lemma~\ref{pkp:tau}(i), (iv), (v) in \S\ref{pkt:sec};
Proposition~\ref{pki:chart} in \S\ref{pkj:sec}.



\providecommand{\auO}{\mathcal O}
\providecommand{\auQlb}{\overline{\Q}_\ell}
\providecommand{\pkpA}{A_{\mathrm{inf}}}
\providecommand{\pkpW}{W(C^\flat)}
\providecommand{\pkpAmb}{\mathbb A}
\providecommand{\pknBp}{B^+_{\mathrm{dR}}}
\providecommand{\pknB}{B_{\mathrm{dR}}}
\providecommand{\pkfAm}{A_{\max}}
\providecommand{\pkfAo}{A_{\max}^{\circ}}
\providecommand{\pkfBp}{B^+}
\providecommand{\pkfBm}{B_{\max}}
\providecommand{\pkfSB}{\mathfrak S_B}
\providecommand{\pkfEnd}{\mathcal G}

\subsection{Existence of frames}\label{pkf:sec}

This section proves that every lattice with jump sets has a frame in the sense of
Definition~\ref{pkn:def-frame}, with $a$ the jump bound (Theorem~\ref{pkf:main}). We use the
conventions of \S\S\ref{pkp:sec}--\ref{pkn:sec}: $\ell$ is the residue characteristic,
$\pkpA=W(O^\flat)$, $\theta\colon\pkpA\to O_C$, $\varphi$ the (bijective) Frobenius of $\pkpA$, a fixed
Kummer chart (Definition~\ref{pkp:def-chart}) with $\mathfrak S=\Lambda[[z]]\subset\pkpA$,
$E=z^eA+\ell$, $p_a=\varphi^a(E)$, the subgroup $\Gamma'$ and $K'=\auQlb^{\,j(\Gamma')}$;
$\varepsilon=(1,\zeta_\ell,\zeta_{\ell^2},\dots)$, $\mu=[\varepsilon]-1$, $\tau=\varphi^{-1}(\mu)=
[\varepsilon^{1/\ell}]-1$. By Lemma~\ref{pkn:mu}, $\mu=E\tau v$ with $v\in\pkpA^\times$; we fix this
$v$. For $x\in\pkpA$ we write $\bar x\in O^\flat=\pkpA/\ell\pkpA$ for its reduction.

\emph{Dependencies.} Theorem~\ref{pkf:main} is proved in full \emph{modulo} the published facts of
$\ell$-adic Hodge theory (P1)--(P3), (P5), (P6) stated in \S\ref{pkf:ss-rings}, cited with author and
work. The eight auxiliary statements (A1)--(A8) about the rings $\pkfAm$, $\pkfBp$ and the chart, and
the statement (P4), are proved in \S\ref{pkr:sec} (which also uses the standard facts (W1)--(W4)
recalled there). Every other step is proved here. (The labels (A1)--(A8) are distinct from the case
inputs (I1)--(I6) of \S\S\ref{pkc:sec}--\ref{pkd:sec}.)

\subsubsection{Period rings, the representation, and the inputs}\label{pkf:ss-rings}

\emph{The rings.} $E$ is distinguished and generates $\ker\theta$
(Corollary~\ref{pkn:E-gen}). Let $R=\pkpA[X]/(\ell X-E)$ and let $\pkfAm$ be its $\ell$-adic
completion; write $E/\ell$ for the image of $X$. Put $\pkfBp=\pkfAm[1/\ell]$ and
$\pkfBm=\pkfBp[1/\mu]$. (Replacing $E$ by another generator $Eu$ of $\ker\theta$, $u\in\pkpA^\times$,
gives an isomorphic ring via $X\mapsto Xu$; so these are Colmez's rings $A_{\max}$, $B^+_{\max}$ and
$B_{\max}$.) We write $\pkfAo$ for the image of $\pkfAm$ in $\pkfBp$.

\begin{enumerate}
\item[(P1)] (Structure; Colmez \cite[\S III.2]{pkfColmez98}, Fontaine \cite[4.1]{pkfFontaineII}.)
$\pkfAm$ is $\ell$-torsion free, $\ell$-adically complete and separated; the maps
$\pkpA\to\pkfAm\to\pkfBp\to\pkfBm\to\pknB$ and $\pkfBp\to\pknBp$ are injective ring maps, the two maps
$\pkfBp\to\pknB$ agree, and the composite $\pkpA\to\pknBp$ is the natural map; $\pkfBp$ and $\pkfBm$
are domains. $\varphi$ and $G_{\Q_\ell}$ act on $\pkfAm$, $\pkfBp$, $\pkfBm$ compatibly with these maps
(the action of $G_{\Q_\ell}$ also on $\pknB$), $\varphi$ is injective and commutes with $G_{\Q_\ell}$.
In particular $\varphi$ and every $g\in G_{\Q_\ell}$ map $\pkfAo$ into itself.
\item[(P2)] (Crystalline representations; Fontaine \cite[5.1.2]{pkfFontaineIII}, Colmez \cite[\S\S III.2, III.4]{pkfColmez98}.) $\pkfBm^{G_{K'}}=K'_0$, the maximal unramified subfield of $K'$; a continuous
representation of $G_{K'}$ on a finite dimensional $\Q_\ell$-vector space $V$ is crystalline iff
$\dim_{K'_0}(\pkfBm\otimes_{\Q_\ell}V)^{G_{K'}}=\dim V$, and then the natural map
$\pkfBm\otimes_{K'_0}(\pkfBm\otimes V)^{G_{K'}}\to\pkfBm\otimes V$ is an isomorphism. Crystalline
representations are de Rham. (Fontaine proves $B_{\mathrm{cris}}^{G_{K'}}=K'_0$ and that $B_{\mathrm{cris}}$ is $G_{K'}$-regular; Colmez notes $\varphi(\pkfBm)\subseteq B_{\mathrm{cris}}\subseteq\pkfBm$ and that $\pkfBm$ and $B_{\mathrm{cris}}$ give the same $D_{\mathrm{cris}}$. Hence $K'_0\subseteq\pkfBm^{G_{K'}}$, and if $b\in\pkfBm^{G_{K'}}$ then $\varphi(b)\in B_{\mathrm{cris}}^{G_{K'}}=K'_0=\varphi(K'_0)$, so $b\in K'_0$ by injectivity of $\varphi$.)
\item[(P3)] (Filtered comparison for de Rham representations; Fontaine \cite[3.8]{pkfFontaineIII}.) If $V$ is de Rham, $D_{\mathrm{dR}}(V)=(\pknB\otimes V)^{G_{K'}}$ with the filtration
$\mathrm{Fil}^kD_{\mathrm{dR}}=D_{\mathrm{dR}}\cap(\mathrm{Fil}^k\pknB\otimes V)$, then
$\pknB\otimes_{K'}D_{\mathrm{dR}}\to\pknB\otimes V$ is an isomorphism carrying
$\sum_{k}\mathrm{Fil}^{-k}\pknB\otimes\mathrm{Fil}^kD_{\mathrm{dR}}$ onto $\pknBp\otimes V$. The
\emph{Hodge--Tate jumps} of $V$ (at the identity embedding) are the $k$ with
$\mathrm{gr}^kD_{\mathrm{dR}}\ne0$; for $V=L[1/\ell]$, $L$ a lattice with jump sets viewed as a
$\Z_\ell$-representation, these are the Hodge--Tate jumps of $L$ at all embeddings, so they lie in
$\bigcup_iU(i)$ (Definition~\ref{pkp:def-orders}; this is the convention of the paper).
\item[(P4)] (Proved in Proposition~\ref{pkr:P4}, from (P1), (P6).) If $b\in\pkfBm$ and $m\ge0$ are such that the image of
$\varphi^n(b)$ in $\pknB$ lies in $\mathrm{Fil}^{-m}$ for every $n\ge0$, then $\mu^mb\in\pkfBp$.
\item[(P5)] (Ax--Sen--Tate; \cite[Theorem, p.~417]{pkfAx70}: for a field $k$ henselian for a non-trivial
valuation, $G_k$ acting on the completion of an algebraic closure of $k$, the fixed field is the completion of
the perfect closure of $k$; here $k=L$, which is henselian as an algebraic extension of $\Q_\ell$, of
characteristic $0$, and the completion of $\auQlb$ is $C$.) For every algebraic extension $L$ of $\Q_\ell$ inside
$\auQlb$, the elements of $C$ fixed by $\operatorname{Gal}(\auQlb/L)$ form the closure $\hat L$ of $L$ in
$C$. It is used only through Lemma~\ref{pkf:fixed}, with $H\subseteq G_{\Q_\ell}$ the group of elements
fixing $\mathfrak S$ pointwise. (The same lemma can be read off from the field of norms of Fontaine and
Wintenberger \cite{pkfWintenberger83} and the tilting of Scholze \cite{pkfScholze12}, but we do not use
them.)
\item[(P6)] $\mathrm{Fil}^n\pknB=E^n\pknBp$ (Lemma~\ref{pkn:BdR}, which rests on Fontaine's theorem
that $\pknBp$ is a discrete valuation ring with uniformiser $t$).
\end{enumerate}

\emph{The representation.} Let $L$ be a lattice with jump sets, $e_1,\dots,e_d$ a $\Z_\ell$-basis,
$\rho(g)\in GL_d(\Z_\ell)$ the matrix of $g\in\Gamma$ ($ge_j=\sum_k\rho(g)_{kj}e_k$). The hypothesis
used is: \emph{$L[1/\ell]$ is crystalline as a representation of $\Gamma'$ (through $j$, of
$G_{K'}$), with Hodge--Tate jumps in $[-a,a]$}, where $a\ge0$ is the \emph{jump bound},
$a=\max\{|u|:u\in\bigcup_iU(i)\}$. For a ring $B$ with an action of $G_{\Q_\ell}$, a matrix $M\in
M_d(B)$ has \emph{invariant columns} if $\rho(g)\,g(M)=M$ for $g\in\Gamma'$ ($g$ acting through $j$
entrywise); equivalently the vectors $\sum_kM_{kj}\otimes e_k$ of $B\otimes L$ are fixed by $g\otimes
g$. If $M$ has invariant columns so does $\varphi^n(M)$ for a ring with a Frobenius commuting with
$\Gamma'$, since $\varphi$ fixes $\rho(g)$.

\emph{Auxiliary inputs.} Write $(\ell,E)$ for the ideal of $\pkpA$ generated by $\ell,E$.
\begin{enumerate}
\item[(A1)] $\pkpA\cap\ell^n\pkfAm=(\ell,E)^n$ for $n\ge0$.
\item[(A2)] If $\delta\in\pkpA$ divides $\mu$, $h\ge0$, $w\in\pkfBp$ and $\delta^hw\in\pkfAo$, then
$\ell^{2h}w\in\pkfAo$.
\item[(A3)] (Contraction.) Let $V\subseteq\pkfAo$, $h\ge0$ and $u\in\pkpA$ be such that every $\alpha\in V$
has some $\alpha'\in V$ with $E^h\alpha=u\,\varphi(\alpha')$. Then $\ell^{h+1}V\subseteq\pkpA$.
\item[(A4)] (Approximation.) For $y\in\pkfBp$ and $h\ge0$ there are $m\ge0$, $\alpha\in\pkpA$ and
$w\in\pkfAo$ with $\ell^my=\alpha+\ell^m\mu^hw$.
\item[(A5)] (Topological basis.) There are a family $(b_i)_{i\in I}$ in $\pkpA$ and $\mathfrak S$-linear
maps $\kappa_i\colon\pkpA\to\mathfrak S$ such that: for every $x\in\pkpA$ and $N\ge0$, $x-\sum_{i\in
G}\kappa_i(x)b_i\in(\ell,z)^N\pkpA$ for all sufficiently large finite $G\subseteq I$ (in particular
$x=0$ if all $\kappa_i(x)=0$); and there is $0\in I$ with $\kappa_0(s)=s$ for $s\in\mathfrak S$.
\item[(A6)] (Freeness.) Let $P\subseteq\pkpA^d$ be an $\mathfrak S$-submodule containing no $d+1$
elements linearly independent over $\mathfrak S$, such that $za\in\ell P$ with $a\in P$ implies
$a\in\ell P$. Then $P$ is a free $\mathfrak S$-module of finite rank ($\le d$).
\item[(A7)] (Closed ideals.) If $\delta,c\in\pkpA$, $\ell\nmid\delta$, and for every $m$ there is
$z_m\in\pkpA$ with $c-\delta z_m\in(\ell,E)^m$, then $\delta\mid c$.
\item[(A8)] (Coordinates.) (a) In the notation of (A5): for $i\in I$, $x\in\pkpA$ and $m\ge0$ with
$\theta(\varphi^{-m}x)\in\ell O_C$, there is $s\in\mathfrak S$ with $p_m\mid\kappa_i(x)-\ell s$ in
$\mathfrak S$. (b) If $b\in\pkfBp$ and $\varphi(\tau)\mid\varphi(b)$ in $\pkfBp$, then $\tau\mid b$ in
$\pkfBp$. (c) Let $G_f$ be the extension of Lemma~\ref{pkf:extend}. For $\lambda\in K'^d$,
$y\in(\pkfBp)^d$ and an $\mathfrak S$-linear $f\colon\pkpA\to\pkpA$, $\sum_i\lambda_iy_i$ divides
$\sum_i\lambda_iG_f(y_i)$ in $\pknBp$ (images in $\pknBp$).
\end{enumerate}

The inputs (A1)--(A8) are proved in \S\ref{pkr:sec}: (A1) is Proposition~\ref{pkr:I1}, (A2)
Proposition~\ref{pkr:I2}, (A3) Proposition~\ref{pkr:I3}, (A4) Proposition~\ref{pkr:I4}, (A5)
Proposition~\ref{pkr:I5}, (A6) Proposition~\ref{pkr:I6}, (A7) Proposition~\ref{pkr:I7}, and (A8)(a),
(b), (c) Proposition~\ref{pkr:I8}. Those proofs use only (P1), (P6), Lemma~\ref{pkf:approx0},
Lemmas~\ref{pkf:SB-basic}, \ref{pkf:extend}, the description $K'=K_0(\pi)$ from the proof of
Lemma~\ref{pkf:invariants}, and \S\S\ref{pkp:sec}--\ref{pkn:sec}; none of them uses
Theorem~\ref{pkf:main} or a statement depending on it.

\begin{lemma}[Elementary facts on $\pkfAm$]\label{pkf:approx0}
(i) Every $\beta\in\pkfAm$ can be written $\beta=\sum_{k\ge0}\alpha_k(E/\ell)^k$ with $\alpha_k\in\pkpA$,
$\alpha_k\to0$ $\ell$-adically; for every $n$ there is a polynomial $P\in\pkpA[X]$ with
$\beta-P(E/\ell)\in\ell^n\pkfAm$. (ii) For $y\in\pkfBp$ and $n\ge0$ there are $m\ge0$, $\alpha\in\pkpA$,
$w\in\pkfAo$ with $\ell^my=\alpha+\ell^{m+n}w$. (iii) $\varphi(E)=\ell\cdot u_E$ with
$u_E\in\pkfAm^\times$; hence $\varphi(E)$ is a unit of $\pkfBp$. (iv) For $g\in G_{\Q_\ell}$, $gE=Er$ and
$g\tau=\tau r'$ with $r,r'\in\pkpA^\times$. (v) $\theta$ extends to $\pkfBp\to C$, and
$\ker(\theta|_{\pkfBp})=E\pkfBp$; consequently, for $x\in\pkfBp$ and $n\ge0$, $E^n\mid x$ in $\pkfBp$
iff the image of $x$ in $\pknBp$ lies in $\mathrm{Fil}^n$.
\end{lemma}

\begin{proof}
(i) The $\ell$-adic completion of $\pkpA[X]$ is the ring of series $\sum\alpha_kX^k$ with $\alpha_k\to0$
($\pkpA$ is $\ell$-adically complete), and it surjects onto the completion of the quotient $R$.
Truncating the series gives $P$.
(ii) $\ell^{m_0}y\in\pkfAo$ for some $m_0$. Write $\ell^{m_0}y=P(E/\ell)+\ell^{m_0+n}w$ by (i) with
$\deg P\le D$; then $\ell^DP(E/\ell)\in\pkpA$, and $m=m_0+D$ works.
(iii) $E=[E_0]+\ell E'$ with $E'\in\pkpA^\times$ (Definition~\ref{pkn:def-dist}). In $\pkfAm$,
$[E_0]=E-\ell E'=\ell(E/\ell-E')$, so $\varphi([E_0])=[E_0]^\ell=\ell^\ell(E/\ell-E')^\ell$ and
$\varphi(E)=\ell\bigl(\varphi(E')+\ell^{\ell-1}(E/\ell-E')^\ell\bigr)$. The bracket is a unit plus $\ell$
times an element, hence a unit of the $\ell$-adically complete ring $\pkfAm$.
(iv) $\theta(gE)=g\theta(E)=0$, so $E\mid gE$ (Corollary~\ref{pkn:E-gen}); likewise $E\mid g^{-1}E$,
i.e.\ $gE\mid E$; as $\pkpA$ is a domain, $gE=Er$ with $r$ a unit. For $\tau$: $\varphi(g\tau)=g\mu$ and
$\theta(\varphi^j(g\mu))=g\theta(\varphi^j\mu)=0$ for $j\ge0$, so $\mu\mid g\mu$ by Lemma~\ref{pkn:mu},
i.e.\ $\varphi(\tau)\mid\varphi(g\tau)$, and $\tau\mid g\tau$ since $\varphi$ is an automorphism of $\pkpA$;
symmetrically $g\tau\mid\tau$.
(v) $\theta$ extends to $\pkpA[X]\to O_C$ by $X\mapsto0$, which kills $\ell X-E$, and by continuity to
$\pkfAm$, then to $\pkfBp$; this is the restriction of $\theta\colon\pknBp\to C$ (P1). By (i),
$\beta=\alpha_0+(E/\ell)\beta'$ with $\beta'\in\pkfAm$, and $\theta(\beta)=\theta(\alpha_0)$. If
$\theta(\beta)=0$ then $\alpha_0=Ec$ (Corollary~\ref{pkn:E-gen}) and $\beta=E(c+\ell^{-1}\beta')\in
E\pkfBp$. Since $\mathrm{Fil}^1\cap\pkfBp=\ker\theta|_{\pkfBp}$, the case $n=1$ follows; for general $n$,
if $x\in\mathrm{Fil}^n$, $n\ge1$, then $x=Ex_1$ with $x_1\in\pkfBp$, and $x_1\in\mathrm{Fil}^{n-1}$
because $E$ is a uniformiser of $\pknBp$ (P6); induct.
\end{proof}

\subsubsection{The ring $\mathfrak S_B$ and extension of $\mathfrak S$-linear maps}\label{pkf:ss-SB}

Let $\pkfEnd$ be the set of $\mathfrak S$-linear maps $G\colon\pkfBp\to\pkfBp$ with
$G(\pkfAo)\subseteq\pkfAo$.

\begin{definition}[The ring $\mathfrak S_B$]\label{pkf:def-SB}
$\pkfSB=\{y\in\pkfBp:\ G(yx)=yG(x)\text{ for all }G\in\pkfEnd,\ x\in\pkfBp\}$.
\end{definition}

$\pkfSB$ is an $\mathfrak S$-subalgebra of $\pkfBp$ containing the image of $\mathfrak S$ (each $G$ is
$\mathfrak S$-linear). This is the bicommutant of \S\ref{pkn:ss-frame}; Proposition~\ref{pkf:closure}
shows that it is the $\mathfrak S$-algebra described there. Concretely (Proposition~\ref{pkf:closure}),
$\pkfSB=\widehat{\mathfrak S[E/\ell]}[1/\ell]$, where the hat denotes the $\ell$-adic closure in $\pkfAm$;
in the unramified case this ring embeds into $\Lambda[[z/\ell]][1/\ell]$ (Lemma~\ref{pkg:Shat}). It enters
only through the properties proved in this section and in \S\ref{pkg:ss-radius}; no comparison with
Breuil--Kisin modules is needed.

\begin{lemma}[Elementary properties]\label{pkf:SB-basic}
(i) If $v'\in\pkfSB\smallsetminus\{0\}$ and $v'u\in\pkfSB$ ($u\in\pkfBp$), then $u\in\pkfSB$. Hence
$\ell^{-1}\in\pkfSB$, $E/\ell\in\pkfSB$, and $\varphi(E)^{-1}\in\pkfSB$.
(ii) If $g\in G_{\Q_\ell}$ fixes $\mathfrak S$ pointwise, then $g$ fixes $\pkfSB$ pointwise.
\end{lemma}

\begin{proof}
(i) $v'G(ux)=G(v'ux)=v'uG(x)$ and $\pkfBp$ is a domain. Apply with $v'=\ell$ and $v'u=1$, with $v'=\ell$
and $v'u=E$, and with $v'=\varphi(E)\in\mathfrak S$ and $v'u=1$ (Lemma~\ref{pkf:approx0}(iii)).
(ii) $g$ is additive, $g(sx)=g(s)g(x)=sg(x)$ for $s\in\mathfrak S$, and $g(\pkfAo)\subseteq\pkfAo$ (P1),
so $g\in\pkfEnd$. For $y\in\pkfSB$: $g(y)=g(y\cdot1)=y\,g(1)=y$.
\end{proof}

\begin{lemma}[Density]\label{pkf:density}
Let $D\colon\pkfBp\to\pkfBp$ be additive with $D(\pkfAo)\subseteq\pkfAo$ and $D(\alpha)=0$ for
$\alpha\in\pkpA$. Then $D=0$.
\end{lemma}

\begin{proof}
Let $y\in\pkfBp$, $n\ge0$, and $\ell^my=\alpha+\ell^{m+n}w$ as in Lemma~\ref{pkf:approx0}(ii). An
additive map commutes with multiplication by integers, so $\ell^mD(y)=\ell^{m+n}D(w)$ and $D(y)=
\ell^nD(w)\in\ell^n\pkfAo$. As $\pkfAm$ is $\ell$-adically separated and $\pkfAm\to\pkfBp$ injective,
$D(y)=0$.
\end{proof}

\begin{lemma}[Extension]\label{pkf:extend}
Every $\mathfrak S$-linear $f\colon\pkpA\to\pkpA$ extends to some $G_f\in\pkfEnd$, i.e.\
$G_f(\alpha)=f(\alpha)$ for $\alpha\in\pkpA$. We fix such a $G_f$ for every $f$. For $y\in\pkfSB$ and
$x\in\pkfBp$, $G_f(yx)=yG_f(x)$.
\end{lemma}

\begin{proof}
\emph{Step 1: if $\alpha\in\pkpA\cap\ell^n\pkfAm$ then $f(\alpha)\in\ell^n\pkfAm$.} By (A1),
$\alpha\in(\ell,E)^n$, an ideal generated by the elements $\ell^iE^{n-i}$ of $\mathfrak S$; by
$\mathfrak S$-linearity $f(\alpha)\in(\ell,E)^n\subseteq\ell^n\pkfAm$.
\emph{Step 2.} Let $\mathcal D=\{\beta\in\pkfAm:\ \ell^k\beta\in\pkpA\text{ for some }k\}$, an
$\mathfrak S$-submodule. For $\beta\in\mathcal D$ with $\ell^k\beta\in\pkpA$, Step 1 gives
$f(\ell^k\beta)=\ell^k\gamma$ with $\gamma\in\pkfAm$, unique as $\pkfAm$ is $\ell$-torsion free, and
independent of $k$ since $f(\ell\cdot)=\ell f(\cdot)$; put $G^\circ(\beta)=\gamma$. $G^\circ$ is $\mathfrak S$-linear
on $\mathcal D$ and extends $f$. If moreover $\beta\in\ell^n\pkfAm$, then $\ell^k\beta\in\pkpA\cap
\ell^{k+n}\pkfAm$, so $\ell^k G^\circ(\beta)\in\ell^{k+n}\pkfAm$ by Step 1 and $G^\circ(\beta)\in\ell^n\pkfAm$.
By Lemma~\ref{pkf:approx0}(i), $\mathcal D$ (which contains all $P(E/\ell)$, $P\in\pkpA[X]$) is
$\ell$-adically dense in $\pkfAm$. As $\pkfAm$ is $\ell$-adically complete and separated, $G^\circ$ extends
uniquely to an $\ell$-adically continuous $\mathfrak S$-linear map $G^\circ\colon\pkfAm\to\pkfAm$.
\emph{Step 3.} $G_f(\ell^{-k}\beta)=\ell^{-k}G^\circ(\beta)$ is well defined on $\pkfBp=\pkfAm[1/\ell]$, is
$\mathfrak S$-linear, extends $f$ and maps $\pkfAo$ into itself. The last assertion is
Definition~\ref{pkf:def-SB}.
\end{proof}

\begin{lemma}[Values in $\mathfrak S_B$]\label{pkf:values}
Let $\sigma$ be either the identity or $\varphi$, and $f$ an $\mathfrak S$-linear map with
$\sigma(f(\alpha))\in\mathfrak S$ for all $\alpha\in\pkpA$. Then $\sigma(G_f(y))\in\pkfSB$ for all
$y\in\pkfBp$. In particular, with $\kappa_i$ as in (A5) and $G_i=G_{\kappa_i}$ (viewing $\kappa_i$ as a
map to $\pkpA$): $G_i(y)\in\pkfSB$ for all $y$; $G_0(y)=y$ for $y\in\pkfSB$; and
$\varphi(\pkfSB)\subseteq\pkfSB$.
\end{lemma}

\begin{proof}
Let $G\in\pkfEnd$ and $x\in\pkfAo$. The map $D(y)=G(\sigma(G_f(y))x)-\sigma(G_f(y))G(x)$ is additive and
maps $\pkfAo$ into $\pkfAo$ ($G_f$, $\sigma$, $G$ preserve $\pkfAo$, a ring). For $\alpha\in\pkpA$,
$s=\sigma(f(\alpha))\in\mathfrak S$ and $D(\alpha)=G(sx)-sG(x)=0$. By Lemma~\ref{pkf:density}, $D=0$. For
general $x\in\pkfBp$ take $n$ with $\ell^nx\in\pkfAo$ and cancel $\ell^n$ ($G$ is additive, $\pkfBp$ a
domain). So $\sigma(G_f(y))\in\pkfSB$. For $G_i$ take $\sigma=1$. For $y\in\pkfSB$: $G_0(y)=y\,G_0(1)=
y\,\kappa_0(1)=y$ by Lemma~\ref{pkf:extend}. Finally, for $y\in\pkfSB$, $\varphi(y)=\varphi(G_0(y))$ and
$\varphi(\kappa_0(\alpha))\in\varphi(\mathfrak S)\subseteq\mathfrak S$.
\end{proof}

\begin{proposition}[Description of $\mathfrak S_B$]\label{pkf:closure}
$\pkfSB$ is the set of $y\in\pkfBp$ for which there is $m$ such that for every $n$,
$\ell^my\in\mathfrak S[E/\ell]+\ell^n\pkfAo$. In particular $\pkfSB$ is the $\mathfrak S$-subalgebra
described in \S\ref{pkn:ss-frame}, and its image in $\pknB$ lies in $\mathcal K_0$.
\end{proposition}

\begin{proof}
If $y$ has the property, let $D=G_0-1$, which maps $\pkfAo$ into itself and is $\pkfSB$-linear
(Lemma~\ref{pkf:extend}) and kills $\pkfSB$ (Lemma~\ref{pkf:values}); since $\mathfrak S[E/\ell]
\subseteq\pkfSB$ (Lemma~\ref{pkf:SB-basic}(i)), $D(\ell^my)=\ell^nD(w)\in\ell^n\pkfAo$ for all $n$, so
$D(\ell^my)=0$ and $\ell^my=G_0(\ell^my)\in\pkfSB$; then $y\in\pkfSB$ by Lemma~\ref{pkf:SB-basic}(i).
Conversely let $y\in\pkfSB$, $\ell^my=\beta\in\pkfAo$, and $\beta=P(E/\ell)+\ell^nw$
(Lemma~\ref{pkf:approx0}(i)). Applying $G_0$: $\ell^my=G_0(P(E/\ell))+\ell^nG_0(w)$, and for a monomial
$G_0(\alpha(E/\ell)^k)=(E/\ell)^k\kappa_0(\alpha)\in\mathfrak S[E/\ell]$ ($E/\ell\in\pkfSB$). The last
sentence is the definition of $\mathcal K_0$ in Definition~\ref{pkn:def-N}.
\end{proof}

\begin{lemma}[Invariants]\label{pkf:invariants}
If $b\in\pkfBm$ is fixed by $j(\Gamma')$, then $b$ is the image of an element of $\pkfSB$.
\end{lemma}

\begin{proof}
By the chart, $j(\Gamma')$ is the group of elements of $G_{\Q_\ell}$ fixing $\Lambda$ and $\theta(z)=\pi$,
i.e.\ $G_{K'}$ with $K'=K_0(\pi)$, $K_0=\Lambda[1/\ell]$ (by Weierstrass preparation $E(X)$ is a
unit of $\Lambda[[X]]$ times an Eisenstein polynomial $P\in\Lambda[X]$ of degree $e$, and $\pi$ is a
root of $P$, so $[K_0(\pi):K_0]=e$); so $K'_0=\Lambda[1/\ell]$. By (P2), $b\in\Lambda[1/\ell]$, i.e.\
$b=\ell^{-m}\lambda$ with $\lambda\in\Lambda\subseteq\mathfrak S$, and $\ell^{-m}\in\pkfSB$.
\end{proof}

\begin{lemma}[Coordinate approximation]\label{pkf:coord-approx}
For $y\in\pkfBp$ and $m\ge0$ there is a finite $G\subseteq I$ with $y-\sum_{i\in G}b_iG_i(y)\in
\ell^m\pkfAo$.
\end{lemma}

\begin{proof}
Write $\ell^ky=\alpha+\ell^{k+m}w$ (Lemma~\ref{pkf:approx0}(ii)). Since $z^e=(E-\ell)A^{-1}\in(\ell,E)$,
$(\ell,z)^{e(k+m)}\subseteq(\ell,E)^{k+m}\subseteq\ell^{k+m}\pkfAm$. By (A5) choose $G$ with
$\alpha-\sum_{i\in G}\kappa_i(\alpha)b_i\in(\ell,z)^{e(k+m)}$. The map $T=1-\sum_{i\in G}b_iG_i$ is
additive and preserves $\pkfAo$, and $T(\alpha)=\alpha-\sum_{G}b_i\kappa_i(\alpha)$. So $\ell^kT(y)=
T(\alpha)+\ell^{k+m}T(w)\in\ell^{k+m}\pkfAo$, and $T(y)\in\ell^m\pkfAo$ ($\pkfAm$ is $\ell$-torsion free).
\end{proof}

\subsubsection{Stage (a): a period matrix}\label{pkf:ss-pm}

\begin{definition}[Period matrix]\label{pkf:def-pm}
Let $h\ge0$ and $v_1\in\pkpA^\times$. A matrix $U\in M_d(\pkfBp)$ is a \emph{period matrix of height $h$
with Frobenius factor $v_1$} if:
\begin{enumerate}
\item[(PM1)] there is $U'\in M_d(\pkfBp)$ with $UU'=(E\tau)^h\cdot1$;
\item[(PM2)] $\varphi(U)=v_1UF$ with $F\in M_d(\pkfSB)$ having a left inverse $F'\in M_d(\pkfBp)$;
\item[(PM3)] for every $g\in G_{\Q_\ell}$ fixing $\Lambda$ pointwise and $\theta(z)$ there is $R_g\in
M_d(\pkpA)$ with $R_g\,g(U)=U$;
\item[(PM4)] for every $\mathfrak S$-linear $f\colon\pkpA\to\pkpA$, $n\ge0$ and $y\in(\pkfBp)^d$: if
$Uy\in E^n(\pkfBp)^d$ then $UG_f(y)\in E^n(\pkfBp)^d$ ($G_f$ applied entrywise).
\end{enumerate}
\end{definition}

\begin{lemma}[Hodge frames]\label{pkf:hodge}
Let $N\in GL_d(\pknB)$ have invariant columns. Then there are $Y\in SL_d(K')$, $W\in GL_d(\pknBp)$ and
integers $k_1,\dots,k_d$, each a Hodge--Tate jump of $L[1/\ell]$, with $NY=W\operatorname{diag}
(t^{k_1},\dots,t^{k_d})$.
\end{lemma}

\begin{proof}
The columns $\delta_j=\sum_kN_{kj}\otimes e_k$ are $G_{K'}$-invariant vectors of $\pknB\otimes L$, linearly
independent over $\pknB$ ($N$ invertible), hence over $K'$; so they lie in $D_{\mathrm{dR}}$, which has
$K'$-dimension $\le d$ (by (P3), or since $L[1/\ell]$ is de Rham by (P2)); they form a $K'$-basis of
$D_{\mathrm{dR}}$. Choose a $K'$-basis $\delta'_1,\dots,\delta'_d$ adapted to the filtration, with
$\delta'_j\in\mathrm{Fil}^{k_j}D_{\mathrm{dR}}$ and $\mathrm{Fil}^kD_{\mathrm{dR}}$ spanned by the
$\delta'_j$ with $k_j\ge k$; the $k_j$ are then the jumps with multiplicity. Let $Y\in GL_d(K')$ be the
change of basis ($\delta'=\delta Y$); rescaling $\delta'_1$ by $\det Y^{-1}\in K'^\times$ we may take
$\det Y=1$. By (P3), $\pknBp\otimes L=\sum_k\mathrm{Fil}^{-k}\pknB\otimes\mathrm{Fil}^kD_{\mathrm{dR}}
=\bigoplus_j t^{-k_j}\pknBp\,\delta'_j$ (the second equality because $\mathrm{Fil}^{-k}\pknB=t^{-k}\pknBp$
and the basis is adapted). Hence the vectors $t^{-k_j}\delta'_j$ form a $\pknBp$-basis of $\pknBp\otimes
L$: their coordinate matrix $W=NY\operatorname{diag}(t^{-k_j})$ lies in $GL_d(\pknBp)$.
\end{proof}

\begin{lemma}[Filtration of the entries]\label{pkf:fil}
In the situation of Lemma~\ref{pkf:hodge}, if $-a'\le k_j\le b'$ for all $j$, then the entries of $N$
lie in $\mathrm{Fil}^{-a'}$ and those of $N^{-1}$ in $\mathrm{Fil}^{-b'}$.
\end{lemma}

\begin{proof}
$N=W\operatorname{diag}(t^{k_j})Y^{-1}$ with $W$ over $\pknBp$ and $Y^{-1}$ over $K'\subset\pknBp$; the
$(r,s)$ entry is $\sum_jW_{rj}t^{k_j}(Y^{-1})_{js}\in\mathrm{Fil}^{\min k_j}$. Likewise
$N^{-1}=Y\operatorname{diag}(t^{-k_j})W^{-1}$.
\end{proof}

\begin{theorem}[Existence of a period matrix]\label{pkf:pm}
Suppose $L[1/\ell]$ is crystalline on $\Gamma'$ with Hodge--Tate jumps in $[-a,b]$ ($a,b\ge0$). Then
there are $M\in GL_d(\pkfBm)$ with invariant columns and $U\in M_d(\pkfBp)$ whose image in $\pkfBm$ is
$(E\tau)^aM$, such that $U$ is a period matrix of height $a+b$ with Frobenius factor $v^a$.
\end{theorem}

\begin{proof}
\emph{Step 1: $M$.} By (P2) applied to $V=L[1/\ell]$, $(\pkfBm\otimes L)^{G_{K'}}$ has a $K'_0$-basis
which is a $\pkfBm$-basis of $\pkfBm\otimes L$; let $M\in GL_d(\pkfBm)$ be the matrix of its coordinates
in the basis $1\otimes e_k$. Then $M$ has invariant columns, and so has $M_n=\varphi^n(M)$ for every
$n\ge0$.

\emph{Step 2: $U$ and $U'$.} Apply Lemmas~\ref{pkf:hodge} and \ref{pkf:fil} to the image of $M_n$ in
$GL_d(\pknB)$: the entries of $\varphi^n(M)$ lie in $\mathrm{Fil}^{-a}$, those of $\varphi^n(M^{-1})=
M_n^{-1}$ in $\mathrm{Fil}^{-b}$, for every $n\ge0$. By (P4), $\mu^aM$ and $\mu^bM^{-1}$ have entries in
$\pkfBp$; as $\mu=E\tau v$ with $v$ a unit, $(E\tau)^aM$ and $(E\tau)^bM^{-1}$ are images of matrices $U$,
$U'$ over $\pkfBp$, unique by injectivity of $\pkfBp\to\pkfBm$. Then $UU'=(E\tau)^{a+b}$ (compute in
$\pkfBm$): (PM1).

\emph{Step 3: (PM2).} Let $X=M^{-1}\varphi(M)\in M_d(\pkfBm)$. For $g\in\Gamma'$, $g(M)=\rho(g)^{-1}M$
and $g(\varphi(M))=\rho(g)^{-1}\varphi(M)$, so $g(X)=M^{-1}\rho(g)\rho(g)^{-1}\varphi(M)=X$. By
Lemma~\ref{pkf:invariants}, $X$ is the image of a matrix over $\pkfSB$, again written $X$; likewise
$X'=\varphi(M)^{-1}M$, and $X'X=1$ (injectivity). Put $F=\varphi(E)^aX$ and $F'=\varphi(E)^{-a}X'$; $F$
is over $\pkfSB$ ($\varphi(E)\in\mathfrak S$), $F'$ over $\pkfSB\subseteq\pkfBp$
(Lemma~\ref{pkf:SB-basic}(i)), $F'F=1$. In $\pkfBm$: $\varphi(U)=\varphi(E)^a\varphi(\tau)^a\varphi(M)=
\varphi(E)^a\mu^aMX=v^a(E\tau)^aM\cdot\varphi(E)^aX=v^aUF$, using $\varphi(\tau)=\mu=E\tau v$. By
injectivity, $\varphi(U)=v^aUF$ in $\pkfBp$.

\emph{Step 4: (PM3).} Let $g$ fix $\Lambda$ and $\theta(z)$. By the chart, $g=j(\gamma)$ with
$\gamma\in\Gamma'$. By Lemma~\ref{pkf:approx0}(iv), $g(E)g(\tau)r''=E\tau$ with $r''\in\pkpA^\times$. Put
$R_g=r''^a\rho(\gamma)\in M_d(\pkpA)$. In $\pkfBm$: $R_g\,g(U)=r''^a(gE\,g\tau)^a\rho(\gamma)g(M)=
(E\tau)^aM=U$.

\emph{Step 5: (PM4).} Let $Y,W,k$ be as in Lemma~\ref{pkf:hodge} for $N=M$ (image in $\pknB$), so that
over $\pknB$, $U=(E\tau)^aW\operatorname{diag}(t^{k_j})Y^{-1}$, with $Y^{-1}=\operatorname{adj}Y$ over $K'$.
Suppose $Uy=E^nw$ with $w$ over $\pkfBp$. By (A8)(c), applied to the rows $\lambda=(Y^{-1})_{r\cdot}$,
there are $s_r\in\pknBp$ with $\sum_i(Y^{-1})_{ri}G_f(y_i)=s_r\sum_i(Y^{-1})_{ri}y_i$, i.e.\
$Y^{-1}G_f(y)=\operatorname{diag}(s)Y^{-1}y$. Hence, diagonal matrices commuting,
\[UG_f(y)=(E\tau)^aW\operatorname{diag}(t^{k_j})\operatorname{diag}(s)Y^{-1}y
=W\operatorname{diag}(s)W^{-1}\cdot(E\tau)^aW\operatorname{diag}(t^{k_j})Y^{-1}y
=W\operatorname{diag}(s)W^{-1}\,E^nw .\]
$W\operatorname{diag}(s)W^{-1}$ is over $\pknBp$, so the entries of $UG_f(y)$ lie in $E^n\pknBp=
\mathrm{Fil}^n$ (P6), and by Lemma~\ref{pkf:approx0}(v) they are divisible by $E^n$ in $\pkfBp$.
\end{proof}

\subsubsection{Stage (b): from a period matrix to a frame}\label{pkf:ss-frame}

Throughout this subsection $U$ is a period matrix of height $h$ with Frobenius factor $v_1$, with $U'$,
$F$, $F'$, $R_g$ as in Definition~\ref{pkf:def-pm}. A \emph{solution} is a pair $(x,y)$ with $x\in
\pkpA^d$, $y\in(\pkfBp)^d$ and $Uy=E^hx$. Vectors over $\pkpA$ are mapped to $\pkfBp$ entrywise.

\begin{lemma}\label{pkf:inj}
$\det U\ne0$; $U$ is injective on $(\pkfBp)^d$; $U'U=(E\tau)^h$; and for every $x'\in\pkpA^d$,
$(\tau^hx',U'x')$ is a solution.
\end{lemma}

\begin{proof}
$\det U\det U'=(E\tau)^{hd}\ne0$ in the domain $\pkfBp$ ($E,\tau\ne0$ in $\pkpA\subseteq\pkfBp$). Over the
fraction field $U'=(E\tau)^hU^{-1}$, so $U'U=(E\tau)^h$. Finally $UU'x'=E^h\tau^hx'$.
\end{proof}

\begin{lemma}[Denominators]\label{pkf:denom}
Let $T\in M_d(\pkfBp)$ and $\delta\mid\mu$ in $\pkpA$. There is $k$ such that: if $Ty=\delta^hw$ with $y$ over
$\pkfAo$ and $w$ over $\pkfBp$, then $\ell^kw$ is over $\pkfAo$.
\end{lemma}

\begin{proof}
Take $k_0$ with $\ell^{k_0}T$ over $\pkfAo$. Then $\delta^h\ell^{k_0}w=\ell^{k_0}Ty$ is over $\pkfAo$, and
(A2) gives $\ell^{2h+k_0}w$ over $\pkfAo$; $k=2h+k_0$.
\end{proof}

\begin{lemma}[Descent along Frobenius]\label{pkf:descent}
Let $(x,y)$ be a solution, $f$ an $\mathfrak S$-linear map $\pkpA\to\pkpA$ and $w$ with
$UG_f(y)=E^hw$. Then there are a solution $(x',y')$, an $\mathfrak S$-linear $f'$ and $w'$ with
$UG_{f'}(y')=E^hw'$ and $\varphi(w')=v_1E^hw$.
\end{lemma}

\begin{proof}
Put $x'=\varphi^{-1}(v_1E^hx)\in\pkpA^d$. \emph{Claim: if $r\in\pkpA$ and $Ua=rx'$ with $a$ over
$\pkfBp$, then $F\varphi(a)=\varphi(r)y$.} Indeed $\varphi(U)\varphi(a)=\varphi(r)v_1E^hx=\varphi(r)v_1Uy$,
and $\varphi(U)=v_1UF$; cancel the unit $v_1$ and use injectivity of $U$.
Apply the claim to $a=U'x'$, $r=(E\tau)^h$: $F\varphi(U'x')=\varphi(E\tau)^hy$, so $\varphi(U'x')=
\varphi(\tau)^h\varphi(E)^hF'y$. By (A8)(b) and induction on $h$ (cancelling $\varphi(\tau)=\mu\ne0$ in the
domain $\pkfBp$ at each step), $U'x'=\tau^hy'$ with $y'$ over $\pkfBp$; then $\tau^hUy'=(E\tau)^hx'$ and
$Uy'=E^hx'$: $(x',y')$ is a solution.
Let $f'=\varphi^{-1}\circ f\circ\varphi$; it is $\mathfrak S$-linear because $\varphi(\mathfrak S)\subseteq
\mathfrak S$: $f'(s\alpha)=\varphi^{-1}(\varphi(s)f(\varphi\alpha))=sf'(\alpha)$. By (PM4) there is $w'$ with
$UG_{f'}(y')=E^hw'$. The map $\varphi\circ G_{f'}-G_f\circ\varphi$ is additive, preserves $\pkfAo$ and
kills $\pkpA$, so it is $0$ (Lemma~\ref{pkf:density}). By the claim with $a=y'$, $r=E^h$:
$F\varphi(y')=\varphi(E)^hy$. Since $G_f$ is $\pkfSB$-linear (Lemma~\ref{pkf:extend}), $F$ is over
$\pkfSB$ and $\varphi(E)\in\mathfrak S$,
\[\varphi(E)^h\varphi(w')=\varphi(UG_{f'}(y'))=v_1UF\,G_f(\varphi(y'))=v_1U\,G_f(F\varphi(y'))
=v_1\varphi(E)^hUG_f(y)=v_1\varphi(E)^hE^hw .\]
Cancel $\varphi(E)^h\ne0$.
\end{proof}

\begin{lemma}[Uniform integrality]\label{pkf:integral}
There is $N\ge0$ such that whenever $(x,y)$ is a solution, $f$ is $\mathfrak S$-linear and
$UG_f(y)=E^hw$, the vector $\ell^Nw$ is (the image of a vector) over $\pkpA$.
\end{lemma}

\begin{proof}
Let $k_1$ be as in Lemma~\ref{pkf:denom} for $T=U'$, $\delta=E\tau$, and $k_2$ for $T=U$, $\delta=E$. Given the data,
$U'(Uy)=(E\tau)^hy$ with $Uy=E^hx$ over $\pkfAo$, so $\ell^{k_1}y$ is over $\pkfAo$, hence so is
$G_f(\ell^{k_1}y)$, and $U(G_f(\ell^{k_1}y))=E^h\ell^{k_1}w$ gives $\ell^{m}w$ over $\pkfAo$, $m=k_1+k_2$.
Let $V\subseteq\pkfAo$ be the set of all entries of all such $\ell^mw$. Apply
Lemma~\ref{pkf:descent}: for $\alpha=\ell^mw_i\in V$ there is $\alpha'=\ell^mw'_i\in V$ with
$\varphi(\alpha')=v_1E^h\alpha$, i.e.\ $E^h\alpha=v_1^{-1}\varphi(\alpha')$. By (A3) with $u=v_1^{-1}$,
$\ell^{h+1}V\subseteq\pkpA$; $N=h+1+m$.
\end{proof}

\begin{lemma}[The spanning matrix]\label{pkf:span}
Let $P=\{x\in\pkpA^d:\ \exists\,y\in\pkfSB^d,\ Uy=E^hx\}$. There are $Q\in M_d(\pkpA)$ and $Z\in
M_d(\pkfSB)$ with $UZ=E^hQ$ and $P=Q\,\mathfrak S^d$.
\end{lemma}

\begin{proof}
$P$ is an $\mathfrak S$-submodule. For $x\in P$ the $y$ is unique (Lemma~\ref{pkf:inj}); call it $y_x$.
For $\alpha_j\in\pkpA$ and $x_j\in P$: $\sum\alpha_jx_j=0$ iff $\sum\alpha_jy_{x_j}=0$, since
$U\sum\alpha_jy_{x_j}=E^h\sum\alpha_jx_j$, $U$ is injective and $E^h\ne0$.
\emph{No $d+1$ elements of $P$ are $\mathfrak S$-independent.} Let $x_1,\dots,x_r\in P$ be
$\mathfrak S$-independent and $\sum_j\alpha_jx_j=0$ with $\alpha_j\in\pkpA$. Then $\sum_j\alpha_jy_{x_j}=0$;
apply $G_i$ entrywise: as the entries of $y_{x_j}$ lie in $\pkfSB$ and $G_i$ is $\pkfSB$-linear,
$\sum_j\kappa_i(\alpha_j)y_{x_j}=0$, so $\sum_j\kappa_i(\alpha_j)x_j=0$ and $\kappa_i(\alpha_j)=0$ for all
$i$; by (A5), $\alpha_j=0$. So the $x_j$ are independent over the domain $\pkpA$, whence $r\le d$.
\emph{Saturation.} If $a,b\in P$ and $za=\ell b$: $\ell$ is prime in $\pkpA$ ($\pkpA/\ell=O^\flat$ is a
domain) and $\ell\nmid z$ ($\bar z=\varpi^\flat\ne0$), so $a=\ell a'$ with $a'\in\pkpA^d$, and $a'\in P$
with $y_{a'}=\ell^{-1}y_a$ ($\ell^{-1}\in\pkfSB$). By (A6), $P$ is free with a basis
$q_1,\dots,q_r$, $r\le d$. Let $Q$ have columns $q_1,\dots,q_r,0,\dots,0$ and $Z$ the columns
$y_{q_1},\dots,y_{q_r},0,\dots,0$.
\end{proof}

\begin{lemma}[Divisibility]\label{pkf:dvd}
Let $Q$ be as in Lemma~\ref{pkf:span}, $(x,y)$ a solution, $\delta\in\pkpA$ with $\ell\nmid\delta$, and
$l,g\in\pkpA^d$ with $l^{\mathsf T}Q=\delta g^{\mathsf T}$. Then $\delta\mid l\cdot x$.
\end{lemma}

\begin{proof}
Let $N$ be as in Lemma~\ref{pkf:integral} and $k$ as in Lemma~\ref{pkf:denom} for $T=U$, $\delta=E$. For
$i\in I$, (PM4) gives $w_i$ with $UG_i(y)=E^hw_i$, and $x'_i=\ell^Nw_i$ is over $\pkpA$. Since
$\ell^NG_i(y)$ is over $\pkfSB$ (Lemma~\ref{pkf:values}) and $U\ell^NG_i(y)=E^hx'_i$, $x'_i\in P$, so
$x'_i=Qs_i$ with $s_i\in\mathfrak S^d$. Fix $m$. By Lemma~\ref{pkf:coord-approx} (for each entry, then
the union of the sets) there are a finite $G$ and $r$ over $\pkfAo$ with $y-\sum_{i\in G}b_iG_i(y)=
\ell^{m+k}r$. Put $t=\ell^Nx-Q\sum_{i\in G}b_is_i\in\pkpA^d$. Then
\[U(\ell^Nr)=\ell^{-m-k}\Bigl(\ell^NUy-\sum_{i\in G}b_i\,\ell^NUG_i(y)\Bigr)=E^h\,\ell^{-m-k}t,\]
so $\ell^{-m}t$ is over $\pkfAo$ (Lemma~\ref{pkf:denom}), i.e.\ $t\in(\pkpA\cap\ell^m\pkfAm)^d=
((\ell,E)^m)^d$ by (A1). Now $l\cdot\ell^Nx-\delta\,g\cdot\sum_Gb_is_i=l\cdot t+l^{\mathsf T}Q\sum_G
b_is_i-\delta g^{\mathsf T}\sum_Gb_is_i=l\cdot t\in(\ell,E)^m$. By (A7), $\delta\mid\ell^N(l\cdot x)$; as
$\ell$ is prime and $\ell\nmid\delta$, $\delta\mid l\cdot x$.
\end{proof}

\begin{lemma}\label{pkf:detQ}
$\det Q\ne0$.
\end{lemma}

\begin{proof}
Otherwise there is $0\ne l\in\pkpA^d$ with $l^{\mathsf T}Q=0$ ($\pkpA$ a domain). For $x'\in\pkpA^d$,
$(\tau^hx',U'x')$ is a solution (Lemma~\ref{pkf:inj}); Lemma~\ref{pkf:dvd} with $\delta=z^n$
($\ell\nmid z^n$), $g=0$ gives $z^n\mid\tau^h\,l\cdot x'$ for all $n$. \emph{$\bigcap_nz^n\pkpA=0$}: if
$c=z^nc_n$ for all $n$, then $\bar c\in\bigcap_n(\varpi^\flat)^nO^\flat=0$, so $c=\ell c'$; then $\ell\mid
z^nc_n$ forces $\ell\mid c_n$ and $c'=z^n(c_n/\ell)$ for all $n$; iterating, $c\in\bigcap_k\ell^k\pkpA=0$.
Hence $\tau^h\,l\cdot x'=0$, so $l\cdot x'=0$ for all $x'$ and $l=0$.
\end{proof}

\begin{lemma}[Frobenius, saturation, Galois]\label{pkf:FSG}
(i) There is $B\in M_d(\mathfrak S)$ with $\varphi(Q)=v_1QB$; and $\ell\nmid\det B$.
(ii) If $s\in\mathfrak S^d$ and $\ell$ divides every entry of $Qs$, then $s\in\ell\mathfrak S^d$.
(iii) If $g\in G_{\Q_\ell}$ fixes $\mathfrak S$ pointwise, there is $\rho_g\in M_d(\pkpA)$ with
$Q=\rho_g\,g(Q)$. More generally, for $g$ fixing $\Lambda$ and $\theta(z)$, $Ug(Z)=E^h\rho_gg(Q)$ for some
$\rho_g\in M_d(\pkpA)$.
\end{lemma}

\begin{proof}
(i) Apply $\varphi$ to $UZ=E^hQ$ and use $\varphi(U)=v_1UF$: $U\cdot F\varphi(Z)=v_1^{-1}\varphi(E)^h
\varphi(Q)$, so $U\bigl(E^h\varphi(E)^{-h}F\varphi(Z)\bigr)=E^h\,v_1^{-1}\varphi(Q)$. The matrix in brackets
is over $\pkfSB$ (Lemmas~\ref{pkf:SB-basic}(i), \ref{pkf:values}), so the columns of $v_1^{-1}\varphi(Q)$
lie in $P=Q\mathfrak S^d$: $v_1^{-1}\varphi(Q)=QB$. Then $\varphi(\det Q)=v_1^d\det Q\det B$. If
$\ell\mid\det B$, induction gives $\ell^n\mid\det Q$ for all $n$ ($\ell^n\mid\det Q\Rightarrow\ell^{n+1}\mid
\varphi(\det Q)\Rightarrow\ell^{n+1}\mid\det Q$, $\varphi$ an automorphism fixing $\ell$), so $\det Q=0$,
contradicting Lemma~\ref{pkf:detQ}.
(ii) Write $Qs=\ell w$. Then $U(Z\ell^{-1}s)=E^hw$ with $Z\ell^{-1}s$ over $\pkfSB$, so $w\in P$, $w=Qs'$
with $s'\in\mathfrak S^d$, and $Q(s-\ell s')=0$; $\det Q\ne0$ gives $s=\ell s'$.
(iii) Let $g$ fix $\Lambda$ and $\theta(z)$, and $gE=Er$ (Lemma~\ref{pkf:approx0}(iv)). Applying $g$ to
$UZ=E^hQ$ and multiplying by $R_g$ (PM3): $Ug(Z)=R_g\,g(E)^hg(Q)=E^h\rho_gg(Q)$ with
$\rho_g=r^hR_g$. If $g$ fixes $\mathfrak S$ pointwise, it fixes $\Lambda$ and $\theta(z)=g\theta(z)$
(as $\theta(gz)=\theta(z)$), and $g(Z)=Z$ by Lemma~\ref{pkf:SB-basic}(ii); so $E^hQ=UZ=E^h\rho_gg(Q)$, and
$Q=\rho_gg(Q)$ ($E^h\ne0$, $\pkpA\to\pkfBp$ injective).
\end{proof}

\begin{lemma}[Rows with nonzero determinant]\label{pkf:rows}
Let $\mathcal R\colon S\to D$ be a ring map to a domain and $\mathcal X\subseteq S^d$ such that for
$t\in S^d$: if $\mathcal R(t\cdot x)=0$ for all $x\in\mathcal X$, then $\mathcal R(t_i)=0$ for all $i$.
Then there is $N\in M_d(S)$ with all rows in $\mathcal X$ and $\mathcal R(\det N)\ne0$.
\end{lemma}

\begin{proof}
Among $N$ with $\mathcal R(\det N)\ne0$ (e.g.\ $N=1$) choose one with the largest number of rows in
$\mathcal X$. Suppose row $k$ is not in $\mathcal X$. For $x\in\mathcal X$, the matrix $N[k\leftarrow x]$
(row $k$ replaced by $x$) has determinant $\sum_ix_i\operatorname{adj}(N)_{ik}$ and more rows in
$\mathcal X$, so by maximality $\mathcal R(x\cdot\operatorname{adj}(N)_{\cdot k})=0$. By hypothesis
$\mathcal R(\operatorname{adj}(N)_{ik})=0$ for all $i$, and $\mathcal R(\det N)=\sum_i\mathcal R(N_{ki})
\mathcal R(\operatorname{adj}(N)_{ik})=0$, a contradiction.
\end{proof}

\begin{lemma}[Minimal support]\label{pkf:artin}
Let $\mathcal V$ be a valuation ring (a domain) and $T\in M_d(\mathcal V)$ with $\det T=0$. There are
$a\in\mathcal V^d$ and $\nu$ with $a_\nu=1$, $Ta=0$, and $g(a)=a$ for every ring endomorphism $g$ of
$\mathcal V$ with $T\,g(a)=0$.
\end{lemma}

\begin{proof}
Let $w\ne0$ be a kernel vector with minimal support. The ideals $w_i\mathcal V$ are totally ordered, so
some $w_\nu\ne0$ divides all $w_i$; put $a=w/w_\nu\in\mathcal V^d$, with $Ta=0$, $a_\nu=1$ and the same
support as $w$. If $Tg(a)=0$, then $g(a)-a$ is a kernel vector with support in that of $a$ and
$(g(a)-a)_\nu=g(1)-1=0$; by minimality $g(a)-a=0$.
\end{proof}

For $m\ge0$ let $c_m\colon O^\flat\to O_C/\ell$, $c_m(b)=(b^{1/\ell^m})^\sharp\bmod\ell$; it is a ring
map (the projection of $O^\flat=\varprojlim_{x\mapsto x^\ell}O_C/\ell$ to the $m$-th factor), with
$c_{m+n}(b^{\ell^n})=c_m(b)$, $c_m(b)=0$ iff $v^\flat(b)\ge\ell^m$, and $c_m(\bar x)=\theta(\varphi^{-m}x)
\bmod\ell$ for $x\in\pkpA$ (Witt coordinates, \S\ref{pkn:ss-dist}).

\begin{lemma}[Approximation of fixed vectors]\label{pkf:fixed}
Let $a\in(O^\flat)^d$ be fixed by the group $H$ of (P5). For every $M\ge0$ there are $n\ge0$ and
$s\in\mathfrak S^d$ with $c_M(a_i)=c_{M+n}(\bar s_i)$ for all $i$.
\end{lemma}

\begin{proof}
\emph{The group $H$.} Let $\pi_n=((\varpi^\flat)^{1/\ell^n})^\sharp\in O_C$, so $\pi_0=\pi=\theta(z)$ is a
uniformiser of $K'$ and $\pi_{n+1}^\ell=\pi_n$. An element $g\in G_{\Q_\ell}$ fixes $\mathfrak S=\Lambda[[z]]$
pointwise iff it fixes $\Lambda$ and $z=[\varpi^\flat]$, i.e.\ $\varpi^\flat$ (Teichm\"uller lifts are
injective and equivariant), i.e.\ all $(\varpi^\flat)^{1/\ell^n}$ (unique $\ell$-power roots in the perfect
ring $O^\flat$), i.e.\ all $\pi_n$ (an element $b\in O^\flat$ is determined by the $(b^{1/\ell^n})^\sharp$,
and $\sharp$ is equivariant); and $g$ fixes $\Lambda=W(\mathbb F_{\ell^f})$ iff it fixes the unramified
extension $K_f=\Lambda[1/\ell]$ of degree $f$. So $H=\operatorname{Gal}(\auQlb/K_\infty)$ with
$K_\infty=\bigcup_nK_n$, $K_n=K_f(\pi_n)$.

\emph{The fields $K_n$.} $K'$ has residue degree $f$ and contains $K_f$ (as $j(\Gamma')$ fixes $\Lambda$), and
$\pi$ is a root of the Eisenstein polynomial $P\in\Lambda[X]$ of degree $e$ that differs from $E(X)$ by a
unit of $\Lambda[[X]]$ (Weierstrass preparation; Definition~\ref{pkp:def-chart}). $P(X^{\ell^n})$ is
Eisenstein over $\Lambda$, so it is irreducible, $K_n/K_f$ is totally ramified of degree $e\ell^n$, $\pi_n$ is
a uniformiser and $O_{K_n}=\Lambda[\pi_n]$ (Serre, \emph{Corps locaux}, Ch.~I, \S6, Proposition~17 and its Corollary; p.~19 of the English edition \cite{SerreLF}: for an Eisenstein polynomial $f$ over a discrete valuation ring $A$, $A[X]/(f)$ is a discrete valuation ring with uniformiser the image of $X$ and the same residue field, $f$ is irreducible, and $A[X]/(f)$ is the integral closure of $A$). Hence every element
of $O_{K_n}$ is congruent modulo $\ell$ to a finite sum $\sum_k[\lambda_k]\pi_n^k$ with
$\lambda_k\in\mathbb F_{\ell^f}$.

\emph{Approximation.} Let $a\in O^\flat$ be fixed by $H$ and $M\ge0$. Then $a^{1/\ell^M}$ is fixed by $H$,
so $x=(a^{1/\ell^M})^\sharp\in O_C$ is fixed by $H$, and by (P5) $x\in O_{\hat K_\infty}$. Choose
$y\in O_{K_\infty}$ with $x-y\in\ell O_C$, say $y\in O_{K_N}$ with $N\ge M$, and write
$y\equiv\sum_k[\lambda_k]\pi_N^k\pmod\ell$. Put $\beta=\sum_k\lambda_k^{\ell^N}(\varpi^\flat)^k\in
\mathbb F_{\ell^f}[\varpi^\flat]\subseteq O^\flat$. The ring map $c_N$ sends $\varpi^\flat$ to
$\pi_N\bmod\ell$ and $\lambda\in\mathbb F_{\ell^f}$ to $(\lambda^{1/\ell^N})^\sharp=[\lambda^{1/\ell^N}]
\bmod\ell$; so $c_N(\beta)=\sum_k[\lambda_k]\pi_N^k\bmod\ell=y\bmod\ell=x\bmod\ell=c_M(a)$. With $n=N-M$,
$c_{M+n}(\beta)=c_M(a)$. For finitely many $a_i$ take $N$ common to all (enlarging $N$ is allowed as
$O_{K_N}\subseteq O_{K_{N+1}}$) and $\beta_i$ as above. Finally $\mathbb F_{\ell^f}[\varpi^\flat]$ is the
image of $\Lambda[z]\subseteq\mathfrak S$ in $O^\flat$; take $s_i\in\Lambda[z]$ lifting $\beta_i$. Then
$c_M(a_i)=c_{M+n}(\bar s_i)$.
\end{proof}

\begin{proposition}[Key point]\label{pkf:key}
$\ell\nmid\det Q$.
\end{proposition}

\begin{proof}
Let $\mathcal R\colon\mathfrak S\to O^\flat$ be the reduction, $\mathcal X=\{x^{(i,j)}:i\in I,\ 1\le j\le d\}$
with $x^{(i,j)}_m=\kappa_i(Q_{jm})$; by $\mathfrak S$-linearity $t\cdot x^{(i,j)}=\kappa_i((Qt)_j)$ for
$t\in\mathfrak S^d$. \emph{Separation:} if $\mathcal R(\kappa_i((Qt)_j))=0$ for all $i,j$, let $c=(Qt)_j$;
for every $n$ and large $G$ (A5), $c\in\sum_G\kappa_i(c)b_i+(\ell,z)^n\subseteq\ell\pkpA+z^n\pkpA$, so
$\bar c\in\bigcap_n(\varpi^\flat)^nO^\flat=0$; thus $\ell$ divides every $(Qt)_j$, and $t\in\ell\mathfrak S^d$
by Lemma~\ref{pkf:FSG}(ii). By Lemma~\ref{pkf:rows} there are $N_\star\in M_d(\mathfrak S)$ and maps
$\sigma,\iota$ with $(N_\star t)_k=\kappa_{\sigma(k)}((Qt)_{\iota(k)})$ for all $t\in\mathfrak S^d$ and
$\mathcal R(\det N_\star)\ne0$. Put $\varkappa=\mathcal R(\det B)\mathcal R(\det N_\star)\ne0$ (Lemma~\ref{pkf:FSG}(i);
$O^\flat$ is a domain).

Suppose $\ell\mid\det Q$; then $\det\bar Q=0$. By Lemma~\ref{pkf:artin} ($O^\flat$ is a valuation ring) there
are $a$, $\nu$ with $a_\nu=1$, $\bar Qa=0$, and $g(a)=a$ whenever $\bar Qg(a)=0$. For $g\in H$:
$\bar Q=\bar\rho_gg(\bar Q)$ (Lemma~\ref{pkf:FSG}(iii)), so $\bar Qg(a)=\bar\rho_gg(\bar Qa)=0$ and
$g(a)=a$.

Let $C_0=1$, $C_{n+1}=B\varphi(C_n)\in M_d(\mathfrak S)$. \emph{Claim: $\bar Q\,\bar C_n\,a^{\ell^n}=0$
and $\varkappa\mathcal R(\det C_n)\mid \varkappa^{\ell^n}$} ($a^{\ell^n}$ entrywise). For $n=0$ this is clear. Applying the
Frobenius $b\mapsto b^\ell$ of $O^\flat$, which is the reduction of $\varphi$, to $\bar Q\bar C_na^{\ell^n}
=0$ and using $\overline{\varphi(Q)}=\bar v_1\bar Q\bar B$ gives $\bar Q\bar B\,\overline{\varphi(C_n)}
a^{\ell^{n+1}}=0$. Also $\mathcal R(\det C_{n+1})=\mathcal R(\det B)\mathcal R(\det C_n)^\ell$, and
$\varkappa\mathcal R(\det B)\mid \varkappa^2\mid \varkappa^\ell$, so $\varkappa\mathcal R(\det C_{n+1})\mid \varkappa^\ell\mathcal R(\det
C_n)^\ell\mid \varkappa^{\ell^{n+1}}$.

Fix $M\ge0$. By Lemma~\ref{pkf:fixed} choose $n$, $s$ with $c_M(a_i)=c_{M+n}(\bar s_i)$; let
$\psi=c_{M+n}\circ\mathcal R\colon\mathfrak S\to O_C/\ell$. Then $\psi(s_\nu)=c_M(1)=1$, and for every
$k$, $\psi((N_\star C_ns)_k)=0$: let $x=(QC_ns)_{\iota(k)}\in\pkpA$; as $c_{M+n}(\bar s_i)=c_M(a_i)=
c_{M+n}(a_i^{\ell^n})$, $c_{M+n}(\bar x)=c_{M+n}\bigl((\bar Q\bar C_na^{\ell^n})_{\iota(k)}\bigr)=0$, i.e.\
$\theta(\varphi^{-M-n}x)\in\ell O_C$. By (A8)(a), $\kappa_{\sigma(k)}(x)=\ell s'+p_{M+n}q$ with
$s',q\in\mathfrak S$; and $(N_\star C_ns)_k=\kappa_{\sigma(k)}(x)$ since $C_ns\in\mathfrak S^d$. Now $\psi(\ell)=0$ and
$\psi(p_{M+n})=\theta(\varphi^{-M-n}\varphi^{M+n}E)\bmod\ell=0$. Hence $\psi(N_\star C_n)\psi(s)=0$ with
$\psi(s_\nu)=1$; multiplying by $\operatorname{adj}\psi(N_\star C_n)$, $\psi(\det(N_\star C_n))=0$. So
$c_{M+n}(\varkappa\mathcal R(\det C_n))=c_{M+n}(\mathcal R\det B)\,\psi(\det N_\star\det C_n)=0$, and since
$\varkappa\mathcal R(\det C_n)\mid \varkappa^{\ell^n}$, $c_M(\varkappa)=c_{M+n}(\varkappa^{\ell^n})=0$. As $M$ is arbitrary,
$v^\flat(\varkappa)\ge\ell^M$ for all $M$ and $\varkappa=0$, a contradiction.
\end{proof}

\begin{theorem}[Frame theorem for a period matrix]\label{pkf:frame-pm}
There are $Q,Q'\in M_d(\pkpA)$, $Z\in M_d(\pkfSB)$, $Y\in M_d(\pkfBp)$ and $B\in M_d(\mathfrak S)$ with
$QQ'=\tau^h$, $UZ=E^hQ$, $YZ=E^h$, $\varphi(Q)=v_1QB$, and, for $g$ fixing $\Lambda$ and $\theta(z)$,
$g(Z)=ZG_g$ with $G_g\in M_d(\pkpA)$.
\end{theorem}

\begin{proof}
Take $Q,Z$ of Lemma~\ref{pkf:span} and $B$ of Lemma~\ref{pkf:FSG}(i).
\emph{Span: if $UY_1=E^hX$ with $X\in M_d(\pkpA)$, $Y_1\in M_d(\pkfBp)$, then $X=QL_X$ with $L_X\in
M_d(\pkpA)$.} The rows $l=\operatorname{adj}(Q)_{i\cdot}$ satisfy $l^{\mathsf T}Q=\det Q\cdot e_i^{\mathsf T}$;
Lemma~\ref{pkf:dvd} with $\delta=\det Q$ (Proposition~\ref{pkf:key}), applied to the columns of $(X,Y_1)$,
gives $\operatorname{adj}(Q)X=\det Q\cdot L_X$, and $\det Q\,(QL_X-X)=Q\operatorname{adj}(Q)X-\det Q\,X=0$.
\emph{$Q'$.} $UU'=E^h\tau^h$: so $\tau^h\cdot1=QQ'$.
\emph{$Y$.} For each entry $y=U_{rs}$, (A4) gives $\ell^{k_y}y=\alpha_y+\ell^{k_y}\mu^hw_y$. Let $n\ge$
all $k_y$, $A=(\ell^{n-k_y}\alpha_y)$ over $\pkpA$ and $W_\circ=(w_y)$ over $\pkfBp$. Then
$A=\ell^n(U-\mu^hW_\circ)=\ell^nU(1-v^{h}_{\phantom{1}}U'W_\circ)$, as $\mu^h=v^h(E\tau)^h$ and $(E\tau)^h=UU'$.
So $U\cdot E^h\ell^n(1-v^hU'W_\circ)=E^hA$ and $A=Q\gamma$ with $\gamma$ over $\pkpA$ by the span property.
Hence $U=\ell^{-n}Q\gamma+v^h(E\tau)^hW_\circ=Q\bigl(\ell^{-n}\gamma+v^hE^hQ'W_\circ\bigr)=:QY$ (using
$\tau^h=QQ'$). Then $Q(YZ-E^h)=UZ-E^hQ=0$, so $YZ=E^h$ ($\det Q\ne0$ in the domain $\pkfBp$).
\emph{Galois.} By Lemma~\ref{pkf:FSG}(iii), $Ug(Z)=E^h\rho_gg(Q)$, and the span property gives
$\rho_gg(Q)=QG_g$; so $Ug(Z)=E^hQG_g=UZG_g$, and $g(Z)=ZG_g$ by injectivity of $U$.
\end{proof}

\begin{theorem}[Existence of frames]\label{pkf:main}
Let $L$ be a lattice with jump sets and $a$ its jump bound. Then $L$ has a frame for $a$
(Definition~\ref{pkn:def-frame}). Consequently Proposition~\ref{pkn:frame} holds unconditionally (modulo
the inputs of \S\ref{pkf:ss-rings}): items (M) (third assertion) and (H) (first assertion, with $H=a$)
of Lemma~\ref{pkm:inputs}.
\end{theorem}

\begin{proof}
The Hodge--Tate jumps of $L[1/\ell]$ lie in $\bigcup_iU(i)\subseteq[-a,a]$ (P3). Theorem~\ref{pkf:pm}
with $b=a$ gives $M\in GL_d(\pkfBm)$ with invariant columns and a period matrix $U$ of height $h=2a$
whose image in $\pkfBm$ is $(E\tau)^aM$. Theorem~\ref{pkf:frame-pm} gives $Q,Q',Z,Y$ with $QQ'=
\tau^{2a}$, $UZ=E^{2a}Q$, $YZ=E^{2a}$, $Z$ over $\pkfSB$. Define $n_i$ by (Fr4) ($\tau$ is a unit of
$\pkpW$). Then (Fr1)--(Fr4) hold: (Fr1) is the invariance of the columns of $M$, (Fr2) is the
relation between $U$ and $M$, (Fr3) and (Fr4) are as stated. The rings $B^+$, $B_{\max}$ and the
maps of \S\ref{pkn:ss-frame} are those of \S\ref{pkf:ss-rings}, which have the properties listed
there by (P1).
\end{proof}

\begin{remark}\label{pkf:rem-extra}
The proof gives more than (Fr1)--(Fr4): a Frobenius matrix $B$ over $\mathfrak S$ with
$\varphi(Q)=v^aQB$ and $\ell\nmid\det B$, and the Galois relation $g(Z)=ZG_g$. These are not used in
\S\ref{pkn:ss-frame}.
\end{remark}



\providecommand{\auO}{\mathcal O}
\providecommand{\auQlb}{\overline{\Q}_\ell}
\providecommand{\pkpA}{A_{\mathrm{inf}}}
\providecommand{\pknBp}{B^+_{\mathrm{dR}}}
\providecommand{\pknB}{B_{\mathrm{dR}}}
\providecommand{\pkfAm}{A_{\max}}
\providecommand{\pkfAo}{A_{\max}^{\circ}}
\providecommand{\pkfBp}{B^+}
\providecommand{\pkfBm}{B_{\max}}
\providecommand{\pkfSB}{\mathfrak S_B}
\providecommand{\pkrI}{(\ell,E)}
\providecommand{\pkrm}{\mathfrak n}

\subsection{The period rings: proofs of the inputs (A1)--(A8)}\label{pkr:sec}

This section proves the eight auxiliary inputs (A1)--(A8) of \S\ref{pkf:ss-rings}, on which the
existence of frames (Theorem~\ref{pkf:main}) rests, and the statement (P4) listed there among the
published facts. We use the conventions and notation of \S\ref{pkf:ss-rings}:
$\pkpA=W(O^\flat)$, $R=\pkpA[X]/(\ell X-E)$, $\pkfAm$ its $\ell$-adic completion, $E/\ell$ the image
of $X$, $\pkfBp=\pkfAm[1/\ell]$, $\pkfBm=\pkfBp[1/\mu]$, $\pkfAo$ the image of $\pkfAm$ in $\pkfBp$,
the Kummer chart with $\mathfrak S=\Lambda[[z]]$, $E=z^eA+\ell$ ($A\in\mathfrak S^\times$),
$p_a=\varphi^a(E)$, $\mu=E\tau v$ with $v\in\pkpA^\times$ fixed. We write $\pkrI$ for the ideal of
$\pkpA$ generated by $\ell$ and $E$, and $\pkrm=(\ell,z)$ for the maximal ideal of $\mathfrak S$.

The following are proved in this section from (P1), (P6) and the standard facts (W1)--(W4)
recalled below: (A1) (Proposition~\ref{pkr:I1}), (A7) (Proposition~\ref{pkr:I7}), (A5)
(Proposition~\ref{pkr:I5}), (A6) (Proposition~\ref{pkr:I6}), (A3) (Proposition~\ref{pkr:I3}), (A2)
(Proposition~\ref{pkr:I2}), (A4) (Proposition~\ref{pkr:I4}), (A8) (Proposition~\ref{pkr:I8}), and
(P4) (Proposition~\ref{pkr:P4}). Of (P1) we use only: $\pkfAm\to\pkfBp\to\pkfBm\to\pknB$ and
$\pkfBp\to\pknBp$ are injective ring maps compatible with $\theta$, $\varphi$ and $G_{\Q_\ell}$, and
$\pkfBp$, $\pkfBm$ are domains (the $\ell$-torsion freeness, completeness and separation of $\pkfAm$
are proved in Lemma~\ref{pkr:Amax}).

\subsubsection{Standard facts and the structure of $\pkfAm$ modulo $\ell$}\label{pkr:ss-prelim}

We use the following standard facts, already used in \S\S\ref{pkn:sec}--\ref{pkf:sec}.
\begin{enumerate}
\item[(W1)] $\pkpA$ is $\ell$-torsion free, $\ell$-adically complete and separated, and
$\pkpA/\ell\pkpA=O^\flat$ via $x\mapsto x_0$ (the zeroth Witt coordinate). As $O^\flat$ is a domain,
$\ell$ is a prime element of $\pkpA$; $x\in\pkpA^\times$ iff $x_0\in O^{\flat\times}$
(\S\ref{pkn:ss-dist}).
\item[(W2)] $O^\flat$ is a valuation ring for $v^\flat$ (\S\ref{pkn:ss-dist}), complete for $v^\flat$
(Scholze \cite[Lemma~3.4(iii)]{pkfScholze12}: the tilt of a perfectoid field is a perfectoid field, hence complete, with ring of integers $\varprojlim_{x\mapsto x^\ell}O$ and valuation $x\mapsto|x^\sharp|$). Hence for $0\ne\varpi\in\mathfrak m^\flat$,
$O^\flat$ is $\varpi$-adically separated ($v^\flat(x)\ge kv^\flat(\varpi)$ for all $k$ forces $x=0$) and
complete (a sequence with $x_{k+1}-x_k\in\varpi^kO^\flat$ is $v^\flat$-Cauchy, and its limit $x$
satisfies $v^\flat(x-x_k)\ge kv^\flat(\varpi)$, i.e.\ $x-x_k\in\varpi^kO^\flat$).
\item[(W3)] $\theta\colon\pkpA\to O_C$ is surjective, and $\varphi$ is a ring automorphism of $\pkpA$
commuting with $G_{\Q_\ell}$.
\item[(W4)] $\mathfrak S=\Lambda[[z]]$ is a complete noetherian local ring with maximal ideal
$\pkrm=(\ell,z)$ and residue field $k=\mathbb F_{\ell^f}$; every ideal $\pkrm^n$ is closed for the
$\pkrm$-adic topology (it contains $\pkrm^m$ for $m\ge n$, so it is open); and
$\varphi(\mathfrak S)\subseteq\mathfrak S$ ($\varphi(z)=z^\ell$, $\varphi(\Lambda)=\Lambda$), so
$p_a\in\mathfrak S$.
\end{enumerate}
Two elementary consequences. Since $E=[E_0]+\ell E'$ (Corollary~\ref{pkn:E-gen},
Definition~\ref{pkn:def-dist}), we have $\pkrI=(\ell,t)$ with $t=[E_0]$, $E_0\in\mathfrak m^\flat$,
$v^\flat(E_0)=1$ (Lemma~\ref{pkn:dist}(i)). Since $z^e=A^{-1}(E-\ell)\in\pkrI$, we have
\begin{equation}\label{pkr:eq-zE}
\pkrm^{ek}\pkpA\subseteq\pkrI^k\qquad(k\ge0):
\end{equation}
a monomial $\ell^iz^j$ with $i+j=ek$ lies in $\ell^i\pkrI^{\lfloor j/e\rfloor}$ and
$i+\lfloor j/e\rfloor\ge\lfloor(i+j)/e\rfloor=k$.

\begin{lemma}[Structure of $\pkfAm$]\label{pkr:Amax}
\begin{enumerate}
\item[(i)] $R$ is $\ell$-torsion free; $\pkfAm=\varprojlim_nR/\ell^nR$ is $\ell$-torsion free; the kernel
of the projection $\pkfAm\to R/\ell^nR$ is $\ell^n\pkfAm$; $\pkfAm$ is $\ell$-adically complete and
separated, and $\pkfAm\to\pkfAo$ is bijective.
\item[(ii)] There is a surjective ring map $\rho\colon\pkfAm\to(\pkpA/\pkrI)[X]$ with kernel
$\ell\pkfAm$, sending $\alpha\in\pkpA$ to its class and $E/\ell$ to $X$.
\item[(iii)] Let $f\in\pkpA[X]$ be monic, $y\in\pkfAm$ and $m\ge0$. If $\ell^m$ divides $f(E/\ell)y$ in
$\pkfAm$, then $\ell^m$ divides $y$.
\end{enumerate}
\end{lemma}

\begin{proof}
(i) Let $\ell g\in(\ell X-E)\pkpA[X]$, say $\ell g=(\ell X-E)h$. Reducing modulo $\ell$ (by (W1),
$\pkpA[X]/\ell=O^\flat[X]$), $-E_0\bar h=0$ in the domain $O^\flat[X]$, and $E_0\ne0$; so $h=\ell h'$
and, $\pkpA[X]$ being $\ell$-torsion free, $g=(\ell X-E)h'$. Thus $R$ is $\ell$-torsion free. Hence for
$y\in R$, $\ell^ny\in\ell^{m+n}R$ implies $y\in\ell^mR$.

Let $x=(x_k)_k\in\pkfAm$ ($x_k\in R/\ell^kR$ compatible) with $x_k=0$ for $k\le n$. For $m\ge0$ lift
$x_{m+n}$ to $\tilde x\in R$; as $\tilde x\in\ell^nR$ write $\tilde x=\ell^n\tilde y$, and let $y_m$ be
the class of $\tilde y$ in $R/\ell^mR$; by the previous sentence $y_m$ does not depend on the
choices, and the $y_m$ are compatible. So $y=(y_m)\in\pkfAm$ and $x=\ell^ny$. Conversely $\ell^n\pkfAm$
maps to $0$ in $R/\ell^nR$. This proves the statement on kernels. If $\ell x=0$ in $\pkfAm$, then
for every $k$ a lift of $x_{k+1}$ times $\ell$ lies in $\ell^{k+1}R$, so the lift lies in $\ell^kR$ and
$x_k=0$: $\pkfAm$ is $\ell$-torsion free. Separation: if $x\in\ell^n\pkfAm$ for all $n$, then $x_n=0$ for
all $n$. Completeness: if $x^{(k+1)}-x^{(k)}\in\ell^k\pkfAm$, the components $x^{(k)}_n$ are constant
for $k\ge n$, and their common values form the limit. Finally $\pkfAm\to\pkfBp=\pkfAm[1/\ell]$ is
injective because $\pkfAm$ is $\ell$-torsion free.

(ii) $R/\ell R=\pkpA[X]/(\ell,\ell X-E)=\pkpA[X]/(\ell,E)=(\pkpA/\pkrI)[X]$; compose with the projection
$\pkfAm\to R/\ell R$, whose kernel is $\ell\pkfAm$ by (i).

(iii) $\pkpA/\pkrI\ne0$, since $\theta(\pkrI)\subseteq\ell O_C\ne O_C$. Case $m=1$: $\rho(f(E/\ell))=\bar f$
is monic in $(\pkpA/\pkrI)[X]$, hence not a zero divisor (the leading coefficient of $\bar f\,q$ is
that of $q$); so $\rho(f(E/\ell)y)=0$ gives $\rho(y)=0$, i.e.\ $\ell\mid y$. Induction: if $\ell^{m+1}
\mid f(E/\ell)y$, then $\ell\mid y$ by the case $m=1$, $y=\ell y'$, and $\ell^m\mid f(E/\ell)y'$ because
$\pkfAm$ is $\ell$-torsion free; so $\ell^m\mid y'$.
\end{proof}

We write $\theta$ also for the extension $\pkfAm\to O_C$ of Lemma~\ref{pkf:approx0}(v) (with
$\theta(E/\ell)=0$) and for its extension to $\pkfBp\to C$.

\begin{proposition}[Input (A1)]\label{pkr:I1}
For $n\ge0$, $\pkpA\cap\ell^n\pkfAm=\pkrI^n$ (intersection inside $\pkfAm$).
\end{proposition}

\begin{proof}
$\pkrI^n\subseteq\ell^n\pkfAm$ because $\ell$ and $E=\ell\cdot(E/\ell)$ lie in $\ell\pkfAm$. Conversely
we show by induction on $n$ that $a\in\pkpA$ and $\ell^n\mid a$ in $\pkfAm$ imply $a\in\pkrI^n$; the case
$n=0$ is empty. Let $\ell^{n+1}\mid a$ in $\pkfAm$. Applying $\theta$, $\theta(a)\in\ell^{n+1}O_C$; by
(W3) write $\theta(a)=\ell^{n+1}\theta(c)$ with $c\in\pkpA$. Then $\theta(a-\ell^{n+1}c)=0$, so
$a-\ell^{n+1}c=Ew$ with $w\in\pkpA$ (Corollary~\ref{pkn:E-gen}), i.e.\ $a=Ew+\ell\cdot\ell^nc$. In $\pkfAm$,
$Ew=\ell\,(E/\ell)w$, so $\ell^{n+1}\mid\ell\bigl((E/\ell)w+\ell^nc\bigr)$, and since $\pkfAm$ is $\ell$-torsion
free (Lemma~\ref{pkr:Amax}(i)), $\ell^n\mid(E/\ell)w$. By Lemma~\ref{pkr:Amax}(iii) with $f=X$,
$\ell^n\mid w$ in $\pkfAm$, so $w\in\pkrI^n$ by induction; also $\ell^nc\in\pkrI^n$. Hence
$a=Ew+\ell\cdot\ell^nc\in\pkrI\cdot\pkrI^n=\pkrI^{n+1}$.
\end{proof}

\subsubsection{Adic completeness of $\pkpA$ and closed principal ideals: input (A7)}\label{pkr:ss-complete}

For an element $s$ of a ring acting on a module $M$, we say $M$ is \emph{$s$-complete} if every sequence
$(x_k)_{k\ge0}$ in $M$ with $x_{k+1}-x_k\in s^kM$ has a limit $x\in M$ with $x-x_k\in s^kM$ for all $k$,
and \emph{$s$-separated} if $\bigcap_ks^kM=0$.

\begin{lemma}[Extensions]\label{pkr:ext}
Let $M'\subseteq M$ be modules over a commutative ring, $M''=M/M'$, and $s$ an element acting
injectively on $M''$. If $M'$ and $M''$ are $s$-complete and $s$-separated, so is $M$, and
$s^kM\cap M'=s^kM'$ for all $k$.
\end{lemma}

\begin{proof}
If $s^km\in M'$, then $s^k\bar m=0$ in $M''$, so $\bar m=0$ and $m\in M'$: this gives $s^kM\cap M'=
s^kM'$. Separation: if $x\in\bigcap_ks^kM$, its image lies in $\bigcap_ks^kM''=0$, so $x\in M'\cap s^kM=
s^kM'$ for all $k$ and $x=0$. Completeness: let $x_{k+1}-x_k\in s^kM$. The images $\bar x_k$ have a
limit $\bar x$ in $M''$; lift it to $x'\in M$. Then $x'-x_k=s^kr_k+m_k$ with $r_k\in M$, $m_k\in M'$.
Now $m_{k+1}-m_k=(x_k-x_{k+1})+s^kr_k-s^{k+1}r_{k+1}\in s^kM\cap M'=s^kM'$, so $(m_k)$ has a limit $m$ in
$M'$ with $m-m_k\in s^kM'$. Put $x=x'-m$; then $x-x_k=s^kr_k+(m_k-m)\in s^kM$.
\end{proof}

\begin{proposition}[$\pkrI$-adic completeness of $\pkpA$]\label{pkr:complete}
(i) $\bigcap_{k\ge0}\pkrI^k=0$ and $\bigcap_k\pkrm^k\pkpA=0$.
(ii) If $(a_k)_{k\ge0}$ in $\pkpA$ satisfies $a_{k+1}-a_k\in\pkrI^k$ for all $k$, there is $b\in\pkpA$
with $b-a_k\in\pkrI^{\lfloor k/2\rfloor}$ for all $k$.
\end{proposition}

\begin{proof}
Let $t=[E_0]$, so that $\pkrI=(\ell,t)$, and $M_n=\pkpA/\ell^n\pkpA$.

\emph{Step 1: $t$ acts injectively on $M_n$, and $M_n$ is $t$-complete and $t$-separated.} If
$tx\in\ell^n\pkpA$, then since $\ell$ is prime (W1) and $\ell\nmid t$ ($t_0=E_0\ne0$), $\ell^n\mid x$: so
$t$ is injective on $M_n$. Completeness and separation by induction on $n$. For $n=1$, $M_1=O^\flat$
with $t$ acting as $E_0\in\mathfrak m^\flat\smallsetminus\{0\}$, and (W2) applies. For $n\ge2$, the
submodule $M'=\ell M_n$ is isomorphic to $M_{n-1}$ via $x\mapsto\ell x$ (as $\pkpA$ is $\ell$-torsion
free), compatibly with $t$, and $M_n/M'=M_1$; Lemma~\ref{pkr:ext} applies.

\emph{Step 2.} For $0\le n\le k$, $\pkrI^k\subseteq\ell^n\pkpA+t^{k-n+1}\pkpA$: in a monomial $\ell^it^{k-i}$
either $i\ge n$ or $k-i\ge k-n+1$. Also $t^jx\in t^{j'}M_n$ with $j\le j'$ implies $x\in t^{j'-j}M_n$
(Step 1, injectivity).

\emph{Step 3: (i).} If $x\in\bigcap_k\pkrI^k$, then by Step 2 the class of $x$ in $M_n$ lies in
$\bigcap_kt^{k-n+1}M_n=0$; so $x\in\bigcap_n\ell^n\pkpA=0$ (W1). The second assertion follows by
\eqref{pkr:eq-zE}.

\emph{Step 4: (ii).} Fix $n\ge1$. For $k\ge n$, Step 2 gives $a_{k+1}-a_k\in\ell^n\pkpA+t^{k-n+1}\pkpA$,
so the classes $\bar a_{k}$ in $M_n$ satisfy $\bar a_{k+1}-\bar a_k\in t^{k-n+1}M_n$. By Step 1 (applied
to the sequence $j\mapsto\bar a_{j+n-1}$) there is $\beta_n\in M_n$ with $\beta_n-\bar a_k\in
t^{k-n+1}M_n$ for $k\ge n$. The image of $\beta_{n+1}$ in $M_n$ satisfies the same condition with
$t^{k-n}$ in place of $t^{k-n+1}$; two elements of $M_n$ both congruent to $\bar a_k$ modulo
$t^{k-n}M_n$ for all large $k$ agree, by $t$-separation. So the $\beta_n$ are compatible, and as
$\pkpA=\varprojlim_n M_n$ (W1) there is $b\in\pkpA$ with image $\beta_n$ in every $M_n$. Thus $b-a_k\in
\ell^n\pkpA+t^{k-n+1}\pkpA$ for $1\le n\le k$. Taking $n=\lceil k/2\rceil$ (for $k\ge1$; the case $k=0$
is trivial): $b-a_k\in\ell^{\lceil k/2\rceil}\pkpA+t^{\lfloor k/2\rfloor+1}\pkpA\subseteq\pkrI^{\lfloor
k/2\rfloor}$.
\end{proof}

\begin{remark}\label{pkr:rem-complete}
Part (ii) gives convergence modulo $\pkrI^{\lfloor k/2\rfloor}$ rather than modulo $\pkrI^k$. This is all
that is used: part (ii) is applied only in Step~2 of the proof of Proposition~\ref{pkr:I3}, where it is
needed only that $b-a_k\in\ell^{c_k}\pkfAm$ with $c_k\to\infty$, and $\pkrI^{\lfloor k/2\rfloor}\subseteq
\ell^{\lfloor k/2\rfloor}\pkfAm$. Part (i) is used in Propositions~\ref{pkr:I5}, \ref{pkr:I6} and at the
beginning of \S\ref{pkr:ss-div}, in the form stated. (The $\pkrI$-adic completeness in the usual sense
also holds, since $\pkrI^{2k}\subseteq\ell^k\pkpA+t^{k+1}\pkpA\subseteq\pkrI^k$ shows that the
$\pkrI$-adic topology is the one defined by the ideals $\ell^k\pkpA+t^{k+1}\pkpA$, and Step~4 constructs
limits for the latter; it is not needed.)
\end{remark}

\begin{lemma}[Saturated ideals]\label{pkr:sat}
Let $J$ be an ideal of $\pkpA$ such that $\ell x\in J$ implies $x\in J$, and let $\delta\in J$ be such
that $\delta_0$ divides $x_0$ in $O^\flat$ for every $x\in J$. Then $J=\delta\pkpA$.
\end{lemma}

\begin{proof}
Let $x\in J$. As $\delta_0\mid x_0$, write $x_0=\delta_0c$ and $y_0=[c]$; then $x-\delta y_0\equiv
[x_0]-[\delta_0][c]=0\pmod{\ell}$ (W1), so $x-\delta y_0=\ell x^{(1)}$, and $x^{(1)}\in J$ by saturation
(as $x-\delta y_0\in J$). Iterating, $x=\delta(y_0+\ell y_1+\dots+\ell^ky_k)+\ell^{k+1}x^{(k+1)}$ for all
$k$. By $\ell$-adic completeness and separation of $\pkpA$, $y=\sum_k\ell^ky_k$ converges and $x=\delta
y$.
\end{proof}

\begin{proposition}[Input (A7)]\label{pkr:I7}
Let $\delta,c\in\pkpA$ with $\ell\nmid\delta$, and suppose that for every $m$ there is $z_m\in\pkpA$
with $c-\delta z_m\in\pkrI^m$. Then $\delta\mid c$.
\end{proposition}

\begin{proof}
\emph{Step 1: for every $m\ge1$ there is $N_m$ with $\pkrI^{N_m}\subseteq(\ell^m,\delta)$.} Since
$\ell\nmid\delta$, $\delta_0\ne0$; and $E_0\in\mathfrak m^\flat$ has $v^\flat(E_0)=1>0$. Choose $r$ with
$r\ge v^\flat(\delta_0)$; then $\delta_0\mid E_0^r$ in the valuation ring $O^\flat$, say $E_0^r=\delta_0c'$,
and $E^r-\delta[c']\equiv0\pmod\ell$ (W1). So $E^r\in(\ell,\delta)$, hence $E^{rm}\in(\ell,\delta)^m\subseteq
(\ell^m,\delta)$. Put $N_m=m(r+1)$: a monomial $\ell^iE^{N_m-i}$ has $i\ge m$ or $N_m-i\ge rm$, so it lies
in $(\ell^m,\delta)$.

\emph{Step 2.} Let $J=\bigcap_{m\ge1}(\ell^m,\delta)$. Then $c\in J$: $c=(c-\delta z_{N_m})+\delta z_{N_m}
\in\pkrI^{N_m}+\delta\pkpA\subseteq(\ell^m,\delta)$. Clearly $\delta\in J$. $J$ is saturated: if
$\ell x\in J$, then for each $m$, $\ell x=\ell^{m+1}b+\delta y$ for some $b,y$, so $\ell\mid\delta y$,
and as $\ell$ is prime with $\ell\nmid\delta$, $y=\ell y'$; cancelling $\ell$ ($\pkpA$ torsion free),
$x=\ell^mb+\delta y'\in(\ell^m,\delta)$. Finally every $x\in J$ lies in $(\ell,\delta)$, so
$x=\ell b+\delta y$ and $x_0=\delta_0y_0$. By Lemma~\ref{pkr:sat}, $J=\delta\pkpA$, so $\delta\mid c$.
\end{proof}

\subsubsection{Topological bases over $\mathfrak S$: inputs (A5) and (A6)}\label{pkr:ss-basis}

For an $\mathfrak S$-module $M$ and $n\ge0$, $\pkrm^nM=\sum_{i+j=n}\ell^iz^jM$. For a set $I$ let
$F_I$ be the free module of finitely supported families $f=(f_i)_{i\in I}$ in $\mathfrak S$.

\begin{proposition}[Topological bases]\label{pkr:basis}
Let $M$ be an $\mathfrak S$-module such that
\begin{enumerate}
\item[(M1)] $\ell$ acts injectively on $M$;
\item[(M2)] if $x\in M$ and $zx\in\ell M$, then $x\in\ell M$;
\item[(M3)] $\bigcap_n\pkrm^nM=0$.
\end{enumerate}
Let $T\subseteq M$ be a subset whose image in the $k$-vector space $M/\pkrm M$ is linearly independent.
Then there are a family $(b_i)_{i\in I}$ in $M$ containing $T$ and $\mathfrak S$-linear maps
$\kappa_i\colon M\to\mathfrak S$ ($i\in I$) such that: (a) $\kappa_i(\sum_jf_jb_j)=f_i$ for $f\in F_I$;
(b) for every $x\in M$ and $N\ge0$ there is a finite $G_0\subseteq I$ with $x-\sum_{i\in G}\kappa_i(x)b_i
\in\pkrm^NM$ for every finite $G\supseteq G_0$; (c) if $\kappa_i(x)=0$ for all $i$, then $x=0$.
\end{proposition}

\begin{proof}
Extend the image of $T$ to a $k$-basis of $M/\pkrm M$ and lift the added vectors: this gives
$(b_i)_{i\in I}\supseteq T$ whose image is a $k$-basis of $M/\pkrm M$. Let $\beta\colon F_I\to M$,
$\beta(f)=\sum_if_ib_i$.

\emph{Step 1 (density): for $x\in M$ and $N\ge0$ there is $f\in F_I$ with $x-\beta(f)\in\pkrm^NM$.}
For $N=1$ this holds because the $b_i$ span $M/\pkrm M$ over $k=\mathfrak S/\pkrm$. If $x-\beta(f)=
\sum_{i+j=N}\ell^iz^jy_{ij}$, choose $g_{ij}$ with $y_{ij}-\beta(g_{ij})\in\pkrm M$; then
$x-\beta(f+\sum\ell^iz^jg_{ij})\in\pkrm^{N+1}M$.

\emph{Step 2 (a criterion).} Let $a,s\in\mathfrak S$ be such that (C1) $\beta(f)\in(a,s)M$ implies
$f\in(a,s)F_I$, and (H) $aw=sy$ with $w,y\in M$ implies $y\in aM$. Then for every $n\ge1$,
$\beta(f)\in(s^n,a)M$ implies $f\in(s^n,a)F_I$. Induction on $n$, the case $n=1$ being (C1). Let
$\beta(f)=s^{n+1}w'+au$ with $w',u\in M$. By (C1), $f=au'+sc$ with $u',c\in F_I$. Then
$a(u-\beta(u'))=s(\beta(c)-s^nw')$, so by (H) $\beta(c)-s^nw'=aw$ with $w\in M$, i.e.\ $\beta(c)\in
(s^n,a)M$. By induction $c=s^nm+am'$ with $m,m'\in F_I$, and $f=s^{n+1}m+a(u'+sm')$.

\emph{Step 3: $\beta(f)\in\pkrm^{2n}M$ implies $f\in\pkrm^nF_I$.} First take $(a,s)=(\ell,z)$: (C1) is
the linear independence of the images of the $b_i$ in $M/\pkrm M$ over $k$, and (H) is (M2). So
$\beta(f)\in(z^n,\ell)M$ implies $f\in(z^n,\ell)F_I$. Next take $(a,s)=(z^n,\ell)$: (C1) is what was
just proved; (H) says that $z^nw=\ell y$ implies $y\in z^nM$, which we prove by induction on $n$ (the
case $n=0$ is trivial): $z(z^{n-1}w)=\ell y$ gives $z^{n-1}w=\ell c$ by (M2); then $\ell(zc-y)=0$, so
$y=zc$ by (M1); and $c\in z^{n-1}M$ by induction, so $y\in z^nM$. Hence $\beta(f)\in(\ell^n,z^n)M$
implies $f\in(\ell^n,z^n)F_I$. Since $\pkrm^{2n}\subseteq(\ell^n,z^n)\subseteq\pkrm^n$, the claim
follows.

\emph{Step 4 (coordinates).} Let $x\in M$. An \emph{approximation of level $N$} of $x$ is an $f\in F_I$
with $x-\beta(f)\in\pkrm^{2N}M$; they exist by Step 1. If $f$ and $f'$ are approximations of levels
$N\le N'$, then $\beta(f-f')\in\pkrm^{2N}M$, so $f_i-f'_i\in\pkrm^N$ for all $i$ (Step 3). Choosing an
approximation $f^{(N)}$ of level $N$ for each $N$, $(f^{(N)}_i)_N$ is a Cauchy sequence in $\mathfrak S$;
let $\kappa_i(x)$ be its limit (W4). Then $\kappa_i(x)-f_i\in\pkrm^N$ for every approximation $f$ of
level $N$ (as $\pkrm^N$ is closed), so $\kappa_i(x)$ does not depend on the choices. If $f,g$ are
approximations of level $N$ of $x,y$ and $s\in\mathfrak S$, then $f+g$ and $sf$ are approximations of
level $N$ of $x+y$ and $sx$; hence $\kappa_i$ is $\mathfrak S$-linear. For $x=\beta(f)$, the constant
choice $f^{(N)}=f$ gives (a).

(b) Let $f$ be an approximation of level $N$ of $x$ and $G_0=\operatorname{supp}f$. For finite
$G\supseteq G_0$,
\[x-\sum_{i\in G}\kappa_i(x)b_i=(x-\beta(f))-\sum_{i\in G}(\kappa_i(x)-f_i)b_i\in\pkrm^{2N}M+\pkrm^NM
\subseteq\pkrm^NM.\]
(c) follows from (b) and (M3).
\end{proof}

\begin{proposition}[Input (A5)]\label{pkr:I5}
There are a family $(b_i)_{i\in I}$ in $\pkpA$ and $\mathfrak S$-linear maps $\kappa_i\colon\pkpA\to
\mathfrak S$ such that for every $x\in\pkpA$ and $N\ge0$, $x-\sum_{i\in G}\kappa_i(x)b_i\in
\pkrm^N\pkpA=(\ell,z)^N\pkpA$ for all sufficiently large finite $G\subseteq I$ (in particular $x=0$ if
all $\kappa_i(x)=0$); and there is $0\in I$ with $b_0=1$ and $\kappa_0(s)=s$ for $s\in\mathfrak S$.
\end{proposition}

\begin{proof}
Apply Proposition~\ref{pkr:basis} to $M=\pkpA$ and $T=\{1\}$. (M1) is (W1). (M2): if $zx=\ell y$, then
$\ell\mid zx$, and $\ell$ is prime with $\ell\nmid z$ ($z_0=\varpi^\flat\ne0$), so $\ell\mid x$. (M3) is
Proposition~\ref{pkr:complete}(i). The image of $1$ in $\pkpA/\pkrm\pkpA$ is nonzero: if $1=\ell
x+zy$, then $1=\varpi^\flat y_0$ in $O^\flat$, so $\varpi^\flat$ would be a unit, while
$\varpi^\flat\in\mathfrak m^\flat$. Let $0\in I$ be the index with $b_0=1$; for $s\in\mathfrak S$,
$s=\beta(f)$ with $f$ supported at $0$, $f_0=s$, so $\kappa_0(s)=s$ by (a).
\end{proof}

\begin{proposition}[Input (A6)]\label{pkr:I6}
Let $P\subseteq\pkpA^d$ be an $\mathfrak S$-submodule containing no $d+1$ elements linearly
independent over $\mathfrak S$, such that $za\in\ell P$ with $a\in P$ implies $a\in\ell P$. Then $P$
is a free $\mathfrak S$-module of finite rank $\le d$.
\end{proposition}

\begin{proof}
Apply Proposition~\ref{pkr:basis} to $M=P$ and $T=\emptyset$. (M1): $\pkpA^d$ is $\ell$-torsion free.
(M2) is the hypothesis. (M3): $\pkrm^nP\subseteq(\pkrm^n\pkpA)^d$, and $\bigcap_n\pkrm^n\pkpA=0$
(Proposition~\ref{pkr:complete}(i)). Let $(b_i)_{i\in I}$, $(\kappa_i)$ be as given there. Any finitely
many distinct $b_i$ are linearly independent over $\mathfrak S$: if $\sum_{i\in G}f_ib_i=0$, then
$f_i=\kappa_i(0)=0$ by (a). Hence $|I|\le d$. For $x\in P$, (b) with $G=I$ gives $x-\sum_{i\in I}
\kappa_i(x)b_i\in\bigcap_N\pkrm^NP=0$. So $(b_i)_{i\in I}$ is a basis of $P$.
\end{proof}

\subsubsection{Division in $\pkfAm$ and approximation by $\pkpA$: inputs (A3), (A2), (A4)}\label{pkr:ss-div}

We identify $\pkfAm$ with $\pkfAo\subseteq\pkfBp$ (Lemma~\ref{pkr:Amax}(i)), and $\pkpA$ with its image
in $\pkfAm$ (injective by (P1), or by Proposition~\ref{pkr:I1} and Proposition~\ref{pkr:complete}(i):
the kernel lies in $\bigcap_n\pkrI^n$).

\begin{lemma}[The element $\tau$]\label{pkr:tau}
$\tau$ divides $E-\ell v^{-1}$ in $\pkpA$. In $\pkfAm$, $E-\ell v^{-1}=\ell(E/\ell-v^{-1})$.
\end{lemma}

\begin{proof}
By Lemma~\ref{pkn:mu}, $\mu=\tau\xi_\varepsilon$ with $\xi_\varepsilon=\sum_{k=0}^{\ell-1}
[\varepsilon^{1/\ell}]^k=Eu$, $u\in\pkpA^\times$; since also $\mu=E\tau v$ and $\pkpA$ is a domain
($\mu\ne0$), $u=v$. Now $\xi_\varepsilon-\ell=\sum_{k<\ell}([\varepsilon^{1/\ell}]^k-1)$ and each term is
divisible by $[\varepsilon^{1/\ell}]-1=\tau$. So $\tau\mid Ev-\ell$, hence $\tau\mid v^{-1}(Ev-\ell)=
E-\ell v^{-1}$.
\end{proof}

\begin{lemma}[Division]\label{pkr:div}
Let $\nu,b\in\pkpA$ with $\nu\mid E-\ell b$. For every $n\ge0$ and $y\in\pkfAm$ there are $a\in\pkpA$
and $w\in\pkfAm$ with $\ell^ny=a+\nu^nw$. This applies to $\nu=E$ ($b=0$) and to $\nu=\tau$
($b=v^{-1}$, Lemma~\ref{pkr:tau}).
\end{lemma}

\begin{proof}
\emph{Step 1: every $y\in\pkfAm$ is $c+(E/\ell-b)w$ with $c\in\pkpA$, $w\in\pkfAm$.} By
Lemma~\ref{pkf:approx0}(i), $y=P(E/\ell)+\ell y_1$ with $P\in\pkpA[X]$, $y_1\in\pkfAm$. Division by the
monic $X-b$ gives $P(X)=P(b)+(X-b)Q(X)$, so $y=c_0+(E/\ell-b)w_0+\ell y_1$ with $c_0=P(b)\in\pkpA$,
$w_0=Q(E/\ell)$. Repeating with $y_1,y_2,\dots$:
$y=\sum_{k<K}\ell^kc_k+(E/\ell-b)\sum_{k<K}\ell^kw_k+\ell^Ky_K$ for all $K$. The series
$c=\sum\ell^kc_k$ converges in $\pkpA$ and $w=\sum\ell^kw_k$ in $\pkfAm$ ((W1), Lemma~\ref{pkr:Amax}(i)),
and the map $\pkpA\to\pkfAm$ carries $\ell^K\pkpA$ into $\ell^K\pkfAm$; so $y-c-(E/\ell-b)w\in\ell^K
\pkfAm$ for every $K$, hence is $0$.

\emph{Step 2.} Write $E-\ell b=\nu g$, $g\in\pkpA$; in $\pkfAm$, $\ell(E/\ell-b)=\nu g$. Induction on
$n$, the case $n=0$ being $a=0$, $w=y$. If $\ell^ny=a+\nu^nw$, write $w=c+(E/\ell-b)w'$ by Step 1; then
$\ell^{n+1}y=\ell a+\ell\nu^nc+\nu^n\ell(E/\ell-b)w'=(\ell a+\ell\nu^nc)+\nu^{n+1}gw'$.
\end{proof}

\begin{lemma}[Divisibility by $E^h$]\label{pkr:Ediv}
If $b\in\pkpA$ and $E^h\mid b$ in $\pkfBp$, then $E^h\mid b$ in $\pkpA$.
\end{lemma}

\begin{proof}
Induction on $h$. If $b=E^hw$ with $h\ge1$, $w\in\pkfBp$, then $\theta(b)=0$, so $b=Eb_1$ with
$b_1\in\pkpA$ (Corollary~\ref{pkn:E-gen}); then $E(b_1-E^{h-1}w)=0$ in the domain $\pkfBp$ (P1), and
$E\ne0$; so $E^{h-1}\mid b_1$ in $\pkfBp$, and by induction in $\pkpA$.
\end{proof}

\begin{proposition}[Input (A3): contraction]\label{pkr:I3}
Let $V\subseteq\pkfAo$, $h\ge0$ and $u\in\pkpA$ be such that every $\alpha\in V$ has some $\alpha'\in V$
with $E^h\alpha=u\,\varphi(\alpha')$. Then $\ell^{h+1}V\subseteq\pkpA$.
\end{proposition}

\begin{proof}
\emph{Step 1: for every $k\ge0$ and $\alpha\in V$ there are $a\in\pkpA$ and $w\in\pkfAm$ with
$a=\ell^{h+1}\alpha-\ell^k\tau^{h+1}w$.} For $k=0$ this is Lemma~\ref{pkr:div} with $\nu=\tau$,
$n=h+1$. Suppose it holds for $k$ (for all elements of $V$), and let $\alpha\in V$, with $\alpha'$ as
in the hypothesis; choose $a',w'$ for $\alpha'$ and $k$: $a'=\ell^{h+1}\alpha'-\ell^k\tau^{h+1}w'$.
Apply $\varphi$ (which acts on $\pkfAm$ compatibly with $\pkpA$ and $\pkfBp$, (P1)) and use
$\varphi(\tau)=\mu=E\tau v$:
$\varphi(a')=\ell^{h+1}\varphi(\alpha')-\ell^kE^{h+1}\tau^{h+1}v^{h+1}\varphi(w')$. Multiply by $u$ and
use $u\varphi(\alpha')=E^h\alpha$ and $E=\ell\cdot E/\ell$:
\[u\varphi(a')=E^h\bigl(\ell^{h+1}\alpha-\ell^{k+1}\tau^{h+1}w\bigr),\qquad w=(E/\ell)\,v^{h+1}u\,
\varphi(w')\in\pkfAm.\]
So $E^h\mid u\varphi(a')$ in $\pkfBp$, and by Lemma~\ref{pkr:Ediv} $u\varphi(a')=E^ha$ with $a\in\pkpA$.
Cancelling $E^h$ in the domain $\pkfBp$: $a=\ell^{h+1}\alpha-\ell^{k+1}\tau^{h+1}w$.

\emph{Step 2.} Fix $\alpha\in V$ and choose $a_k,w_k$ as in Step 1 for every $k$. For $k\le l$,
$a_k-a_l=\ell^k\tau^{h+1}(\ell^{l-k}w_l-w_k)\in\pkpA\cap\ell^k\pkfAm=\pkrI^k$
(Proposition~\ref{pkr:I1}). By Proposition~\ref{pkr:complete}(ii) there is $b\in\pkpA$ with
$b-a_k\in\pkrI^{\lfloor k/2\rfloor}\subseteq\ell^{\lfloor k/2\rfloor}\pkfAm$. Then
$\ell^{h+1}\alpha-b=(a_k-b)+\ell^k\tau^{h+1}w_k\in\ell^{\lfloor k/2\rfloor}\pkfAm$ for every $k$, so
$\ell^{h+1}\alpha=b\in\pkpA$ by separation (Lemma~\ref{pkr:Amax}(i)).
\end{proof}

\begin{proposition}[Input (A2)]\label{pkr:I2}
If $\delta\in\pkpA$ divides $\mu$, $h\ge0$, $w\in\pkfBp$ and $\delta^hw\in\pkfAo$, then $\ell^{2h}w\in\pkfAo$.
\end{proposition}

\begin{proof}
By Lemma~\ref{pkr:tau}, $E-\ell v^{-1}=\tau s$ with $s\in\pkpA$, so $E(E-\ell v^{-1})v=E\tau v\,s=\mu s$
is divisible by $\delta$: write $E(E-\ell v^{-1})v=\delta\,r$ with $r\in\pkpA$. In $\pkfAm$,
$E(E-\ell v^{-1})v=\ell^2g$ with $g=v\,(E/\ell)(E/\ell-v^{-1})=v\,f(E/\ell)$, $f=X(X-v^{-1})$ monic.
Choose $m$ with $\ell^m\ell^{2h}w=y\in\pkfAo$, and let $\delta^hw=c\in\pkfAo$. In $\pkfBp$,
\[g^hy=\ell^m(\ell^2g)^hw=\ell^m(\delta r)^hw=\ell^mr^hc,\]
an identity in $\pkfAm$ (injectivity). So $\ell^m\mid v^hf(E/\ell)^hy$, hence (as $v$ is a unit)
$\ell^m\mid f^h(E/\ell)\,y$ with $f^h$ monic, and $\ell^m\mid y$ by Lemma~\ref{pkr:Amax}(iii): $y=\ell^my'$
with $y'\in\pkfAm$, and $\ell^{2h}w=y'$ ($\ell$ is a unit of $\pkfBp$).
\end{proof}

\begin{proposition}[Input (A4): approximation]\label{pkr:I4}
For $y\in\pkfBp$ and $h\ge0$ there are $m\ge0$, $\alpha\in\pkpA$ and $w\in\pkfAo$ with
$\ell^my=\alpha+\ell^m\mu^hw$.
\end{proposition}

\begin{proof}
By Lemma~\ref{pkf:approx0}(ii) with $n=2h$: $\ell^my=\alpha_0+\ell^{m+2h}w_0$ with $\alpha_0\in\pkpA$,
$w_0\in\pkfAo$. By Lemma~\ref{pkr:div} with $\nu=\tau$: $\ell^hw_0=a_1+\tau^hw_1$; with $\nu=E$:
$\ell^hw_1=a_2+E^hw_2$ ($a_1,a_2\in\pkpA$, $w_1,w_2\in\pkfAm$). Then
\[\ell^{2h}w_0=\ell^ha_1+\tau^ha_2+(E\tau)^hw_2=\ell^ha_1+\tau^ha_2+\mu^hv^{-h}w_2,\]
and $\ell^my=\alpha+\ell^m\mu^hw$ with $\alpha=\alpha_0+\ell^m(\ell^ha_1+\tau^ha_2)\in\pkpA$ and
$w=v^{-h}w_2\in\pkfAo$.
\end{proof}

\subsubsection{Divisibility by $\mu$ and $\tau$, the chart facts (A8), and (P4)}\label{pkr:ss-chart}

\begin{lemma}[Divisibility by $\mu$ in $\pkfBp$]\label{pkr:mu-Bplus}
Let $c\in\pkfBp$ with $\theta(\varphi^nc)=0$ for all $n\ge0$. Then $\mu\mid c$ in $\pkfBp$.
\end{lemma}

\begin{proof}
By Proposition~\ref{pkr:I4} with $h=1$, $\ell^mc=\alpha+\ell^m\mu w$ with $\alpha\in\pkpA$, $w\in\pkfAo$.
For $n\ge0$, $\theta(\varphi^n\mu)=\theta([\varepsilon^{\ell^n}])-1=(\varepsilon^{\ell^n})^\sharp-1=0$, and
$\theta\circ\varphi^n$ is a ring map on $\pkfBp$ extending the one on $\pkpA$ (P1); so
$\theta(\varphi^n\alpha)=\ell^m\theta(\varphi^nc)-\ell^m\theta(\varphi^n\mu)\theta(\varphi^nw)=0$ for all
$n\ge0$. By Lemma~\ref{pkn:mu}, $\mu\mid\alpha$ in $\pkpA$. Hence $\mu\mid\ell^mc$, and $\mu\mid c$ since
$\ell$ is a unit of $\pkfBp$.
\end{proof}

\emph{The field $K'$ inside $\pknBp$.} As in \S\ref{pkf:sec}, $K'$ (indeed $\auQlb$) is regarded as a
subfield of $\pknBp$ through Fontaine's embedding $\iota$, characterised by $\theta\circ\iota=\mathrm{id}$;
and $K_0=\Lambda[1/\ell]$ is regarded as a subfield of $\auQlb$ through $\theta$. By the proof of
Lemma~\ref{pkf:invariants}, $K'=K_0(\pi)$ with $\pi=\theta(z)$.

\begin{lemma}[Hensel uniqueness and Newton approximation in $\pknBp$]\label{pkr:newton}
(i) Let $P$ be a separable polynomial with coefficients in a subfield $F\subseteq\pknBp$ with $\theta|_F$
injective, and $x,x'\in\pknBp$ roots of $P$ with $\theta(x)=\theta(x')$. Then $x=x'$.
(ii) The image of $K_0=\Lambda[1/\ell]$ in $\pknBp$ through $\pkpA[1/\ell]$ is $\iota(K_0)$, elementwise:
for $c\in\Lambda[1/\ell]$, the image of $c$ equals $\iota(\theta(c))$.
(iii) For every $\lambda\in K'$ and $m\ge0$ there is $s\in\pkfSB$ with $E^m\mid s-\lambda$ in $\pknBp$.
\end{lemma}

\begin{proof}
Recall that $\pknBp$ is a local ring with maximal ideal $\ker\theta=E\pknBp$ (Lemma~\ref{pkn:BdR}), so
an element with $\theta\ne0$ is a unit.
(i) By Taylor's formula, $P(x')-P(x)=(x'-x)\bigl(P'(x)+(x'-x)r\bigr)$ with $r\in\pknBp$. The second
factor has $\theta$-image $P^\theta{}'(\theta x)\ne0$, where $P^\theta$ is $P$ with $\theta$ applied to the
coefficients (separable, with the simple root $\theta x$). So it is a unit and $x'=x$.
(ii) $\Lambda=W(\mathbb F_{\ell^f})$ contains a primitive $(\ell^f-1)$-th root of unity $\zeta$ (a
Teichm\"uller representative), and $K_0=\Q_\ell(\zeta)$. The image of $\zeta$ and $\iota(\theta(\zeta))$
are roots of $X^{\ell^f-1}-1$ (separable over $\Q_\ell$) with the same $\theta$; by (i) they are equal.
Both maps are $\Q_\ell$-algebra maps on $K_0=\Q_\ell[\zeta]$, so they agree.
(iii) Write $\lambda=\sum_{j<e}c_j\pi^j$ with $c_j\in K_0$ ($K'=K_0(\pi)$ has degree $e$ over $K_0$: $K_0$ is the maximal unramified subextension of $K'$, which has ramification index $e$, Lemma~\ref{pkf:invariants}; equivalently, $\pi$ is a root of the Eisenstein factor of $E(X)$, Definition~\ref{pkp:def-chart}) and put $s_0=\sum_jc_jz^j\in\pkfBp$. It lies in
$\pkfSB$, since $z\in\mathfrak S$, $\Lambda\subseteq\mathfrak S$ and $\ell^{-1}\in\pkfSB$
(Lemma~\ref{pkf:SB-basic}(i)). By (ii), $\theta(s_0)=\sum_jc_j\pi^j=\lambda=\theta(\iota\lambda)$, so
$E\mid s_0-\lambda$ in $\pknBp$. Let $P\in K_0[X]$ be the minimal polynomial of $\lambda$ (separable,
characteristic $0$), and $g,g_1\in K_0[X]$ with $g_1P+gP'=1$. Put $\Phi=X-gP\in K_0[X]$. Then
$\Phi(\lambda)=\lambda$ and $\Phi'(\lambda)=1-g'(\lambda)P(\lambda)-g(\lambda)P'(\lambda)=g_1(\lambda)
P(\lambda)=0$, so by Taylor's formula $\Phi(x)-\lambda=(x-\lambda)^2r_x$ with $r_x\in\pknBp$ for every
$x\in\pknBp$ (the coefficients of $\Phi$ lie in $K_0\subseteq\pknBp$, compatibly by (ii)). Put
$s_{n+1}=\Phi(s_n)$; as $\pkfSB$ is a ring containing $K_0$, $s_n\in\pkfSB$, and by induction
$E^{2^n}\mid s_n-\lambda$. Take $s=s_n$ with $2^n\ge m$.
\end{proof}

\begin{proposition}[Input (A8)]\label{pkr:I8}
(a) In the notation of (A5): for $i\in I$, $x\in\pkpA$ and $m\ge0$ with $\theta(\varphi^{-m}x)\in\ell
O_C$, there is $s\in\mathfrak S$ with $p_m\mid\kappa_i(x)-\ell s$ in $\mathfrak S$.
(b) If $b\in\pkfBp$ and $\varphi(\tau)\mid\varphi(b)$ in $\pkfBp$, then $\tau\mid b$ in $\pkfBp$.
(c) Let $G_f$ be the extension of Lemma~\ref{pkf:extend}. For $\lambda\in K'^d$, $y\in(\pkfBp)^d$ and an
$\mathfrak S$-linear $f\colon\pkpA\to\pkpA$, $\sum_i\lambda_iy_i$ divides $\sum_i\lambda_iG_f(y_i)$ in
$\pknBp$ (images in $\pknBp$).
\end{proposition}

\begin{proof}
(a) Write $\theta_m=\theta\circ\varphi^{-m}\colon\pkpA\to O_C$, a surjective ring map (W3) with kernel
$p_m\pkpA$ (Corollary~\ref{pkn:E-gen}), and $\theta_m(\ell)=\ell$. By hypothesis $\theta_m(x)=\ell c$ with
$c\in O_C$; choose $a\in\pkpA$ with $\theta_m(a)=c$. Then $\theta_m(x-\ell a)=0$, so $x-\ell a=p_mq$ with
$q\in\pkpA$. As $\kappa_i$ is $\mathfrak S$-linear and $\ell,p_m\in\mathfrak S$ (W4),
$\kappa_i(x)-\ell\kappa_i(a)=p_m\kappa_i(q)$; take $s=\kappa_i(a)$.

(b) Write $\varphi(b)=\mu c$ with $c\in\pkfBp$ ($\varphi(\tau)=\mu$). For $n\ge0$,
$\theta(\varphi^{n+1}b)=\theta(\varphi^n\mu)\theta(\varphi^nc)=0$ (proof of Lemma~\ref{pkr:mu-Bplus}). By
Proposition~\ref{pkr:I4} with $h=1$, $\ell^mb=\alpha+\ell^m\mu w$; as in the proof of
Lemma~\ref{pkr:mu-Bplus}, $\theta(\varphi^{n+1}\alpha)=0$ for all $n\ge0$, i.e.\ $\varphi(\alpha)$
satisfies the hypothesis of Lemma~\ref{pkn:mu}, so $\varphi(\tau)=\mu\mid\varphi(\alpha)$ in $\pkpA$ and,
$\varphi$ being an automorphism of $\pkpA$ (W3), $\tau\mid\alpha$. Also $\tau\mid\mu$. Hence
$\tau\mid\ell^mb$ in $\pkfBp$, and $\tau\mid b$.

(c) Put $x=\sum_i\lambda_iy_i$ and $x'=\sum_i\lambda_iG_f(y_i)$ in $\pknBp$. As $\pknBp$ is a discrete
valuation ring with uniformiser $E$ (Lemma~\ref{pkn:BdR}), it suffices to show: for every $m\ge0$,
$E^m\mid x$ implies $E^m\mid x'$ (if $x\ne0$ take $m$ its valuation; if $x=0$, $x'\in\bigcap_mE^m\pknBp=0$).
Fix $m$ and choose by Lemma~\ref{pkr:newton}(iii) $s_i\in\pkfSB$ with $E^m\mid s_i-\lambda_i$ in
$\pknBp$. Then $E^m$ divides $x-\sum_is_iy_i$ and $x'-\sum_is_iG_f(y_i)$ in $\pknBp$. If $E^m\mid x$, then
$E^m$ divides the image of $\sum_is_iy_i\in\pkfBp$ in $\pknBp$, hence divides it in $\pkfBp$
(Lemma~\ref{pkf:approx0}(v) and (P6)): $\sum_is_iy_i=E^mq$, $q\in\pkfBp$. Since $G_f$ is
$\pkfSB$-linear (Lemma~\ref{pkf:extend}) and $s_i,E\in\pkfSB$,
$\sum_is_iG_f(y_i)=G_f(\sum_is_iy_i)=E^mG_f(q)$. Hence $E^m\mid x'$ in $\pknBp$.
\end{proof}

\begin{proposition}[(P4) proved]\label{pkr:P4}
If $b\in\pkfBm$ and $m\ge0$ are such that the image of $\varphi^n(b)$ in $\pknB$ lies in
$\mathrm{Fil}^{-m}$ for every $n\ge0$, then $\mu^mb\in\pkfBp$.
\end{proposition}

\begin{proof}
As $\pkfBm=\pkfBp[1/\mu]$, there are $N\ge0$ and $c\in\pkfBp$ with $\mu^{m+N}b=c$. We show by induction
on $N$ that $\mu^mb\in\pkfBp$. If $N=0$ there is nothing to prove. Let $N\ge1$. For $n\ge0$,
$\varphi^n(c)=\varphi^n(b)\varphi^n(\mu)^{m+N}$ in $\pkfBm$, hence in $\pknB$ (P1). Since
$\theta(\varphi^n\mu)=0$, $\varphi^n(\mu)\in E\pkpA$ (Corollary~\ref{pkn:E-gen}), whose image lies in
$\mathrm{Fil}^1$; so $\varphi^n(c)\in\mathrm{Fil}^{-m}\cdot\mathrm{Fil}^{m+N}\subseteq\mathrm{Fil}^N
\subseteq\mathrm{Fil}^1$ ($\mathrm{Fil}^k=E^k\pknBp$, Lemma~\ref{pkn:BdR}). As $\theta$ on $\pkfBp$ is the
restriction of $\theta$ on $\pknBp$, whose kernel is $\mathrm{Fil}^1$, $\theta(\varphi^nc)=0$ for all
$n\ge0$. By Lemma~\ref{pkr:mu-Bplus}, $c=\mu d$ with $d\in\pkfBp$, and $\mu^{m+N-1}b=d$ in the domain
$\pkfBm$ (P1). Conclude by induction.
\end{proof}

With Propositions~\ref{pkr:I1}--\ref{pkr:P4}, the inputs (A1)--(A8) of \S\ref{pkf:ss-rings} and the
statement (P4) are proved; the existence of frames (Theorem~\ref{pkf:main}) therefore rests only on
(P1)--(P3), (P5), (P6) and the standard facts (W1)--(W4).



\section{Lattices with orders}\label{C4:g4}
\providecommand{\auO}{\mathcal O}
\providecommand{\auQlb}{\overline{\Q}_\ell}
\providecommand{\pkpA}{A_{\mathrm{inf}}}
\providecommand{\pkpW}{W(C^\flat)}
\providecommand{\pkpAmb}{\mathbb A}
\providecommand{\pknMinf}{\mathfrak M_{\mathrm{inf}}}
\providecommand{\pknK}{\mathcal K_0}
\providecommand{\pknBp}{B^+_{\mathrm{dR}}}
\providecommand{\pknB}{B_{\mathrm{dR}}}
\providecommand{\pkfBp}{B^+}
\providecommand{\pkfBm}{B_{\max}}
\providecommand{\pkfSB}{\mathfrak S_B}
\providecommand{\pklAN}{\pkpA\mathcal N}
\providecommand{\pklL}{\pkpA\otimes L}
\providecommand{\pklstd}{\mathrm{std}}

\subsection{Lattices with orders: standard orders, Galois stability and the retraction}\label{pkl:sec}

This section continues \S\S\ref{pkn:sec}--\ref{pkr:sec} and uses their conventions: a fixed Kummer
chart (Definition~\ref{pkp:def-chart}) with $\mathfrak S=\Lambda[[z]]$, $E$, $p_a=\varphi^a(E)$, $\Gamma'$,
$K'=\auQlb^{\,j(\Gamma')}$; the modules $\pknMinf(L)\subseteq\pklAN(L)\supseteq\mathcal N(L)$ and
$L_i(o)$ of Definition~\ref{pkn:def-N}; $\mathrm{Fil}^n=t^n\pknBp=E^n\pknBp$ (Lemma~\ref{pkn:BdR}); the
rings $\pkfBp\subseteq\pkfBm$ and $\pkfSB$ of \S\ref{pkf:ss-rings}--\S\ref{pkf:ss-SB}; frames
(Definition~\ref{pkn:def-frame}), which exist with $a$ the jump bound (Theorem~\ref{pkf:main}). We write
$\pklL$ for the $\pkpA$-span of the $1\otimes y$ in $\pkpAmb(L)=\pkpW\otimes_{\Z_\ell}L$, and
$\pkpAmb(L)_i$ for the letter part (Definition~\ref{pkp:def-letters}); $\sigma(i)=i^\ell$. We prove the
items of Lemma~\ref{pkm:inputs} not proved there: (F), (H), (D), (C)(d), (R) and (G).

\emph{Dependencies.} Proved in full, modulo the results cited from \S\S\ref{pkn:sec}--\ref{pkr:sec} (hence
modulo the published inputs listed there) and the reading of ``Hodge--Tate jumps'' recorded in
\S\ref{pkl:ss-gap}: Proposition~\ref{pkl:std} and (F) (Corollary~\ref{pkl:F}); (H), first assertion in
full (Corollary~\ref{pkl:H1}), second and third assertions by
Proposition~\ref{pkl:H2} and Corollary~\ref{pkl:H3}; the $\Gamma'$-stability part of (G), first assertion
(Lemma~\ref{pkl:G1}); (D) (Proposition~\ref{pkl:D}); (C)(d) (Proposition~\ref{pkl:Cd}); (R)
(Proposition~\ref{pkl:R}); $\varphi(\mathcal N)\subseteq\mathcal N$ (Corollary~\ref{pkl:Phi}); the second assertion of (G) and (G$_0$) (Proposition~\ref{pkl:G0}).

\subsubsection{Preliminaries}\label{pkl:ss-pre}

Throughout, $L$ is a lattice with jump sets, $e_1,\dots,e_d$ a $\Z_\ell$-basis of $L$, and $\lambda_0\colon
\pkpW\to\pknB$ a fixed $\pkpA$-linear map with $\lambda_0(1)=1$ (it exists by the proof of
Lemma~\ref{pkn:frame-id}(iii)). For $k\le d$ let $\lambda_k$ be the $\pkpA$-linear map
$\lambda_k(\sum_sw_s\otimes e_s)=\lambda_0(w_k)$ on $\pkpAmb(L)=\bigoplus_s\pkpW\otimes e_s$; so
$\lambda_k(1\otimes e_s)=\delta_{ks}$. Recall that $\pkpA\to\pknBp$ is injective (\S\ref{pkn:ss-frame}),
so nonzero elements of $\pkpA$ are invertible in the field $\pknB$; and that an element of $\pknB$ fixed by
$j(\Gamma')$ lies in $K'$ (proof of Lemma~\ref{pkn:frame-id}(i)).

\begin{lemma}\label{pkl:basic}
\begin{enumerate}
\item[(i)] An element $x\in\pkpA$ maps to a unit of $\pknBp$ iff $\theta(x)\ne0$; for $r\in\pkpA$ and $n\ge0$,
$r\in\mathrm{Fil}^n$ iff $E^n\mid r$ in $\pkpA$. The images of $p_c$ ($c\ge1$), of $\varphi^{-b}(\tau)$ and
$\varphi^{-b}(z)$ ($b\ge0$) are units of $\pknBp$; for $b\ge1$, $\varphi^b(\tau)\in\mathrm{Fil}^1$. For
$l\ge0$, $\varphi^l(\tau)=\tau\prod_{b<l}p_b\cdot w_l$ with $w_l\in\pkpA^\times$.
\item[(ii)] If $x\in\pkpAmb(L)$ has $\tau^nx\in\pklL$ and $x\in\pknMinf(L)$, then $E^n\varphi(x)\in
\pknMinf(L)$. Hence $\varphi$ maps $\pknMinf(L)$ and $\pklAN(L)$ into $\pklAN(L)$.
\item[(iii)] For $x\in\pklAN(L)$, a form $\mu$ and $m\ge1$: $\mu(\varphi^mx)\in\pknBp$.
\item[(iv)] $\pknMinf(L)$, $\pklAN(L)$ and $\mathcal N(L)$ are stable under $1\otimes\zeta$ and under the
projectors $e_i$.
\item[(v)] (Shift.) If $x\in L_{i^{(-l)}}(o)$, then $\varphi^lx\in L_i(\mathrm{sh}_{l,0}(o))$.
\end{enumerate}
\end{lemma}

\begin{proof}
(i) $\pknBp$ is local with maximal ideal $\mathrm{Fil}^1=\ker\theta$ (Lemma~\ref{pkn:BdR}). If $r\in\pkpA\cap
\mathrm{Fil}^1$, then $\theta(r)=0$ and $r=Er'$, $r'\in\pkpA$ (Corollary~\ref{pkn:E-gen}); as $E$ is a uniformiser,
$r\in\mathrm{Fil}^n$ forces $r'\in\mathrm{Fil}^{n-1}$; induct. $\theta(p_c)\ne0$ for $c\ne0$
(Corollary~\ref{pkn:E-gen}); $\theta(\varphi^{-b}\tau)=\zeta_{\ell^{b+1}}-1\ne0$ and $\theta(\varphi^{-b}z)$ is an
$\ell^b$-th root of $\pi\ne0$, as $\theta([c])=c^\sharp$ and $\varphi^{-b}[c]=[c^{1/\ell^b}]$;
$\theta(\varphi^b\tau)=\zeta_\ell^{\ell^b}-1=0$ for $b\ge1$. Finally, with $\xi_\varepsilon=Eu$ as in
Lemma~\ref{pkn:mu}, $\varphi^{b+1}(\tau)=\varphi^b(\tau\xi_\varepsilon)=\varphi^b(\tau)p_b\varphi^b(u)$; induct
on $l$.

(ii) $\varphi(\tau)=\tau Eu$ with $u$ a unit (Lemma~\ref{pkn:mu}), and $\varphi(b\otimes y)=\varphi(b)\otimes y$
preserves $\pklL$; so $\tau^nE^n\varphi(x)=u^{-n}\varphi(\tau^nx)\in\pklL$. For $m\ge1$ and a form $\mu$,
$\mu(\varphi^m(E^n\varphi x))=\varphi^m(E)^n\mu(\varphi^{m+1}x)\in\pknBp$. So $E^n\varphi(x)\in\pknMinf$ and
$\varphi(x)\in\pklAN$ (with $r=p_0^n$). If $x\in\pklAN$ with $rx\in\pknMinf$, $r=\prod_kp_{b_k}$, then
$\varphi(r)\varphi(x)\in\pklAN$, $\varphi(r)=\prod_kp_{b_k+1}$, and $\pklAN$ is stable under $p_b^{-1}$
(Lemma~\ref{pkn:M}(iii)).

(iii) With $r$ as above, $\varphi^m(r)\mu(\varphi^mx)=\mu(\varphi^m(rx))\in\pknBp$ and $\varphi^m(r)=
\prod_kp_{b_k+m}$ is a unit of $\pknBp$ by (i).

(iv) $1\otimes\zeta$ is $\pkpW$-linear, commutes with $\varphi$ and with $\Gamma$ (the action on $L$ is
$\auO$-linear) and preserves $\pklL$; for a form $\mu$, $\mu\circ(1\otimes\zeta)$ is a form. This gives the
stability of the three modules (for $\mathcal N$ the value condition is preserved since
$\mu\circ(1\otimes\zeta)$ is a form). $e_i$ is a polynomial in $1\otimes\zeta$ with coefficients in
$\Lambda\subseteq\mathfrak S$, and the three modules are $\mathfrak S$-modules (Lemma~\ref{pkn:M}(i)).

(v) $\varphi^lx\in\pklAN$ by (ii) and $\varphi^lx\in\pkpAmb(L)_i$ (Definition~\ref{pkp:def-letters}). Let $a\ge0$,
$u\in U(i^{(-a)})$ and $\mu$ of type $u$ at $i^{(-a)}$; the value is $\mu(\varphi^{l-a}x)$. If $a<l$ it lies in
$\pknBp=\mathrm{Fil}^0$ by (iii). If $a\ge l$, then $i^{(-a)}=(i^{(-l)})^{(-(a-l))}$ and the value
$\mu(\varphi^{-(a-l)}x)$ lies in $\mathrm{Fil}^{o(a-l,u)}$.
\end{proof}

\subsubsection{Hodge--Tate jumps, the gap rule and the Hodge expansion}\label{pkl:ss-gap}

For $n\in\Z$ let $C(n)$ be $C$ with $g\in G_{K'}$ acting by $\chi(g)^ng$, $\chi$ the cyclotomic character;
$g(t)=\chi(g)t$, so $\mathrm{gr}^n\pknB\cong C(n)$. We assume $\auO[1/\ell]$ is a field, so that
$C\otimes_{\Z_\ell}\auO=\prod_\gamma C$ over the embeddings $\gamma\colon\auO\to C$, and
$C\otimes_{\Z_\ell}L=\bigoplus_\gamma C\otimes_{\auO,\gamma}L$, a $G_{K'}$-stable decomposition.

\emph{Reading of Definition~\ref{pkp:def-orders}.} ``$U(i)$ contains the Hodge--Tate jumps of $L$ at all
embeddings of letter $i$'' is used in the following form: \emph{if $\gamma$ is of letter $i$ and there is a
nonzero $C$-linear $G_{K'}$-equivariant map $C\otimes_{\auO,\gamma}L\to C(-u)$, then $u\in U(i)$.} By the
Hodge--Tate decomposition $C\otimes_{\auO,\gamma}L\cong\bigoplus_kC(-k)^{m_{\gamma,k}}$ of the de Rham
representation $L[1/\ell]$ and Tate's theorem $C(n)^{G_{K'}}=0$ for $n\ne0$ (Tate \cite[\S3, Theorem~2, p.~176]{pklTate67}; also quoted in \cite[1.5.7]{pkfFontaineII}), such a map exists iff $m_{\gamma,u}\ne0$; with the normalisation of (P3) of \S\ref{pkf:ss-rings}
($\pknBp\otimes V=\bigoplus_jt^{-k_j}\pknBp\delta'_j$, whence $C\otimes V\cong\bigoplus_jC(-k_j)$) the $k$
with $m_{\gamma,k}\ne0$ are the jumps $k$ of (P3). This is the convention of Theorem~\ref{pkf:main}.

\begin{lemma}[Gap]\label{pkl:gap}
Let $j$ be a letter, $u\in\Z\smallsetminus U(j)$ and $\mu$ a form of type $u$ at $j$. Then $\mu$ is of
type $u-1$ at $j$.
\end{lemma}

\begin{proof}
For $y\in L$, $t^u\mu(e_j(1\otimes y))\in\pknBp$; put $f(y)=\theta(t^u\mu(e_j(1\otimes y)))$, a $\Z_\ell$-linear
map $L\to C$. For $g\in\Gamma'$: $e_j(1\otimes gy)=g(e_j(1\otimes y))$ ($g$ fixes $\Lambda$ and commutes with
$1\otimes\zeta$), and $\mu(gw)=g\mu(w)$ for $w\in\pklL$ ($\mu(g(b\otimes y))=g(b)\,g\mu(1\otimes y)$); so
$t^u\mu(e_j(1\otimes gy))=\chi(g)^{-u}g\bigl(t^u\mu(e_j(1\otimes y))\bigr)$ and $f(gy)=\chi(g)^{-u}gf(y)$. Also
$(1\otimes\zeta)e_j=[j]e_j$, so $f(\zeta y)=\theta([j])f(y)$. The $C$-linear extension $F(c\otimes y)=cf(y)$
of $f$ to $C\otimes_{\Z_\ell}L$ is $G_{K'}$-equivariant into $C(-u)$ (here $G_{K'}=j(\Gamma')$: $j(\Gamma')$ is compact, hence closed, and $K'$ is its fixed field) and satisfies $F(\zeta w)=\theta([j])F(w)$.
On the summand $C\otimes_{\auO,\gamma}L$, $\zeta$ acts by $\gamma(\zeta)$; if $\gamma(\zeta)\ne\theta([j])$, then
$(\gamma(\zeta)-\theta([j]))F(w)=0$ and $F$ vanishes there. If $\gamma(\zeta)=\theta([j])$, $\gamma$ is of letter
$j$ (Definition~\ref{pkp:def-letters}) and $F$ vanishes there because $u\notin U(j)$. So $f=0$, i.e.\
$t^u\mu(e_j(1\otimes y))\in\ker\theta=\mathrm{Fil}^1$ and $\mu(e_j(1\otimes y))\in\mathrm{Fil}^{1-u}$.
\end{proof}

\begin{lemma}[Gap rule]\label{pkl:gap-rule}
Let $x\in L_i(o)$, $a\ge0$, $j=i^{(-a)}$, $\nu$ a form of type $\kappa$ at $j$, and $n\in\Z$ with
$o(a,v)\ge n$ for every $v\in U(j)$ with $v\le\kappa$. Then $\nu(\varphi^{-a}x)\in\mathrm{Fil}^n$.
\end{lemma}

\begin{proof}
Suppose not. Then $\nu$ is of type $v$ at $j$ for every $v\le\kappa$: by downward induction, if $\nu$ is of type
$v\le\kappa$ and $v\in U(j)$, then $\nu(\varphi^{-a}x)\in\mathrm{Fil}^{o(a,v)}\subseteq\mathrm{Fil}^n$, which is
excluded; so $v\notin U(j)$ and Lemma~\ref{pkl:gap} gives type $v-1$. Take $r\ne0$ with $rx\in\pklL$
(Lemma~\ref{pkn:M}(ii)). Then $w=\varphi^{-a}(rx)\in\pklL\cap\pkpAmb(L)_j$; writing $w=\sum b_k\otimes y_k$,
$w=e_jw=\sum b_ke_j(1\otimes y_k)$, so $\nu(w)\in\mathrm{Fil}^{-v}$ for all $v\le\kappa$, and $\nu(w)\in
\bigcap_m\mathrm{Fil}^m=0$. As $\nu(w)=\varphi^{-a}(r)\nu(\varphi^{-a}x)$ with $\varphi^{-a}(r)\ne0$ in the field
$\pknB$, $\nu(\varphi^{-a}x)=0\in\mathrm{Fil}^n$, a contradiction.
\end{proof}

\begin{lemma}[Hodge expansion]\label{pkl:expand}
There are forms $\nu_1,\dots,\nu_d$ and integers $k_1,\dots,k_d$ with:
\begin{enumerate}
\item[(i)] $\nu_m(\pklL)\subseteq\mathrm{Fil}^{-k_m}$; in particular $\nu_m$ is of type $k_m$ at every letter;
\item[(ii)] for every $\pkpA$-linear $\lambda\colon\pkpAmb(L)\to\pknB$ there are $D_m(\lambda)\in\pknB$ with
$\lambda(x)=\sum_mD_m(\lambda)\nu_m(x)$ for every $x\in\pkpAmb(L)$ such that $rx\in\pklL$ for some
$r\in\pkpA\smallsetminus\{0\}$ (e.g.\ $x\in\pklAN(L)$); if $\lambda(1\otimes L)\subseteq\mathrm{Fil}^0$ then
$D_m(\lambda)\in\mathrm{Fil}^{k_m}$, and if $\lambda$ is a form then $D_m(\lambda)\in K'$.
\end{enumerate}
\end{lemma}

\begin{proof}
Let $M$ be the matrix of a frame (Theorem~\ref{pkf:main}) and $N$ its image in $GL_d(\pknB)$; $N$ has invariant
columns by (Fr1). Lemma~\ref{pkf:hodge} gives $H=NY=W\operatorname{diag}(t^{k_1},\dots,t^{k_d})$ with $Y\in
SL_d(K')$, $W\in GL_d(\pknBp)$; $H$ has invariant columns ($Y$ is fixed by $G_{K'}$), so $g(H^{-1})=
H^{-1}\rho(g)$. Put $\nu_m(\sum_sw_s\otimes e_s)=\sum_s\lambda_0(w_s)(H^{-1})_{ms}$; as in the proof of
Lemma~\ref{pkn:frame-id}(iii), $\nu_m$ is a form, and $\nu_m(1\otimes e_s)=(H^{-1})_{ms}=t^{-k_m}(W^{-1})_{ms}\in
\mathrm{Fil}^{-k_m}$. By $\pkpA$-linearity $\nu_m(\pklL)\subseteq\mathrm{Fil}^{-k_m}$, and $e_j(1\otimes L)\subseteq
\pklL$. (ii) Put $c_s=\lambda(1\otimes e_s)$ and $D_m(\lambda)=\sum_sc_sH_{sm}$. Then
$\sum_mD_m(\lambda)\nu_m(1\otimes e_s)=(cHH^{-1})_s=c_s$, so both sides of the identity agree on $\pklL$ by
linearity, and on $x$ with $rx\in\pklL$ after multiplication by $r$, which is invertible in $\pknB$. Since
$H_{sm}=W_{sm}t^{k_m}\in\mathrm{Fil}^{k_m}$, $c_s\in\mathrm{Fil}^0$ gives $D_m(\lambda)\in\mathrm{Fil}^{k_m}$.
If $\lambda$ is a form, $g(c)=c\rho(g)$ for $g\in\Gamma'$ (proof of Lemma~\ref{pkn:frame-id}(i)) and
$g(H)=\rho(g)^{-1}H$, so $g(cH)=cH$ and $D_m(\lambda)\in K'$.
\end{proof}

\subsubsection{The standard orders; items (F) and (H)}\label{pkl:ss-std}

\begin{proposition}\label{pkl:std}
For $l\ge0$ and a letter $i$: $L_i(\pklstd_l)=\varphi^l\bigl(\pknMinf(L)\cap\pkpAmb(L)_{i^{(-l)}}\bigr)
=\varphi^l(\pknMinf(L))\cap\pkpAmb(L)_i$.
\end{proposition}

\begin{proof}
The two right sides agree because $\varphi^l$ is bijective on $\pkpAmb(L)$ and maps $\pkpAmb(L)_{i^{(-l)}}$ onto
$\pkpAmb(L)_i$.

$\supseteq$. Let $x\in\pknMinf\cap\pkpAmb(L)_{i^{(-l)}}$, $\tau^nx\in\pklL$. We show $x\in L_{i^{(-l)}}(o)$ with
$o(b,u)=-u$; then $\varphi^lx\in L_i(\mathrm{sh}_{l,0}(o))=L_i(\pklstd_l)$ by Lemma~\ref{pkl:basic}(v). Clearly
$x\in\pklAN$. Let $b\ge0$, $j=i^{(-l-b)}$, $u\in U(j)$, $\mu$ of type $u$ at $j$. The vector
$w=\varphi^{-b}(\tau^nx)$ lies in $\pklL\cap\pkpAmb(L)_j$, so $w=e_jw$ is an $\pkpA$-combination of elements of
$e_j(1\otimes L)$ and $\mu(w)\in\mathrm{Fil}^{-u}$. Since $\mu(w)=\varphi^{-b}(\tau)^n\mu(\varphi^{-b}x)$ and
$\varphi^{-b}(\tau)$ is a unit of $\pknBp$ (Lemma~\ref{pkl:basic}(i)), $\mu(\varphi^{-b}x)\in\mathrm{Fil}^{-u}$.

$\subseteq$. Let $y\in L_i(\pklstd_l)$ and $x=\varphi^{-l}y\in\pkpAmb(L)_{i^{(-l)}}$. We show $x\in\pknMinf$.

\emph{Values.} Let $m\ge1$, $\mu$ a form, and $z=\varphi^my$. By Lemma~\ref{pkl:basic}(v) (at the letter
$\sigma^m(i)$, with $\sigma^m(i)^{(-m)}=i$), $z\in L_{\sigma^m(i)}(o')$ with $o'=\mathrm{sh}_{m,0}(\pklstd_l)$,
and $o'(l,v)=0$ for all $v$: if $l<m$ by definition, and if $l\ge m$ because $o'(l,v)=\pklstd_l(l-m,v)$ with
$l-m<l$. By Lemma~\ref{pkl:expand}, $\mu(\varphi^mx)=\mu(\varphi^{-l}z)=\sum_kD_k(\mu)\nu_k(\varphi^{-l}z)$ (note
$r\varphi^{-l}z\in\pklL$ for $r=\varphi^{-l}(r')$, $r'z\in\pklL$), with $D_k(\mu)\in K'\subseteq\mathrm{Fil}^0$, and
$\nu_k(\varphi^{-l}z)\in\mathrm{Fil}^0$ by Lemma~\ref{pkl:gap-rule} (with $a=l$, $n=0$, $\kappa=k_k$). So
$\mu(\varphi^mx)\in\pknBp$.

\emph{Denominators.} Put $\tau_l=\varphi^l(\tau)$. We show: \emph{every $y\in L_i(\pklstd_l)$ has $\tau_l^Ny\in
\pklL$ for some $N$}; applying $\varphi^{-l}$ then gives $\tau^Nx\in\pklL$. First a sublemma: \emph{if $y\in
L_i(\pklstd_l)$, $a\ge l$, $s\in\pkpA$ and $sp_ay\in\pklL$, then $sy\in\pklL$.} Write $sp_ay=\sum_kc_k\otimes e_k$.
By Lemma~\ref{pkl:expand} applied to $\lambda_k$ (with $\lambda_k(1\otimes L)\subseteq\Z_\ell$), $\lambda_k(\varphi^{-a}
y)=\sum_mD_m(\lambda_k)\nu_m(\varphi^{-a}y)$ with $D_m(\lambda_k)\in\mathrm{Fil}^{k_m}$, and $\nu_m(\varphi^{-a}y)\in
\mathrm{Fil}^{-k_m}$ by Lemma~\ref{pkl:gap-rule} (as $\pklstd_l(a,v)=-v\ge-k_m$ for $v\le k_m$); so $\lambda_k
(\varphi^{-a}y)\in\mathrm{Fil}^0$. Then $\varphi^{-a}(c_k)=\lambda_k(\varphi^{-a}(sp_ay))=\varphi^{-a}(s)\,E\,
\lambda_k(\varphi^{-a}y)\in\mathrm{Fil}^1$, so $\theta(\varphi^{-a}c_k)=0$ and $p_a\mid c_k$ (Corollary~\ref{pkn:E-gen});
cancelling the unit $p_a$ of $\pkpW$, $sy\in\pklL$. Now let $r=p_{a_1}\cdots p_{a_k}$ with $ry\in\pknMinf$; we
induct on $k$ (for all $y$ at once). If $k=0$, $\tau^Ny\in\pklL$ for some $N$ and $\tau\mid\tau_l$
(Lemma~\ref{pkl:basic}(i)). If $k\ge1$, $y'=p_{a_1}y$ has orders $\pklstd_l+[a=a_1]\ge\pklstd_l$
(Lemma~\ref{pkn:E}), so $y'\in L_i(\pklstd_l)$ (Lemma~\ref{pkm:le}(i)) with $k-1$ factors, and
$\tau_l^Np_{a_1}y\in\pklL$. If $a_1<l$, then $p_{a_1}\mid\tau_l$ (Lemma~\ref{pkl:basic}(i)) and
$\tau_l^{N+1}y\in\pklL$; if $a_1\ge l$, the sublemma with $s=\tau_l^N$ gives $\tau_l^Ny\in\pklL$.
\end{proof}

\begin{corollary}[Lemma~\ref{pkm:inputs}(F)]\label{pkl:F}
$L_i(\pklstd_0)=\pknMinf(L)\cap\pkpAmb(L)_i$, and for $l\ge0$, $\varphi^l$ maps $L_i(\pklstd_0)$ bijectively onto
$L_{\sigma^l(i)}(\pklstd_l)$.
\end{corollary}

\begin{proof}
The first assertion is Proposition~\ref{pkl:std} with $l=0$. By Proposition~\ref{pkl:std} at the letter
$\sigma^l(i)$, whose $l$ steps back is $i$, $L_{\sigma^l(i)}(\pklstd_l)=\varphi^l(L_i(\pklstd_0))$, and $\varphi^l$
is injective.
\end{proof}

\begin{corollary}[Lemma~\ref{pkm:inputs}(H), first assertion]\label{pkl:H1}
With $a$ the jump bound of $L$: $\tau^a(1\otimes y)\in\pknMinf(L)\subseteq\pklAN(L)$ and
$e_i\tau^a(1\otimes y)\in L_i(\pklstd_0)$ for all $y\in L$ and letters $i$.
\end{corollary}

\begin{proof}
The first part is Proposition~\ref{pkn:frame}(v) with Theorem~\ref{pkf:main}. Then $e_i\tau^a(1\otimes y)\in
\pknMinf\cap\pkpAmb(L)_i$ by Lemma~\ref{pkl:basic}(iv), which is $L_i(\pklstd_0)$ by Corollary~\ref{pkl:F}.
\end{proof}

\emph{Remark.} With Definition~\ref{pkn:def-N}, $\mathcal N(L)$ is stable under $p_b^{-1}$
(Lemma~\ref{pkn:M}(iii)). When $\Gamma$ acts trivially on $L\ne0$, $\mathfrak S\otimes L\subseteq\mathcal N(L)$
by Proposition~\ref{pkl:H2}(ii), and then the inclusion is strict ($E^{-1}(1\otimes y)\in\mathcal N(L)$ for
$0\ne y\in L$, but $E^{-1}\notin\pkpA$); this is why (H) is stated with the intersection
$\mathcal N(L)\cap(\pklL)$.

\begin{proposition}[Lemma~\ref{pkm:inputs}(H), trivial action]\label{pkl:H2}
Suppose $\Gamma$ acts trivially on $L$. Then:
\begin{enumerate}
\item[(i)] $\pknMinf(L)=\pklL$ and $L_i(\pklstd_0)=\pklL\cap\pkpAmb(L)_i\subseteq\pklL$;
\item[(ii)] $\mathfrak S\otimes L\subseteq\mathcal N(L)$, and $e_i(1\otimes y)\in\mathcal N_i(\pklstd_l)$ for all
$y\in L$, letters $i$ and $l\ge0$.
\end{enumerate}
\end{proposition}

\begin{proof}
For a form $\mu$ and $y\in L$, $g\mu(1\otimes y)=\mu(1\otimes gy)=\mu(1\otimes y)$, so $\mu(1\otimes y)\in K'$.

(i) $\supseteq$: for $x=\sum_sb_s\otimes e_s$ and $m\ge1$, $\mu(\varphi^mx)=\sum_s\varphi^m(b_s)\mu(1\otimes e_s)\in
\pknBp$. $\subseteq$: let $x=\sum_sw_s\otimes e_s\in\pknMinf$ with $\tau^nx\in\pklL$, i.e.\ $\tau^nw_s\in\pkpA$. Each
$\lambda_k$ is a form ($\lambda_k(1\otimes gy)=\lambda_k(1\otimes y)\in\Z_\ell$). Hence, for $m\ge1$,
$\lambda_k(\varphi^mx)=\lambda_0(\varphi^mw_k)\in\pknBp$, and $\varphi^m(\tau)^n\lambda_0(\varphi^mw_k)=
\varphi^m(\tau^nw_k)$. We show by induction on $n$: \emph{if $w\in\pkpW$, $\tau^nw\in\pkpA$ and $q_m:=
\varphi^m(\tau^nw)/\varphi^m(\tau)^n\in\pknBp$ for all $m\ge1$, then $w\in\pkpA$.} For $n=0$ there is nothing
to prove. For $n\ge1$, $t=\tau^nw\in\pkpA$ has $\varphi^m(t)=\varphi^m(\tau)^nq_m\in\mathrm{Fil}^1$ for $m\ge1$
(Lemma~\ref{pkl:basic}(i)), i.e.\ $\theta(\varphi^j(\varphi t))=0$ for all $j\ge0$; by Lemma~\ref{pkn:mu},
$\varphi(\tau)\mid\varphi(t)$, so $t=\tau t'$ with $t'=\tau^{n-1}w\in\pkpA$, and the quotients for $(w,n-1)$ are the
same $q_m$. So $w_k\in\pkpA$ and $x\in\pklL$. The second assertion follows from Corollary~\ref{pkl:F}.

(ii) $s\otimes y\in\pklL=\pknMinf\subseteq\pklAN$ and $\mu(s\otimes y)=s\mu(1\otimes y)\in\pknK$ ($\pknK\supseteq
K'$ and the image of $\mathfrak S$). Next, $e_i(1\otimes y)\in\Lambda\otimes L\subseteq\mathfrak S\otimes L$ lies in
$\mathcal N\cap\pkpAmb(L)_i$. For $a\ge0$, $\varphi^{-a}e_i(1\otimes y)=e_{i^{(-a)}}(1\otimes y)$ ($\varphi^{-a}[i']=
[i'^{1/\ell^a}]$ and $i'\mapsto i'^{1/\ell^a}$ permutes the letters). For a form $\mu$ its value lies in
$\pkpA\cdot K'\subseteq\mathrm{Fil}^0$, and in $\mathrm{Fil}^{-u}$ if $\mu$ is of type $u$ at $i^{(-a)}$ (by
definition). As $\pklstd_l(a,u)\in\{0,-u\}$, $e_i(1\otimes y)\in L_i(\pklstd_l)$.
\end{proof}

\subsubsection{Galois stability and the level-one gain (D)}\label{pkl:ss-D}

\emph{Remark.} The $\Gamma'$-stability of $\mathcal N(L)$ itself is neither claimed nor needed: $g\in\Gamma'$ does
not fix $\mathfrak S$, and the field $\pknK$ need not be $g$-stable.

\begin{lemma}[Lemma~\ref{pkm:inputs}(G), stability]\label{pkl:G1}
For $g\in\Gamma'$: $g\pknMinf(L)\subseteq\pknMinf(L)$, $g\,\pklAN(L)\subseteq\pklAN(L)$ and $gL_i(o)\subseteq L_i(o)$
for all letters $i$ and orders $o$.
\end{lemma}

\begin{proof}
$g$ preserves $\pklL$ ($g(b\otimes y)=g(b)\otimes gy$) and $g\tau=\tau r'$ with $r'\in\pkpA^\times$
(Lemma~\ref{pkf:approx0}(iv)); so $\tau^nx\in\pklL$ implies $\tau^ngx=r'^{-n}g(\tau^nx)\in\pklL$. For a form $\mu$
the map $\mu^g(w)=g^{-1}\mu(gw)$ is $\pkpA$-linear and is a form: $\mu^g(1\otimes hy)=g^{-1}\mu(1\otimes
(ghg^{-1})gy)=h\mu^g(1\otimes y)$ for $h\in\Gamma'$. As $g$ commutes with $\varphi$, $\mu(\varphi^mgx)=
g\mu^g(\varphi^mx)\in g\pknBp=\pknBp$ for $x\in\pknMinf$, $m\ge1$. So $g\pknMinf\subseteq\pknMinf$. If
$rx\in\pknMinf$ with $r=\prod_kp_{b_k}$, then $g(r)gx\in\pknMinf$ and $g(p_b)=\varphi^b(gE)=p_b\varphi^b(r_E)$ with
$r_E$ a unit (Lemma~\ref{pkf:approx0}(iv)), so $rgx\in\pknMinf$. Finally let $x\in L_i(o)$: $gx\in\pklAN$, $gx\in
\pkpAmb(L)_i$ ($g$ commutes with $1\otimes\zeta$ and fixes $[i]$), and by Lemma~\ref{pkm:equiv}
$\mu(\varphi^{-a}gx)=g\mu(\varphi^{-a}x)\in g\mathrm{Fil}^n=\mathrm{Fil}^n$ ($g(t)=\chi(g)t$).
\end{proof}

The next lemma transports relations through $\mathfrak S$-linear maps of $\pkpA$; it rests on input (A8)(c),
proved as Proposition~\ref{pkr:I8}(c).

\begin{lemma}[Retraction of values]\label{pkl:valret}
Let $f\colon\pkpA\to\pkpA$ be $\mathfrak S$-linear, $v_1,\dots,v_n\in\pknK$, $r_1,\dots,r_n\in\pkpA$, and $W\subseteq
\pknB$ a $\pknBp$-submodule. If $\sum_lr_lv_l\in W$, then $\sum_lf(r_l)v_l\in W$.
\end{lemma}

\begin{proof}
Let $R\subseteq\pknB$ be the $K'$-span of the image of $\pkfSB$; it is a ring ($K'$ consists of constants and
$\pkfSB$ is a ring) containing the image of $\mathfrak S\subseteq\pkfSB$. The generators $bs$ ($b\in\pkfSB$,
$s\in\mathfrak S$) of $\pknK$ lie in the image of $\pkfSB$, and $s=1$ gives all of it; so $\pknK$ is the field of
fractions of the domain $R$ inside $\pknB$. Choose $d\in R\smallsetminus\{0\}$ with $dv_l=w_l\in R$, and write
$w_l=\sum_ic_{li}b_{li}$ (finite sums, $c_{li}\in K'$, $b_{li}\in\pkfSB$). By Proposition~\ref{pkr:I8}(c), applied to
the coefficients $c_{li}$ and the elements $b_{li}r_l\in\pkfBp$, there is $u\in\pknBp$ with $\sum_{l,i}c_{li}
G_f(b_{li}r_l)=u\sum_{l,i}c_{li}b_{li}r_l$, and $G_f(b_{li}r_l)=b_{li}f(r_l)$ by Lemma~\ref{pkf:extend}. Hence
$\sum_lw_lf(r_l)=u\sum_lw_lr_l$, and dividing by $d$, $\sum_lf(r_l)v_l=u\sum_lr_lv_l\in W$.
\end{proof}

\begin{proposition}[Lemma~\ref{pkm:inputs}(D)]\label{pkl:D}
If $x\in\mathcal N_i(\pklstd_1)$ and $g\in\Gamma'$, then $gx-x\in L_i(o_1)$, where $o_1(0,u)=1$ and $o_1(a,u)=-u$
for $a\ge1$.
\end{proposition}

\begin{proof}
By Lemma~\ref{pkl:G1}, $gx-x\in L_i(\pklstd_1)$, and $\pklstd_1(a,u)=o_1(a,u)$ for $a\ge1$; it remains to treat
$a=0$. Let $u\in U(i)$ and $\mu$ a form of type $u$ at $i$, and put $y=\mu(x)$. Then $y\in\pknK$ ($x\in\mathcal N$)
and $y\in\mathrm{Fil}^{\pklstd_1(0,u)}=\pknBp$. As $\theta\colon\pkpA\to O_C$ is onto, there are $e\ge0$ and
$r\in\pkpA$ with $\theta(r)=\ell^e\theta(y)$, i.e.\ $\ell^ey-r\in\mathrm{Fil}^1$. Lemma~\ref{pkl:valret} with
$f=\kappa_0$ (the coordinate of (A5), Proposition~\ref{pkr:I5}, $\mathfrak S$-linear with values in $\mathfrak S$
and $\kappa_0(\ell^e)=\ell^e$), the values $(y,-1)$ and the coefficients $(\ell^e,r)$, gives $\ell^ey-s\in
\mathrm{Fil}^1$ with $s=\kappa_0(r)\in\mathfrak S$. Now $g$ preserves $\mathrm{Fil}^1$, and $gs-s\in z\tau E\pkpA
\subseteq\mathrm{Fil}^1$ (Proposition~\ref{pkn:Ca}). Hence $\ell^e(gy-y)=g(\ell^ey-s)-(\ell^ey-s)+(gs-s)\in
\mathrm{Fil}^1$, and $gy-y\in\mathrm{Fil}^1$. By Lemma~\ref{pkm:equiv}, $\mu(gx-x)=gy-y$.
\end{proof}

\subsubsection{The depth step (C)(d)}\label{pkl:ss-Cd}

\begin{lemma}\label{pkl:different}
Let $P\in\Lambda[X]$ be the distinguished polynomial with $E=P(z)V(z)$, $V\in\Lambda[[X]]^\times$ (Weierstrass
preparation in $\Lambda[[X]]\cong\mathfrak S$). Then $P$ is Eisenstein of degree $e$, $P(\pi)=0$, and
$v_\ell(P'(\pi))\le d_v$.
\end{lemma}

\begin{proof}
Modulo $\ell$, $E\equiv z^eA$ with $A$ a unit, so $P\equiv X^e$ and $\deg P=e$; $P(0)=E(0)/V(0)=\ell/V(0)$ has
valuation $1$. So $P$ is Eisenstein, and $P(\pi)=\theta(P(z))=\theta(E)\theta(V(z))^{-1}=0$. If $e=1$, $P'=1$
and $d_v=0$. If $e>1$, write $P=\sum_{k=0}^ec_kX^k$ with $c_e=1$ and $v_\ell(c_k)\ge1$ for $k<e$. The terms
$kc_k\pi^{k-1}$ ($1\le k\le e$) of $P'(\pi)$ have valuations $v_\ell(k)+v_\ell(c_k)+(k-1)/e$, pairwise distinct
(the fractional parts $(k-1)/e$ are distinct), so $v_\ell(P'(\pi))$ is their minimum, at most
$v_\ell(e)+(e-1)/e<v_\ell(e)+1=d_v$.
\end{proof}

\begin{proposition}[Lemma~\ref{pkm:inputs}(C)(d)]\label{pkl:Cd}
For every $n\ge1$ there are $g\in\Gamma'$, $\omega\in\pkpA^\times$ and $\upsilon,q\in\pkpA$ with $gE=\omega E$,
$E\mid\upsilon q-\ell^{v_\ell(n)+d_v}$ and $z\tau q=\omega^n-1$.
\end{proposition}

\begin{proof}
Let $g=t_1$ with $gz=[\varepsilon]z$ ((C)(b)), and $gE=\omega E$ with $\omega\in\pkpA^\times$
(Lemma~\ref{pkf:approx0}(iv)). Write $\delta=z\tau$; $\theta(\delta)=\pi(\zeta_\ell-1)\ne0$.

\emph{Step 1: $\delta\mid\omega-1$, hence $\delta\mid\omega^n-1=\delta q$ for a unique $q\in\pkpA$.} By
Proposition~\ref{pkn:Ca}, $\delta E\mid gE-E=E(\omega-1)$; cancel $E$.

\emph{Step 2: $\theta(\omega-1)\mid\ell^{d_v}\theta(\delta)$ in $O_C$.} With $P,V$ as in Lemma~\ref{pkl:different},
$P(z)=EG$ with $G=V(z)^{-1}\in\mathfrak S$. By Lemma~\ref{pkn:mu}, $gz-z=z([\varepsilon]-1)=z\tau Eu=Eh$ with
$h=\delta u$, $u\in\pkpA^\times$. Taylor's formula for the polynomial $P$ gives $c\in\pkpA$ with
$P(gz)=P(z)+P'(z)Eh+E^2h^2c$. As $g$ fixes $\Lambda$, $P(gz)=g(P(z))=g(E)g(G)=\omega Eg(G)$; cancelling $E$,
$\omega g(G)=G+P'(z)h+Eh^2c$. Apply $\theta$: $\theta(g(G))=g\theta(G)=\theta(G)$, because $\theta(G)$ is a
convergent series in $\pi$ with coefficients in $\theta(\Lambda)$, fixed by $j(\Gamma')$. So
$\theta(\omega-1)\theta(G)=P'(\pi)\theta(\delta)\theta(u)$, and $\theta(G)=\theta(V(z))^{-1}$ is a unit of $O_C$
($V(z)\in\mathfrak S^\times$). By Lemma~\ref{pkl:different}, $\theta(\omega-1)\mid\ell^{d_v}\theta(\delta)$; in particular
$\theta(\omega-1)\ne0$.

\emph{Step 3: $v_\ell(\theta(\omega^n-1))=v_\ell(n)+v_\ell(\theta(\omega-1))$.} For $r=\delta b\in\delta\pkpA$,
$\ell\,\theta(r)\theta(z)\mid\theta(r)^\ell$: indeed $y=\theta(\tau)=\zeta_\ell-1$ satisfies $(1+y)^\ell=1$, so
$y^\ell=-\sum_{0<k<\ell}\binom\ell ky^k$ is divisible by $\ell y$; and $\theta(z)^2\mid\theta(z)^\ell$,
$\theta(b)\mid\theta(b)^\ell$. Let $x_0=\theta(\omega)$ and $y_0=x_0-1=\theta(\omega-1)\ne0$, of positive valuation
($\theta(z)\mid y_0$). Write $n=m\ell^v$ with $\ell\nmid m$. First $x_0^m-1=y_0(1+x_0+\dots+x_0^{m-1})$ and the
second factor is $\equiv m\not\equiv0$ modulo the maximal ideal, so $v_\ell(x_0^m-1)=v_\ell(y_0)$. Next, if
$x'=\theta(\omega')$ with $\omega'$ a power of $\omega$ and $y'=x'-1\ne0$, then $\delta\mid\omega'-1$ (Step 1), and
$x'^\ell-1=\ell y'+\sum_{1<k<\ell}\binom\ell ky'^k+y'^\ell$, where the middle terms have valuation $\ge1+2v_\ell(y')$
and the last $\ge1+v_\ell(y')+v_\ell(\pi)$ by the divisibility just shown; so $v_\ell(x'^\ell-1)=1+v_\ell(y')$.
Applying this $v$ times starting from $\omega'=\omega^m$ gives the claim.

\emph{Step 4.} By Steps 2 and 3, $\theta(\delta)\theta(q)=\theta(\omega^n-1)$ divides $\ell^{v_\ell(n)}\theta(\omega-1)$,
which divides $\ell^{v_\ell(n)+d_v}\theta(\delta)$. Cancelling $\theta(\delta)\ne0$ in the domain $O_C$,
$\ell^{v_\ell(n)+d_v}=\theta(q)c$ with $c\in O_C$. Choose $\upsilon\in\pkpA$ with $\theta(\upsilon)=c$; then
$\theta(\upsilon q-\ell^{v_\ell(n)+d_v})=0$ and $E\mid\upsilon q-\ell^{v_\ell(n)+d_v}$ (Corollary~\ref{pkn:E-gen}).
\end{proof}

\subsubsection{The retraction (R)}\label{pkl:ss-R}

\begin{lemma}[Relations]\label{pkl:rel}
Let $v_1,\dots,v_n\in\pklAN(L)$ with $\mu(v_i)\in\pknK$ for every form $\mu$ (e.g.\ $v_i\in\mathcal N(L)$), let
$c_1,\dots,c_n\in\pkpA$ with $\sum_ic_iv_i=0$, and let $f\colon\pkpA\to\pkpA$ be $\mathfrak S$-linear. Then
$\sum_if(c_i)v_i=0$.
\end{lemma}

\begin{proof}
Fix a frame (Theorem~\ref{pkf:main}). Put $w=\sum_if(c_i)v_i\in\pklAN$ and choose $r$, a product of a power of
$\tau$ and of elements $p_b$, with $rw\in\pklL$ (Lemma~\ref{pkn:M}(ii)). As in the proof of
Proposition~\ref{pkn:frame}(ii), $\tau^arw=\sum_jt_jn_j$ with $t=Q'\alpha\in\pkpA^d$, where $rw=\sum_j\alpha_j
(1\otimes e_j)$. For every form $\mu$, $\sum_ic_i\mu(v_i)=\mu(0)=0$, so Lemma~\ref{pkl:valret} (with $W=0$) gives
$\mu(w)=\sum_if(c_i)\mu(v_i)=0$, hence $\mu(\sum_jt_jn_j)=\tau^ar\mu(w)=0$. By the coefficient test
(Lemma~\ref{pkn:frame-id}(iv) with $W=0$, $m=0$), $E^{2a}t_j=0$, so $t_j=0$, $\tau^arw=0$, and $w=0$ since
$\tau^ar$ is a unit of $\pkpW$.
\end{proof}

\begin{lemma}\label{pkl:phin}
For a frame of $L$ and $l\ge0$, $\varphi^l(n_j)\in\mathcal N(L)$ for every $j$.
\end{lemma}

\begin{proof}
$n_j\in\pknMinf$ (Proposition~\ref{pkn:frame}(i)), so $\varphi^l(n_j)\in\pklAN$ by Lemma~\ref{pkl:basic}(ii). For a form
$\mu$, Lemma~\ref{pkn:frame-id}(ii) with $m=l$ and $t$ the $j$-th unit vector gives $p_l^a\mu(\varphi^ln_j)=
c^{(l)}(\mu)\,\varphi^l(Z)_{\cdot j}$, with $c^{(l)}(\mu)$ over $K'$ (Lemma~\ref{pkn:frame-id}(i)) and the entries of
$\varphi^l(Z)$ in $\pkfSB$ (Lemma~\ref{pkf:values}); their images lie in $\pknK$, and $p_l$ is a nonzero element of
$\pknK$. So $\mu(\varphi^ln_j)\in\pknK$.
\end{proof}

\begin{proposition}[Lemma~\ref{pkm:inputs}(R)]\label{pkl:R}
Let $\kappa_0\colon\pkpA\to\mathfrak S$ be the coordinate of (A5) (Proposition~\ref{pkr:I5}); it plays the role of the
map $\lambda_0$ of Lemma~\ref{pkm:inputs}(R) (not the map $\lambda_0$ of \S\ref{pkl:ss-pre}), and it is
$\mathfrak S$-linear and the identity on $\mathfrak S$. For every lattice with jump sets $L$ there is a unique map
$\mathrm{ret}\colon\pklAN(L)\to\mathcal N(L)$ with $\mathrm{ret}(\sum_ib_ix_i)=\sum_i\kappa_0(b_i)x_i$ for $b_i\in\pkpA$,
$x_i\in\mathcal N(L)$. It is $\mathfrak S$-linear, is the identity on $\mathcal N(L)$, commutes with $1\otimes\zeta$ and
the $e_i$, maps $L_i(\pklstd_l)$ into $\mathcal N_i(\pklstd_l)$ for all $i$ and $l\ge0$, and, for an extension $X_S$,
$\mathrm{pr}\circ\mathrm{ret}=\mathrm{ret}\circ\mathrm{pr}$.
\end{proposition}

\begin{proof}
$\pklAN(L)$ is the $\pkpA$-span of $\mathcal N(L)$ (Proposition~\ref{pkn:frame}(iv) with Theorem~\ref{pkf:main}), so
every element has the form $\sum_ib_ix_i$. If $\sum_ib_ix_i=\sum_jb'_jx'_j$, Lemma~\ref{pkl:rel} applied to the
relation $\sum_ib_ix_i-\sum_jb'_jx'_j=0$ with $f=\kappa_0$ shows $\sum_i\kappa_0(b_i)x_i=\sum_j\kappa_0(b'_j)x'_j$; so
$\mathrm{ret}$ is well defined, with values in the $\mathfrak S$-module $\mathcal N$. It is $\mathfrak S$-linear because
$\kappa_0$ is, and $\mathrm{ret}(x)=\kappa_0(1)x=x$ for $x\in\mathcal N$. Since $1\otimes\zeta$ preserves $\mathcal N$
(Lemma~\ref{pkl:basic}(iv)), $(1\otimes\zeta)\mathrm{ret}(\sum b_ix_i)=\sum\kappa_0(b_i)(1\otimes\zeta)x_i=
\mathrm{ret}((1\otimes\zeta)\sum b_ix_i)$; as $e_i$ is a polynomial in $1\otimes\zeta$ with coefficients in
$\Lambda\subseteq\mathfrak S$, $\mathrm{ret}$ commutes with $e_i$.

Let $x\in L_i(\pklstd_l)$. By Proposition~\ref{pkl:std} and Proposition~\ref{pkn:frame}(ii), $x=\varphi^l(\sum_jt_jn_j)
=\sum_j\varphi^l(t_j)\varphi^l(n_j)$ with $t_j\in\pkpA$, and $\varphi^l(n_j)\in\mathcal N$ (Lemma~\ref{pkl:phin}); so
$\mathrm{ret}(x)=\sum_j\kappa_0(\varphi^lt_j)\varphi^l(n_j)=\varphi^l\bigl(\sum_j\varphi^{-l}(\kappa_0(\varphi^lt_j))n_j
\bigr)\in\varphi^l(\pknMinf)$. Also $\mathrm{ret}(x)=\mathrm{ret}(e_ix)=e_i\mathrm{ret}(x)\in\pkpAmb(L)_i$. By
Proposition~\ref{pkl:std}, $\mathrm{ret}(x)\in L_i(\pklstd_l)$, and it lies in $\mathcal N$.

For an extension $X_S$, $\mathrm{pr}$ is the $\pkpW$-linear extension of the equivariant $\auO$-linear projection
$X_S\to C_S$; it commutes with $\varphi$, maps $\pkpA\otimes X_S$ onto $\pkpA\otimes C_S$, and $\nu\circ\mathrm{pr}$ is a
form on $X_S$ for every form $\nu$ on $C_S$. Hence $\mathrm{pr}$ maps $\pknMinf(X_S)$, $\pklAN(X_S)$ and $\mathcal N(X_S)$
into the corresponding modules of $C_S$, and $\mathrm{pr}(\mathrm{ret}(\sum b_ix_i))=\sum\kappa_0(b_i)\mathrm{pr}(x_i)=
\mathrm{ret}(\sum b_i\mathrm{pr}(x_i))=\mathrm{ret}(\mathrm{pr}(\sum b_ix_i))$.
\end{proof}


\begin{corollary}[Frobenius on $\mathcal N$]\label{pkl:Phi}
$\varphi(\mathcal N(L))\subseteq\mathcal N(L)$.
\end{corollary}

\begin{proof}
Let $x\in\mathcal N$ and $r=\prod_kp_{b_k}$ with $rx\in\pknMinf$; $rx\in\mathcal N$ (Lemma~\ref{pkn:M}(iii)). By
Proposition~\ref{pkn:frame}(ii), $rx=\sum_jt_jn_j$ with $t_j\in\pkpA$; as $rx$ and the $n_j$ lie in $\mathcal N$
(Proposition~\ref{pkn:frame}(i)), $rx=\mathrm{ret}(rx)=\sum_j\kappa_0(t_j)n_j$, and by linear independence
$t_j=\kappa_0(t_j)\in\mathfrak S$. Then $\varphi(r)\varphi(x)=\sum_j\varphi(t_j)\varphi(n_j)$ with $\varphi(t_j)\in\mathfrak S$
and $\varphi(n_j)\in\mathcal N$ (Lemma~\ref{pkl:phin}), so $\varphi(r)\varphi(x)\in\mathcal N$; $\varphi(r)=\prod_kp_{b_k+1}$,
and $\mathcal N$ is stable under $p_b^{-1}$, so $\varphi(x)\in\mathcal N$.
\end{proof}

\begin{corollary}[(H), second assertion]\label{pkl:H3}
If $\Gamma$ acts trivially on $L$, then $\mathcal N(L)\cap\pklL=\mathfrak S\otimes L$, and $\mathcal N(L)$ is the set of
$x$ with $rx\in\mathfrak S\otimes L$ for some finite product $r$ of elements $p_b$. Moreover
$L_i(\pklstd_0)\subseteq\pklL$ (Proposition~\ref{pkl:H2}(i)).
\end{corollary}

\begin{proof}
$\supseteq$ in the first claim is Proposition~\ref{pkl:H2}(ii). If $x=\sum_sb_s(1\otimes e_s)\in\mathcal N$, then, as
$1\otimes e_s\in\mathcal N$, $x=\mathrm{ret}(x)=\sum_s\kappa_0(b_s)(1\otimes e_s)\in\mathfrak S\otimes L$. For the second
claim: if $x\in\mathcal N$, then $rx\in\pklL$ for some product $r$ of a power of $\tau$ and elements $p_b$
(Lemma~\ref{pkn:M}(ii)); by Proposition~\ref{pkl:H2}(i) and the definition of $\pklAN$ one may take $r$ a product of
$p_b$'s only ($rx\in\pknMinf=\pklL$ by definition of $\pklAN$), and $rx\in\mathcal N$ (Lemma~\ref{pkn:M}(iii)), so
$rx\in\mathfrak S\otimes L$. Conversely $\mathfrak S\otimes L\subseteq\mathcal N$ and $\mathcal N$ is stable under
$p_b^{-1}$.
\end{proof}

\subsubsection{The second assertion of (G)}\label{pkl:ss-open}

The second assertion of (G) (Lemma~\ref{pkp:tau}(iv)) reduces, by Lemma~\ref{pkl:basic}(i) and the argument
below, to a divisibility statement. Let $x\in\mathcal N(L)$ and $g\in\Gamma'$ with $gx-x\in L_i(o)$. If
$y=(z\tau)^{-1}(gx-x)$ lies in $\pklAN(L)$, then $y\in L_i(o)$: it lies at $i$, and for a form $\mu$ of type $u$
at $i^{(-a)}$, $\varphi^{-a}(z\tau)\mu(\varphi^{-a}y)=\mu(\varphi^{-a}(gx-x))\in\mathrm{Fil}^{o(a,u)}$ with
$\varphi^{-a}(z\tau)$ a unit of $\pknBp$ (Lemma~\ref{pkl:basic}(i)). Moreover, if $y\in\pklAN$ for every
$x\in\mathcal N\cap\pknMinf$, then $y\in\pklAN$ for every $x\in\mathcal N$: if $x\in\mathcal N$ and $p_bx$ has the
property, then $g(p_bx)-p_bx=p_b(gx-x)+(gp_b-p_b)gx$ with $gp_b-p_b\in z\tau E\pkpA$ (Proposition~\ref{pkn:Ca}), so
$p_b\,y=(z\tau)^{-1}(g(p_bx)-p_bx)-E\beta\,gx\in\pklAN$ for some $\beta\in\pkpA$, and $y\in\pklAN$
(Lemma~\ref{pkn:M}(iii), Lemma~\ref{pkl:G1}); one inducts on the number of factors $p_b$ needed to bring $x$ into
$\pknMinf$ (note $p_bx\in\mathcal N$). For $x\in\mathcal N\cap\pknMinf$ one has $x=\sum_jt_jn_j$ with $t_j\in\mathfrak S$
(apply $\mathrm{ret}$ and Proposition~\ref{pkn:frame}(i),(ii)), so that $E^a\sigma_\nu(x)=(Zt)_\nu\in\pkfSB$ for the
forms $\sigma_\nu$ of Lemma~\ref{pkn:frame-id}(iii) ($m=0$). Also, if $z\tau$ divides $E^hr$ in $\pkfBp$
($r\in\pkpA$), then $\tau\mid r$ in $\pkpA$: from $E^hr=\tau w$, $\theta(p_n)^h\theta(\varphi^nr)=
\theta(\varphi^n\tau)\theta(\varphi^nw)=0$ for $n\ge1$, so $\varphi(\tau)\mid\varphi(r)$ by Lemma~\ref{pkn:mu}.

We now prove the two facts about $\pkfBp$ needed to finish, and then (G0).

\begin{lemma}[Integral divisibility]\label{pkl:intdiv}
Let $\nu,\beta,x\in\pkpA$ with $\ell\nmid\nu$ and $\nu\mid E-\ell\beta$ in $\pkpA$. If $\nu$ divides $x$ in $\pkfBp$, then
$\nu\mid x$ in $\pkpA$. In particular this holds for $\nu=z$ ($\beta=1$, as $E-\ell=z^eA$).
\end{lemma}

\begin{proof}
Write $x=\nu w$ with $w\in\pkfBp$ and choose $k$ with $\ell^kw\in\pkfAo$; so $\ell^kx=\nu w_0$ in $\pkfAm$. Fix $n\ge1$
and let $Q=\pkpA/(\nu,\ell^n)$. The ring map $\pkpA[X]\to Q$, $X\mapsto\bar\beta$, kills $\ell X-E$ (as
$E-\ell\beta\in\nu\pkpA$), so it factors through $R=\pkpA[X]/(\ell X-E)$; since $\ell^n=0$ in $Q$, it factors through
$R/\ell^nR$, hence through $\pkfAm\to R/\ell^nR$ (Lemma~\ref{pkr:Amax}(i)). The resulting map $F\colon\pkfAm\to Q$
restricts to the projection on $\pkpA$, and $F(\ell^kx)=F(\nu)F(w_0)=0$. So $\ell^kx\in(\nu,\ell^n)\subseteq\nu\pkpA+
\pkrI^n$ for every $n$, and Proposition~\ref{pkr:I7} (with $\delta=\nu$, $c=\ell^kx$) gives $\nu\mid\ell^kx$. As $\ell$ is a
prime element of $\pkpA$ not dividing $\nu$ ((W1) of \S\ref{pkr:ss-prelim}), $\nu\mid x$. For $\nu=z$: $z_0=\varpi^\flat\ne0$,
so $\ell\nmid z$.
\end{proof}

\begin{lemma}\label{pkl:SBdec}
For $b\in\pkfSB$ and $M\ge0$ there are $k\ge0$, $s\in\mathfrak S$ and $w\in\pkfSB$ with $\ell^kb=s+E^Mw$, and there are
$k'\ge0$, $\lambda\in\Lambda$ and $w'\in\pkfSB$ with $\ell^{k'}b=\lambda+zw'$.
\end{lemma}

\begin{proof}
Let $\ell^jb=y\in\pkfAo$. Lemma~\ref{pkr:div} with $\nu=E$ ($b=0$), resp.\ $\nu=z$ ($b=1$; $z\mid E-\ell$), gives
$\ell^Ny=a+\nu^Nw_1$ with $a\in\pkpA$, $w_1\in\pkfAo$. Apply $G_0$ (Lemmas~\ref{pkf:extend}, \ref{pkf:values}): $G_0$ fixes
$\ell^Ny\in\pkfSB$, is $\pkfSB$-linear, and $G_0(w_1)\in\pkfSB$; so $\ell^{N+j}b=\kappa_0(a)+\nu^NG_0(w_1)$. With $\nu=E$,
$N=M$ this is the first claim. With $\nu=z$, $N=1$, write $\kappa_0(a)=\lambda+zs'$ ($\lambda\in\Lambda$, $s'\in\mathfrak S$)
and $w'=s'+G_0(w_1)$.
\end{proof}

\begin{lemma}[Galois action on $\pkfSB$]\label{pkl:galSB}
For $g\in\Gamma'$ and $b\in\pkfSB$, $z\tau$ divides $gb-b$ in $\pkfBp$.
\end{lemma}

\begin{proof}
\emph{$z$ divides.} By Lemma~\ref{pkl:SBdec}, $\ell^{k'}b=\lambda+zw'$. $g$ fixes $\lambda$, and $gz=z[\eta]$ with
$\eta\in O^{\flat\times}$ (proof of Proposition~\ref{pkn:Ca}); so $\ell^{k'}(gb-b)=z([\eta]g(w')-w')$, and $\ell$ is a unit of
$\pkfBp$. Write $gb-b=zc$, $c\in\pkfBp$.

\emph{$\tau$ divides $c$.} First, $\theta(gb'-b')=0$ for every $b'\in\pkfSB$: by Lemma~\ref{pkl:SBdec} with $M=1$,
$\ell^kb'=s+Ew$, and $\theta(g(\ell^kb'))=g\theta(s)+g\theta(E)\theta(gw)=\theta(s)=\theta(\ell^kb')$, since $\theta$ commutes with
$g$ on $\pkfBp$ (P1), $\theta(E)=0$, and $\theta(s)$, a convergent series in $\pi$ over $\theta(\Lambda)$, is fixed by
$j(\Gamma')$. For $n\ge0$, $\varphi^n(b)\in\pkfSB$ (Lemma~\ref{pkf:values}) and $g$ commutes with $\varphi$, so
$0=\theta(\varphi^n(gb-b))=\theta(\varphi^nz)\theta(\varphi^nc)$ with $\theta(\varphi^nz)=\pi^{\ell^n}\ne0$. Hence
$\theta(\varphi^n(\varphi c))=0$ for all $n\ge0$, $\mu=\varphi(\tau)$ divides $\varphi(c)$ (Lemma~\ref{pkr:mu-Bplus}), and
$\tau\mid c$ in $\pkfBp$ (Proposition~\ref{pkr:I8}(b)).
\end{proof}

\begin{lemma}[Cancellation]\label{pkl:cancel}
If $r\in\pkpA$, $h\ge0$ and $z\tau$ divides $E^hr$ in $\pkfBp$, then $z\tau\mid r$ in $\pkpA$.
\end{lemma}

\begin{proof}
By the argument above (before this lemma), $r=\tau r'$ with $r'\in\pkpA$. Then $E^h\tau r'=z\tau w$ in the domain $\pkfBp$
gives $E^hr'=zw$, i.e.\ $z$ divides $E^hr'$ in $\pkfBp$, and $z\mid E^hr'$ in $\pkpA$ by Lemma~\ref{pkl:intdiv}. As
$\theta(z)\ne0$ and $\ker\theta=E\pkpA$ is prime ($\pkpA/E\pkpA\cong O_C$), $E\nmid z$; $z\mid E^hr'$ then gives
$z\mid r'$: by induction on $h$, if $E^hr'=zq$ then $\theta(z)\theta(q)=0$, $q=Eq'$, and $E^{h-1}r'=zq'$.
\end{proof}

\begin{proposition}[(G0); Lemma~\ref{pkm:inputs}(G), second assertion]\label{pkl:G0}
For $x\in\mathcal N(L)$ and $g\in\Gamma'$, $(z\tau)^{-1}(gx-x)\in\pklAN(L)$. Consequently, if $x\in\mathcal N(L)$ and
$gx-x\in L_i(o)$, then $(z\tau)^{-1}(gx-x)\in L_i(o)$.
\end{proposition}

\begin{proof}
The consequence and the reduction to $x\in\mathcal N\cap\pknMinf$ were shown at the beginning of this subsection. Let
$x\in\mathcal N\cap\pknMinf$; fix a frame. As in the proof of Corollary~\ref{pkl:Phi}, $x=\sum_jt_jn_j$ with
$t_j\in\mathfrak S$, so $X_\nu:=E^a\sigma_\nu(x)=(Zt)_\nu\in\pkfSB$ (Lemma~\ref{pkn:frame-id}(ii), $m=0$, with the forms
$\sigma_\nu$ of Lemma~\ref{pkn:frame-id}(iii)). Write $gE=\omega E$, $\omega\in\pkpA^\times$; $z\tau\mid\omega-1$ (proof of
Proposition~\ref{pkl:Cd}, Step 1, valid for any $g\in\Gamma'$), hence $z\tau\mid\omega^a-1$. By Lemma~\ref{pkm:equiv},
$\sigma_\nu(gx-x)=g\sigma_\nu(x)-\sigma_\nu(x)$, and from $E^a\sigma_\nu(x)=X_\nu$, $\omega^aE^ag\sigma_\nu(x)=gX_\nu$. So
\[
\omega^aE^{2a}\bigl(g\sigma_\nu(x)-\sigma_\nu(x)\bigr)=E^a(gX_\nu-X_\nu)-(\omega^a-1)E^aX_\nu\in W:=z\tau\pkfBp
\]
(image in $\pknB$, a $\pkfBp$-submodule), by Lemma~\ref{pkl:galSB}; as $\omega$ is a unit, $\sigma_\nu(E^{2a}(gx-x))\in W$ for
all $\nu$. By Lemma~\ref{pkl:G1} and Proposition~\ref{pkn:frame}(ii), $gx-x=\sum_js_jn_j$ with $s_j\in\pkpA$. The proof of
the coefficient test Lemma~\ref{pkn:frame-id}(iv) uses its hypothesis only for the forms $\sigma_\nu$; applied to
$E^{2a}(gx-x)=\sum_jE^{2a}s_jn_j$ it gives $E^{4a}s_j\in W$, i.e.\ (as $\pkfBp\to\pknB$ is injective) $z\tau$ divides
$E^{4a}s_j$ in $\pkfBp$. By Lemma~\ref{pkl:cancel}, $s_j=z\tau s'_j$ with $s'_j\in\pkpA$, and $(z\tau)^{-1}(gx-x)=\sum_js'_jn_j\in
\pknMinf\subseteq\pklAN$.
\end{proof}



\providecommand{\auO}{\mathcal O}
\providecommand{\auQlb}{\overline{\Q}_\ell}
\providecommand{\pkpA}{A_{\mathrm{inf}}}
\providecommand{\pkpW}{W(C^\flat)}
\providecommand{\pkpm}{W(\mathfrak m^\flat)}
\providecommand{\pkpM}{M_m}
\providecommand{\pkpsm}{\mathrm{Sm}_m}
\providecommand{\pkpAmb}{\mathbb A}
\providecommand{\pknBp}{B^+_{\mathrm{dR}}}
\providecommand{\pknB}{B_{\mathrm{dR}}}
\providecommand{\pknMinf}{\mathfrak M_{\mathrm{inf}}}
\providecommand{\pknK}{\mathcal K_0}
\providecommand{\pkcPt}{\mathrm{Pt}}
\providecommand{\pkcsh}{\mathrm{sh}}
\providecommand{\pktV}{V_{\mathrm{dR}}}
\providecommand{\pktFr}{\mathrm{Fr}}

\subsection{Orders add, Frobenius and the points from killers}\label{pkt:sec}

This section proves Lemma~\ref{pkp:orders-add} (Lemma~\ref{pkt:add}), the assertions (i), (iv), (v)
of Lemma~\ref{pkp:tau} (in the form of \S\ref{pkn:ss-def}), and the inputs (I1) and (I3) of
\S\ref{pkc:ss-inputs}. We use the conventions and labels of \S\S\ref{pkp:sec}--\ref{pkn:sec}: the
modules $\pknMinf(L)\subseteq\pkpA\mathcal N(L)$ and $\mathcal N(L)$ are those of
Definition~\ref{pkn:def-N}, forms are as defined there, and $L_i(o)$ is the module of elements with
orders $o$ at $i$ (Definition~\ref{pkn:def-N}, last sentence). We write $\mathrm{Fil}^n=t^n\pknBp$,
$K'=\auQlb^{\,j(\Gamma')}$, and for a letter $j$ we write $j=i^{(-a)}$ when $j^{\ell^a}=i$.

Besides the published facts (P2), (P3) of \S\ref{pkf:ss-rings} (Fontaine) and Tate's theorem on
$C(n)^{G_{K'}}$ (already used in Lemma~\ref{pkm:BdR}), we use Theorem~\ref{pkf:main},
Proposition~\ref{pkn:frame}, and items of Lemma~\ref{pkm:inputs}, cited by letter; these are proved in
\S\S\ref{pkn:sec}, \ref{pkl:sec}, and the precise references are collected in \S\ref{pkt:ss-open}.

Two statements about the modules $\mathcal N$ are used besides the items of
Lemma~\ref{pkm:inputs} (they are listed there too, and proved in \S\ref{pkl:sec}):
\begin{enumerate}
\item[(G$_0$)] For $x\in\mathcal N(L)$ and $g\in\Gamma'$: $(z\tau)^{-1}(gx-x)\in\pkpA\mathcal N(L)$.
\item[($\Phi$)] $\varphi(\mathcal N(L))\subseteq\mathcal N(L)$.
\end{enumerate}
(G$_0$) is the membership half of the second assertion of Lemma~\ref{pkm:inputs}(G): by
Lemma~\ref{pkt:tau-iv} below, (G$_0$) implies that assertion. The inclusion $\varphi(\pkpA\mathcal
N(L))\subseteq\pkpA\mathcal N(L)$ is proved below (Lemma~\ref{pkt:frob}); ($\Phi$) adds the condition on the
values in $\pknK$.

\subsubsection{The Hodge expansion}\label{pkt:ss-hodge}

Let $L$ be a lattice with jump sets. Put $\pktV=\pknB\otimes_{\Z_\ell}L$ and $V^+=\pknBp\otimes_{\Z_\ell}L$;
$\Gamma$ acts on both diagonally (through $j$ on $\pknB$), semilinearly over $\pknB$. The ring
$\Lambda\otimes_{\Z_\ell}\auO$ acts on $\pktV$ ($\Lambda\subset\pkpA\to\pknBp$ on the first factor, $\auO$ on
$L$), so the projectors $e_j$ of Definition~\ref{pkp:def-letters} act on $\pktV$ and $V^+$; put
$V_j=e_j\pktV$, $V_j^+=e_jV^+$. Since $\Gamma'$ fixes $\Lambda$ and commutes with $\auO$, each $e_j$ commutes
with $\Gamma'$, and $\pktV=\bigoplus_jV_j$, $V^+=\bigoplus_jV_j^+$.

\emph{Convention on jumps.} For an embedding $\gamma\colon\auO\to C$, an integer $\kappa$ is a jump of $L$
at $\gamma$ if $C(-\kappa)$ occurs in the Hodge--Tate decomposition of $C\otimes_{\auO,\gamma}L$ (with
multiplicity the number of times it occurs). This is the convention of Lemma~\ref{pkp:tau} and of
(P3) in \S\ref{pkf:ss-rings}: $\auO(1)$ has the single jump $-1$. By Definition~\ref{pkp:def-orders} the
jumps at the embeddings of letter $j$ lie in $U(j)$.

\emph{Elements with a denominator.} Let $\pkpAmb(L)^{\mathrm{fr}}$ be the set of $x\in\pkpAmb(L)$ with
$rx\in\pkpA\otimes L$ for some nonzero $r\in\pkpA$ ($\pkpA\otimes L$ the $\pkpA$-span of the $1\otimes y$,
a free $\pkpA$-module on $1\otimes e_1,\dots,1\otimes e_d$ for a $\Z_\ell$-basis $(e_k)$ of $L$). The map
$\pkpA\to\pknBp$ is injective (\S\ref{pkn:ss-frame}), so a nonzero $r$ is invertible in $\pknB$, and
$\hat x=r^{-1}\iota(rx)\in\pktV$, with $\iota\colon\pkpA\otimes L\to\pktV$ the natural map, does not
depend on $r$. The set $\pkpAmb(L)^{\mathrm{fr}}$ is an $\pkpA$-module, stable under $\varphi^{\pm1}$ and
$\Gamma$, and $x\mapsto\hat x$ is $\pkpA$-linear, $\Gamma$-equivariant and commutes with the $e_j$. By
Lemma~\ref{pkn:M}(ii), $\pkpA\mathcal N(L)\subseteq\pkpAmb(L)^{\mathrm{fr}}$; if $x$ lies in the letter-$i$
part, $\widehat{\varphi^{-a}x}\in V_{i^{(-a)}}$ (Definition~\ref{pkp:def-letters}).

For an $\pkpA$-linear $\mu\colon\pkpAmb(L)\to\pknB$ let $\hat\mu\colon\pktV\to\pknB$ be the $\pknB$-linear
map with $\hat\mu(1\otimes y)=\mu(1\otimes y)$. Then $\mu(x)=\hat\mu(\hat x)$ for
$x\in\pkpAmb(L)^{\mathrm{fr}}$: indeed $r\mu(x)=\mu(rx)=\hat\mu(\iota(rx))=r\hat\mu(\hat x)$, and $r\ne0$ in
the field $\pknB$.

\begin{lemma}[Forms and equivariant functionals]\label{pkt:forms}
(i) $\mu$ is a form iff $\hat\mu(gv)=g\hat\mu(v)$ for all $g\in\Gamma'$, $v\in\pktV$; it has type $u$ at
$j$ iff $\hat\mu(V_j^+)\subseteq\mathrm{Fil}^{-u}$. (ii) Every $\pknB$-linear $\lambda\colon\pktV\to\pknB$
with $\lambda(gv)=g\lambda(v)$ for $g\in\Gamma'$ is $\hat\mu$ for some form $\mu$. (iii) If $\lambda$ is as
in (ii) and $d\in\pktV$ is fixed by $\Gamma'$, then $\lambda(d)\in K'$; a nonzero element of $K'$ is a
unit of $\pknBp$.
\end{lemma}

\begin{proof}
(i) If $\mu$ is a form, $\hat\mu(g(b\otimes y))=\hat\mu(gb\otimes gy)=g(b)\,g(\mu(1\otimes y))=
g(\hat\mu(b\otimes y))$; the converse is the case $v=1\otimes y$. $V_j^+$ is the $\pknBp$-span of
$e_j(1\otimes L)$, whence the second claim. (ii) Let $\lambda_0\colon\pkpW\to\pknB$ be an $\pkpA$-linear
map with $\lambda_0(1)=1$ (it exists: proof of Lemma~\ref{pkn:frame-id}(iii)); then
$\mu(\sum_kw_k\otimes e_k)=\sum_k\lambda_0(w_k)\lambda(1\otimes e_k)$ is $\pkpA$-linear with $\hat\mu=
\lambda$, and it is a form by (i). (iii) $g(\lambda(d))=\lambda(gd)=\lambda(d)$ for $g\in\Gamma'$, and the
elements of $\pknB$ fixed by $G_{K'}$ lie in $K'$ (proof of Lemma~\ref{pkn:frame-id}(i)). A nonzero
$c\in K'$ has $\theta(c)=c\ne0$, so it is not in the maximal ideal $\mathrm{Fil}^1$ of $\pknBp$.
\end{proof}

\begin{lemma}[Hodge bases at a letter]\label{pkt:hodge}
For every letter $j$ there are elements $d_{j,1},\dots,d_{j,n_j}\in V_j$ fixed by $\Gamma'$ and integers
$\kappa_{j,1},\dots,\kappa_{j,n_j}\in U(j)$ such that $(t^{-\kappa_{j,m}}d_{j,m})_m$ is a $\pknBp$-basis of
$V_j^+$ (hence $(d_{j,m})_m$ is a $\pknB$-basis of $V_j$).
\end{lemma}

\begin{proof}
$L[1/\ell]$ is crystalline, hence de Rham, as a representation of $G_{K'}$ (P2). Let
$D=\pktV^{G_{K'}}$ with $\mathrm{Fil}^kD=D\cap t^kV^+$. By (P3), $\pknB\otimes_{K'}D\to\pktV$ is an
isomorphism carrying $\sum_kt^{-k}\pknBp\otimes\mathrm{Fil}^kD$ onto $V^+$. The $e_j$ have coefficients in
$\Lambda[1/\ell]\otimes\auO\subseteq K'\otimes\auO$, commute with $G_{K'}$, preserve $V^+$ and commute with
$t$; so $D=\bigoplus_jD_j$ with $D_j=e_jD$, and $\mathrm{Fil}^kD=\bigoplus_j(\mathrm{Fil}^kD\cap D_j)$.
Choose a $K'$-basis $d_{j,1},\dots,d_{j,n_j}$ of $D_j$ adapted to the filtration: $d_{j,m}\in
\mathrm{Fil}^{\kappa_{j,m}}D_j$ and, for each $k$, the $d_{j,m}$ with $\kappa_{j,m}\ge k$ form a basis of
$\mathrm{Fil}^kD_j$. Then $\sum_kt^{-k}\pknBp\otimes\mathrm{Fil}^kD_j=\bigoplus_mt^{-\kappa_{j,m}}\pknBp\,
d_{j,m}$ (an element $t^{-k}b\otimes d_{j,m}$ with $k\le\kappa_{j,m}$ lies in $t^{-\kappa_{j,m}}\pknBp
d_{j,m}$), and applying $e_j$ to the isomorphism of (P3) gives that $(t^{-\kappa_{j,m}}d_{j,m})_m$ is a
$\pknBp$-basis of $V_j^+$.

It remains to see $\kappa_{j,m}\in U(j)$. Reduce modulo $t$: $V_j^+/tV_j^+=e_j(C\otimes_{\Z_\ell}L)$. Since
$\auO[1/\ell]$ is a finite separable extension of $\Q_\ell$, $C\otimes_{\Z_\ell}\auO=\prod_\gamma C$ over the
embeddings $\gamma\colon\auO\to C$, so $C\otimes_{\Z_\ell}L=\bigoplus_\gamma C\otimes_{\auO,\gamma}L$, and
$e_j$ acts on the factor $\gamma$ by the scalar $\prod_{i'\ne j}(\gamma(\zeta)-\theta[i'])/(\theta[j]-
\theta[i'])$. As $\zeta^{\ell^f-1}=1$, $\gamma(\zeta)$ is the Teichm\"uller representative $\theta[i']$ of
its residue $i'$, the letter of $\gamma$; so the scalar is $1$ if $\gamma$ has letter $j$ and $0$
otherwise. Hence $V_j^+/tV_j^+=\bigoplus_{\gamma\text{ of letter }j}C\otimes_{\auO,\gamma}L$. On the
other hand $V_j^+/tV_j^+=\bigoplus_mC\cdot\overline{t^{-\kappa_{j,m}}d_{j,m}}$, and $g\in\Gamma'$ acts on
$\overline{t^{-\kappa}d}$ ($d$ fixed) through $\chi(g)^{-\kappa}$, so this summand is $C(-\kappa_{j,m})$. By
the uniqueness of the Hodge--Tate decomposition (Tate's theorem $C(n)^{G_{K'}}=0$ for $n\ne0$, as in
Lemma~\ref{pkm:BdR}), the multiset of the $\kappa_{j,m}$ is the union over the embeddings $\gamma$ of letter
$j$ of the multisets of jumps of $L$ at $\gamma$; these lie in $U(j)$.
\end{proof}

Fix such bases. Let $\delta_{j,m}\colon\pktV\to\pknB$ be the $\pknB$-linear map vanishing on $V_{j'}$
for $j'\ne j$ and sending $v\in V_j$ to its coefficient on $d_{j,m}$. It is $\Gamma'$-equivariant (the
$d_{j,m}$ are fixed and the decomposition $\bigoplus V_{j'}$ is $\Gamma'$-stable) and $\delta_{j,m}(V_{j'}^+)
\subseteq\mathrm{Fil}^{-\kappa_{j,m}}$ for all $j'$; so by Lemma~\ref{pkt:forms}, $\delta_{j,m}=\hat\mu$ for a
form $\mu$ of type $\kappa_{j,m}$ at every letter. We call $\delta_{j,m}$ the \emph{coordinate forms}.

\begin{lemma}[Expansion]\label{pkt:expand}
Let $\lambda\colon\pktV\to\pknB$ be $\pknB$-linear and $v\in V_j$. Then
$\lambda(v)=\sum_m\lambda(d_{j,m})\,\delta_{j,m}(v)$. If $\lambda(V_j^+)\subseteq\mathrm{Fil}^{s}$, then
$\lambda(d_{j,m})\in\mathrm{Fil}^{s+\kappa_{j,m}}$. If $\lambda$ is $\Gamma'$-equivariant and $\lambda(V_j^+)
\subseteq\mathrm{Fil}^{-u}$, then $\lambda(d_{j,m})\in K'$, and $\lambda(d_{j,m})=0$ unless $\kappa_{j,m}\le u$.
\end{lemma}

\begin{proof}
The first claim is linearity, $v=\sum_m\delta_{j,m}(v)d_{j,m}$. Next, $\lambda(t^{-\kappa_{j,m}}d_{j,m})\in
\mathrm{Fil}^s$ gives $\lambda(d_{j,m})\in\mathrm{Fil}^{s+\kappa_{j,m}}$. If $\lambda$ is equivariant,
$\lambda(d_{j,m})\in K'$ by Lemma~\ref{pkt:forms}(iii); being in $\mathrm{Fil}^{\kappa_{j,m}-u}$, it is $0$ if
$\kappa_{j,m}-u\ge1$ (a nonzero element of $K'$ is a unit of $\pknBp$).
\end{proof}

\begin{lemma}[Values at non-jumps]\label{pkt:gap}
Let $x\in L_i(c)$, $a\ge0$, $j=i^{(-a)}$, $u\in\Z$ (not necessarily in $U(j)$), and $\mu$ a form of type
$u$ at $j$. Then $\mu(\varphi^{-a}x)\in\sum_{v}\mathrm{Fil}^{c(a,v)}$, the sum over the $v\in U(j)$ with
$v\le u$ (the value is $0$ if there is no such $v$).
\end{lemma}

\begin{proof}
$y=\widehat{\varphi^{-a}x}\in V_j$ and $\mu(\varphi^{-a}x)=\hat\mu(y)=\sum_m\hat\mu(d_{j,m})\delta_{j,m}(y)$ by
Lemma~\ref{pkt:expand}. The coordinate form $\delta_{j,m}$ has type $\kappa_{j,m}\in U(j)$ at $j$, so
$\delta_{j,m}(y)\in\mathrm{Fil}^{c(a,\kappa_{j,m})}$ by the definition of orders. By
Lemmas~\ref{pkt:forms}(i) and \ref{pkt:expand}, $\hat\mu(d_{j,m})\in K'\subset\pknBp$ vanishes unless
$\kappa_{j,m}\le u$.
\end{proof}

\subsubsection{Orders add}\label{pkt:ss-add}

\begin{lemma}[Lemma~\ref{pkp:orders-add}]\label{pkt:add}
Let $L,L',L''$ be lattices with jump sets $U,U',U''$ and $b\colon L\times L'\to L''$ an $\auO$-bilinear
map with $b(gy,gy')=gb(y,y')$ for $g\in\Gamma'$; let $\tilde b\colon\pkpAmb(L)\times\pkpAmb(L')\to
\pkpAmb(L'')$ be its $\pkpW$-bilinear extension. (i) If $w\in\pknMinf(L)$, $y\in\pknMinf(L')$ then
$\tilde b(w,y)\in\pknMinf(L'')$; if $w\in\pkpA\mathcal N(L)$, $y\in\pkpA\mathcal N(L')$ then
$\tilde b(w,y)\in\pkpA\mathcal N(L'')$. (ii) Let $w\in L_i(o)$, $y\in L'_i(o')$, and let $o''$ satisfy
\[
o''(a,v)\le o(a,u)+o'(a,u')\quad\text{whenever }u+u'\le v,
\]
for all $a\ge0$, $u\in U(i^{(-a)})$, $u'\in U'(i^{(-a)})$, $v\in U''(i^{(-a)})$. Then $\tilde b(w,y)\in
L''_i(o'')$. In particular, for $L''=L\otimes_\auO L'$ with jump sets $U(i)+U'(i)$ and $b=\otimes$, $w\otimes
y\in(L\otimes L')_i(o'')$.
\end{lemma}

(Part (ii) is Lemma~\ref{pkp:orders-add}; part (i) is the membership statement used in its proof. For
$L\otimes L'$, the sets $U(i)+U'(i)$ contain the jumps: if $(d_{j,m})$, $(d'_{j,n})$ are Hodge bases of
$L,L'$ at $j$, the $e_j(d_{j,m}\otimes d'_{j,n})$ span $e_j(\pktV\otimes\pktV')$ and the argument of
Lemma~\ref{pkt:hodge} shows that the jumps at letter $j$ are sums $\kappa_{j,m}+\kappa'_{j,n}$.)

\begin{proof}
Let $\hat b\colon\pktV\times\pktV'\to\pktV''$ be the $\pknB$-bilinear extension of $b$; then
$\widehat{\tilde b(w,y)}=\hat b(\hat w,\hat y)$ for $w,y$ with denominators (multiply by the two
denominators), and $\hat b$ is $\Gamma'$-equivariant. Since $b$ is $\auO$-bilinear, $\tilde b(e_jw,y)=
e_j\tilde b(w,y)=\tilde b(w,e_jy)$, so $\hat b(V_j,V'_{j'})=0$ for $j\ne j'$ and $\hat b(V_j,V'_j)\subseteq
V''_j$; also $\hat b(V^+,V'^+)\subseteq V''^+$. Fix Hodge bases $(d'_{j,n},\kappa'_{j,n})$ of $L'$.

\emph{Key identity.} Let $\eta$ be a form on $L''$, $j$ a letter and $n$ an index. The map
$\lambda_{j,n}\colon v\mapsto\hat\eta(\hat b(v,d'_{j,n}))$ is $\pknB$-linear and $\Gamma'$-equivariant (as
$d'_{j,n}$ is fixed), so $\lambda_{j,n}=\hat\sigma_{j,n}$ for a form $\sigma_{j,n}$ on $L$ (Lemma~\ref{pkt:forms}(ii)).
For $v\in\pktV$ and $v'\in V'_j$, by Lemma~\ref{pkt:expand} applied to $v'\mapsto\hat\eta(\hat b(v,v'))$,
\[
\hat\eta(\hat b(v,v'))=\sum_n\hat\sigma_{j,n}(v)\,\delta'_{j,n}(v').\tag{$*$}
\]

(i) If $\tau^Nw,\tau^{N'}y\in\pkpA\otimes(\cdot)$ then $\tau^{N+N'}\tilde b(w,y)\in\pkpA\otimes L''$. Let
$k\ge1$ and $\eta$ a form on $L''$. As $\varphi$ acts on the first factor of $\pkpW\otimes(\cdot)$,
$\varphi^k\tilde b(w,y)=\tilde b(\varphi^kw,\varphi^ky)$. Write $\widehat{\varphi^ky}=\sum_jv'_j$ with
$v'_j\in V'_j$; by $(*)$,
$\eta(\varphi^k\tilde b(w,y))=\sum_j\sum_n\sigma_{j,n}(\varphi^kw)\,\delta'_{j,n}(\widehat{\varphi^k y})$. Each
factor is the value of a form at $\varphi^k$ of an element of $\pknMinf$, hence lies in $\pknBp$. So
$\tilde b(w,y)\in\pknMinf(L'')$. If $rw\in\pknMinf(L)$, $r'y\in\pknMinf(L')$ with $r,r'$ products of
the $p_a$, then $rr'\tilde b(w,y)=\tilde b(rw,r'y)\in\pknMinf(L'')$.

(ii) By (i), $\tilde b(w,y)\in\pkpA\mathcal N(L'')$, and it lies in the letter-$i$ part as
$e_i\tilde b(w,y)=\tilde b(e_iw,y)=\tilde b(w,y)$. Let $a\ge0$, $j=i^{(-a)}$, $v\in U''(j)$ and $\eta$ a form
of type $v$ at $j$. Put $\omega=\widehat{\varphi^{-a}w}\in V_j$ and $\upsilon=\widehat{\varphi^{-a}y}\in V'_j$;
then $\eta(\varphi^{-a}\tilde b(w,y))=\hat\eta(\hat b(\omega,\upsilon))=\sum_n\hat\sigma_{j,n}(\omega)
\delta'_{j,n}(\upsilon)$ by $(*)$. The coordinate form $\delta'_{j,n}$ has type $\kappa'_{j,n}\in U'(j)$ at
$j$, so $\delta'_{j,n}(\upsilon)\in\mathrm{Fil}^{o'(a,\kappa'_{j,n})}$. The form $\sigma_{j,n}$ has type
$v-\kappa'_{j,n}$ at $j$: $\hat\sigma_{j,n}(V_j^+)=t^{\kappa'_{j,n}}\hat\eta(\hat b(V_j^+,t^{-\kappa'_{j,n}}
d'_{j,n}))\subseteq t^{\kappa'_{j,n}}\hat\eta(V''^+_j)\subseteq\mathrm{Fil}^{\kappa'_{j,n}-v}$. By
Lemma~\ref{pkt:gap}, $\hat\sigma_{j,n}(\omega)\in\sum_u\mathrm{Fil}^{o(a,u)}$ over $u\in U(j)$ with $u\le
v-\kappa'_{j,n}$. Each product thus lies in some $\mathrm{Fil}^{o(a,u)+o'(a,u')}$ with $u\in U(j)$,
$u'=\kappa'_{j,n}\in U'(j)$, $u+u'\le v$, hence in $\mathrm{Fil}^{o''(a,v)}$.
\end{proof}

\subsubsection{Frobenius, the standard orders and Lemma~\ref{pkp:tau}}\label{pkt:ss-tau}

\begin{lemma}[Frobenius on orders]\label{pkt:frob}
(i) $\varphi(\pknMinf(L))\subseteq E^{-N}\pknMinf(L)$ for elements $x$ with $\tau^Nx\in\pkpA\otimes L$; hence
$\varphi^l(\pkpA\mathcal N(L))\subseteq\pkpA\mathcal N(L)$ for $l\ge0$. (ii) For $x\in\pkpA\mathcal N(L)$, $m\ge0$
and every form $\mu$, $\mu(\varphi^{m+1}x)\in\pknBp$. (iii) If $x\in L_{i^{(-l)}}(o)$, then $\varphi^lx\in
L_i(\pkcsh_{l,0}(o))$.
\end{lemma}

\begin{proof}
(i) Let $x\in\pknMinf$, $\tau^Nx\in\pkpA\otimes L$. Since $\varphi(\tau)=[\varepsilon]-1=E\tau v$ with
$v\in\pkpA^\times$ (Lemma~\ref{pkn:mu}) and $\varphi$ fixes $1\otimes y$, $\varphi(\tau^Nx)=E^N\tau^Nv^N\varphi(x)
\in\pkpA\otimes L$, so $\tau^N(E^N\varphi x)\in\pkpA\otimes L$; and for $k\ge1$, $\mu(\varphi^k(E^N\varphi
x))=\varphi^k(E)^N\mu(\varphi^{k+1}x)\in\pknBp$. So $E^N\varphi(x)\in\pknMinf$. If $rx\in\pknMinf$ with $r$ a
product of the $p_b$, then $\varphi(r)$ is a product of the $p_{b+1}$ and $E^N\varphi(r)\varphi(x)\in\pknMinf$,
so $\varphi(x)\in\pkpA\mathcal N$; iterate.

(ii) If $x\in\pknMinf$ this is the definition. In general let $x=r^{-1}x_0$ with $x_0\in\pknMinf$ and $r$ a
product of $p_b$'s; then $\varphi^{m+1}(r)$ is a product of $p_{b+m+1}$'s, whose images in $\pknBp$ are units
(Lemma~\ref{pkn:BdR}: $\varphi^{c}(E)$ with $c\ge1$), and $\mu(\varphi^{m+1}x)=\varphi^{m+1}(r)^{-1}\mu(\varphi^{m+1}
x_0)\in\pknBp$.

(iii) By (i), $\varphi^lx\in\pkpA\mathcal N$, and $\varphi^l$ maps the letter part at $i^{(-l)}$ to that at $i$
(Definition~\ref{pkp:def-letters}). Let $a\ge0$, $u\in U(i^{(-a)})$, $\mu$ a form of type $u$ at $i^{(-a)}$.
If $a<l$, $\mu(\varphi^{-a}\varphi^lx)=\mu(\varphi^{l-a}x)\in\pknBp=\mathrm{Fil}^0$ by (ii). If $a\ge l$,
$\mu(\varphi^{-a}\varphi^lx)=\mu(\varphi^{-(a-l)}x)\in\mathrm{Fil}^{o(a-l,u)}$, since $i^{(-a)}=(i^{(-l)})^{(-(a-l))}$.
\end{proof}

\begin{lemma}[Standard orders of level $0$]\label{pkt:std0}
If $x\in L_i(\mathrm{std}_0)$, then $x\in\pknMinf(L)$.
\end{lemma}

\begin{proof}
\emph{Denominator.} Let $T$ be the set of $z\in\pkpAmb(L)$ with $\tau^nz\in\pkpA\otimes L$ for some $n$. We
show: if $x'\in L_i(\mathrm{std}_0)$, $b\ge0$ and $p_bx'\in T$, then $x'\in T$. Let $\tau^np_bx'=\sum_kc_k
(1\otimes e_k)$, $c_k\in\pkpA$. Let $\mu_k$ be an $\pkpA$-linear map $\pkpAmb(L)\to\pknB$ with
$\mu_k(1\otimes y)=e_k^*(y)$ (the coordinate; construct it as in Lemma~\ref{pkt:forms}(ii)); then
$\hat\mu_k(V^+)\subseteq\pknBp$. With $j=i^{(-b)}$ and $\omega=\widehat{\varphi^{-b}x'}\in V_j$,
Lemma~\ref{pkt:expand} gives $\hat\mu_k(\omega)=\sum_m\hat\mu_k(d_{j,m})\delta_{j,m}(\omega)$ with
$\hat\mu_k(d_{j,m})\in\mathrm{Fil}^{\kappa_{j,m}}$ and $\delta_{j,m}(\omega)\in\mathrm{Fil}^{\mathrm{std}_0(b,
\kappa_{j,m})}=\mathrm{Fil}^{-\kappa_{j,m}}$. So $\mu_k(\varphi^{-b}x')\in\pknBp$. Applying $\mu_k\circ\varphi^{-b}$
to the expansion: $\varphi^{-b}(c_k)=\varphi^{-b}(\tau)^n\,E\,\mu_k(\varphi^{-b}x')\in\mathrm{Fil}^1$, since
$\varphi^{-b}(p_b)=E$. Hence $\theta(\varphi^{-b}c_k)=0$, i.e.\ $c_k\in p_b\pkpA$ (Definition~\ref{pkp:def-chart}),
say $c_k=p_bc'_k$, and $\tau^nx'=\sum_kc'_k(1\otimes e_k)$ ($p_b$ is a unit of $\pkpW$). Now $x\in\pkpA\mathcal
N$, so $rx\in\pknMinf\subseteq T$ for a product $r=p_{b_1}\cdots p_{b_k}$; by Lemma~\ref{pkn:E}, $p_{b_{s+1}}
\cdots p_{b_k}x\in L_i(o_s)$ with $o_s\ge\mathrm{std}_0$, so these elements lie in $L_i(\mathrm{std}_0)$ (orders
only increase), and removing the factors one at a time gives $x\in T$.

\emph{Values.} Let $n\ge1$ and $\mu$ a form. By Lemma~\ref{pkt:frob}(iii), $\varphi^nx\in
L_{i'}(\pkcsh_{n,0}(\mathrm{std}_0))$ with $i'=i^{\ell^n}$; at level $0<n$ these orders are $0$. With
$y=\widehat{\varphi^nx}\in V_{i'}$, $\mu(\varphi^nx)=\sum_m\hat\mu(d_{i',m})\delta_{i',m}(y)$, where
$\hat\mu(d_{i',m})\in K'\subset\pknBp$ (Lemma~\ref{pkt:forms}(iii)) and $\delta_{i',m}(y)\in\mathrm{Fil}^0$
(coordinate form of type $\kappa_{i',m}\in U(i')$, orders $0$ at level $0$). So $\mu(\varphi^nx)\in\pknBp$.
\end{proof}

\begin{proposition}[Lemma~\ref{pkp:tau}(i)]\label{pkt:tau-i}
For the lattice $\auO(1)$ with jump sets $\{-1\}$, every $\xi\in\auO(1)_i(\mathrm{std}_0)$ has $\tau\xi\in
\pkpA\otimes\auO(1)$.
\end{proposition}

\begin{proof}
By Lemma~\ref{pkt:std0}, $\xi\in\pknMinf(\auO(1))$. The jump bound of $\auO(1)$ (Theorem~\ref{pkf:main}) is
$\max\{|u|:u\in\{-1\}\}=1$; by Theorem~\ref{pkf:main} it has a frame for $a=1$, and by
Proposition~\ref{pkn:frame}(iii), $\tau\pknMinf\subseteq\pkpA\otimes\auO(1)$.
\end{proof}

Lemma~\ref{pkp:tau}(ii) is Lemma~\ref{pkp:AS}(vii) and (iii) is Lemma~\ref{pkm:tau-iii}.

\begin{lemma}[Lemma~\ref{pkp:tau}(iv)]\label{pkt:tau-iv}
Let $x\in\mathcal N(L)$ lie in the letter-$i$ part and $gx-x\in L_i(o')$ for all $g\in\Gamma'$. Then
$(z\tau)^{-1}(gx-x)\in L_i(o')$ for $g\in\Gamma'$. If $L=X_S$ and $\mathrm{pr}(x)$ is fixed by $\Gamma'$, then
$(z\tau)^{-1}(gx-x)$ lies over $0$. More generally, for $y\in\pkpA\mathcal N(L)$ and $r\in\pkpA$ not divisible by any $p_a$: if $ry\in
L_i(o)$ then $y\in L_i(o)$.
\end{lemma}

\begin{proof}
The last assertion first. $y$ lies in the letter-$i$ part as $r$ commutes with $e_i$ and is a nonzerodivisor on
$\pkpAmb(L)=\pkpW\otimes_{\Z_\ell}L$ ($\pkpW$ is a domain, $L$ is free, and $r\ne0$ since $p_0\nmid r$), so
$r\bigl(e_iy-y\bigr)=e_i(ry)-ry=0$ gives $e_iy=y$. For $a\ge0$, $\theta(\varphi^{-a}r)\ne0$ (else $r\in p_a\pkpA$), so the image of $\varphi^{-a}(r)$
is a unit of $\pknBp$, and $\mu(\varphi^{-a}y)=\varphi^{-a}(r)^{-1}\mu(\varphi^{-a}(ry))$ has the same filtration
as $\mu(\varphi^{-a}(ry))$. Now take $r=z\tau$ and $y=(z\tau)^{-1}(gx-x)$, which lies in $\pkpA\mathcal N$ by
(G$_0$). No $p_a$ divides $z\tau$: $\theta(\varphi^{-a}z)=\theta([\varpi^{\flat\,1/\ell^a}])$ is an
$\ell^a$-th root of $\pi\ne0$, and $\theta(\varphi^{-a}\tau)=\zeta_{\ell^{a+1}}-1\ne0$. So $y\in L_i(o')$.
Finally $\mathrm{pr}$ is $\pkpW$-linear and $\Gamma$-equivariant, so $\mathrm{pr}(y)=(z\tau)^{-1}(g\,\mathrm{pr}(x)-
\mathrm{pr}(x))=0$.
\end{proof}

\begin{lemma}[Lemma~\ref{pkp:tau}(v)]\label{pkt:tau-v}
If $x\in L_{i^{(-l)}}(o)$ then $\varphi^lx\in L_i(\pkcsh_{l,0}(o))$. If $x$ is a point of $S$ over
$\ell^ec$ at $i^{(-l)}$ with orders $o$ and depth $o'$, then $\varphi^lx$ is a point over $\ell^ec$ at $i$
with orders $\pkcsh_{l,0}(o)$ and depth $\pkcsh_{l,1}(o')$.
\end{lemma}

\begin{proof}
The first assertion is Lemma~\ref{pkt:frob}(iii). Let $x$ be a point. $\varphi^lx\in\mathcal N$ by ($\Phi$);
$\mathrm{pr}(\varphi^lx)=\varphi^l(\ell^ee_{i^{(-l)}}(1\otimes c))=\ell^ee_i(1\otimes c)$
(Definition~\ref{pkp:def-letters}); the orders are the first assertion. Depth: for $g\in\Gamma'$ put $y=
(z\tau)^{-1}(gx-x)$; $\mathrm{pr}(x)$ is fixed by $\Gamma'$, so $y\in L_{i^{(-l)}}(o')$ by
Lemma~\ref{pkt:tau-iv}, and $\varphi^ly\in L_i(\pkcsh_{l,0}(o'))$. As $g$ commutes with $\varphi$,
$g\varphi^lx-\varphi^lx=\varphi^l(z\tau)\cdot\varphi^ly$, which lies in $L_i(\pkcsh_{l,0}(o'))$. At a level $a<l$
and a form $\mu$ of type $u$ there, $\mu(\varphi^{-a}(g\varphi^lx-\varphi^lx))=\varphi^{l-a}(z\tau)\,
\mu(\varphi^{-a}\varphi^ly)$, where the second factor is in $\mathrm{Fil}^0$ and $\varphi^{l-a}(\tau)=
\varphi^{l-a-1}([\varepsilon]-1)$ has $\theta(\varphi^{l-a-1}([\varepsilon]-1))=1-1=0$, so the first factor is in
$\mathrm{Fil}^1$. Hence the value lies in $\mathrm{Fil}^1=\mathrm{Fil}^{\pkcsh_{l,1}(o')(a,u)}$; at levels $a\ge l$
the orders are $\pkcsh_{l,0}(o')=\pkcsh_{l,1}(o')$.
\end{proof}

\subsubsection{Points from killers: input (I1)}\label{pkt:ss-I1}

\emph{The two readings of ``killing the cokernel''.} Definition~\ref{pkp:def-points} asks for a lift
$n\in(X_S)_i(\mathrm{std}_0)\subseteq\pkpA\mathcal N(X_S)$, not necessarily in $\mathcal N(X_S)$. This
weaker reading is the definition (it is the one
proved by Propositions~\ref{pkm:kill-small}--\ref{pkm:kill-J}, through the lift ideals of
Definition~\ref{pkm:def-J}). For a multiplier in $\mathfrak S$ the two readings are equivalent, by the
retraction (R); this is the content of the next lemma, and it is why (I1) produces points in $\mathcal N$
itself. So (I1) as stated in \S\ref{pkc:ss-inputs} (a point in $\mathcal N(X_S)$, with depth) holds
(Remark~\ref{pkm:rem-kill}).

\begin{lemma}[Lifts in $\mathcal N$]\label{pkt:lift}
Let $S$ be a side, $i$ a letter, $l\ge0$, $f\in\mathfrak S$, and suppose that for every $c\in C_S$ there is
$x\in(X_S)_i(\mathrm{std}_l)$ with $\mathrm{pr}(x)=f\,e_i(1\otimes c)$. Then for every $c$ there is such an $x$
in $\mathcal N(X_S)$.
\end{lemma}

\begin{proof}
Let $\mathrm{ret}$ be the retractions of Lemma~\ref{pkm:inputs}(R) for $X_S$ and for the trivial lattice
$C_S$, and $q=\mathrm{ret}(x)$. By (R), $q\in\mathcal N(X_S)\cap(X_S)_i(\mathrm{std}_l)$ and $\mathrm{pr}(q)=
\mathrm{ret}(\mathrm{pr}(x))=\mathrm{ret}(f\cdot e_i(1\otimes c))$. By Lemma~\ref{pkm:inputs}(H) (second
assertion, $\Gamma$ acting trivially on $C_S$), $e_i(1\otimes c)\in\mathcal N(C_S)$; so by (R),
$\mathrm{ret}(f\,e_i(1\otimes c))=\lambda_0(f)\,e_i(1\otimes c)=f\,e_i(1\otimes c)$, as $\lambda_0$ is the
identity on $\mathfrak S$.
\end{proof}

\begin{proposition}[Input (I1)]\label{pkt:I1}
If $\ell^\kappa E^t$ kills the cokernel of a side $S$ at $i$ (in the sense of \S\ref{pkn:ss-def}), then
for every $c\in C_S$ there is a point of $S$ over $\ell^\kappa c$ at $i$, i.e.\ an element $\nu\in\mathcal
N(X_S)$, with orders $\rho_{-t}$ and depth $\rho_{-t}$: $\pkcPt_S(c,i,\kappa;\rho_{-t},\rho_{-t})\ne\emptyset$.
\end{proposition}

\begin{proof}
By Lemma~\ref{pkt:lift} with $l=0$ and $f=\ell^\kappa E^t\in\mathfrak S$, there is $q\in\mathcal N(X_S)\cap
(X_S)_i(\mathrm{std}_0)$ with $\mathrm{pr}(q)=\ell^\kappa E^te_i(1\otimes c)$. Put $\nu=E^{-t}q$. By
Lemma~\ref{pkn:M}(iii), $\nu\in\mathcal N(X_S)$; by Lemma~\ref{pkn:E} (with $b=0$, $n=-t$), $\nu\in
(X_S)_i(\mathrm{std}_0-t[a=0])=(X_S)_i(\rho_{-t})$; $\mathrm{pr}(\nu)=\ell^\kappa e_i(1\otimes c)$ as $\mathrm{pr}$
is $\pkpW$-linear. Depth: for $g\in\Gamma'$, $g\nu\in(X_S)_i(\rho_{-t})$ by Lemma~\ref{pkm:inputs}(G) (first
assertion), so $g\nu-\nu\in(X_S)_i(\rho_{-t})$.
\end{proof}

\subsubsection{The rule with a dual point: input (I3)}\label{pkt:ss-I3}

\begin{lemma}[Orders in the kernel]\label{pkt:sub}
Let $S$ be a side and $y\in\pkpAmb(Y)$ with image $\bar y=(y,0)\in(X_S)_i(o)$. Then $y\in Y_i(o)$.
\end{lemma}

\begin{proof}
Write $\iota\colon Y\to X_S$ for the inclusion (equivariant) and $\mathrm{fst}$ for the projection
($\auO$-linear, not equivariant). \emph{Extension of forms.} Let $\mu$ be a form on $Y$ of type $u$ at a
letter $j$ (possibly $u=\infty$, i.e.\ no condition). Let $D_Y\subseteq D_X$ be the spaces of
Lemma~\ref{pkt:hodge} for $Y$, $X_S$, and $D_{Y,j}=e_jD_Y\subseteq D_{X,j}$. Since $X_S=Y\oplus C_S$ as
$\auO$-modules, $V_Y\cap t^kV^+_X=t^kV^+_Y$, so the filtration of $D_{Y,j}$ is induced by that of $D_{X,j}$.
Hence an adapted basis of $D_{Y,j}$ extends to an adapted basis $(d_m)$ of $D_{X,j}$ (extend bases of the
graded pieces and lift), and by the proof of Lemma~\ref{pkt:hodge} the $t^{-\kappa_m}d_m$ form a
$\pknBp$-basis of $V^+_{X,j}$. Let $\lambda\colon\pktV(X_S)\to\pknB$ be the $\pknB$-linear map equal to
$\hat\mu$ on $V_{Y,j}$, zero on the added $d_m$, and zero on $V_{X,j'}$ for $j'\ne j$. It is
$\Gamma'$-equivariant (the added $d_m$ are fixed), $\lambda\circ\hat\iota=\hat\mu\circ e_j$, and
$\lambda(V^+_{X,j})=\hat\mu(V^+_{Y,j})\subseteq\mathrm{Fil}^{-u}$. By Lemma~\ref{pkt:forms}(ii), $\lambda=\hat\nu$
for a form $\nu$ on $X_S$, of type $u$ at $j$.

\emph{Membership.} If $r\bar y\in\pknMinf(X_S)$ with $r$ a product of $p_a$'s and $\tau^N r\bar y\in
\pkpA\otimes X_S$, then $\tau^Nry=\mathrm{fst}(\tau^Nr\bar y)\in\pkpA\otimes Y$, and for a form $\mu$ on $Y$
and $k\ge1$, $\mu(\varphi^k(ry))=\sum_j\hat\mu(e_j\widehat{\varphi^k(ry)})=\sum_j\nu_j(\varphi^k(r\bar y))
\in\pknBp$, with $\nu_j$ the extension of $\mu$ at $j$ just built. So $y\in\pkpA\mathcal N(Y)$; it lies in the
letter-$i$ part since $\iota$ commutes with $e_i$. \emph{Orders.} For $a\ge0$, $u\in U(i^{(-a)})$ and $\mu$ of
type $u$ at $j=i^{(-a)}$: $\widehat{\varphi^{-a}y}\in V_{Y,j}$, so $\mu(\varphi^{-a}y)=\nu(\varphi^{-a}\bar y)\in
\mathrm{Fil}^{o(a,u)}$ (the jump sets of $X_S$ are those of $Y$, Definition~\ref{pkp:def-side}).
\end{proof}

\begin{proposition}[Input (I3)]\label{pkt:I3}
Let $P$ pair sides $S,S'$ (Definition~\ref{pkp:def-pairing}), $T_P$ its Tate-dual side, $l,\kappa\ge0$, and
let $\nu$ be a point of $T_P$ over $\ell^\kappa\cdot1$ at $i^{(-l)}$ with orders $o_\nu$, where $o_\nu(a,\cdot)$
is nonincreasing for each $a$. Let $w\in(X_S)_i(o_w)$ lie over $0$ and let $y$ be a point of $S'$ over some
$\ell^{e'}c'$ at $i$ with depth $d_y$. If for all $a\ge0$, $u\in U_S(i^{(-a)})$, $u'\in U_{S'}(i^{(-a)})$
\[
0<\pkcsh_{l,0}(o_\nu)(a,-u-u'-1)+o_w(a,u)+d_y(a,u'),
\]
then the move of the first side by $w$ with $y$ at rest is $\kappa$-safe. Symmetrically, if $x$ is a point of
$S$ over $\ell^ec$ at $i$ with depth $d_x$, $w'\in(X_{S'})_i(o_{w'})$ lies over $0$, and $0<\pkcsh_{l,0}(o_\nu)
(a,-u-u'-1)+d_x(a,u)+o_{w'}(a,u')$ for all such $a,u,u'$, then the move of the second side by $w'$ with $x$ at
rest is $\kappa$-safe.
\end{proposition}

\begin{proof}
\emph{Step 1: safe pairs.} We show: if $w\in(X_S)_i(o_1)$ and $v\in(X_{S'})_i(o_2)$ both lie over $0$ and
$0<o_3(a,-u-u'-1)+o_1(a,u)+o_2(a,u')$ for all $a,u,u'$ as above, where $o_3=\pkcsh_{l,0}(o_\nu)$, then
$z\tau\,\ell^\kappa P(\mathrm{fst}\,w,\mathrm{fst}\,v)\in\pkpsm$. Since $w,v$ lie over $0$ they lie in
$\pkpAmb(Y)$, $\pkpAmb(Y')$ (the kernels of $\mathrm{pr}$), and by Lemma~\ref{pkt:sub} they have orders $o_1$,
$o_2$ there. By
Lemma~\ref{pkt:tau-v}, $\varphi^l\nu$ lies over $\ell^\kappa\cdot1$ at $i$ with orders $o_3$; $o_3(a,\cdot)$ is
nonincreasing. Let $o''(a,v_0)=1-o_3(a,-v_0-1)$. By Lemma~\ref{pkt:add}(ii) with $b=\otimes$, $w\otimes v\in
(Y\otimes Y')_i(o'')$: if $u+u'\le v_0$ then $-v_0-1\le-u-u'-1$, so $o''(a,v_0)\le1-o_3(a,-u-u'-1)\le
o_1(a,u)+o_2(a,u')$, since $0<o_3(a,-u-u'-1)+o_1(a,u)+o_2(a,u')$ is an inequality between integers. Let $\epsilon$ be the equivariant form of
Lemma~\ref{pkp:extform} on $X_{T_P}\times(Y\otimes Y')$ with values in $\auO(1)$ (jump set $\{-1\}$; the jump
sets of $T_P$ are $\{0\}\cup\{-s-1:s\in U_S(\cdot)+U_{S'}(\cdot)\}$). By Lemma~\ref{pkt:add}(ii), $\xi=
\tilde\epsilon(\varphi^l\nu,w\otimes v)\in\auO(1)_i(\mathrm{std}_0)$: we need $1=\mathrm{std}_0(a,-1)\le o_3(a,u_3)+
o''(a,v_0)$ whenever $u_3+v_0\le-1$; and then $u_3\le-v_0-1$ gives $o_3(a,u_3)\ge o_3(a,-v_0-1)=1-o''(a,v_0)$.
By Proposition~\ref{pkt:tau-i}, $\tau\xi\in\pkpA\otimes\auO(1)$, and by Lemma~\ref{pkp:AS}(vii) the image of
$z\tau\xi$ in $\pkpM$ lies in $\pkpsm$. By Lemma~\ref{pkp:extform} this image is $z\tau\,P(\mathrm{pr}(\varphi^l
\nu)(w\otimes v))$ with $\mathrm{pr}(\varphi^l\nu)=\ell^\kappa e_i(1\otimes1)$; since $e_iw=w$, this is $z\tau\,\ell^\kappa P(w,v)$
(exactly as in the proof of Proposition~\ref{pkp:rule}). So $z\tau\ell^\kappa P(\mathrm{fst}\,w,
\mathrm{fst}\,v)\in\pkpsm$.

\emph{Step 2: the first move.} Let $g\in\Gamma'$. By Lemma~\ref{pkt:tau-iv}, $v=(z\tau)^{-1}(gy-y)\in(X_{S'})_i
(d_y)$ lies over $0$ ($\mathrm{pr}(y)=\ell^{e'}e_i(1\otimes c')$ is fixed by $\Gamma'$). Step 1 with $o_1=o_w$,
$o_2=d_y$ gives $\ell^\kappa P(\mathrm{fst}\,w,\mathrm{fst}(gy-y))=z\tau\ell^\kappa P(\mathrm{fst}\,w,\mathrm{fst}\,
v)\in\pkpsm$, i.e.\ $(w,gy-y)$ is $\kappa$-safe (Definition~\ref{pkp:def-safe}).

\emph{Step 3: the second move.} Let $g\in\Gamma'$. $v=(z\tau)^{-1}(gx-x)\in(X_S)_i(d_x)$ lies over $0$
(Lemma~\ref{pkt:tau-iv}); $gw'\in(X_{S'})_i(o_{w'})$ (Lemma~\ref{pkm:inputs}(G), first assertion) and lies
over $0$. Step 1 with the roles of the two sides exchanged in the bookkeeping ($o_1=d_x$ for $v$, $o_2=
o_{w'}$ for $gw'$) gives $\ell^\kappa P(\mathrm{fst}(gx-x),\mathrm{fst}\,gw')\in\pkpsm$.
\end{proof}

\begin{remark}[Proposition~\ref{pkp:rule} as a special case]\label{pkt:rem-rule}
With $\nu$ given by Proposition~\ref{pkt:I1} (orders $\rho_{-t}$, nonincreasing in $u$) and the identity
\eqref{pkp:eq-ruleid}, the hypothesis $\mathrm{rule}_{l,t}(a,u+u')\le o(a,u)+d'(a,u')$ of
Proposition~\ref{pkp:rule} is the hypothesis of Proposition~\ref{pkt:I3}. Thus Proposition~\ref{pkp:rule} also follows
from Proposition~\ref{pkt:I3}; its direct proof uses the inequalities $u+u'\le v$ and $u+u''\le-1$, in
accordance with Lemma~\ref{pkt:add}.
\end{remark}

\subsubsection{References for the inputs}\label{pkt:ss-open}

The statements (G$_0$), ($\Phi$) used in Lemmas~\ref{pkt:tau-iv}, \ref{pkt:tau-v} are
Proposition~\ref{pkl:G0} and Corollary~\ref{pkl:Phi}; the items (R), (H), (G) of Lemma~\ref{pkm:inputs}
used in Proposition~\ref{pkt:I1} are Propositions~\ref{pkl:R}, \ref{pkl:H2}, Corollaries~\ref{pkl:H1},
\ref{pkl:H3} and Lemma~\ref{pkl:G1}. The input (I2) is Theorem~\ref{pke:I2}, (I4) is
Theorem~\ref{pkh:I4}, (I5) is Theorem~\ref{pkg:I5}, and Proposition~\ref{pki:chart} is
Theorem~\ref{pkj:chart}.



\section{Costs, the long gap and the colength of the deep constants}\label{C4:g5}
\providecommand{\auO}{\mathcal O}
\providecommand{\auQlb}{\overline{\Q}_\ell}
\providecommand{\panot}[1]{\par\smallskip{\raggedright\noindent\textbf{[Not proved here:} #1\textbf{]}\par}\smallskip}
\providecommand{\pkpA}{A_{\mathrm{inf}}}
\providecommand{\pkpW}{W(C^\flat)}
\providecommand{\pkpAmb}{\mathbb A}
\providecommand{\pknBp}{B^+_{\mathrm{dR}}}
\providecommand{\pknB}{B_{\mathrm{dR}}}
\providecommand{\pknMinf}{\mathfrak M_{\mathrm{inf}}}
\providecommand{\pknK}{\mathcal K_0}
\providecommand{\pkfBp}{B^+}
\providecommand{\pkfSB}{\mathfrak S_B}
\providecommand{\pkcPt}{\mathrm{Pt}}
\providecommand{\pkcimp}{\mathrm{imp}}
\providecommand{\pkeK}[1]{\mathcal K_{#1}}
\providecommand{\pkeAN}{\pkpA\mathcal N}
\providecommand{\pkecost}{\mathrm{cost}}
\providecommand{\pkeDp}{\mathrm{Dp}}

\subsection{Deepening with costs: the input (I2)}\label{pke:sec}

This section proves the input (I2) of \S\ref{pkc:ss-inputs} (deepening of points, with the costs
$\pkecost(a,n)=v_\ell(n)+d_v+3a$, $v_\ell(0)=0$). On the way it introduces the value fields
$\pkeK a$ and proves their properties (\S\ref{pke:ss-fields}), a version of item (C)(d) of
Lemma~\ref{pkm:inputs} at the level $a$, for all $n\ne0$ (\S\ref{pke:ss-chart}), and the depth step at
the level $a$ (\S\ref{pke:ss-step}).

\emph{Dependencies.} Theorem~\ref{pke:I2} ((I2)) is proved in full, modulo results cited by label from
\S\S\ref{pkp:sec}--\ref{pkt:sec} and \S\ref{pkl:sec}: the frame theorem (Theorem~\ref{pkf:main},
Proposition~\ref{pkn:frame}), Corollary~\ref{pkl:H1} (item (H)), Proposition~\ref{pkl:R} (item (R)),
Corollary~\ref{pkl:Phi} ($\varphi(\mathcal N)\subseteq\mathcal N$), Lemma~\ref{pkl:G1} (stability under
$\Gamma'$), Proposition~\ref{pkr:I8}(c) and Lemma~\ref{pkf:extend} (extension of
$\mathfrak S$-linear maps), Lemma~\ref{pkm:BdR} (Tate), and the statement (G$_0$) of \S\ref{pkt:sec}
together with the second assertion of Lemma~\ref{pkm:inputs}(G), which follows from (G$_0$) by
Lemma~\ref{pkt:tau-iv}. (G$_0$) is Proposition~\ref{pkl:G0}. Items (I4), (I5) and
Proposition~\ref{pki:chart} are proved elsewhere (\S\ref{pke:ss-open}).

\emph{Conventions.} We use those of \S\S\ref{pkn:sec}, \ref{pkt:sec}, \ref{pkl:sec}: the Kummer chart
(Definition~\ref{pkp:def-chart}) with $\mathfrak S=\Lambda[[z]]$, $E$, $p_a=\varphi^a(E)\in\mathfrak S$,
$\Gamma'$, $K'=\auQlb^{\,j(\Gamma')}$; $j(\Gamma')=G_{K'}$ (Definition~\ref{pkp:def-chart}: every element
of $G_{\Q_\ell}$ fixing $\Lambda$ and $\theta(z)$ lies in $j(\Gamma')$, and $K'=K_0(\pi)$); the modules
$\pknMinf(L)\subseteq\pkeAN(L)\supseteq\mathcal N(L)$ and $L_i(o)$ of Definition~\ref{pkn:def-N};
$\mathrm{Fil}^n=t^n\pknBp=E^n\pknBp$ (Lemma~\ref{pkn:BdR}). The map $\pkpA\to\pknBp$ is injective
(\S\ref{pkn:ss-frame}), and $\varphi$ is an automorphism of $\pkpA$ commuting with $\Gamma$. For
$r\in\pkpA$ and $a\ge0$ we write $r^{(a)}$ for the image of $\varphi^{-a}(r)$ in $\pknBp$; $r\mapsto r^{(a)}$
is an injective ring map, and $g(r^{(a)})=(gr)^{(a)}$ for $g\in\Gamma'$. Put $\theta_a=\theta\circ
\varphi^{-a}\colon\pkpA\to O_C$; it is a surjective, $\Gamma'$-equivariant ring map with kernel $p_a\pkpA$
(Corollary~\ref{pkn:E-gen}), and $r^{(a)}\in\mathrm{Fil}^1$ iff $\theta_a(r)=0$ iff $p_a\mid r$. Note
$p_a^{(a)}=E$, a generator of $\mathrm{Fil}^1$. For a form $\mu$ on a lattice $L$ and $x\in\pkpAmb(L)$
put $\mu_a(x)=\mu(\varphi^{-a}x)$; then $\mu_a(rx)=r^{(a)}\mu_a(x)$ for $r\in\pkpA$, and
$\mu_a(gx)=g\,\mu_a(x)$ for $x\in\pkeAN(L)$, $g\in\Gamma'$ (Lemma~\ref{pkm:equiv}). On $O_C$, $v$ is the
valuation with $v(\ell)=1$; for $x,y\in O_C$, $x\mid y$ iff $v(x)\le v(y)$. An element $x\in\pkeAN(L)$ is
\emph{deep for $o$ at $i$} if $gx-x\in L_i(o)$ for all $g\in\Gamma'$. Finally $\delta=z\tau$; its images
$\delta^{(a)}$ are units of $\pknBp$ (Lemma~\ref{pkl:basic}(i)), and $\delta$ is a unit of $\pkpW$.

\subsubsection{The value fields}\label{pke:ss-fields}

\begin{definition}[Value fields]\label{pke:def-K}
For $a\ge0$ let $R_a\subseteq\pknB$ be the $K'$-span of the elements $b\,s^{(a)}$ with $b\in\pkfSB$ (its
image in $\pknBp$) and $s\in\mathfrak S$, and let $\pkeK a$ be the subfield of $\pknB$ generated over
$K'$ by the same elements.
\end{definition}

For $a=0$ this is the field $\pknK$ of Definition~\ref{pkn:def-N}.

\begin{lemma}\label{pke:K-basic}
(i) $R_a$ is a subring of $\pknB$ and $\pkeK a$ is its field of fractions: every element of $\pkeK a$ is
$w/d$ with $w,d\in R_a$, $d\ne0$. (ii) $K'\subseteq\pkeK a$; $s^{(a)}\in\pkeK a$ for $s\in\mathfrak S$, in
particular $E=p_a^{(a)}\in\pkeK a$; and $\pknK\subseteq\pkeK a$.
\end{lemma}

\begin{proof}
(i) $(b\,s^{(a)})(b'\,s'^{(a)})=(bb')(ss')^{(a)}$ with $bb'\in\pkfSB$ (a ring) and $ss'\in\mathfrak S$; and
$K'$ consists of scalars. So $R_a$ is a $K'$-algebra, a domain inside the field $\pknB$, and the field it
generates is $\{w/d\}$. (ii) Take $b=1$, resp.\ $s=1$, $b=1$; $p_a\in\mathfrak S$ as $\varphi(\mathfrak S)
\subseteq\mathfrak S$. The generators $b\,s$ ($b\in\pkfSB$, $s\in\mathfrak S$) of $\pknK$ satisfy
$bs=(bs)\cdot1^{(a)}$ with $bs\in\pkfSB$, since the image of $\mathfrak S$ lies in $\pkfSB$
(\S\ref{pkf:ss-SB}).
\end{proof}

\begin{lemma}[Retraction of values at level $a$]\label{pke:valret}
Let $a\ge0$, $f\colon\pkpA\to\pkpA$ $\mathfrak S$-linear, $v_1,\dots,v_n\in\pkeK a$, $r_1,\dots,r_n\in\pkpA$,
and $W\subseteq\pknB$ a $\pknBp$-submodule. If $\sum_lr_l^{(a)}v_l\in W$, then $\sum_lf(r_l)^{(a)}v_l\in W$.
\end{lemma}

\begin{proof}
Put $f'(x)=\varphi^{-a}(f(\varphi^ax))$. It is additive, and for $s\in\mathfrak S$, $f'(sx)=
\varphi^{-a}(\varphi^a(s)f(\varphi^ax))=sf'(x)$ because $\varphi^a(s)\in\mathfrak S$; so $f'$ is $\mathfrak
S$-linear, and Lemma~\ref{pkf:extend} gives its extension $G_{f'}$ to $\pkfBp$, with $G_{f'}(bx)=bG_{f'}(x)$
for $b\in\pkfSB$, $x\in\pkfBp$. By Lemma~\ref{pke:K-basic}(i) choose $d\in R_a\smallsetminus\{0\}$ with
$w_l=dv_l\in R_a$, and write $w_l=\sum_ic_{li}b_{li}s_{li}^{(a)}$ (finite sums, $c_{li}\in K'$,
$b_{li}\in\pkfSB$, $s_{li}\in\mathfrak S$). Then
\[
\textstyle\sum_lw_lr_l^{(a)}=\sum_{l,i}c_{li}\,y_{li},\qquad y_{li}=b_{li}\,\varphi^{-a}(s_{li}r_l)\in\pkfBp,
\]
where $\varphi^{-a}(s_{li}r_l)\in\pkpA\subseteq\pkfBp$ and images are taken in $\pknBp$. Now
$G_{f'}(y_{li})=b_{li}f'(\varphi^{-a}(s_{li}r_l))=b_{li}\varphi^{-a}(f(s_{li}r_l))=b_{li}\varphi^{-a}(s_{li})
\varphi^{-a}(f(r_l))$, by Lemma~\ref{pkf:extend} and $\mathfrak S$-linearity of $f$. By
Proposition~\ref{pkr:I8}(c) there is $u\in\pknBp$ with $\sum c_{li}G_{f'}(y_{li})=u\sum c_{li}y_{li}$, i.e.\
$\sum_lw_lf(r_l)^{(a)}=u\sum_lw_lr_l^{(a)}$. Dividing by $d$: $\sum_lf(r_l)^{(a)}v_l=u\sum_lr_l^{(a)}v_l\in
W$.
\end{proof}

\begin{lemma}[Values lie in $\pkeK a$]\label{pke:values}
Let $L$ be a lattice with jump sets, $x\in\mathcal N(L)$, $\mu$ a form and $a\ge0$. Then $\mu_a(x)\in\pkeK
a$.
\end{lemma}

\begin{proof}
By Lemma~\ref{pkn:M}(ii) there is $r'\ne0$ with $r'x\in\pkpA\otimes L$; then $r\varphi^{-a}x\in\pkpA\otimes L$
with $r=\varphi^{-a}(r')\ne0$. By Corollary~\ref{pkl:H1}, $\tau^H(\pkpA\otimes L)\subseteq\pkeAN(L)$ ($H$ the
jump bound), and $\pkeAN(L)$ is the $\pkpA$-span of $\mathcal N(L)$ (Proposition~\ref{pkn:frame}(iv) with
Theorem~\ref{pkf:main}). So $\tau^Hr\varphi^{-a}x=\sum_{l=1}^n t_lv_l$ with $t_l\in\pkpA$, $v_l\in\mathcal N(L)$.
Apply $\varphi^a$: with $\rho=\varphi^a(\tau^Hr)\ne0$,
\[
\textstyle-\rho\,x+\sum_l\varphi^a(t_l)\,\varphi^a(v_l)=0,
\]
a relation with coefficients in $\pkpA$ among elements of $\mathcal N(L)$: $x\in\mathcal N(L)$ and
$\varphi^a(v_l)\in\mathcal N(L)$ (Corollary~\ref{pkl:Phi}). By Proposition~\ref{pkr:I5} some coordinate
$\kappa_j$ has $\kappa_j(\rho)\ne0$; it is $\mathfrak S$-linear with values in $\mathfrak S$. By
Lemma~\ref{pkl:rel} (with $f=\kappa_j$), $\kappa_j(\rho)x=\sum_l\kappa_j(\varphi^at_l)\varphi^a(v_l)$. Apply
$\mu_a$: $\kappa_j(\rho)^{(a)}\mu_a(x)=\sum_l\kappa_j(\varphi^at_l)^{(a)}\mu(v_l)$. Here $\mu(v_l)\in\pknK
\subseteq\pkeK a$ (as $v_l\in\mathcal N$), $\kappa_j(\cdot)^{(a)}\in\pkeK a$ (values in $\mathfrak S$,
Lemma~\ref{pke:K-basic}(ii)), and $\kappa_j(\rho)^{(a)}\ne0$. So $\mu_a(x)\in\pkeK a$.
\end{proof}

\begin{lemma}[Galois action on integral values]\label{pke:fix}
Let $a\ge0$, $g,t_1,\dots,t_m\in\Gamma'$, and suppose $p_a$ divides $g(\sum_lt_ls)-\sum_lt_ls$ in $\pkpA$
for every $s\in\mathfrak S$. If $y\in\pkeK a\cap\mathrm{Fil}^0$, then $g(\sum_lt_ly)-\sum_lt_ly\in
\mathrm{Fil}^1$.
\end{lemma}

\begin{proof}
$\theta(y)\in C$, so $\ell^e\theta(y)\in O_C$ for some $e\ge0$; as $\theta\colon\pkpA\to O_C$ is onto,
choose $r\in\pkpA$ with $\theta(r)=\ell^e\theta(y)$, i.e.\ $\ell^ey-r\in\ker(\theta|\pknBp)=\mathrm{Fil}^1$.
Since the image of $r$ in $\pknBp$ is $(\varphi^ar)^{(a)}$ and $(\ell^e)^{(a)}=\ell^e$, this reads
$(\ell^e)^{(a)}y+(\varphi^ar)^{(a)}\cdot(-1)\in\mathrm{Fil}^1$. Lemma~\ref{pke:valret} with $f=\kappa_0$ (Proposition~\ref{pkr:I5}; $\kappa_0(\ell^e)=
\ell^e$), the values $y,-1\in\pkeK a$, the coefficients $\ell^e,\varphi^a(r)$ and $W=\mathrm{Fil}^1$ gives
$\ell^ey-s^{(a)}\in\mathrm{Fil}^1$ with $s=\kappa_0(\varphi^ar)\in\mathfrak S$. Elements of $\Gamma'$ preserve
$\mathrm{Fil}^1$ ($g(t)=\chi(g)t$), so $\sum_lt_l(\ell^ey-s^{(a)})$ and its image under $g$ lie in
$\mathrm{Fil}^1$. Also $g(\sum_lt_ls^{(a)})-\sum_lt_ls^{(a)}=\bigl(g(\sum_lt_ls)-\sum_lt_ls\bigr)^{(a)}\in
(p_a\pkpA)^{(a)}\subseteq E\pknBp=\mathrm{Fil}^1$. Adding, $\ell^e\bigl(g(\sum_lt_ly)-\sum_lt_ly\bigr)\in
\mathrm{Fil}^1$, and $\ell$ is a unit of $\pknBp$.
\end{proof}

\begin{lemma}[The retraction and values]\label{pke:ret}
Let $\mathrm{ret}\colon\pkeAN(L)\to\mathcal N(L)$ be the retraction of Proposition~\ref{pkl:R}. Let $y\in
\pkeAN(L)$, $\mu$ a form, $a\ge0$, $v\in\pkeK a$ and $W\subseteq\pknB$ a $\pknBp$-submodule. If
$\mu_a(y)-v\in W$, then $\mu_a(\mathrm{ret}\,y)-v\in W$. Consequently $\mathrm{ret}(L_i(o))\subseteq
\mathcal N(L)\cap L_i(o)$ for every letter $i$ and all orders $o$.
\end{lemma}

\begin{proof}
Write $y=\sum_sr_sw_s$ with $r_s\in\pkpA$, $w_s\in\mathcal N(L)$ (Proposition~\ref{pkn:frame}(iv)); then
$\mathrm{ret}\,y=\sum_s\kappa_0(r_s)w_s$ (Proposition~\ref{pkl:R}). The values $v_s=\mu_a(w_s)$ lie in
$\pkeK a$ (Lemma~\ref{pke:values}), and $\mu_a(y)-v=\sum_sr_s^{(a)}v_s+1^{(a)}(-v)\in W$. Lemma~\ref{pke:valret}
with $f=\kappa_0$ ($\kappa_0(1)=1$) gives $\sum_s\kappa_0(r_s)^{(a)}v_s-v=\mu_a(\mathrm{ret}\,y)-v\in W$. For
the consequence let $y\in L_i(o)$: $\mathrm{ret}\,y\in\mathcal N(L)\subseteq\pkeAN(L)$, it lies in the
letter-$i$ part as $\mathrm{ret}$ commutes with $e_i$ (Proposition~\ref{pkl:R}), and for $\mu$ of type $u$
at $i^{(-a)}$ take $v=0$, $W=\mathrm{Fil}^{o(a,u)}$.
\end{proof}

\subsubsection{Chart identities at the level $a$}\label{pke:ss-chart}

\begin{lemma}\label{pke:triv}
Let $a\ge0$, $m\ge0$ with $\ell^a\mid m$, and $g\in\Gamma'$ with $gz=[\varepsilon]^mz$. Then $p_a\mid
gs-s$ for every $s\in\mathfrak S$.
\end{lemma}

\begin{proof}
$gz-z=z([\varepsilon]^m-1)$, and $\varphi^{-a}([\varepsilon]^m)=[\varepsilon^{m/\ell^a}]$ has
$\theta_a([\varepsilon]^m)=(\varepsilon^\sharp)^{m/\ell^a}=1$. So $\theta_a(gz-z)=0$ and $p_a\mid gz-z$. By the
argument of Step~2 of the proof of Proposition~\ref{pkn:Ca} (with $\eta=\varepsilon^m$), $gz-z$ divides
$gs-s$ for every $s\in\mathfrak S$.
\end{proof}

\begin{lemma}[Averaging]\label{pke:avg}
Let $a\ge0$, $c\ge a$, $N=\ell^c$, and for $0\le l<N$ let $t_l\in\Gamma'$ with $t_lz=[\varepsilon]^lz$
(Definition~\ref{pkp:def-chart}). Then for every $g\in\Gamma'$ and $s\in\mathfrak S$, $p_a$ divides
$g(\sum_{l<N}t_ls)-\sum_{l<N}t_ls$.
\end{lemma}

\begin{proof}
It suffices that $\theta_a(\sum_lt_ls)$ is fixed by $j(\Gamma')=G_{K'}$, since $\theta_a$ is
$\Gamma'$-equivariant with kernel $p_a\pkpA$. Write $s=\sum_k\lambda_kz^k$, $\lambda_k\in\Lambda$. Put
$\varpi_a=\theta_a(z)$, an $\ell^a$-th root of $\pi$, and $\zeta_a=\theta_a([\varepsilon])=
\theta([\varepsilon^{1/\ell^a}])$, a primitive $\ell^a$-th root of unity. As $t_l$ fixes $\Lambda$ and is
continuous, $t_ls=\sum_k\lambda_k[\varepsilon]^{lk}z^k$ (as in Step~2 of the proof of
Proposition~\ref{pkn:Ca}); $\theta_a$ is continuous for the $(\ell,z)$-adic topology ($v(\varpi_a)>0$), and
$\theta_a(\lambda_k)=\theta(\varphi^{-a}\lambda_k)\in W(\mathbb F_{\ell^f})\subseteq K'$. Hence
\[
\textstyle\theta_a\bigl(\sum_{l<N}t_ls\bigr)=\sum_k\theta_a(\lambda_k)\varpi_a^k\sum_{l<N}\zeta_a^{lk}
=N\sum_{j\ge0}\theta_a(\lambda_{j\ell^a})\pi^j,
\]
because $\sum_{l<N}\zeta_a^{lk}$ is $N$ if $\ell^a\mid k$ and $0$ otherwise ($\ell^a\mid N$). The last series
converges in $O_C$ with partial sums in the complete ring $O_{K'}$, so its sum lies in $K'$ and is fixed by
$G_{K'}$.
\end{proof}

\begin{proposition}[Item (C)(d) at the level $a$]\label{pke:Cd}
Let $a\ge0$ and $n\in\Z\smallsetminus\{0\}$, and put $c=\pkecost(a,n)=v_\ell(n)+d_v+3a$. There are $g\in
\Gamma'$, $\omega'\in\pkpA^\times$ and $\upsilon,X\in\pkpA$ with
\[
gp_a=\omega'p_a,\qquad p_a\mid gs-s\ \ (s\in\mathfrak S),\qquad p_a\mid\upsilon X-\ell^c,\qquad
\delta X=\omega'^n-1 .
\]
\end{proposition}

\begin{proof}
Let $g\in\Gamma'$ with $gz=[\varepsilon]^{\ell^a}z$ (Definition~\ref{pkp:def-chart}), and $gE=\omega E$ with
$\omega\in\pkpA^\times$ (Lemma~\ref{pkf:approx0}(iv)). Put $\omega'=\varphi^a(\omega)$; then $gp_a=
\varphi^a(gE)=\omega'p_a$ ($g$ commutes with $\varphi$). The second property is Lemma~\ref{pke:triv}.

\emph{Step 1: $\delta\mid\omega-1$.} By Proposition~\ref{pkn:Ca}, $\delta E\mid gE-E=E(\omega-1)$; cancel $E$.
Hence $\delta\mid\omega^k-1$ for $k\ge1$, and $\omega^{-k}-1=-\omega^{-k}(\omega^k-1)$; so $\delta\mid
\omega^n-1$.

\emph{Step 2: $v(\theta(\omega-1))\le d_v+a+v(\theta(\delta))$, and $\theta(\omega-1)\ne0$.} Let $P,V$ be as in
Lemma~\ref{pkl:different}, $P(z)=EG$ with $G=V(z)^{-1}\in\mathfrak S^\times$. By Lemma~\ref{pkn:mu},
$gz-z=z([\varepsilon]-1)w=z\tau Euw=Eh$, with $w=\sum_{k<\ell^a}[\varepsilon]^k$, $u\in\pkpA^\times$ and $h=
\delta uw$. Taylor's formula for the polynomial $P$ gives $c'\in\pkpA$ with $P(gz)=P(z)+P'(z)Eh+E^2h^2c'$. As $g$
fixes $\Lambda$, $P(gz)=g(EG)=\omega Eg(G)$; cancelling $E$, $\omega g(G)=G+P'(z)h+Eh^2c'$. Apply $\theta$:
$\theta(gG)=\theta(G)$ (as in Step~2 of the proof of Proposition~\ref{pkl:Cd}), $\theta(E)=0$, $\theta(w)=\ell^a$
($\theta([\varepsilon])=1$). So $\theta(\omega-1)\theta(G)=P'(\pi)\theta(\delta)\theta(u)\ell^a$, with
$\theta(G),\theta(u)$ units, $P'(\pi)\ne0$ ($P$ is separable) and $v(P'(\pi))\le d_v$
(Lemma~\ref{pkl:different}); and $\theta(\delta)=\pi(\zeta_\ell-1)\ne0$.

\emph{Step 3: $v(\theta(\omega^n-1))=v_\ell(n)+v(\theta(\omega-1))$.} For $n\ge1$ this is Step~3 of the proof
of Proposition~\ref{pkl:Cd}, whose argument uses only $\delta\mid\omega-1$ and $\theta(\omega-1)\ne0$ (Steps 1
and 2 here). For $n<0$, $\omega^n-1=-\omega^n(\omega^{|n|}-1)$ and $\omega^n$ is a unit.

\emph{Step 4: $X$ and $\upsilon$.} By Lemma~\ref{pkl:basic}(i), $\varphi^a(\delta)=z^{\ell^a}\varphi^a(\tau)=
z^{\ell^a}\tau\prod_{b<a}p_b\cdot w_a$, so $\delta\mid\varphi^a(\delta)$. With $\omega^n-1=\delta q$ (Step 1),
$\omega'^n-1=\varphi^a(\delta)\varphi^a(q)=\delta X$, $X=(\varphi^a(\delta)/\delta)\varphi^a(q)\in\pkpA$. Apply
$\theta_a$: $\theta_a(\delta)\theta_a(X)=\theta_a(\varphi^a(\omega^n-1))=\theta(\omega^n-1)$. Also $v(\theta(\delta))
\le2a+v(\theta_a(\delta))$: for $a=0$ trivially, and for $a\ge1$ because $v(\theta(\delta))=v(\pi)+v(\zeta_\ell-1)
\le2$. By Steps 2 and 3,
\[
v(\theta_a(\delta))+v(\theta_a(X))\le v_\ell(n)+d_v+a+v(\theta(\delta))\le c+v(\theta_a(\delta)),
\]
and $\theta_a(\delta)\ne0$ (Lemma~\ref{pkl:basic}(i)); so $\theta_a(X)\mid\ell^c$ in $O_C$, say
$\ell^c=\theta_a(X)\gamma$. Choose $\upsilon_0\in\pkpA$ with $\theta(\upsilon_0)=\gamma$ and put $\upsilon=
\varphi^a(\upsilon_0)$, so $\theta_a(\upsilon)=\gamma$. Then $\theta_a(\upsilon X-\ell^c)=0$, i.e.\ $p_a\mid\upsilon X-
\ell^c$.
\end{proof}

\subsubsection{The depth step at the level $a$}\label{pke:ss-step}

For a lattice with jump sets $L$ (Definition~\ref{pkp:def-orders}), $\upsilon\in\pkpA$, $e\ge0$ and $t_1,\dots,t_m\in\Gamma'$ put, for $x\in\mathcal N(L)$,
\[
D_{e,\upsilon,t}(x)=\ell^ex+\upsilon\sum_{l=1}^m\delta^{-1}(t_lx-x)\ \in\pkeAN(L);
\]
it lies in $\pkeAN(L)$ by (G$_0$) of \S\ref{pkt:sec} and Lemma~\ref{pkn:M}(i). Note
$\mu_a(\delta^{-1}w)=(\delta^{(a)})^{-1}\mu_a(w)$ for $w\in\pkpAmb(L)$ ($\mu$ is $\pkpA$-linear and
$\delta^{(a)}\ne0$).

\begin{proposition}[The depth step]\label{pke:step}
Let $a\ge0$, $n\in\Z$, $c=\pkecost(a,n)$. There are $m\ge1$, $t_1,\dots,t_m\in\Gamma'$ and $\upsilon\in\pkpA$,
depending only on $(a,n)$ and the chart, such that for every lattice with jump sets $L$, every $x\in\mathcal N(L)$ and every
form $\mu$ on $L$: if $\mu_a(x)-\gamma\in\mathrm{Fil}^n$ for some $\gamma\in K'$, then
$\mu_a(D_{c,\upsilon,t}(x))-\gamma'\in\mathrm{Fil}^{n+1}$ for some $\gamma'\in K'$.
\end{proposition}

\begin{proof}
\emph{Case $n\ne0$.} Take $g,\omega',\upsilon_1,X$ from Proposition~\ref{pke:Cd}, $m=1$, $t_1=g$ and
$\upsilon=-\upsilon_1$. Put $y=\mu_a(x)-\gamma\in\mathrm{Fil}^n$; $y\in\pkeK a$ by Lemmas~\ref{pke:values},
\ref{pke:K-basic}(ii). \emph{Eigenvalue.} $w=E^{-n}y$ lies in $\mathrm{Fil}^0$ (Lemma~\ref{pkn:BdR}) and in
$\pkeK a$ ($E\in\pkeK a$). By Lemma~\ref{pke:fix} ($m=1$, $t_1=1$; the hypothesis is the second property in
Proposition~\ref{pke:Cd}), $gw-w\in\mathrm{Fil}^1$. Since $E=p_a^{(a)}$, $gE=(gp_a)^{(a)}=\omega'^{(a)}E$, so
$gy=(\omega'^{(a)})^nE^ng(w)$ and
\[
gy-(\omega'^n)^{(a)}y=(\omega'^{(a)})^nE^n(gw-w)\in\mathrm{Fil}^{n+1}.
\]
\emph{Cancellation.} Let $w'=\delta^{-1}(gx-x)\in\pkeAN(L)$. Then $\delta^{(a)}\mu_a(w')=\mu_a(gx-x)=g\mu_a(x)-
\mu_a(x)=gy-y$ ($g$ fixes $\gamma\in K'$), and $\delta^{(a)}X^{(a)}=(\omega'^n)^{(a)}-1$. Hence
$\delta^{(a)}\bigl(\mu_a(w')-X^{(a)}y\bigr)=gy-(\omega'^n)^{(a)}y\in\mathrm{Fil}^{n+1}$, and as $\delta^{(a)}$ is a
unit of $\pknBp$, $\mu_a(w')-X^{(a)}y\in\mathrm{Fil}^{n+1}$. \emph{Conclusion.} $D(x)=\ell^cx-\upsilon_1w'$ and,
with $\gamma'=\ell^c\gamma\in K'$,
\[
\mu_a(D(x))-\gamma'=\ell^cy-\upsilon_1^{(a)}\mu_a(w')=-\upsilon_1^{(a)}\bigl(\mu_a(w')-X^{(a)}y\bigr)-y\,
(\upsilon_1X-\ell^c)^{(a)} .
\]
The first term lies in $\mathrm{Fil}^{n+1}$ ($\upsilon_1^{(a)}\in\mathrm{Fil}^0$); in the second, $y\in
\mathrm{Fil}^n$ and $(\upsilon_1X-\ell^c)^{(a)}\in(p_a\pkpA)^{(a)}\subseteq\mathrm{Fil}^1$.

\emph{Case $n=0$.} Let $N=\ell^c$ ($c=d_v+3a\ge a$), $m=N$, and $t_{l+1}$ ($0\le l<N$) the elements $t_l$ of Lemma~\ref{pke:avg}, and
$\upsilon=\delta$. Then $D(x)=\ell^cx+\sum_l(t_lx-x)=\sum_lt_lx$. The hypothesis gives $\mu_a(x)\in\mathrm{Fil}^0$
($\gamma\in K'\subseteq\mathrm{Fil}^0$), and $\mu_a(x)\in\pkeK a$ (Lemma~\ref{pke:values}). As
$\mu_a(D(x))=\sum_lt_l\mu_a(x)$, Lemmas~\ref{pke:avg} and \ref{pke:fix} give $g\mu_a(D(x))-\mu_a(D(x))\in
\mathrm{Fil}^1$ for every $g\in\Gamma'$, and Lemma~\ref{pkm:BdR} gives $\gamma'\in K'$ with
$\mu_a(D(x))-\gamma'\in\mathrm{Fil}^1$.
\end{proof}

\subsubsection{Operators of a cost, and the proof of (I2)}\label{pke:ss-ops}

Fix a side $S$, with extension $X_S$ and projection $\mathrm{pr}$; $\mathcal N=\mathcal N(X_S)$, $L_i(o)=
(X_S)_i(o)$, and $L_i^0(o)$ the elements of $L_i(o)$ over $0$. A \emph{constant at $i$} is an element
$e_i(1\otimes c)$, $c\in C_S$; these form a group stable under multiplication by integers, and $\Gamma'$ fixes
them (Definition~\ref{pkp:def-letters}, $\Gamma$ acting trivially on $C_S$). For a letter $i$ and orders $o$ let
$\pkeDp_i(o)$ be the set of $x\in\mathcal N$ that lie over a constant at $i$ and are deep for $o$ at $i$. An
additive map $A\colon\mathcal N\to\mathcal N$ is an \emph{operator of cost $e$} if for every letter $i$, all
orders $o$ and every $x\in\pkeDp_i(o)$: $Ax-\ell^ex\in L_i^0(o)$.

\begin{lemma}\label{pke:op-basic}
(i) An operator $A$ of cost $e$ maps $\pkeDp_i(o)$ into itself, and $\mathrm{pr}(Ax)=\ell^e\mathrm{pr}(x)$
there. (ii) If $A$, $B$ are operators of costs $e$, $e'$, then $B\circ A$ is an operator of cost $e+e'$.
(iii) The identity is an operator of cost $0$.
\end{lemma}

\begin{proof}
(i) Let $x\in\pkeDp_i(o)$ and $w=Ax-\ell^ex\in L_i^0(o)$. Then $\mathrm{pr}(Ax)=\ell^e\mathrm{pr}(x)$, a constant
at $i$, and $g(Ax)-Ax=\ell^e(gx-x)+(gw-w)\in L_i(o)$ for $g\in\Gamma'$, as $L_i(o)$ is a $\Gamma'$-stable
$\pkpA$-module (Lemma~\ref{pkl:G1}). (ii) $BAx-\ell^{e+e'}x=(B(Ax)-\ell^{e'}Ax)+\ell^{e'}(Ax-\ell^ex)$; the
first term is in $L_i^0(o)$ by (i), the second by hypothesis. (iii) is clear.
\end{proof}

\begin{lemma}[One step]\label{pke:B}
For every $a\ge0$ and $n\in\Z$ there is an operator $B$ of cost $\pkecost(a,n)$ with the following property:
for $x\in\mathcal N$ and a form $\mu$ on $X_S$, if $\mu_a(gx-x)\in\mathrm{Fil}^n$ for all $g\in\Gamma'$, then
$\mu_a(gBx-Bx)\in\mathrm{Fil}^{n+1}$ for all $g\in\Gamma'$.
\end{lemma}

\begin{proof}
Let $c,m,t,\upsilon$ be as in Proposition~\ref{pke:step} and put $Bx=\mathrm{ret}\bigl(D_{c,\upsilon,t}(x)\bigr)$;
$B$ is additive with values in $\mathcal N$. Since $\mathrm{ret}$ is $\mathfrak S$-linear and the identity on
$\mathcal N$ (Proposition~\ref{pkl:R}), $Bx-\ell^cx=\sum_l\mathrm{ret}\bigl(\upsilon\delta^{-1}(t_lx-x)\bigr)$.

\emph{Cost.} Let $x\in\pkeDp_i(o)$. By Lemma~\ref{pkm:inputs}(G) (second assertion; Lemma~\ref{pkt:tau-iv}),
$\delta^{-1}(t_lx-x)\in L_i(o)$, hence $\upsilon\delta^{-1}(t_lx-x)\in L_i(o)$, and by Lemma~\ref{pke:ret} its
retraction lies in $\mathcal N\cap L_i(o)$. It lies over $0$: $\mathrm{pr}$ is $\pkpW$-linear,
$\Gamma'$-equivariant and commutes with $\mathrm{ret}$ (Proposition~\ref{pkl:R}), and $t_l$ fixes the constant
$\mathrm{pr}(x)$. So $Bx-\ell^cx\in L_i^0(o)$.

\emph{Raising.} Let $x,\mu$ be as in the statement. By Lemma~\ref{pkm:equiv}, $g\mu_a(x)-\mu_a(x)=\mu_a(gx-x)\in
\mathrm{Fil}^n$ for all $g$, so by Lemma~\ref{pkm:BdR} $\mu_a(x)-\gamma\in\mathrm{Fil}^n$ for some $\gamma\in K'$.
By Proposition~\ref{pke:step}, $\mu_a(D_{c,\upsilon,t}(x))-\gamma'\in\mathrm{Fil}^{n+1}$ with $\gamma'\in K'\subseteq
\pkeK a$; by Lemma~\ref{pke:ret} ($W=\mathrm{Fil}^{n+1}$), $\mu_a(Bx)-\gamma'\in\mathrm{Fil}^{n+1}$. As $g$ fixes
$\gamma'$ and preserves $\mathrm{Fil}^{n+1}$, $\mu_a(gBx-Bx)=g(\mu_a(Bx)-\gamma')-(\mu_a(Bx)-\gamma')\in
\mathrm{Fil}^{n+1}$ (Lemma~\ref{pkm:equiv}).
\end{proof}

For orders $o,o'$ and a letter $i$ say that a set $I$ of pairs $(a,n)\in\Z_{\ge0}\times\Z$ \emph{covers
$o\to o'$ at $i$} if it contains every $(a,n)$ with $o(a,u)\le n<o'(a,u)$ for some $u\in U(i^{(-a)})$.

\begin{proposition}[Deepening operators]\label{pke:A}
For every finite set $I$ of pairs $(a,n)$ there is an operator $A_I$ of cost $\sum_{(a,n)\in I}\pkecost(a,n)$
such that $A_Ix\in\pkeDp_i(o')$ whenever $x\in\pkeDp_i(o)$ and $I$ covers $o\to o'$ at $i$.
\end{proposition}

\begin{proof}
Induction on $|I|$. If $I=\emptyset$, take $A_I=\mathrm{id}$ (Lemma~\ref{pke:op-basic}(iii)); if $\emptyset$
covers $o\to o'$, then $o'(a,u)\le o(a,u)$ on the jumps, and $\pkeDp_i(o)\subseteq\pkeDp_i(o')$
(Lemma~\ref{pkm:le}(i)). Let $I\ne\emptyset$, choose $(a,n)\in I$ with $n$ maximal among the second coordinates
of the elements of $I$, put $I'=I\smallsetminus\{(a,n)\}$, let $A_{I'}$ be given by induction and $B$ by
Lemma~\ref{pke:B} for $(a,n)$, and $A_I=B\circ A_{I'}$, an operator of the right cost by
Lemma~\ref{pke:op-basic}(ii).

Let $x\in\pkeDp_i(o)$ with $I$ covering $o\to o'$. Define $o_-(b,u)=n$ if $b=a$ and $o'(b,u)=n+1$, and
$o_-(b,u)=o'(b,u)$ otherwise. \emph{$I'$ covers $o\to o_-$.} Let $o(b,u)\le k<o_-(b,u)$ with $u\in
U(i^{(-b)})$. Since $o_-\le o'$, $(b,k)\in I$. If $(b,k)=(a,n)$, then $o_-(a,u)>n$ forces $o'(a,u)\ne n+1$, and
$o'(a,u)\ge o_-(a,u)>n$, so $o'(a,u)\ge n+2$; then $o(a,u)\le n+1<o'(a,u)$ and $(a,n+1)\in I$, contradicting the
maximality of $n$. So $(b,k)\in I'$. Hence $y=A_{I'}x\in\pkeDp_i(o_-)$, and $By\in\pkeDp_i(o_-)$
(Lemma~\ref{pke:op-basic}(i)). \emph{Raising.} Let $g\in\Gamma'$; $gBy-By\in L_i(o_-)$, and it remains to see
$\mu_b(gBy-By)\in\mathrm{Fil}^{o'(b,u)}$ where $o'(b,u)>o_-(b,u)$, i.e.\ $b=a$, $o'(a,u)=n+1$, for $\mu$ of type
$u\in U(i^{(-a)})$ at $i^{(-a)}$. Since $y$ is deep for $o_-$ and $o_-(a,u)=n$, $\mu_a(\gamma y-y)\in\mathrm{Fil}^n$
for all $\gamma\in\Gamma'$, and Lemma~\ref{pke:B} gives $\mu_a(gBy-By)\in\mathrm{Fil}^{n+1}$.
\end{proof}

\begin{theorem}[Input (I2)]\label{pke:I2}
Let $x\in\pkcPt_S(c,i,e;o,d)$ and let $I$ be a finite set of pairs $(a,n)\in\Z_{\ge0}\times\Z$ containing every
$(a,n)$ with $d(a,u)\le n<d'(a,u)$ for some $u\in U(i^{(-a)})$. If $e+\sum_{(a,n)\in I}\pkecost(a,n)\le e'$,
then $\pkcPt_S(c,i,e';\pkcimp(o,d'),d')\ne\emptyset$.
\end{theorem}

\begin{proof}
Let $\Sigma=\sum_I\pkecost(a,n)$, $s=e'-e-\Sigma\ge0$, $y=\ell^sx$ and $A=A_I$ (Proposition~\ref{pke:A}). Then
$y\in\mathcal N$, $\mathrm{pr}(y)=\ell^{e+s}e_i(1\otimes c)$, $y\in L_i(o)$, and $y$ is deep for $d$ and, as
$L_i(o)$ is $\Gamma'$-stable, for $o$; so $y\in\pkeDp_i(d)\cap\pkeDp_i(o)$. By Proposition~\ref{pke:A}, $Ay\in
\pkeDp_i(d')$. As $A$ has cost $\Sigma$ and $y\in\pkeDp_i(o)$, $Ay-\ell^\Sigma y\in L_i(o)$, so $Ay\in L_i(o)$;
and $\mathrm{pr}(Ay)=\ell^\Sigma\mathrm{pr}(y)=\ell^{e'}e_i(1\otimes c)$ (Lemma~\ref{pke:op-basic}(i)). By
Lemma~\ref{pkm:improved}, $Ay\in L_i(\pkcimp(o,d'))$. So $Ay\in\pkcPt_S(c,i,e';\pkcimp(o,d'),d')$.
\end{proof}

\begin{remark}\label{pke:rem-use}
The operators depend only on the chart and $I$ (through the elements $t_l$, $\upsilon$ of
Proposition~\ref{pke:step}), not on $c$, $i$, $o$; this is not needed for (I2). The statement of (I2) in
\S\ref{pkc:ss-inputs} is proved as stated. Lemma~\ref{pkc:free} follows from it as there.
\end{remark}

\subsubsection{The other inputs}\label{pke:ss-open}

The input (I4) of \S\ref{pkc:ss-inputs} is Theorem~\ref{pkh:I4}; the input (I5) of \S\ref{pkd:ss-tools}
is Theorem~\ref{pkg:I5}; Proposition~\ref{pki:chart} is Theorem~\ref{pkj:chart}.



\providecommand{\auO}{\mathcal O}
\providecommand{\auQlb}{\overline{\Q}_\ell}
\providecommand{\panot}[1]{\par\smallskip{\raggedright\noindent\textbf{[Not proved here:} #1\textbf{]}\par}\smallskip}
\providecommand{\pkpA}{A_{\mathrm{inf}}}
\providecommand{\pkpW}{W(C^\flat)}
\providecommand{\pkpAmb}{\mathbb A}
\providecommand{\pknBp}{B^+_{\mathrm{dR}}}
\providecommand{\pknB}{B_{\mathrm{dR}}}
\providecommand{\pknMinf}{\mathfrak M_{\mathrm{inf}}}
\providecommand{\pknK}{\mathcal K_0}
\providecommand{\pkfAm}{A_{\max}}
\providecommand{\pkfAo}{A_{\max}^{\circ}}
\providecommand{\pkfBp}{B^+}
\providecommand{\pkfSB}{\mathfrak S_B}
\providecommand{\pkcPt}{\mathrm{Pt}}
\providecommand{\pkgAN}{\pkpA\mathcal N}
\providecommand{\pkgP}{\mathcal P}
\providecommand{\pkgC}{\mathcal C}
\providecommand{\pkgV}{\mathcal V}
\providecommand{\pkgSh}{\widehat{\mathfrak S}}
\providecommand{\pkgstd}{\mathrm{std}}

\subsection{The long gap: the input (I5)}\label{pkg:sec}

This section proves the input (I5) of \S\ref{pkd:ss-tools} (long gap), in the form printed in
\S\ref{pkd:ss-tools}.

\emph{Dependencies.} Theorem~\ref{pkg:I5} ((I5)) is proved in full, modulo results cited by label from
\S\S\ref{pkp:sec}--\ref{pke:sec}: the frame theorem with its Galois relation
(Theorems~\ref{pkf:main}, \ref{pkf:frame-pm}, Lemma~\ref{pkn:frame-id}, Proposition~\ref{pkn:frame}),
the facts on $\pkfAm$, $\pkfBp$, $\pkfSB$ of \S\S\ref{pkf:sec}--\ref{pkr:sec} and \S\ref{pkl:ss-open},
item (I1) of \S\ref{pkc:ss-inputs} (Proposition~\ref{pkt:I1}), the retraction (Proposition~\ref{pkl:R}),
improved orders (Lemma~\ref{pkm:improved}) and the lift ideals (Definition~\ref{pkm:def-J},
Lemma~\ref{pkm:J-basic}). The input (I4) and Proposition~\ref{pki:chart} are proved elsewhere
(\S\ref{pkg:ss-open}).

\emph{Conventions.} Those of \S\ref{pke:sec}. Throughout, the chart is \emph{unramified}: $e=1$ and
$E=z+\ell$ (Definition~\ref{pkp:def-chart}), so $p_a=\varphi^a(E)=z^{\ell^a}+\ell=(E-\ell)^{\ell^a}+\ell$.
For orders $o$ and a function $c\colon\Z\to\Z$ we write $o+c[a=0]$ for the orders
$(a,u)\mapsto o(a,u)+c(u)[a=0]$. For a lattice with jump sets $L$, a letter $i$ and $l\ge0$ put
\[
\pkgP_l(i)=\mathcal N(L)\cap L_i(\pkgstd_l),\qquad
\pkgC_{l,A}(i)=\{\xi\in\pkgP_l(i):\ g\xi-\xi\in L_i(\pkgstd_l+A[a=0])\ \text{for all }g\in\Gamma'\}
\]
($A\ge0$), the \emph{constants modulo $E^A$}. Both are $\Lambda$-modules, $\pkgP_l(i)$ an
$\mathfrak S$-module: $\mathcal N(L)$ is an $\mathfrak S$-module, $L_i(o)$ an $\pkpA$-module, and for
$w\in\Lambda$, $g(w\xi)-w\xi=w(g\xi-\xi)$ since $\Gamma'$ fixes $\Lambda$. The trivial lattice $C_S$ of a
side is given the jump set $\{0\}$ at every letter; since $\pkgstd_l(a,0)=0=\pkgstd_0(a,0)$,
$(C_S)_i(\pkgstd_l)=(C_S)_i(\pkgstd_0)$ for all $l$.

\subsubsection{Constants modulo $E^A$}\label{pkg:ss-const}

\begin{lemma}\label{pkg:const}
Let $l\ge1$ and $A\ge0$.
\begin{enumerate}
\item[(i)] For a side $S$, $\mathrm{pr}$ maps $\mathcal N(X_S)$, $(X_S)_i(o)$, $\pkgP_l(i)$, $\pkgC_{l,A}(i)$
into the corresponding objects of $C_S$; and $e_i(1\otimes c)\in\pkgC_{l,A}(i)$ for $C_S$, $c\in C_S$.
\item[(ii)] Let $x\in\mathcal N(L)$ with $\varphi^{-l}x\in\pknMinf(L)$, and suppose $g\mu(x)-\mu(x)\in
\mathrm{Fil}^A$ for every form $\mu$ and every $g\in\Gamma'$. Then $e_ix\in\pkgC_{l,A}(i)$.
\end{enumerate}
\end{lemma}

\begin{proof}
(i) By the last paragraph of the proof of Proposition~\ref{pkl:R}, $\mathrm{pr}$ maps $\pkgAN(X_S)$ and
$\mathcal N(X_S)$ into $\pkgAN(C_S)$, $\mathcal N(C_S)$, commutes with $\varphi$, and $\nu\circ\mathrm{pr}$ is
a form on $X_S$ for every form $\nu$ on $C_S$; it commutes with $1\otimes\zeta$, hence with the $e_i$, and
with $\Gamma$. A form $\nu$ of type $0$ at $j$ on $C_S$ gives $\nu\circ\mathrm{pr}$ of type $0$ at $j$ on
$X_S$, and $0\in U(j)$ for the side (Definition~\ref{pkp:def-side}); so $x\in(X_S)_i(o)$ gives
$\nu(\varphi^{-a}\mathrm{pr}\,x)=(\nu\circ\mathrm{pr})(\varphi^{-a}x)\in\mathrm{Fil}^{o(a,0)}$, i.e.\
$\mathrm{pr}\,x\in(C_S)_i(o)$. The assertions on $\pkgP$ and $\pkgC$ follow, as $\mathrm{pr}(gx-x)=
g\,\mathrm{pr}(x)-\mathrm{pr}(x)$. Finally $e_i(1\otimes c)\in\mathcal N(C_S)\cap(C_S)_i(\pkgstd_l)$ by
Proposition~\ref{pkl:H2}(ii), and it is fixed by $\Gamma'$ ($\Gamma$ acts trivially on $C_S$ and fixes
$\Lambda$).

(ii) $e_ix\in\mathcal N(L)$ (Lemma~\ref{pkl:basic}(iv)). As $\varphi^{-l}e_i=e_{i^{(-l)}}\varphi^{-l}$
(Definition~\ref{pkp:def-letters}), $\varphi^{-l}(e_ix)=e_{i^{(-l)}}\varphi^{-l}x\in\pknMinf(L)\cap
\pkpAmb(L)_{i^{(-l)}}$ (Lemma~\ref{pkl:basic}(iv)), so $e_ix\in L_i(\pkgstd_l)$ by
Proposition~\ref{pkl:std}. Let $g\in\Gamma'$; $g(e_ix)-e_ix\in L_i(\pkgstd_l)$ by Lemma~\ref{pkl:G1}. As
$l\ge1$, $\pkgstd_l(0,u)+A=A$, and it remains to see $\mu(g e_ix-e_ix)\in\mathrm{Fil}^A$ for a form $\mu$ of
type $u$ at $i$. The map $\mu'=\mu\circ e_i$ is $\pkpA$-linear, and a form: $e_i$ is a polynomial in
$1\otimes\zeta$ with coefficients in $\Lambda$, and $\zeta\in\auO$ commutes with $\Gamma$. Since $g$
commutes with $e_i$ and $x\in\pkgAN(L)$, Lemma~\ref{pkm:equiv} gives $\mu(ge_ix-e_ix)=\mu'(gx)-\mu'(x)=
g\mu'(x)-\mu'(x)\in\mathrm{Fil}^A$.
\end{proof}

\begin{lemma}[Peeling an expansion]\label{pkg:peel}
Let $l\ge1$, $A,N\ge0$, $\xi_0,\dots,\xi_N\in\pkgC_{l,A}(i)$, $y\in L_i(\pkgstd_l)$, and $o\colon\Z\to\Z$
with $o\le A$. If $\sum_{n\le N}E^n\xi_n+E^{N+1}y\in L_i(\pkgstd_l+o[a=0])$, then
$\xi_N\in L_i(\pkgstd_l+A[o(u)>N][a=0])$.
\end{lemma}

\begin{proof}
\emph{Case $N=0$.} By Lemma~\ref{pkn:E}, $Ey\in L_i(\pkgstd_l+1[a=0])$, so $\xi_0\in L_i(o_1)$ with
$o_1=\pkgstd_l+\min(1,o)[a=0]$ (Lemma~\ref{pkm:le}(i)). Since $\xi_0$ is a constant modulo $E^A$,
Lemma~\ref{pkm:improved} gives $\xi_0\in L_i(\mathrm{imp}(o_1,\pkgstd_l+A[a=0]))$. At $a=0$ (where
$\pkgstd_l=0$ as $l\ge1$) this order is $A$ if $o(u)\ge1$ and $\min(A,0)=0$ otherwise; at $a\ge1$ both
orders equal $\pkgstd_l(a,u)$, and $\mathrm{imp}$ is $\pkgstd_l(a,u)$ in both cases. This is the claim.

\emph{Induction.} Let $N\ge1$ and write the sum as $\xi_0+Ey'$ with $y'=\sum_{n<N}E^n\xi_{n+1}+E^Ny$, which
lies in $L_i(\pkgstd_l)$ (Lemma~\ref{pkn:E}, Lemma~\ref{pkm:le}(i)). By the case $N=0$, $\xi_0\in
L_i(\pkgstd_l+A[o>0][a=0])\subseteq L_i(\pkgstd_l+o[a=0])$, as $A[o>0]\ge o$ ($o\le A$). Hence
$Ey'\in L_i(\pkgstd_l+o[a=0])$ and $y'=E^{-1}(Ey')\in L_i(\pkgstd_l+(o-1)[a=0])$ (Lemma~\ref{pkn:E}). By
induction applied to $\xi_1,\dots,\xi_N$, $y$ and $o-1\le A$, $\xi_N\in L_i(\pkgstd_l+A[o(u)-1>N-1][a=0])$.
\end{proof}

\begin{lemma}[Digits in powers of $E$]\label{pkg:digits}
For $s\in\mathfrak S$ there are $w_n\in\Lambda$ ($n\ge0$) such that for every $N$,
$s-\sum_{n<N}E^nw_n\in E^N\mathfrak S$.
\end{lemma}

\begin{proof}
For $r=\sum_k\lambda_kz^k\in\mathfrak S$ put $r(-\ell)=\sum_k\lambda_k(-\ell)^k\in\Lambda$ ($\ell$-adically
convergent). Then $r-r(-\ell)=\sum_{k\ge1}\lambda_k(z^k-(-\ell)^k)=(z+\ell)q$ with
$q=\sum_{j\ge0}z^j\sum_{k>j}\lambda_k(-\ell)^{k-1-j}\in\mathfrak S$ (inner sums converge). So
$r=r(-\ell)+Eq(r)$. Put $r_0=s$, $w_n=r_n(-\ell)$, $r_{n+1}=q(r_n)$; by induction
$s=\sum_{n<N}E^nw_n+E^Nr_N$.
\end{proof}

\begin{lemma}[The polynomial $f$]\label{pkg:f}
Let $s,g,l\ge0$ and $N=s+g(1+\ell+\dots+\ell^{l-1})$. There is a monic $f\in\Z[X]$ of degree $N$ with
$f(E)=E^s\prod_{a<l}p_a^g$. If $y\in L_i(\rho_{-s})$, then $f(E)y\in L_i(o)$ with $o(a,u)=g[a<l]-u$.
\end{lemma}

\begin{proof}
$f=X^s\prod_{a<l}((X-\ell)^{\ell^a}+\ell)^g$, monic of degree $N$. By Lemma~\ref{pkn:E}, multiplication by
$E^s$ adds $s$ at $a=0$, and by $p_a^g$ adds $g$ at the level $a$; starting from $\rho_{-s}=-u-s[a=0]$
one gets $o$.
\end{proof}

\subsubsection{The radius lemma for frames}\label{pkg:ss-radius}

Put $T=z/\ell=E/\ell-1\in\pkfAm$ (image in $\pkfBp$; $\pkfAm$ is $\ell$-torsion free, so we identify
$\pkfAm$ with $\pkfAo$). The subring $\mathfrak S[T]\subseteq\pkfAm$ lies in $\mathfrak S[1/\ell]\subseteq
\pkfBp$, and $\mathfrak S[1/\ell]=\Lambda[[z]][1/\ell]\subseteq K_0[[z]]$; inside $K_0[[z]]$ also
$\Lambda[[T]]\subseteq K_0[[z]]$ ($T=z/\ell$) and $\mathfrak S[T]\subseteq\Lambda[[T]]$, as
$\mathfrak S=\Lambda[[\ell T]]$. Let $\pkgSh$ be the $\ell$-adic closure of $\mathfrak S[T]$ in $\pkfAm$.

\begin{lemma}\label{pkg:Shat}
(i) $\mathfrak S[T]\cap\ell^n\pkfAm=\ell^n\mathfrak S[T]$ and $\pkgSh\cap\ell^n\pkfAm=\ell^n\pkgSh$; and
$\pkgSh\subseteq\pkfSB$, $\pkfSB=\pkgSh[1/\ell]$.
(ii) The inclusion $\mathfrak S[T]\subseteq\Lambda[[T]]$ extends to an injective ring map
$\psi\colon\pkgSh\to\Lambda[[T]]$.
(iii) If $y\in\pkfSB$ and $z^N$ divides $y$ in $\pkfSB$ for every $N$, then $y=0$.
\end{lemma}

\begin{proof}
(i) \emph{Claim: if $P\in\mathfrak S[T]$ and $P\in\ell^n\Lambda[[T]]$, then $P\in\ell^n\mathfrak S[T]$.}
Let $P=\sum_{k\le D}s_kT^k$, $s_k\in\Lambda[[\ell T]]$; the coefficient of $T^m$ in $P$ is divisible by
$\ell^{m-D}$ for $m>D$. So $Q=\ell^{-n}P\in\Lambda[[T]]$ has coefficients $q_m$ with $\ell^{m-D-n}\mid q_m$ for
$m>D+n$, and $Q=\sum_{m\le D+n}q_mT^m+T^{D+n}\sum_{j\ge1}(\ell^{-j}q_{D+n+j})(\ell T)^j\in\mathfrak S[T]$.
Now let $P\in\mathfrak S[T]\cap\ell^n\pkfAm$ of degree $\le D$ in $T$. Then $\ell^DP\in\mathfrak S\cap
\ell^{n+D}\pkfAm$, which by Proposition~\ref{pkr:I1} is contained in $\mathfrak S\cap(\ell,E)^{n+D}\pkpA$;
applying $\kappa_0$ (Proposition~\ref{pkr:I5}; $\mathfrak S$-linear, identity on $\mathfrak S$, and
$\ell,E\in\mathfrak S$) to an expression $\sum_j\ell^jE^{n+D-j}a_j$ shows $\ell^DP\in(\ell,E)^{n+D}\mathfrak S=
(\ell,z)^{n+D}\mathfrak S$ (the two ideals coincide because $E=z+\ell$), and $(\ell,z)^{n+D}\mathfrak S
\subseteq\ell^{n+D}\Lambda[[T]]$ because $z=\ell T$. So $P\in\ell^n\Lambda[[T]]$ and the claim applies. For
$y\in\pkgSh\cap\ell^n\pkfAm$ choose $P_k\in\mathfrak S[T]$ with $y-P_k\in\ell^{n+k}\pkfAm$; then
$P_k\in\ell^n\mathfrak S[T]$, the $\ell^{-n}P_k$ converge, and their limit $\ell^{-n}y$ lies in $\pkgSh$
($\pkfAm$ is complete, separated and $\ell$-torsion free). By Proposition~\ref{pkf:closure} (with $m=0$;
$\mathfrak S[T]=\mathfrak S[E/\ell]$), $\pkgSh\subseteq\pkfSB$, and every $y\in\pkfSB$ has $\ell^my\in\pkgSh$
for some $m$; with $\ell^{-1}\in\pkfSB$ (Lemma~\ref{pkf:SB-basic}(i)), $\pkfSB=\pkgSh[1/\ell]$.

(ii) If $P_k\to y$ in $\pkgSh$, then $P_k-P_{k'}\in\mathfrak S[T]\cap\ell^m\pkfAm=\ell^m\mathfrak S[T]$ for
large $k,k'$, so $(P_k)$ is $\ell$-adically Cauchy in $\Lambda[[T]]$; put $\psi(y)=\lim P_k$, which is well
defined (apply the same to the alternated sequence) and a ring map. Injectivity: let $\psi(y)=0$. Choose
$P\in\mathfrak S[T]$ with $y-P\in\ell\pkfAm\cap\pkgSh=\ell\pkgSh$; then $\psi(P)\in\ell\Lambda[[T]]$, so
$P\in\ell\mathfrak S[T]$ by the claim, and $y=\ell y_1$ with $y_1\in\pkgSh$, $\psi(y_1)=0$ ($\Lambda[[T]]$ is
torsion free). Hence $y\in\bigcap_n\ell^n\pkfAm=0$.

(iii) Extend $\psi$ to $\pkfSB=\pkgSh[1/\ell]\to\Lambda[[T]][1/\ell]$, still injective. If $y=z^Nc_N$
with $c_N\in\pkfSB$, then $\psi(y)=\ell^NT^N\psi(c_N)\in T^N\Lambda[[T]][1/\ell]$ for all $N$, so $\psi(y)=0$
and $y=0$.
\end{proof}

Fix $l\ge1$ with $\ell^l\ge3$ and put $X_l=z^{\ell^l}/\ell^2=\ell^{\ell^l-2}T^{\ell^l}\in\ell\pkgSh$. As
$\pkgSh$ is closed in the complete ring $\pkfAm$, the $\mathfrak S$-algebra map
$\mathrm{ev}\colon\mathfrak S[[X]]\to\pkgSh$, $\sum_kf_kX^k\mapsto\sum_kf_kX_l^k$, is well defined.

\begin{lemma}[Approximation]\label{pkg:approx}
Let $A,\delta\ge0$ with $A\le\ell^l\delta$ and $F\in\mathfrak S[[X]]$. There is $s\in\mathfrak S$ with
$s-\ell^{2\delta}\mathrm{ev}(F)\in E^A\pkfSB$.
\end{lemma}

\begin{proof}
Write $F=\sum_{k<\delta}f_kX^k+X^\delta F'$. Then $\ell^{2\delta}\mathrm{ev}(F)=s_0+z^{\ell^l\delta}\mathrm{ev}(F')$
with $s_0=\sum_{k<\delta}f_k\ell^{2\delta-2k}z^{\ell^lk}\in\mathfrak S$, and
$z^{\ell^l\delta}\mathrm{ev}(F')=\ell^Ab$ with $b=T^Az^{\ell^l\delta-A}\mathrm{ev}(F')\in\pkgSh\subseteq\pkfSB$.
By Lemma~\ref{pkf:approx0}(i), $b=\sum_k\alpha_k(E/\ell)^k$ with $\alpha_k\in\pkpA$, $\alpha_k\to0$, so
$\ell^Ab=\alpha+E^Aw$ with $\alpha=\sum_{k<A}\alpha_k\ell^{A-k}E^k\in\pkpA$ and $w\in\pkfAm$. Let $G_0$ be the
extension of $\kappa_0$ (Lemma~\ref{pkf:extend}); it is $\pkfSB$-linear, the identity on $\pkfSB$ and has
values in $\pkfSB$ (Lemma~\ref{pkf:values}). So $\ell^Ab=G_0(\ell^Ab)=\kappa_0(\alpha)+E^AG_0(w)$, and
$s=s_0+\kappa_0(\alpha)$ works.
\end{proof}

Now fix a lattice with jump sets $L$ with jump bound $a$, and the frame of Theorem~\ref{pkf:main}: by its
proof and Theorem~\ref{pkf:frame-pm}, $Z\in M_d(\pkfSB)$, $Y\in M_d(\pkfBp)$ with $YZ=E^h$ ($h=2a$), and
$g(Z)=ZG_g$ with $G_g\in M_d(\pkpA)$ for every $g\in\Gamma'$ (such $g$ fix $\Lambda$ and $\theta(z)$). Fix
$g_1\in\Gamma'$ with $g_1z=[\varepsilon]z$ (Definition~\ref{pkp:def-chart}) and write $G=G_{g_1}$.

\begin{lemma}\label{pkg:BW}
(i) There are $B\in M_d(K_0)$, $K_0=\Lambda[1/\ell]$, and $W\in M_d(\pkfSB)$ with $BZ=1+zW$.
(ii) $z\tau$ divides every entry of $G-1$ in $\pkpA$; and $g_1(BZ)=BZG$.
(iii) If $r\in\pkpA$, $n\ge0$ and $\tau z^{n+1}$ divides $r$ in $\pkfBp$, then it divides $r$ in $\pkpA$.
\end{lemma}

\begin{proof}
(i) Let $Y_0=G_0(Y)$ entrywise; $Y_0\in M_d(\pkfSB)$ and, $G_0$ being additive and $\pkfSB$-linear,
$Y_0Z=G_0(YZ)=E^h$. Write $E^h=(z+\ell)^h=\ell^h+zq$ with $q\in\mathfrak S$. By Lemma~\ref{pkl:SBdec} each
entry $b$ of $\ell^{-h}Y_0$ has $\ell^kb=\lambda+zw'$ ($\lambda\in\Lambda$, $w'\in\pkfSB$); let $B$, $W_1$ be the
matrices of the $\ell^{-k}\lambda$, $\ell^{-k}w'$. Then $BZ=\ell^{-h}Y_0Z-zW_1Z=1+z(\ell^{-h}q-W_1Z)$.
(ii) $E^h(G-1)=YZG-YZ=Y(g_1Z-Z)$. By Lemma~\ref{pkl:galSB}, $z\tau$ divides the entries of $g_1Z-Z$ in
$\pkfBp$, hence those of $E^h(G-1)$, and Lemma~\ref{pkl:cancel} gives the first claim. $g_1$ fixes $K_0$, so
$g_1(BZ)=Bg_1(Z)=BZG$.
(iii) By Lemma~\ref{pkl:cancel} ($h=0$), $r=z\tau r'$ with $r'\in\pkpA$; cancelling $z\tau$ in the domain
$\pkfBp$, $z^n$ divides $r'$ in $\pkfBp$, and $n$ applications of Lemma~\ref{pkl:intdiv} (each followed by a
cancellation of $z$) give $z^n\mid r'$ in $\pkpA$.
\end{proof}

Put $Y_1=BZ=1+zW$. A \emph{datum} for a row index $i$ is a triple $(n,R,q)$ with $n\ge0$, $R\in\pkfSB^d$ and
$q\in\Z_{>0}$ such that $q(Y_1)_{ij}-z^{n+1}R_j\in\pkpA$ for all $j$.

\begin{lemma}[One step]\label{pkg:step}
Let $(n,R,q)$ be a datum and $v=v_\ell(n+1)$. There are $s\in\mathfrak S^d$ and $R'\in\pkfSB^d$ with
$\ell^{v+1}R=s+zR'$, and $(n+1,R',\ell^{v+1}q)$ is a datum.
\end{lemma}

\begin{proof}
Let $\alpha_j=q(Y_1)_{ij}-z^{n+1}R_j\in\pkpA$. By Lemma~\ref{pkg:BW}(ii), $\sum_k(Y_1)_{ik}G_{kj}=
g_1((Y_1)_{ij})$, and $g_1(z^{n+1})=[\varepsilon]^{n+1}z^{n+1}$; hence
\[
\textstyle\sum_k\alpha_kG_{kj}-g_1(\alpha_j)=z^{n+1}\Bigl(([\varepsilon]^{n+1}-1)R_j+[\varepsilon]^{n+1}
(g_1R_j-R_j)-\sum_kR_k(G-1)_{kj}\Bigr)=z^{n+1}\bigl(([\varepsilon]^{n+1}-1)R_j-z\tau D_j\bigr)
\]
with $D_j\in\pkfBp$, by Lemma~\ref{pkl:galSB} and Lemma~\ref{pkg:BW}(ii). By Proposition~\ref{pkn:Cc} there are
$\alpha',\beta\in\pkpA$ with $\tau E(\ell^v+\beta z)=\alpha'([\varepsilon]^{n+1}-1)$, and
$E(\ell^v+\beta z)=\ell^{v+1}+z(\ell^v+\beta E)$ as $E=\ell+z$. Multiplying by $\alpha'$,
\[
\textstyle\alpha'\bigl(\sum_k\alpha_kG_{kj}-g_1\alpha_j\bigr)=\tau z^{n+1}\bigl(\ell^{v+1}R_j-zR''_j\bigr),\qquad
R''_j=\alpha'D_j-(\ell^v+\beta E)R_j\in\pkfBp .
\]
The left side lies in $\pkpA$, so by Lemma~\ref{pkg:BW}(iii) it equals $\tau z^{n+1}x_j$ with $x_j\in\pkpA$,
and $x_j=\ell^{v+1}R_j-zR''_j$ ($\pkfBp$ is a domain). Applying $G_0$:
$\kappa_0(x_j)=\ell^{v+1}R_j-zG_0(R''_j)$. Put $s_j=\kappa_0(x_j)$, $R'_j=G_0(R''_j)\in\pkfSB$. Then
$\ell^{v+1}\alpha_j+z^{n+1}s_j=\ell^{v+1}q(Y_1)_{ij}-z^{n+2}R'_j\in\pkpA$.
\end{proof}

\begin{lemma}\label{pkg:series}
There is $F\in M_d(\mathfrak S[[X]])$ with $\ell^2\varphi^l(W)=\mathrm{ev}(F)$ (entrywise).
\end{lemma}

\begin{proof}
Fix $i,j$. $(0,(W_{ik})_k,1)$ is a datum, since $(Y_1)_{ik}-zW_{ik}=\delta_{ik}$. Applying
Lemma~\ref{pkg:step} repeatedly gives data $(n,R^{(n)},q_n)$ and $s^{(n)}\in\mathfrak S^d$ with
$\ell^{v_n+1}R^{(n)}=s^{(n)}+zR^{(n+1)}$, $v_n=v_\ell(n+1)$. Put $\gamma_n=\ell^{n-v_\ell(n!)+2}\varphi^l(R^{(n)}_j)
\in\pkfSB$ (Lemma~\ref{pkf:values}) and $a_n=n+1-v_\ell((n+1)!)\ge0$ ($v_\ell(m!)<m/(\ell-1)\le m$). Apply
$\varphi^l$ ($\varphi^l(z)=z^{\ell^l}=\ell^2X_l$) and multiply by $\ell^{a_n}$; as $a_n+v_n+1=n-v_\ell(n!)+2$,
\[
\gamma_n=\ell^{a_n}\varphi^l(s^{(n)}_j)+X_l\gamma_{n+1},\qquad\text{so}\qquad
\gamma_0=\textstyle\sum_{k<N}c_kX_l^k+X_l^N\gamma_N\quad(N\ge0),
\]
with $c_k=\ell^{a_k}\varphi^l(s^{(k)}_j)\in\mathfrak S$. Let $F_{ij}=\sum_kc_kX^k$. Then
$\gamma_0-\mathrm{ev}(F_{ij})=X_l^N\bigl(\gamma_N-\mathrm{ev}(\sum_{k\ge N}c_kX^{k-N})\bigr)\in X_l^N\pkfSB\subseteq
z^N\pkfSB$ for all $N$ ($X_l=z\cdot\ell^{-2}z^{\ell^l-1}$). By Lemma~\ref{pkg:Shat}(iii),
$\gamma_0=\ell^2\varphi^l(W_{ij})$ equals $\mathrm{ev}(F_{ij})$.
\end{proof}

\begin{proposition}[Radius lemma]\label{pkg:radius}
There are $D\in GL_d(K_0)$ and $V\in GL_d(\pkgSh)$ with $\varphi^l(Z)=p_l^aDV$, such that for all $A,\delta\ge0$
with $A\le\ell^l\delta$: the set $\mathcal Y$ of $y\in\mathfrak S^d$ with $Vy-r\in E^A(\pkfBp)^d$ for some
$r\in\Z^d$ spans an $\mathfrak S$-submodule containing $\ell^{4\delta}\mathfrak S^d$.
\end{proposition}

\begin{proof}
Let $F$ be as in Lemma~\ref{pkg:series}, $\mathcal V'=1+XF\in GL_d(\mathfrak S[[X]])$ (its determinant has
constant term $1$), $\varepsilon_X=1+\ell X\in\mathfrak S[[X]]^\times$, $\mathcal V=\varepsilon_X^{-a}\mathcal V'$ and
$V=\mathrm{ev}(\mathcal V)\in GL_d(\pkgSh)$. As $\varphi^l$ is a ring endomorphism of $\pkfBp$ (P1),
$\varphi^l(B)\varphi^l(Z)=\varphi^l(Y_1)=1+z^{\ell^l}\varphi^l(W)=1+X_l\ell^2\varphi^l(W)=\mathrm{ev}(\mathcal V')$, and
$p_l=z^{\ell^l}+\ell=\ell\,\mathrm{ev}(\varepsilon_X)$. Since $\mathrm{ev}(\mathcal V')$ is invertible,
$\det\varphi^l(B)\ne0$; as $\varphi^l(B)\in M_d(K_0)$, it is invertible over $K_0$, and
$\varphi^l(Z)=\varphi^l(B)^{-1}\mathrm{ev}(\varepsilon_X)^aV=p_l^aDV$ with $D=\ell^{-a}\varphi^l(B)^{-1}$.

For each entry of $\mathcal V$ and of $\mathcal V^{-1}$ choose by Lemma~\ref{pkg:approx} an $s\in\mathfrak S$ as
there, and let $M'$, $M$ be the resulting matrices: modulo the ideal $E^A\pkfBp$ (entrywise),
$M'\equiv\ell^{2\delta}V$ and $M\equiv\ell^{2\delta}V^{-1}$. Hence $VM\equiv\ell^{2\delta}\cdot1$ and
$V(\ell^{4\delta}-MM')\equiv\ell^{4\delta}V-\ell^{2\delta}M'\equiv0$; so the columns of $M$ and of
$\ell^{4\delta}-MM'$ lie in $\mathcal Y$, and for $x\in\mathfrak S^d$,
$\ell^{4\delta}x=(\ell^{4\delta}-MM')x+M(M'x)$ is an $\mathfrak S$-combination of them.
\end{proof}

\begin{theorem}[Spanning by almost fixed vectors]\label{pkg:span}
Let the chart be unramified, $l\ge1$, $\ell^l\ge3$, $A\le\ell^l\delta$, and
$\pkgV_l=\{x\in\mathcal N(L):\varphi^{-l}x\in\pknMinf(L)\}$. For $x\in\pkgV_l$, $\ell^{4\delta}x$ lies in the
$\mathfrak S$-span of the $y\in\pkgV_l$ with $g\mu(y)-\mu(y)\in\mathrm{Fil}^A$ for all forms $\mu$ and
$g\in\Gamma'$.
\end{theorem}

\begin{proof}
\emph{$\pkgV_l$ is the $\mathfrak S$-span of the $\varphi^l(n_j)$.} These lie in $\mathcal N$
(Lemma~\ref{pkl:phin}) and in $\varphi^l\pknMinf$ (Proposition~\ref{pkn:frame}(i)), and $\pkgV_l$ is an
$\mathfrak S$-module ($\varphi^{-l}(\mathfrak S)\subseteq\pkpA$). Conversely, for $x\in\pkgV_l$,
$\varphi^{-l}x=\sum_jt_jn_j$ (Proposition~\ref{pkn:frame}(ii)), so $x=\sum_j\varphi^l(t_j)\varphi^l(n_j)$ and,
applying $\mathrm{ret}$ (Proposition~\ref{pkl:R}), $x=\sum_j\kappa_0(\varphi^lt_j)\varphi^l(n_j)$.

\emph{Conclusion.} Let $x=\sum_jc_j\varphi^l(n_j)$, $c\in\mathfrak S^d$, and write $\ell^{4\delta}c=\sum_ks_ky^{(k)}$
with $y^{(k)}\in\mathcal Y$ (Proposition~\ref{pkg:radius}). Then $\ell^{4\delta}x=\sum_ks_kx_k$ with
$x_k=\sum_jy^{(k)}_j\varphi^l(n_j)\in\pkgV_l$. Let $y\in\mathcal Y$, $Vy=r+E^Aw$, $x_y=\sum_jy_j\varphi^l(n_j)$, and
$\mu$ a form. By Lemma~\ref{pkn:frame-id}(ii), $p_l^a\mu(x_y)=c^{(l)}(\mu)\varphi^l(Z)y=p_l^ac^{(l)}(\mu)DVy$; as
$p_l\ne0$ in $\pknB$, $\mu(x_y)=c^{(l)}(\mu)Dr+E^Ac^{(l)}(\mu)Dw$. The first term lies in $K'$ (entries of
$c^{(l)}(\mu)$ in $K'$, Lemma~\ref{pkn:frame-id}(i); $D$ over $K_0\subseteq K'$) and is fixed by $\Gamma'$; the
second lies in $\mathrm{Fil}^A$, which $\Gamma'$ preserves. So $g\mu(x_y)-\mu(x_y)\in\mathrm{Fil}^A$.
\end{proof}

\subsubsection{Expansion in powers of $E$ and the proof of (I5)}\label{pkg:ss-I5}

\begin{proposition}[Taylor expansion]\label{pkg:taylor}
Let the chart be unramified, $l\ge1$, $\ell^l\ge3$, $A\le\ell^l\delta$, and $x\in\pkgP_l(i)$. There are
$\xi_n\in\pkgC_{l,A}(i)$ ($n\ge0$) such that for every $N$ there is $y_N\in\pkgP_l(i)$ with
$\ell^{4\delta}x=\sum_{n<N}E^n\xi_n+E^Ny_N$.
\end{proposition}

\begin{proof}
By Proposition~\ref{pkl:std}, $\varphi^{-l}x\in\pknMinf$, so $x\in\pkgV_l$. By Theorem~\ref{pkg:span},
$\ell^{4\delta}x=\sum_kr_kv_k$ with $r_k\in\mathfrak S$ and $v_k$ as there; applying $e_i$ ($x=e_ix$),
$\ell^{4\delta}x=\sum_kr_kb_k$ with $b_k=e_iv_k\in\pkgC_{l,A}(i)$ (Lemma~\ref{pkg:const}(ii)). By
Lemma~\ref{pkg:digits}, $r_k=\sum_{n<N}E^nw_{k,n}+E^Nt_{k,N}$; put $\xi_n=\sum_kw_{k,n}b_k\in\pkgC_{l,A}(i)$ and
$y_N=\sum_kt_{k,N}b_k\in\pkgP_l(i)$.
\end{proof}

\begin{theorem}[Input (I5)]\label{pkg:I5}
Let the chart be unramified, $l\ge1$ with $\ell^l\ge3$, and $A,\delta,\kappa,s,g\ge0$ with $A\le\ell^l\delta$.
Let $\ell^\kappa$ kill the cokernel of a side $S$ at $i^{(-l)}$, and $E^s$ kill it at $i$. Let every jump of
$S$ at $i^{(-a)}$, $0\le a<l$, be $\le g$, put $N=s+g(1+\ell+\dots+\ell^{l-1})$, and assume $N<A$. Then for
every $c\in C_S$, $\pkcPt_S(c,i,\kappa+4\delta;o^{\mathrm{gap}}_{l,A,N},d^{\mathrm{gap}}_{l,A})\ne\emptyset$,
where $o^{\mathrm{gap}}_{l,A,N}=\pkgstd_l+A[u<-N][a=0]$ and $d^{\mathrm{gap}}_{l,A}=\pkgstd_l+A[a=0]$.
\end{theorem}

\begin{proof}
We work in $L=X_S$; put $\epsilon=e_i(1\otimes c)$.

\emph{(1) A point and the polynomial.} By Proposition~\ref{pkt:I1} (killer $\ell^0E^s$) there is
$y\in\mathcal N(X_S)\cap(X_S)_i(\rho_{-s})$ with $\mathrm{pr}(y)=\epsilon$. Let $f$ be as in Lemma~\ref{pkg:f}
(degree $N$); $f(E)y\in(X_S)_i(o)$ with $o(a,u)=g[a<l]-u$. On the jumps $u$ at $i^{(-a)}$, $o\ge\pkgstd_l$
($g-u\ge0$ for $a<l$), so $f(E)y\in\pkgP_l(i)$.

\emph{(2) The coefficient $\xi_N$.} By Proposition~\ref{pkg:taylor}, $\ell^{4\delta}f(E)y=\sum_{n\le N}E^n\xi_n+
E^{N+1}y'$ with $\xi_n\in\pkgC_{l,A}(i)$, $y'\in\pkgP_l(i)$. The left side lies in
$(X_S)_i(\pkgstd_l+o_1[a=0])$ with $o_1(u)=\min(A,g-u)\le A$ (at $a=0$, $o(0,u)=g-u\ge o_1(u)$; at $a\ge1$,
$o\ge\pkgstd_l$). By Lemma~\ref{pkg:peel}, $\xi_N\in(X_S)_i(\pkgstd_l+A[o_1(u)>N][a=0])$. If $u<-N$ then
$g-u>N$ and $A>N$, so $o_1(u)>N$; hence $\xi_N\in(X_S)_i(o^{\mathrm{gap}}_{l,A,N})$.

\emph{(3) Its image.} Write $f=\sum_{n\le N}f_nX^n$, $f_N=1$, and $\eta_n=\mathrm{pr}(\xi_n)-\ell^{4\delta}f_n
\epsilon\in\pkgC_{l,A}(i)$ (for $C_S$; Lemma~\ref{pkg:const}(i)). Applying $\mathrm{pr}$ ($\pkpW$-linear),
$\sum_{n\le N}E^n\eta_n+E^{N+1}\mathrm{pr}(y')=\ell^{4\delta}f(E)\epsilon-\ell^{4\delta}f(E)\epsilon=0$, with
$\mathrm{pr}(y')\in(C_S)_i(\pkgstd_l)$. Lemma~\ref{pkg:peel} with $o\equiv A>N$ gives
$w:=\eta_N=\mathrm{pr}(\xi_N)-\ell^{4\delta}\epsilon\in(C_S)_i(d^{\mathrm{gap}}_{l,A})$, and $w\in\mathcal N(C_S)$.

\emph{(4) Lifting the error.} By Lemma~\ref{pkn:E} and Lemma~\ref{pkn:M}(iii), $E^{-A}w\in\mathcal N(C_S)\cap
(C_S)_i(\pkgstd_l)$, and $(C_S)_i(\pkgstd_l)=(C_S)_i(\pkgstd_0)\subseteq\pkpA\otimes C_S$
(Proposition~\ref{pkl:H2}(i)); so $E^{-A}w=e_iE^{-A}w=\sum_kb_ke_i(1\otimes c_k)$ with $b_k\in\pkpA$ and $(c_k)$ an
$\auO$-basis of $C_S$ (refined to a $\Z_\ell$-basis). By Definition~\ref{pkm:def-J}, $\ell^\kappa\in I_0(i^{(-l)})$,
hence $\ell^\kappa=\varphi^l(\ell^\kappa)\in I_l(i)$ (Lemma~\ref{pkm:J-basic}(iii) with $r=0$; $\sigma^l(i^{(-l)})=i$):
there are $x_k\in(X_S)_i(\pkgstd_l)$ with $\mathrm{pr}(x_k)=\ell^\kappa e_i(1\otimes c_k)$. Let $x''=
\mathrm{ret}(\sum_kb_kx_k)$. By Proposition~\ref{pkl:R}, $x''\in\mathcal N(X_S)\cap(X_S)_i(\pkgstd_l)$ and
$\mathrm{pr}(x'')=\mathrm{ret}(\ell^\kappa E^{-A}w)=\ell^\kappa E^{-A}w$, as $\ell^\kappa E^{-A}w\in\mathcal N(C_S)$.

\emph{(5) The point.} Put $\nu=\ell^\kappa\xi_N-E^Ax''\in\mathcal N(X_S)\cap\pkpAmb(X_S)_i$. Then
$\mathrm{pr}(\nu)=\ell^\kappa(w+\ell^{4\delta}\epsilon)-\ell^\kappa w=\ell^{\kappa+4\delta}\epsilon$. By
Lemma~\ref{pkn:E}, $E^Ax''\in(X_S)_i(d^{\mathrm{gap}}_{l,A})\subseteq(X_S)_i(o^{\mathrm{gap}}_{l,A,N})$, so
$\nu$ has orders $o^{\mathrm{gap}}_{l,A,N}$. For $g\in\Gamma'$, $g\xi_N-\xi_N\in(X_S)_i(d^{\mathrm{gap}}_{l,A})$
($\xi_N$ is a constant modulo $E^A$) and $g(E^Ax'')-E^Ax''\in(X_S)_i(d^{\mathrm{gap}}_{l,A})$
(Lemma~\ref{pkl:G1}); so $\nu$ is deep for $d^{\mathrm{gap}}_{l,A}$, and
$\nu\in\pkcPt_S(c,i,\kappa+4\delta;o^{\mathrm{gap}}_{l,A,N},d^{\mathrm{gap}}_{l,A})$.
\end{proof}

\begin{remark}
(I5) is proved as stated in \S\ref{pkd:ss-tools}; the proof of Theorem~\ref{pkd:scheme} applies it without
change. Only $\ell^l\ge3$ is used to make $X_l\in\ell\pkgSh$.
\end{remark}

\subsubsection{The other inputs}\label{pkg:ss-open}

The input (I4) of \S\ref{pkc:ss-inputs} is Theorem~\ref{pkh:I4}; Proposition~\ref{pki:chart} is
Theorem~\ref{pkj:chart}.



\providecommand{\auO}{\mathcal O}
\providecommand{\auQlb}{\overline{\Q}_\ell}
\providecommand{\pkpA}{A_{\mathrm{inf}}}
\providecommand{\pkpW}{W(C^\flat)}
\providecommand{\pkpAmb}{\mathbb A}
\providecommand{\pknBp}{B^+_{\mathrm{dR}}}
\providecommand{\pknB}{B_{\mathrm{dR}}}
\providecommand{\pknK}{\mathcal K_0}
\providecommand{\pktV}{V_{\mathrm{dR}}}
\providecommand{\pkcPt}{\mathrm{Pt}}
\providecommand{\pkcimp}{\mathrm{imp}}
\providecommand{\pkcNM}{\mathrm{NM}}
\providecommand{\pkcnw}{\mathrm{nw}}
\providecommand{\pkcSig}{\Sigma}
\providecommand{\pkiWin}{\operatorname{Win}}
\providecommand{\pglen}{\operatorname{len}}
\providecommand{\pkecost}{\mathrm{cost}}
\providecommand{\pkeDp}{\mathrm{Dp}}
\providecommand{\pkhD}{\mathcal D}
\providecommand{\pkhK}{\mathrm{Const}}
\providecommand{\pkhX}{\mathcal X}
\providecommand{\pkhup}[3]{#1^{+}_{#2,#3}}
\providecommand{\pkhlv}[2]{#1^{+}_{#2}}
\providecommand{\pkhrl}{\operatorname{rl}}

\subsection{The colength of the deep constants: the input (I4)}\label{pkh:sec}

This section proves the input (I4) of \S\ref{pkc:ss-inputs} (``few low jumps: the deep constants''),
exactly as stated there; it is Theorem~\ref{pkh:I4} below. With it, Proposition~\ref{pki:generic}(ii)
is proved from (I1)--(I4), as in \S\ref{pkc:ss-generic}, and (I1)--(I3) are proved in
\S\S\ref{pkt:sec}, \ref{pke:sec}.

\emph{Dependencies.} Proved in full, modulo results cited by label: Lemma~\ref{pkt:hodge},
Lemma~\ref{pkt:expand}, Lemma~\ref{pkt:forms} (Hodge bases and expansion), Lemma~\ref{pkt:sub} and its
proof (forms on the kernel), Proposition~\ref{pkl:R} and Lemma~\ref{pke:ret} (the retraction),
Lemma~\ref{pkl:G1} ($\Gamma'$-stability), Lemmas~\ref{pkn:M} and \ref{pkn:E} (multiplication by $p_b^{\pm1}$),
Lemma~\ref{pkm:improved} (improved orders), Proposition~\ref{pke:A} and Lemma~\ref{pke:op-basic}
(operators of a cost), Proposition~\ref{pkm:kill-unram}, Proposition~\ref{pkt:I1}, Lemmas~\ref{pkc:win} and
\ref{pkc:free}, Lemma~\ref{pkm:le}, Remark~\ref{pkc:mono}, and Definitions~\ref{pkp:def-side} and
\ref{pkp:def-points} (sides and points), Definition~\ref{pkn:def-N} (the modules $\mathcal N$), and the
definition of $\pkcimp$ in \S\ref{pkm:ss-state}. Besides these and published results cited by name (Tate's
theorem on the Galois invariants of $C(n)$), no other input is used; all numerical inequalities are checked in the text.

\emph{Conventions.} Those of \S\S\ref{pkc:sec}, \ref{pkt:sec}, \ref{pke:sec}. Throughout, the Kummer
chart is \emph{unramified}: $e=1$, $E=z+\ell$, $d_v=0$, so $\pkecost(a,n)=v_\ell(n)+3a$ with $v_\ell(0)=0$.
Recall $z=[\varpi^\flat]$, so $\varphi(z)=z^\ell$ and $p_a=\varphi^a(E)=z^{\ell^a}+\ell$. Orders are functions
$o\colon\Z_{\ge0}\times\Z\to\Z$; for orders $o$, a level $a\ge0$ and $u\in\Z$ put
\[
\pkhup oau(a',v)=o(a',v)+[(a',v)=(a,u)],\qquad \pkhlv oa(a',v)=o(a',v)+[a'=a]
\]
(raise the single value $o(a,u)$, resp.\ all values at the level $a$, by one). As in \S\ref{pkc:sec},
inequalities between orders are required only on the jumps, i.e.\ at the $(a',v)$ with $v\in
U(i^{(-a')})$ for the letter $i$ in question. ``Length'' means length of $\Z_\ell$-modules, with values in
$\Z_{\ge0}\cup\{\infty\}$; it is additive in short exact sequences. For $\Z_\ell$-submodules $A,B$ of a
module put
\[
\pkhrl(A,B)=\pglen(A/(A\cap B)).
\]

\subsubsection{Two algebraic lemmas}\label{pkh:ss-alg}

\begin{lemma}[Relative lengths]\label{pkh:len}
Let $A,B,B'$ and $A_1,\dots,A_k$ be $\Z_\ell$-submodules of a $\Z_\ell$-module $M$.
\begin{enumerate}
\item[(i)] $\pkhrl(A,B')\le\pkhrl(A,B)+\pkhrl(B,B')$; if $A\subseteq B$ then $\pkhrl(A,B)=0$.
\item[(ii)] $\pglen\bigl(M/\bigcap_jA_j\bigr)\le\sum_j\pglen(M/A_j)$.
\item[(iii)] If $f\colon M\to M'$ is $\Z_\ell$-linear and $A'\subseteq M'$ is a submodule with $f(B)\subseteq
A'$, then $\pglen\bigl(f(A)/(f(A)\cap A')\bigr)\le\pkhrl(A,B)$.
\item[(iv)] Let $t,m\ge0$. Suppose $\ell^tA\subseteq B$, and that for any $x_0,\dots,x_m\in A$ with $\ell
x_j\in B$ for all $j$ there are $\beta_0,\dots,\beta_m\in\Z_\ell$, one of them equal to $1$, with
$\sum_j\beta_jx_j\in B$. Then $\pkhrl(A,B)\le tm$.
\end{enumerate}
\end{lemma}

\begin{proof}
(i) In $\bar A=A/(A\cap B')$ let $\bar B=((A\cap B)+(A\cap B'))/(A\cap B')\cong(A\cap B)/(A\cap B\cap B')$. This
is isomorphic to a submodule of $B/(B\cap B')$, so $\pglen\bar B\le\pkhrl(B,B')$; and $\bar A/\bar B=
A/((A\cap B)+(A\cap B'))$ is a quotient of $A/(A\cap B)$, so $\pglen(\bar A/\bar B)\le\pkhrl(A,B)$. Additivity
gives the claim. The second assertion is clear.

(ii) The map $M/\bigcap_jA_j\to\bigoplus_jM/A_j$ is injective.

(iii) $f$ maps $A$ onto $f(A)$ and $A\cap B$ into $f(A)\cap A'$ (as $f(B)\subseteq A'$), so $f(A)/(f(A)\cap
A')$ is a quotient of $A/(A\cap B)$.

(iv) Put $Q=A/(A\cap B)$; it is killed by $\ell^t$. Let $Q[\ell]$ be its $\ell$-torsion, an $\mathbb
F_\ell$-vector space. \emph{Claim: $\dim Q[\ell]\le m$.} Let $q_0,\dots,q_m\in Q[\ell]$ and lift them to
$x_j\in A$; then $\ell x_j\in A\cap B$, so the hypothesis gives $\beta_j\in\Z_\ell$, $\beta_{j_0}=1$, with
$\sum_j\beta_jx_j\in B$, i.e.\ $\sum_j\bar\beta_jq_j=0$ in $Q$ with $\bar\beta_{j_0}=1\ne0$ in $\mathbb F_\ell$.
So any $m+1$ elements of $Q[\ell]$ are linearly dependent. \emph{Induction on $t$}: we show that a
$\Z_\ell$-module $Q$ killed by $\ell^t$ with $\dim Q[\ell]\le m$ has length $\le tm$. For $t=0$, $Q=0$. For
$t\ge1$ the sequence $0\to Q[\ell]\to Q\xrightarrow{\ \ell\ }\ell Q\to0$ is exact, $\ell Q$ is killed by
$\ell^{t-1}$ and $(\ell Q)[\ell]\subseteq Q[\ell]$; so $\pglen Q=\dim Q[\ell]+\pglen\ell Q\le m+(t-1)m$.
\end{proof}

\begin{lemma}[A primitive relation over $\Lambda$]\label{pkh:wrel}
Let $d\le n$ and $w_0,\dots,w_{fn}\in\Lambda^d$, where $\Lambda=W(\mathbb F_{\ell^f})$. There are
$\beta_0,\dots,\beta_{fn}\in\Z_\ell$, one of them equal to $1$, with $\sum_k\beta_kw_k=0$.
\end{lemma}

\begin{proof}
$\Lambda$ is free of rank $f$ over $\Z_\ell$, so $\Lambda^d$ is free of rank $fd\le fn<fn+1$, and the $fn+1$
vectors $w_k$ are linearly dependent over $\Q_\ell$; clearing denominators, $\sum_k\alpha_kw_k=0$ with
$\alpha_k\in\Z_\ell$ not all $0$. Choose $k_0$ with $v_\ell(\alpha_{k_0})$ minimal (so $\alpha_{k_0}\ne0$) and put
$\beta_k=\alpha_k/\alpha_{k_0}\in\Z_\ell$. Then $\alpha_{k_0}\sum_k\beta_kw_k=0$, and $\Lambda^d$ is torsion free, so
$\sum_k\beta_kw_k=0$; $\beta_{k_0}=1$.
\end{proof}

\begin{lemma}[Division by $E$]\label{pkh:Ediv}
(i) Every $s\in\mathfrak S$ is $s=w+Ey$ with $w\in\Lambda$, $y\in\mathfrak S$. (ii) For $a\ge0$,
$E^{\ell^a}-p_a\in\ell\mathfrak S$.
\end{lemma}

\begin{proof}
(i) Write $s=\sum_{k\ge0}\lambda_kz^k$, $\lambda_k\in\Lambda$. Put $w=\sum_k\lambda_k(-\ell)^k$, convergent in
$\Lambda$, and $y=\sum_{j\ge0}y_jz^j$ with $y_j=\sum_{k>j}\lambda_k(-\ell)^{k-1-j}$, convergent in $\Lambda$. Since
$z^k-(-\ell)^k=(z+\ell)\sum_{j<k}z^j(-\ell)^{k-1-j}$, comparing coefficients gives $(z+\ell)y=s-w$: the constant
coefficient of $(z+\ell)y$ is $\ell y_0=-\sum_{k\ge1}\lambda_k(-\ell)^k=\lambda_0-w$, and for $j\ge1$ the coefficient
of $z^j$ is $y_{j-1}+\ell y_j=\lambda_j+\sum_{k>j}\lambda_k(-\ell)^{k-j}-\sum_{k>j}\lambda_k(-\ell)^{k-j}=\lambda_j$.
(ii) $E^{\ell^a}=(z+\ell)^{\ell^a}
\equiv z^{\ell^a}\equiv p_a\pmod{\ell\mathfrak S}$, since $\ell\mid\binom{\ell^a}{k}$ for $0<k<\ell^a$ and
$p_a=z^{\ell^a}+\ell$.
\end{proof}

\subsubsection{Leading terms and primitive relations}\label{pkh:ss-lead}

For a lattice $L$ with jump sets, a letter $j$ and $u\in\Z$ let $m^L_j(u)$ be the sum, over the
embeddings $\gamma$ of letter $j$, of the multiplicity of the jump $u$ of $L$ at $\gamma$. By
Lemma~\ref{pkt:hodge} (last paragraph of its proof), in a Hodge basis $(d_{j,m},\kappa_{j,m})_m$ of $L$ at $j$
exactly $m^L_j(u)$ indices have $\kappa_{j,m}=u$. For a side $S$ with lattice $Y$ we write $m_j(u)=m^Y_j(u)$;
this is the $m_i(u)$ of (I4).

\begin{lemma}[Extension and kernel]\label{pkh:kernel}
Let $S$ be a side with lattice $Y$ and extension $X_S$, and $\iota\colon\pkpAmb(Y)\to\pkpAmb(X_S)$, $y\mapsto
(y,0)$, the $\pkpW$-linear extension of the inclusion. (i) For every letter $i$ and orders $o$, $\iota$ is a
bijection from $\mathcal N(Y)\cap Y_i(o)$ onto the set of $x\in\mathcal N(X_S)\cap(X_S)_i(o)$ with
$\mathrm{pr}(x)=0$. (ii) For $u\ne0$ and every letter $j$, $m^{X_S}_j(u)=m_j(u)$.
\end{lemma}

\begin{proof}
(i) $\iota$ is injective, and its image is the kernel of $\mathrm{pr}$; $\iota$ is $\Gamma$-equivariant (the
inclusion $Y\to X_S$ is) and commutes with $\varphi$, with $1\otimes\zeta$ and with the $e_j$.
\emph{Into.} Let $y\in\mathcal N(Y)\cap Y_i(o)$. For a form $\mu$ on $X_S$, $\mu\circ\iota$ is
$\pkpA$-linear and $\mu(\iota(1\otimes gy'))=\mu(1\otimes g\iota y')=g\mu(1\otimes\iota y')$, so it is a form
on $Y$; if $\mu$ has type $u$ at $j$, so has $\mu\circ\iota$, as $\iota(e_j(1\otimes Y))\subseteq
e_j(1\otimes X_S)$. If $ry\in\mathfrak M_{\mathrm{inf}}(Y)$ with $r$ a product of $p_a$'s and $\tau^Nry\in
\pkpA\otimes Y$, then $\tau^Nr\iota y=\iota(\tau^Nry)\in\pkpA\otimes X_S$ and $\mu(\varphi^k(r\iota
y))=(\mu\circ\iota)(\varphi^k(ry))\in\pknBp$ for $k\ge1$; so $\iota y\in\pkpA\mathcal N(X_S)$
(Definition~\ref{pkn:def-N}). Likewise $\mu(\iota y)=(\mu\circ\iota)(y)\in\pknK$, so $\iota y\in\mathcal N(X_S)$;
$\iota y$ lies in the letter-$i$ part, and $\mu(\varphi^{-a}\iota y)=(\mu\circ\iota)(\varphi^{-a}y)\in
\mathrm{Fil}^{o(a,u)}$ for $\mu$ of type $u\in U(i^{(-a)})$ at $i^{(-a)}$ (the jump sets of $X_S$ and $Y$ are
the same, Definition~\ref{pkp:def-side}). So $\iota y\in\mathcal N(X_S)\cap(X_S)_i(o)$, over $0$.
\emph{Onto.} Let $x\in\mathcal N(X_S)\cap(X_S)_i(o)$ with $\mathrm{pr}(x)=0$, and $y=\mathrm{fst}(x)$, so
$x=\iota y$. By Lemma~\ref{pkt:sub}, $y\in Y_i(o)\subseteq\pkpA\mathcal N(Y)$. Let $\mu$ be a form on $Y$. The
proof of Lemma~\ref{pkt:sub} (\emph{extension of forms}, with no type condition) gives for every letter $j$ a
form $\nu_j$ on $X_S$ with $\hat\nu_j\circ\hat\iota=\hat\mu\circ e_j$ on $\pktV(Y)$. Hence
$\mu(y)=\hat\mu(\hat y)=\sum_j\hat\mu(e_j\hat y)=\sum_j\hat\nu_j(\widehat{\iota y})=\sum_j\nu_j(x)\in\pknK$,
since $x\in\mathcal N(X_S)$ and $\pknK$ is a field. So $y\in\mathcal N(Y)$.

(ii) Let $\gamma$ be an embedding of letter $j$. Applying $C\otimes_{\auO,\gamma}-$ to the $\auO$-split exact
sequence $0\to Y\to X_S\to C_S\to0$ of $\auO[\Gamma']$-modules gives an exact sequence $0\to W'\to W\to W''\to0$
of $C$-semilinear representations of $G_{K'}$. They are Hodge--Tate ($Y[1/\ell]$, $X_S[1/\ell]$ are
crystalline on $\Gamma'$, and $C_S$ is trivial), so each is a direct sum of copies of $C(-\kappa)$, and by
Tate's theorem ($C(n)^{G_{K'}}$ is $K'$ for $n=0$ and $0$ otherwise) the multiplicity of the jump $\kappa$ is
$h_\kappa=\dim_{K'}(\,\cdot\,(\kappa))^{G_{K'}}$, with $\sum_\kappa h_\kappa$ equal to the $C$-dimension. Taking
invariants is left exact, so $h_\kappa(W)\le h_\kappa(W')+h_\kappa(W'')$ for every $\kappa$, and summing over
$\kappa$ gives equality of the totals; hence $h_\kappa(W)=h_\kappa(W')+h_\kappa(W'')$ for every $\kappa$. As
$W''$ is a sum of copies of $C$, $h_u(W'')=0$ for $u\ne0$. Sum over the $\gamma$ of letter $j$.
\end{proof}

\begin{lemma}[Leading terms]\label{pkh:lead}
Let $L$ be a lattice with jump sets $U$, $l$ a letter, $b\ge0$, $j=l^{(-b)}$, $u\in\Z$, and $o$ orders with
$o(b,u)<o(b,v)$ for every $v\in U(j)$ with $v<u$. Let $x_1,\dots,x_N\in L_l(o)$. There are $d\le m^L_j(u)$,
indices $s_1,\dots,s_d\in\{1,\dots,N\}$ and $A_{k,q}\in\pkpA$ ($1\le k\le N$, $1\le q\le d$) with
\[
x_k-\textstyle\sum_{q=1}^dA_{k,q}x_{s_q}\in L_l(\pkhup obu)\qquad(1\le k\le N).
\]
\end{lemma}

\begin{proof}
If $u\notin U(j)$, then $L_l(\pkhup obu)=L_l(o)$ (the value at $(b,u)$ is not tested) and $d=0$ works. Let
$u\in U(j)$. Fix a Hodge basis $(d_{j,m},\kappa_{j,m})_m$ of $L$ at $j$ (Lemma~\ref{pkt:hodge}) with coordinate
forms $\delta_{j,m}$, and let $m_1,\dots,m_r$ be the indices with $\kappa_{j,m}=u$; $r=m^L_j(u)$. For a vector
$x$ in the letter-$l$ part of $\pkpA\mathcal N(L)$ write $\hat x_b=\widehat{\varphi^{-b}x}\in V_j$
(\S\ref{pkt:ss-hodge}). Each $\delta_{j,m}$ is $\hat\mu$ for a form $\mu$ of type $\kappa_{j,m}\in U(j)$ at $j$,
so for $x\in L_l(o)$, $\delta_{j,m}(\hat x_b)\in\mathrm{Fil}^{o(b,\kappa_{j,m})}$.

\emph{Leading coefficients.} Write $\delta_{j,m_p}((\hat x_k)_b)=t^{o(b,u)}\tau_{k,p}$ with $\tau_{k,p}\in\pknBp$,
and let $T=(\theta(\tau_{k,p}))\in C^{N\times r}$, of rank $d\le r$. If $d=0$ put $\alpha=0$. Otherwise choose
$d$ rows $s_1,\dots,s_d$ and $d$ columns $J$ such that the minor $\Delta=\det T_{s,J}$ has maximal absolute
value among all $d\times d$ minors of $T$; $\Delta\ne0$. The rows $s_q$ are linearly independent, hence span
the row space, and every row is $T_k=\sum_q\alpha_{k,q}T_{s_q}$ with $\alpha_{k,q}\in C$; restricting to the
columns $J$ and using Cramer's rule, $\alpha_{k,q}=\Delta_{k,q}/\Delta$, where $\Delta_{k,q}$ is the determinant
of $T_{s,J}$ with the row $s_q$ replaced by $T_k|_J$. Up to sign $\Delta_{k,q}$ is a $d\times d$ minor of $T$
(or $0$, if $k$ is another $s_{q'}$), so $|\alpha_{k,q}|\le1$ and $\alpha_{k,q}\in O_C$. As $\theta\colon\pkpA\to
O_C$ is onto, choose $A'_{k,q}\in\pkpA$ with $\theta(A'_{k,q})=\alpha_{k,q}$, and put $A_{k,q}=\varphi^b(A'_{k,q})$.
Then $\theta\bigl(\tau_{k,p}-\sum_qA'_{k,q}\tau_{s_q,p}\bigr)=0$, i.e.\
$\tau_{k,p}-\sum_qA'_{k,q}\tau_{s_q,p}\in\mathrm{Fil}^1$, for all $k,p$.

\emph{The difference.} Fix $k$ and put $x'=x_k-\sum_qA_{k,q}x_{s_q}\in L_l(o)$ ($L_l(o)$ is an $\pkpA$-module).
Since $\varphi^{-b}(A_{k,q}x)=A'_{k,q}\varphi^{-b}x$ and $\delta_{j,m}$ is $\pknB$-linear,
\[
\delta_{j,m_p}(\hat x'_b)=t^{o(b,u)}\bigl(\tau_{k,p}-\textstyle\sum_qA'_{k,q}\tau_{s_q,p}\bigr)\in
\mathrm{Fil}^{o(b,u)+1}.
\]
Now let $\mu$ be a form of type $u$ at $j$. By Lemmas~\ref{pkt:forms}(i) and \ref{pkt:expand},
$\mu(\varphi^{-b}x')=\hat\mu(\hat x'_b)=\sum_m\hat\mu(d_{j,m})\delta_{j,m}(\hat x'_b)$ with $\hat\mu(d_{j,m})\in
K'\subseteq\pknBp$, and $\hat\mu(d_{j,m})=0$ unless $\kappa_{j,m}\le u$. The terms with $\kappa_{j,m}=u$ lie in
$\mathrm{Fil}^{o(b,u)+1}$ by the display; a term with $\kappa_{j,m}=v<u$ has $v\in U(j)$ and lies in
$\mathrm{Fil}^{o(b,v)}\subseteq\mathrm{Fil}^{o(b,u)+1}$, as $o(b,v)\ge o(b,u)+1$. So $\mu(\varphi^{-b}x')\in
\mathrm{Fil}^{\pkhup obu(b,u)}$; at all other $(a,v)$ the orders $\pkhup obu$ and $o$ agree. Hence $x'\in
L_l(\pkhup obu)$.
\end{proof}

\begin{lemma}[Primitive relation]\label{pkh:prim}
Let the chart be unramified. Let $L$, $l$, $b$, $j$, $u$, $o$ be as in Lemma~\ref{pkh:lead}, $n\ge m^L_j(u)$,
and $x_0,\dots,x_{fn}\in\mathcal N(L)\cap L_l(o)$. There are $\beta_0,\dots,\beta_{fn}\in\Z_\ell$, one of them
equal to $1$, with
\[
\textstyle\sum_k\beta_kx_k=w+z,\qquad w\in\mathcal N(L)\cap L_l(\pkhlv o0),\quad z\in\mathcal N(L)\cap
L_l(\pkhup obu).
\]
\end{lemma}

\begin{proof}
Lemma~\ref{pkh:lead} gives $d\le m^L_j(u)\le n$, indices $s_q$ and $A_{k,q}$ with $y_k=x_k-\sum_qA_{k,q}x_{s_q}\in
L_l(\pkhup obu)$. Let $\mathrm{ret}$ be the retraction of Proposition~\ref{pkl:R}, with coordinate $\kappa_0\colon
\pkpA\to\mathfrak S$. As the $x$'s lie in $\mathcal N(L)$, $\mathrm{ret}(y_k)=x_k-\sum_q\kappa_0(A_{k,q})x_{s_q}$, and
by Lemma~\ref{pke:ret} it lies in $\mathcal N(L)\cap L_l(\pkhup obu)$. By Lemma~\ref{pkh:Ediv}(i) write
$\kappa_0(A_{k,q})=w_{k,q}+Ey_{k,q}$ with $w_{k,q}\in\Lambda$, $y_{k,q}\in\mathfrak S$. Lemma~\ref{pkh:wrel},
applied to the vectors $(w_{k,q})_{q\le d}\in\Lambda^d$, gives $\beta_k\in\Z_\ell$, one equal to $1$, with
$\sum_k\beta_kw_{k,q}=0$ for every $q$. Then
\[
\textstyle\sum_k\beta_kx_k=\underbrace{\sum_k\beta_k\,\mathrm{ret}(y_k)}_{z}+\underbrace{E\sum_k\beta_k\sum_q
y_{k,q}x_{s_q}}_{w}+\sum_q\Bigl(\sum_k\beta_kw_{k,q}\Bigr)x_{s_q},
\]
and the last sum vanishes. $z\in\mathcal N(L)\cap L_l(\pkhup obu)$, a $\Z_\ell$-module. The element
$w'=\sum_k\beta_k\sum_qy_{k,q}x_{s_q}$ lies in the $\mathfrak S$-module $\mathcal N(L)$ (Lemma~\ref{pkn:M}(i)) and
in $L_l(o)$, so $w=Ew'=p_0w'\in\mathcal N(L)\cap L_l(\pkhlv o0)$ by Lemma~\ref{pkn:E} ($b=0$, $n=1$).
\end{proof}

\subsubsection{Deep constants and the cost of one step}\label{pkh:ss-step}

Fix a side $S$ with lattice $Y$, extension $X_S$, jump sets $U$; write $\mathcal N=\mathcal N(X_S)$,
$L_i(o)=(X_S)_i(o)$, $L^0_i(o)$ for its elements over $0$, and $\pkeDp_i(o)$ for the set of $x\in\mathcal N$
lying over a constant at $i$ and deep for $o$ at $i$ (\S\ref{pke:ss-ops}). The \emph{constants at $i$} form
the $\Z_\ell$-module $\pkhK_i=\{\varepsilon_i(c):c\in C_S\}$, $\varepsilon_i(c)=e_i(1\otimes c)$.

\begin{definition}[Deep constants]\label{pkh:def-D}
For a letter $i$, $t\in\Z$ and orders $o$ put $\pkhX_i(t;o)=\mathcal N\cap L_i(\rho_t)\cap\pkeDp_i(o)$ and
$\pkhD_i(t;o)=\mathrm{pr}\bigl(\pkhX_i(t;o)\bigr)\subseteq\pkhK_i$.
\end{definition}

\begin{remark}\label{pkh:rem-D}
(a) $\pkeDp_i(o)$ is a $\Z_\ell$-module: sums and $\Z_\ell$-multiples of vectors over constants at $i$ lie over
constants at $i$, and $L_i(o)$ is an $\pkpA$-module. Hence $\pkhX_i(t;o)$ and $\pkhD_i(t;o)$ are
$\Z_\ell$-modules. (b) $\mathcal N\cap L^0_i(o)\subseteq\pkeDp_i(o)$: such a vector lies over the constant
$\varepsilon_i(0)=0$, and $gx-x\in L_i(o)$ for $g\in\Gamma'$ by Lemma~\ref{pkl:G1}. (c) For $c\in C_S$,
$\varepsilon_i(c)\in\pkhD_i(t;o)$ if and only if $\pkcPt_S(c,i,0;\rho_t,o)\ne\emptyset$ (Definition of
$\pkcPt$ in \S\ref{pkc:sec}: a point over $c$ at $i$ lies in $\pkeDp_i(o)$, and conversely a vector of
$\pkhX_i(t;o)$ with $\mathrm{pr}=\varepsilon_i(c)$ lies over $c$ at $i$, being in the letter-$i$ part). (d) If
$o'\le o$ on the jumps then $\pkeDp_i(o)\subseteq\pkeDp_i(o')$ and $\pkhD_i(t;o)\subseteq\pkhD_i(t;o')$
(Lemma~\ref{pkm:le}(i)).
\end{remark}

For orders $c$ and $a\ge0$ we say that \emph{$c$ is $a$-free at $i$} if no multiple of $\ell$ lies strictly
between $c(0,v)$ and $c(0,v)+\ell^a$, for every $v\in U(i)$. Every $c$ is $0$-free.

\begin{lemma}[Division by $\ell$]\label{pkh:div}
Let $a\ge0$, $u\in\Z$, and $c$ orders that are $a$-free at $i$. Let $x\in\pkeDp_i(c)$ and $w\in\mathcal N\cap
L^0_i(\pkhlv c0)$ with $\ell x-w\in\pkeDp_i(\pkhup cau)$. Then there is $m\in\mathcal N\cap L^0_i(c)$ with
$x-m\in\pkeDp_i(\pkhup cau)$.
\end{lemma}

\begin{proof}
\emph{An operator of cost $0$.} Let $c''=c+\ell^a[a'=0]$ and $I=\{(0,y):c(0,v)<y<c(0,v)+\ell^a\text{ for some }
v\in U(i)\}$, a finite set. $I$ covers $\pkhlv c0\to c''$ at $i$ (\S\ref{pke:ss-ops}): at the level $0$ the
pairs $(0,y)$ with $c(0,v)+1\le y<c(0,v)+\ell^a$ are in $I$, and at the levels $a'\ge1$ there are none, as
$\pkhlv c0=c''=c$ there. For $(0,y)\in I$, $\ell\nmid y$ (so $y\ne0$), hence $\pkecost(0,y)=v_\ell(y)=0$.
Proposition~\ref{pke:A} gives an operator $A$ of cost $0$ with $A(\pkeDp_i(\pkhlv c0))\subseteq\pkeDp_i(c'')$.

\emph{The element $w'$.} By Remark~\ref{pkh:rem-D}(b), $w\in\pkeDp_i(\pkhlv c0)$, so $Aw\in\pkeDp_i(c'')$, and
$Aw-w\in L^0_i(\pkhlv c0)$; so $Aw\in\mathcal N\cap L^0_i(\pkhlv c0)$. By Lemma~\ref{pkm:improved}, $Aw\in
L_i(\pkcimp(\pkhlv c0,c''))$. We claim $\pkcimp(\pkhlv c0,c'')\ge c''$ on the jumps. At a level $a'\ge1$ both
arguments equal $c$, and $\pkcimp(c,c)=c$ (it is $c$ if $c>0$ and $\min(c,0)=c$ otherwise). At the level $0$,
for $v\in U(i)$: if $c(0,v)+1>0$ the value is $c''(0,v)$; otherwise $c(0,v)<0$, and since $0$ is a multiple of
$\ell$ not strictly between $c(0,v)$ and $c(0,v)+\ell^a$, we get $c(0,v)+\ell^a\le0$, so the value
$\min(c''(0,v),0)$ equals $c''(0,v)$. Hence $Aw\in L_i(c'')$. Put $w'=E^{-\ell^a}Aw$. By Lemma~\ref{pkn:M}(iii),
$w'\in\mathcal N$; by Lemma~\ref{pkn:E} ($b=0$, $n=-\ell^a$), $w'\in L_i(c)$; and $\mathrm{pr}(w')=0$. So
$w'\in\mathcal N\cap L^0_i(c)$.

\emph{The element $m$.} By Lemma~\ref{pkh:Ediv}(ii) write $E^{\ell^a}-p_a=\ell r$ with $r\in\mathfrak S$, and put
$m=rw'-(Ax-x)$. As $x\in\pkeDp_i(c)$ and $A$ has cost $0$, $Ax-x\in L^0_i(c)$, and $Ax\in\mathcal N$; with
$rw'\in\mathcal N\cap L^0_i(c)$ this gives $m\in\mathcal N\cap L^0_i(c)$. Then $x-m\in\pkeDp_i(c)$
(Remark~\ref{pkh:rem-D}(a),(b)), and it lies over the constant $\mathrm{pr}(x)$.

\emph{Depth.} We have
\[
\ell(x-m)=A(\ell x-w)+(Aw-\ell rw'),\qquad Aw-\ell rw'=(E^{\ell^a}-\ell r)w'=p_aw'.
\]
(Indeed $A(\ell x-w)+Aw-\ell rw'=\ell Ax-\ell rw'=\ell x-\ell m$.) By Lemma~\ref{pke:op-basic}(i), $A(\ell
x-w)\in\pkeDp_i(\pkhup cau)$. By Lemma~\ref{pkn:E}, $p_aw'\in L_i(c+[a'=a])\subseteq L_i(\pkhup cau)$, and this
module is $\Gamma'$-stable (Lemma~\ref{pkl:G1}). So for $g\in\Gamma'$, $\ell\bigl(g(x-m)-(x-m)\bigr)\in L_i(\pkhup
cau)$. Finally $y=g(x-m)-(x-m)\in L_i(c)$ lies in the letter-$i$ part of $\pkpA\mathcal N(X_S)$, and for a form
$\mu$, $\mu(\varphi^{-a'}(\ell y))=\ell\mu(\varphi^{-a'}y)$ with $\ell$ a unit of $\pknBp$; so $y\in L_i(\pkhup
cau)$. Thus $x-m\in\pkeDp_i(\pkhup cau)$.
\end{proof}

\begin{lemma}[A relation that gains one step]\label{pkh:raw}
Let the chart be unramified, $a\ge0$, $u\in\Z$, $n\ge m_{i^{(-a)}}(u)$, and let $c$ be $a$-free at $i$ with
$c(a,u)<c(a,v)$ for all $v\in U(i^{(-a)})$, $v<u$. Let $x_0,\dots,x_{fn}\in\pkeDp_i(c)$ with $\ell x_k=x'_k+m_k$,
$x'_k\in\pkeDp_i(\pkhup cau)$, $m_k\in\mathcal N\cap L^0_i(c)$. There are $\beta_k\in\Z_\ell$, one equal to $1$,
and $m\in\mathcal N\cap L^0_i(c)$ with $\sum_k\beta_kx_k-m\in\pkeDp_i(\pkhup cau)$.
\end{lemma}

\begin{proof}
By Lemma~\ref{pkh:kernel}(i), $m_k=\iota(y_k)$ with $y_k\in\mathcal N(Y)\cap Y_i(c)$. Lemma~\ref{pkh:prim} for
the lattice $Y$ (jump sets $U$), $l=i$, $b=a$, $j=i^{(-a)}$, the orders $c$ and $n\ge m^Y_j(u)$ gives
$\beta_k$, one equal to $1$, with $\sum_k\beta_ky_k=w_Y+z_Y$, $w_Y\in\mathcal N(Y)\cap Y_i(\pkhlv c0)$, $z_Y\in
\mathcal N(Y)\cap Y_i(\pkhup cau)$. Put $w=\iota w_Y\in\mathcal N\cap L^0_i(\pkhlv c0)$ and $z=\iota z_Y\in
\mathcal N\cap L^0_i(\pkhup cau)\subseteq\pkeDp_i(\pkhup cau)$ (Lemma~\ref{pkh:kernel}(i),
Remark~\ref{pkh:rem-D}(b)). Then $x=\sum_k\beta_kx_k\in\pkeDp_i(c)$ and
\[
\textstyle\ell x-w=\sum_k\beta_kx'_k+\iota\bigl(\sum_k\beta_ky_k\bigr)-w=\sum_k\beta_kx'_k+z\in
\pkeDp_i(\pkhup cau).
\]
Lemma~\ref{pkh:div} gives $m$.
\end{proof}

\begin{lemma}[One step]\label{pkh:one}
Let the chart be unramified, $a\ge0$, $u\in\Z$, and let $c$ be $a$-free at $i$, with $c\ge\rho_0$ on the jumps
($c(l,v)\ge-v$ for $v\in U(i^{(-l)})$) and $c(a,u)<c(a,v)$ for all $v\in U(i^{(-a)})$, $v<u$. Then
\[
\pkhrl\bigl(\pkhD_i(0;c),\pkhD_i(0;\pkhup cau)\bigr)\le\pkecost(a,c(a,u))\cdot f\,m_{i^{(-a)}}(u).
\]
\end{lemma}

\begin{proof}
Put $e=\pkecost(a,c(a,u))$, $n=m_{i^{(-a)}}(u)$, $X(o)=\pkhX_i(0;o)$ and $M=\mathcal N\cap L^0_i(c)$. Since
$c\ge\rho_0$ on the jumps, $M\subseteq L_i(\rho_0)$ (Lemma~\ref{pkm:le}(i)), so $M\subseteq X(c)$
(Remark~\ref{pkh:rem-D}(b)); also $X(\pkhup cau)\subseteq X(c)$. Let $B=X(\pkhup cau)+M\subseteq X(c)$. As
$\mathrm{pr}(M)=0$, $\mathrm{pr}(B)\subseteq\pkhD_i(0;\pkhup cau)$, and Lemma~\ref{pkh:len}(iii) (with $f=\mathrm{pr}$,
$A=X(c)$) gives $\pkhrl(\pkhD_i(0;c),\pkhD_i(0;\pkhup cau))\le\pkhrl(X(c),B)$. We bound the latter by
Lemma~\ref{pkh:len}(iv) with $t=e$, $m=fn$.

\emph{$\ell^eX(c)\subseteq B$.} The set $\{(a,c(a,u))\}$ covers $c\to\pkhup cau$ at $i$; Proposition~\ref{pke:A}
gives an operator $A$ of cost $e$ with $A(\pkeDp_i(c))\subseteq\pkeDp_i(\pkhup cau)$. For $x\in X(c)$, $Ax-\ell^e
x\in L^0_i(c)\cap\mathcal N=M$ ($A$ has values in $\mathcal N$), so $Ax\in L_i(\rho_0)$ and $Ax\in X(\pkhup
cau)$; hence $\ell^ex=Ax-(Ax-\ell^ex)\in B$.

\emph{Relations.} Let $x_0,\dots,x_{fn}\in X(c)$ with $\ell x_k\in B$: $\ell x_k=x'_k+m_k$ with $x'_k\in
X(\pkhup cau)$, $m_k\in M$. Lemma~\ref{pkh:raw} gives $\beta_k$, one equal to $1$, and $m\in M$ with
$y=\sum_k\beta_kx_k-m\in\pkeDp_i(\pkhup cau)$. Also $y\in\mathcal N\cap L_i(\rho_0)$, so $y\in X(\pkhup cau)$
and $\sum_k\beta_kx_k=y+m\in B$.
\end{proof}

\begin{lemma}[All steps at one level]\label{pkh:level}
Let the chart be unramified, $a\ge0$, and let $c,c'$ be orders such that $c$ is $a$-free at $i$, $c\ge\rho_0$ on
the jumps, $c'(l,v)\le c(l,v)$ for $l\ne a$ and $v\in U(i^{(-l)})$, and $c'(a,\cdot)$ is nonincreasing on
$U(i^{(-a)})$. Then
\[
\pkhrl\bigl(\pkhD_i(0;c),\pkhD_i(0;c')\bigr)\le\sum_{u\in U(i^{(-a)})}\ \sum_{c(a,u)\le y<c'(a,u)}
\pkecost(a,y)\,f\,m_{i^{(-a)}}(u).
\]
\end{lemma}

\begin{proof}
Write $U_a=U(i^{(-a)})$ and induct on $N(c)=\sum_{u\in U_a}\max(0,c'(a,u)-c(a,u))$, the number of steps
$(u,y)$, $u\in U_a$, $c(a,u)\le y<c'(a,u)$. If $N(c)=0$, then $c'\le c$ on all jumps, so $\pkhD_i(0;c)\subseteq
\pkhD_i(0;c')$ (Remark~\ref{pkh:rem-D}(d)) and the left side is $0$. Otherwise let $u$ be the least element of
$U_a$ with $c(a,u)<c'(a,u)$. For $v\in U_a$, $v<u$: $c'(a,v)\le c(a,v)$ by minimality, and
$c(a,u)<c'(a,u)\le c'(a,v)$ as $c'(a,\cdot)$ is nonincreasing; so $c(a,u)<c(a,v)$. Put $\tilde c=\pkhup cau$. It
is $a$-free at $i$: if $a\ne0$ it agrees with $c$ at the level $0$, and all orders are $0$-free. Also
$\tilde c\ge c\ge\rho_0$, $c'(l,\cdot)\le\tilde c(l,\cdot)$ for $l\ne a$, and the steps for $(\tilde c,c')$ are
those for $(c,c')$ except $(u,c(a,u))$; so $N(\tilde c)=N(c)-1$. By Lemmas~\ref{pkh:len}(i), \ref{pkh:one} and
induction,
\[
\pkhrl(\pkhD_i(0;c),\pkhD_i(0;c'))\le\pkecost(a,c(a,u))\,f\,m_{i^{(-a)}}(u)+\sum_{(v,y)\ne(u,c(a,u))}
\pkecost(a,y)\,f\,m_{i^{(-a)}}(v),
\]
which is the sum over all steps.
\end{proof}

\begin{corollary}[Level bound with budgets]\label{pkh:levelcor}
In the situation of Lemma~\ref{pkh:level}, let $\eta\ge0$, $C_{\mathrm{all}},C_{\mathrm{low}}\ge0$, and let
$n_j,r_j$ be integers with $n_j\ge f\sum_{v\in[-\eta,\eta]}m_j(v)$ and $r_j\ge f\sum_{v\in U(j)}m_j(v)$ for
every letter $j$. If $\sum_{c(a,u)\le y<c'(a,u)}\pkecost(a,y)\le C_{\mathrm{all}}+C_{\mathrm{low}}[|u|\le\eta]$
for every $u\in U(i^{(-a)})$, then $\pkhrl(\pkhD_i(0;c),\pkhD_i(0;c'))\le C_{\mathrm{all}}r_{i^{(-a)}}+
C_{\mathrm{low}}n_{i^{(-a)}}$.
\end{corollary}

\begin{proof}
With $j=i^{(-a)}$ the bound of Lemma~\ref{pkh:level} is at most $\sum_{u\in U(j)}(C_{\mathrm{all}}+C_{\mathrm{low}}
[|u|\le\eta])fm_j(u)=C_{\mathrm{all}}f\sum_{u\in U(j)}m_j(u)+C_{\mathrm{low}}f\sum_{u\in U(j)\cap[-\eta,\eta]}
m_j(u)\le C_{\mathrm{all}}r_j+C_{\mathrm{low}}n_j$.
\end{proof}

\subsubsection{Few low jumps}\label{pkh:ss-few}

In this subsection the side $S$ satisfies the hypotheses of (I4) (there the side is called $S'$): the chart is
unramified, $\ell\ge5$, $a,\eta,W\ge0$, $\eta\ge1$, and for every letter $j$ an integer $b_j\ge0$ and a sign
$\varepsilon_j$ are given with
\[
\begin{gathered}b_j+W+\eta\le(\ell-3)a,\qquad n_0:=6(a+\eta+W)<b_j,\qquad\pkcNM(\ell,n_0,b_j),\\ U(j)\subseteq
\pkiWin(\eta,\varepsilon_j,b_j,W)=[-\eta,\eta]\cup\varepsilon_j[b_j,b_j+W];\end{gathered}
\]
and $n_j,r_j$ are integers with $n_j\ge f\sum_{v\in[-\eta,\eta]}m_j(v)$, $r_j\ge f\sum_{v\in U(j)}m_j(v)$. Recall
that $\pkcNM(\ell,n_0,b)$ means that no nonzero element of $\pkcSig_{n_0}(b)=\{y:|y|\le n_0\text{ or }|y\mp b|\le
n_0\}$ is divisible by $\ell$. A jump $u\in U(j)$ is \emph{low} if $|u|\le\eta$ and \emph{far} otherwise; a far
jump has $b_j\le|u|\le b_j+W$, so $|u|>n_0\ge a+\eta$. Every jump is $\le b_j+W+\eta\le(\ell-3)a$.

\begin{lemma}[Deep points over multiples of constants]\label{pkh:points}
Let $0\le s\le a$, $i$ a letter and $c\in C_S$. Then
$\pkcPt_S\bigl(c,i,s;\pkcimp(\rho_{-(a-s)},\rho_{2a+\eta}),\rho_{2a+\eta}\bigr)\ne\emptyset$, where
\[
\pkcimp(\rho_{-(a-s)},\rho_{2a+\eta})(0,u)=\begin{cases}2a+\eta-u&u<-(a-s),\\ \min(2a+\eta-u,0)&u\ge-(a-s),
\end{cases}\qquad \pkcimp(\rho_{-(a-s)},\rho_{2a+\eta})(a',u)=-u\ (a'\ge1).
\]
\end{lemma}

\begin{proof}
Every jump of $S$ is $\le(\ell-3)a=(\ell-3)(s+(a-s))$, so by Proposition~\ref{pkm:kill-unram} ($\ell\ge5$),
$\ell^sE^{a-s}$ kills the cokernel of $S$ at every letter, and Proposition~\ref{pkt:I1} gives
$x\in\pkcPt_S(c,i,s;\rho_{-(a-s)},\rho_{-(a-s)})$. We deepen it to $\rho_{2a+\eta}$ by Lemma~\ref{pkc:free}
(unramified case) with $n=n_0$, $b=b_i$. The pairs $(a',y)$ with $\rho_{-(a-s)}(a',u)\le y<\rho_{2a+\eta}(a',u)$
for some $u\in U(i^{(-a')})$ have $a'=0$ (at the levels $a'\ge1$ both orders are $-u$) and
$-(a-s)-u\le y<2a+\eta-u$ with $u\in U(i)$; by Lemma~\ref{pkc:win}(iii) (with $s'=a-s$, $t'=2a+\eta$, and
$n_0=6(a+\eta+W)\ge\max(a-s,2a+\eta)+\eta+W=2a+2\eta+W$) they satisfy $y\in\pkcSig_{n_0}(b_i)$. Since
$\pkcNM(\ell,n_0,b_i)$ holds, Lemma~\ref{pkc:free} gives a point in
$\pkcPt_S(c,i,s;\pkcimp(\rho_{-(a-s)},\rho_{2a+\eta}),\rho_{2a+\eta})$. The formula is the definition of
$\pkcimp$ (\S\ref{pkm:ss-state}): at the level $0$, $\rho_{-(a-s)}(0,u)=-(a-s)-u>0$ iff $u<-(a-s)$; at the levels
$a'\ge1$ both arguments are $-u$ and $\pkcimp=-u$ in both cases.
\end{proof}

\begin{lemma}[Removing the pole]\label{pkh:pole}
For every letter $i$, $\pkhrl\bigl(\pkhK_i,\pkhD_i(0;\rho_{2a+\eta})\bigr)\le a\,n_i$.
\end{lemma}

\begin{proof}
Write $s_0=2a+\eta$ and $\pkhD(t)=\pkhD_i(t;\rho_{s_0})$.

\emph{Start: $\pkhK_i\subseteq\pkhD(-\eta)$.} Let $c\in C_S$ and $x$ the point of Lemma~\ref{pkh:points} with
$s=0$, orders $o=\pkcimp(\rho_{-a},\rho_{s_0})$. We check $o\ge\rho_{-\eta}$ on the jumps. At levels $a'\ge1$
both are $-u$. At the level $0$, for $u\in U(i)$: if $u<-a$, $o(0,u)=2a+\eta-u\ge-\eta-u$. If $u\ge-a$, then
$u$ is low or $u>0$ (a far jump with $u\ge-a$ has $|u|>a$, hence $u>a\ge0$); in both cases $-\eta-u\le0$, and
$-\eta-u\le2a+\eta-u$, so $o(0,u)=\min(2a+\eta-u,0)\ge-\eta-u$. By Remark~\ref{pkc:mono}, $x\in\pkcPt_S(c,i,0;
\rho_{-\eta},\rho_{s_0})$, so $\varepsilon_i(c)\in\pkhD(-\eta)$ (Remark~\ref{pkh:rem-D}(c)).

\emph{Kill: $\ell^a\pkhK_i\subseteq\pkhD(t)$ for $t\le0$.} Let $x$ be the point of Lemma~\ref{pkh:points} with
$s=a$, orders $o=\pkcimp(\rho_0,\rho_{s_0})$; $\mathrm{pr}(x)=\ell^a\varepsilon_i(c)=\varepsilon_i(\ell^ac)$. At
the level $0$: if $u<0$, $o(0,u)=2a+\eta-u\ge t-u$; if $u\ge0$, $o(0,u)=\min(2a+\eta-u,0)\ge t-u$ as $t\le0\le
u$. At levels $\ge1$ both are $-u$. So $x\in\pkcPt_S(\ell^ac,i,0;\rho_t,\rho_{s_0})$ and
$\ell^a\varepsilon_i(c)\in\pkhD(t)$.

\emph{Step: for $-\eta\le t\le-1$, $\pkhrl(\pkhD(t),\pkhD(t+1))\le a\,f\,m_i(t)$.} Apply
Lemma~\ref{pkh:len}(iv) with $A=\pkhD(t)$, $B=\pkhD(t+1)$, exponent $a$ and $m=fn$, $n=m_i(t)$; $\ell^aA\subseteq
B$ by the kill step ($t+1\le0$). Let $k_0,\dots,k_{fn}\in\pkhD(t)$ and choose $x_k\in\pkhX_i(t;\rho_{s_0})$ with
$\mathrm{pr}(x_k)=k_k$. Since $t\ne0$, $m^{X_S}_i(t)=m_i(t)=n$ (Lemma~\ref{pkh:kernel}(ii)). The orders
$\rho_t$ satisfy $\rho_t(0,t)=0<t-v=\rho_t(0,v)$ for $v<t$. Lemma~\ref{pkh:prim} for the lattice $X_S$,
$l=j=i$, $b=0$, $u=t$, $o=\rho_t$ gives $\beta_k\in\Z_\ell$, one equal to $1$, with $x=\sum_k\beta_kx_k=w+z$,
$w\in L_i(\pkhlv{(\rho_t)}0)=L_i(\rho_{t+1})$, $z\in L_i(\pkhup{(\rho_t)}0t)$. As $\pkhup{(\rho_t)}0t\le\rho_{t+1}$
(at $(0,v)$: $t-v+[v=t]\le t+1-v$; elsewhere equal), $x\in L_i(\pkhup{(\rho_t)}0t)$; and $x\in\pkeDp_i(\rho_{s_0})$.
By Lemma~\ref{pkm:improved}, $x\in L_i(\pkcimp(\pkhup{(\rho_t)}0t,\rho_{s_0}))$. We check that this is $\ge
\rho_{t+1}$ on the jumps. At levels $\ge1$ both arguments are $-v$ and $\pkcimp=-v$. At the level $0$, let
$\lambda=t-v+[v=t]$: if $\lambda>0$ the value is $s_0-v\ge t+1-v$ (as $s_0\ge0>t$); if $\lambda\le0$ then
$v\ne t$ (else $\lambda=1$) and $t-v\le0$, so $v\ge t+1$, and the value $\min(s_0-v,0)$ is $\ge t+1-v$. Hence
$x\in\pkhX_i(t+1;\rho_{s_0})$ and $\sum_k\beta_kk_k=\mathrm{pr}(x)\in B$.

\emph{Conclusion.} By Lemma~\ref{pkh:len}(i) applied repeatedly,
\[
\pkhrl(\pkhK_i,\pkhD(0))\le\pkhrl(\pkhK_i,\pkhD(-\eta))+\sum_{t=-\eta}^{-1}\pkhrl(\pkhD(t),\pkhD(t+1))\le
0+a\,f\sum_{t=-\eta}^{-1}m_i(t)\le a\,n_i.\qedhere
\]
\end{proof}

\begin{lemma}[Costs at the level $0$]\label{pkh:cost0}
Let $i$ be a letter with $\varepsilon_i=-1$, and $q\ell$ the least multiple of $\ell$ with $q\ell\ge
b_i+W+a+\eta$. For $u\in U(i)$,
\[
\sum_{2a+\eta-u\le y<q\ell}\pkecost(0,y)\le(a+1)[|u|\le\eta].
\]
\end{lemma}

\begin{proof}
Here $\pkecost(0,y)=v_\ell(y)\ge0$, and $q\ell\le b_i+W+a+\eta+\ell-1$. \emph{$u$ low.} The steps lie in
$[0,q\ell]$ as $2a+\eta-u\ge2a\ge0$. With $N=q\ell$, $\sum_{0\le y\le N}v_\ell(y)=v_\ell(N!)=\sum_{k\ge1}\lfloor
N/\ell^k\rfloor\le N/(\ell-1)$ (Legendre). And $N\le b_i+W+\eta+a+\ell-1\le(\ell-3)a+a+\ell-1\le(\ell-1)a+\ell-1=
(\ell-1)(a+1)$; so the sum is $\le a+1$. \emph{$u$ far.} As $\varepsilon_i=-1$, $-u\in[b_i,b_i+W]$. Let $y$ be a
step with $\ell\mid y$. Then $y\le q\ell-\ell<b_i+W+a+\eta$ (minimality of $q\ell$) and $y\ge2a+\eta-u\ge
b_i+2a+\eta$; so $0\le y-b_i\le W+a+\eta\le n_0$ and $y\ge b_i>0$, contradicting $\pkcNM(\ell,n_0,b_i)$. So every
step has $v_\ell(y)=0$ and the sum is $0$.
\end{proof}

\begin{lemma}[Costs at the level $1$]\label{pkh:cost1}
Let $i$ be a letter, $m=i^{(-1)}$, and $\pkcnw_i$ the depth of (I4). For $u\in U(m)$,
\[
\sum_{-u\le y<\pkcnw_i(1,u)}\pkecost(1,y)\le3(W+1)+6a[|u|\le\eta].
\]
\end{lemma}

\begin{proof}
$\pkcnw_i(1,u)=W+1-u+2a[|u|\le\eta]$, so there are $W+1+2a[|u|\le\eta]$ steps $y\in[-u,-u+W+2a[|u|\le\eta]]$,
and $\pkecost(1,y)=v_\ell(y)+3$. Each step lies in $\pkcSig_{n_0}(b_m)$: if $u$ is low, $|y|\le\eta+W+2a\le n_0$;
if $u=b_m+k$ ($0\le k\le W$), then $y+b_m\in[-k,W-k]$, so $|y+b_m|\le W$; if $u=-(b_m+k)$, then $y-b_m\in[k,k+W]$,
so $|y-b_m|\le2W\le n_0$. By $\pkcNM(\ell,n_0,b_m)$, $y=0$ or $\ell\nmid y$; in both cases $v_\ell(y)=0$. So
the sum is $3(W+1+2a[|u|\le\eta])$.
\end{proof}

\begin{lemma}[The negative window]\label{pkh:neg}
Let $i$ be a letter with $\varepsilon_i=-1$ and $m=i^{(-1)}$. Then
\[
\pkhrl\bigl(\pkhD_i(0;\rho_{2a+\eta}),\pkhD_i(0;\pkcnw_i)\bigr)\le(a+1)n_i+3(W+1)r_m+6a\,n_m .
\]
\end{lemma}

\begin{proof}
Let $q\ell$ be as in Lemma~\ref{pkh:cost0} and $d(0,u)=q\ell$, $d(a',u)=-u$ for $a'\ge1$. By
Lemma~\ref{pkh:len}(i) it suffices to bound $\pkhrl(\pkhD_i(0;\rho_{2a+\eta}),\pkhD_i(0;d))$ by $(a+1)n_i$ and
$\pkhrl(\pkhD_i(0;d),\pkhD_i(0;\pkcnw_i))$ by $3(W+1)r_m+6an_m$.

\emph{Level $0$.} Corollary~\ref{pkh:levelcor} with level $0$ (that is, with its level parameter, called $a$ there, set to $0$; the $a$ of (I4) plays no role in this choice), $c=\rho_{2a+\eta}$, $c'=d$: $c$ is $0$-free;
$c\ge\rho_0$; for $l\ge1$, $c'(l,u)=-u=c(l,u)$; $c'(0,\cdot)$ is constant; and by Lemma~\ref{pkh:cost0} the
budgets $C_{\mathrm{all}}=0$, $C_{\mathrm{low}}=a+1$ work. The bound is $(a+1)n_i$.

\emph{Level $1$.} Corollary~\ref{pkh:levelcor} with level $1$ (its level parameter set to $1$), $c=d$, $c'=\pkcnw_i$. $d$ is $1$-free at $i$: no
multiple of $\ell$ lies strictly between $q\ell$ and $q\ell+\ell$. $d\ge\rho_0$ on the jumps: for $u\in U(i)$,
$-u\le\max(\eta,b_i+W)\le b_i+W+a+\eta\le q\ell$, and at levels $\ge1$ equality. For $l\ne1$: $\pkcnw_i(0,u)=
b_i+W+a+\eta\le q\ell=d(0,u)$, and $\pkcnw_i(l,u)=-u=d(l,u)$ for $l\ge2$. \emph{$\pkcnw_i(1,\cdot)$ is
nonincreasing on $U(m)$}: let $u<v$ in $U(m)$. If both are low or both far, $\pkcnw_i(1,u)-\pkcnw_i(1,v)=v-u>0$.
If $u$ is low and $v$ far, the difference is $v-u+2a>0$. If $u$ is far and $v$ low, then $u<v\le\eta<b_m\le|u|$
forces $u\le-b_m$, and the difference is $v-u-2a\ge-\eta+b_m-2a>n_0-\eta-2a\ge0$. By Lemma~\ref{pkh:cost1}
(the steps start at $d(1,u)=-u$) the budgets $C_{\mathrm{all}}=3(W+1)$, $C_{\mathrm{low}}=6a$ work. The bound
is $3(W+1)r_m+6an_m$.
\end{proof}

\begin{lemma}[Colength at one letter]\label{pkh:const}
For every letter $i$, with $\delta_i$ as in (I4),
\[
\pkhrl\bigl(\pkhK_i,\pkhD_i(0;\delta_i)\bigr)\le(2a+1)n_i+3(W+1)r_{i^{(-1)}}+6a\,n_{i^{(-1)}}.
\]
\end{lemma}

\begin{proof}
If $\varepsilon_i=1$, $\delta_i=\rho_{2a+\eta}$ and Lemma~\ref{pkh:pole} gives $\le an_i$. If $\varepsilon_i=-1$,
$\delta_i=\pkcnw_i$, and Lemmas~\ref{pkh:len}(i), \ref{pkh:pole}, \ref{pkh:neg} give $\le an_i+(a+1)n_i+
3(W+1)r_{i^{(-1)}}+6an_{i^{(-1)}}$. All $n_j,r_j\ge0$.
\end{proof}

\subsubsection{The input (I4)}\label{pkh:ss-I4}

\begin{theorem}[Input (I4)]\label{pkh:I4}
The statement (I4) of \S\ref{pkc:ss-inputs} holds as stated: under its hypotheses, the set $V$ of $c'\in
C_{S'}$ with $\pkcPt_{S'}(c',i,0;\rho_0,\delta_i)\ne\emptyset$ for every letter $i$ is a $\Z_\ell$-submodule of
$C_{S'}$ with $\pglen_{\Z_\ell}(C_{S'}/V)\le11\sum_i\bigl(n_i(a+1)+r_i(W+1)\bigr)$.
\end{theorem}

\begin{proof}
Apply \S\ref{pkh:ss-few} to $S=S'$. By Remark~\ref{pkh:rem-D}(c), $V=\bigcap_i\varepsilon_i^{-1}(\pkhD_i(0;\delta_i))$,
an intersection of $\Z_\ell$-submodules (Remark~\ref{pkh:rem-D}(a); $\varepsilon_i$ is $\Z_\ell$-linear). Since
$\varepsilon_i\colon C_S\to\pkhK_i$ is onto, $C_S/\varepsilon_i^{-1}(\pkhD_i(0;\delta_i))\cong\pkhK_i/\pkhD_i(0;\delta_i)$.
By Lemmas~\ref{pkh:len}(ii) and \ref{pkh:const},
\[
\pglen(C_S/V)\le\sum_i\bigl((2a+1)n_i+3(W+1)r_{i^{(-1)}}+6an_{i^{(-1)}}\bigr)=\sum_i\bigl((8a+1)n_i+3(W+1)r_i\bigr),
\]
because $i\mapsto i^{(-1)}$ is a bijection of the set of letters. Finally $8a+1\le11(a+1)$ and $3(W+1)\le
11(W+1)$, and $n_i,r_i\ge0$.
\end{proof}

\begin{remark}[Use in Proposition~\ref{pki:generic}(ii)]\label{pkh:rem-use}
The statement is proved as stated, so the proof of Proposition~\ref{pki:generic}(ii) in \S\ref{pkc:ss-generic}
goes through unchanged. The proof gives the slightly better bound $\sum_i((8a+1)n_i+3(W+1)r_i)$; the
hypothesis $\eta\ge1$ is not used.
\end{remark}



\providecommand{\auO}{\mathcal O}
\providecommand{\auQlb}{\overline{\Q}_\ell}
\providecommand{\pkilog}{\operatorname{lg}}
\providecommand{\pkjG}{G_{\Q_\ell}}
\providecommand{\pkjcl}{\operatorname{cl}}
\providecommand{\pkjPl}{\mathsf P}
\providecommand{\pkjPb}{\mathsf P_b}

\section{The chart for the data}\label{pkj:sec}

This section proves Proposition~\ref{pki:chart}: for the data, at every prime $\ell$ and every place
$v\mid\ell$, the letters read thickness in the sense of Definition~\ref{pki:def-reads}. The proof
uses only the arithmetic of the field $E$ of the data ({\cite[\S13.1]{P46}}), the two properties of
the set $\pkjPb$ of $b$-places required in (D5) of {\cite[Definition~13.3]{P46}} (class $-1$ modulo
$q_0$, and sparseness for the primes up to $5r(r+1-m)=kt$, {\cite[Definition~13.2]{P46}}), the
sizes of {\cite[Lemma~13.4]{P46}}, and elementary facts on roots of unity in local fields. No
Kummer chart is needed: Definition~\ref{pki:def-reads} involves only $e$, $f$ (those of $E_v$), the
root $\zeta$ and the letters.

We prove Proposition~\ref{pki:chart} (Theorem~\ref{pkj:chart}) for either of
the two natural readings of ``the complex embedding attached to $\gamma$'' (Convention~\ref{pkj:conv});
reading (R$^-$) is the one used ({\cite[\S5.1]{P55}}, \S\ref{pkd:ss-data}).

\subsection{Setting and conventions}\label{pkj:ss-setting}

Let $(q_0,d)$ be the parameters and $D$ the admissible data ({\cite[Definition~13.3]{P46}}), with
$E=K(\zeta_{q_1})$, $K=\Q(\sqrt{-d})$, $q_1=q_1(s)$, $f_E=[E:K]=\varphi(q_1)$ (written $f$ in
{\cite[\S13]{P46}}; here $f$ is reserved for the residue degree of $E_v$), $\pkjPl$ the set of
infinite places of $E$, $\pkjPb\subseteq\pkjPl$ the $b$-places, $Q=kt=5r(r+1-m)$ and $h=r^8$. We
write $\zeta_{q_1}\in E$ for the fixed primitive $q_1$-th root of unity, and put
\[
N_E=q_1d .
\]

\emph{The CM type and the signs.} As in {\cite[\S13.1]{P46}}, an embedding $K\subset\C$ is fixed;
let $\delta_0\in\C$ be the image of $\sqrt{-d}$. The CM type $\Phi$ of the data is the set of
embeddings $\sigma\colon E\to\C$ extending $K\subset\C$; each place $\tau\in\pkjPl$ is induced by
exactly one $\Phi(\tau)\in\Phi$, the other embedding inducing $\tau$ being $\overline{\Phi(\tau)}$.
For $\sigma\colon E\to\C$ put $\varepsilon(\sigma)=1$ if $\sigma\in\Phi$ and $-1$ otherwise. Then
\begin{equation}\label{pkj:eq-sign}
\sigma(\sqrt{-d})=\varepsilon(\sigma)\,\delta_0,\qquad
\Phi(\tau_\sigma)(\zeta_{q_1})=\sigma(\zeta_{q_1})^{\varepsilon(\sigma)},
\end{equation}
where $\tau_\sigma$ is the place of $\sigma$: if $\sigma\notin\Phi$ then $\sigma|_K$ is the conjugate
embedding and $\Phi(\tau_\sigma)=\bar\sigma$, and complex conjugation inverts roots of unity.

\emph{Ratios of places.} For $\tau,\tau'\in\pkjPl$, $\Phi(\tau)$ and $\Phi(\tau')$ both extend
$K\subset\C$, so $\Phi(\tau)=\Phi(\tau')\circ\phi$ for a unique $\phi\in\operatorname{Gal}(E/K)$; if
$\phi(\zeta_{q_1})=\zeta_{q_1}^{j}$ we write $\tau/\tau'=j\in(\Z/q_1)^\times$. Thus
\begin{equation}\label{pkj:eq-ratio}
\Phi(\tau)(\zeta_{q_1})=\Phi(\tau')(\zeta_{q_1})^{\tau/\tau'},
\end{equation}
and $(\tau/\tau')(\tau'/\tau'')=\tau/\tau''$. Fixing a base place $\tau_0$, $\tau\mapsto\tau/\tau_0$ is
the identification $\pkjPl\cong(\Z/q_1)^\times$ of {\cite[\S13.1]{P46}}, and for $m'\mid q_1$ the
reduction $\pkjcl_{m'}$ satisfies $\pkjcl_{m'}(\tau/\tau')=\pkjcl_{m'}(\tau)\pkjcl_{m'}(\tau')^{-1}$,
independently of $\tau_0$. This direction of the ratio, \eqref{pkj:eq-ratio}, is the one used
throughout. (If the identification were read with the inverse ratio, every statement
below holds with the exponent $\varpi$ of Lemma~\ref{pkj:ratio} replaced by $-\varpi$; the proofs of
Lemmas~\ref{pkj:signs} and~\ref{pkj:arc} cover both values of $\varpi$.)

\emph{The local data.} Fix the isomorphism $\iota\colon\C\to\auQlb$ and let $v\mid\ell$. Let
$e_v\colon E\to\auQlb$ be an embedding inducing $v$ (for instance $e_v=\iota\circ\sigma$ for a complex
embedding $\sigma$ with this property), $E_v=\Q_\ell\,e_v(E)\subset\auQlb$ (a finite extension of
$\Q_\ell$, hence complete), and identify $\Gamma=D_v$ with
\[
\Gamma=\operatorname{Gal}(\auQlb/E_v)=\{g\in\pkjG:\ g(e_v(x))=e_v(x)\ \text{for all }x\in E\},
\qquad \pkjG=\operatorname{Gal}(\auQlb/\Q_\ell).
\]
Let $e$ and $f$ be the ramification index and residue degree of $E_v/\Q_\ell$, as in
\S\ref{pki:ss-arch}. Since $E/\Q$ is abelian, $e_v(E)$ is stable under $\pkjG$. Every $g\in\pkjG$ is an
isometry of $\auQlb$ (uniqueness of the extension of $|\cdot|_\ell$), so it preserves the valuation
ring and induces an automorphism $\bar g$ of the residue field $\overline{\mathbb F}_\ell$.

\begin{definition}[Convention]\label{pkj:conv}
For $g\in\pkjG$ let
\[
\sigma_g=\iota^{-1}\circ g^{-1}\circ e_v\quad\text{(reading (R$^-$))}\qquad\text{or}\qquad
\sigma_g=\iota^{-1}\circ g\circ e_v\quad\text{(reading (R$^+$))}.
\]
The paper uses reading (R$^-$); we prove everything for both readings, putting $\varpi=-1$ for (R$^-$)
and $\varpi=+1$ for (R$^+$), one of the two being fixed throughout a proof. Write
$P(g)\in\pkjPl$ for the place of $\sigma_g$ and $\varepsilon(g)=\varepsilon(\sigma_g)$; call $g$
\emph{thick} if $P(g)\in\pkjPb$. A coefficient ring $\auO\subset\auQlb$ is \emph{large} if it contains
the ring of integers $\auO_{E_v}$ of $E_v$. For such $\auO$, every $\Z_\ell$-algebra embedding $\gamma\colon\auO\to
\C_\ell$ has image in $\auQlb$ and is the restriction of some $g\in\pkjG$ (extend $\gamma$ $\Q_\ell$-linearly to the algebraic extension $\auO[1/\ell]$ of $\Q_\ell$, whose image is algebraic over $\Q_\ell$, and then to an automorphism of $\auQlb$); the embeddings of Proposition~\ref{pki:chart} are of this kind; the restriction of $g$ to
$e_v(E)\subset\auO[1/\ell]$ is determined by $\gamma$, hence so is $\sigma_g$ (in (R$^-$) because
$g^{-1}|_{e_v(E)}=(g|_{e_v(E)})^{-1}$). We put $\sigma_\gamma=\sigma_g$, and $\gamma$ is thick,
resp.\ has sign $\varepsilon(\gamma)$, if $g$ is thick, resp.\ $\varepsilon(\gamma)=\varepsilon(g)$.
This is ``the complex embedding of $E$ attached to $\gamma$ through $\iota$ and $v$'' of
Proposition~\ref{pki:chart}.
\end{definition}

(The statements below are proved for all $g\in\pkjG$, which is stronger than for the embeddings
$\gamma$. Reading (R$^-$) is the one fixed for the paper, in {\cite[\S5.1]{P55}} and
\S\ref{pkd:ss-data}; (R$^+$) is the other natural one, and \S\ref{pkd:ss-data} uses only
$\sigma_\gamma$, $\varepsilon(\gamma)$ and thickness, so either would do there.)

\subsection{Elementary inputs}\label{pkj:ss-inputs}

\begin{lemma}[Sizes]\label{pkj:sizes}
\begin{enumerate}
\item[(i)] If $\ell\mid N_E$, then $\ell\mid q_0$ or $\ell=d$, and $\ell\le Q$.
\item[(ii)] If $\ell\nmid N_E$, then $\ell\ge13$.
\item[(iii)] $[E_v:\Q_\ell]\le[E:\Q]\le h$; in particular $e\le h$.
\item[(iv)] For every prime $\ell$, $2\,\pkilog_\ell h^3\le24s<l_0(s)$.
\end{enumerate}
\end{lemma}

\begin{proof}
We use {\cite[Lemma~13.4]{P46}}: (S1) $\log x\ge10^{14}(q_0+2d)$ and $x\ge(\log b)^{100}\ge128^{100}$;
(S2) $x^{11}\le2^s$ and $r/2\le k$; (S3) $f_E\le x^{12}$; and (D3) $x^5\le r\le3x^5$.

(i) The primes dividing $q_1$ divide $q_0$ ({\cite[\S13.1]{P46}}) and $d$ is prime; so $\ell\mid q_0$
or $\ell=d$, and $\ell\le q_0+2d\le\log x<x\le r\le5rk=Q$, since $k\ge1$.

(ii) $18480=2^4\cdot3\cdot5\cdot7\cdot11$ divides $q_0$, and $q_0\mid q_1$ since $q_1/q_0=3\,l_0(s)\,2^s$ ({\cite[\S13.1]{P46}}); so $2,3,5,7,11$ divide $N_E$.

(iii) $E_v=\Q_\ell e_v(E)$ has degree at most $[E:\Q]=2f_E\le2x^{12}\le x^{40}\le r^8=h$; and
$e\le[E_v:\Q_\ell]$.

(iv) $h^3=r^{24}\le3^{24}x^{120}\le x^{132}\le2^{12s}$ (as $x\ge9$ and $2^s\ge x^{11}$). Hence
$\pkilog_\ell h^3\le\log h^3/\log2\le12s$, and $24s<27s<l_0(s)$ ({\cite[\S13.1]{P46}}).
\end{proof}

\begin{lemma}[Roots of unity of order prime to $\ell$]\label{pkj:roots}
\begin{enumerate}
\item[(i)] Let $L\subset\auQlb$ be a finite extension of $\Q_\ell$ with residue field of $\ell^{f'}$
elements. If $w\in L$ is a root of unity of order $n$ prime to $\ell$, then $n\mid\ell^{f'}-1$.
Moreover, for every $n$ prime to $\ell$, reduction is injective on the group $\mu_n\subset\auQlb$ of
$n$-th roots of unity.
\item[(ii)] If $\ell\nmid n$ and $\theta\in\pkjG$, there is $c\ge0$ with $\theta(w)=w^{\ell^c}$ for all $w\in\mu_n$.
\item[(iii)] $\mu_{\ell^f-1}\subset\auO_{E_v}$, and $\mu_{\ell^f-1}$ is generated by any of its
elements of exact order $\ell^f-1$.
\item[(iv)] If $\ell\nmid N_E$, then $e=1$, and every $\theta\in\pkjG$ that fixes $\mu_{\ell^f-1}$ pointwise
lies in $\Gamma$.
\end{enumerate}
\end{lemma}

\begin{proof}
(i) Let $w\in\mu_n$, $w\ne1$, of order $n'\mid n$. Evaluating $(X^{n'}-1)/(X-1)=\prod_{j=1}^{n'-1}(X-w^j)$
at $X=1$ gives $\prod_{j=1}^{n'-1}(1-w^j)=n'$, a unit since $\ell\nmid n'$; as all $w^j$ are integral,
$1-w$ is a unit, i.e.\ $w\not\equiv1$ modulo the maximal ideal. Hence reduction is injective on
$\mu_n$. If $w\in L$ has order $n$, its reduction lies in the residue field of $L$ and has order $n$,
so $n$ divides the order $\ell^{f'}-1$ of the multiplicative group of that field.

(ii) Let $w_0$ be a primitive $n$-th root of unity; $\mathbb F=\mathbb F_\ell(\bar w_0)$ is a finite
field, and $\bar\theta$ restricts to an automorphism of $\mathbb F$, hence to $x\mapsto x^{\ell^c}$ for some
$c\ge0$. Since $\theta(w_0)\in\mu_n$, its reduction $\bar\theta(\bar w_0)=\bar w_0^{\ell^c}$ is the reduction of
$w_0^{\ell^c}$, so $\theta(w_0)=w_0^{\ell^c}$ by (i); and $\mu_n=\langle w_0\rangle$.

(iii) The residue field of $E_v$ is $\mathbb F_{\ell^f}$, over which $X^{\ell^f-1}-1$ has $\ell^f-1$
simple roots; by Hensel's lemma \cite[II, \S4, Proposition~7, p.~34]{SerreLF} each lifts to a
root in $\auO_{E_v}$, and these lifts are distinct. So $\auO_{E_v}$ contains $\ell^f-1$ roots of
$X^{\ell^f-1}-1$, i.e.\ all of $\mu_{\ell^f-1}$. A cyclic group is generated by any element of maximal
order.

(iv) Since $d\equiv7\pmod8$, $d\equiv3\pmod4$ and $\sqrt{-d}\in\Q(\mu_d)$ (quadratic Gauss sum: for the odd quadratic character $\chi$ of conductor $d$,
$\tau(\chi)^2=\chi(-1)|\tau(\chi)|^2=-d$ by \cite[Lemmas~4.7 and~4.8]{Washington}); so $E=K(\zeta_{q_1})\subseteq\Q(\mu_{N_E})$, and
$E_v\subseteq\Q_\ell(\mu_{N_E})$ (with $\mu_{N_E}\subset\auQlb$). For $\ell\nmid N_E$ the extension
$\Q_\ell(\mu_{N_E})/\Q_\ell$ is unramified \cite[IV, \S4, Proposition~16, p.~77]{SerreLF}, hence so is
$E_v$: $e=1$ and $[E_v:\Q_\ell]=f$. By (iii), $F_0=\Q_\ell(\mu_{\ell^f-1})\subseteq E_v$. By (i) the
reductions of the elements of $\mu_{\ell^f-1}$ are $\ell^f-1$ distinct elements of the residue field of
$F_0$, which therefore has at least $\ell^f$ elements; so $[F_0:\Q_\ell]\ge f=[E_v:\Q_\ell]$ (the degree
of a local field is at least its residue degree), and $F_0=E_v$. If $\theta$ fixes $\mu_{\ell^f-1}$
pointwise, it fixes $F_0=E_v\supseteq e_v(E)$, so $\theta\in\Gamma$.
\end{proof}

\begin{lemma}[Letters and the action on roots of unity]\label{pkj:transfer}
Let $\zeta\in\auQlb$ have exact order $\ell^f-1$, $g_0,g_1\in\pkjG$ and $c\in\Z$ with
$g_1(\zeta)=g_0(\zeta)^c$. Then $g_1(w)=g_0(w)^c$ for all $w\in\mu_{\ell^f-1}$. Let $v_\ell=v_\ell(q_1)$,
$q'=q_1/\ell^{v_\ell}$ (the largest divisor of $q_1$ prime to $\ell$) and $y=e_v(\zeta_{q_1})^{\ell^{v_\ell}}$.
Then $y$ is a primitive $q'$-th root of unity, $\mu_{q'}=\langle y\rangle\subseteq\mu_{\ell^f-1}$, and
$g_1(w)=g_0(w)^c$ for all $w\in\mu_{q'}$.
\end{lemma}

\begin{proof}
Every $w\in\mu_{\ell^f-1}$ is $\zeta^a$ (Lemma~\ref{pkj:roots}(iii)), and $g_1(\zeta^a)=g_1(\zeta)^a=
g_0(\zeta)^{ca}=g_0(\zeta^a)^c$. The root $y$ has order $q_1/\ell^{v_\ell}=q'$, prime to $\ell$, and lies
in $e_v(E)\subseteq E_v$; by Lemma~\ref{pkj:roots}(i) with $L=E_v$ (residue degree $f$), $q'\mid\ell^f-1$.
\end{proof}

\subsection{The three properties of the data}\label{pkj:ss-three}

\begin{lemma}[Letters determine the embedding off the conductor]\label{pkj:letter-det}
Let $\ell\nmid N_E$, $\zeta\in\auQlb$ of exact order $\ell^f-1$, and $g,g'\in\pkjG$ with $g(\zeta)=g'(\zeta)$.
Then $\sigma_g=\sigma_{g'}$; in particular $g$ is thick iff $g'$ is, and $\varepsilon(g)=\varepsilon(g')$.
\end{lemma}

\begin{proof}
By Lemma~\ref{pkj:transfer} (with $c=1$), $g(w)=g'(w)$ for all $w\in\mu_{\ell^f-1}$. Applying this to
$w=g'^{-1}(u)$ and to $w=g^{-1}(u)$ ($u\in\mu_{\ell^f-1}$, a $\pkjG$-stable set) shows that $g g'^{-1}$
and $g'g^{-1}$, and likewise $g'^{-1}g$ (apply $g'^{-1}$ to $g(w)=g'(w)$), fix $\mu_{\ell^f-1}$
pointwise. By Lemma~\ref{pkj:roots}(iv) they lie in $\Gamma$. In reading (R$^-$),
$\sigma_g=\sigma_{g'}$ means $g^{-1}(e_v x)=g'^{-1}(e_v x)$ for all $x$, i.e.\ $g'g^{-1}\in\Gamma$; in
reading (R$^+$) it means $g(e_vx)=g'(e_vx)$, i.e.\ $g'^{-1}g\in\Gamma$.
\end{proof}

\begin{lemma}[Ratios of places]\label{pkj:ratio}
Let $q'$ and $y$ be as in Lemma~\ref{pkj:transfer}, $g_0,g_1\in\pkjG$ and $c\in\Z$ with
$g_1(w)=g_0(w)^c$ for all $w\in\mu_{q'}$. Then $c$ is prime to $q'$ and, with $\varpi$ as in
Convention~\ref{pkj:conv},
\[
\pkjcl_{q'}\bigl(P(g_1)/P(g_0)\bigr)=\varepsilon(g_0)\,\varepsilon(g_1)\,c^{\varpi}\qquad\text{in }(\Z/q')^\times .
\]
\end{lemma}

\begin{proof}
$g_1(y)$ and $g_0(y)$ are primitive $q'$-th roots of unity and $g_1(y)=g_0(y)^c$, so $c$ is a unit
modulo $q'$. Write $\varepsilon_i=\varepsilon(g_i)$, $z=e_v(\zeta_{q_1})$, so $y=z^{\ell^{v_\ell}}$. By
\eqref{pkj:eq-sign} and \eqref{pkj:eq-ratio},
\begin{equation}\label{pkj:eq-places}
\sigma_{g_1}(\zeta_{q_1})^{\varepsilon_1}=\sigma_{g_0}(\zeta_{q_1})^{\varepsilon_0j_{10}},\qquad
j_{10}=P(g_1)/P(g_0).
\end{equation}
\emph{Reading (R$^+$).} Applying $\iota$ to \eqref{pkj:eq-places} gives $g_1(z)^{\varepsilon_1}=
g_0(z)^{\varepsilon_0j_{10}}$; raising to the power $\ell^{v_\ell}$, $g_1(y)^{\varepsilon_1}=g_0(y)^{\varepsilon_0j_{10}}$.
With $g_1(y)=g_0(y)^c$ and $g_0(y)$ primitive of order $q'$: $c\varepsilon_1\equiv\varepsilon_0j_{10}$, i.e.\
$j_{10}\equiv\varepsilon_0\varepsilon_1c\pmod{q'}$.

\emph{Reading (R$^-$).} Applying $\iota$ gives $g_1^{-1}(z)^{\varepsilon_1}=g_0^{-1}(z)^{\varepsilon_0j_{10}}$,
hence $g_1^{-1}(y)^{\varepsilon_1}=g_0^{-1}(y)^{\varepsilon_0j_{10}}$. Put $u=g_0^{-1}(y)\in\mu_{q'}$; then
$g_1(u)=g_0(u)^c=y^c$, so $g_1^{-1}(y)^c=u=g_0^{-1}(y)$. Hence $g_1^{-1}(y)^{\varepsilon_1}=
g_1^{-1}(y)^{c\varepsilon_0j_{10}}$, and $g_1^{-1}(y)$ is primitive of order $q'$: $\varepsilon_1\equiv
c\varepsilon_0j_{10}$, i.e.\ $j_{10}\equiv\varepsilon_0\varepsilon_1c^{-1}\pmod{q'}$.
\end{proof}

\begin{lemma}[One sign on the $b$-places at primes of the conductor]\label{pkj:signs}
Let $\ell\mid N_E$ and let $g,g'\in\pkjG$ be thick. Then $\varepsilon(g)=\varepsilon(g')$.
\end{lemma}

\begin{proof}
Suppose $\varepsilon(g)=-\varepsilon(g')$. By Lemma~\ref{pkj:sizes}(i), $\ell\mid q_0$ or $\ell=d$.

\emph{Case $\ell\mid q_0$.} By \eqref{pkj:eq-sign}, $\sigma_g(\sqrt{-d})=-\sigma_{g'}(\sqrt{-d})$. Applying
$\iota$ and then $g'$ (reading (R$^-$)), resp.\ $g'^{-1}$ (reading (R$^+$)), we get an element
$\theta\in\pkjG$ ($\theta=g'g^{-1}$, resp.\ $\theta=g'^{-1}g$) with $\theta(e_v\sqrt{-d})=-e_v\sqrt{-d}$. Since $d\equiv3
\pmod4$, $\beta=(1+\sqrt{-d})/2\in E$ satisfies $\beta(1-\beta)=(1+d)/4$. Put $x=e_v(\beta)\in\auQlb$;
then $\theta(x)=1-x$, so $|1-x|=|x|$ ($\theta$ is an isometry), and
$|x|^2=|x(1-x)|=|(d+1)/4|_\ell<1$, because $4\ell\mid d+1$ for every prime $\ell\mid q_0$
({\cite[\S13.1]{P46}}). So $|x|<1$ and $|1-x|<1$, contradicting $1=|x+(1-x)|\le\max(|x|,|1-x|)$.

\emph{Case $\ell=d$.} Then $\ell\nmid q_1$ (as $d\nmid q_0$ and the primes of $q_1$ divide $q_0$), so
$q'=q_1$. By Lemma~\ref{pkj:roots}(ii), $g'g^{-1}$ acts on $\mu_{q_1}$ as $w\mapsto w^{\ell^c}$ for some
$c\ge0$; as $g(\mu_{q_1})=\mu_{q_1}$, $g'(w)=g'g^{-1}(g(w))=g(w)^{\ell^c}$ for $w\in\mu_{q_1}$. By
Lemma~\ref{pkj:ratio} (with $g_0=g$, $g_1=g'$), $\pkjcl_{q_1}(P(g')/P(g))=\varepsilon(g)\varepsilon(g')
d^{\varpi c}=-d^{\varpi c}$. Reducing modulo $q_0\mid q_1$, and using $\pkjcl_{q_0}(P(g))=\pkjcl_{q_0}(P(g'))
=-1$ (both are $b$-places, (D5)), we get $1=-d^{\varpi c}$ in $(\Z/q_0)^\times$, hence $d^{c}\equiv-1
\pmod{q_0}$. But $d\equiv q_0/2-1\pmod{q_0}$ and $4\mid q_0$ give $d^2\equiv q_0\cdot\frac{q_0}4-q_0+1
\equiv1$, so $d^c\equiv1$ or $d^c\equiv d\equiv q_0/2-1\pmod{q_0}$, and neither is $-1$ (as $q_0>2$ and
$q_0/2\not\equiv0$). Contradiction.
\end{proof}

\begin{lemma}[No three thick letters on a short arc]\label{pkj:arc}
Let $\ell\le Q$, $\zeta\in\auQlb$ of exact order $\ell^f-1$, $a,a'\in[1,\pkilog_\ell h^3]$, and
$g_0,g_1,g_2\in\pkjG$ with $g_1(\zeta)=g_0(\zeta)^{\ell^a}$ and $g_2(\zeta)=g_0(\zeta)^{\ell^{a+a'}}$. Then
$g_0$, $g_1$, $g_2$ are not all thick.
\end{lemma}

\begin{proof}
Suppose they are; put $P_i=P(g_i)\in\pkjPb$ and let $q'$ be as in Lemma~\ref{pkj:transfer}. By
Lemma~\ref{pkj:transfer}, $g_1(w)=g_0(w)^{\ell^a}$ and $g_2(w)=g_0(w)^{\ell^{a+a'}}$ for $w\in\mu_{q'}$;
by Lemma~\ref{pkj:ratio}, with signs $\pm$ (each $+1$ or $-1$),
\[
\pkjcl_{q'}(P_1/P_0)=\pm\ell^{\varpi a},\qquad \pkjcl_{q'}(P_2/P_0)=\pm\ell^{\varpi(a+a')} .
\]
If $\varpi=1$, take $\sigma=P_0$, $\tau=P_1$, $\rho=P_2$, $i=a$, $j=a+a'$: then $\pkjcl_{q'}(\tau/\sigma)=
\pm\ell^i$ and $\pkjcl_{q'}(\rho/\sigma)=\pm\ell^j$. If $\varpi=-1$, take $\sigma=P_2$, $\tau=P_1$,
$\rho=P_0$, $i=a'$, $j=a+a'$: then $\pkjcl_{q'}(\rho/\sigma)=\pkjcl_{q'}(P_2/P_0)^{-1}=\pm\ell^{j}$ and
$\pkjcl_{q'}(\tau/\sigma)=\pkjcl_{q'}(P_1/P_0)\,\pkjcl_{q'}(P_2/P_0)^{-1}=\pm\ell^{a+a'-a}=\pm\ell^{i}$.
In both cases $\sigma,\tau,\rho\in\pkjPb$ and $0<i<j\le2\pkilog_\ell h^3<l_0(s)$ by
Lemma~\ref{pkj:sizes}(iv). Since $\ell\le Q=5r(r+1-m)$ and $q'$ is the largest divisor of $q_1$
prime to $\ell$, this contradicts the sparseness of $\pkjPb$ ((D5), {\cite[Definition~13.2]{P46}}).
\end{proof}

\subsection{Proof of Proposition~\ref{pki:chart}}\label{pkj:ss-proof}

\begin{theorem}[Proposition~\ref{pki:chart}]\label{pkj:chart}
Let $Q=kt$, $h=r^8$, and let thickness and $\varepsilon(\gamma)$ be as in Convention~\ref{pkj:conv}
(either reading). For every prime $\ell$, every $v\mid\ell$, every large $\auO$ and every
$\zeta\in\auO$ of exact order $\ell^f-1$ (such $\zeta$ exist), there are a set $T$ of letters and a sign
$s$ on letters for which the letters read thickness for $(Q,h)$ (Definition~\ref{pki:def-reads}).
\end{theorem}

\begin{proof}
Elements $\zeta\in\auO$ of exact order $\ell^f-1$ exist by Lemma~\ref{pkj:roots}(iii), as
$\auO\supseteq\auO_{E_v}$. For an embedding $\gamma\colon\auO\to\C_\ell$ choose $g_\gamma\in\pkjG$
restricting to $\gamma$ (Convention~\ref{pkj:conv}); then $\gamma(\zeta)=g_\gamma(\zeta)$,
$\sigma_\gamma=\sigma_{g_\gamma}$, and the letter $\gamma(\zeta)$ is an $(\ell^f-1)$-th root of unity. Put
\[
T=\{\gamma(\zeta):\ \gamma\ \text{thick}\}.
\]
\emph{Case $\ell\nmid N_E$.} For a letter $z$ let $s(z)=\varepsilon(\gamma)$ if $z=\gamma(\zeta)$ for some
$\gamma$, and $s(z)=1$ otherwise; by Lemma~\ref{pkj:letter-det} (applied to $g_\gamma$, $g_{\gamma'}$),
$\gamma(\zeta)=\gamma'(\zeta)$ implies $\sigma_\gamma=\sigma_{\gamma'}$, so $s$ is well defined. For every
$\gamma$: if $\gamma$ is thick, $\gamma(\zeta)\in T$ by definition; if $\gamma(\zeta)\in T$, then
$\gamma(\zeta)=\gamma'(\zeta)$ with $\gamma'$ thick, so $\sigma_\gamma=\sigma_{\gamma'}$ and $\gamma$ is thick;
and $\varepsilon(\gamma)=s(\gamma(\zeta))$. Moreover $e=1$ (Lemma~\ref{pkj:roots}(iv)) and $\ell\ge13$
(Lemma~\ref{pkj:sizes}(ii)).

\emph{Case $\ell\mid N_E$.} Then $\ell\le Q$ (Lemma~\ref{pkj:sizes}(i)). Let $s$ be the constant sign
$\varepsilon(\gamma_0)$ for a thick $\gamma_0$ if one exists, and $s\equiv1$ otherwise. For thick
$\gamma$: $\gamma(\zeta)\in T$ (so ``$\gamma$ thick $\Leftrightarrow\gamma(\zeta)\in T$'' holds, both sides
being true), and $\varepsilon(\gamma)=\varepsilon(\gamma_0)=s(\gamma(\zeta))$ by Lemma~\ref{pkj:signs}. Since
$\ell\le Q$, the clause ``for every $\gamma$ if $\ell>Q$'' is void, and $\ell\le Q$ with $s$ constant.

\emph{The remaining clauses, in both cases.} $e\le h$ by Lemma~\ref{pkj:sizes}(iii). Let $\ell\le Q$ and
suppose a letter $z$ has $z$, $z^{\ell^a}$, $z^{\ell^{a+a'}}\in T$ with $a,a'\in[1,\pkilog_\ell h^3]$:
$z=\gamma_0(\zeta)$, $z^{\ell^a}=\gamma_1(\zeta)$, $z^{\ell^{a+a'}}=\gamma_2(\zeta)$ with $\gamma_0,\gamma_1,
\gamma_2$ thick. Then $g_i=g_{\gamma_i}$ are thick, $g_1(\zeta)=g_0(\zeta)^{\ell^a}$ and
$g_2(\zeta)=g_0(\zeta)^{\ell^{a+a'}}$, contradicting Lemma~\ref{pkj:arc}.
\end{proof}

\begin{remark}[The uses of Proposition~\ref{pki:chart}]\label{pkj:uses}
Proposition~\ref{pki:chart} is used in \S\ref{pki:ss-deduce} (Lemma~\ref{pki:letterwin},
Theorem~\ref{pki:crys} and the proof of {\cite[Theorem~5.8]{P55}}), only through the clauses of
Definition~\ref{pki:def-reads} for the thickness and signs of Convention~\ref{pkj:conv}; and
\S\ref{pkd:ss-data} (Theorem~\ref{pkd:data}) uses the same $\sigma_\gamma$, $\varepsilon(\gamma)$ and
thickness, through the jumps $J_{\sigma_\gamma}$. Theorem~\ref{pkj:chart} gives these clauses for every
large $\auO$ and every admissible choice of $\zeta$, as the deduction requires; the quantities $e$, $f$
are those of $E_v$, as in \S\ref{pki:ss-arch} (with $\Gamma'=\Gamma$ in Definition~\ref{pkp:def-chart},
Lemma~\ref{pkp:chart-exists}, they are also those of the Kummer chart). The theorem holds for both
readings of $\sigma_\gamma$, so it applies whichever reading \S\ref{pkd:ss-data} and the Hodge--Tate
convention of $L_{\mathrm{main}}$ adopt, as long as the two agree; both use (R$^-$) ({\cite[\S5.1]{P55}}, \S\ref{pkd:ss-data}).
\end{remark}



\section{The four cases}\label{C4:g7}
\providecommand{\auO}{\mathcal O}
\providecommand{\auQlb}{\overline{\Q}_\ell}
\providecommand{\pllZf}{Z_f}
\providecommand{\pllrk}{\operatorname{rk}}
\providecommand{\plldef}{\delta^{\vee}}
\providecommand{\pglen}{\operatorname{len}}
\providecommand{\pkiWin}{\operatorname{Win}}
\providecommand{\pkiC}{\mathcal C}
\providecommand{\pkilog}{\operatorname{lg}}
\providecommand{\pkpA}{A_{\mathrm{inf}}}
\providecommand{\pkpAmb}{\mathbb A}
\providecommand{\pkcPt}{\mathrm{Pt}}
\providecommand{\pkcimp}{\mathrm{imp}}
\providecommand{\pkcpd}{\mathrm{pd}}
\providecommand{\pkcdep}{\mathrm{dep}}
\providecommand{\pkcsh}{\mathrm{sh}}
\providecommand{\pkcNM}{\mathrm{NM}}
\providecommand{\pkcnw}{\mathrm{nw}}
\providecommand{\pkcSig}{\Sigma}

\subsection{The four cases of the local-loss engine}\label{pkc:sec}

This section proves the core propositions of \S\ref{pki:ss-arch}, one per subsection, in the order
Proposition~\ref{pki:generic} (\S\ref{pkc:ss-generic}), Proposition~\ref{pki:tame}
(\S\ref{pkc:ss-tame}), Proposition~\ref{pki:sparse} (\S\ref{pkc:ss-sparse}),
Proposition~\ref{pki:any} (\S\ref{pkc:ss-any}), and then treats the remaining assertions of
Proposition~\ref{pki:data} (\S\ref{pkd:ss-data}, Theorem~\ref{pkd:data}; the membership $L_c^*\in\pkiC(U',K^{-1})$ needed in
case (a) uses the input (I6), which is Lemma~\ref{pkw:dual}). The tools are those of \S\ref{pkp:sec},
\S\ref{pkm:sec} and \S\ref{pkn:sec}, which we cite by label and do not reprove. Where a case needs a
statement about points, lift ideals or charts that is not among the statements of those
sections, we state the exact form needed in \S\ref{pkc:ss-inputs} and prove the case from it; these
statements, the inputs (I1)--(I4) and (T1)--(T4), are proved in \S\S\ref{pkt:sec}--\ref{pkw:sec}.

\emph{Dependencies.} Proved in full from the statements of \S\S\ref{pkp:sec}--\ref{pkn:sec} and the inputs
(I1), (I2) of \S\ref{pkc:ss-inputs}: Proposition~\ref{pki:generic}(i). Proved from these and the
inputs (I3), (I4): Proposition~\ref{pki:generic}(ii). Proved from (I1), (I2) and the tame inputs
(T1)--(T4): Proposition~\ref{pki:tame}. The inputs are proved elsewhere: (I1) and (I3) in
\S\ref{pkt:sec} (Propositions~\ref{pkt:I1}, \ref{pkt:I3}), (I2) in \S\ref{pke:sec}
(Theorem~\ref{pke:I2}), (I4) in \S\ref{pkh:sec} (Theorem~\ref{pkh:I4}), (T1)--(T4) in \S\ref{pkw:sec}.
Propositions~\ref{pki:sparse}, \ref{pki:any} and the rest of Proposition~\ref{pki:data} are proved in
\S\ref{pkd:sec}.

\emph{Conventions.} As in \S\ref{pkp:sec}: a Kummer chart of $\Gamma=D_v$ is fixed
(Definition~\ref{pkp:def-chart}), with ramification index $e$, residue degree $f$, subgroup $\Gamma'$
and $d_v$ as there; letters, the Frobenius $\sigma(i)=i^\ell$ on letters and $i^{(-a)}=\sigma^{-a}(i)$
are as in Definition~\ref{pkp:def-letters}; orders $o\colon\Z_{\ge0}\times\Z\to\Z$, the modules
$L_i(o)$, and $\rho_t$, $\mathrm{std}_l$, $\pkcsh_{l,n}$ are as in Definition~\ref{pkp:def-orders} with
the modules $\mathcal N$ of Definition~\ref{pkn:def-N}; sides, extensions $X_S$, pairings modulo
$\ell^m$, Tate-dual sides $T_P$, safe moves and $\mathrm{SCC}_N$ are as in
Definitions~\ref{pkp:def-side}--\ref{pkp:def-scc}; $\pkcimp$ and $\pkcpd$ are as in
\S\ref{pkm:ss-state}, and $\pkcdep_{t,\eta,l}$, ``standard point'' and ``stop'' as in
Definition~\ref{pkp:def-points}. For a side $S$, a letter $i$, $c\in C_S$, $e\ge0$ and orders $o,o'$
we write
\[
\pkcPt_S(c,i,e;\,o,o')
\]
for the set of points of $S$ over $\ell^ec$ at $i$ with orders $o$ and depth $o'$
(Definition~\ref{pkp:def-points}). All inequalities between orders are required only at the pairs
$(a,u)$ with $u\in U(i^{(-a)})$, the jump set of the side at the letter $a$ steps before $i$; ``at
level $a$'' refers to the first argument. Thus, explicitly,
\[
\rho_t(a,u)=t[a=0]-u,\qquad \pkcdep_{t,\eta,1}(a,u)=-u+(t+\eta)[a=0],\qquad
\pkcpd_{t,\eta,1,\eta_0}(a,u)=-u+(t+\eta)[a=0]+\eta_0[a=1].
\]

\subsubsection{Windows, steps, and the inputs used}\label{pkc:ss-inputs}

\paragraph*{Elementary facts}

\begin{remark}[Monotonicity]\label{pkc:mono}
If $o_1\le o$ and $d_1\le d$ (at the jumps), then $\pkcPt_S(c,i,e;o,d)\subseteq\pkcPt_S(c,i,e;o_1,d_1)$,
since $\mathrm{Fil}^n$ decreases in $n$ and so $L_i(o)\subseteq L_i(o_1)$. In particular, as
$\pkcpd_{t,\eta,1,\eta_0}\ge\pkcdep_{t,\eta,1}$, a point with orders $\rho_{-\eta}$ (or larger) and
depth $\pkcpd_{t,\eta,1,\eta_0}$ (or larger) is a standard point for $(t,\eta,1)$.
\end{remark}

For $n\ge0$ and $b\in\Z$ let $\pkcSig_n(b)=\{y\in\Z:\ |y|\le n\text{ or }|y-b|\le n\text{ or
}|y+b|\le n\}$; it has at most $3(2n+1)$ elements, and $\pkcNM(\ell,n,b)$ (Definition preceding
Definition~\ref{pki:def-reads}) says exactly that no nonzero element of $\pkcSig_n(b)$ is divisible by
$\ell$.

\begin{lemma}[Windows]\label{pkc:win}
Let $\eta,b,W\ge0$, a sign $\varepsilon$, $u\in\pkiWin(\eta,-\varepsilon,b,W)$ and
$u'\in\pkiWin(\eta,\varepsilon,b,W)$. Then
\begin{enumerate}
\item[(i)] $u\le b+W+\eta$ and $u'\le b+W+\eta$;
\item[(ii)] $-(u+u')\le b+W+2\eta$, and if $\eta\le b$ then $-u-u'-1\le b+W+\eta$;
\item[(iii)] if $s,t'\ge0$, $n\ge\max(s,t')+\eta+W$ and $-u-s\le y<t'-u$, then
$y\in\pkcSig_n(b)$; the same holds for $u'$.
\end{enumerate}
\end{lemma}

\begin{proof}
(i) Elements of the low window are $\le\eta$, those of a far window are $\le b+W$. (ii) If both
$u,u'$ are in the low window, $-u-u'\le2\eta$, which is $\le b+W+2\eta$, and $\le b+\eta$ if $\eta\le b$.
If exactly one is in a far window, say $u'$, then $-u'\le b+W$ (if the far window of $u'$ is negative)
or $-u'\le-b\le0$ (if positive), and $-u\le\eta$; so $-u-u'\le b+W+\eta$. If both are in far windows,
the windows have opposite signs, so $u=\delta x$, $u'=-\delta x'$ with $\delta=\pm1$, $x,x'\in[b,b+W]$,
and $|u+u'|=|x-x'|\le W$. In all cases $-u-u'\le b+W+2\eta$, and if $\eta\le b$ then $-u-u'\le
b+W+\eta$, a fortiori $-u-u'-1\le b+W+\eta$. (iii) If $|u|\le\eta$
then $-\eta-s\le y<t'+\eta$, so $|y|\le\max(s,t')+\eta\le n$. If $u\in[b,b+W]$ then $-W-s\le
y+b<t'$, so $|y+b|\le n$. If $-u\in[b,b+W]$ then $-s\le y-b<W+t'$, so $|y-b|\le n$.
\end{proof}

\begin{lemma}[Jumps of the Tate-dual side]\label{pkc:dualjumps}
Let $P$ pair sides $S,S'$, and at a letter $i$ let $U_S(i)\subseteq\pkiWin(\eta,-\varepsilon,b,W)$ and
$U_{S'}(i)\subseteq\pkiWin(\eta,\varepsilon,b,W)$ with $\eta\le b$. Then every jump of $T_P$ at $i$ is
$\le b+W+\eta$. If moreover $6(a+\eta+W)<b$, $a\ge0$, $\eta\ge1$, then every jump $s$ of $T_P$ at $i$
satisfies $-a-s\le y<a-s\Rightarrow y\in\pkcSig_{6(a+\eta+W)}(b)$.
\end{lemma}

\begin{proof}
The jumps of $T_P$ at $i$ are $0$ and the $s=-u-u'-1$ with $u\in U_S(i)$, $u'\in U_{S'}(i)$
(Definition~\ref{pkp:def-pairing}); the first claim is Lemma~\ref{pkc:win}(ii). For the second put
$n=6(a+\eta+W)$. If $s=0$, $|y|\le a\le n$. Otherwise the case analysis of Lemma~\ref{pkc:win}(ii)
shows: if both $u,u'$ are low, $|s|\le2\eta+1$; if both are far, $|s|\le W+1$; if exactly one is far,
$s=\mp x-v-1$ with $x\in[b,b+W]$ and $|v|\le\eta$, so $|s\pm b|\le W+\eta+1$. Hence $s$ lies within
$m:=2\eta+W+1$ of $0$, $b$ or $-b$, and $y\in[-a-s,a-s)$ lies within $a+m\le n$ (as $\eta\ge1$) of
$0$, $-b$ or $b$.
\end{proof}

\paragraph*{The inputs}

The cases use, besides the statements of \S\S\ref{pkp:sec}--\ref{pkn:sec} cited by label, the following
statements. (I1) is the form of Lemma~\ref{pkp:tau}(iii) that the cases need. The four inputs are proved
elsewhere: (I1) is Proposition~\ref{pkt:I1}, (I2) is Theorem~\ref{pke:I2}, (I3) is
Proposition~\ref{pkt:I3}, (I4) is Theorem~\ref{pkh:I4}. For $n\in\Z$ let $v_\ell(n)$ be the $\ell$-adic valuation, with
$v_\ell(0)=0$, and put $\mathrm{cost}(a,n)=v_\ell(n)+d_v+3a$.

\begin{enumerate}
\item[(I1)] \emph{Points from killers.} If $\ell^\kappa E^t$ kills the cokernel of a side $S$ at $i$
(Definition~\ref{pkm:def-J}), then for every $c\in C_S$ the set $\pkcPt_S(c,i,\kappa;\rho_{-t},\rho_{-t})$
is nonempty.
\item[(I2)] \emph{Deepening.} Let $x\in\pkcPt_S(c,i,e;o,d)$ and let $I$ be a finite set of pairs
$(a,n)\in\Z_{\ge0}\times\Z$ containing every $(a,n)$ with $d(a,u)\le n<d'(a,u)$ for some $u\in
U(i^{(-a)})$. If $e+\sum_{(a,n)\in I}\mathrm{cost}(a,n)\le e'$, then
$\pkcPt_S(c,i,e';\pkcimp(o,d'),d')\ne\emptyset$.
\item[(I3)] \emph{The rule with a dual point.} Let $P$ pair $S,S'$, $l,\kappa\ge0$, and let $\nu$ be a
point of $T_P$ over $\ell^\kappa\cdot1$ at $i^{(-l)}$ with orders $o_\nu$, where $o_\nu(a,\cdot)$ is
nonincreasing for each $a$. Let $w\in L_i(o_w)$ (for $X_S$) lie over $0$ and let $y$ be a point of
$S'$ over some $\ell^{e'}c'$ at $i$ with depth $d_y$. If for all $a\ge0$, $u\in U_S(i^{(-a)})$,
$u'\in U_{S'}(i^{(-a)})$
\[
0<\pkcsh_{l,0}(o_\nu)(a,-u-u'-1)+o_w(a,u)+d_y(a,u'),
\]
then the move of the first side by $w$ with $y$ at rest is $\kappa$-safe.
\item[(I4)] \emph{Few low jumps: the deep constants.} Let the chart be unramified, $\ell\ge5$,
$a,\eta,W\ge0$, $\eta\ge1$, and for each letter $i$ let $b_i\ge0$ and a sign $\varepsilon_i$ satisfy
$b_i+W+\eta\le(\ell-3)a$, $6(a+\eta+W)<b_i$, $\pkcNM(\ell,6(a+\eta+W),b_i)$. Let $S'$ be a side with
$U_{S'}(i)\subseteq\pkiWin(\eta,\varepsilon_i,b_i,W)$, and let $n_i$, $r_i$ be integers with $n_i\ge
f\sum_{u\in[-\eta,\eta]}m_i(u)$ and $r_i\ge f\sum_{u\in U_{S'}(i)}m_i(u)$, where $m_i(u)$ is the sum
over the embeddings $\gamma$ of letter $i$ of the multiplicity of the jump $u$ of the lattice of $S'$
at $\gamma$. Define the depth
\[
\delta_i=\rho_{2a+\eta}\ \text{if }\varepsilon_i=1,\qquad
\delta_i(a',u)=\pkcnw_i(a',u)=\begin{cases}b_i+W+a+\eta&a'=0\\ W+1-u+2a[|u|\le\eta]&a'=1\\
-u&a'\ge2\end{cases}\ \text{if }\varepsilon_i=-1 .
\]
Then the set $V$ of $c'\in C_{S'}$ such that $\pkcPt_{S'}(c',i,0;\rho_0,\delta_i)\ne\emptyset$ for
every letter $i$ is a $\Z_\ell$-submodule of $C_{S'}$ with
$\pglen_{\Z_\ell}(C_{S'}/V)\le11\sum_i\bigl(n_i(a+1)+r_i(W+1)\bigr)$.
\end{enumerate}

\begin{lemma}[Free steps]\label{pkc:free}
Let $x\in\pkcPt_S(c,i,e;o,d)$, $n\ge0$, $b\in\Z$ with $\pkcNM(\ell,n,b)$, and let $d'$ be orders such
that every pair $(a,y)$ with $d(a,u)\le y<d'(a,u)$ for some $u\in U(i^{(-a)})$ has $a=0$ and $y\in
\pkcSig_n(b)$. Then $\pkcPt_S(c,i,e+3(2n+1)d_v;\pkcimp(o,d'),d')\ne\emptyset$. If the chart is
unramified, $\pkcPt_S(c,i,e;\pkcimp(o,d'),d')\ne\emptyset$.
\end{lemma}

\begin{proof}
Apply (I2) with $I=\{0\}\times\pkcSig_n(b)$. For $y\in\pkcSig_n(b)$, $v_\ell(y)=0$: if $y=0$ by
convention, otherwise by $\pkcNM(\ell,n,b)$. So $\mathrm{cost}(0,y)=d_v$ and $\sum_I\mathrm{cost}\le
3(2n+1)d_v$. If the chart is unramified, $d_v=0$.
\end{proof}

\begin{lemma}[From sides to lattices]\label{pkc:cup}
Let $\mathcal C$ be a class of lattices $Y$, each crystalline on $\Gamma'$ with jumps at each embedding
$\gamma$ in a set $U(\gamma)$, and let $L$ be crystalline on $\Gamma'$ with jumps in $U'(\gamma)$. For a
lattice $Y$ with jump sets $U$ let $S(Y)$ be its \emph{universal side}: lattice $Y$, $C_{S(Y)}=
\pllZf(Y)$ with trivial action, $\alpha=\mathrm{id}$, and jump set $U_i=\{0\}\cup\bigcup_{\gamma}U(\gamma)$
at the letter $i$ (union over the embeddings $\gamma$ of letter $i$). Let $N\ge0$ and $x\in\pllZf(L)$.
If $\mathrm{SCC}_N(c,x)$ holds for every $Y\in\mathcal C$, every pairing modulo $\ell^m$ (any $m$) of
$S(Y)$ and $S(L)$ and every $c\in\pllZf(Y)$, then $\ell^Nx\in L^{\perp\mathcal C}$.
In particular, if this holds for all $x\in\pllZf(L)$, then $\mathcal C$ has cup exponent $N$ on $L$,
and if it holds with $N=0$ for all $x$ in an $\auO$-submodule (or $\Z_\ell$-submodule) $V$ of
$\pllZf(L)$, then $\plldef_{\mathcal C}(L)\le\pglen(\pllZf(L)/V)$.
\end{lemma}

\begin{proof}
$\pllZf(Y)$ and $\pllZf(L)$ are finite free ({\cite[Lemma~5.6]{P55}}), and $U_i$ contains $0$ and all
jumps at the embeddings of letter $i$, so $S(Y)$, $S(L)$ are sides (Definition~\ref{pkp:def-side}).
Let $Y\in\mathcal C$, $m\ge1$, $(\psi,c_\psi)$ a pairing modulo $\ell^m$ of $Y$ and $L$ in the sense of
{\cite[Definition~5.1(4)]{P55}} (with $(Y\otimes L)^*=\operatorname{Hom}(Y\otimes L,\auO(1))$, as the
values in $\Z/\ell^m(1)$ require), and $a\in\pllZf(Y)$. The same data is a pairing $P$ modulo $\ell^m$
of $S(Y)$ and $S(L)$ in the sense of Definition~\ref{pkp:def-pairing}, with $P=b_\psi$. Put
$\beta(g,h)=b_\psi(a(g),g\,\ell^Nx(h))$, the cup product of the continuous cocycles $a$ and $\ell^Nx$
through the $\Gamma$-equivariant pairing $b_\psi$; it is a continuous $2$-cocycle with values in
$\Z/\ell^m(1)\cong\mu_{\ell^m}$. By $\mathrm{SCC}_N(a,x)$ its image in $C^2(\Gamma,\auQlb^\times)$ is the
coboundary of a continuous cochain. The Kummer sequence $1\to\mu_{\ell^m}\to\auQlb^\times
\xrightarrow{\ \ell^m\ }\auQlb^\times\to1$ and $H^1(\Gamma,\auQlb^\times)=0$ (Hilbert's Theorem~90 for
the absolute Galois group of $E_v$) show that $H^2(\Gamma,\mu_{\ell^m})\to H^2(\Gamma,\auQlb^\times)$ is
injective. So the class of $\beta$ vanishes, i.e.\ $\beta$ is a coboundary in $C^2(\Gamma,
\Z/\ell^m(1))$. As $Y,m,\psi,a$ were arbitrary, $\ell^Nx\in L^{\perp\mathcal C}$ ({\cite[Definition~5.1(5)]{P55}}).
The last two assertions follow from the definitions of cup exponent and of $\plldef_{\mathcal C}$:
in the second case $V\subseteq\pllZf(L)\cap L^{\perp\mathcal C}$.
\end{proof}

\paragraph*{Two walks}

The two lemmas below are the common core of Propositions~\ref{pki:generic}(i) and \ref{pki:tame}.
Throughout, $P$ pairs sides $S$, $S'$, and at every letter $i$ an integer $b_i\ge0$ and a sign
$\varepsilon_i$ are given with
\begin{equation}\label{pkc:eq-win}
U_S(i)\subseteq\pkiWin(\eta,-\varepsilon_i,b_i,W),\qquad U_{S'}(i)\subseteq\pkiWin(\eta,\varepsilon_i,
b_i,W),\qquad \pkcNM(\ell,n,b_i),\qquad n=6(a+\eta+W),
\end{equation}
for fixed integers $a,\eta,W\ge0$. If $\varepsilon_i=1$, every jump of $S$ at $i$ is $\le\eta$ and
every jump of $S'$ at $i$ is $\ge-\eta$; if $\varepsilon_i=-1$, the reverse. (This is how a far
window is ``negative'' or ``positive''.)

\begin{lemma}[The walk with a gap below zero]\label{pkc:gap}
Assume \eqref{pkc:eq-win}, that for every letter $i$ the element $\ell^\kappa E^a$ kills the cokernels of
$S$, $S'$ and $T_P$ at $i$, and that no jump of $S$ or $S'$ at any letter lies in $[-a,0)$. Then
$\mathrm{SCC}_N(c,c')$ for all $c\in C_S$, $c'\in C_{S'}$ and all $N\ge3\kappa+2\cdot3(2n+1)d_v$.
\end{lemma}

\begin{proof}
Put $\sigma=3(2n+1)d_v$ and $e=\kappa+\sigma$.

\emph{Points without pole.} Fix $i$. By (I1), $\pkcPt_S(c,i,\kappa;\rho_{-a},\rho_{-a})\ne\emptyset$.
Let $d'=\pkcpd_{a,\eta,1,0}$. A pair $(a',y)$ with $\rho_{-a}(a',u)\le y<d'(a',u)$ has $a'=0$ (for
$a'\ge1$ both sides are $-u$) and $-u-a\le y<a+\eta-u$, so $y\in\pkcSig_n(b_i)$ by
Lemma~\ref{pkc:win}(iii) (with $s=a$, $t'=a+\eta$, and $a+\eta+\eta+W\le n$). By
Lemma~\ref{pkc:free} there is $x(i)\in\pkcPt_S(c,i,e;\pkcimp(\rho_{-a},d'),d')$. We claim
$\pkcimp(\rho_{-a},d')\ge\rho_0$ at the jumps. At level $a'\ge1$: $\rho_{-a}(a',u)=-u$; if $-u>0$ the
improved order is $d'(a',u)=-u$, otherwise $\min(-u,0)=-u$. At level $0$: if $\rho_{-a}(0,u)=-a-u>0$,
the improved order is $d'(0,u)=a+\eta-u\ge-u$; otherwise $u\ge-a$, so $u\ge0$ by the gap, and
$\min(a+\eta-u,0)\ge-u$. Hence $x(i)\in\pkcPt_S(c,i,e;\rho_0,\pkcpd_{a,\eta,1,0})$ (Remark~\ref{pkc:mono}).
In the same way, with $\ell^\kappa E^a$ killing the cokernel of $S'$, there is
$x'(i)\in\pkcPt_{S'}(c',i,e;\rho_0,\pkcpd_{a,\eta,1,0})$. By Remark~\ref{pkc:mono}, $x(i)$ and $x'(i)$
are standard points for $(a,\eta,1)$ (as $\rho_0\ge\rho_{-\eta}$).

\emph{Stops.} Fix $i$. If $\varepsilon_i=1$, the jumps of $S$ at $i$ are $\le\eta$ and those of $S'$
at $i$ are $\ge-\eta$. Apply Lemma~\ref{pkm:plain} with $l=1$, $t=a$, $\eta$, $\eta_0=0$, the killer
$\ell^\kappa E^a$ of $T_P$ at $i$, $x=x(i)$ (orders $\rho_0\ge\rho_{-\eta}$, depth
$\pkcpd_{a,\eta,1,0}$), $x'=x'(i)$ (orders $\rho_0\ge\rho_{-\eta}$) and $x'_0=x'(i^{(-1)})$ (orders
$\rho_0=\rho_{-\eta_0}$): $x(i)$ is a stop for $(\eta,1)$ and the move of the second side by
$\varphi x'(i^{(-1)})-x'(i)$ with $x(i)$ at rest is $\kappa$-safe. If $\varepsilon_i=-1$, the same
lemma with the roles of $S$, $S'$ exchanged gives that $x'(i)$ is a stop and the move of the first side
by $\varphi x(i^{(-1)})-x(i)$ with $x'(i)$ at rest is $\kappa$-safe. So (St) of
Theorem~\ref{pkm:walk} holds at every letter with $l(i)=1$.

\emph{The walk.} Corollary~\ref{pkm:walk-all} with $t=a$, $\eta$, $e=e'=\kappa+\sigma$ and the killer
$\ell^\kappa E^a$ of $T_P$ gives $\mathrm{SCC}_N(c,c')$ for $N\ge e+e'+\kappa=3\kappa+2\sigma$.
\end{proof}

\begin{lemma}[The walk with nonpositive jumps on one side]\label{pkc:nonpos}
Let the chart be unramified. Assume \eqref{pkc:eq-win}, that $\ell^\kappa E^a$ kills the cokernels of $S$
and of $T_P$ at every letter, and that $U_{S'}(i)\subseteq[-\eta,0]$ for every letter $i$. Then
$\mathrm{SCC}_N(c,c')$ for all $c,c'$ and all $N\ge2\kappa$.
\end{lemma}

\begin{proof}
Put $t=a+\eta$; note $a+t+\eta+W\le n$ and $a+t+\eta\le n$. As the chart is unramified, $d_v=0$, and
Lemma~\ref{pkc:free} costs nothing.

\emph{Points of $S$.} By (I1), $\pkcPt_S(c,i,\kappa;\rho_{-a},\rho_{-a})\ne\emptyset$. Let
$d'=\pkcpd_{a,t,1,0}$. Steps occur only at level $0$, with $-u-a\le y<a+t-u$, hence in
$\pkcSig_n(b_i)$ by Lemma~\ref{pkc:win}(iii) ($s=a$, $t'=a+t$, $a+t+\eta+W\le n$). Lemma~\ref{pkc:free}
gives $x(i)\in\pkcPt_S(c,i,\kappa;\pkcimp(\rho_{-a},d'),d')$. At level $\ge1$ the improved orders are
$-u$; at level $0$ they are $a+t-u\ge-t-u$ if $u<-a$, and $\min(a+t-u,0)\ge-t-u$ if $u\ge-a$ (as then
$u\ge-t$). So $x(i)\in\pkcPt_S(c,i,\kappa;\rho_{-t},\pkcpd_{a,t,1,0})$.

\emph{Points of $S'$.} All jumps of $S'$ are $\le0$, so $1$ kills the cokernel of $S'$ at every letter
(Proposition~\ref{pkm:kill-small} with $\omega=0$: $e\cdot0\le\ell-2$). By (I1),
$\pkcPt_{S'}(c',i,0;\rho_0,\rho_0)\ne\emptyset$. Steps to $d'=\pkcpd_{a,t,1,0}$ are at level $0$ with
$-u\le y<a+t-u$ and $u\in[-\eta,0]$, so $0\le y<a+t+\eta\le n$ and $y\in\pkcSig_n(b_i)$. We get
$x'(i)\in\pkcPt_{S'}(c',i,0;\pkcimp(\rho_0,d'),d')$; at level $0$ the improved order is
$a+t-u\ge-u$ if $u<0$ and $\min(a+t,0)=0=-u$ if $u=0$, and at level $\ge1$ it is $-u$. So
$x'(i)\in\pkcPt_{S'}(c',i,0;\rho_0,\pkcpd_{a,t,1,0})$.

Both $x(i)$ and $x'(i)$ are standard points for $(a,t,1)$ (Remark~\ref{pkc:mono}).

\emph{Stops.} If $\varepsilon_i=1$: the jumps of $S$ at $i$ are $\le\eta\le t$, those of $S'$ are
$\ge-\eta\ge-t$. Lemma~\ref{pkm:plain} (with $l=1$, killer exponent $a$, radius $t$, $\eta_0=0$,
$x=x(i)$, $x'=x'(i)$, $x'_0=x'(i^{(-1)})$ of orders $\rho_0$) gives (St).

If $\varepsilon_i=-1$: the jumps of $S$ at $i$ are $\ge-\eta\ge-t$ and those of $S'$ are $\le0\le t$.
We construct a point $q'$ of $S'$ over $c'$ at $i$ with orders $\ge\rho_{-t}$ and depth
$\pkcpd_{a,t,1,t}$. By Lemma~\ref{pkp:tau}(v) in the form of \S\ref{pkn:ss-def},
$\varphi x'(i^{(-1)})\in\pkcPt_{S'}(c',i,0;\pkcsh_{1,0}(\rho_0),\pkcsh_{1,1}(\pkcpd_{a,t,1,0}))$. The old
depth is $1$ at level $0$, $a+t-u$ at level $1$, $-u$ above; the new depth $\pkcpd_{a,t,1,t}$ is
$a+t-u$, $t-u$, $-u$. So steps occur only at level $0$, with $1\le y<a+t-u\le a+t+\eta\le n$, in
$\pkcSig_n(b_i)$. Lemma~\ref{pkc:free} gives $q'\in\pkcPt_{S'}(c',i,0;\pkcimp(\pkcsh_{1,0}(\rho_0),
\pkcpd_{a,t,1,t}),\pkcpd_{a,t,1,t})$. The orders $\pkcsh_{1,0}(\rho_0)$ are $0$ at level $0$ and $-u$ at
level $a'\ge1$. Improved: at level $0$, $\min(a+t-u,0)=0\ge-t-u$ (as $-\eta\le u\le0$); at level
$a'\ge1$, $-u\ge0$, and the improved order is $\pkcpd_{a,t,1,t}(a',u)\ge-u$ if $u<0$ and
$\min(\pkcpd_{a,t,1,t}(a',0),0)=0$ if $u=0$. So the orders of $q'$ are $\ge\rho_{-t}$. Now
Lemma~\ref{pkm:plain} with the roles exchanged ($l=1$, killer exponent $a$, radius $t$, $\eta_0=t$,
the point of the side that is down $q'$ with depth $\pkcpd_{a,t,1,t}$, $x=x(i)$ of orders
$\rho_{-t}$ and $x_0=x(i^{(-1)})$ of orders $\rho_{-t}=\rho_{-\eta_0}$) gives that $q'$ is a stop
(for $t,1$) and the move of the first side by $\varphi x(i^{(-1)})-x(i)$ with $q'$ at rest is
$\kappa$-safe. This is (St).

\emph{The walk.} Corollary~\ref{pkm:walk-all} with $t$ there equal to $a$, radius $t$, $e=\kappa$,
$e'=0$ gives $\mathrm{SCC}_N(c,c')$ for $N\ge\kappa+0+\kappa$.
\end{proof}

\subsubsection{Generic primes: Proposition~\ref{pki:generic}}\label{pkc:ss-generic}

\begin{theorem}[Proposition~\ref{pki:generic}]\label{pkc:generic}
Let $e=1$, $\ell\ge5$, and $a,\eta,W\ge0$ with $\eta\ge1$, such that for every letter $z$:
$b_z+W+\eta\le(\ell-3)a$; $6(a+\eta+W)<b_z$; $\pkcNM(\ell,6(a+\eta+W),b_z)$; $U_z\subseteq
\pkiWin(\eta,-\varepsilon_z,b_z,W)$ and $U'_z\subseteq\pkiWin(\eta,\varepsilon_z,b_z,W)$. Let $L$ be
crystalline with jumps in $U'$.
\begin{enumerate}
\item[(i)] If no element of any $U(\gamma)\cup U'(\gamma)$ lies in $[-a,0)$, or if every
$U'(\gamma)\subseteq[-\eta,0]$, then $\plldef_{\pkiC(U)}(L)=0$.
\item[(ii)] If $L$ has at most $\mu$ jumps (with multiplicity) in $[-\eta,\eta]$ at each $\gamma$, then
$\plldef_{\pkiC(U)}(L)\le11f\bigl(\mu(a+1)+\pllrk_\auO L\cdot(W+1)\bigr)$.
\end{enumerate}
\end{theorem}

\emph{Dependencies:} (i) is proved from (I1), (I2); (ii) from (I1)--(I4).

\begin{proof}[Setting up]
As $e=1$, $E_v/\Q_\ell$ is unramified, so $j(\Gamma)$ contains the inertia group and
Lemma~\ref{pkp:chart-exists} gives an unramified Kummer chart ($E=z+\ell$, $d_v=0$) with $\Gamma'=
\Gamma$. Fix $\zeta\in\auO$ of exact order $\ell^f-1$ as in \S\ref{pki:ss-arch}; the letter of $\gamma$
is the residue of $\gamma(\zeta)$, and the letter jump sets $U_z$, $U'_z$ of \S\ref{pki:ss-arch} are the
jump sets of the universal sides $S(Y)$, $Y\in\pkiC(U)$, and $S(L)$ of Lemma~\ref{pkc:cup} (every
member of $\pkiC(U)$ is crystalline on $\Gamma=\Gamma'$ with jumps in $U$). Fix $Y\in\pkiC(U)$, a
pairing $P$ modulo some $\ell^m$ of $S=S(Y)$ and $S'=S(L)$, and put $n=6(a+\eta+W)$. Then
\eqref{pkc:eq-win} holds with $b_i=b_z$, $\varepsilon_i=\varepsilon_z$ at each letter $i=z$. Since
$\eta<n<b_z$, Lemmas~\ref{pkc:win}(i) and~\ref{pkc:dualjumps} show that every jump of $S$, of $S'$ and
of $T_P$ at the letter $z$ is $\le b_z+W+\eta\le(\ell-3)a$. Hence, by Proposition~\ref{pkm:kill-unram}
(with $s=0$, $t=a$; $\ell\ge5$), $E^a=\ell^0E^a$ kills the cokernels of $S$, $S'$ and $T_P$ at every
letter.
\end{proof}

\begin{proof}[Proof of (i)]
If no element of any $U(\gamma)\cup U'(\gamma)$ lies in $[-a,0)$, then no jump of $S$ or $S'$ at any
letter lies in $[-a,0)$ (the jump sets are $\{0\}$ and unions of the $U(\gamma)$, resp.\
$U'(\gamma)$), and Lemma~\ref{pkc:gap} with $\kappa=0$ and $d_v=0$ gives $\mathrm{SCC}_0(c,x)$ for all
$c\in\pllZf(Y)$, $x\in\pllZf(L)$. If every $U'(\gamma)\subseteq[-\eta,0]$, then every
$U'_z\subseteq[-\eta,0]$, and Lemma~\ref{pkc:nonpos} with $\kappa=0$ gives the same. As $Y$ and $P$ were
arbitrary, Lemma~\ref{pkc:cup} gives $\pllZf(L)\subseteq L^{\perp\pkiC(U)}$, i.e.\
$\plldef_{\pkiC(U)}(L)=0$.
\end{proof}

\begin{proof}[Proof of (ii)]
\emph{The submodule.} For a letter $i$ let $d_i$ be the number of embeddings $\gamma$ of letter $i$.
Embeddings of different letters are different (the Teichm\"uller lifts of different letters have
different images), so $\sum_id_i\le[\auO:\Z_\ell]$. At each $\gamma$, $L$ has at most $\mu$ jumps in
$[-\eta,\eta]$ and $\pllrk_\auO L$ jumps in all, with multiplicity; hence (I4), applied to $S'=S(L)$ with
$b_i=b_z$, $\varepsilon_i=\varepsilon_z$, $n_i=fd_i\mu$ and $r_i=fd_i\pllrk_\auO L$, gives a
$\Z_\ell$-submodule $V\subseteq\pllZf(L)$ with
\[
\pglen_{\Z_\ell}(\pllZf(L)/V)\le11\sum_ifd_i\bigl(\mu(a+1)+\pllrk L\,(W+1)\bigr)\le
11f\bigl(\mu(a+1)+\pllrk L\,(W+1)\bigr)[\auO:\Z_\ell],
\]
such that for every $x\in V$ and every letter $i$ there is $x'(i)\in\pkcPt_{S'}(x,i,0;\rho_0,\delta_i)$,
with $\delta_i$ as in (I4). We show that $\mathrm{SCC}_0(c,x)$ holds for $c\in\pllZf(Y)$, $x\in V$.
Then Lemma~\ref{pkc:cup} gives $V\subseteq\pllZf(L)\cap L^{\perp\pkiC(U)}$, so
$\pllZf(L)/(\pllZf(L)\cap L^{\perp\pkiC(U)})$ is a quotient of $\pllZf(L)/V$; its $\Z_\ell$-length is
at most the bound above, and dividing by $\pglen_{\Z_\ell}(\auO/\ell)=[\auO:\Z_\ell]$ (the
normalisation $\pglen\auO/\ell=1$) gives (ii).

\emph{Points of $S$.} Fix $c$, $x$ and a letter $i$; write $b=b_i$. By (I1) (with the killer $E^a$ of
$S$), $\pkcPt_S(c,i,0;\rho_{-a},\rho_{-a})\ne\emptyset$. Deepen to $\rho_{2a+\eta}$: the steps are at
level $0$ with $-u-a\le y<2a+\eta-u$, hence in $\pkcSig_n(b)$ by Lemma~\ref{pkc:win}(iii) ($s=a$,
$t'=2a+\eta$, $2a+2\eta+W\le n$). Lemma~\ref{pkc:free} gives $x(i)\in\pkcPt_S(c,i,0;o,\rho_{2a+\eta})$
with $o=\pkcimp(\rho_{-a},\rho_{2a+\eta})$, that is,
\[
o(0,u)=2a+\eta-u\ (u<-a),\qquad o(0,u)=\min(2a+\eta-u,0)\ (u\ge-a),\qquad o(a',u)=-u\ (a'\ge1).
\]
Two consequences. (a) $o(0,u)\ge-a-u$ for every $u$ (if $u\ge-a$ both terms of the minimum are $\ge-a-u$).
(b) $o\ge\rho_{-\eta}$ at the jumps of $S$: at level $0$ with $u<-a$, $2a+\eta-u\ge-\eta-u$; with $u\ge-a$
the jump $u$ is either low, so $u\ge-\eta$, or far, so $|u|\ge b>a$ and hence $u\ge b>0$; in both cases
$\min(2a+\eta-u,0)\ge-\eta-u$. Moreover $\rho_{2a+\eta}\ge\pkcpd_{a,\eta,1,0}$. So, by
Remark~\ref{pkc:mono}, $x(i)\in\pkcPt_S(c,i,0;\rho_{-\eta},\pkcpd_{a,\eta,1,0})$ and $x(i)$ is a standard
point for $(a,\eta,1)$.

\emph{Points of $S'$.} The point $x'(i)$ has orders $\rho_0\ge\rho_{-\eta}$ and depth
$\delta_i\ge\pkcdep_{a,\eta,1}$: if $\varepsilon_i=1$, $\rho_{2a+\eta}(0,u)=2a+\eta-u\ge a+\eta-u$; if
$\varepsilon_i=-1$, at level $0$ we have $b+W+a+\eta\ge a+\eta-u$ because every jump of $S'$ at $i$ is
$\ge-(b+W)$, at level $1$ we have $W+1-u+2a[|u|\le\eta]\ge-u$, and above both are $-u$. So $x'(i)$ is a
standard point for $(a,\eta,1)$.

\emph{Stops, $\varepsilon_i=1$.} The jumps of $S$ at $i$ are $\le\eta$, those of $S'$ are $\ge-\eta$.
Lemma~\ref{pkm:plain} ($l=1$, $t=a$, radius $\eta$, $\eta_0=0$, killer $E^a$ of $T_P$ at $i$, $x=x(i)$,
$x'=x'(i)$, $x'_0=x'(i^{(-1)})$ of orders $\rho_0$) gives (St) of Theorem~\ref{pkm:walk} at $i$.

\emph{Stops, $\varepsilon_i=-1$.} The jumps of $S'$ at $i$ are $\le\eta$, so the standard point $x'(i)$
is a stop for $(\eta,1)$ by Lemma~\ref{pkm:down}. It remains to show that the move of the first side by
$w=\varphi x(i^{(-1)})-x(i)$ with $x'(i)$ at rest is $0$-safe; we use (I3) with $l=1$, $\kappa=0$.

The dual point: $E^a$ kills the cokernel of $T_P$ at $j=i^{(-1)}$, so (I1) gives a point of $T_P$ over $1$
at $j$ with orders and depth $\rho_{-a}$; deepening it to $\rho_a$ involves only level-$0$ steps
$-a-s\le y<a-s$ at the jumps $s$ of $T_P$ at $j$, which lie in $\pkcSig_n(b_j)$ by
Lemma~\ref{pkc:dualjumps}. Lemma~\ref{pkc:free} gives a point $\nu$ of $T_P$ over $1$ at $j$ with orders
$o_\nu=\pkcimp(\rho_{-a},\rho_a)$:
\[
o_\nu(0,s)=a-s\ (s<-a),\qquad o_\nu(0,s)=\min(a-s,0)\ (s\ge-a),\qquad o_\nu(a',s)=-s\ (a'\ge1).
\]
Each $o_\nu(a',\cdot)$ is nonincreasing (at level $0$: both branches are, and at the junction
$o_\nu(0,-a-1)=2a+1\ge0\ge o_\nu(0,-a)$). Note $o_\nu(0,s)\ge-a-s$ and $o_\nu(0,s)\ge\min(a-s,-(W+1)-s)$
when $s\ge-W-1$ (if $s\ge-a$, $\min(a-s,0)\ge-(W+1)-s$).

The move: $\varphi x(i^{(-1)})\in L_i(\pkcsh_{1,0}(o))$ lies over $c$ at $i$ (Lemma~\ref{pkp:tau}(v)), and
so does $x(i)\in L_i(o)$. Hence $w$ lies over $0$ and $w\in L_i(o_w)$ with
$o_w=\min(\pkcsh_{1,0}(o),o)$: $o_w(0,u)=\min(0,o(0,u))$, $o_w(1,u)=\min(o(0,u),-u)$, $o_w(a',u)=-u$ for
$a'\ge2$. The point at rest is $y=x'(i)$, with depth $d_y=\pkcnw_i$.

We check the inequality of (I3) at $a'\ge0$, $u\in U_S(i^{(-a')})$, $u'\in U_{S'}(i^{(-a')})$, $s=-u-u'-1$.
At level $a'$ the windows are those of the letter $i^{(-a')}$, with radius $\eta$, width $W$, height
$\beta=b_{i^{(-a')}}>n$ and some sign $\varepsilon'$ (opposite for $S$ and $S'$).

\emph{Level $0$} (letter $i$, $\varepsilon_i=-1$, so $U_S(i)\subseteq\pkiWin(\eta,1,b,W)$). The sum is
$0+\min(0,o(0,u))+b+W+a+\eta$. If $u$ is low, then $o(0,u)\ge0$ ($2a+\eta-u\ge2a\ge0$), and the sum is
$b+W+a+\eta>0$. If $u\in[b,b+W]$, then $u>2a+\eta$ and $o(0,u)=2a+\eta-u$, so the sum is
$\ge 2a+\eta-(b+W)+b+W+a+\eta=3a+2\eta>0$.

\emph{Level $1$.} The sum is $o_\nu(0,s)+\min(o(0,u),-u)+W+1-u'+2a[|u'|\le\eta]$, and
$\min(o(0,u),-u)\ge-a-u$ by (a).
(A) $u'$ low: with $o_\nu(0,s)\ge-a-s=-a+u+u'+1$ the sum is $\ge(-a+u+u'+1)+(-a-u)+(W+1-u'+2a)=W+2$.
(B) $u'$ far, $u$ low: then $|u'|\ge\beta$ and $|u|\le\eta$, so either $s\le\eta-\beta-1<-a$ or
$s\ge\beta-\eta-1>a$ (as $\beta>n\ge a+\eta+1$); in both cases $o_\nu(0,s)=a-s$. The sum is
$\ge(a+u+u'+1)+(-a-u)+(W+1-u')=W+2$.
(C) $u'$ and $u$ far: they lie in far windows of opposite signs, so $|u+u'|\le W$ and $s\ge-W-1$, whence
$o_\nu(0,s)\ge-(W+1)-s=u+u'-W$. Also $|u|\ge\beta>2a+\eta$, so $o(0,u)=2a+\eta-u>-u$ (if $u<-a$ this
is the formula for $o$; if $u\ge-a$, then $u\ge\beta>2a+\eta$ and $o(0,u)=\min(2a+\eta-u,0)=2a+\eta-u$), and $\min(o(0,u),-u)=-u$. The sum is $\ge u+u'-W-u+W+1-u'=1$.
In all cases the sum is positive.

\emph{Level $a'\ge2$.} $\pkcsh_{1,0}(o_\nu)(a',s)=o_\nu(a'-1,s)=-s$, $o_w(a',u)=-u$, $d_y(a',u')=-u'$;
the sum is $u+u'+1-u-u'=1>0$.

By (I3), the move is $0$-safe, which is (St) at $i$.

\emph{The walk.} At every letter, $E^a$ kills the cokernel of $T_P$, $x(i)$ and $x'(i)$ are standard
points for $(a,\eta,1)$ over $\ell^0c$, $\ell^0x$, and (St) holds with $l(i)=1$.
Corollary~\ref{pkm:walk-all} with $\kappa=e=e'=0$ gives $\mathrm{SCC}_0(c,x)$.
\end{proof}

\subsubsection{Tame types: Proposition~\ref{pki:tame}}\label{pkc:ss-tame}

Let $E_v=\Q_\ell$, so that $j\colon\Gamma\to G_{\Q_\ell}$ is bijective. Let $r\ge1$ and
$e_0=\operatorname{lcm}_{1\le i\le r}(\ell^i-1)$; note $\ell\nmid e_0$. For a potentially crystalline
lattice $Y$ and $s\in\Gamma$ in the inertia group, $s$ acts $\auO$-linearly on $D_{\mathrm{pcris}}(Y)$
and ``$P$ kills $s$ on $D_{\mathrm{pcris}}(Y)$'' is as in Definition~\ref{pkw:def-pcris}. The case
uses the following inputs, proved in \S\ref{pkw:sec}.

\begin{enumerate}
\item[(T1)] \emph{Tame chart.} If $r<\ell$, there is a tame chart (Definition~\ref{pkw:def-tame}) with
ramification index $e_0$, subgroup $\Gamma'$ and tame generator $s_0$ (Proposition~\ref{pkw:T1}).
\item[(T2)] \emph{Descent.} Every potentially crystalline lattice $Y$ such that $X^{e_0}-1$ kills every
inertial $s$ on $D_{\mathrm{pcris}}(Y)$ is crystalline on $\Gamma'$, with the same Hodge--Tate jumps
(Theorem~\ref{pkw:T2}).
\item[(T3)] \emph{Tame killer.} Let $S$ be a side over the chart of (T1) whose jumps at all letters are
$\le n_0$, and let a monic $q\in\Z_\ell[X]$ kill $s_0$ on $D_{\mathrm{pcris}}$ of the lattice of $S$. If
$e_0\le\ell^{m_0}$, $m_0(\deg q+1)\le\kappa$ and $n_0+\kappa+2<\ell$, then $\ell^\kappa=\ell^\kappa E^0$
kills the cokernel of $S$ at every letter (Theorem~\ref{pkw:T3}).
\item[(T4)] \emph{Tate duals of tensor products.} If monic $q,q'\in\Z_\ell[X]$ kill $s_0^{-1}$ on
$D_{\mathrm{pcris}}(Y)$ and on $D_{\mathrm{pcris}}(L)$ respectively, then some monic $q''\in\Z_\ell[X]$
of degree $\le\deg q\cdot\deg q'$ kills $s_0$ on
$D_{\mathrm{pcris}}(\operatorname{Hom}(Y\otimes L,\auO(1)))$ (Theorem~\ref{pkw:T4}).
\end{enumerate}

In (T3) we use the letters with the root $\zeta=1$: all embeddings have letter $1$, the letter-$1$ part
is everything, and the other letter parts are $0$ (Definition~\ref{pkp:def-letters} uses $\zeta$ only
through $\zeta^{\ell^f-1}=1$; Lemma~\ref{pkw:letters}). This avoids asking that $\auO$ contain roots of
unity of order $\ell^f-1$ for the residue degree $f$ of the tame chart.

\begin{lemma}\label{pkc:lcm}
$e_0\le\ell^{r^2}$.
\end{lemma}

\begin{proof}
$e_0$ divides $\prod_{i=1}^r(\ell^i-1)<\prod_{i=1}^r\ell^i=\ell^{r(r+1)/2}\le\ell^{r^2}$.
\end{proof}

\begin{theorem}[Proposition~\ref{pki:tame}]\label{pkc:tame}
Let $E_v=\Q_\ell$, $r,N\ge1$, $e_0$ as above. For each inertial $s\in\Gamma$ let $K(s)$ be a set of
polynomials over $\auO$ containing $X^{e_0}-1$ and the image of a monic polynomial over $\Z_\ell$ of
degree $\le N$; let $\pkiC(U,K)$ be the class of potentially crystalline lattices with jumps in $U$ on
whose $D_{\mathrm{pcris}}$ each $P\in K(s)$ kills $s$, for every inertial $s$. Let $L\in\pkiC(U',K^{-1})$,
$K^{-1}(s)=K(s^{-1})$. Let $\eta,W,b\ge0$ and a sign $\varepsilon$ with
$U(\gamma)\subseteq\pkiWin(\eta,-\varepsilon,b,W)$ and $U'(\gamma)\subseteq\pkiWin(\eta,\varepsilon,b,W)$
for all $\gamma$, and $\ell>b+6(W+\eta)+r^2(N+1)^2+2$. Then $\pkiC(U,K)$ has cup exponent
$3r^2(N+1)^2+150(\eta+W+1)$ on $L$.
\end{theorem}

\emph{Dependencies:} proved from (I1), (I2) and (T1)--(T4); these are proved in \S\S\ref{pkt:sec},
\ref{pke:sec}, \ref{pkw:sec}.

\begin{proof}
Put $\kappa=r^2(N+1)^2$, $n=6(\eta+W)$, $n_0=b+n$ and $H=3\kappa+150(\eta+W+1)$. Since
$\ell>\kappa+2\ge r+2$, (T1) gives a chart with ramification index $e_0$, subgroup $\Gamma'$ and tame
generator $s_0$; as $\ell\nmid e_0$, $d_v=0$ if $e_0=1$ and $d_v=v_\ell(e_0)+1=1$ otherwise, so $d_v\le1$.
Use $\zeta=1$: every letter jump set is $\{0\}$ or $\{0\}\cup\bigcup_\gamma U(\gamma)$ (resp.\ with
$U'$), hence lies in $\pkiWin(\eta,-\varepsilon,b,W)$ (resp.\ $\pkiWin(\eta,\varepsilon,b,W)$), as $0$ is
in the low window.

\emph{Descent.} Each $Y\in\pkiC(U,K)$ is killed at every inertial $s$ by $X^{e_0}-1\in K(s)$, and $L$ by
$X^{e_0}-1\in K(s^{-1})=K^{-1}(s)$; by (T2) all of them are crystalline on $\Gamma'$, with jumps in $U$,
resp.\ $U'$. So Lemma~\ref{pkc:cup} applies to the class $\pkiC(U,K)$ and $L$ over this chart, and it
suffices to prove $\mathrm{SCC}_H(c,x)$ for every $Y\in\pkiC(U,K)$, every pairing $P$ modulo some
$\ell^m$ of $S=S(Y)$ and $S'=S(L)$, $c\in\pllZf(Y)$ and $x\in\pllZf(L)$.

\emph{Killers.} By hypothesis there are monic $q_0\in\Z_\ell[X]$ with image in $K(s_0)$ and $q'_0$ with
image in $K(s_0^{-1})=K^{-1}(s_0)$, of degrees $\le N$; so $q_0$ kills $s_0$ on $D_{\mathrm{pcris}}(Y)$
and $q'_0$ kills $s_0$ on $D_{\mathrm{pcris}}(L)$; and also $q'_0$ kills $s_0^{-1}$ on
$D_{\mathrm{pcris}}(Y)$ (as $Y\in\pkiC(U,K)$) and $q_0$ kills $s_0^{-1}$ on $D_{\mathrm{pcris}}(L)$ (as
$K^{-1}(s_0^{-1})=K(s_0)$). By (T4), applied with $q=q'_0$ and $q'=q_0$, a monic
$q''$ of degree $\le N^2$ kills $s_0$ on $D_{\mathrm{pcris}}$ of $\operatorname{Hom}(Y\otimes L,\auO(1))$,
the lattice of $T_P$. All jumps of $S$, $S'$ are $\le b+W+\eta\le n_0$ (Lemma~\ref{pkc:win}(i)), and all
jumps of $T_P$ (namely $0$ and the $-u-u'-1$) are $\le b+W+2\eta\le n_0$ (Lemma~\ref{pkc:win}(ii), first
assertion, which does not need $\eta\le b$). Take $m_0=r^2$: $e_0\le\ell^{r^2}$ by Lemma~\ref{pkc:lcm};
$r^2(\deg q_0+1)\le r^2(N+1)\le\kappa$, the same for $q'_0$, and $r^2(\deg q''+1)\le r^2(N^2+1)\le\kappa$;
and $n_0+\kappa+2<\ell$ is the hypothesis. By (T3), $\ell^\kappa=\ell^\kappa E^0$ kills the cokernels of
$S$, $S'$ and $T_P$ at every letter.

\emph{The walk.} Apply Lemma~\ref{pkc:gap} with $a=0$ (so the gap condition, on $[-0,0)=\emptyset$, is
empty), $b_i=b$, $\varepsilon_i=\varepsilon$ at every letter, and $n=6(0+\eta+W)$. The condition
$\pkcNM(\ell,n,b)$ holds because every element of $\pkcSig_n(b)$ has absolute value $\le b+n<\ell$, while
a nonzero multiple of $\ell$ has absolute value $\ge\ell$. Lemma~\ref{pkc:gap} gives
$\mathrm{SCC}_{N'}(c,x)$ for every $N'\ge3\kappa+6(2n+1)d_v$, in particular for $N'=H$, since
\[
3\kappa+6(2n+1)d_v\le3\kappa+72(\eta+W)+6\le3\kappa+150(\eta+W+1)=H .
\]
By Lemma~\ref{pkc:cup}, $\ell^Hx\in L^{\perp\pkiC(U,K)}$ for every $x\in\pllZf(L)$.
\end{proof}

\begin{remark}[On (T4)]\label{pkc:rem-T4}
(T4) is Theorem~\ref{pkw:T4}: a Cayley--Hamilton polynomial on
$\Z_\ell[x]/(q)\otimes\Z_\ell[x']/(q')$ kills $s_0^{-1}$ on the tensor product, and passing to the Tate
dual exchanges $s_0^{-1}$ and $s_0$ (Lemma~\ref{pkw:dual}).
\end{remark}

\subsubsection{Sparse $b$-letters: Proposition~\ref{pki:sparse}}\label{pkc:ss-sparse}

Proposition~\ref{pki:sparse} is proved in \S\ref{pkd:ss-sparse} (Theorem~\ref{pkd:sparse}), from
(I1), (I2) and, when the sign is not constant, the input (I5) of \S\ref{pkd:ss-tools}.

\subsubsection{Any pattern of $b$-letters: Proposition~\ref{pki:any}}\label{pkc:ss-any}

Proposition~\ref{pki:any} is proved in \S\ref{pkd:ss-any} (Theorem~\ref{pkd:any}), from (I1), (I2),
(I5).




\providecommand{\auO}{\mathcal O}
\providecommand{\panot}[1]{\par\smallskip{\raggedright\noindent\textbf{[Not proved here:} #1\textbf{]}\par}\smallskip}
\providecommand{\pllZf}{Z_f}
\providecommand{\plldef}{\delta^{\vee}}
\providecommand{\pglen}{\operatorname{len}}
\providecommand{\pkiWin}{\operatorname{Win}}
\providecommand{\pkiC}{\mathcal C}
\providecommand{\pkilog}{\operatorname{lg}}
\providecommand{\pkcPt}{\mathrm{Pt}}
\providecommand{\pkcimp}{\mathrm{imp}}
\providecommand{\pkcpd}{\mathrm{pd}}
\providecommand{\pkcdep}{\mathrm{dep}}
\providecommand{\pkcsh}{\mathrm{sh}}
\providecommand{\pkdco}{\mathrm{co}}
\providecommand{\pkdlg}{\mathrm{gd}}
\providecommand{\pkdstop}{\mathrm{stop}}
\providecommand{\pkdstd}{\mathrm{std}}
\providecommand{\pkdcost}{\mathrm{cost}}

\subsection{The cases with cuts: sparse and arbitrary patterns of $b$-letters}\label{pkd:sec}

This section proves Proposition~\ref{pki:sparse} (\S\ref{pkd:ss-sparse}) and
Proposition~\ref{pki:any} (\S\ref{pkd:ss-any}), the two cases not treated in \S\ref{pkc:sec}, and
the remaining assertions of Proposition~\ref{pki:data} (\S\ref{pkd:ss-data}). Both cases are instances of
one scheme (Theorem~\ref{pkd:scheme} in \S\ref{pkd:ss-scheme}): at a chosen set of letters, the
\emph{cuts}, both sides receive standard points (from a cokernel killer and deepening), and at every
cut the walk of Theorem~\ref{pkm:walk} needs a stop, which is either the standard point of the
side that is ``down'' (a \emph{plain} cut, Lemma~\ref{pkm:plain}) or a point produced by a
\emph{long gap} (an expansion in powers of $E$; input (I5) below) together with the rule
(Proposition~\ref{pkp:rule}). We use the conventions and notation of \S\ref{pkc:sec}
(points $\pkcPt_S(c,i,e;o,d)$, the orders $\rho_t$, $\pkcdep$, $\pkcpd$, $\pkcimp$, $\pkcsh$,
$\mathrm{std}_l$, $\mathrm{stop}_{\eta,l}$, $\mathrm{rule}_{l,t}$, the cost
$\mathrm{cost}(a,n)=v_\ell(n)+d_v+3a$ with $v_\ell(0)=0$), and cite the inputs (I1), (I2) of
\S\ref{pkc:ss-inputs} as stated there.

Propositions~\ref{pki:sparse} and~\ref{pki:any} are proved below as Theorems~\ref{pkd:sparse}
and~\ref{pkd:any} from \S\S\ref{pkp:sec}--\ref{pkn:sec} and the inputs (I1), (I2), (I5)
((I5) is not needed when the sign is constant); Proposition~\ref{pki:data} is Theorem~\ref{pkd:data},
whose tame membership $L_c^*\in\pkiC(U',K^{-1})$ uses the input (I6). The inputs are proved elsewhere: (I1)
(Proposition~\ref{pkt:I1}), (I2) (Theorem~\ref{pke:I2}), (I5) (Theorem~\ref{pkg:I5}), (I6)
(Lemma~\ref{pkw:dual}); Proposition~\ref{pki:chart} is Theorem~\ref{pkj:chart}.

\subsubsection{Tools}\label{pkd:ss-tools}

\begin{remark}[Scaling]\label{pkd:scale}
If $x\in\pkcPt_S(c,i,e;o,d)$ and $k\ge0$, then $\ell^kx\in\pkcPt_S(c,i,e+k;o,d)$: $\mathcal N(X_S)$
and the modules $L_i(o)$, $L_i(d)$ are modules over $\Z_\ell$, $\mathrm{pr}$ is linear, and
$g(\ell^kx)-\ell^kx=\ell^k(gx-x)$. So a point over $\ell^ec$ gives points over $\ell^{e'}c$ for all
$e'\ge e$, with the same orders and depth.
\end{remark}

\begin{lemma}[Valuations in an interval]\label{pkd:val}
Let $c\in\Z$, $r\ge0$ and $B\ge1$ with $|c|+r\le B$. Then
$\sum_{|y-c|\le r}v_\ell(y)\le(2r+1)/(\ell-1)+\pkilog_\ell B$.
\end{lemma}

\begin{proof}
For $k\ge1$ let $n_k$ be the number of $y\neq0$ with $|y-c|\le r$ and $\ell^k\mid y$; then
$\sum_{|y-c|\le r}v_\ell(y)=\sum_{k\ge1}n_k$ (the term $y=0$ contributes $0$). A nonzero $y$ in the
interval has $|y|\le B$, so $n_k=0$ unless $\ell^k\le B$, i.e.\ unless $k\le\pkilog_\ell B$. Among
$2r+1$ consecutive integers at most $(2r+1)/\ell^k+1$ are divisible by $\ell^k$. Hence the sum is at
most $\sum_{k=1}^{\pkilog_\ell B}\bigl((2r+1)\ell^{-k}+1\bigr)\le(2r+1)/(\ell-1)+\pkilog_\ell B$.
\end{proof}

\paragraph*{The setting of this section}

Throughout, $P$ pairs sides $S$, $S'$ over a Kummer chart, $\eta\ge1$, $V\ge0$, and at every letter
$i$ an integer $\beta_i\ge0$ and a sign $\varepsilon_i$ are given with
\begin{equation}\label{pkd:eq-win}
U_S(i)\subseteq\pkiWin(\eta,-\varepsilon_i,\beta_i,V),\qquad
U_{S'}(i)\subseteq\pkiWin(\eta,\varepsilon_i,\beta_i,V).
\end{equation}
We say that a side \emph{has the negative window at $i$} if it is $S$ and $\varepsilon_i=1$, or $S'$
and $\varepsilon_i=-1$; then all its jumps at $i$ are $\le\eta$ and all jumps of the other side at
$i$ are $\ge-\eta$. A letter $i$ is \emph{small} if $\beta_i+V\le\eta$; then all jumps of $S$ and $S'$
at $i$ lie in $[-\eta,\eta]$.

\begin{lemma}[Jump bounds]\label{pkd:jumps}
At every letter $i$, every jump of $S$ and of $S'$ is $\le\beta_i+V+\eta$ and $\ge-(\beta_i+V+\eta)$,
and every jump of $T_P$ is $\le\beta_i+V+2\eta$. A jump $u$ of $S$ or $S'$ at $i$ with
$u>-\beta_i$ satisfies $u\ge-\eta$.
\end{lemma}

\begin{proof}
The first assertion is clear from the windows. The jumps of $T_P$ at $i$ are $0$ and the
$s=-u-u'-1$ with $u\in U_S(i)$, $u'\in U_{S'}(i)$ (Definition~\ref{pkp:def-pairing}). If both are low,
$s\le2\eta-1$. If exactly one is far, say $u=\delta x$ with $x\in[\beta_i,\beta_i+V]$ and $|u'|\le\eta$,
then $s\le\beta_i+V+\eta-1$ (if $\delta=-1$) or $s\le\eta-1$ (if $\delta=1$). If both are far they lie in
windows of opposite signs, so $|u+u'|\le V$ and $s\le V-1$. For the last assertion: a jump in a far
window with $u>-\beta_i$ is in the positive far window, so $u\ge\beta_i\ge0$; otherwise $|u|\le\eta$.
\end{proof}

\paragraph*{Cut orders and the long-gap depth}

For $s,t,A\ge0$ and $l\ge1$ with $t<l$ define
\[
\pkdco_{s,l}=\pkcimp\bigl(\rho_{-s},\pkcdep_{s,\eta,l+1}\bigr),\qquad
\pkdlg_{l,t,A}(a,u)=\begin{cases}A&a=t,\\ 1&a<l,\ a\ne t,\\ -u&a\ge l.\end{cases}
\]
Explicitly, $\pkcdep_{s,\eta,l+1}(a,u)=-u+s[a=0]+\eta[a\le l]$, and $\pkdco_{s,l}(a,u)$ is
$\pkcdep_{s,\eta,l+1}(a,u)$ where $\rho_{-s}(a,u)>0$ and $\min(\pkcdep_{s,\eta,l+1}(a,u),0)$ elsewhere.

\begin{lemma}[Cut orders]\label{pkd:co}
Let $s\ge0$, $l\ge1$. For all $a\ge0$ and $u\in\Z$:
\begin{enumerate}
\item[(i)] if $a\le l$, then $\pkdco_{s,l}(a,u)\ge\min(\eta-u,0)$;
\item[(ii)] if $a\ge1$, or $a=0$ and $s=0$, then $\pkdco_{s,l}(a,u)\ge-u$;
\item[(iii)] if $u\ge-s$ implies $u\ge-\eta$, then $\pkdco_{s,l}(0,u)\ge-\eta-u$.
\end{enumerate}
In particular, if no jump of a side $R$ at $i$ lies in $[-s,-\eta)$, then $\pkdco_{s,l}\ge\rho_{-\eta}$
at the jumps of $R$ (at level $a$, the jumps at $i^{(-a)}$).
\end{lemma}

\begin{proof}
Write $D=\pkcdep_{s,\eta,l+1}$. (i) At level $0$, $D(0,u)=s+\eta-u\ge\eta-u$; at level $1\le a\le l$,
$D(a,u)=\eta-u$. In both cases $\pkdco_{s,l}(a,u)\in\{D(a,u),\min(D(a,u),0)\}$ is
$\ge\min(\eta-u,0)$. (ii) If $a\ge1$, then $\rho_{-s}(a,u)=-u$ and $D(a,u)\ge-u$; if $-u>0$ the value is
$D(a,u)\ge-u$, otherwise it is $\min(D(a,u),0)\ge-u$ since both terms are $\ge-u$. If $a=0$ and
$s=0$ the same argument applies with $\rho_0(0,u)=-u$. (iii) If $u<-s$, then $\rho_{-s}(0,u)=-s-u>0$
and the value is $s+\eta-u\ge-\eta-u$. If $u\ge-s$, then $u\ge-\eta$, and $\min(s+\eta-u,0)\ge-\eta-u$.
The last assertion follows from (ii) and (iii), as $\rho_{-\eta}(a,u)=-u$ for $a\ge1$.
\end{proof}

\paragraph*{The input (I5)}

The long-gap points are produced by the following statement about the modules $\mathcal N$, which is
not among the statements of \S\S\ref{pkp:sec}--\ref{pkn:sec}; it is proved as Theorem~\ref{pkg:I5}.
For $l,A\ge0$ and $N\in\Z$ put
\[
o^{\mathrm{gap}}_{l,A,N}(a,u)=\mathrm{std}_l(a,u)+A[a=0][u<-N],\qquad
d^{\mathrm{gap}}_{l,A}(a,u)=\mathrm{std}_l(a,u)+A[a=0].
\]

\begin{enumerate}
\item[(I5)] \emph{Long gap.} Let the chart be unramified, $l\ge1$ with $\ell^l\ge3$, and
$A,\delta,\kappa,s,g\ge0$ with $A\le\ell^l\delta$. Let $\ell^\kappa$ kill the cokernel of a side $S$
at $i^{(-l)}$, and $E^s$ kill it at $i$. Let every jump of $S$ at $i^{(-a)}$, $0\le a<l$, be $\le g$,
put $N=s+g(1+\ell+\dots+\ell^{l-1})$, and assume $N<A$. Then for every $c\in C_S$,
$\pkcPt_S(c,i,\kappa+4\delta;o^{\mathrm{gap}}_{l,A,N},d^{\mathrm{gap}}_{l,A})\ne\emptyset$.
\end{enumerate}

\subsubsection{The scheme of cuts}\label{pkd:ss-scheme}

\begin{definition}[Cuts, plain letters, long gaps]\label{pkd:def-cut}
Let $\mathrm{cut}$ be a set of letters such that every letter $i$ has some $a>0$ with
$i^{(-a)}\in\mathrm{cut}$, and let $l(i)$ be the least such $a$. Let shifts $s_i\ge0$ ($i\in\mathrm{cut}$)
be given with $s_i=0$ or $s_i<\beta_i$. A letter $i$ (with $l=l(i)$) is
\begin{enumerate}
\item[(a)] \emph{plain} if for some sign $\varrho$ every letter $i^{(-a)}$, $0\le a<l$, has
$\varepsilon_{i^{(-a)}}=\varrho$ or is small;
\item[(b)] a \emph{long gap with pivot $t$}, for data $(J',s_0,\omega,A,\delta)$ of nonnegative
integers, if: the chart is unramified and $\ell\ge5$; $t+J'+1=l$; $s_{i^{(-l)}}=0$; for $a<l$,
$a\ne t$, either $\varepsilon_{i^{(-a)}}\ne\varepsilon_{i^{(-t)}}$ or $i^{(-a)}$ is small; $\eta\le\omega$
and $\beta_{i^{(-a)}}+V\le\omega$ for $t<a<l$; $N:=s_0+\omega(1+\ell+\dots+\ell^{J'})<\beta_{i^{(-t)}}$;
$E^{s_0}$ kills the cokernels of $S$ and $S'$ at $i^{(-t)}$; and
$\beta_{i^{(-t)}}+V\le A\le\ell^{J'+1}\delta$.
\end{enumerate}
\end{definition}

For a cut $i$ and a side $R\in\{S,S'\}$ let $\mathcal I_i(R)$ be the set of pairs $(a,y)$ with $0\le
a\le l(i)$ and $\rho_{-s_i}(a,u)\le y<\pkcdep_{s_i,\eta,l(i)+1}(a,u)$ for some $u\in U_R(i^{(-a)})$.

\begin{lemma}[Steps at a cut]\label{pkd:steps}
Let $i$ be a cut, $l=l(i)$, $R\in\{S,S'\}$, and let $B\ge1$ with $B\ge\beta_j+s_i+V+2\eta$ for all
letters $j$. Put $r_0=s_i+V+2\eta$ and $r_a=V+2\eta$ for $a\ge1$. Then
\[
\sum_{(a,y)\in\mathcal I_i(R)}\pkdcost(a,y)\le\sum_{a=0}^{l}\Bigl(\frac{2(2r_a+1)}{\ell-1}
+2\pkilog_\ell B+2(2r_a+1)(d_v+3a)\Bigr).
\]
\end{lemma}

\begin{proof}
Fix $a\le l$, let $j=i^{(-a)}$ and let $\sigma$ be the sign of the far window of $R$ at $j$. At level
$0$ the condition is $-u-s_i\le y<-u+s_i+\eta$; at level $1\le a\le l$ it is $-u\le y<-u+\eta$. If
$|u|\le\eta$ this gives $|y|\le r_a$. If $u=\sigma x$ with $x\in[\beta_j,\beta_j+V]$, it gives
$|y+u|\le s_i[a=0]+\eta$ and so $|y-c_a|\le r_a$ with $c_a=-\sigma\beta_j$. Hence the pairs at level
$a$ have $y$ in the union of the intervals $|y|\le r_a$ and $|y-c_a|\le r_a$, each of $2r_a+1$
integers and contained in $[-B,B]$. Lemma~\ref{pkd:val} bounds the sum of $v_\ell(y)$ over each
interval by $(2r_a+1)/(\ell-1)+\pkilog_\ell B$, and the remaining part $d_v+3a$ of the cost is summed
over at most $2(2r_a+1)$ values.
\end{proof}

\begin{theorem}[The scheme]\label{pkd:scheme}
In the setting \eqref{pkd:eq-win}, let $\mathrm{cut}$, $l$ and the shifts be as in
Definition~\ref{pkd:def-cut}, and let $\kappa,e^*\ge0$. Assume:
\begin{enumerate}
\item[(K)] for every cut $i$, $\ell^\kappa E^{s_i}$ kills the cokernels of $S$, $S'$ and $T_P$ at $i$;
\item[(C)] for every cut $i$ and $R\in\{S,S'\}$, $\kappa+\sum_{(a,y)\in\mathcal I_i(R)}\pkdcost(a,y)\le e^*$;
\item[(G)] every cut $i$ is plain, or is a long gap with some pivot $t$ and data
$(J',s_0,\omega,A,\delta)$ such that $\kappa+4\delta+3\sum_{t<a<l(i)}a\le e^*$.
\end{enumerate}
Then $\mathrm{SCC}_N(c,c')$ for all $c\in C_S$, $c'\in C_{S'}$ and $N\ge2e^*+\kappa$.
\end{theorem}

\begin{proof}
Fix $c$, $c'$. We verify the hypotheses of Theorem~\ref{pkm:walk} with $t(i)=s_i$ and
$e=e'=e^*$.

\emph{Standard points.} Let $i$ be a cut, $l=l(i)$, $R\in\{S,S'\}$ and $d\in\{c,c'\}$ the
corresponding element. By (K) and (I1), $\pkcPt_R(d,i,\kappa;\rho_{-s_i},\rho_{-s_i})\ne\emptyset$. By
(I2) with $I=\mathcal I_i(R)$ (which contains all steps from $\rho_{-s_i}$ to
$\pkcdep_{s_i,\eta,l+1}$: at levels $a>l$ the two orders agree) and (C), and by
Remark~\ref{pkd:scale}, there is a point
\[
x_R(i)\in\pkcPt_R\bigl(d,i,e^*;\pkdco_{s_i,l},\pkcdep_{s_i,\eta,l+1}\bigr).
\]
No jump $u$ of $R$ at $i$ lies in $[-s_i,-\eta)$: if $s_i=0$ this is empty, and if $s_i<\beta_i$,
then $u\ge-s_i$ gives $u>-\beta_i$ and so $u\ge-\eta$ by Lemma~\ref{pkd:jumps}. By
Lemma~\ref{pkd:co} and Remark~\ref{pkc:mono}, $x_R(i)$ has orders $\rho_{-\eta}$. As
$\pkcdep_{s_i,\eta,l+1}=\pkcpd_{s_i,\eta,l,\eta}\ge\pkcdep_{s_i,\eta,l}$, $x_R(i)$ is a standard point
(for $s_i,\eta,l$) and has the depth $\pkcpd_{s_i,\eta,l,\eta}$. Write $x(i)=x_S(i)$, $x'(i)=x_{S'}(i)$.
Then (P) of Theorem~\ref{pkm:walk} holds, and (K) there is part of our (K).

\emph{Plain letters.} Let $i$ be a cut that is plain, with sign $\varrho$, $l=l(i)$, and let $R$ be $S$ if $\varrho=1$
and $S'$ if $\varrho=-1$, $R'$ the other side. For $a<l$ the letter $i^{(-a)}$ either has
$\varepsilon=\varrho$, so $R$ has the negative window there, or is small; in both cases the jumps of $R$
at $i^{(-a)}$ are $\le\eta$ and those of $R'$ are $\ge-\eta$. Lemma~\ref{pkm:plain} (with $l$,
$t=s_i$, $\eta$, $\eta_0=\eta$, the killer $\ell^\kappa E^{s_i}$ of $T_P$ at $i$, the point $x_R(i)$ of
orders $\rho_{-\eta}$ and depth $\pkcpd_{s_i,\eta,l,\eta}$, the point $x_{R'}(i)$ of orders $\rho_{-\eta}$
and the point $x_{R'}(i^{(-l)})$ at $i^{(-l)}$ of orders $\rho_{-\eta}=\rho_{-\eta_0}$; for $R=S'$ in
the form with the roles exchanged) shows that $x_R(i)$ is a stop (for $\eta,l$) over $\ell^{e^*}d$ and
that the move of $R'$ by $\varphi^lx_{R'}(i^{(-l)})-x_{R'}(i)$ with $x_R(i)$ at rest is $\kappa$-safe.
This is (St) at $i$.

\emph{Long gaps.} Let $i$ be a cut that is a long gap, with pivot $t$ and data $(J',s_0,\omega,A,\delta)$, $l=l(i)$,
$i_t=i^{(-t)}$. Let $R$ be the side with the negative window at $i_t$, $R'$ the other side, and $d$, $d'$
the corresponding elements of $\{c,c'\}$. For $a<l$, $a\ne t$, either $\varepsilon_{i^{(-a)}}\ne
\varepsilon_{i_t}$, so that $R'$ has the negative window at $i^{(-a)}$, or $i^{(-a)}$ is small; hence
\begin{equation}\label{pkd:eq-updown}
\text{for }a<l,\ a\ne t:\quad U_R(i^{(-a)})\subseteq[-\eta,\infty),\qquad
U_{R'}(i^{(-a)})\subseteq(-\infty,\eta].
\end{equation}

\emph{(1) A point at the pivot.} We apply (I5) to $R$ at the letter $i_t$, with $l$ there equal to
$J'+1$, $s=s_0$, $g=\omega$, and $A,\delta,\kappa$. Note $i_t^{(-(J'+1))}=i^{(-l)}$ is a cut with shift
$0$, so $\ell^\kappa=\ell^\kappa E^0$ kills the cokernel of $R$ there by (K); $E^{s_0}$ kills it at
$i_t$; $\ell^{J'+1}\ge\ell\ge5$; $A\le\ell^{J'+1}\delta$. The jumps of $R$ at $i_t$ are $\le\eta\le\omega$ (the
negative window), and at $i_t^{(-a)}=i^{(-(t+a))}$ with $1\le a\le J'$ they are $\le\max(\eta,
\beta_{i^{(-(t+a))}}+V)\le\omega$, since $t<t+a<l$. Finally $N<\beta_{i_t}\le A$. So there is
$x_{i_t}\in\pkcPt_R(d,i_t,\kappa+4\delta;o^{\mathrm{gap}}_{J'+1,A,N},d^{\mathrm{gap}}_{J'+1,A})$.

\emph{(2) Transport and deepening.} By Lemma~\ref{pkp:tau}(v), $\varphi^tx_{i_t}$ is a point of $R$ over
$\ell^{\kappa+4\delta}d$ at $i$ with orders $o_1=\pkcsh_{t,0}(o^{\mathrm{gap}}_{J'+1,A,N})$ and depth
$d_1=\pkcsh_{t,1}(d^{\mathrm{gap}}_{J'+1,A})$. Explicitly, since $\mathrm{std}_{J'+1}(b,u)=0$ for
$b\le J'$ and $=-u$ for $b\ge J'+1$, and $t+J'+1=l$:
\[
o_1(a,u)=\begin{cases}0&a<t\text{ or }t<a<l,\\ A[u<-N]&a=t,\\ -u&a\ge l,\end{cases}\qquad
d_1(a,u)=\begin{cases}1&a<t,\\ A&a=t,\\ 0&t<a<l,\\ -u&a\ge l.\end{cases}
\]
Thus $d_1(a,u)\le y<\pkdlg_{l,t,A}(a,u)$ only for $t<a<l$ and $y=0$. By (I2) with
$I=\{(a,0):t<a<l\}$, whose cost is $\sum_{t<a<l}(0+0+3a)$ ($d_v=0$), by (G) and by
Remark~\ref{pkd:scale}, there is a point
$q\in\pkcPt_R(d,i,e^*;\pkcimp(o_1,\pkdlg_{l,t,A}),\pkdlg_{l,t,A})$.

\emph{(3) $q$ is a stop.} We check, at the jumps $u$ of $R$ at $i^{(-a)}$, that
$\pkcimp(o_1,\pkdlg_{l,t,A})\ge\mathrm{stop}_{\eta,l}$ and $\pkdlg_{l,t,A}\ge\mathrm{stop}_{\eta,l}$,
where $\mathrm{stop}_{\eta,l}(a,u)=\max(0,-u-\eta)$ for $a<l$ and $=-u$ for $a\ge l$.
For $a\ge l$: $o_1(a,u)=-u$, so the improved order is $-u$ if $-u>0$ and $\min(-u,0)=-u$ otherwise;
and $\pkdlg_{l,t,A}(a,u)=-u$. For $a<l$, $a\ne t$: $o_1=0$, the improved order is $\min(1,0)=0$, the
depth is $1$, and $\max(0,-u-\eta)=0$ by \eqref{pkd:eq-updown}. For $a=t$: the jumps of $R$ at $i_t$ lie
in $[-\eta,\eta]$ or in $-[\beta_{i_t},\beta_{i_t}+V]$, so $-u-\eta\le\beta_{i_t}+V\le A$ and the depth $A$ is
$\ge\max(0,-u-\eta)$. If $u<-N$ the improved order is $A$; if $u\ge-N$, then $u>-\beta_{i_t}$ (as
$N<\beta_{i_t}$), so $u$ is not in the negative far window and $u\ge-\eta$; the improved order is
$\min(A,0)=0=\max(0,-u-\eta)$. By Remark~\ref{pkc:mono}, $q$ is a stop (for $\eta,l$) of $R$ over
$\ell^{e^*}d$ at $i$, and it has depth $\pkdlg_{l,t,A}$.

\emph{(4) The move is safe.} Let $y'=\varphi^lx_{R'}(i^{(-l)})-x_{R'}(i)$. The cut $i^{(-l)}$ has shift
$0$, and $l(i^{(-l)})=:l_0$; so $x_{R'}(i^{(-l)})$ has orders $\pkdco_{0,l_0}$, and by
Lemma~\ref{pkp:tau}(v) $\varphi^lx_{R'}(i^{(-l)})$ is a point over $\ell^{e^*}d'$ at $i$ with orders
$\pkcsh_{l,0}(\pkdco_{0,l_0})$. As $x_{R'}(i)$ lies over the same element and has orders
$\pkdco_{s_i,l}$, $y'$ lies over $0$ and $y'\in L^0_i(o')$ with
$o'=\min(\pkdco_{s_i,l},\pkcsh_{l,0}(\pkdco_{0,l_0}))$ (each $L_i(o)$ is a module, monotone in $o$). We
apply Proposition~\ref{pkp:rule} (in the form for the move of $R'$ with the point $q$ of $R$ at rest,
depth $\pkdlg_{l,t,A}$), with $l$, $t=0$ and the killer $\ell^\kappa E^0$ of $T_P$ at $i^{(-l)}$
given by (K). We must show, for all $a\ge0$, $v\in U_R(i^{(-a)})$, $u\in U_{R'}(i^{(-a)})$:
\[
\mathrm{rule}_{l,0}(a,u+v)\le\pkdlg_{l,t,A}(a,v)+o'(a,u).
\]
\emph{$a\ge l$.} Here $\mathrm{rule}_{l,0}(a,u+v)=-u-v$ and $\pkdlg_{l,t,A}(a,v)=-v$. By
Lemma~\ref{pkd:co}(ii), $\pkdco_{s_i,l}(a,u)\ge-u$ (as $a\ge l\ge1$) and
$\pkcsh_{l,0}(\pkdco_{0,l_0})(a,u)=\pkdco_{0,l_0}(a-l,u)\ge-u$ (shift $0$). So $o'(a,u)\ge-u$.
\emph{$a<l$.} Here $\mathrm{rule}_{l,0}=1$ and $\pkcsh_{l,0}(\cdot)(a,u)=0$, so
$o'(a,u)=\min(\pkdco_{s_i,l}(a,u),0)\ge\min(\eta-u,0)$ by Lemma~\ref{pkd:co}(i). If $a\ne t$, then
$u\le\eta$ by \eqref{pkd:eq-updown}, so $o'(a,u)=0$ and the right side is $1$. If $a=t$, the jumps of
$R'$ at $i_t$ lie in $[-\eta,\eta]$ or in $[\beta_{i_t},\beta_{i_t}+V]$ (the positive window), so $u\le
\beta_{i_t}+V\le A$ and $\eta-u\ge\eta-A\ge1-A$; as $A\ge\beta_{i_t}+V>N\ge0$, $A\ge1$ and
$o'(t,u)\ge\min(1-A,0)=1-A$. The right side is $\ge A+1-A=1$.

By Proposition~\ref{pkp:rule} the move of $R'$ by $y'$ with $q$ at rest is $\kappa$-safe. If $R=S$
this is the first alternative of (St) at $i$, and if $R=S'$ the second.

\emph{The walk.} (K), (P) and (St) of Theorem~\ref{pkm:walk} hold at every cut, with $e=e'=e^*$; so $\mathrm{SCC}_N(c,c')$ for $N\ge2e^*+\kappa$.
\end{proof}

\begin{remark}
The hypothesis
$s_{i^{(-l)}}=0$ in a long gap is used twice: for the killer $\ell^\kappa$ in (I5) and for the
killer $\ell^\kappa E^0$ in the rule.
\end{remark}

\subsubsection*{Passing to lattices}

\begin{lemma}[Set-up for the cases]\label{pkd:setup}
Let $U$, $U'$ be jump sets and $L$ crystalline with jumps in $U'$. Fix a Kummer chart with
$\Gamma'=\Gamma$ (Lemma~\ref{pkp:chart-exists}), unramified if $e=1$, and $\zeta$ as in
\S\ref{pki:ss-arch}. Suppose that for every $Y\in\pkiC(U)$, every pairing $P$ modulo some $\ell^m$ of
$S=S(Y)$ and $S'=S(L)$ (the universal sides of Lemma~\ref{pkc:cup}), every $c\in\pllZf(Y)$ and every
$x\in\pllZf(L)$ one has $\mathrm{SCC}_H(c,x)$. Then $\pkiC(U)$ has cup exponent $H$ on $L$.
\end{lemma}

\begin{proof}
Every member of $\pkiC(U)$ and $L$ is crystalline on $\Gamma=\Gamma'$; the jump sets of $S(Y)$, $S(L)$
at the letter $z$ are the letter jump sets $U_z$, $U'_z$ of \S\ref{pki:ss-arch} (as in the proof of
Theorem~\ref{pkc:generic}). Lemma~\ref{pkc:cup} gives $\ell^Hx\in L^{\perp\pkiC(U)}$ for all $x$.
\end{proof}

\subsubsection{Sparse $b$-letters: Proposition~\ref{pki:sparse}}\label{pkd:ss-sparse}

Let $B$ be a set of letters and $J\ge1$. Write $\sigma(i)=i^\ell$. A letter $i$ has \emph{thin
successors} if $\sigma^n(i)\notin B$ for $1\le n\le J$. $B$ is \emph{sparse} if there is no letter $z$
with $z,\sigma^s(z),\sigma^{s+t}(z)\in B$ for some $s,t\in[1,J]$.

\begin{lemma}\label{pkd:thin-add}
Let $B$ be sparse, $j\in B$ a letter without thin successors. Then $j^{(-(J+1))}$ has thin successors.
\end{lemma}

\begin{proof}
Let $k=j^{(-(J+1))}$ and suppose $m=\sigma^n(k)\in B$ with $1\le n\le J$. Put $s=J+1-n\in[1,J]$; then
$\sigma^s(m)=\sigma^{J+1}(k)=j\in B$. As $j$ has no thin successors, $\sigma^t(j)=\sigma^{s+t}(m)\in B$
for some $t\in[1,J]$. So $m,\sigma^s(m),\sigma^{s+t}(m)\in B$, contradicting sparseness.
\end{proof}

\begin{lemma}\label{pkd:thin-cut}
Let $B$ be sparse and $\mathrm{cut}$ the set of letters with thin successors. Then every letter $i$ has
$l(i)\le2J+1$ (in particular $l(i)$ is defined). If $0<r<l(i)$ and $i^{(-r)}\in B$, then
$r+J+1=l(i)$.
\end{lemma}

\begin{proof}
If $i^{(-(J+1))}$ has thin successors, $l(i)\le J+1$. Otherwise some
$\sigma^n(i^{(-(J+1))})=i^{(-r)}$ with $1\le n\le J$ lies in $B$, where $r=J+1-n\in[1,J]$. If
$i^{(-r)}$ has thin successors, $l(i)\le r$; otherwise, by Lemma~\ref{pkd:thin-add},
$i^{(-(r+J+1))}$ has thin successors and $l(i)\le r+J+1\le2J+1$.

Let $l=l(i)$, $0<r<l$, $i^{(-r)}\in B$. Since $i^{(-l)}$ has thin successors and
$\sigma^{l-r}(i^{(-l)})=i^{(-r)}\in B$ with $l-r\ge1$, we get $l-r\ge J+1$. If $r+J+1<l$, then
$i^{(-r)}$ is not a cut (as $0<r<l$), i.e.\ it is in $B$ without thin successors, and
Lemma~\ref{pkd:thin-add} makes $i^{(-(r+J+1))}$ a cut with $0<r+J+1<l$, contradicting the minimality
of $l$. So $r+J+1=l$.
\end{proof}

\begin{theorem}[Proposition~\ref{pki:sparse}]\label{pkd:sparse}
Let $\eta,J\ge1$ and $V,G,b,M,\vartheta\ge0$ with $G+V\le\eta$; let $B$ be a set of letters, $b_z=b$
for $z\in B$ and $b_z=G$ otherwise, and $U_z\subseteq\pkiWin(\eta,-\varepsilon_z,b_z,V)$,
$U'_z\subseteq\pkiWin(\eta,\varepsilon_z,b_z,V)$ for every letter $z$. Assume that $B$ is sparse; that
either (A) $e=1$, $\ell\ge7$ and $10\ell^J\eta\le b$, or (B) $\varepsilon_z$ is constant and
$8e\le\ell^{J+1}$; that $2e(b+5\eta)\le\ell^{J+1}\vartheta$; and that
$20(J+1+d_v)\eta+\pkilog_\ell(2(b+5\eta))\le M$. Then $\pkiC(U)$ has cup exponent
$10^3(\vartheta+(J+1)M)$ on $L$.
\end{theorem}

The proof uses (I1), (I2), and in case (A) also (I5).

\begin{proof}
By Lemma~\ref{pkd:setup} it suffices to prove $\mathrm{SCC}_H(c,x)$ with $H=10^3(\vartheta+(J+1)M)$
for $S=S(Y)$, $S'=S(L)$, a pairing $P$, $c\in C_S$, $x\in C_{S'}$, over a chart with $\Gamma'=\Gamma$,
unramified in case (A). The windows \eqref{pkd:eq-win} hold with $\beta_i=b_i$. We apply
Theorem~\ref{pkd:scheme} with $\mathrm{cut}$ the letters with thin successors (so $l(i)\le2J+1$ by
Lemma~\ref{pkd:thin-cut}), all shifts $0$, and
\[
\kappa=4\vartheta+(J+1)M,\qquad e^*=8\vartheta+6(J+1)M .
\]
Then $2e^*+\kappa=20\vartheta+13(J+1)M\le H$. Put $B_0=2(b+5\eta)\ge1$.

\emph{Jump sizes.} By Lemma~\ref{pkd:jumps}, at a letter $j$ the jumps of $S$, $S'$, $T_P$ are $\le
b_j+V+2\eta$. At a letter not in $B$ this is $\le G+V+2\eta\le3\eta$; at every letter it is
$\le b+3\eta$ (if $b_j=G$, it is $\le3\eta$). Every letter not in $B$ is small ($G+V\le\eta$).

\emph{(K).} Let $i$ be a cut. We apply Proposition~\ref{pkm:kill-J} to each of $S$, $S'$, $T_P$ with
$\omega=b+3\eta$ and $\omega'=3\eta$ (the letters $\sigma^n(i)$, $1\le n\le J$, are not in $B$). In
case (A), $8e=8<49\le\ell^{J+1}$; in case (B) $8e\le\ell^{J+1}$ is assumed. Then
\[
2Jd_v\omega'+\frac{(2J+4)\omega'}{\ell-1}+\frac{8e\omega}{\ell^{J+1}}
\le6(J+1)d_v\eta+12(J+1)\eta+\frac{4\cdot2e(b+5\eta)}{\ell^{J+1}}
\le(J+1)M+4\vartheta=\kappa,
\]
using $\ell-1\ge1$, $2J+4\le4(J+1)$, $b+3\eta\le b+5\eta$ and $6d_v\eta+12\eta\le20(J+1+d_v)\eta\le M$.
So $\ell^\kappa=\ell^\kappa E^0$ kills the three cokernels at $i$.

\emph{(C).} Let $i$ be a cut, $l=l(i)\le2J+1$. As $s_i=0$, Lemma~\ref{pkd:steps} applies with $B_0$
(indeed $b_j+V+2\eta\le b+3\eta\le B_0$), $r_a=V+2\eta\le3\eta$, $2r_a+1\le7\eta$:
\[
\sum_{\mathcal I_i(R)}\pkdcost\le\sum_{a=0}^{2J+1}\bigl(14\eta+2\pkilog_\ell B_0+14\eta(d_v+3a)\bigr)
=(J+1)\bigl(28\eta(1+d_v)+4\pkilog_\ell B_0\bigr)+42\eta(2J+1)(J+1).
\]
This is $(J+1)(84J\eta+70\eta+28d_v\eta+4\pkilog_\ell B_0)\le5(J+1)M$, since
$5M\ge100(J+1+d_v)\eta+5\pkilog_\ell B_0$. So $\kappa+\sum\pkdcost\le\kappa+5(J+1)M\le e^*$.

\emph{(G).} Let $i$ be a letter, $l=l(i)$. If every $i^{(-a)}$, $a<l$, has
$\varepsilon_{i^{(-a)}}=\varepsilon_i$ or is small, $i$ is plain. In case (B) this always holds.
Otherwise we are in case (A), and there is $t<l$ with $\varepsilon_{i^{(-t)}}\ne\varepsilon_i$ and
$i^{(-t)}$ not small; then $i^{(-t)}\in B$ and $t\ne0$, so $t+J+1=l$ by Lemma~\ref{pkd:thin-cut}, and
every $i^{(-a)}$ with $0<a<l$, $a\ne t$, is not in $B$, hence small. We show that $i$ is a long gap
with pivot $t$ and data $(J,s_0,\eta,b+V,\vartheta)$, where $s_0=\lceil(b+2\eta)/(\ell-3)\rceil$.
The chart is unramified and $\ell\ge7$; $t+J+1=l$; $s_{i^{(-l)}}=0$. For $a<l$, $a\ne t$: $a=0$ has
$\varepsilon_i\ne\varepsilon_{i^{(-t)}}$, the others are small. $\eta\le\omega=\eta$, and for $t<a<l$
the letter $i^{(-a)}$ is small, so $\beta+V\le\eta$. Next,
$1+\ell+\dots+\ell^J<\ell^{J+1}/(\ell-1)\le\tfrac76\ell^J$ and $s_0\le(b+2\eta)/4+1$, so
\[
N=s_0+\eta(1+\dots+\ell^J)<\frac b4+\frac\eta2+1+\frac76\ell^J\eta
\le b\Bigl(\frac14+\frac3{140}+\frac7{60}\Bigr)<b=\beta_{i^{(-t)}},
\]
using $b\ge10\ell^J\eta\ge70\eta$ and $\eta/2+1\le3\eta/2$. The jumps of $S$ and $S'$ at every letter are
$\le b+2\eta\le(\ell-3)s_0$ (Lemma~\ref{pkd:jumps}, with $V\le\eta$), so by
Proposition~\ref{pkm:kill-unram} ($s=0$, $t=s_0$) $E^{s_0}$ kills both cokernels at $i^{(-t)}$. Finally
$\beta_{i^{(-t)}}+V=b+V=A\le b+5\eta\le\ell^{J+1}\vartheta/2\le\ell^{J+1}\vartheta$ ($e=1$). The cost
condition: $\sum_{t<a<l}a\le\sum_{a=1}^{2J}a=J(2J+1)$ and $3J(2J+1)\le6J(J+1)\le(J+1)M$ (as $M\ge20(J+1)$),
so $\kappa+4\vartheta+3\sum_{t<a<l}a\le8\vartheta+2(J+1)M\le e^*$.

Theorem~\ref{pkd:scheme} gives $\mathrm{SCC}_H(c,x)$.
\end{proof}

\subsubsection{Any pattern of $b$-letters: Proposition~\ref{pki:any}}\label{pkd:ss-any}

\begin{theorem}[Proposition~\ref{pki:any}]\label{pkd:any}
Let $e=1$, $\ell\ge7$, $\eta\ge1$, and $V,M,G,b,K\ge0$ with $V+4\eta\le M\le K$, $0<G\le b$,
$b<K\ell^2$ and $2G<K\ell$. Let $B$ be a set of letters, $b_z=b$ for $z\in B$ and $b_z=G$ otherwise,
$U_z\subseteq\pkiWin(\eta,-\varepsilon_z,b_z,V)$ and $U'_z\subseteq\pkiWin(\eta,\varepsilon_z,b_z,V)$.
Then $\pkiC(U)$ has cup exponent $10^4(K+\pkilog_\ell(2(b+M)))$ on $L$.
\end{theorem}

The proof uses (I1), (I2), (I5).

\begin{proof}
By Lemma~\ref{pkd:setup} (unramified chart, $d_v=0$, $\Gamma'=\Gamma$) it suffices to prove
$\mathrm{SCC}_H(c,x)$, $H=10^4(K+\pkilog_\ell B_0)$ with $B_0=2(b+M)$, for $S=S(Y)$, $S'=S(L)$, a pairing
$P$ and $c$, $x$; \eqref{pkd:eq-win} holds with $\beta_i=b_i$. Note $K\ge1$ (as $0<2G<K\ell$), $b\ge1$,
and $V+2\eta\le M-2\eta<M$. By Lemma~\ref{pkd:jumps} all jumps of $S$, $S'$, $T_P$ at every letter are
$\le b+V+2\eta\le b+M$, and those at a letter not in $B$ are $\le G+M$.

\emph{Shifts and cuts.} Put $q=\lfloor b/\ell\rfloor$ and $s^+=\max(q-7K,0)$, and for every letter $i$
let $s_i=s^+$ if $\sigma(i)\in B$ and $s_i=0$ otherwise. Let $\mathrm{cut}=\{i:s_i<\beta_i\}$. Every
$i\in B$ is a cut ($s_i\le q<b$, as $b\ge1$), and every $i$ with $\sigma(i)\notin B$ is a cut
($0<G\le\beta_i$). Hence for every letter $i$: if $i^{(-1)}\in B$ then $l(i)=1$; otherwise
$\sigma(i^{(-2)})=i^{(-1)}\notin B$, so $i^{(-2)}$ is a cut and $l(i)\le2$. The shifts satisfy the
condition of Definition~\ref{pkd:def-cut}. Put
\[
\kappa=14K,\qquad e^*=76K+6\pkilog_\ell B_0,
\]
so that $2e^*+\kappa=166K+12\pkilog_\ell B_0\le H$.

\emph{(K).} Let $i$ be a cut. If $\sigma(i)\notin B$, then $s_i=0$, and Proposition~\ref{pkm:kill-J}
with $J=1$, $\omega=b+M$, $\omega'=G+M$ (the jumps at $\sigma(i)$) applies to $S$, $S'$, $T_P$:
$8e=8\le\ell^2$, $d_v=0$, and
\[
\frac{6(G+M)}{\ell-1}+\frac{8(b+M)}{\ell^2}<\frac{3K\ell}{\ell-1}+K+8K+K\le\tfrac72K+10K<14K,
\]
using $G<K\ell/2$, $\ell/(\ell-1)\le7/6$, $6M/(\ell-1)\le M\le K$, $b<K\ell^2$ and $8M/\ell^2<K$. If
$\sigma(i)\in B$, we use Proposition~\ref{pkm:kill-unram} with $s=\kappa$, $t=s_i$: it suffices that
$b+M\le(\ell-3)(\kappa+s_i)$. Now $\kappa+s_i\ge7K+q$, and $(\ell-3)q\ge(\ell-3)(b/\ell-1)=b-3b/\ell-\ell+3
>b-3K\ell-\ell+3$, while $(\ell-3)\cdot7K\ge4K\ell$ (as $\ell-3\ge4\ell/7$). So
$(\ell-3)(7K+q)>b+K\ell-\ell+3\ge b+K+2\ge b+M$, since $K\ell-\ell-K+1=(K-1)(\ell-1)\ge0$.

\emph{(C).} Let $i$ be a cut, $l=l(i)\le2$. In Lemma~\ref{pkd:steps} take $B_0$: $b_j+s_i+V+2\eta\le
b+q+M\le B_0$. At level $0$, $r_0=s_i+V+2\eta\le s_i+M$ and $s_i\le q<K\ell$, so
$(2r_0+1)/(\ell-1)\le 2s_i/(\ell-1)+(2M+1)/(\ell-1)<\tfrac73K+\tfrac12K<3K$ (using $2M+1\le3K$); the
level-$0$ term is $<6K+2\pkilog_\ell B_0$ (the factor $d_v+3a$ is $0$). At levels $a=1,2$,
$2r_a+1\le2M+1\le3K$, so the term is $\le K+2\pkilog_\ell B_0+18aK$. The total is
$<6K+2K+54K+6\pkilog_\ell B_0=62K+6\pkilog_\ell B_0$, and $\kappa+62K+6\pkilog_\ell B_0\le e^*$.

\emph{(G).} Let $i$ be a letter, $l=l(i)\in\{1,2\}$. If $l=1$, or $l=2$ and
$\varepsilon_{i^{(-1)}}=\varepsilon_i$, then $i$ is plain with $\varrho=\varepsilon_i$. Let $l=2$ and
$\varepsilon_{i^{(-1)}}\ne\varepsilon_i$. Then $i^{(-1)}$ is not a cut, so $i^{(-1)}\notin B$ (thus
$\beta_{i^{(-1)}}=G$) and $\sigma(i^{(-1)})=i\in B$; and $s_{i^{(-1)}}=s^+\ge G>0$, so $G+7K\le q$. Also
$\sigma(i^{(-2)})=i^{(-1)}\notin B$ gives $s_{i^{(-2)}}=0$. We show that $i$ is a long gap with pivot
$t=0$ and data $(1,s_0,\omega,A,\delta)=(1,q+12K,G+M,b+V,2K)$. The chart is unramified, $\ell\ge7$,
$0+1+1=l$, $s_{i^{(-2)}}=0$. The only $a<2$, $a\ne0$, is $a=1$, where
$\varepsilon_{i^{(-1)}}\ne\varepsilon_i$. $\eta\le G+M$, and $\beta_{i^{(-1)}}+V=G+V\le G+M$. Next, from
$G+M+6K\le G+7K\le q$:
\[
N=q+12K+(G+M)(1+\ell)\le q+12K+(q-6K)(1+\ell)=(\ell+2)q-6K(\ell-1)
<b+2K\ell-6K\ell+6K<b=\beta_i,
\]
using $\ell q\le b$ and $q\le b/\ell<K\ell$. $E^{s_0}$ kills the cokernels of $S$, $S'$ at $i$ by
Proposition~\ref{pkm:kill-unram} ($s=0$, $t=s_0$): $(\ell-3)(q+12K)>b-3K\ell-\ell+3+6K\ell\ge b+M$, as
above. $\beta_i+V=b+V=A$, and $A\le b+M<K\ell^2+K\le2K\ell^2=\ell^2\delta$. The cost: $\sum_{0<a<2}a=1$,
and $\kappa+4\delta+3=22K+3\le e^*$.

Theorem~\ref{pkd:scheme} gives $\mathrm{SCC}_H(c,x)$.
\end{proof}

\subsubsection{Windows, jumps and characters of the data: proof of Proposition~\ref{pki:data}}\label{pkd:ss-data}

We use the data of {\cite[\S2]{P55}} and of {\cite[(36)]{P46}}: $h=r^8$, $\Delta=2h$, $t=5r$,
$k=r+1-m$, $\sqrt x\le m\le2\sqrt x$, $x^5\le r\le3x^5$, the weight $\kappa(\tau)=b$ at the $b$-places
and $\kappa(\tau)=G$ otherwise, and $\kappa_j(\tau)=\kappa(\tau)+(j+\frac12)\Delta=\kappa(\tau)+j\Delta+h$.
The summands $L_c$ are as in {\cite[Proposition~5.2]{P55}} and the paragraph before it: the main summand
$L_{\mathrm{main}}$, of rank $r'\le r$, has at $\gamma$ the jumps $J_{\sigma_\gamma}$ of
{\cite[Definition~2.6(M2)]{P55}} (with multiplicity), where $\sigma_\gamma$ is the complex embedding
attached to $\gamma$ through $\iota$ and $v$ (Convention~\ref{pkj:conv}, reading (R$^-$)), and the characters $L_e$, $0\le e<t'\le t$, have the single jump $e$. $U(\gamma)$ is
the union of the jumps at $\gamma$ of all summands. For a lattice $L$ we use the convention of
Lemma~\ref{pkp:tau} and Definition~\ref{pkp:def-pairing}: $\auO(1)$ has the jump $-1$, so
$L^*=\operatorname{Hom}(L,\auO(1))$ has the jump $-u-1$ (with the multiplicity of $u$) for each jump
$u$ of $L$. Let $\varepsilon(\gamma)=1$ if $\sigma_\gamma$ is in the CM type of the data and $-1$
otherwise, $\eta=h^2+1$, $W=4t$, $\delta_j=j\Delta+h-2t$, and $a_\gamma=\kappa(\sigma_\gamma)$ ($=b$
if $\gamma$ is thick, $=G$ otherwise, in the sense of Proposition~\ref{pki:chart}).

\begin{lemma}[Elementary sizes]\label{pkd:sizes}
$t\le h/3$, $4m h\le h^2$, $t\le h^2$, and $h\ge4$.
\end{lemma}

\begin{proof}
$r\ge x^5\ge10^2$ (as in the proof of Lemma~\ref{pki:data-sizes}), so $h=r^8\ge15r=3t$; $4m\le
8\sqrt x\le8r\le h$; $t\le h\le h^2$.
\end{proof}

\begin{theorem}[Proposition~\ref{pki:data}, remaining assertions]\label{pkd:data}
Let $j\le m$ be the index of {\cite[Definition~2.6(M2)]{P55}} for $L_{\mathrm{main}}$, and let
$U'(\gamma)$ be the jump set of $L_c^*$.
\begin{enumerate}
\item[(i)] For every summand $L_c$ and every $\gamma$: $U(\gamma)\subseteq\pkiWin(\eta,-\varepsilon(\gamma),
a_\gamma+\delta_j,W)$ and $U'(\gamma)\subseteq\pkiWin(\eta,\varepsilon(\gamma),a_\gamma+\delta_j,W)$.
\item[(ii)] $L_{\mathrm{main}}^*$ has few low jumps for $(h,\lfloor r'm/k\rfloor,\eta)$ and the jump sets
$U$, $U'$ (with $c=\mathrm{main}$).
\item[(iii)] For a character $L_e$, $U'(\gamma)=\{-e-1\}\subseteq[-\eta,0]$.
\item[(iv)] Let $\ell=p_i$ and $v\mid\ell$. Then $E_v=\Q_\ell$, $\sigma_\gamma$, $\varepsilon(\gamma)$ and
$a_\gamma$ do not depend on $\gamma$, and $\ell>a+\delta_j+6(W+\eta)+r^2(N+1)^2+2$ for $a\in\{b,G\}$,
$N=\max(r+t,2r)\,r$ $(=(r+t)r=6r^2$, as $t=5r)$. If, for every inertial $s$, $q_s$ denotes the monic polynomial over $\Z_\ell$ of degree
$\le N$ of the definition of $\mathcal C(U,\mathrm{tame})$ ({\cite[\S5]{P55}}), which depends only on $s$, and
$K(s)=\{X^{e'}-1,q_s\}$, then $\mathcal C_v\subseteq\pkiC(U,K)$; and $L_c^*\in\pkiC(U',K^{-1})$, by
the input (I6) below.
\end{enumerate}
\end{theorem}

The inclusion in (iv) uses the definition of $\mathcal C(U,\mathrm{tame})$ with polynomials $q_s$
depending only on $s$, which (G4) provides (Proposition~\ref{pkw:G4}); the last assertion of (iv)
uses (I6), which is Lemma~\ref{pkw:dual}.

\begin{enumerate}
\item[(I6)] \emph{Inertia on duals.} If $Y$ is a potentially crystalline lattice, $s$ is inertial and
$P\in\auO[X]$ kills $s^{-1}$ on $D_{\mathrm{pcris}}(Y)$, then $P$ kills $s$ on
$D_{\mathrm{pcris}}(\operatorname{Hom}(Y,\auO(1)))$. This is Lemma~\ref{pkw:dual}; the jumps of
$\operatorname{Hom}(Y,\auO(1))$ are the $-u-1$ by Proposition~\ref{pkw:HT}(ii).
\end{enumerate}

\begin{proof}
Write $S_\gamma=a_\gamma+\delta_j$, so $\kappa_j(\sigma_\gamma)=S_\gamma+2t$, and $\varepsilon=\varepsilon(\gamma)$.

\emph{The jumps of $L_{\mathrm{main}}$.} By (M2) each $u\in J_{\sigma_\gamma}$ is low,
$\Delta/4\le|u|\le2m\Delta=4mh$, or high, $|u+\varepsilon\kappa_j(\sigma_\gamma)|\le t$ (the sign $+$
for $\sigma$ in the CM type). If low, $|u|\le h^2<\eta$ (Lemma~\ref{pkd:sizes}). If high,
$-\varepsilon u\in[S_\gamma+t,S_\gamma+3t]$, so $u\in-\varepsilon[S_\gamma,S_\gamma+W]$. In both cases
$|u|\ge h/2$: $\Delta/4=h/2$, and $\kappa_j-t\ge h-t\ge h/2$.

(i) The characters have jumps $e\in[0,t)\subseteq[-\eta,\eta]$. So
$U(\gamma)\subseteq\pkiWin(\eta,-\varepsilon,S_\gamma,W)$. For $U'(\gamma)$: if $c$ is a character,
$-e-1\in[-t,-1]\subseteq[-\eta,\eta]$. If $c=\mathrm{main}$ and $u$ is low, $|-u-1|\le4mh+1\le\eta$. If
$u$ is high and $\varepsilon=1$, then $-u\in[S_\gamma+t,S_\gamma+3t]$ and $-u-1\in[S_\gamma,S_\gamma+W]$
(as $t\ge1$); if $\varepsilon=-1$, then $u\in[S_\gamma+t,S_\gamma+3t]$ and
$-u-1\in-[S_\gamma+t+1,S_\gamma+3t+1]\subseteq-[S_\gamma,S_\gamma+W]$. So
$U'(\gamma)\subseteq\pkiWin(\eta,\varepsilon,S_\gamma,W)$.

(ii) A negative element of $U(\gamma)$ is a jump of $L_{\mathrm{main}}$ (the characters' jumps are
$\ge0$), so $2u\le-h$. A negative element of $U'(\gamma)$ is $-u-1$ with $u\in J_{\sigma_\gamma}$,
$u\ge0$, hence $u\ge h/2$ and $2(-u-1)\le-h-2$. A jump $-u-1$ of $L_{\mathrm{main}}^*$ lies in
$[-\eta,\eta]$ only if $|u|\le\eta+1=h^2+2\le2h^2$; by (M2) at most $r'm/k$ elements of
$J_{\sigma_\gamma}$, with multiplicity, satisfy this, and the number is an integer, so at most
$\lfloor r'm/k\rfloor$.

(iii) Clear, as $t\le\eta$.

(iv) $p_i\equiv1\pmod{q_2^2}$ with $q_1d\mid q_2$ splits completely in $E$ ({\cite[\S2]{P55}}), so
$E_v=\Q_\ell$; every $\gamma$ restricts to the unique embedding of $E_v$, so the complex embedding
$\sigma_\gamma$ attached to $\gamma$ through $\iota$ and $v$, and with it $\varepsilon(\gamma)$ and
$a_\gamma$, is independent of $\gamma$. Size: $a\le b$, $\delta_j\le2mh+h\le h^2$, $W+\eta\le2h^2$,
$r^2(N+1)^2\le r^2(6r^2+1)^2\le49r^6\le h$ (here $N=\max(r+t,2r)r=6r^2$ and $r\ge10^2$), so the right side is at most $b+14h^2+3\le2b<b^{10}\le p_i=\ell$
(as $h^3\le b$ by (S2)). Inclusion: a member of $\mathcal C_v=\mathcal C(U,\mathrm{tame})$ is potentially
crystalline with jumps in $U$, and on its $D_{\mathrm{pcris}}$ every inertial $s$ is killed by
$X^{e'}-1$ and by $q_s$, i.e.\ by every element of $K(s)$; and $K(s)$ contains $X^{e'}-1$ and the
image of the monic $q_s$ of degree $\le N$, as Proposition~\ref{pki:tame} requires. Finally, if
$L_c\in\mathcal C_v$, then for inertial $s$ every element of $K(s^{-1})$ kills $s^{-1}$ on
$D_{\mathrm{pcris}}(L_c)$, hence by (I6) kills $s$ on $D_{\mathrm{pcris}}(L_c^*)$; and $L_c^*$ is
potentially crystalline with jumps in $U'$ by definition. So $L_c^*\in\pkiC(U',K^{-1})$.
\end{proof}

\begin{remark}[Use in the deduction of {\cite[Theorem~5.8]{P55}}]\label{pkd:rem-deduce}
In the proof of {\cite[Theorem~5.8]{P55}} given Propositions~\ref{pki:generic}--\ref{pki:data}
(\S\ref{pki:ss-deduce}), case (a) applies Proposition~\ref{pki:tame} to $L=L_c^*$, which requires
$L_c^*\in\pkiC(U',K^{-1})$: this is the last assertion of Theorem~\ref{pkd:data}(iv), and it is part of the last clause of
Proposition~\ref{pki:data}. A cup exponent of
$\pkiC(U,K)$ on $L$ is one of every subclass, in particular of $\mathcal C_v$. With
Theorems~\ref{pkd:sparse}, \ref{pkd:any}, \ref{pkd:data} and \S\ref{pkc:sec}, the deduction of
{\cite[Theorem~5.8]{P55}} in \S\ref{pki:ss-deduce} rests only on Proposition~\ref{pki:chart} (Theorem~\ref{pkj:chart}) and on the
inputs (I1)--(I6), (T1)--(T4), all proved in \S\S\ref{pkt:sec}--\ref{pkw:sec}.
\end{remark}



\providecommand{\auO}{\mathcal O}
\providecommand{\auQlb}{\overline{\Q}_\ell}
\providecommand{\pkpA}{A_{\mathrm{inf}}}
\providecommand{\pkpW}{W(C^\flat)}
\providecommand{\pkpAmb}{\mathbb A}
\providecommand{\pkmJ}[3]{J^{#1}_{#2}(#3)}
\providecommand{\pkmI}[2]{I_{#1}(#2)}
\providecommand{\pkiC}{\mathcal C}
\providecommand{\pkwB}{B_{\mathrm{cris}}}
\providecommand{\pkwD}{D_{\mathrm{pcris}}}
\providecommand{\pkwur}{\Q_\ell^{\mathrm{ur}}}
\providecommand{\pkwI}{I_{\Q_\ell}}

\section{Tame charts, inertia on $D_{\mathrm{pcris}}$, and the tame inputs}\label{pkw:sec}

This section proves the inputs (T1)--(T4) of \S\ref{pkc:ss-tame} and the input (I6) of
\S\ref{pkd:ss-data}, and settles two points of convention used in \S\ref{pkd:sec}: the jumps of
Tate duals, and the shape of the class $\mathcal C(U,\mathrm{tame})$. We use the conventions and the
labels of \S\S\ref{pkp:sec}--\ref{pkd:sec}.

\emph{Dependencies.} Proposition~\ref{pkw:T1} (input (T1)) and Lemma~\ref{pkw:letters} use only
elementary local field theory. Theorem~\ref{pkw:T2} (input (T2)), Theorem~\ref{pkw:T4} (input (T4)),
Lemma~\ref{pkw:dual} (input (I6)), Proposition~\ref{pkw:HT} and Proposition~\ref{pkw:G4} use the published facts
of $p$-adic Hodge theory and of local ramification theory stated in \S\ref{pkw:ss-pcris} and
\S\ref{pkw:ss-checks}. Theorem~\ref{pkw:Z} and Theorem~\ref{pkw:T3} (input (T3)) use the results of
\S\S\ref{pkn:sec}--\ref{pkl:sec} and (F6).

\subsection{$D_{\mathrm{pcris}}$ and the operators $\beta_s$}\label{pkw:ss-pcris}

$\Gamma=D_v$ is an open subgroup of $G_{\Q_\ell}$ through $j$; $\pkwI\subseteq G_{\Q_\ell}$ is the inertia
group, i.e.\ $\operatorname{Gal}(\auQlb/\pkwur)$ with $\pkwur$ the maximal unramified extension of
$\Q_\ell$ in $\auQlb$; $s\in\Gamma$ is \emph{inertial} if $j(s)\in\pkwI$. $\pkwB$ is Fontaine's
crystalline period ring, with its action of $G_{\Q_\ell}$, $t\in\pkwB$ the element with $gt=
\chi_{\mathrm{cyc}}(g)t$ ($\chi_{\mathrm{cyc}}$ the cyclotomic character), a unit of $\pkwB$, and
$P_0=\widehat{\pkwur}\subseteq\pkwB$ the completion of $\pkwur$ (the fraction field of
$W(\overline{\mathbb F}_\ell)$). For an open subgroup $H\subseteq G_{\Q_\ell}$ we write $K_H$ for its fixed
field and $K_{H,0}$ for the maximal unramified subfield of $K_H$. For a lattice $Y$ (a finite free
$\auO$-module with a continuous $\auO$-linear action of $\Gamma$) we put $\pkwB(Y)=\pkwB\otimes_{\Z_\ell}Y$,
with $g(b\otimes y)=gb\otimes gy$ for $g\in\Gamma$ and $\auO$ acting through $Y$.

We use the following published facts, with their hypotheses.

\begin{enumerate}
\item[(F1)] (Fontaine \cite[1.4.1--1.4.2, 5.1.2 and 5.6.1]{pkfFontaineIII}, \emph{Repr\'esentations $p$-adiques semi-stables}.) For every open subgroup $H\subseteq G_{\Q_\ell}$, $\pkwB^{H}=K_{H,0}$; and for
every continuous representation of $H$ on a finite dimensional $\Q_\ell$-vector space $V$, the natural
map $\pkwB\otimes_{K_{H,0}}(\pkwB\otimes_{\Q_\ell}V)^H\to\pkwB\otimes_{\Q_\ell}V$ is injective. $V$ is
\emph{crystalline} (as a representation of $H$) if this map is also surjective, equivalently
$\dim_{K_{H,0}}(\pkwB\otimes V)^H=\dim_{\Q_\ell}V$.
\item[(F2)] (Proved in Lemma~\ref{pkw:F2} below from (F1); it is the form, for crystalline periods, of
Fontaine's criterion \cite[1.8.5]{pkfFontaineIII} for potentially admissible representations,
together with Galois descent.) Let $H\subseteq
G_{\Q_\ell}$ be open and $V$ a representation of $H$ that is crystalline on some open subgroup of $H$. If
$H\cap\pkwI$ acts trivially on $\bigcup_{H'}(\pkwB\otimes V)^{H'}$ ($H'$ running over the open subgroups
of $H$), then $V$ is crystalline as a representation of $H$.
\item[(F3)] (Serre \cite[IV, \S2, Proposition~7 and Corollary~1, p.~67]{SerreLF}, the structure of tamely ramified extensions.) Every finite extension of $\pkwur$ inside $\auQlb$ whose degree $m$ is prime to $\ell$ is
$\pkwur(\varpi_m)$ for an element $\varpi_m$ with $\varpi_m^m=\ell$. ($\pkwur$ is henselian, discretely
valued with uniformiser $\ell$ and algebraically closed residue field, so such an extension $L$ is totally
and tamely ramified. If $\pi_L$ is a uniformiser of $L$, then $\pi_L^m=\ell u$ with $u$ a unit of the ring of
integers of $L$; the residue field of $L$ is algebraically closed and $\ell\nmid m$, so by Hensel's lemma
$u=w^m$ with $w\in L$, and $\varpi_m=\pi_L/w$ satisfies $\varpi_m^m=\ell$ and generates $L$, as
$X^m-\ell$ is an Eisenstein polynomial of degree $m=[L:\pkwur]$. The cited statements of Serre give the
cyclicity of the tame inertia quotient.)
\end{enumerate}

\begin{lemma}[(F2)]\label{pkw:F2}
Statement (F2) holds.
\end{lemma}

\begin{proof}
Let $H'_0\subseteq H$ be open with $V$ crystalline on $H'_0$, and $H'=\bigcap_{h\in H}hH'_0h^{-1}$, a
finite intersection, hence open and normal in $H$. $V$ is crystalline on $H'$: by (F1) the map
$\pkwB\otimes_{K_{H'_0,0}}(\pkwB\otimes V)^{H'_0}\to\pkwB\otimes V$ is surjective, and
$(\pkwB\otimes V)^{H'_0}\subseteq(\pkwB\otimes V)^{H'}$, so $\pkwB\otimes_{K_{H',0}}(\pkwB\otimes V)^{H'}
\to\pkwB\otimes V$ is surjective, and injective by (F1). Put $L=K_{H'}$, $F'=K_{H',0}=\pkwB^{H'}$ (F1),
$D'=(\pkwB\otimes V)^{H'}$, an $F'$-vector space.

Since $H'$ is normal in $H$, $H$ preserves $F'$ and $D'$, and acts on $D'$ semilinearly with respect to its
action on $F'$. We claim that $H\cap G_{F'}=(H\cap\pkwI)H'$. The inclusion $\supseteq$ holds as $\pkwI$
fixes the unramified field $F'$ and $H'=G_L$ fixes $L\supseteq F'$. Conversely let $g\in H\cap G_{F'}$.
As $L\cap\pkwur=F'$ ($F'$ is the maximal unramified subfield of $L$) and $\pkwur/\Q_\ell$ is Galois, restriction
maps $G_L$ onto $\operatorname{Gal}(\pkwur/F')$; so some $h\in G_L=H'$ agrees with $g$ on $\pkwur$, and
$h^{-1}g\in\pkwI\cap H$. This proves the claim.

By hypothesis $H\cap\pkwI$ acts trivially on $D'\subseteq\bigcup_{H''}(\pkwB\otimes V)^{H''}$, and $H'$ acts
trivially on $D'$. Hence the action of $H$ on $D'$ and on $F'$ factors through the finite group
$\Delta=H/(H\cap G_{F'})$, which acts faithfully on $F'$. By Artin's theorem $F'/F$ is Galois with group
$\Delta$, where $F=F'^{\Delta}=(\pkwB^{H'})^H=\pkwB^H=K_{H,0}$ (F1).

\emph{Galois descent.} We show that $F'\otimes_FD'^{\Delta}\to D'$ is surjective (it is all we need). Let
$D_1$ be the $F'$-span of $D'^\Delta$ and suppose $D_1\ne D'$; choose an $F'$-linear $f\colon D'\to F'$ with
$f(D_1)=0$ and $f(d)\ne0$ for some $d\in D'$. For every $\lambda\in F'$, $\sum_{\delta\in\Delta}\delta(\lambda d)
\in D'^\Delta$, so $0=\sum_\delta\delta(\lambda)f(\delta d)$ for all $\lambda\in F'$. The distinct field
automorphisms $\delta$ of $F'$ are linearly independent over $F'$ (Dedekind--Artin), so $f(\delta d)=0$ for
every $\delta$, in particular $f(d)=0$, a contradiction.

Now $D'^\Delta=(\pkwB\otimes V)^H$, and the composite
$\pkwB\otimes_{K_{H,0}}(\pkwB\otimes V)^H\to\pkwB\otimes_{F'}D'\to\pkwB\otimes V$ is surjective, the
first map by Galois descent and the second because $V$ is crystalline on $H'$. By (F1), $V$ is crystalline
on $H$.
\end{proof}

\begin{definition}[$D_{\mathrm{pcris}}$, $\beta_s$, killing]\label{pkw:def-pcris}
Let $Y$ be a lattice and $V=Y[1/\ell]$. Put $\pkwD(Y)=\bigcup_H\pkwB(Y)^H$, the union over the open
subgroups $H$ of $\Gamma$. $Y$ is \emph{potentially crystalline} if $\pkwB(Y)^H$ spans $\pkwB(Y)$ over
$\pkwB$ for some open $H\subseteq\Gamma$ (by (F1), equivalently $V$ is crystalline on $H$). For such $Y$
and an inertial $s\in\Gamma$, $\beta_s=\beta_s^Y$ denotes the unique $\pkwB$-linear endomorphism of
$\pkwB(Y)$ that agrees with $s$ on $\pkwD(Y)$ (Lemma~\ref{pkw:beta}(i)). For $P=\sum_kP_kX^k\in\auO[X]$,
``$P$ \emph{kills $s$ on $\pkwD(Y)$}'' means $P(\beta_s)=\sum_kP_k\beta_s^k=0$ on $\pkwB(Y)$,
equivalently (Lemma~\ref{pkw:beta}(i)) $\sum_kP_ks^k=0$ on $\pkwD(Y)$.
\end{definition}

This is the meaning of the phrase in Proposition~\ref{pki:tame}, in \S\ref{pkc:ss-tame} and in
Theorem~\ref{pkd:data}, and in the definition of $\mathcal C(U,\mathrm{tame})$ in {\cite[\S5.1]{P55}}.

\begin{lemma}[The operators $\beta_s$]\label{pkw:beta}
Let $Y$, $Y'$ be potentially crystalline lattices and $s,s'\in\Gamma$ inertial.
\begin{enumerate}
\item[(i)] $\beta_s$ exists and is unique; it commutes with $\auO$; $\beta_{ss'}=\beta_s\beta_{s'}$ and
$\beta_1=1$, so $\beta_s$ is invertible with inverse $\beta_{s^{-1}}$. If $P\in\auO[X]$, then
$P(\beta_s)=0$ if and only if $\sum_kP_ks^k$ vanishes on $\pkwD(Y)$. If $s\in H$ for an open $H$ such that
$\pkwB(Y)^H$ spans $\pkwB(Y)$, then $\beta_s=1$.
\item[(ii)] If $f\colon Y\to Y'$ is $\auO$-linear and $\Gamma$-equivariant, then $(1\otimes f)\beta^Y_s=
\beta^{Y'}_s(1\otimes f)$. If $f$ is injective and $P$ kills $s$ on $\pkwD(Y')$, then $P$ kills $s$ on
$\pkwD(Y)$.
\item[(iii)] If $\Gamma$ acts trivially on $Y$, then $Y$ is potentially crystalline and $\beta_s=1$. The
lattice $\auO(1)=\auO\otimes_{\Z_\ell}\Z_\ell(1)$ is potentially crystalline and $\beta_s=1$ on it.
\end{enumerate}
\end{lemma}

\begin{proof}
(i) Let $H_0\subseteq\Gamma$ be open with $\pkwB(Y)^{H_0}$ spanning, and let $H\subseteq H_0$ be an open
subgroup normal in $\Gamma$ (the intersection of the finitely many conjugates of $H_0$). Then
$D_H=\pkwB(Y)^H\supseteq\pkwB(Y)^{H_0}$ spans, so by (F1) the map $\pkwB\otimes_{K_{H,0}}D_H\to\pkwB(Y)$
is bijective. As $H$ is normal, $s$ maps $D_H$ into itself; as $s$ is inertial it fixes
$K_{H,0}\subseteq\widehat{\pkwur}$, so $s|_{D_H}$ is $K_{H,0}$-linear, and also $\auO$-linear. Define
$\beta_s$ as $1\otimes s|_{D_H}$ transported by this bijection; it is $\pkwB$-linear, commutes with
$\auO$ and agrees with $s$ on $D_H$. If $x\in\pkwD(Y)$, say $x$ fixed by an open $H'$, let $H''\subseteq
H\cap H'$ be open and normal in $\Gamma$; the operator built from $H''$ agrees with $s$ on $D_{H''}\supseteq
D_H\cup\{x\}$, hence with $\beta_s$ on the spanning set $D_H$, hence everywhere; so $\beta_sx=sx$.
Uniqueness holds because $\pkwD(Y)$ spans. $\beta_s\beta_{s'}$ and $\beta_{ss'}$ are $\pkwB$-linear and
agree with $ss'$ on $\pkwD(Y)$ (which $s'$ maps into itself: $s'$ maps the vectors fixed by $H'$ to
vectors fixed by $s'H's'^{-1}$); so they are equal. For $P\in\auO[X]$, $P(\beta_s)$ is $\pkwB$-linear and
agrees with $\sum_kP_ks^k$ on the spanning set $\pkwD(Y)$. If $s\in H$ with $\pkwB(Y)^H$ spanning, $\beta_s$
and $1$ agree on $\pkwB(Y)^H$.

(ii) $1\otimes f$ is $\pkwB$-linear and equivariant, so it maps $\pkwD(Y)$ into $\pkwD(Y')$; both sides of
the identity are $\pkwB$-linear and agree on $\pkwD(Y)$. If $f$ is injective, so is $1\otimes f$
($\pkwB$ is flat over $\Z_\ell$), and $(1\otimes f)P(\beta^Y_s)=P(\beta^{Y'}_s)(1\otimes f)=0$.

(iii) If the action is trivial, $1\otimes Y\subseteq\pkwB(Y)^\Gamma$ spans, so $\beta_s=1$ by (i). For
$\auO(1)$ let $e_1$ be a basis of $\Z_\ell(1)$, so $ge_1=\chi_{\mathrm{cyc}}(g)e_1$; then
$g(t^{-1}\otimes a\otimes e_1)=\chi_{\mathrm{cyc}}(g)^{-1}t^{-1}\otimes a\otimes\chi_{\mathrm{cyc}}(g)e_1=
t^{-1}\otimes a\otimes e_1$ for $a\in\auO$, and these vectors span ($t$ is a unit); so $\Gamma$ itself is an
admissible $H$ and $\beta_s=1$ by (i).
\end{proof}

\begin{lemma}[Extensions]\label{pkw:ext}
Let $S$ be a side (Definition~\ref{pkp:def-side}) whose lattice $Y$ is potentially crystalline, and
let $q\in\auO[X]$ kill an inertial $s$ on $\pkwD(Y)$. Then the extension $X_S$ is potentially crystalline
and $q(X)\cdot(X-1)$ kills $s$ on $\pkwD(X_S)$.
\end{lemma}

\begin{proof}
$X_S$ is crystalline on $\Gamma'$ (Definition~\ref{pkp:def-side}: it is a lattice with jump sets),
hence potentially crystalline. The inclusion $\iota\colon Y\to X_S$ and the projection $\pi\colon X_S\to
C_S$ are equivariant, and $0\to\pkwB(Y)\to\pkwB(X_S)\to\pkwB(C_S)\to0$ is exact. By
Lemma~\ref{pkw:beta}(ii),(iii), $(1\otimes\pi)(\beta_s-1)=(\beta^{C_S}_s-1)(1\otimes\pi)=0$, so for
$x\in\pkwB(X_S)$, $(\beta_s-1)x=(1\otimes\iota)y$ with $y\in\pkwB(Y)$, and
$q(\beta_s)(\beta_s-1)x=(1\otimes\iota)q(\beta^Y_s)y=0$.
\end{proof}

\subsection{(T1): tame charts}\label{pkw:ss-T1}

\emph{The field of a chart.} In Definition~\ref{pkp:def-chart} the ramification index $e$, the residue
degree $f$ and the uniformiser $\pi=\theta(z)$ are those of the fixed field
$K'=\auQlb^{\,j(\Gamma')}$ of the subgroup $\Gamma'$ of the chart, not of $\auQlb^{\,j(\Gamma)}$: this is how
the chart is used throughout (the conditions on $\Gamma'$ say exactly that $j(\Gamma')=G_{K'}$ with
$K'=K'_0(\pi)$, $K'_0=\Lambda[1/\ell]$, as in Lemma~\ref{pkf:invariants}), and the two readings agree
when $\Gamma'=\Gamma$, as in Lemma~\ref{pkp:chart-exists}. A tame chart has $\Gamma'\ne\Gamma$ in general,
and needs this reading, which is the one stated in Definition~\ref{pkp:def-chart}.

\begin{definition}[Tame chart]\label{pkw:def-tame}
Let $E_v=\Q_\ell$, so that $j\colon\Gamma\to G_{\Q_\ell}$ is bijective. A \emph{tame chart} is a Kummer chart
(Definition~\ref{pkp:def-chart}, read with $K'$ as above) with ramification index $e$ and residue degree
$f$, together with an element $\varpi\in\auQlb$ and an element $s_0\in\Gamma$ (the \emph{tame generator}),
such that: $\ell\nmid e$; $\varpi^e=\ell$, $\theta(z)=\varpi$ and $E=\ell-z^e$; $\Gamma'$ is normal in $\Gamma$; $s_0$ is
inertial and fixes $\Lambda$ pointwise; and $s_0z=\xi z$ with $\xi=[a]\in\Lambda$ the Teichm\"uller lift of
a primitive $e$-th root of unity $a\in\mathbb F_{\ell^f}^\times$. Then $s_0^kz=\xi^kz$ for all $k\ge0$
($s_0$ fixes $\xi$), and $d_v\le1$ ($d_v=v_\ell(e)+1=1$ if $e>1$).
\end{definition}

\begin{lemma}[Letters for any root of $X^{\ell^f-1}-1$]\label{pkw:letters}
In Definition~\ref{pkp:def-letters} let $\zeta\in\auO$ be any element with $\zeta^{N}=1$, $N=\ell^f-1$.
Then the $e_i$ defined there are orthogonal idempotents of $\operatorname{End}(\pkpAmb(L))$ with
$\sum_ie_i=1$ and $(1\otimes\zeta)e_i=[i]e_i$; the image of $e_i$ is the letter part
$\pkpAmb(L)_i=\{x:(1\otimes\zeta)x=[i]x\}$, so $\pkpAmb(L)=\bigoplus_i\pkpAmb(L)_i$; $\varphi e_i=
e_{i^\ell}\varphi$, and $ge_i=e_ig$ for every $g\in G_{\Q_\ell}$ fixing $\Lambda$ (in particular for
$g\in\Gamma'$). For $\zeta=1$: $\pkpAmb(L)_1=\pkpAmb(L)$ and $\pkpAmb(L)_i=0$ for $i\ne1$, and every
embedding $\gamma\colon\auO\to C$ has letter $1$.
\end{lemma}

\begin{proof}
Write $T=1\otimes\zeta$ and $\lambda_i=[i]$. In $\Lambda[X]$, $X^N-1=\prod_{i}(X-\lambda_i)$, the product
over the $N$ letters (the $\lambda_i$ are $N$ distinct roots of $X^N-1$ in the domain $\Lambda$); hence
$\prod_i(T-\lambda_i)=T^N-1=0$. Let $e_i(X)=\prod_{i'\ne i}(X-\lambda_{i'})/(\lambda_i-\lambda_{i'})\in
\Lambda[X]$ (the denominators are units). The polynomial $\sum_ie_i(X)-1$ has degree $<N$ and vanishes at
the $N$ points $\lambda_{i'}$, whose differences are units, so it is $0$ (the Vandermonde determinant is a
unit). $(X-\lambda_i)e_i(X)$ is a unit multiple of $\prod_{i'}(X-\lambda_{i'})$, so $(T-\lambda_i)e_i(T)=0$,
i.e.\ $Te_i=\lambda_ie_i$. Hence $e_j(T)e_i(T)=e_j(\lambda_i)e_i(T)=\delta_{ij}e_i(T)$, and with
$\sum_ie_i=1$ the $e_i=e_i(T)$ are orthogonal idempotents summing to $1$. If $Tx=\lambda_ix$ then
$x=\sum_je_jx=\sum_je_j(\lambda_i)x=e_ix$; so the image of $e_i$ is the letter part at $i$. $\varphi$ fixes
$T$ and maps $\lambda_{i'}$ to $\lambda_{i'^\ell}$, permuting the letters, so $\varphi e_i(T)=
e_{i^\ell}(T)\varphi$; an element $g$ fixing $\Lambda$ commutes with $T$ (as $\zeta$ acts through the
$\auO$-module structure, which commutes with $G_{\Q_\ell}$) and with the coefficients of $e_i(X)$. For
$\zeta=1$, $T=1$ and $(1-\lambda_i)$ is a unit for $i\ne1$, so $\pkpAmb(L)_i=0$ for $i\ne1$; and the
residue of $\gamma(1)=1$ is the letter $1$.
\end{proof}

Only the properties listed in Lemma~\ref{pkw:letters} are used in \S\S\ref{pkp:sec}--\ref{pkd:sec} about
the letters (the exact order of $\zeta$ is never used), so Definition~\ref{pkp:def-letters} may be stated
with $\zeta^{\ell^f-1}=1$; with $\zeta=1$ the jump set at the letter $1$ is $\{0\}\cup\bigcup_\gamma
U(\gamma)$ and every other jump set is the set $\{0\}$ that sides are required to contain (the other letter
parts being $0$, any finite set containing $0$ is admissible there). This avoids requiring $\auO$ to contain
roots of unity of order $\ell^f-1$ for the residue degree $f$ of a tame chart.

\begin{proposition}[(T1): existence of tame charts]\label{pkw:T1}
Let $E_v=\Q_\ell$, $f\ge1$ with $\ell\nmid f$, and $e\ge1$ with $e\mid\ell^f-1$. There is a tame chart with
ramification index $e$ and residue degree $f$. In particular, if $1\le r<\ell$ and
$e'=\operatorname{lcm}_{1\le a\le r}(\ell^a-1)$, there is a tame chart with ramification index $e'$ (and
residue degree $r!$); its $d_v$ is at most $1$, and $e'\le\ell^{r^2}$ (Lemma~\ref{pkc:lcm}).
\end{proposition}

\begin{proof}
\emph{The field.} Let $N=\ell^f-1$, fix $\varpi\in\auQlb$ with $\varpi^e=\ell$, and let
$K'=\Q_\ell(\mu_N,\varpi)$. $\Q_\ell(\mu_N)=\Q_{\ell^f}$ is unramified of degree $f$ with residue field
$\mathbb F_{\ell^f}$, and $X^e-\ell$ is Eisenstein over its ring of integers, so $K'$ has ramification index
$e$, residue degree $f$, uniformiser $\varpi$, and $[K':\Q_\ell]=ef$, which is prime to $\ell$ ($e\mid N$
and $\ell\nmid fN$). $K'$ contains $\mu_e\subseteq\mu_N$, hence all roots $\zeta\varpi$ of $X^e-\ell$; so $K'$
is the splitting field of $(X^N-1)(X^e-\ell)$ and is Galois over $\Q_\ell$. Put $\Gamma'=j^{-1}(G_{K'})$: an
open normal subgroup of $\Gamma$ of index $ef$, prime to $\ell$.

\emph{Fixing $\Lambda$.} For $g\in G_{\Q_\ell}$: $g$ fixes $\Lambda=W(\mathbb F_{\ell^f})\subseteq\pkpA$
pointwise iff $g$ fixes $\mu_N\subseteq\auQlb$ pointwise. Indeed $\Lambda$ is the $\ell$-adic closure of the
$\Z_\ell$-span of the Teichm\"uller lifts $[x]$, $x\in\mathbb F_{\ell^f}\subseteq O^\flat$, and $g[x]=[gx]$;
for $x\ne0$, $x=(x_n\bmod\ell)_n$ where $x_n\in\mu_N$ is the unique $\ell^n$-th root of $x_0=\omega(x)$
in $\mu_N$, and $x_n$ is a power of $x_0$; distinct elements of $\mu_N$ have distinct residues; so $gx=x$
iff $g$ fixes every $x_n$ iff $g$ fixes $\omega(x)$; and $\omega$ maps $\mathbb F_{\ell^f}^\times$ onto
$\mu_N$. Hence the elements of $G_{\Q_\ell}$ fixing $\Lambda$ pointwise and $\varpi$ are exactly those of
$G_{K'}$.

\emph{The chart.} Choose $\varpi_n\in\auQlb$ with $\varpi_0=\varpi$, $\varpi_{n+1}^\ell=\varpi_n$, put
$\varpi^\flat=(\varpi_n\bmod\ell)_n\in\mathfrak m^\flat$, $z=[\varpi^\flat]$, so $\theta(z)=\varpi$, and
$E=\ell-z^e=z^eA+\ell$ with $A=-1\in\mathfrak S^\times$. Then $\theta(E)=\ell-\varpi^e=0$, and $E$
generates $\ker\theta$ by Corollary~\ref{pkn:E-gen} (whose proof uses only $E=z^eA+\ell$ with $A$ a unit,
$z=[\varpi^\flat]$ and $\theta(E)=0$). By the previous paragraph $j(\Gamma')=G_{K'}$ fixes $\Lambda$ and
$\theta(z)$, and every element of $G_{\Q_\ell}$ fixing both lies in $j(\Gamma')$. Kummer elements: fix
$l\ge0$ and $\zeta_{\ell^n}$ with $\varepsilon=(\zeta_{\ell^n}\bmod\ell)_n$; the sets
$X_n=\{g\in G_{K'}:g\varpi_n=\zeta_{\ell^n}^l\varpi_n\}$ are closed, decreasing (raise to the $\ell$-th power)
and nonempty ($X^{\ell^n}-\varpi$ is Eisenstein over $O_{K'}$, hence irreducible, and $G_{K'}$ is transitive
on its roots); by compactness some $g\in G_{K'}$ lies in all $X_n$, and then $g\varpi^\flat=
\varepsilon^l\varpi^\flat$, $gz=[\varepsilon]^lz$. This is a Kummer chart (with $K'$ as above), with
ramification index $e$ and residue degree $f$; $\Gamma'$ is normal and $\ell\nmid e$.

\emph{The tame generator.} As $e\mid N$, $\mathbb F_{\ell^f}^\times$ contains a primitive $e$-th root of
unity $a$; let $a_n\in\mu_N$ be the unique $\ell^n$-th root of $\omega(a)$ in $\mu_N$, so
$a=(a_n\bmod\ell)_n\in O^\flat$, and put $y_n=a_n\varpi_n$. Then $y_{n+1}^\ell=y_n$ and
$y_n^{e\ell^n}=a_0^e\varpi^e=\ell$. The polynomial $X^{e\ell^n}-\ell$ is Eisenstein over the discrete
valuation ring of $\pkwur$, hence irreducible over $\pkwur$, and $\varpi_n$, $y_n$ are two of its roots;
so the closed sets $Y_n=\{g\in\pkwI:g\varpi_n=y_n\}$ are nonempty, and decreasing (raise to the $\ell$-th
power). $\pkwI$ is compact, so some $s\in\pkwI$ lies in all $Y_n$. Then $s\varpi^\flat=(a_n\varpi_n\bmod\ell)_n
=a\varpi^\flat$ and $sz=[a][\varpi^\flat]=\xi z$ with $\xi=[a]$. Being inertial, $s$ fixes $\mu_N\subseteq
\pkwur$, hence $\Lambda$. Put $s_0=j^{-1}(s)$.

\emph{The index $e'$.} If $1\le r<\ell$, then $\ell\nmid r!$, and for $1\le a\le r$, $a\mid r!$ gives
$\ell^a-1\mid\ell^{r!}-1$; so $e'\mid\ell^{r!}-1$, and the first part applies with $f=r!$, $e=e'$.
\end{proof}

\subsection{(T2): descent to crystalline over $\Gamma'$}\label{pkw:ss-T2}

\begin{lemma}[Tame inertia]\label{pkw:tame-inertia}
Let $e\ge1$ with $\ell\nmid e$, $\varpi\in\auQlb$ with $\varpi^e=\ell$, and let $\phi\colon\pkwI\to G$ be a
homomorphism to a group, with open kernel, such that $\phi(s)^e=1$ for all $s\in\pkwI$. Then $\phi(s)=1$ for
every $s\in\pkwI$ with $s\varpi=\varpi$.
\end{lemma}

\begin{proof}
$N=\ker\phi$ is open and normal in $\pkwI=\operatorname{Gal}(\auQlb/\pkwur)$; its fixed field $M$ is a finite
Galois extension of $\pkwur$ with group $\pkwI/N\cong\phi(\pkwI)$, in which every element has order dividing
$e$. Hence $m=[M:\pkwur]$ is prime to $\ell$ (otherwise, by Cauchy's theorem, some element would have order
$\ell\nmid e$). By (F3), $M=\pkwur(\varpi_m)$ with $\varpi_m^m=\ell$. Since $\mu_m\subseteq\pkwur$ ($\ell\nmid m$),
Kummer theory shows that $\operatorname{Gal}(M/\pkwur)$ is cyclic of order $m$; as its elements have order
dividing $e$, $m\mid e$. Both $\varpi_m$ and $\varpi^{e/m}$ are $m$-th roots of $\ell$, so
$\varpi_m=\zeta\varpi^{e/m}$ with $\zeta\in\mu_m\subseteq\pkwur$. Hence $M\subseteq\pkwur(\varpi)$, and an
$s\in\pkwI$ fixing $\varpi$ fixes $M$, i.e.\ $s\in N$.
\end{proof}

\begin{theorem}[(T2): descent]\label{pkw:T2}
Fix a tame chart (Definition~\ref{pkw:def-tame}) with ramification index $e$. Let $Y$ be a potentially
crystalline lattice such that $X^e-1$ kills every inertial $s\in\Gamma$ on $\pkwD(Y)$. Then $Y$ is
crystalline on $\Gamma'$. Its Hodge--Tate jumps are unchanged (they are those of $Y[1/\ell]$, which do not
depend on the open subgroup on which they are computed), so if $Y$ has jumps in $U$ it is a lattice with
jump sets over the chart (Definition~\ref{pkp:def-orders}).
\end{theorem}

\begin{proof}
As $j$ is bijective, $s\mapsto\beta_s$ is defined on $j^{-1}(\pkwI)\cong\pkwI$, and by
Lemma~\ref{pkw:beta}(i) it is a homomorphism $\phi\colon\pkwI\to\operatorname{Aut}_{\pkwB}(\pkwB(Y))$ whose
kernel contains $H\cap\pkwI$ for an open $H$ with $\pkwB(Y)^H$ spanning; so $\ker\phi$ is open. By
hypothesis $\phi(s)^e=1$. By Lemma~\ref{pkw:tame-inertia}, $\beta_s=1$ for every inertial $s$ with
$j(s)\varpi=\varpi$, in particular (Definition~\ref{pkw:def-tame}: $j(\Gamma')$ fixes $\theta(z)=\varpi$) for
every inertial $s\in\Gamma'$. So $j(\Gamma')\cap\pkwI=I_{K'}$ acts trivially on $\pkwD(Y)$ (which is the union
in (F2) for $H=j(\Gamma')$, as the open subgroups of $\Gamma'$ are cofinal among those of $\Gamma$), and by (F2)
$Y[1/\ell]$ is crystalline as a representation of $\Gamma'$.
\end{proof}

\subsection{(T3): the tame cokernel killer}\label{pkw:ss-T3}

Fix a tame chart with ramification index $e$, residue degree $f$, tame generator $s_0$ and $\xi=[a]$ as in
Definition~\ref{pkw:def-tame}, and letters with any $\zeta$, $\zeta^{\ell^f-1}=1$ (Lemma~\ref{pkw:letters}).
For a side $S$, $\pkmI li$ and $\pkmJ rli$ are the lift ideals of Definition~\ref{pkm:def-J}, and
$\varepsilon_i(c)=e_i(1\otimes c)$. For a lattice with jump sets $X$ we write
$\mathcal N_i(\mathrm{std}_0)=\mathcal N(X)\cap X_i(\mathrm{std}_0)$ (Lemma~\ref{pkm:inputs}).

The argument rests on the following statement (Z), proved in Theorem~\ref{pkw:Z} below.

\begin{enumerate}
\item[(Z)] \emph{Division by $z$.} Let $X$ be a lattice with jump sets over the tame chart, $i$ a letter and
$Q=\sum_kQ_kX^k\in\Z_\ell[X]$ with $Q(\beta_{s_0})=0$ on $\pkwB(X)$. Then for every
$x\in\mathcal N_i(\mathrm{std}_0)$ the vector $z^{-1}\sum_kQ_ks_0^kx$ ($z$ is a unit of $\pkpW$) lies in
$X_i(\mathrm{std}_0)$.
\end{enumerate}

We use one more published fact.
\begin{enumerate}
\item[(F6)] (Colmez \cite[\S III.2]{pkfColmez98}, Fontaine \cite[5.1.2]{pkfFontaineIII}.) $\pkwB\subseteq\pkfBm$,
compatibly with $G_{\Q_\ell}$, with $\varphi$ and with the maps to $\pknB$; and $\pkfBm^H=K_{H,0}$ for every open
$H\subseteq G_{\Q_\ell}$. (Colmez gives $\varphi(\pkfBm)\subseteq\pkwB\subseteq\pkfBm$; with $\pkwB^H=K_{H,0}$ (F1) and injectivity of $\varphi$ this gives $\pkfBm^H=K_{H,0}$, as in (P2) of \S\ref{pkf:ss-rings}.)
\end{enumerate}

\begin{lemma}[$s_0$ on $\pknMinf$ and on $\mathfrak S_B$]\label{pkw:s0}
Let $X$ be a lattice with jump sets over the tame chart and $k\ge0$.
\begin{enumerate}
\item[(i)] For a form $\mu$ (Definition~\ref{pkn:def-N}), $\mu_k(w)=s_0^{-k}\mu(s_0^kw)$ is a form; if $\mu$ maps
$1\otimes X$ into (the image of) $\pkfBm$, so does $\mu_k$.
\item[(ii)] $s_0^k$ maps $\pknMinf(X)$ and each letter part $\pkpAmb(X)_i$ into themselves.
\item[(iii)] $s_0^kE=E$.
\item[(iv)] For $b\in\pkfSB$, $z$ divides $s_0^kb-b$ in $\pkfBp$.
\end{enumerate}
\end{lemma}

\begin{proof}
(i) $\mu_k$ is $\pkpA$-linear ($s_0$ acts semilinearly on $\pkpAmb(X)$ and on $\pknB$). For $h\in\Gamma'$ put
$h'=s_0^khs_0^{-k}\in\Gamma'$ ($\Gamma'$ is normal); then $\mu_k(1\otimes hy)=s_0^{-k}\mu(1\otimes h's_0^ky)=
s_0^{-k}h'\mu(1\otimes s_0^ky)=h\,\mu_k(1\otimes y)$. $s_0$ preserves $\pkfBm$ (P1).
(ii) As in the proof of Lemma~\ref{pkl:G1}: $s_0$ preserves $\pkpA\otimes X$, $s_0\tau=\tau r'$ with $r'\in\pkpA^\times$
(Lemma~\ref{pkf:approx0}(iv)), and $\mu(\varphi^ms_0^kx)=s_0^k\mu_k(\varphi^mx)\in\pknBp$ by (i). The letter parts:
Lemma~\ref{pkw:letters}.
(iii) $E=\ell-z^e$ (Definition~\ref{pkw:def-tame}; the chart of Proposition~\ref{pkw:T1} has it), so
$s_0^kE=\ell-\xi^{ke}z^e=E$ as $\xi^e=[a^e]=1$.
(iv) By Lemma~\ref{pkl:SBdec}, $\ell^{k'}b=\lambda+zw'$ with $\lambda\in\Lambda$, $w'\in\pkfSB\subseteq\pkfBp$. $s_0$ fixes
$\lambda$ and $s_0^kz=\xi^kz$, and $s_0^k$ preserves $\pkfBp$ (P1); so $\ell^{k'}(s_0^kb-b)=z(\xi^ks_0^k(w')-w')$, and $\ell$
is a unit of $\pkfBp$.
\end{proof}

\begin{lemma}[The relation among conjugate forms]\label{pkw:rel}
Let $X$ be a lattice with jump sets over the tame chart, $Q=\sum_kQ_kX^k\in\Z_\ell[X]$ with $Q(\beta_{s_0})=0$ on
$\pkwB(X)$, and $\mu$ a form mapping $1\otimes X$ into $\pkfBm$. Then $\sum_kQ_k\mu_k(x)=0$ for every $x\in\pknMinf(X)$.
\end{lemma}

\begin{proof}
Let $l\colon X\to\pkfBm$, $l(y)=\mu(1\otimes y)$; it is $\Z_\ell$-linear and $l(hy)=hl(y)$ for $h\in\Gamma'$ (the map
$\pkfBm\to\pknB$ is injective and equivariant, (P1)). Let $\lambda\colon\pkfBm\otimes_{\Z_\ell}X\to\pkfBm$,
$b\otimes y\mapsto bl(y)$, and $\lambda_k(b\otimes y)=b\,s_0^{-k}l(s_0^ky)=s_0^{-k}\lambda(s_0^k(b\otimes y))$; these are
$\pkfBm$-linear, $\lambda$ is $\Gamma'$-equivariant, and $\lambda_k(1\otimes y)=\mu_k(1\otimes y)$. Regard
$\pkwB(X)\subseteq\pkfBm\otimes X$ (F6). Let $d\in\pkwD(X)$. Then $s_0^kd=\beta^k_{s_0}d\in\pkwD(X)$ is fixed by an open
$H'\subseteq\Gamma'$, so $\lambda(s_0^kd)\in\pkfBm^{H'}=K_{H',0}\subseteq\pkwur$ (F6), which the inertial $s_0^{-k}$ fixes;
hence $\lambda_k(d)=\lambda(\beta_{s_0}^kd)$ and $\sum_kQ_k\lambda_k(d)=\lambda(Q(\beta_{s_0})d)=0$. As $\pkwD(X)$ spans
$\pkwB(X)$ over $\pkwB$, hence $\pkfBm\otimes X$ over $\pkfBm$, the $\pkfBm$-linear map $\sum_kQ_k\lambda_k$ vanishes; in
particular $\sum_kQ_k\mu_k$ vanishes on $1\otimes X$, hence on $\pkpA\otimes X$. If $x\in\pknMinf$ then $\tau^nx\in
\pkpA\otimes X$ for some $n$, and $\tau^n\sum_kQ_k\mu_k(x)=0$ with $\tau\ne0$ in the field $\pknB$.
\end{proof}

\begin{theorem}[(Z)]\label{pkw:Z}
Statement (Z) holds.
\end{theorem}

\begin{proof}
Let $x\in\mathcal N_i(\mathrm{std}_0)$ and $w=\sum_kQ_ks_0^kx$. By Corollary~\ref{pkl:F}, $X_i(\mathrm{std}_0)=
\pknMinf(X)\cap\pkpAmb(X)_i$; so $x\in\mathcal N\cap\pknMinf$, and $w\in\pknMinf\cap\pkpAmb(X)_i$ by
Lemma~\ref{pkw:s0}(ii). Fix a frame of $X$ for its jump bound $a$ (Theorem~\ref{pkf:main}), with vectors $n_j$ and
matrix $Z$ over $\pkfSB$, and let $\sigma_\nu$ be the forms of Lemma~\ref{pkn:frame-id}(iii) for $m=0$; by their
construction $\sigma_\nu(1\otimes e_j)=(M^{-1})_{\nu j}\in\pkfBm$, so $\sigma_\nu$ maps $1\otimes X$ into $\pkfBm$.

\emph{Values at $x$.} As in the proof of Corollary~\ref{pkl:Phi}, $x=\sum_jt_jn_j$ with $t_j\in\mathfrak S$. Let $\mu$ be a
form mapping $1\otimes X$ into $\pkfBm$. By Lemma~\ref{pkn:frame-id}(ii) ($m=0$), $E^a\mu(x)=c^{(0)}(\mu)Zt$, and the
entries of $c^{(0)}(\mu)=\mu(\mathbf e)M$ lie in $\pkfBm$ and are fixed by $j(\Gamma')$ (proof of
Lemma~\ref{pkn:frame-id}(i)), hence lie in $\pkfBm^{G_{K'}}=K'_0=\Lambda[1/\ell]\subseteq\pkfSB$ (F6, and
Lemma~\ref{pkf:SB-basic}(i)). So $E^a\mu(x)\in\pkfSB$.

\emph{Values at $w$.} Fix $\nu$, put $\mu=\sigma_\nu$, $b_k=\mu_k(x)$ and $X_k=E^ab_k\in\pkfSB$ (the previous paragraph
applied to the form $\mu_k$, Lemma~\ref{pkw:s0}(i)). Then $\mu(s_0^kx)=s_0^kb_k$, and by Lemma~\ref{pkw:rel}
$\sum_kQ_kb_k=0$, so $\mu(w)=\sum_kQ_k(s_0^kb_k-b_k)$. By Lemma~\ref{pkw:s0}(iii), $E^a s_0^kb_k=s_0^kX_k$, hence
\[
E^{2a}\sigma_\nu(w)=\sum_kQ_kE^a\bigl(s_0^kX_k-X_k\bigr)\in W:=\text{image of }z\pkfBp\text{ in }\pknB
\]
by Lemma~\ref{pkw:s0}(iv); $W$ is a $\pkfBp$-submodule.

\emph{Division.} By Proposition~\ref{pkn:frame}(ii), $w=\sum_js_jn_j$ with $s_j\in\pkpA$. The proof of the coefficient test
Lemma~\ref{pkn:frame-id}(iv) uses its hypothesis only for the forms $\sigma_\nu$; applied to $E^{2a}w=\sum_jE^{2a}s_jn_j$ it
gives $E^{4a}s_j\in W$, i.e.\ (as $\pkfBp\to\pknB$ is injective, (P1)) $z$ divides $E^{4a}s_j$ in $\pkfBp$. By
Lemma~\ref{pkl:intdiv}, $z\mid E^{4a}s_j$ in $\pkpA$, and then $z\mid s_j$ (as in the proof of Lemma~\ref{pkl:cancel}:
$E\nmid z$ and $E\pkpA$ is prime). So $z^{-1}w=\sum_j(s_j/z)n_j\in\pknMinf$, and it lies in $\pkpAmb(X)_i$ since $w$ does
and $z$ commutes with $e_i$. By Corollary~\ref{pkl:F}, $z^{-1}w\in X_i(\mathrm{std}_0)$.
\end{proof}

\begin{lemma}[Counting bad steps]\label{pkw:count}
Let $Q\in\Z_\ell[X]$ be monic and $\nu<e$. Call $k$ \emph{bad} if the residue of $Q(\xi^k)\in\Lambda$ in
$\mathbb F_{\ell^f}$ is $0$. Then at most $\deg Q$ of the integers $1\le k\le\nu$ are bad, and for $k$ not bad
$Q(\xi^k)$ is a unit of $\Lambda$, hence of $\pkpA$.
\end{lemma}

\begin{proof}
The residue of $Q(\xi^k)$ is $\bar Q(a^k)$, $\bar Q\in\mathbb F_\ell[X]$ the reduction, monic of degree $\deg Q$.
The $a^k$, $1\le k\le\nu<e$, are pairwise distinct ($a$ has order $e$), and each bad $k$ gives a root $a^k$ of
$\bar Q\ne0$. An element of $\Lambda=W(\mathbb F_{\ell^f})$ with nonzero residue is a unit.
\end{proof}

\begin{lemma}[The tame step]\label{pkw:step}
Let $S$ be a side whose lattice $Y$ is potentially crystalline, $Q\in\Z_\ell[X]$ with $Q(1)=0$
killing $s_0$ on $\pkwD(X_S)$, $i$ a letter and $r,t,k\ge0$. If $\ell^tz^{k+1}\in\pkmJ r0i$, then
$Q(\xi^{k+1})\,\ell^tz^k\in\pkmJ r0i$.
\end{lemma}

\begin{proof}
Put $f_0=\ell^tz^{k+1}\in\mathfrak S\cap\pkmJ r0i$. By Lemma~\ref{pkm:J-basic}(iv) there is $u\in\mathfrak S$
such that for every $c\in C_S$ some $x_c\in\mathcal N_i(\mathrm{std}_0)$ (for $X_S$) has
$\mathrm{pr}(x_c)=(f_0-ru)\varepsilon_i(c)$. Write $u=w+zv$ with $w\in\Lambda$ and $v\in\mathfrak S$
($\mathfrak S=\Lambda[[z]]$). Fix $c$. By (Z), $y_c=z^{-1}\sum_jQ_js_0^jx_c\in(X_S)_i(\mathrm{std}_0)$.

$\mathrm{pr}$ is induced by the equivariant projection $X_S\to C_S$, so it commutes with $s_0$; and
$s_0\varepsilon_i(c)=\varepsilon_i(c)$, because $\Gamma$ acts trivially on $C_S$ and $s_0$ fixes $\Lambda$, hence
commutes with $e_i$ (Lemma~\ref{pkw:letters}). Since $s_0^jz=\xi^jz$ and $s_0$ fixes $\ell$, $\xi$, $w$,
\[
\mathrm{pr}(s_0^jx_c)=s_0^j(f_0-ru)\,\varepsilon_i(c)=a_j\,\varepsilon_i(c),\qquad
a_j=\ell^t\xi^{j(k+1)}z^{k+1}-r\bigl(w+\xi^jz\,s_0^j(v)\bigr),
\]
with $s_0^j(v)\in\pkpA$. Using $\sum_jQ_j=Q(1)=0$,
\[
\sum_jQ_ja_j=z\,b,\qquad b=Q(\xi^{k+1})\,\ell^tz^k-r\sum_jQ_j\xi^js_0^j(v)\in\pkpA .
\]
Hence $z\,\mathrm{pr}(y_c)=\mathrm{pr}(\sum_jQ_js_0^jx_c)=zb\,\varepsilon_i(c)$, and as $z$ is a unit of $\pkpW$
and $\mathrm{pr}$ is $\pkpW$-linear, $\mathrm{pr}(y_c)=b\,\varepsilon_i(c)$. As $b$ does not depend on $c$, this
says $b\in\pkmI0i$, and $Q(\xi^{k+1})\ell^tz^k=b+r\sum_jQ_j\xi^js_0^j(v)\in\pkmJ r0i$.
\end{proof}

\begin{theorem}[(T3): the tame killer]\label{pkw:T3}
Fix a tame chart with ramification index $e$. Let $S$ be a side whose jumps at all letters are
$\le n_0$, whose lattice $Y$ is potentially crystalline, and let a monic $q\in\Z_\ell[X]$ kill $s_0$ on
$\pkwD(Y)$. If $e\le\ell^{m_0}$, $m_0(\deg q+1)\le\kappa$ and $n_0+\kappa+2<\ell$, then $\ell^\kappa=\ell^\kappa E^0$
kills the cokernel of $S$ at every letter.
\end{theorem}

\begin{proof}
Fix a letter $i$ and write $J=\pkmJ{\ell^{\kappa+1}}0i$. By Lemma~\ref{pkm:zpow} (with $\omega=n_0$) there is
$\nu\ge0$ with $(\ell-1)\nu\le e(n_0+\kappa)$ and $z^\nu\in J$; as $n_0+\kappa<\ell-1$, $\nu<e$. Let
$Q=q(X)(X-1)$, monic of degree $\deg q+1$ with $Q(1)=0$; it kills $s_0$ on $\pkwD(X_S)$ by
Lemma~\ref{pkw:ext}. Call $k$ bad as in Lemma~\ref{pkw:count} and let $b(k)$ be the number of bad
$j\in[1,k]$. We show, by induction on $k$:

\emph{if $k<e$, $t\ge0$, $m_0b(k)+t\le\kappa$ and $\ell^tz^k\in J$, then $\ell^\kappa\in\pkmI0i$.}

$k=0$: $t\le\kappa$, so $\ell^\kappa\in J=\pkmJ{\ell^{\kappa+1}}0i$, and $\ell^\kappa\in\pkmI0i$ by
Lemma~\ref{pkm:J-basic}(i). Step $k\to k+1$ (with $k+1<e$, $\ell^tz^{k+1}\in J$, $m_0b(k+1)+t\le\kappa$). If $k+1$
is bad, Lemma~\ref{pkm:lower}(i) gives $\ell^{t+v_\ell(k+1)+1}z^k\in J$; as $k+1<e\le\ell^{m_0}$,
$v_\ell(k+1)+1\le m_0$, so $\ell^{t+m_0}z^k\in J$, and $m_0b(k)+(t+m_0)=m_0b(k+1)+t\le\kappa$: apply the induction
hypothesis with $t+m_0$. If $k+1$ is not bad, Lemma~\ref{pkw:step} gives $Q(\xi^{k+1})\ell^tz^k\in J$, and
$Q(\xi^{k+1})$ is a unit of $\pkpA$ (Lemma~\ref{pkw:count}), so $\ell^tz^k\in J$ with $b(k)=b(k+1)$: apply the
induction hypothesis with $t$.

Apply this with $k=\nu$, $t=0$: by Lemma~\ref{pkw:count}, $b(\nu)\le\deg Q=\deg q+1$, so
$m_0b(\nu)\le m_0(\deg q+1)\le\kappa$. Hence $\ell^\kappa\in\pkmI0i$, i.e.\ $\ell^\kappa$ kills the cokernel at $i$
(Definition~\ref{pkm:def-J}).
\end{proof}



\subsection{(I6) and (T4): Tate duals and tensor products}\label{pkw:ss-T4}

For a lattice $Y$, $Y^*=\operatorname{Hom}_\auO(Y,\auO(1))$ with $(g\phi)(y)=g\,\phi(g^{-1}y)$ and
$(a\phi)(y)=a\phi(y)=\phi(ay)$. Fix, as in {\cite[Definition~5.1(4)]{P55}}, a $\Z_\ell$-linear $\mathrm f\colon
\auO\to\Z_\ell$ for which $(a,b)\mapsto\mathrm f(ab)$ is a perfect pairing, and write also $\mathrm f$ for
$\mathrm f\otimes1\colon\auO(1)\to\Z_\ell(1)$, which is equivariant. Then
\[
Y\times Y^*\to\Z_\ell(1),\qquad (y,\phi)\mapsto\mathrm f(\phi(y)),
\]
is $\Z_\ell$-bilinear, $\Gamma$-equivariant, satisfies $\mathrm f((a\phi)(y))=\mathrm f(\phi(ay))$, and is perfect:
$Y$ is free over $\auO$, and for $Y=\auO$ the map $\auO\to\operatorname{Hom}_{\Z_\ell}(\auO,\Z_\ell)$, $b\mapsto
\mathrm f(b\,\cdot)$, is bijective by the choice of $\mathrm f$. Its $\pkwB$-bilinear extension
$\langle\ ,\ \rangle\colon\pkwB(Y)\times\pkwB(Y^*)\to\pkwB\otimes_{\Z_\ell}\Z_\ell(1)$ is therefore a perfect
pairing of free $\pkwB$-modules, $\Gamma$-equivariant: $\langle gx,g\phi\rangle=g\langle x,\phi\rangle$.

\begin{lemma}[(I6): inertia on duals]\label{pkw:dual}
Let $Y$ be a potentially crystalline lattice. Then $Y^*$ is potentially crystalline, and for every
inertial $s$ and $x\in\pkwB(Y)$, $\phi\in\pkwB(Y^*)$: $\langle\beta^Y_sx,\beta^{Y^*}_s\phi\rangle=\langle
x,\phi\rangle$. Consequently, for $P\in\auO[X]$: $P$ kills $s$ on $\pkwD(Y^*)$ if and only if $P$ kills $s^{-1}$ on
$\pkwD(Y)$.
\end{lemma}

\begin{proof}
Let $H\subseteq\Gamma$ be open with $D_H=\pkwB(Y)^H$ spanning, $F=K_{H,0}$. By (F1), $\pkwB\otimes_FD_H\cong
\pkwB(Y)$, so an $F$-basis $d_1,\dots,d_n$ of $D_H$ is a $\pkwB$-basis of $\pkwB(Y)$. Let $e_1$ be a basis of
$\Z_\ell(1)$ and $\epsilon=t^{-1}\otimes e_1$, a $\pkwB$-basis of $\pkwB\otimes\Z_\ell(1)$ fixed by $G_{\Q_\ell}$
(Lemma~\ref{pkw:beta}(iii)). By perfectness there are $d^1,\dots,d^n\in\pkwB(Y^*)$, a $\pkwB$-basis, with
$\langle d_i,d^k\rangle=\delta_{ik}\epsilon$. For $h\in H$, $\langle d_i,hd^k\rangle=h\langle h^{-1}d_i,d^k\rangle=
h(\delta_{ik}\epsilon)=\delta_{ik}\epsilon$, so $hd^k=d^k$: the $H$-invariants span $\pkwB(Y^*)$, and $Y^*$ is
potentially crystalline. Shrinking $H$, we may take it normal, and then $d_i\in\pkwD(Y)$, $d^k\in\pkwD(Y^*)$,
so $\langle\beta_sd_i,\beta_sd^k\rangle=\langle sd_i,sd^k\rangle=s(\delta_{ik}\epsilon)=\delta_{ik}\epsilon$; by
bilinearity $\langle\beta_sx,\beta_s\phi\rangle=\langle x,\phi\rangle$ for all $x$, $\phi$.

Hence $\langle x,\beta_s\phi\rangle=\langle\beta_s\beta_{s^{-1}}x,\beta_s\phi\rangle=\langle\beta_{s^{-1}}x,\phi\rangle$
(Lemma~\ref{pkw:beta}(i)), and $\langle x,a\phi\rangle=\langle ax,\phi\rangle$ for $a\in\auO$. So for
$P=\sum_kP_kX^k$, $\langle x,P(\beta^{Y^*}_s)\phi\rangle=\langle P(\beta^Y_{s^{-1}})x,\phi\rangle$ for all $x,\phi$.
By perfectness, $P(\beta^{Y^*}_s)=0$ if and only if $P(\beta^Y_{s^{-1}})=0$.
\end{proof}

This is input (I6) of \S\ref{pkd:ss-data} (with the meaning of ``kills'' of
Definition~\ref{pkw:def-pcris}). The jumps of $Y^*$ are treated in Proposition~\ref{pkw:HT}.

\begin{lemma}[Tensor products]\label{pkw:tensor}
Let $Y$, $L$ be potentially crystalline lattices and $s$ inertial. Then $Y\otimes_\auO L$ is potentially
crystalline, and the natural surjection $\pi\colon\pkwB(Y)\otimes_{\pkwB}\pkwB(L)\to\pkwB(Y\otimes_\auO L)$
satisfies $\pi\circ(\beta^Y_s\otimes\beta^L_s)=\beta^{Y\otimes L}_s\circ\pi$.
\end{lemma}

\begin{proof}
$\pkwB(Y)\otimes_{\pkwB}\pkwB(L)=\pkwB\otimes_{\Z_\ell}(Y\otimes_{\Z_\ell}L)$, and $\pi$ is induced by the
equivariant surjection $Y\otimes_{\Z_\ell}L\to Y\otimes_\auO L$, so it is $\pkwB$-linear, surjective and
equivariant. If $d\in\pkwD(Y)$ and $d'\in\pkwD(L)$ are fixed by open $H$, $H'$, then $\pi(d\otimes d')$ is fixed
by $H\cap H'$; the $d\otimes d'$ span the source, so the $\pi(d\otimes d')$, which lie in a single
$\pkwB(Y\otimes_\auO L)^{H''}$ for finitely many spanning ones, span the target: $Y\otimes_\auO L$ is potentially
crystalline. Both sides of the identity are $\pkwB$-linear and send $d\otimes d'$ to
$\pi(sd\otimes sd')=s\,\pi(d\otimes d')=\beta^{Y\otimes L}_s\pi(d\otimes d')$, as $\pi(d\otimes d')\in
\pkwD(Y\otimes_\auO L)$.
\end{proof}

\begin{lemma}[A Cayley--Hamilton polynomial]\label{pkw:CH}
Let $q,q'\in\Z_\ell[X]$ be monic, $M$ a module over a commutative ring containing $\Z_\ell$, and $A,A'$ commuting
endomorphisms of $M$ with $q(A)=0$, $q'(A')=0$. Then the characteristic polynomial $q''$ of multiplication by
$x\otimes x'$ on $R=\Z_\ell[x]/(q)\otimes_{\Z_\ell}\Z_\ell[x']/(q')$ is monic of degree $\deg q\cdot\deg q'$ and
satisfies $q''(AA')=0$.
\end{lemma}

\begin{proof}
$R$ is free over $\Z_\ell$ with basis $x^a\otimes x'^b$ ($a<\deg q$, $b<\deg q'$), so $q''$ is monic of degree
$\deg q\deg q'$. By Cayley--Hamilton $q''(m)=0$ for the multiplication $m$ by $\rho=x\otimes x'$; as $q''(m)$ is
multiplication by $q''(\rho)$, applying it to $1$ gives $q''(\rho)=0$ in $R$. Since $A$, $A'$ commute and
$q(A)=q'(A')=0$, there is a ring homomorphism $R\to\operatorname{End}(M)$ with $x\mapsto A$, $x'\mapsto A'$ (the
tensor product of two commutative $\Z_\ell$-algebras mapping to commuting subrings); it maps $\rho$ to $AA'$ and
$q''(\rho)$ to $q''(AA')$.
\end{proof}

\begin{theorem}[(T4)]\label{pkw:T4}
Let $Y$, $L$ be potentially crystalline lattices, $s$ inertial, and let monic $q,q'\in\Z_\ell[X]$ kill $s^{-1}$ on
$\pkwD(Y)$ and on $\pkwD(L)$ respectively. Then some monic $q''\in\Z_\ell[X]$ of degree $\deg q\cdot\deg q'$ kills
$s$ on $\pkwD(\operatorname{Hom}_\auO(Y\otimes_\auO L,\auO(1)))$.
\end{theorem}

\begin{proof}
On $M=\pkwB(Y)\otimes_{\pkwB}\pkwB(L)$ the endomorphisms $A=\beta^Y_{s^{-1}}\otimes1$ and $A'=1\otimes
\beta^L_{s^{-1}}$ commute and satisfy $q(A)=0$, $q'(A')=0$. Lemma~\ref{pkw:CH} gives $q''$ monic of degree
$\deg q\deg q'$ with $q''(AA')=0$, and by Lemma~\ref{pkw:tensor}
$q''(\beta^{Y\otimes L}_{s^{-1}})\pi=\pi\,q''(AA')=0$ with $\pi$ surjective; so $q''$ kills $s^{-1}$ on
$\pkwD(Y\otimes_\auO L)$, and by Lemma~\ref{pkw:dual} it kills $s$ on $\pkwD((Y\otimes_\auO L)^*)$.
\end{proof}

\begin{remark}[Use of (T4) in Theorem~\ref{pkc:tame}]\label{pkw:rem-T4}
Theorem~\ref{pkw:T4} is applied in the proof of Theorem~\ref{pkc:tame} with $s=s_0$, $q=q'_0$, $q'=q_0$: $q'_0$
kills $s_0^{-1}$ on $\pkwD(Y)$ (as $Y\in\pkiC(U,K)$) and $q_0$ kills $s_0^{-1}$ on $\pkwD(L)$ (as
$L\in\pkiC(U',K^{-1})$ and $K^{-1}(s_0^{-1})=K(s_0)$), so the resulting $q''$ has degree
$\deg q_0\deg q'_0\le N^2$. The hypotheses concern $s_0^{-1}$, not $s_0$: by Lemma~\ref{pkw:dual} the dual is
controlled by the polynomials killing $s_0^{-1}$.
\end{remark}

\subsection{Jumps of Tate duals; uniform tame polynomials}\label{pkw:ss-checks}

\subsubsection*{(a) The jumps of $L^*=\operatorname{Hom}_\auO(L,\auO(1))$}

We use two more published facts. Let $K'$ be a finite extension of $\Q_\ell$ and $V$ a de Rham representation of
$G_{K'}$ on a finite free module over $E_\auO=\auO[1/\ell]$.
\begin{enumerate}
\item[(F4)] (Tate \cite[\S3, Theorems~1 and~2, p.~176]{pklTate67}; also quoted in Fontaine \cite[1.5.7]{pkfFontaineII}; Fontaine \cite[3.8]{pkfFontaineIII}.) $C(a)^{G_{K'}}=0$ for $a\ne0$ and
$C^{G_{K'}}=K'$; $V$ is Hodge--Tate: for every embedding $\gamma\colon E_\auO\to C$ over $\Q_\ell$ there are integers
$m_\gamma(u)\ge0$, determined by $V$, with $C\otimes_{E_\auO,\gamma}V\cong\bigoplus_uC(-u)^{m_\gamma(u)}$
$G_{K'}$-equivariantly.
\item[(F5)] (Fontaine \cite[3.7--3.8]{pkfFontaineIII}.) $\mathrm{gr}^kD_{\mathrm{dR}}(V)\cong
(C(k)\otimes_{\Q_\ell}V)^{G_{K'}}$, compatibly with the action of $E_\auO$.
\end{enumerate}
By (F4)--(F5), $u$ is a jump of $V$ at $\gamma$ in the sense of $\mathrm{gr}^uD_{\mathrm{dR}}\ne0$ (the convention
fixed in (P3) of \S\ref{pkf:ss-rings}) exactly when $C(-u)$ occurs in $C\otimes_{E_\auO,\gamma}V$, with multiplicity
$m_\gamma(u)$; this does not depend on the open subgroup $G_{K'}$ on which it is computed.

\begin{proposition}[Jumps of Tate duals; consistency of conventions]\label{pkw:HT}
\begin{enumerate}
\item[(i)] $\auO(n)=\auO\otimes\Z_\ell(n)$ has the single jump $-n$ at every $\gamma$. In particular $\auO(1)$ has the
jump $-1$, as in Lemma~\ref{pkp:tau} and Definition~\ref{pkp:def-pairing}, and the character $\chi^e=
\chi_{\mathrm{cyc}}^{-e}$ of {\cite[\S5]{P55}} has the single jump $e$.
\item[(ii)] If $L$ is potentially crystalline, the jumps of $L^*$ at $\gamma$ are the $-u-1$, $u$ running over the
jumps of $L$ at $\gamma$, with the same multiplicities.
\item[(iii)] With the forms of Definition~\ref{pkn:def-N}: on $\auO(1)$ every form $\mu$ satisfies
$\mu(1\otimes\auO(1))\subseteq tK'\subseteq t^{-u}\pknBp$ with $u=-1$; so the jump $-1$ of $\auO(1)$ is also its type
in the sense of Definition~\ref{pkp:def-orders}.
\end{enumerate}
\end{proposition}

\begin{proof}
(i) $C\otimes_{E_\auO,\gamma}E_\auO(n)=C(n)=C(-(-n))$. For $\chi^e$: $C(-e)$, jump $e$.

(ii) $L^*[1/\ell]=\operatorname{Hom}_{E_\auO}(V,E_\auO(1))$ with $V=L[1/\ell]$ free over $E_\auO$; it is de Rham
(potentially crystalline by Lemma~\ref{pkw:dual}). Base change along $\gamma$ commutes with $\operatorname{Hom}$ of free
modules, so $C\otimes_{E_\auO,\gamma}L^*[1/\ell]\cong\operatorname{Hom}_C(C\otimes_{E_\auO,\gamma}V,C(1))\cong
\bigoplus_u\operatorname{Hom}_C(C(-u),C(1))^{m_\gamma(u)}=\bigoplus_uC(u+1)^{m_\gamma(u)}$, equivariantly; and
$C(u+1)=C(-(-u-1))$.

(iii) For $y=1\otimes a\otimes e_1$ and $g\in\Gamma'$, $\mu(1\otimes gy)=g\mu(y)$ reads
$\chi_{\mathrm{cyc}}(g)\mu(y)=g\mu(y)$, so $t^{-1}\mu(y)\in\pknB^{G_{K'}}=K'$ (Lemma~\ref{pkn:frame-id}(i), proof).
\end{proof}

\begin{remark}
The convention ``the jumps of $L^*$ are $-u-1$'' used in
Definition~\ref{pkp:def-pairing}, in Theorem~\ref{pkd:data} (i)--(iii) and in the deduction of the window
statements, agrees with the definition of jumps by $\mathrm{gr}^uD_{\mathrm{dR}}\ne0$ used in \S\ref{pkf:sec}
(property (P3)), with the jump $-1$ of $\auO(1)$ in Lemma~\ref{pkp:tau}, with the types of forms in
Definition~\ref{pkp:def-orders}, and with the statement of {\cite[\S5]{P55}} that $\chi^e=\chi_{\mathrm{cyc}}^{-e}$ has
the jump $e$. It is also the convention under which $\dim D_{\mathrm{dR}}/\mathrm{Fil}^0$ is the number of
negative jumps (as used in {\cite[\S2]{P55}}).
\end{remark}

\subsubsection*{(b) One tame polynomial for each inertial element}

\begin{proposition}[Uniform tame polynomials]\label{pkw:G4}
Let $E_v=\Q_\ell$, $r,N_0\ge1$, $e'=\operatorname{lcm}_{1\le a\le r}(\ell^a-1)$. For each inertial $s\in\Gamma$ let
$M_s$ be a multiset of at most $N_0$ elements of $\auQlb$, each satisfying $y^{\ell^a-1}=1$ for some $1\le a\le r$,
and let $\mathcal T$ be a set of potentially crystalline lattices such that for every $T\in\mathcal T$ and every
inertial $s$, all eigenvalues of $\beta_s^T$ lie in $M_s$. (Here an eigenvalue of $\beta_s^T$ is a
$y\in\auQlb$ such that $\beta_s-y$ has nonzero kernel on $\pkwB(T)\otimes_{\Z_\ell}\Z_\ell[y]$; since $\beta_s$ has
finite order $n$ (see the proof), each such $y$ lies in $\mu_n$, and $\pkwB(T)\otimes_{\Z_\ell}\Q_\ell(\mu_n)$ is
the direct sum of the kernels of $\beta_s-y$, $y\in\mu_n$.) Put $q_s=\prod_{y\in M_s}m_y$, $m_y$ the minimal
polynomial of $y$ over $\Q_\ell$. Then $q_s\in\Z_\ell[X]$ is monic of degree $\le N_0r$, depends only on $s$, and
$q_s$ and $X^{e'}-1$ kill $s$ on $\pkwD(T)$ for every $T\in\mathcal T$, and on $\pkwD(T')$ for every lattice $T'$
with an injective equivariant map into a finite direct sum of members of $\mathcal T$.
\end{proposition}

\begin{proof}
If $y^{\ell^a-1}=1$, $y$ lies in the unramified extension of degree $a$ and is integral, so $m_y\in\Z_\ell[X]$ is
monic of degree $\le a\le r$; hence $q_s$ is monic of degree $\le N_0r$. By Lemma~\ref{pkw:beta}(i), $\beta_{s'}=1$ for the
inertial $s'$ in some open subgroup $H$. As $\Gamma$ is profinite and $H$ is open, $s^n\in H\cap j^{-1}(\pkwI)$
for some $n\ge1$, so $\beta_s^n=\beta_{s^n}=1$ by Lemma~\ref{pkw:beta}(i). Since $X^n-1=\prod_{y\in\mu_n}(X-y)$
with pairwise distinct roots in $\Q_\ell(\mu_n)$, the idempotents $\prod_{y'\ne y}(\beta_s-y')/(y-y')$ decompose
$\pkwB(T)\otimes_{\Z_\ell}\Q_\ell(\mu_n)$ into the kernels of $\beta_s-y$, $y\in\mu_n$; the nonzero ones are the
eigenvalues, which lie in $M_s$ by hypothesis. Every $y\in M_s$ is a root of $q_s$ and of $X^{e'}-1$
($\ell^a-1\mid e'$), so $q_s(\beta_s)$ and $\beta_s^{e'}-1$ vanish on each eigenspace, hence on
$\pkwB(T)\otimes\Q_\ell(\mu_n)$ and on $\pkwB(T)$, which embeds in it. For $T'$ use
Lemma~\ref{pkw:beta}(ii) (and the evident compatibility of $\beta_s$ with finite direct sums).
\end{proof}

\begin{remark}
The definition of $\mathcal C(U,\mathrm{tame})$ before
{\cite[Proposition~5.2]{P55}} needs, for Proposition~\ref{pki:tame} and Theorem~\ref{pkd:data}(iv), a set $K(s)$ of
polynomials that depends on $s$ only. (G4) ({\cite[Proposition~6.6]{P55}}, with {\cite[Lemma~6.4]{P55}}) gives this, with the sizes read off from the
composition of $\mathcal C_v$ in the paper. The statements of {\cite[Proposition~2.28]{P55}} and {\cite[Proposition~4.5(b)]{P55}} say to which lattices the class is applied, and only this is taken from them here (their conclusions, which rest on {\cite[Theorem~5.8]{P55}}, are not used): at
$\ell=p_i$ the members of $\mathcal C_v$ are the restrictions to $D_v$ of the lattices $u_x$ in the representations
$\rho_x$ of {\cite[Proposition~4.6(c)]{P55}} and of the $t'+1$ summands of $u_\pi=T\oplus\bigoplus_{e<t'}\auO(-e)$; and
at a chosen place $w\nmid\ell$ {\cite[Proposition~4.6(c)]{P55}} says that inertia acts on $\rho_x$ ``through the roots of the
rescaled label times $(X-1)^{t'}$''. At the chosen place $w\mid\ell$ this is replaced by (G4) for $w\mid\ell$
({\cite[Proposition~6.6]{P55}}): $\rho_x$ is potentially crystalline at $w$, and the same data (the roots of the rescaled
label, and $t'$ trivial characters; the Hecke character $\eta$ is trivial on $\auO_{E,w}^\times$ by {\cite[Lemma~8.27]{P55}}, so $\eta_O$ is
crystalline at $w$ and contributes no inertial type) give the inertial type on $\pkwD$ (Definition~\ref{pkw:def-pcris}), whose eigenvalues on an
inertial $s$ are as follows. {\cite[Lemma~6.4]{P55}} gives them on an $s$ with $\theta_d(s)$ a generator of
$\mathbb F_{q^d}^\times$; since the restriction to inertia is semisimple and factors through $\theta_d$, whose
image is cyclic, it is a sum of the characters $\omega\circ\theta_d^{\,e_kq^b}$, so on an arbitrary inertial $s$
the eigenvalues are the $\omega(\theta_d(s)^{e_kq^b})$ attached to the roots of the label $c_w$ (multiplied by the
element of $\mathbb F_\ell^\times$ naming the tame part of the twist), together with $t'$ eigenvalues $1$:
they depend on $s$, $c_w$ and the twist only. Sizes: the label has degree $n_W$, which is $r$ for the pairs
$O_j(\vartheta)$ and $1$ for the pairs $O_{\nu,z}(\vartheta)$ ({\cite[\S2.1]{P55}}), so each irreducible
factor has degree $d_k\le n_W\le r$, and each eigenvalue lies in $\mathbb F_{\ell^{d_k}}^\times$ (the rescaling
factor lies in $\mathbb F_\ell^\times$) and is a root of unity of order dividing $\ell^{d_k}-1$; and
$|M_s|=n_W+t'=n_V$ is the rank of $\rho_x$. For the pairs $O_j(\vartheta)$, $n_V=r+t$; for these
Proposition~\ref{pkw:G4} applies with $N_0=r+t$, to the members $u_x|_{D_v}$ and to the summands of
$u_\pi|_{D_v}$ (which embed equivariantly), giving degree $\le(r+t)r$ as in the definition. For the pairs
$O_{\nu,z}(\vartheta)$ one has $n_V=2z$ with $m<z\le r$, and the bound obtained is $N_0r=2zr\le2r^2$. The
definition of $\mathcal C(U,\mathrm{tame})$ therefore uses the degree bound $\max(r+t,2r)\,r$, which covers both
kinds of pairs; as $t=5r$ it equals $(r+t)r=6r^2$, and Theorem~\ref{pkd:data}(iv) uses $N=\max(r+t,2r)r$ (the size
check there, $r^2(N+1)^2\le h$, is written for this $N$). With it, $K(s)=
\{X^{e'}-1,q_s\}$ in Theorem~\ref{pkd:data}(iv), $\mathcal C_v\subseteq\pkiC(U,K)$ holds as proved there, and
$L_c^*\in\pkiC(U',K^{-1})$ follows from Lemma~\ref{pkw:dual} (input (I6)) and Proposition~\ref{pkw:HT}(ii).
\end{remark}


\appendix

\section{Statements used from \texorpdfstring{\cite{P46}}{[P46]}, \texorpdfstring{\cite{P55}}{[P55]}}\label{C4:quoted}

For the convenience of the reader we reproduce here, without proofs, the statements of other papers that are used in this one, in the notation of the paper quoted. Each is headed by its place in that paper.

\begin{quoteddef}{\cite[Definition~13.2]{P46}}{}
Let $L\ge1$. A set $\mathsf T\subseteq\anPl_s$ is \emph{sparse for the primes
up to $L$} if for every prime $\ell\le L$ and all integers $0<i<j<l_0(s)$
there are no $\sigma,\tau,\rho\in\mathsf T$ and signs
$\epsilon,\epsilon'\in\{\pm1\}$ with
\[
 \ancl_{q'}(\tau/\sigma)=\epsilon\,\ell^{\,i},\qquad \ancl_{q'}(\rho/\sigma)=\epsilon'\,\ell^{\,j}\qquad\text{in }(\mathbb Z/q')^{\times},
\]
where $q'$ is the largest divisor of $q_1(s)$ prime to $\ell$.
\end{quoteddef}

\begin{quoteddef}{\cite[Definition~13.3]{P46}}{ (admissible data)}
Let $(q_0,d)$ be parameters. \emph{Admissible data} $D$ for $(q_0,d)$ consist
of the following objects, subject to the conditions \textup{(D1)--(D8)}.
\begin{enumerate}
\item[(D1)] Integers $1\le x\le x_1\le b$, the \emph{scales}.
\item[(D2)] They satisfy
\begin{gather*}
 10^{12}\,\anCP\le\log\log b,\qquad 100\log\log b\le\log x,\qquad 10^{6}\log x\le\log x_1,\\
 10^{6}\log x_1\le\log b,\qquad 6\log\log x_1\le\log\log b .
\end{gather*}
\item[(D3)] Integers $m$, $r$, $s$: $m$ is a power of $2$ with
$\sqrt x\le m\le2\sqrt x$; $r$ is a power of $3$ with $x^{5}\le r\le3x^{5}$;
$2^{s}=m^{22}$; and $s\le32\log x$. The field of the data is $E=E_s$, with
$q_1=q_1(s)$, $q_2=q_2(s)$, $f=[E:K]$, $\anPl=\anPl_s$.
\item[(D4)] $q_1x_1$ divides $b$.
\item[(D5)] A set $\anPb\subseteq\anPl$ (the \emph{$b$-places}) with
$\ancl_{q_0}(\sigma)=-1$ for all $\sigma\in\anPb$, with
$T(s)\le\#\anPb\cdot\log b$, and sparse for the primes up to $5r(r+1-m)$.
\item[(D6)] The weight $\kappa\colon\anPl\to\mathbb N$: $\kappa(\tau)=b$ for
$\tau\in\anPb$ and $\kappa(\tau)=xb/x_1$ otherwise.
\item[(D7)] Two distinct primes $p_0,p_1\equiv1\pmod{q_2^{2}}$ in the interval
$[b^{10},b^{11}]$ and an integer $\pi_M$ of $K$ with
$\pi_M\overline{\pi_M}=p_0p_1=:M$.
\item[(D8)] A prime $N\equiv1\pmod{q_2^{2}}$ with
\[
 \log N_-\le\log N\le\log N_-+\sqrt{\log b},\qquad
 \log N_-:=\frac mf\sum_{\tau\in\anPl}\log\kappa(\tau)+(m-1)\log\frac{M\,\anrd_E}{2\pi}-\frac{15}2\log b ,
\]
where $\anrd_E=|\operatorname{disc}E|^{1/[E:\mathbb Q]}$ is the root
discriminant, and an integer $\pi_N$ of $E$ with $\pi_N\pi_N^{*}=N$.
\end{enumerate}
We write $a=xb$ and $k=r+1-m$; the sets $\anPa$ and $\anPb$ are disjoint.
\end{quoteddef}

\begin{quoteddef}{\cite[Definition~13.3]{P46}}{ (admissible data)}
Let $(q_0,d)$ be parameters. \emph{Admissible data} $D$ for $(q_0,d)$ consist
of the following objects, subject to the conditions \textup{(D1)--(D8)}.
\begin{enumerate}
\item[(D1)] Integers $1\le x\le x_1\le b$, the \emph{scales}.
\item[(D2)] They satisfy
\begin{gather*}
 10^{12}\,\anCP\le\log\log b,\qquad 100\log\log b\le\log x,\qquad 10^{6}\log x\le\log x_1,\\
 10^{6}\log x_1\le\log b,\qquad 6\log\log x_1\le\log\log b .
\end{gather*}
\item[(D3)] Integers $m$, $r$, $s$: $m$ is a power of $2$ with
$\sqrt x\le m\le2\sqrt x$; $r$ is a power of $3$ with $x^{5}\le r\le3x^{5}$;
$2^{s}=m^{22}$; and $s\le32\log x$. The field of the data is $E=E_s$, with
$q_1=q_1(s)$, $q_2=q_2(s)$, $f=[E:K]$, $\anPl=\anPl_s$.
\item[(D4)] $q_1x_1$ divides $b$.
\item[(D5)] A set $\anPb\subseteq\anPl$ (the \emph{$b$-places}) with
$\ancl_{q_0}(\sigma)=-1$ for all $\sigma\in\anPb$, with
$T(s)\le\#\anPb\cdot\log b$, and sparse for the primes up to $5r(r+1-m)$.
\item[(D6)] The weight $\kappa\colon\anPl\to\mathbb N$: $\kappa(\tau)=b$ for
$\tau\in\anPb$ and $\kappa(\tau)=xb/x_1$ otherwise.
\item[(D7)] Two distinct primes $p_0,p_1\equiv1\pmod{q_2^{2}}$ in the interval
$[b^{10},b^{11}]$ and an integer $\pi_M$ of $K$ with
$\pi_M\overline{\pi_M}=p_0p_1=:M$.
\item[(D8)] A prime $N\equiv1\pmod{q_2^{2}}$ with
\[
 \log N_-\le\log N\le\log N_-+\sqrt{\log b},\qquad
 \log N_-:=\frac mf\sum_{\tau\in\anPl}\log\kappa(\tau)+(m-1)\log\frac{M\,\anrd_E}{2\pi}-\frac{15}2\log b ,
\]
where $\anrd_E=|\operatorname{disc}E|^{1/[E:\mathbb Q]}$ is the root
discriminant, and an integer $\pi_N$ of $E$ with $\pi_N\pi_N^{*}=N$.
\end{enumerate}
We write $a=xb$ and $k=r+1-m$; the sets $\anPa$ and $\anPb$ are disjoint.
\end{quoteddef}

\begin{quotedthm}{\cite[Lemma~13.4]{P46}}{}
For admissible data the following hold.
\begin{enumerate}
\item[(S1)] $x^{10^{6}}\le x_1$, $x_1^{10^{6}}\le b$,
$4\log x_1\le(\log x_1)^{6}\le\log b$,
$(\log b)^{100}\le x$, $\log b\ge128$, and
$\log x\ge10^{14}(q_0+2d+|\aneu|)$.
\item[(S2)] $x^{11}\le2^{s}\le2^{22}x^{11}$; $l_0(s)\le2592\log x$;
$2m\le r$, so $r/2\le k\le r$.
\item[(S3)] $x^{11}\le T(s)=\#\anPa\le f\le q_1\le x^{12}$, and
$f=\varphi(q_0)\,\#\anPa$; moreover $1\le\#\anPa/\#\anPb\le\log b$.
\item[(S4)] $\anrd_E=\sqrt d\;q_1\prod_{p\mid q_0}p^{-1/(p-1)}$, and
$x^{10}\le\anrd_E\le x^{12}$.
\item[(S5)] $1\le xb/x_1\le\kappa(\tau)\le b\le a\le b^{2}$ for all $\tau$;
$b^{20}\le M\le b^{22}$;
$20\log b\le\log(M\anrd_E/2\pi)\le23\log b$; and
$m\log b\le\log N\le25\,m\log b$.
\item[(S6)] With $\Omega:=e^{\aneu}\cdot2\pi/\anrd_E$: $0<\Omega\le1$ and
$20\log b\le\log(M/\Omega)\le23\log b$.
\item[(S7)] With $n=m/2$ and $\varrho=n-1$ (so that $r=k+2\varrho+1$): $r$ is odd; every prime
dividing $q_1$ divides $q_0$; $q_0^{\,q_0}\le x$; $\varrho\ge10^{5}$;
$k\ge512000\,n^{9}$ and $k\ge1000\,(2\varrho+2)^{9}$; $k\ge2\varrho^{2}+3$; $18\varrho f\le b^{8}$;
$f\le b/1000$; $\anrd_E\le b$; and $N\le b^{60n}$. Moreover, with $\Delta=2r^{8}$:
\begin{itemize}
\item[(a)] $\varrho\ge10^{8}$ (so $n\ge100$); $n^{7}\le k$; $12n^{3}+500\le k$; $k-3n-6>10^{9}n^{3}$;
\item[(b)] $r\ge100$ (so $2r^{8}\ge500r^{3}+10r^{5}$); $3r^{3}f\le r^{7}$; $8nr\le b$; $kf\ge40$; $f\ge4000\,n$;
\item[(c)] $2f\le2^{2\varrho}$ and $2^{-\varrho}\le b^{-6}$;
\item[(d)] $(200\,f\,k\,\anrd_E)^{1000}\le b$ and $(25\,n^{3}f\,\anrd_E)^{1000}\le b$;
\item[(e)] $\sqrt d\le\anrd_E$, and $n,r,f,\sqrt d\le b$; $M\le b^{23}$;
\item[(f)] $\binom{a+k-1}{k-1}^{f}\le b^{2kf}$ (this bounds the dimension $d_\infty=\dim W_{\lambda^X}$ of
{\cite[\S14.3]{P46}}, a product of at most $f$ such binomial coefficients);
\item[(g)] $m\Delta+r\le\kappa(\tau)$ for all $\tau$; $k+2b+2m\Delta\le b^{2}$ (this bounds every
entry of $2\varsigma_\tau$, so that $L_\tau:=\max_i|2\varsigma_{\tau,i}|\le b^{2}$); and
$(r+1)^{4}(b^{2}+r+1+2n^{2})\le b^{5}$.
\end{itemize}
\end{enumerate}
\end{quotedthm}

\begin{quotedthm}{\cite[Lemma~2.5]{P55}}{}
Assume that $\pgcodd$ can be taken with $\log\pgcodd\le x^{800}$ ({\cite[Proposition~2.27]{P55}}). Then
for every finite set $T$ of primes
\[
\sum_{\ell\in T}e'_\ell\log\ell\le x_1,\qquad
\sum_{\ell\in T}\bigl(\beta_\ell+\pgv(h_1)\bigr)\log\ell\le\aubud=\frac{f\,a\,(\log b)^3}{m\,k}.
\]
\end{quotedthm}

\begin{quoteddef}{\cite[Definition~2.6]{P55}}{ (Main summands)}
Let $n\ge1$ and $t'\ge1$. A continuous representation $R\colon\pgG\to\auGL_n(\auQlb)$ is a
\emph{main summand of rank $n$ and string length $t'$} if:
\begin{enumerate}
\item[(M1)] $R$ is semisimple and $\pgtr R(g^c)=\chi(g)^{t'-1}\,\pgtr R(g^{-1})$ for all $g$;
\item[(M2)] $R$ is unramified outside $\ell p_0p_1$; at $v\mid\ell$ it is crystalline if
$v\nmid p_0p_1$ and potentially crystalline otherwise; its Hodge--Tate weights $J_\sigma$, for
the embeddings $\sigma\colon E\to\C$ (transported by $\iota$), satisfy:
$1\le n\le r$, $t'\le t$; for one $j\le m$ every $u\in J_\sigma$ is \emph{low},
$\Delta/4\le|u|\le2m\Delta$, or \emph{high}, $|u\pm\kappa_j(\sigma)|\le t$ with the sign $+$ for
$\sigma$ in the CM type of the data and $-$ otherwise; at most $nm/k$ elements of $J_\sigma$ have
$|u|\le2h^2$; and $J_\sigma$ and $J_{\aubar\sigma}$ together have at most $n$ negative elements;
\item[(M3)] at every $v\mid p_0p_1$, $v\nmid\ell$: there is $N_v$ prime to $v$ with $R(s^{N_v})=1$
for $s\in I_v$, and an element $\gamma\in I_v$, equal to $\gamma_v$ if $v$ is chosen, with
$\pgv(\det(R(\gamma)-1))\le n\,\pgv(\pgNch)$ and $\pgv(\det(R(\gamma^{-c})-1))\le n\,\pgv(\pgNch)$;
\item[(M4)] for all large $\auO$ and every $\pgG$-stable $\auO$-lattice, with matrices $u(g)$: the
$u(g)$ span $M_n(\auO)$ over $\auO$, and there are $g\in\pgG$ and $i\in\{0,1\}$ with
$\pgcyc(g)=p_i^{\,n}$ and $u(g)$ a scalar;
\item[(M5)] there are a place $v\mid p_0p_1$, $v\nmid\ell$, and $s\in I_v$ such that $R(s)$ is a
scalar $y$ with $y^{\ell^aj}-1$ a unit for all $a\ge0$ and $1\le j\le4n+2$.
\end{enumerate}
\end{quoteddef}

\begin{quotedthm}{\cite[Proposition~2.28]{P55}}{ (The local loss)}
Let $v\mid\ell$, let $\auO$ be large, and let $\mathcal C_v$ be a class of lattices in potentially
crystalline representations of $D_v$ of the following kind: crystalline with Hodge--Tate weights in
prescribed sets, or, if $\ell=p_i$, potentially crystalline with tame inertial type of the shape
allowed by (G4). Let $L$ be a lattice of $\mathcal C_v$ that is the restriction of a lattice in a
main summand twisted as in {\cite[Theorem~2.10]{P55}}, or of a character $\chi^e$, $e<t'$, where all
weights of all members of $\mathcal C_v$ are weights of these summands. Then for every finite family $u_x$ of global lattices whose restrictions to $D_v$ lie in $\mathcal C_v$, the module of cocycles $\pgG\to L/\ell^M$ obeying all $\auO$-linear relations among the families $(u_x(g)-1)_x$ has, modulo those obeying the crystalline relations of $L$ at $v$, length at most the \emph{defect} $\delta_v(L)$ (the length of the crystalline cocycles of the Tate dual of $L$ modulo those cup-orthogonal to $\mathcal C_v$; see {\cite[Proposition~4.5]{P55}}),
and
\[
\sum_{v\mid\ell}\ \sum_{L}\delta_v(L)\le r'\cdot\pgloc_\ell ,
\]
the inner sum over the $t'+1$ summands. In particular there is no loss if $\ell>2b$ and $\ell\nmid M$ ({\cite[Theorem~5.8(d)]{P55}}; for $\ell\ge4b+2$, $\ell\nmid q_1dM$ see also {\cite[Remark~5.7]{P55}} for another route).
\end{quotedthm}

\begin{quotedthm}{\cite[Proposition~4.5]{P55}}{ (The local loss, form used)}
Let $\auO$ be a large coefficient ring, $u_x\colon\pgG\to\auGL_n(\auO)$ ($x\in Y$, finite)
continuous, $\rho\colon\pgG\to\auGL_{d}(\auO)$, $M\ge1$. Let $Z^{\mathrm R}_M(\rho)$ be the module of
cocycles $z\colon\pgG\to(\auO/\ell^M)^d$ for $\rho$ satisfying $\sum_ia_iz(g_i)=0$ whenever
$a_i\in\auO$ and $\sum_ia_i(u_x(g_i)-1)=0$ for all $x$ (\emph{Ribet cocycles}), and
$Z^{\mathrm{cr}}_M(\rho)$ the module of functions obeying the crystalline relations of $\rho$
({\cite[Definition~2.7]{P55}}). Put $\prics_M(\rho)=\prirl\bigl(Z^{\mathrm R}_M(\rho),
Z^{\mathrm{cr}}_M(\rho)\bigr)$.
\begin{enumerate}
\item[(a)] If, at each $v\mid\ell$, the restrictions to $D_v$ of all $u_x$ and of $\rho$ lie in a
class $\mathcal C_v$ of local lattices (as in {\cite[Proposition~2.28]{P55}}), then
$\prics_M(\rho)\le\sum_{v\mid\ell}\delta_v(\rho)$, where the \emph{defect} $\delta_v(\rho)$ is the
length of the crystalline cocycles of the Tate dual of $\rho|_{D_v}$ modulo those whose cup product
with every crystalline cocycle of every member of $\mathcal C_v$, under every pairing to
$\Z/\ell^m(1)$, is a coboundary.
\item[(b)] Let $c$ be in lift setting, $R$ a main summand of rank $r'$ and string length $t'$,
and $\rho_x$ ($x\in Y$) representations of rank $n$ satisfying the conclusion (c) of
{\cite[Proposition~4.6]{P55}} for $c$. For all large $\auO$, all lattices $u_x$ in $\rho_x$, all
lattices $T_0$ in $R$ and $T$ equal to $T_0$ on the $D_v$, $v\mid\ell$, and one $\pi\in Y$ with
$u_\pi=T\oplus\bigoplus_{e<t'}\auO(-e)$, there are classes $\mathcal C_v$ as in (a) containing all
$u_x|_{D_v}$ and the restrictions of the $t'+1$ summands $L$ of $u_\pi$, with
\[
\sum_{L}\sum_{v\mid\ell}\delta_v(L)/e_\auO\le r'\,\pgloc_\ell .
\]
\end{enumerate}
\end{quotedthm}

\begin{quotedthm}{\cite[Proposition~4.6]{P55}}{ (Base level of a label)}
Let $c$ be in lift setting for an auxiliary $\vartheta$. There are Hodge--Tate data $J'$ and a
space of forms of ``base level'' (the space $\mathcal F(O)$ of {\cite[Definition~8.42]{P55}}: level
$J_O$ and character $\chi_O$; these forms are of principal level at the chosen places) with: (a) $e_c\mathcal F(O)$ consists of forms of base level; (b) if
$y$ is integral at $\iota$, so is $N\,e_cy$; (c) for every finite set $U$ of places off
$\pgS(\ell)\cup Q\cup\overline Q$ and every system $x$ of Hecke eigenvalues occurring in base level
with principal level at $U$, the normalised Galois representation $\rho_x=S_x\otimes\eta_O$ is
conjugate self-dual with string length $t'$, has Hodge--Tate data $J'$, is unramified outside $\ell
p_0p_1$, $Q\cup\overline Q$, $U\cup\overline U$, inertia at a chosen $w\nmid\ell$ acts through the roots of
the rescaled label times $(X-1)^{t'}$ and semisimply (the monodromy of $\operatorname{rec}(\Pi_w)$
is $0$), so that $o_w(\rho_x(s))(\rho_x(s)-1)=0$ for $s\in I_w$, $o_w$ the monic polynomial whose
roots are the (images of the) roots of the rescaled label, and $(\rho_x(s)-\psi(s))(\rho_x(s)-1)=0$
for $s$ in inertia above $Q\cup\overline Q$.
\end{quotedthm}

\begin{quoteddef}{\cite[Definition~5.1]{P55}}{}
Let $T$ be a lattice and $\pllB$ Fontaine's crystalline period ring.
\begin{enumerate}
\item The \emph{crystalline cocycles} of $T$ are the continuous maps $c\colon\Gamma\to T$ for which some
$y\in\pllB\otimes_{\Z_\ell}T$ satisfies $1\otimes c(g)=gy-y$ for all $g\in\Gamma$; they form an
$\auO$-module $\pllZf(T)$, which contains the coboundaries $g\mapsto gx-x$, $x\in T$. (These are
the continuous cocycles whose class lies in $H^1_f(\Gamma,T)$, the preimage of the Bloch--Kato
subspace $H^1_f(\Gamma,V)=\ker(H^1(\Gamma,V)\to H^1(\Gamma,\pllB\otimes V))$.)
\item A map $s\colon\Gamma\to T/\ell^M$ \emph{obeys the crystalline relations} of $T$ if every
finite $\auO$-linear relation $\sum_g\lambda_g\,c(g)=0$ ($\lambda_g\in\operatorname{End}_\auO T$)
that holds for all $c\in\pllZf(T)$ also holds for $s$. The reductions modulo $\ell^M$ of the
elements of $\pllZf(T)$ obey them; we write $\overline{\pllZf}(T)_M$ for the module of these
reductions.
\item Let $\mathcal C$ be a class of lattices. $s\colon\Gamma\to T/\ell^M$ \emph{comes from
$\mathcal C$} if $s=\mathsf t\circ a$ with $Y\in\mathcal C$, $a\in\pllZf(Y)$ and $\mathsf t\colon
Y\to T/\ell^M$ $\auO$-linear and $\Gamma$-equivariant. For a submodule $N$ of maps
$\Gamma\to T/\ell^M$ all of whose elements come from $\mathcal C$, its \emph{defect} is the length
of $N/(N\cap\{\text{maps obeying the crystalline relations of }T\})$.
\item A \emph{pairing modulo $\ell^m$} of lattices $Y$, $L$ is a pair $(\psi,c_\psi)$ with
$\psi\in(Y\otimes_\auO L)^*$ and $c_\psi\in\pllZf((Y\otimes L)^*)$ such that $g\psi-\psi=
\ell^mc_\psi(g)$ for all $g$; with a fixed $\Z_\ell$-linear $\mathrm{f}\colon\auO\to\Z_\ell$ for
which $(x,y)\mapsto\mathrm f(xy)$ is a perfect pairing, it gives the $\Gamma$-equivariant pairing
$b_\psi(y,l)=\mathrm f(\psi(y\otimes l))\bmod\ell^m$ with values in $\Z/\ell^m(1)$.
\item The \emph{cup-orthogonal} $L^{\pllperp\mathcal C}$ of $\mathcal C$ in $L$ is the module of
maps $c\colon\Gamma\to L$ such that for every $Y\in\mathcal C$, every $m$, every pairing
$(\psi,c_\psi)$ modulo $\ell^m$ of $Y$ and $L$ and every $a\in\pllZf(Y)$, the $2$-cochain
$(g,h)\mapsto b_\psi(a(g),g\,c(h))$ is a coboundary in $C^2(\Gamma,\Z/\ell^m(1))$. The \emph{dual
defect} is
\[
\plldef_{\mathcal C}(L)=\pglen\bigl(\pllZf(L)/(\pllZf(L)\cap L^{\pllperp\mathcal C})\bigr),
\]
and $\mathcal C$ has \emph{cup exponent} $\kappa$ on $L$ if $\ell^\kappa\pllZf(L)\subseteq L^{\pllperp
\mathcal C}$.
\end{enumerate}
\end{quoteddef}

\begin{quotedthm}{\cite[Proposition~5.2]{P55}}{ (Exact form of {\cite[Proposition~2.28]{P55}})}
In the situation of {\cite[Proposition~2.28]{P55}}, with $\mathcal C_v$ as above and $\auO$ large, let
$N_v(L_c)$ be the module of maps $\Gamma\to L_c/\ell^M$ that come from $\mathcal C_v$. Then its
defect $\delta_v(L_c)$ satisfies (here $N_v(L_c)$ is indeed a module: the classes
$\mathcal C(U)$ and $\mathcal C(U,\mathrm{tame})$ are stable under finite direct sums, since the
jump condition and the conditions on $X^{e'}-1$ and $q_s$ are checked summand by summand, and
$\mathsf t_1\circ a_1+\mathsf t_2\circ a_2=(\mathsf t_1,\mathsf t_2)\circ(a_1,a_2)$ with
$(a_1,a_2)\in\pllZf(Y_1\oplus Y_2)$)
\begin{equation*}\tag*{{\cite[(9)]{P55}}}
\delta_v(L_c)\le\plldef_{\mathcal C_v}(L_c^*)\le[E_v:\Q_\ell]\cdot\frac{\pllrk_\auO L_c}{x_c}
\cdot\pllD_\ell,\qquad
\pllD_\ell=10^6\,\frac mk\bigl(b\,\pgv(M)+\lambda_\ell\bigr),
\end{equation*}
with $x_{\mathrm{main}}=1$, $x_e=t$ for the characters, and $\lambda_\ell$ as in
{\cite[\S2.1]{P55}}. Consequently $\sum_{v\mid\ell}\sum_c\delta_v(L_c)\le r'\pgloc_\ell$.
\end{quotedthm}

\begin{quotedthm}{\cite[Proposition~5.4]{P55}}{}
Let $\mathcal C$ be a class of lattices containing $T$, all of whose members are potentially
crystalline, and let $N$ be a module of maps $\Gamma\to T/\ell^M$ that come from $\mathcal C$. Then
\[
\pglen\bigl(N/(N\cap\overline{\pllZf}(T)_M)\bigr)\le\plldef_{\mathcal C}(T^*).
\]
In particular the defect of $N$ is at most $\plldef_{\mathcal C}(T^*)$.
\end{quotedthm}

\begin{quotedthm}{\cite[Lemma~5.5]{P55}}{}
Let $\mathcal C$ be a class of lattices over $\Gamma$, stable under finite direct sums and under
passage to $\Gamma$-stable sublattices, all of whose members become crystalline over a fixed finite
Galois extension $E'/E_v$. Let $u_x$ ($x$ in a finite set $X$) be members of $\mathcal C$, of ranks
$n_x$, and let $z\colon\Gamma\to T/\ell^M$ be a continuous cocycle such that $\sum_i\lambda_i\,
z(g_i)=0$ whenever $\lambda_i\in\auO$, $g_i\in\Gamma$ and $\sum_i\lambda_i(u_x(g_i)-1)=0$ for every
$x$. Then $z$ comes from $\mathcal C$. The same holds if $\Gamma$ is replaced by a group $G$ in which
$\Gamma$ is a subgroup, the $u_x$ and $z$ are defined on $G$ and only the restrictions to $\Gamma$
are asserted to lie in $\mathcal C$ and to come from $\mathcal C$.
\end{quotedthm}

\begin{quotedthm}{\cite[Lemma~5.6]{P55}}{}
Let $L$ be a potentially crystalline lattice. Then $\pllZf(L)$ is a free $\auO$-module of rank at
most $2[E_v:\Q_\ell]\pllrk_\auO L$. Consequently, if $\mathcal C$ has cup exponent $\kappa$ on $L$, then
$\plldef_{\mathcal C}(L)\le2\kappa\,[E_v:\Q_\ell]\,\pllrk_\auO L$.
\end{quotedthm}

\begin{quotedthm}{\cite[Theorem~5.8]{P55}}{}
Let $\ell$ be any prime, $v\mid\ell$, $\auO$ large, $\mathcal C_v$ and the summands $L_c$ as in
{\cite[Proposition~5.2]{P55}}, with $L_c\in\mathcal C_v$. Then $\mathcal C_v$ has, on $L_c^*$, a cup
exponent $k_c$ with $2k_c\le\pllD_\ell/x_c$; or, in the sub-case below where no exponent is
claimed, directly $\plldef_{\mathcal C_v}(L_c^*)\le[E_v:\Q_\ell]\pllrk(L_c)\pllD_\ell/x_c$. In
detail, with $h=r^8$, $Q=kt$ and $\lambda=\lambda_\ell$:
\begin{enumerate}
\item[(a)] (tame, $\ell=p_i$) exponent $3r^2(N+1)^2+150(h^2+4t+2)$, $N=\max(r+t,2r)r$;
\item[(b)] (crystalline, $\ell\le Q$) exponent $10^4\lfloor2\lambda/Q\rfloor(\lfloor\log_\ell
h^3\rfloor+1)$;
\item[(c)] (crystalline, $\ell>Q$, $\lambda\ge h^3$) exponent $10^4\lfloor2\lambda/Q\rfloor m$;
\item[(d)] (crystalline, $\ell>Q$, $\lambda<h^3$) $\plldef=0$ if $\lambda<\lfloor h/5\rfloor$ or
$L_c$ is a character, and in all cases $\plldef_{\mathcal C_v}(L_c^*)\le
[E_v:\Q_\ell]\cdot10^5\bigl(\pllrk(L_c)\lfloor2\lambda/Q\rfloor m+\mu\lambda\bigr)$ with
$\mu=\lfloor\pllrk(L_c)m/(kx_c)\rfloor$.
\end{enumerate}
\end{quotedthm}

\begin{quotedthm}{\cite[Lemma~5.10]{P55}}{}
If {\cite[(9)]{P55}} holds for all $v\mid\ell$ and all $c$, then $\sum_{v\mid\ell}\sum_c
\delta_v(L_c)\le r'\pgloc_\ell$.
\end{quotedthm}

\begin{quotedthm}{\cite[Lemma~6.4]{P55}}{ (Depth zero)}
Let $\Pi$ be an irreducible smooth representation of $\auGL_n(K)$ whose vectors fixed by the
principal congruence subgroup $K(1)=1+\varpi M_n(\auO_K)$ contain $\sigma_c$, inflated from
$\auGL_n(\auO_K)\to\auGL_n(\auF_q)$, for a label $c$. Then $\Pi\simeq\pi_1\times\dots\times\pi_s$
(irreducible normalised parabolic induction), where $\pi_k$ is a depth-zero supercuspidal
representation of $\auGL_{d_k}(K)$ containing $\sigma_{f_k}$; $\pgarec(\Pi)$ has monodromy $0$;
and $\pgarec(\Pi)|_{I_K}$ is semisimple, factors through the tame quotient, and, on an element
$s\in I_K$ with $\theta_d(s)$ a generator of $\auF_{q^d}^\times$, $d=\operatorname{lcm}(d_k)$, it
has the eigenvalues $\omega(\xi_k^{\,q^a})^{\epsilon}$, $0\le a<d_k$, where the roots $\xi_k$ are
normalised by $\xi_k=\theta_{d}(s)^{e_k}$ through the naming of $\sigma_c$, for a sign $\epsilon\in\{\pm1\}$ depending only on the normalisations of
$\pgarec$, $\mathrm{Art}$ and of Green's parametrisation, not on $c$ or $\Pi$.
\end{quotedthm}

\begin{quotedthm}{\cite[Proposition~6.6]{P55}}{ ((G4), (G5))}
Let $y$, $S_y$ be as in {\cite[Proposition~6.2]{P55}} and $w$ a split
place, $w\mid p$, $v$ the place of $E^+$ below.
\begin{enumerate}
\item[(G4)] If $y$ is fixed by the principal congruence subgroup at $w$ and lies in the
$\sigma_c$-isotypic part for a label $c$, then $S_y|_{I_w}$ is tame and, if $w\nmid\ell$, of finite
order, semisimple, with the eigenvalues $\omega(\xi^{\epsilon})$, $\xi$ running over the roots of
$c$, on a generator $s$ of tame inertia with $\theta_n(s)$ a generator; if $w\mid\ell$, $S_y$ is
potentially crystalline at $w$, crystalline over $E_w(\varpi^{1/(q^d-1)})$, $d=\operatorname{lcm}_k(d_k)$,
with this inertial type. ({\cite[Lemma~6.4]{P55}} is phrased with $\theta_d(s)$; since the norm
$\auF_{q^n}^\times\to\auF_{q^d}^\times$ is surjective, $\theta_n(s)$ a generator implies
$\theta_d(s)$ a generator.)
\item[(G5)] If $w\nmid\ell$ and $y$ is an eigenvector of the parahoric $\mathcal P$ of type
$(n_W,t)$ with character $\psi_{\mathcal P}$, $\psi=\psi_w\ne1$, then $(S_y(s)-\psi_w(s)^{
\epsilon})(S_y(s)-1)=0$ for $s\in I_w$, $\psi_w$ being viewed on $I_w$ through $\mathrm{Art}_w$.
\end{enumerate}
The sign $\epsilon$ is the same in (G4) and (G5) and for all $w$.
\end{quotedthm}

\begin{quotedthm}{\cite[Lemma~8.27]{P55}}{ ($\eta_O$ is unramified away from $\ell$)}
Let $\eta$ be the algebraic Hecke character $y\mapsto(\chi_W(y/\bar y)\allowbreak \upsilon(y)^{-r'}\allowbreak |y|^{-r'/2})^{-1}$
of {\cite[Proposition~6.11]{P55}}, whose $\ell$-adic character is $\eta_O$. Then $\eta$ is trivial on
$\auO_{E,w}^\times$ for every finite place $w$ of $E$. Hence $\eta_O$ is unramified at every finite
place $w\nmid\ell$; in particular at the places of $Q\cup\overline Q$, at the places of $U\cup
\overline U$ in {\cite[Proposition~4.6(c)]{P55}}, and at the chosen places not above $\ell$.
\end{quotedthm}

\begin{thebibliography}{P46}
\bibitem[P46]{P46} D.~A.~Trejo Pizzo, \emph{The sign of the radial intersection form on the idele class group}, preprint, 2026, \url{https://github.com/dtrejopizzo/riemann/tree/main/04-papers/46-row-d-sign-radial-intersection-form}.
\bibitem[P55]{P55} D.~A.~Trejo Pizzo, \emph{The Covolume Theorem: regular families of twists, Selmer groups and Ribet's method}, preprint, 2026, \url{https://github.com/dtrejopizzo/riemann/tree/main/04-papers/55-covolume-theorem-regular-families-of-twists}.
\bibitem{Washington} L.~C.~Washington, \emph{Introduction to Cyclotomic Fields}, 2nd ed.,
Graduate Texts in Mathematics \textbf{83}, Springer, New York, 1997.
\bibitem{SerreLF} J.-P.~Serre, \emph{Local Fields}, Graduate Texts in Mathematics 67, Springer,
New York, 1979.
\bibitem{BlochKato} S.~Bloch and K.~Kato, \emph{$L$-functions and Tamagawa numbers of motives},
in: The Grothendieck Festschrift, Vol.~I, Progr.\ Math.\ \textbf{86}, Birkh\"auser, 1990,
333--400.
\bibitem{FontaineLaffaille} J.-M.~Fontaine and G.~Laffaille, \emph{Construction de
repr\'esentations $p$-adiques}, Ann.\ Sci.\ \'Ecole Norm.\ Sup.\ (4) \textbf{15} (1982), 547--608.
\bibitem{pkfColmez98} P.~Colmez, \emph{Th\'eorie d'Iwasawa des repr\'esentations de de Rham d'un corps local}, Ann. of Math. (2) \textbf{148} (1998), 485--571.
\bibitem{pkfFontaineII} J.-M.~Fontaine, \emph{Le corps des p\'eriodes $p$-adiques}, expos\'e~II in: P\'eriodes $p$-adiques, Ast\'erisque \textbf{223} (1994), 59--101 (numbers cited are those of the expos\'e).
\bibitem{pkfFontaineIII} J.-M.~Fontaine, \emph{Repr\'esentations $p$-adiques semi-stables}, expos\'e~III in: P\'eriodes $p$-adiques, Ast\'erisque \textbf{223} (1994), 113--184 (numbers cited are those of the expos\'e).
\bibitem{pkfWintenberger83} J.-P.~Wintenberger, \emph{Le corps des normes de certaines extensions infinies de corps locaux; applications}, Ann. Sci. \'Ecole Norm. Sup. (4) \textbf{16} (1983), 59--89.
\bibitem{pkfScholze12} P.~Scholze, \emph{Perfectoid spaces}, Publ. Math. IH\'ES \textbf{116} (2012), 245--313.
\bibitem{pkfAx70} J.~Ax, \emph{Zeros of polynomials over local fields---the Galois action},
J.\ Algebra \textbf{15} (1970), 417--428, DOI 10.1016/0021-8693(70)90069-4.

\bibitem{pklTate67} J.~Tate, \emph{$p$-divisible groups}, in: Proc.\ Conf.\ Local Fields (Driebergen, 1966), Springer, 1967, 158--183.








\end{thebibliography}

\end{document}
